% A_closing.tex -- the final section of the online appendix: the models estimated on the closing half hour alone
% (decision D3). Created 2026-10-26 (integration pass for JFEc; APPENDIX agent). \input at the end of
% rehaul/10_statistical_tests.tex, so it is the last appendix section. Bullets only (author, 2026-10-26).
% It replaces the appendix section "Closing-interval models" of rehaul/10_appendix.tex, built on the master-table
% closing-interval model set under the old data rules (W = 2,000 trading days, its own predictor set and tuning; audit
% numbers-10), and its subsection "Training on the last 30 minutes" (sec:rh_econ_last30): both are parked, whole, at the
% end of 10_appendix.tex. Every label of that section that live text cites is kept here: app:rh_closing,
% app:rh_closing_sample, app:rh_closing_predictors, app:rh_closing_bench, sec:rh_econ_last30, sec:rh_results_pooled.
% sec:rh_close_scale (its standardization display) is not kept: no live text cites it, and Appendix app:rh_dp_scale
% (rehaul/A_data_predictors.tex) states the scaler. Table tab:rh_oos_last30 (the old model set) is parked with the
% section; its label is aliased below to Table tab:rh_train_arms, so a stale \ref still resolves.
% Sources: results/ci_har_rebuild_2026-10-26/RULE.txt and measure/ (the rebuilt closing-interval HAR benchmark; the
% NUMBERS agent switches the generators to it); drafts/2026-10-26_train_arms.tex (the CLOSE sample; \rhta...,
% writeup/make_rehaul_train_arms.py, whose claims() asserts the class names, directions and "all / every / none");
% drafts/2026-10-25_close_accuracy.tex, first item (\rhca..., writeup/make_rehaul_close_accuracy.py);
% drafts/2026-10-25_mcs.tex, Parts B and C (\rhmcs..., writeup/make_rehaul_mcs.py, draft_claims()); the CI-HAR
% strategy macros of writeup/make_rehaul_econ_main.py (\rhem...CiHar, released by ECON for this section).
% Numbers: macros only. The \rhci... macros of the rebuild are not generated yet: they print \rhPending until NUMBERS
% defines them (values in measure/accuracy.csv, dm_old_vs_new.csv and summary.json: first forecast 14 June 2006;
% closing-half-hour QLIKE on T_close 0.1033 (old) and 0.1025 (rebuilt), DM -0.86 (p = 0.39); the sign strategy changes
% position on 18 of the 614 post-warmup days).
% Evaluation spans as in Methods (05_methods.tex): T_oos, Table 1's span; T_close, its 2020-2024 subset.
% FOMC fixes (D13): the CLOSE sample ran on the design before them; marked in the text (DATA relay, app:rh_dp_fomc).
\ifdefined\rhtaNdays\else% AUTO-GENERATED by writeup/make_rehaul_train_arms.py -- do not edit; re-run the script.
% \rhta... macros: the pre-registered comparison of training samples (tab:rh_train_arms; draft
% writeup/drafts/2026-10-26_train_arms.tex).  Names: <Stat><Model><Arm><Bar>; Arm Rth (RTH13, the main design),
% All (ALL48), Cl (CLOSE), ClNat (CLOSE_NAT), AllCt / RthCt (penalties chosen at 16:00), AllThreeK / RthThreeK /
% AllThirtyNineK (the last 3,000 / 39,000 rows); Bar Close (16:00) or Bar (every bar); X... = first arm vs second
% (full set); Sh... = Sharpe ratio (D19: the paper's pooled ratio; p by Ledoit-Wolf); Ws... = the close-only
% window sweep.
% Sources (git show, pinned): insurance-selling-2026-09-19-e34-train3way 0106ccb; insurance-selling-2026-09-19-e34-train3way-trees 3501ab7; insurance-selling-2026-09-19-e34-closeonly-wsweep f4abb69.
\providecommand{\rhtaNdays}{}\renewcommand{\rhtaNdays}{2{,}579}  % trading days scored, 2 January 2014 .. 30 April 2024 (early closes dropped)
\providecommand{\rhtaNbars}{}\renewcommand{\rhtaNbars}{33{,}527}  % regular-hours bars scored (every bar; RV > 0)
\providecommand{\rhtaStart}{}\renewcommand{\rhtaStart}{2 January 2014}  % first evaluation day
\providecommand{\rhtaEnd}{}\renewcommand{\rhtaEnd}{30 April 2024}  % last evaluation day
\providecommand{\rhtaHacLag}{}\renewcommand{\rhtaHacLag}{8}  % Newey-West lag of the daily Diebold-Mariano tests
\providecommand{\rhtaNIntervals}{}\renewcommand{\rhtaNIntervals}{48}  % half-hour intervals of the futures' day (ALL48)
\providecommand{\rhtaNBarsDay}{}\renewcommand{\rhtaNBarsDay}{13}  % regular-hours bars of a day, bar ends 10:00 .. 16:00 ET (RTH13)
\providecommand{\rhtaWinDays}{}\renewcommand{\rhtaWinDays}{3{,}000}  % rolling window of the equal-calendar arms, trading days
\providecommand{\rhtaRowsAll}{}\renewcommand{\rhtaRowsAll}{144{,}000}  % training rows of ALL48 (3,000 days x 48 intervals)
\providecommand{\rhtaRowsRth}{}\renewcommand{\rhtaRowsRth}{39{,}000}  % training rows of RTH13 (3,000 days x 13 bars)
\providecommand{\rhtaRowsCl}{}\renewcommand{\rhtaRowsCl}{3{,}000}  % training rows of CLOSE (3,000 closes)
\providecommand{\rhtaRowsSmall}{}\renewcommand{\rhtaRowsSmall}{3{,}000}  % training rows of the small equal-row arms (ALL48_N3k, RTH13_N3k; CLOSE)
\providecommand{\rhtaRowsLarge}{}\renewcommand{\rhtaRowsLarge}{39{,}000}  % training rows of the large equal-row arm ALL48_N39k (= RTH13's)
\providecommand{\rhtaSpanAllSmall}{}\renewcommand{\rhtaSpanAllSmall}{62.5}  % days of 48 intervals spanned by 3,000 rows (about)
\providecommand{\rhtaSpanRthSmall}{}\renewcommand{\rhtaSpanRthSmall}{231}  % sessions of 13 bars spanned by 3,000 rows (about)
\providecommand{\rhtaSpanClSmall}{}\renewcommand{\rhtaSpanClSmall}{12}  % years of closes spanned by 3,000 closes (about)
\providecommand{\rhtaSpanAllLarge}{}\renewcommand{\rhtaSpanAllLarge}{812.5}  % days of 48 intervals spanned by 39,000 rows (about)
\providecommand{\rhtaValYear}{}\renewcommand{\rhtaValYear}{2013}  % validation year of the pre-sample fold on which every arm's linear penalties are chosen
\providecommand{\rhtaOptDays}{}\renewcommand{\rhtaOptDays}{614}  % post-warmup option days of the Sharpe ratio (Section 7; D19 pooled ratio)
\providecommand{\rhtaNClasses}{}\renewcommand{\rhtaNClasses}{6}  % model classes (ridge, lasso, elastic net, LightGBM, XGBoost, random forest)
\providecommand{\rhtaWsMin}{}\renewcommand{\rhtaWsMin}{1{,}000}  % smallest window of the close-only sweep, closes
\providecommand{\rhtaWsMax}{}\renewcommand{\rhtaWsMax}{3{,}500}  % largest window of the close-only sweep, closes
\providecommand{\rhtaPctRidgeRthClose}{}\renewcommand{\rhtaPctRidgeRthClose}{0.8}  % QLIKE % change, full set vs HAR benchmark: Ridge, RTH13 (main design: the 13 regular-hours bars, 3,000 days), 16:00 bar
\providecommand{\rhtaDmRidgeRthClose}{}\renewcommand{\rhtaDmRidgeRthClose}{0.31}  % ... its Diebold-Mariano statistic (negative: full more accurate)
\providecommand{\rhtaPRidgeRthClose}{}\renewcommand{\rhtaPRidgeRthClose}{\ensuremath{p = 0.76}}  % ... its p statement
\providecommand{\rhtaQHarRidgeRthClose}{}\renewcommand{\rhtaQHarRidgeRthClose}{0.1269}  % ... mean QLIKE, HAR benchmark
\providecommand{\rhtaQFullRidgeRthClose}{}\renewcommand{\rhtaQFullRidgeRthClose}{0.1279}  % ... mean QLIKE, full set
\providecommand{\rhtaPctLassoRthClose}{}\renewcommand{\rhtaPctLassoRthClose}{3.3}  % QLIKE % change, full set vs HAR benchmark: Lasso, RTH13 (main design: the 13 regular-hours bars, 3,000 days), 16:00 bar
\providecommand{\rhtaDmLassoRthClose}{}\renewcommand{\rhtaDmLassoRthClose}{1.04}  % ... its Diebold-Mariano statistic (negative: full more accurate)
\providecommand{\rhtaPLassoRthClose}{}\renewcommand{\rhtaPLassoRthClose}{\ensuremath{p = 0.30}}  % ... its p statement
\providecommand{\rhtaQHarLassoRthClose}{}\renewcommand{\rhtaQHarLassoRthClose}{0.1269}  % ... mean QLIKE, HAR benchmark
\providecommand{\rhtaQFullLassoRthClose}{}\renewcommand{\rhtaQFullLassoRthClose}{0.1310}  % ... mean QLIKE, full set
\providecommand{\rhtaPctEnetRthClose}{}\renewcommand{\rhtaPctEnetRthClose}{\ensuremath{-}0.5}  % QLIKE % change, full set vs HAR benchmark: Elastic net, RTH13 (main design: the 13 regular-hours bars, 3,000 days), 16:00 bar
\providecommand{\rhtaDmEnetRthClose}{}\renewcommand{\rhtaDmEnetRthClose}{\ensuremath{-}0.21}  % ... its Diebold-Mariano statistic (negative: full more accurate)
\providecommand{\rhtaPEnetRthClose}{}\renewcommand{\rhtaPEnetRthClose}{\ensuremath{p = 0.83}}  % ... its p statement
\providecommand{\rhtaQHarEnetRthClose}{}\renewcommand{\rhtaQHarEnetRthClose}{0.1269}  % ... mean QLIKE, HAR benchmark
\providecommand{\rhtaQFullEnetRthClose}{}\renewcommand{\rhtaQFullEnetRthClose}{0.1262}  % ... mean QLIKE, full set
\providecommand{\rhtaPctLgbmRthClose}{}\renewcommand{\rhtaPctLgbmRthClose}{\ensuremath{-}2.7}  % QLIKE % change, full set vs HAR benchmark: LightGBM, RTH13 (main design: the 13 regular-hours bars, 3,000 days), 16:00 bar
\providecommand{\rhtaDmLgbmRthClose}{}\renewcommand{\rhtaDmLgbmRthClose}{\ensuremath{-}1.81}  % ... its Diebold-Mariano statistic (negative: full more accurate)
\providecommand{\rhtaPLgbmRthClose}{}\renewcommand{\rhtaPLgbmRthClose}{\ensuremath{p = 0.071}}  % ... its p statement
\providecommand{\rhtaQHarLgbmRthClose}{}\renewcommand{\rhtaQHarLgbmRthClose}{0.1235}  % ... mean QLIKE, HAR benchmark
\providecommand{\rhtaQFullLgbmRthClose}{}\renewcommand{\rhtaQFullLgbmRthClose}{0.1201}  % ... mean QLIKE, full set
\providecommand{\rhtaPctXgbRthClose}{}\renewcommand{\rhtaPctXgbRthClose}{\ensuremath{-}3.4}  % QLIKE % change, full set vs HAR benchmark: XGBoost, RTH13 (main design: the 13 regular-hours bars, 3,000 days), 16:00 bar
\providecommand{\rhtaDmXgbRthClose}{}\renewcommand{\rhtaDmXgbRthClose}{\ensuremath{-}2.18}  % ... its Diebold-Mariano statistic (negative: full more accurate)
\providecommand{\rhtaPXgbRthClose}{}\renewcommand{\rhtaPXgbRthClose}{\ensuremath{p = 0.030}}  % ... its p statement
\providecommand{\rhtaQHarXgbRthClose}{}\renewcommand{\rhtaQHarXgbRthClose}{0.1234}  % ... mean QLIKE, HAR benchmark
\providecommand{\rhtaQFullXgbRthClose}{}\renewcommand{\rhtaQFullXgbRthClose}{0.1192}  % ... mean QLIKE, full set
\providecommand{\rhtaPctRfRthClose}{}\renewcommand{\rhtaPctRfRthClose}{\ensuremath{-}1.4}  % QLIKE % change, full set vs HAR benchmark: Random forest, RTH13 (main design: the 13 regular-hours bars, 3,000 days), 16:00 bar
\providecommand{\rhtaDmRfRthClose}{}\renewcommand{\rhtaDmRfRthClose}{\ensuremath{-}0.92}  % ... its Diebold-Mariano statistic (negative: full more accurate)
\providecommand{\rhtaPRfRthClose}{}\renewcommand{\rhtaPRfRthClose}{\ensuremath{p = 0.36}}  % ... its p statement
\providecommand{\rhtaQHarRfRthClose}{}\renewcommand{\rhtaQHarRfRthClose}{0.1266}  % ... mean QLIKE, HAR benchmark
\providecommand{\rhtaQFullRfRthClose}{}\renewcommand{\rhtaQFullRfRthClose}{0.1248}  % ... mean QLIKE, full set
\providecommand{\rhtaPctRidgeRthBar}{}\renewcommand{\rhtaPctRidgeRthBar}{\ensuremath{-}5.4}  % QLIKE % change, full set vs HAR benchmark: Ridge, RTH13 (main design: the 13 regular-hours bars, 3,000 days), every bar
\providecommand{\rhtaDmRidgeRthBar}{}\renewcommand{\rhtaDmRidgeRthBar}{\ensuremath{-}5.63}  % ... its Diebold-Mariano statistic (negative: full more accurate)
\providecommand{\rhtaPRidgeRthBar}{}\renewcommand{\rhtaPRidgeRthBar}{\ensuremath{p < 0.001}}  % ... its p statement
\providecommand{\rhtaQHarRidgeRthBar}{}\renewcommand{\rhtaQHarRidgeRthBar}{0.1893}  % ... mean QLIKE, HAR benchmark
\providecommand{\rhtaQFullRidgeRthBar}{}\renewcommand{\rhtaQFullRidgeRthBar}{0.1791}  % ... mean QLIKE, full set
\providecommand{\rhtaPctLassoRthBar}{}\renewcommand{\rhtaPctLassoRthBar}{\ensuremath{-}6.2}  % QLIKE % change, full set vs HAR benchmark: Lasso, RTH13 (main design: the 13 regular-hours bars, 3,000 days), every bar
\providecommand{\rhtaDmLassoRthBar}{}\renewcommand{\rhtaDmLassoRthBar}{\ensuremath{-}4.72}  % ... its Diebold-Mariano statistic (negative: full more accurate)
\providecommand{\rhtaPLassoRthBar}{}\renewcommand{\rhtaPLassoRthBar}{\ensuremath{p < 0.001}}  % ... its p statement
\providecommand{\rhtaQHarLassoRthBar}{}\renewcommand{\rhtaQHarLassoRthBar}{0.1893}  % ... mean QLIKE, HAR benchmark
\providecommand{\rhtaQFullLassoRthBar}{}\renewcommand{\rhtaQFullLassoRthBar}{0.1775}  % ... mean QLIKE, full set
\providecommand{\rhtaPctEnetRthBar}{}\renewcommand{\rhtaPctEnetRthBar}{\ensuremath{-}6.0}  % QLIKE % change, full set vs HAR benchmark: Elastic net, RTH13 (main design: the 13 regular-hours bars, 3,000 days), every bar
\providecommand{\rhtaDmEnetRthBar}{}\renewcommand{\rhtaDmEnetRthBar}{\ensuremath{-}6.98}  % ... its Diebold-Mariano statistic (negative: full more accurate)
\providecommand{\rhtaPEnetRthBar}{}\renewcommand{\rhtaPEnetRthBar}{\ensuremath{p < 0.001}}  % ... its p statement
\providecommand{\rhtaQHarEnetRthBar}{}\renewcommand{\rhtaQHarEnetRthBar}{0.1893}  % ... mean QLIKE, HAR benchmark
\providecommand{\rhtaQFullEnetRthBar}{}\renewcommand{\rhtaQFullEnetRthBar}{0.1778}  % ... mean QLIKE, full set
\providecommand{\rhtaPctLgbmRthBar}{}\renewcommand{\rhtaPctLgbmRthBar}{\ensuremath{-}14.0}  % QLIKE % change, full set vs HAR benchmark: LightGBM, RTH13 (main design: the 13 regular-hours bars, 3,000 days), every bar
\providecommand{\rhtaDmLgbmRthBar}{}\renewcommand{\rhtaDmLgbmRthBar}{\ensuremath{-}10.32}  % ... its Diebold-Mariano statistic (negative: full more accurate)
\providecommand{\rhtaPLgbmRthBar}{}\renewcommand{\rhtaPLgbmRthBar}{\ensuremath{p < 0.001}}  % ... its p statement
\providecommand{\rhtaQHarLgbmRthBar}{}\renewcommand{\rhtaQHarLgbmRthBar}{0.1875}  % ... mean QLIKE, HAR benchmark
\providecommand{\rhtaQFullLgbmRthBar}{}\renewcommand{\rhtaQFullLgbmRthBar}{0.1613}  % ... mean QLIKE, full set
\providecommand{\rhtaPctXgbRthBar}{}\renewcommand{\rhtaPctXgbRthBar}{\ensuremath{-}13.8}  % QLIKE % change, full set vs HAR benchmark: XGBoost, RTH13 (main design: the 13 regular-hours bars, 3,000 days), every bar
\providecommand{\rhtaDmXgbRthBar}{}\renewcommand{\rhtaDmXgbRthBar}{\ensuremath{-}9.94}  % ... its Diebold-Mariano statistic (negative: full more accurate)
\providecommand{\rhtaPXgbRthBar}{}\renewcommand{\rhtaPXgbRthBar}{\ensuremath{p < 0.001}}  % ... its p statement
\providecommand{\rhtaQHarXgbRthBar}{}\renewcommand{\rhtaQHarXgbRthBar}{0.1888}  % ... mean QLIKE, HAR benchmark
\providecommand{\rhtaQFullXgbRthBar}{}\renewcommand{\rhtaQFullXgbRthBar}{0.1627}  % ... mean QLIKE, full set
\providecommand{\rhtaPctRfRthBar}{}\renewcommand{\rhtaPctRfRthBar}{\ensuremath{-}11.1}  % QLIKE % change, full set vs HAR benchmark: Random forest, RTH13 (main design: the 13 regular-hours bars, 3,000 days), every bar
\providecommand{\rhtaDmRfRthBar}{}\renewcommand{\rhtaDmRfRthBar}{\ensuremath{-}9.85}  % ... its Diebold-Mariano statistic (negative: full more accurate)
\providecommand{\rhtaPRfRthBar}{}\renewcommand{\rhtaPRfRthBar}{\ensuremath{p < 0.001}}  % ... its p statement
\providecommand{\rhtaQHarRfRthBar}{}\renewcommand{\rhtaQHarRfRthBar}{0.1895}  % ... mean QLIKE, HAR benchmark
\providecommand{\rhtaQFullRfRthBar}{}\renewcommand{\rhtaQFullRfRthBar}{0.1685}  % ... mean QLIKE, full set
\providecommand{\rhtaPctRidgeAllClose}{}\renewcommand{\rhtaPctRidgeAllClose}{\ensuremath{-}2.2}  % QLIKE % change, full set vs HAR benchmark: Ridge, ALL48 (all 48 intervals, 3,000 days), 16:00 bar
\providecommand{\rhtaDmRidgeAllClose}{}\renewcommand{\rhtaDmRidgeAllClose}{\ensuremath{-}1.18}  % ... its Diebold-Mariano statistic (negative: full more accurate)
\providecommand{\rhtaPRidgeAllClose}{}\renewcommand{\rhtaPRidgeAllClose}{\ensuremath{p = 0.24}}  % ... its p statement
\providecommand{\rhtaQHarRidgeAllClose}{}\renewcommand{\rhtaQHarRidgeAllClose}{0.1262}  % ... mean QLIKE, HAR benchmark
\providecommand{\rhtaQFullRidgeAllClose}{}\renewcommand{\rhtaQFullRidgeAllClose}{0.1234}  % ... mean QLIKE, full set
\providecommand{\rhtaPctLassoAllClose}{}\renewcommand{\rhtaPctLassoAllClose}{\ensuremath{-}2.6}  % QLIKE % change, full set vs HAR benchmark: Lasso, ALL48 (all 48 intervals, 3,000 days), 16:00 bar
\providecommand{\rhtaDmLassoAllClose}{}\renewcommand{\rhtaDmLassoAllClose}{\ensuremath{-}1.45}  % ... its Diebold-Mariano statistic (negative: full more accurate)
\providecommand{\rhtaPLassoAllClose}{}\renewcommand{\rhtaPLassoAllClose}{\ensuremath{p = 0.15}}  % ... its p statement
\providecommand{\rhtaQHarLassoAllClose}{}\renewcommand{\rhtaQHarLassoAllClose}{0.1264}  % ... mean QLIKE, HAR benchmark
\providecommand{\rhtaQFullLassoAllClose}{}\renewcommand{\rhtaQFullLassoAllClose}{0.1232}  % ... mean QLIKE, full set
\providecommand{\rhtaPctEnetAllClose}{}\renewcommand{\rhtaPctEnetAllClose}{\ensuremath{-}1.8}  % QLIKE % change, full set vs HAR benchmark: Elastic net, ALL48 (all 48 intervals, 3,000 days), 16:00 bar
\providecommand{\rhtaDmEnetAllClose}{}\renewcommand{\rhtaDmEnetAllClose}{\ensuremath{-}0.96}  % ... its Diebold-Mariano statistic (negative: full more accurate)
\providecommand{\rhtaPEnetAllClose}{}\renewcommand{\rhtaPEnetAllClose}{\ensuremath{p = 0.34}}  % ... its p statement
\providecommand{\rhtaQHarEnetAllClose}{}\renewcommand{\rhtaQHarEnetAllClose}{0.1259}  % ... mean QLIKE, HAR benchmark
\providecommand{\rhtaQFullEnetAllClose}{}\renewcommand{\rhtaQFullEnetAllClose}{0.1236}  % ... mean QLIKE, full set
\providecommand{\rhtaPctLgbmAllClose}{}\renewcommand{\rhtaPctLgbmAllClose}{\ensuremath{-}6.1}  % QLIKE % change, full set vs HAR benchmark: LightGBM, ALL48 (all 48 intervals, 3,000 days), 16:00 bar
\providecommand{\rhtaDmLgbmAllClose}{}\renewcommand{\rhtaDmLgbmAllClose}{\ensuremath{-}4.45}  % ... its Diebold-Mariano statistic (negative: full more accurate)
\providecommand{\rhtaPLgbmAllClose}{}\renewcommand{\rhtaPLgbmAllClose}{\ensuremath{p < 0.001}}  % ... its p statement
\providecommand{\rhtaQHarLgbmAllClose}{}\renewcommand{\rhtaQHarLgbmAllClose}{0.1250}  % ... mean QLIKE, HAR benchmark
\providecommand{\rhtaQFullLgbmAllClose}{}\renewcommand{\rhtaQFullLgbmAllClose}{0.1174}  % ... mean QLIKE, full set
\providecommand{\rhtaPctXgbAllClose}{}\renewcommand{\rhtaPctXgbAllClose}{\ensuremath{-}6.5}  % QLIKE % change, full set vs HAR benchmark: XGBoost, ALL48 (all 48 intervals, 3,000 days), 16:00 bar
\providecommand{\rhtaDmXgbAllClose}{}\renewcommand{\rhtaDmXgbAllClose}{\ensuremath{-}5.03}  % ... its Diebold-Mariano statistic (negative: full more accurate)
\providecommand{\rhtaPXgbAllClose}{}\renewcommand{\rhtaPXgbAllClose}{\ensuremath{p < 0.001}}  % ... its p statement
\providecommand{\rhtaQHarXgbAllClose}{}\renewcommand{\rhtaQHarXgbAllClose}{0.1251}  % ... mean QLIKE, HAR benchmark
\providecommand{\rhtaQFullXgbAllClose}{}\renewcommand{\rhtaQFullXgbAllClose}{0.1170}  % ... mean QLIKE, full set
\providecommand{\rhtaPctRfAllClose}{}\renewcommand{\rhtaPctRfAllClose}{\ensuremath{-}5.6}  % QLIKE % change, full set vs HAR benchmark: Random forest, ALL48 (all 48 intervals, 3,000 days), 16:00 bar
\providecommand{\rhtaDmRfAllClose}{}\renewcommand{\rhtaDmRfAllClose}{\ensuremath{-}4.11}  % ... its Diebold-Mariano statistic (negative: full more accurate)
\providecommand{\rhtaPRfAllClose}{}\renewcommand{\rhtaPRfAllClose}{\ensuremath{p < 0.001}}  % ... its p statement
\providecommand{\rhtaQHarRfAllClose}{}\renewcommand{\rhtaQHarRfAllClose}{0.1265}  % ... mean QLIKE, HAR benchmark
\providecommand{\rhtaQFullRfAllClose}{}\renewcommand{\rhtaQFullRfAllClose}{0.1194}  % ... mean QLIKE, full set
\providecommand{\rhtaPctRidgeAllBar}{}\renewcommand{\rhtaPctRidgeAllBar}{\ensuremath{-}4.9}  % QLIKE % change, full set vs HAR benchmark: Ridge, ALL48 (all 48 intervals, 3,000 days), every bar
\providecommand{\rhtaDmRidgeAllBar}{}\renewcommand{\rhtaDmRidgeAllBar}{\ensuremath{-}6.71}  % ... its Diebold-Mariano statistic (negative: full more accurate)
\providecommand{\rhtaPRidgeAllBar}{}\renewcommand{\rhtaPRidgeAllBar}{\ensuremath{p < 0.001}}  % ... its p statement
\providecommand{\rhtaQHarRidgeAllBar}{}\renewcommand{\rhtaQHarRidgeAllBar}{0.1919}  % ... mean QLIKE, HAR benchmark
\providecommand{\rhtaQFullRidgeAllBar}{}\renewcommand{\rhtaQFullRidgeAllBar}{0.1825}  % ... mean QLIKE, full set
\providecommand{\rhtaPctLassoAllBar}{}\renewcommand{\rhtaPctLassoAllBar}{\ensuremath{-}5.1}  % QLIKE % change, full set vs HAR benchmark: Lasso, ALL48 (all 48 intervals, 3,000 days), every bar
\providecommand{\rhtaDmLassoAllBar}{}\renewcommand{\rhtaDmLassoAllBar}{\ensuremath{-}7.14}  % ... its Diebold-Mariano statistic (negative: full more accurate)
\providecommand{\rhtaPLassoAllBar}{}\renewcommand{\rhtaPLassoAllBar}{\ensuremath{p < 0.001}}  % ... its p statement
\providecommand{\rhtaQHarLassoAllBar}{}\renewcommand{\rhtaQHarLassoAllBar}{0.1917}  % ... mean QLIKE, HAR benchmark
\providecommand{\rhtaQFullLassoAllBar}{}\renewcommand{\rhtaQFullLassoAllBar}{0.1820}  % ... mean QLIKE, full set
\providecommand{\rhtaPctEnetAllBar}{}\renewcommand{\rhtaPctEnetAllBar}{\ensuremath{-}5.2}  % QLIKE % change, full set vs HAR benchmark: Elastic net, ALL48 (all 48 intervals, 3,000 days), every bar
\providecommand{\rhtaDmEnetAllBar}{}\renewcommand{\rhtaDmEnetAllBar}{\ensuremath{-}6.68}  % ... its Diebold-Mariano statistic (negative: full more accurate)
\providecommand{\rhtaPEnetAllBar}{}\renewcommand{\rhtaPEnetAllBar}{\ensuremath{p < 0.001}}  % ... its p statement
\providecommand{\rhtaQHarEnetAllBar}{}\renewcommand{\rhtaQHarEnetAllBar}{0.1916}  % ... mean QLIKE, HAR benchmark
\providecommand{\rhtaQFullEnetAllBar}{}\renewcommand{\rhtaQFullEnetAllBar}{0.1817}  % ... mean QLIKE, full set
\providecommand{\rhtaPctLgbmAllBar}{}\renewcommand{\rhtaPctLgbmAllBar}{\ensuremath{-}11.2}  % QLIKE % change, full set vs HAR benchmark: LightGBM, ALL48 (all 48 intervals, 3,000 days), every bar
\providecommand{\rhtaDmLgbmAllBar}{}\renewcommand{\rhtaDmLgbmAllBar}{\ensuremath{-}9.84}  % ... its Diebold-Mariano statistic (negative: full more accurate)
\providecommand{\rhtaPLgbmAllBar}{}\renewcommand{\rhtaPLgbmAllBar}{\ensuremath{p < 0.001}}  % ... its p statement
\providecommand{\rhtaQHarLgbmAllBar}{}\renewcommand{\rhtaQHarLgbmAllBar}{0.1898}  % ... mean QLIKE, HAR benchmark
\providecommand{\rhtaQFullLgbmAllBar}{}\renewcommand{\rhtaQFullLgbmAllBar}{0.1685}  % ... mean QLIKE, full set
\providecommand{\rhtaPctXgbAllBar}{}\renewcommand{\rhtaPctXgbAllBar}{\ensuremath{-}12.0}  % QLIKE % change, full set vs HAR benchmark: XGBoost, ALL48 (all 48 intervals, 3,000 days), every bar
\providecommand{\rhtaDmXgbAllBar}{}\renewcommand{\rhtaDmXgbAllBar}{\ensuremath{-}9.70}  % ... its Diebold-Mariano statistic (negative: full more accurate)
\providecommand{\rhtaPXgbAllBar}{}\renewcommand{\rhtaPXgbAllBar}{\ensuremath{p < 0.001}}  % ... its p statement
\providecommand{\rhtaQHarXgbAllBar}{}\renewcommand{\rhtaQHarXgbAllBar}{0.1908}  % ... mean QLIKE, HAR benchmark
\providecommand{\rhtaQFullXgbAllBar}{}\renewcommand{\rhtaQFullXgbAllBar}{0.1679}  % ... mean QLIKE, full set
\providecommand{\rhtaPctRfAllBar}{}\renewcommand{\rhtaPctRfAllBar}{\ensuremath{-}7.0}  % QLIKE % change, full set vs HAR benchmark: Random forest, ALL48 (all 48 intervals, 3,000 days), every bar
\providecommand{\rhtaDmRfAllBar}{}\renewcommand{\rhtaDmRfAllBar}{\ensuremath{-}9.54}  % ... its Diebold-Mariano statistic (negative: full more accurate)
\providecommand{\rhtaPRfAllBar}{}\renewcommand{\rhtaPRfAllBar}{\ensuremath{p < 0.001}}  % ... its p statement
\providecommand{\rhtaQHarRfAllBar}{}\renewcommand{\rhtaQHarRfAllBar}{0.1905}  % ... mean QLIKE, HAR benchmark
\providecommand{\rhtaQFullRfAllBar}{}\renewcommand{\rhtaQFullRfAllBar}{0.1773}  % ... mean QLIKE, full set
\providecommand{\rhtaPctRidgeClClose}{}\renewcommand{\rhtaPctRidgeClClose}{\ensuremath{-}6.1}  % QLIKE % change, full set vs HAR benchmark: Ridge, CLOSE (the closing half hour only, 3,000 closes), 16:00 bar
\providecommand{\rhtaDmRidgeClClose}{}\renewcommand{\rhtaDmRidgeClClose}{\ensuremath{-}3.62}  % ... its Diebold-Mariano statistic (negative: full more accurate)
\providecommand{\rhtaPRidgeClClose}{}\renewcommand{\rhtaPRidgeClClose}{\ensuremath{p < 0.001}}  % ... its p statement
\providecommand{\rhtaQHarRidgeClClose}{}\renewcommand{\rhtaQHarRidgeClClose}{0.1223}  % ... mean QLIKE, HAR benchmark
\providecommand{\rhtaQFullRidgeClClose}{}\renewcommand{\rhtaQFullRidgeClClose}{0.1148}  % ... mean QLIKE, full set
\providecommand{\rhtaPctLassoClClose}{}\renewcommand{\rhtaPctLassoClClose}{\ensuremath{-}5.4}  % QLIKE % change, full set vs HAR benchmark: Lasso, CLOSE (the closing half hour only, 3,000 closes), 16:00 bar
\providecommand{\rhtaDmLassoClClose}{}\renewcommand{\rhtaDmLassoClClose}{\ensuremath{-}5.19}  % ... its Diebold-Mariano statistic (negative: full more accurate)
\providecommand{\rhtaPLassoClClose}{}\renewcommand{\rhtaPLassoClClose}{\ensuremath{p < 0.001}}  % ... its p statement
\providecommand{\rhtaQHarLassoClClose}{}\renewcommand{\rhtaQHarLassoClClose}{0.1223}  % ... mean QLIKE, HAR benchmark
\providecommand{\rhtaQFullLassoClClose}{}\renewcommand{\rhtaQFullLassoClClose}{0.1156}  % ... mean QLIKE, full set
\providecommand{\rhtaPctEnetClClose}{}\renewcommand{\rhtaPctEnetClClose}{\ensuremath{-}6.5}  % QLIKE % change, full set vs HAR benchmark: Elastic net, CLOSE (the closing half hour only, 3,000 closes), 16:00 bar
\providecommand{\rhtaDmEnetClClose}{}\renewcommand{\rhtaDmEnetClClose}{\ensuremath{-}4.59}  % ... its Diebold-Mariano statistic (negative: full more accurate)
\providecommand{\rhtaPEnetClClose}{}\renewcommand{\rhtaPEnetClClose}{\ensuremath{p < 0.001}}  % ... its p statement
\providecommand{\rhtaQHarEnetClClose}{}\renewcommand{\rhtaQHarEnetClClose}{0.1223}  % ... mean QLIKE, HAR benchmark
\providecommand{\rhtaQFullEnetClClose}{}\renewcommand{\rhtaQFullEnetClClose}{0.1144}  % ... mean QLIKE, full set
\providecommand{\rhtaPctLgbmClClose}{}\renewcommand{\rhtaPctLgbmClClose}{\ensuremath{-}7.0}  % QLIKE % change, full set vs HAR benchmark: LightGBM, CLOSE (the closing half hour only, 3,000 closes), 16:00 bar
\providecommand{\rhtaDmLgbmClClose}{}\renewcommand{\rhtaDmLgbmClClose}{\ensuremath{-}4.50}  % ... its Diebold-Mariano statistic (negative: full more accurate)
\providecommand{\rhtaPLgbmClClose}{}\renewcommand{\rhtaPLgbmClClose}{\ensuremath{p < 0.001}}  % ... its p statement
\providecommand{\rhtaQHarLgbmClClose}{}\renewcommand{\rhtaQHarLgbmClClose}{0.1313}  % ... mean QLIKE, HAR benchmark
\providecommand{\rhtaQFullLgbmClClose}{}\renewcommand{\rhtaQFullLgbmClClose}{0.1221}  % ... mean QLIKE, full set
\providecommand{\rhtaPctXgbClClose}{}\renewcommand{\rhtaPctXgbClClose}{\ensuremath{-}8.5}  % QLIKE % change, full set vs HAR benchmark: XGBoost, CLOSE (the closing half hour only, 3,000 closes), 16:00 bar
\providecommand{\rhtaDmXgbClClose}{}\renewcommand{\rhtaDmXgbClClose}{\ensuremath{-}5.95}  % ... its Diebold-Mariano statistic (negative: full more accurate)
\providecommand{\rhtaPXgbClClose}{}\renewcommand{\rhtaPXgbClClose}{\ensuremath{p < 0.001}}  % ... its p statement
\providecommand{\rhtaQHarXgbClClose}{}\renewcommand{\rhtaQHarXgbClClose}{0.1290}  % ... mean QLIKE, HAR benchmark
\providecommand{\rhtaQFullXgbClClose}{}\renewcommand{\rhtaQFullXgbClClose}{0.1180}  % ... mean QLIKE, full set
\providecommand{\rhtaPctRfClClose}{}\renewcommand{\rhtaPctRfClClose}{\ensuremath{-}3.7}  % QLIKE % change, full set vs HAR benchmark: Random forest, CLOSE (the closing half hour only, 3,000 closes), 16:00 bar
\providecommand{\rhtaDmRfClClose}{}\renewcommand{\rhtaDmRfClClose}{\ensuremath{-}3.66}  % ... its Diebold-Mariano statistic (negative: full more accurate)
\providecommand{\rhtaPRfClClose}{}\renewcommand{\rhtaPRfClClose}{\ensuremath{p < 0.001}}  % ... its p statement
\providecommand{\rhtaQHarRfClClose}{}\renewcommand{\rhtaQHarRfClClose}{0.1267}  % ... mean QLIKE, HAR benchmark
\providecommand{\rhtaQFullRfClClose}{}\renewcommand{\rhtaQFullRfClClose}{0.1220}  % ... mean QLIKE, full set
\providecommand{\rhtaPctRidgeClNatClose}{}\renewcommand{\rhtaPctRidgeClNatClose}{\ensuremath{-}6.5}  % QLIKE % change, full set vs HAR benchmark: Ridge, CLOSE_NAT (closing rows, design built natively on them; linear only), 16:00 bar
\providecommand{\rhtaDmRidgeClNatClose}{}\renewcommand{\rhtaDmRidgeClNatClose}{\ensuremath{-}3.76}  % ... its Diebold-Mariano statistic (negative: full more accurate)
\providecommand{\rhtaPRidgeClNatClose}{}\renewcommand{\rhtaPRidgeClNatClose}{\ensuremath{p < 0.001}}  % ... its p statement
\providecommand{\rhtaQHarRidgeClNatClose}{}\renewcommand{\rhtaQHarRidgeClNatClose}{0.1226}  % ... mean QLIKE, HAR benchmark
\providecommand{\rhtaQFullRidgeClNatClose}{}\renewcommand{\rhtaQFullRidgeClNatClose}{0.1146}  % ... mean QLIKE, full set
\providecommand{\rhtaPctLassoClNatClose}{}\renewcommand{\rhtaPctLassoClNatClose}{\ensuremath{-}5.4}  % QLIKE % change, full set vs HAR benchmark: Lasso, CLOSE_NAT (closing rows, design built natively on them; linear only), 16:00 bar
\providecommand{\rhtaDmLassoClNatClose}{}\renewcommand{\rhtaDmLassoClNatClose}{\ensuremath{-}4.95}  % ... its Diebold-Mariano statistic (negative: full more accurate)
\providecommand{\rhtaPLassoClNatClose}{}\renewcommand{\rhtaPLassoClNatClose}{\ensuremath{p < 0.001}}  % ... its p statement
\providecommand{\rhtaQHarLassoClNatClose}{}\renewcommand{\rhtaQHarLassoClNatClose}{0.1226}  % ... mean QLIKE, HAR benchmark
\providecommand{\rhtaQFullLassoClNatClose}{}\renewcommand{\rhtaQFullLassoClNatClose}{0.1159}  % ... mean QLIKE, full set
\providecommand{\rhtaPctEnetClNatClose}{}\renewcommand{\rhtaPctEnetClNatClose}{\ensuremath{-}6.6}  % QLIKE % change, full set vs HAR benchmark: Elastic net, CLOSE_NAT (closing rows, design built natively on them; linear only), 16:00 bar
\providecommand{\rhtaDmEnetClNatClose}{}\renewcommand{\rhtaDmEnetClNatClose}{\ensuremath{-}4.48}  % ... its Diebold-Mariano statistic (negative: full more accurate)
\providecommand{\rhtaPEnetClNatClose}{}\renewcommand{\rhtaPEnetClNatClose}{\ensuremath{p < 0.001}}  % ... its p statement
\providecommand{\rhtaQHarEnetClNatClose}{}\renewcommand{\rhtaQHarEnetClNatClose}{0.1226}  % ... mean QLIKE, HAR benchmark
\providecommand{\rhtaQFullEnetClNatClose}{}\renewcommand{\rhtaQFullEnetClNatClose}{0.1145}  % ... mean QLIKE, full set
\providecommand{\rhtaPctRidgeAllCtClose}{}\renewcommand{\rhtaPctRidgeAllCtClose}{\ensuremath{-}1.9}  % QLIKE % change, full set vs HAR benchmark: Ridge, ALL48@16 (ALL48, penalties chosen on the 16:00 validation rows; linear only), 16:00 bar
\providecommand{\rhtaDmRidgeAllCtClose}{}\renewcommand{\rhtaDmRidgeAllCtClose}{\ensuremath{-}0.97}  % ... its Diebold-Mariano statistic (negative: full more accurate)
\providecommand{\rhtaPRidgeAllCtClose}{}\renewcommand{\rhtaPRidgeAllCtClose}{\ensuremath{p = 0.33}}  % ... its p statement
\providecommand{\rhtaQHarRidgeAllCtClose}{}\renewcommand{\rhtaQHarRidgeAllCtClose}{0.1262}  % ... mean QLIKE, HAR benchmark
\providecommand{\rhtaQFullRidgeAllCtClose}{}\renewcommand{\rhtaQFullRidgeAllCtClose}{0.1238}  % ... mean QLIKE, full set
\providecommand{\rhtaPctLassoAllCtClose}{}\renewcommand{\rhtaPctLassoAllCtClose}{\ensuremath{-}1.1}  % QLIKE % change, full set vs HAR benchmark: Lasso, ALL48@16 (ALL48, penalties chosen on the 16:00 validation rows; linear only), 16:00 bar
\providecommand{\rhtaDmLassoAllCtClose}{}\renewcommand{\rhtaDmLassoAllCtClose}{\ensuremath{-}0.53}  % ... its Diebold-Mariano statistic (negative: full more accurate)
\providecommand{\rhtaPLassoAllCtClose}{}\renewcommand{\rhtaPLassoAllCtClose}{\ensuremath{p = 0.59}}  % ... its p statement
\providecommand{\rhtaQHarLassoAllCtClose}{}\renewcommand{\rhtaQHarLassoAllCtClose}{0.1253}  % ... mean QLIKE, HAR benchmark
\providecommand{\rhtaQFullLassoAllCtClose}{}\renewcommand{\rhtaQFullLassoAllCtClose}{0.1239}  % ... mean QLIKE, full set
\providecommand{\rhtaPctEnetAllCtClose}{}\renewcommand{\rhtaPctEnetAllCtClose}{\ensuremath{-}0.8}  % QLIKE % change, full set vs HAR benchmark: Elastic net, ALL48@16 (ALL48, penalties chosen on the 16:00 validation rows; linear only), 16:00 bar
\providecommand{\rhtaDmEnetAllCtClose}{}\renewcommand{\rhtaDmEnetAllCtClose}{\ensuremath{-}0.40}  % ... its Diebold-Mariano statistic (negative: full more accurate)
\providecommand{\rhtaPEnetAllCtClose}{}\renewcommand{\rhtaPEnetAllCtClose}{\ensuremath{p = 0.69}}  % ... its p statement
\providecommand{\rhtaQHarEnetAllCtClose}{}\renewcommand{\rhtaQHarEnetAllCtClose}{0.1246}  % ... mean QLIKE, HAR benchmark
\providecommand{\rhtaQFullEnetAllCtClose}{}\renewcommand{\rhtaQFullEnetAllCtClose}{0.1236}  % ... mean QLIKE, full set
\providecommand{\rhtaPctRidgeAllCtBar}{}\renewcommand{\rhtaPctRidgeAllCtBar}{\ensuremath{-}5.2}  % QLIKE % change, full set vs HAR benchmark: Ridge, ALL48@16 (ALL48, penalties chosen on the 16:00 validation rows; linear only), every bar
\providecommand{\rhtaDmRidgeAllCtBar}{}\renewcommand{\rhtaDmRidgeAllCtBar}{\ensuremath{-}6.22}  % ... its Diebold-Mariano statistic (negative: full more accurate)
\providecommand{\rhtaPRidgeAllCtBar}{}\renewcommand{\rhtaPRidgeAllCtBar}{\ensuremath{p < 0.001}}  % ... its p statement
\providecommand{\rhtaQHarRidgeAllCtBar}{}\renewcommand{\rhtaQHarRidgeAllCtBar}{0.1919}  % ... mean QLIKE, HAR benchmark
\providecommand{\rhtaQFullRidgeAllCtBar}{}\renewcommand{\rhtaQFullRidgeAllCtBar}{0.1820}  % ... mean QLIKE, full set
\providecommand{\rhtaPctLassoAllCtBar}{}\renewcommand{\rhtaPctLassoAllCtBar}{\ensuremath{-}5.3}  % QLIKE % change, full set vs HAR benchmark: Lasso, ALL48@16 (ALL48, penalties chosen on the 16:00 validation rows; linear only), every bar
\providecommand{\rhtaDmLassoAllCtBar}{}\renewcommand{\rhtaDmLassoAllCtBar}{\ensuremath{-}6.18}  % ... its Diebold-Mariano statistic (negative: full more accurate)
\providecommand{\rhtaPLassoAllCtBar}{}\renewcommand{\rhtaPLassoAllCtBar}{\ensuremath{p < 0.001}}  % ... its p statement
\providecommand{\rhtaQHarLassoAllCtBar}{}\renewcommand{\rhtaQHarLassoAllCtBar}{0.1915}  % ... mean QLIKE, HAR benchmark
\providecommand{\rhtaQFullLassoAllCtBar}{}\renewcommand{\rhtaQFullLassoAllCtBar}{0.1814}  % ... mean QLIKE, full set
\providecommand{\rhtaPctEnetAllCtBar}{}\renewcommand{\rhtaPctEnetAllCtBar}{\ensuremath{-}5.5}  % QLIKE % change, full set vs HAR benchmark: Elastic net, ALL48@16 (ALL48, penalties chosen on the 16:00 validation rows; linear only), every bar
\providecommand{\rhtaDmEnetAllCtBar}{}\renewcommand{\rhtaDmEnetAllCtBar}{\ensuremath{-}7.11}  % ... its Diebold-Mariano statistic (negative: full more accurate)
\providecommand{\rhtaPEnetAllCtBar}{}\renewcommand{\rhtaPEnetAllCtBar}{\ensuremath{p < 0.001}}  % ... its p statement
\providecommand{\rhtaQHarEnetAllCtBar}{}\renewcommand{\rhtaQHarEnetAllCtBar}{0.1922}  % ... mean QLIKE, HAR benchmark
\providecommand{\rhtaQFullEnetAllCtBar}{}\renewcommand{\rhtaQFullEnetAllCtBar}{0.1817}  % ... mean QLIKE, full set
\providecommand{\rhtaPctRidgeRthCtClose}{}\renewcommand{\rhtaPctRidgeRthCtClose}{0.8}  % QLIKE % change, full set vs HAR benchmark: Ridge, RTH13@16 (RTH13, penalties chosen on the 16:00 validation rows; linear only), 16:00 bar
\providecommand{\rhtaDmRidgeRthCtClose}{}\renewcommand{\rhtaDmRidgeRthCtClose}{0.31}  % ... its Diebold-Mariano statistic (negative: full more accurate)
\providecommand{\rhtaPRidgeRthCtClose}{}\renewcommand{\rhtaPRidgeRthCtClose}{\ensuremath{p = 0.76}}  % ... its p statement
\providecommand{\rhtaQHarRidgeRthCtClose}{}\renewcommand{\rhtaQHarRidgeRthCtClose}{0.1269}  % ... mean QLIKE, HAR benchmark
\providecommand{\rhtaQFullRidgeRthCtClose}{}\renewcommand{\rhtaQFullRidgeRthCtClose}{0.1279}  % ... mean QLIKE, full set
\providecommand{\rhtaPctLassoRthCtClose}{}\renewcommand{\rhtaPctLassoRthCtClose}{\ensuremath{-}1.8}  % QLIKE % change, full set vs HAR benchmark: Lasso, RTH13@16 (RTH13, penalties chosen on the 16:00 validation rows; linear only), 16:00 bar
\providecommand{\rhtaDmLassoRthCtClose}{}\renewcommand{\rhtaDmLassoRthCtClose}{\ensuremath{-}0.87}  % ... its Diebold-Mariano statistic (negative: full more accurate)
\providecommand{\rhtaPLassoRthCtClose}{}\renewcommand{\rhtaPLassoRthCtClose}{\ensuremath{p = 0.38}}  % ... its p statement
\providecommand{\rhtaQHarLassoRthCtClose}{}\renewcommand{\rhtaQHarLassoRthCtClose}{0.1262}  % ... mean QLIKE, HAR benchmark
\providecommand{\rhtaQFullLassoRthCtClose}{}\renewcommand{\rhtaQFullLassoRthCtClose}{0.1239}  % ... mean QLIKE, full set
\providecommand{\rhtaPctEnetRthCtClose}{}\renewcommand{\rhtaPctEnetRthCtClose}{1.5}  % QLIKE % change, full set vs HAR benchmark: Elastic net, RTH13@16 (RTH13, penalties chosen on the 16:00 validation rows; linear only), 16:00 bar
\providecommand{\rhtaDmEnetRthCtClose}{}\renewcommand{\rhtaDmEnetRthCtClose}{0.59}  % ... its Diebold-Mariano statistic (negative: full more accurate)
\providecommand{\rhtaPEnetRthCtClose}{}\renewcommand{\rhtaPEnetRthCtClose}{\ensuremath{p = 0.56}}  % ... its p statement
\providecommand{\rhtaQHarEnetRthCtClose}{}\renewcommand{\rhtaQHarEnetRthCtClose}{0.1243}  % ... mean QLIKE, HAR benchmark
\providecommand{\rhtaQFullEnetRthCtClose}{}\renewcommand{\rhtaQFullEnetRthCtClose}{0.1262}  % ... mean QLIKE, full set
\providecommand{\rhtaPctRidgeRthCtBar}{}\renewcommand{\rhtaPctRidgeRthCtBar}{\ensuremath{-}5.4}  % QLIKE % change, full set vs HAR benchmark: Ridge, RTH13@16 (RTH13, penalties chosen on the 16:00 validation rows; linear only), every bar
\providecommand{\rhtaDmRidgeRthCtBar}{}\renewcommand{\rhtaDmRidgeRthCtBar}{\ensuremath{-}5.63}  % ... its Diebold-Mariano statistic (negative: full more accurate)
\providecommand{\rhtaPRidgeRthCtBar}{}\renewcommand{\rhtaPRidgeRthCtBar}{\ensuremath{p < 0.001}}  % ... its p statement
\providecommand{\rhtaQHarRidgeRthCtBar}{}\renewcommand{\rhtaQHarRidgeRthCtBar}{0.1893}  % ... mean QLIKE, HAR benchmark
\providecommand{\rhtaQFullRidgeRthCtBar}{}\renewcommand{\rhtaQFullRidgeRthCtBar}{0.1791}  % ... mean QLIKE, full set
\providecommand{\rhtaPctLassoRthCtBar}{}\renewcommand{\rhtaPctLassoRthCtBar}{\ensuremath{-}6.2}  % QLIKE % change, full set vs HAR benchmark: Lasso, RTH13@16 (RTH13, penalties chosen on the 16:00 validation rows; linear only), every bar
\providecommand{\rhtaDmLassoRthCtBar}{}\renewcommand{\rhtaDmLassoRthCtBar}{\ensuremath{-}8.49}  % ... its Diebold-Mariano statistic (negative: full more accurate)
\providecommand{\rhtaPLassoRthCtBar}{}\renewcommand{\rhtaPLassoRthCtBar}{\ensuremath{p < 0.001}}  % ... its p statement
\providecommand{\rhtaQHarLassoRthCtBar}{}\renewcommand{\rhtaQHarLassoRthCtBar}{0.1896}  % ... mean QLIKE, HAR benchmark
\providecommand{\rhtaQFullLassoRthCtBar}{}\renewcommand{\rhtaQFullLassoRthCtBar}{0.1777}  % ... mean QLIKE, full set
\providecommand{\rhtaPctEnetRthCtBar}{}\renewcommand{\rhtaPctEnetRthCtBar}{\ensuremath{-}6.9}  % QLIKE % change, full set vs HAR benchmark: Elastic net, RTH13@16 (RTH13, penalties chosen on the 16:00 validation rows; linear only), every bar
\providecommand{\rhtaDmEnetRthCtBar}{}\renewcommand{\rhtaDmEnetRthCtBar}{\ensuremath{-}7.94}  % ... its Diebold-Mariano statistic (negative: full more accurate)
\providecommand{\rhtaPEnetRthCtBar}{}\renewcommand{\rhtaPEnetRthCtBar}{\ensuremath{p < 0.001}}  % ... its p statement
\providecommand{\rhtaQHarEnetRthCtBar}{}\renewcommand{\rhtaQHarEnetRthCtBar}{0.1911}  % ... mean QLIKE, HAR benchmark
\providecommand{\rhtaQFullEnetRthCtBar}{}\renewcommand{\rhtaQFullEnetRthCtBar}{0.1778}  % ... mean QLIKE, full set
\providecommand{\rhtaPctRidgeAllThreeKClose}{}\renewcommand{\rhtaPctRidgeAllThreeKClose}{\ensuremath{-}4.7}  % QLIKE % change, full set vs HAR benchmark: Ridge, ALL48_N3k (ALL48, the last 3,000 rows), 16:00 bar
\providecommand{\rhtaDmRidgeAllThreeKClose}{}\renewcommand{\rhtaDmRidgeAllThreeKClose}{\ensuremath{-}4.44}  % ... its Diebold-Mariano statistic (negative: full more accurate)
\providecommand{\rhtaPRidgeAllThreeKClose}{}\renewcommand{\rhtaPRidgeAllThreeKClose}{\ensuremath{p < 0.001}}  % ... its p statement
\providecommand{\rhtaQHarRidgeAllThreeKClose}{}\renewcommand{\rhtaQHarRidgeAllThreeKClose}{0.1341}  % ... mean QLIKE, HAR benchmark
\providecommand{\rhtaQFullRidgeAllThreeKClose}{}\renewcommand{\rhtaQFullRidgeAllThreeKClose}{0.1278}  % ... mean QLIKE, full set
\providecommand{\rhtaPctLassoAllThreeKClose}{}\renewcommand{\rhtaPctLassoAllThreeKClose}{\ensuremath{-}1.6}  % QLIKE % change, full set vs HAR benchmark: Lasso, ALL48_N3k (ALL48, the last 3,000 rows), 16:00 bar
\providecommand{\rhtaDmLassoAllThreeKClose}{}\renewcommand{\rhtaDmLassoAllThreeKClose}{\ensuremath{-}1.25}  % ... its Diebold-Mariano statistic (negative: full more accurate)
\providecommand{\rhtaPLassoAllThreeKClose}{}\renewcommand{\rhtaPLassoAllThreeKClose}{\ensuremath{p = 0.21}}  % ... its p statement
\providecommand{\rhtaQHarLassoAllThreeKClose}{}\renewcommand{\rhtaQHarLassoAllThreeKClose}{0.1322}  % ... mean QLIKE, HAR benchmark
\providecommand{\rhtaQFullLassoAllThreeKClose}{}\renewcommand{\rhtaQFullLassoAllThreeKClose}{0.1301}  % ... mean QLIKE, full set
\providecommand{\rhtaPctEnetAllThreeKClose}{}\renewcommand{\rhtaPctEnetAllThreeKClose}{0.5}  % QLIKE % change, full set vs HAR benchmark: Elastic net, ALL48_N3k (ALL48, the last 3,000 rows), 16:00 bar
\providecommand{\rhtaDmEnetAllThreeKClose}{}\renewcommand{\rhtaDmEnetAllThreeKClose}{0.27}  % ... its Diebold-Mariano statistic (negative: full more accurate)
\providecommand{\rhtaPEnetAllThreeKClose}{}\renewcommand{\rhtaPEnetAllThreeKClose}{\ensuremath{p = 0.78}}  % ... its p statement
\providecommand{\rhtaQHarEnetAllThreeKClose}{}\renewcommand{\rhtaQHarEnetAllThreeKClose}{0.1331}  % ... mean QLIKE, HAR benchmark
\providecommand{\rhtaQFullEnetAllThreeKClose}{}\renewcommand{\rhtaQFullEnetAllThreeKClose}{0.1338}  % ... mean QLIKE, full set
\providecommand{\rhtaPctLgbmAllThreeKClose}{}\renewcommand{\rhtaPctLgbmAllThreeKClose}{\ensuremath{-}8.0}  % QLIKE % change, full set vs HAR benchmark: LightGBM, ALL48_N3k (ALL48, the last 3,000 rows), 16:00 bar
\providecommand{\rhtaDmLgbmAllThreeKClose}{}\renewcommand{\rhtaDmLgbmAllThreeKClose}{\ensuremath{-}3.83}  % ... its Diebold-Mariano statistic (negative: full more accurate)
\providecommand{\rhtaPLgbmAllThreeKClose}{}\renewcommand{\rhtaPLgbmAllThreeKClose}{\ensuremath{p < 0.001}}  % ... its p statement
\providecommand{\rhtaQHarLgbmAllThreeKClose}{}\renewcommand{\rhtaQHarLgbmAllThreeKClose}{0.1409}  % ... mean QLIKE, HAR benchmark
\providecommand{\rhtaQFullLgbmAllThreeKClose}{}\renewcommand{\rhtaQFullLgbmAllThreeKClose}{0.1295}  % ... mean QLIKE, full set
\providecommand{\rhtaPctXgbAllThreeKClose}{}\renewcommand{\rhtaPctXgbAllThreeKClose}{\ensuremath{-}6.6}  % QLIKE % change, full set vs HAR benchmark: XGBoost, ALL48_N3k (ALL48, the last 3,000 rows), 16:00 bar
\providecommand{\rhtaDmXgbAllThreeKClose}{}\renewcommand{\rhtaDmXgbAllThreeKClose}{\ensuremath{-}4.14}  % ... its Diebold-Mariano statistic (negative: full more accurate)
\providecommand{\rhtaPXgbAllThreeKClose}{}\renewcommand{\rhtaPXgbAllThreeKClose}{\ensuremath{p < 0.001}}  % ... its p statement
\providecommand{\rhtaQHarXgbAllThreeKClose}{}\renewcommand{\rhtaQHarXgbAllThreeKClose}{0.1348}  % ... mean QLIKE, HAR benchmark
\providecommand{\rhtaQFullXgbAllThreeKClose}{}\renewcommand{\rhtaQFullXgbAllThreeKClose}{0.1260}  % ... mean QLIKE, full set
\providecommand{\rhtaPctRfAllThreeKClose}{}\renewcommand{\rhtaPctRfAllThreeKClose}{\ensuremath{-}6.9}  % QLIKE % change, full set vs HAR benchmark: Random forest, ALL48_N3k (ALL48, the last 3,000 rows), 16:00 bar
\providecommand{\rhtaDmRfAllThreeKClose}{}\renewcommand{\rhtaDmRfAllThreeKClose}{\ensuremath{-}6.46}  % ... its Diebold-Mariano statistic (negative: full more accurate)
\providecommand{\rhtaPRfAllThreeKClose}{}\renewcommand{\rhtaPRfAllThreeKClose}{\ensuremath{p < 0.001}}  % ... its p statement
\providecommand{\rhtaQHarRfAllThreeKClose}{}\renewcommand{\rhtaQHarRfAllThreeKClose}{0.1355}  % ... mean QLIKE, HAR benchmark
\providecommand{\rhtaQFullRfAllThreeKClose}{}\renewcommand{\rhtaQFullRfAllThreeKClose}{0.1260}  % ... mean QLIKE, full set
\providecommand{\rhtaPctRidgeAllThreeKBar}{}\renewcommand{\rhtaPctRidgeAllThreeKBar}{\ensuremath{-}2.7}  % QLIKE % change, full set vs HAR benchmark: Ridge, ALL48_N3k (ALL48, the last 3,000 rows), every bar
\providecommand{\rhtaDmRidgeAllThreeKBar}{}\renewcommand{\rhtaDmRidgeAllThreeKBar}{\ensuremath{-}7.15}  % ... its Diebold-Mariano statistic (negative: full more accurate)
\providecommand{\rhtaPRidgeAllThreeKBar}{}\renewcommand{\rhtaPRidgeAllThreeKBar}{\ensuremath{p < 0.001}}  % ... its p statement
\providecommand{\rhtaQHarRidgeAllThreeKBar}{}\renewcommand{\rhtaQHarRidgeAllThreeKBar}{0.1976}  % ... mean QLIKE, HAR benchmark
\providecommand{\rhtaQFullRidgeAllThreeKBar}{}\renewcommand{\rhtaQFullRidgeAllThreeKBar}{0.1924}  % ... mean QLIKE, full set
\providecommand{\rhtaPctLassoAllThreeKBar}{}\renewcommand{\rhtaPctLassoAllThreeKBar}{\ensuremath{-}2.1}  % QLIKE % change, full set vs HAR benchmark: Lasso, ALL48_N3k (ALL48, the last 3,000 rows), every bar
\providecommand{\rhtaDmLassoAllThreeKBar}{}\renewcommand{\rhtaDmLassoAllThreeKBar}{\ensuremath{-}4.17}  % ... its Diebold-Mariano statistic (negative: full more accurate)
\providecommand{\rhtaPLassoAllThreeKBar}{}\renewcommand{\rhtaPLassoAllThreeKBar}{\ensuremath{p < 0.001}}  % ... its p statement
\providecommand{\rhtaQHarLassoAllThreeKBar}{}\renewcommand{\rhtaQHarLassoAllThreeKBar}{0.1967}  % ... mean QLIKE, HAR benchmark
\providecommand{\rhtaQFullLassoAllThreeKBar}{}\renewcommand{\rhtaQFullLassoAllThreeKBar}{0.1927}  % ... mean QLIKE, full set
\providecommand{\rhtaPctEnetAllThreeKBar}{}\renewcommand{\rhtaPctEnetAllThreeKBar}{\ensuremath{-}1.8}  % QLIKE % change, full set vs HAR benchmark: Elastic net, ALL48_N3k (ALL48, the last 3,000 rows), every bar
\providecommand{\rhtaDmEnetAllThreeKBar}{}\renewcommand{\rhtaDmEnetAllThreeKBar}{\ensuremath{-}2.57}  % ... its Diebold-Mariano statistic (negative: full more accurate)
\providecommand{\rhtaPEnetAllThreeKBar}{}\renewcommand{\rhtaPEnetAllThreeKBar}{\ensuremath{p = 0.0101}}  % ... its p statement
\providecommand{\rhtaQHarEnetAllThreeKBar}{}\renewcommand{\rhtaQHarEnetAllThreeKBar}{0.1972}  % ... mean QLIKE, HAR benchmark
\providecommand{\rhtaQFullEnetAllThreeKBar}{}\renewcommand{\rhtaQFullEnetAllThreeKBar}{0.1936}  % ... mean QLIKE, full set
\providecommand{\rhtaPctLgbmAllThreeKBar}{}\renewcommand{\rhtaPctLgbmAllThreeKBar}{\ensuremath{-}5.9}  % QLIKE % change, full set vs HAR benchmark: LightGBM, ALL48_N3k (ALL48, the last 3,000 rows), every bar
\providecommand{\rhtaDmLgbmAllThreeKBar}{}\renewcommand{\rhtaDmLgbmAllThreeKBar}{\ensuremath{-}8.86}  % ... its Diebold-Mariano statistic (negative: full more accurate)
\providecommand{\rhtaPLgbmAllThreeKBar}{}\renewcommand{\rhtaPLgbmAllThreeKBar}{\ensuremath{p < 0.001}}  % ... its p statement
\providecommand{\rhtaQHarLgbmAllThreeKBar}{}\renewcommand{\rhtaQHarLgbmAllThreeKBar}{0.2062}  % ... mean QLIKE, HAR benchmark
\providecommand{\rhtaQFullLgbmAllThreeKBar}{}\renewcommand{\rhtaQFullLgbmAllThreeKBar}{0.1941}  % ... mean QLIKE, full set
\providecommand{\rhtaPctXgbAllThreeKBar}{}\renewcommand{\rhtaPctXgbAllThreeKBar}{\ensuremath{-}5.9}  % QLIKE % change, full set vs HAR benchmark: XGBoost, ALL48_N3k (ALL48, the last 3,000 rows), every bar
\providecommand{\rhtaDmXgbAllThreeKBar}{}\renewcommand{\rhtaDmXgbAllThreeKBar}{\ensuremath{-}8.32}  % ... its Diebold-Mariano statistic (negative: full more accurate)
\providecommand{\rhtaPXgbAllThreeKBar}{}\renewcommand{\rhtaPXgbAllThreeKBar}{\ensuremath{p < 0.001}}  % ... its p statement
\providecommand{\rhtaQHarXgbAllThreeKBar}{}\renewcommand{\rhtaQHarXgbAllThreeKBar}{0.2011}  % ... mean QLIKE, HAR benchmark
\providecommand{\rhtaQFullXgbAllThreeKBar}{}\renewcommand{\rhtaQFullXgbAllThreeKBar}{0.1893}  % ... mean QLIKE, full set
\providecommand{\rhtaPctRfAllThreeKBar}{}\renewcommand{\rhtaPctRfAllThreeKBar}{\ensuremath{-}5.0}  % QLIKE % change, full set vs HAR benchmark: Random forest, ALL48_N3k (ALL48, the last 3,000 rows), every bar
\providecommand{\rhtaDmRfAllThreeKBar}{}\renewcommand{\rhtaDmRfAllThreeKBar}{\ensuremath{-}8.83}  % ... its Diebold-Mariano statistic (negative: full more accurate)
\providecommand{\rhtaPRfAllThreeKBar}{}\renewcommand{\rhtaPRfAllThreeKBar}{\ensuremath{p < 0.001}}  % ... its p statement
\providecommand{\rhtaQHarRfAllThreeKBar}{}\renewcommand{\rhtaQHarRfAllThreeKBar}{0.1994}  % ... mean QLIKE, HAR benchmark
\providecommand{\rhtaQFullRfAllThreeKBar}{}\renewcommand{\rhtaQFullRfAllThreeKBar}{0.1894}  % ... mean QLIKE, full set
\providecommand{\rhtaPctRidgeRthThreeKClose}{}\renewcommand{\rhtaPctRidgeRthThreeKClose}{\ensuremath{-}1.9}  % QLIKE % change, full set vs HAR benchmark: Ridge, RTH13_N3k (RTH13, the last 3,000 rows), 16:00 bar
\providecommand{\rhtaDmRidgeRthThreeKClose}{}\renewcommand{\rhtaDmRidgeRthThreeKClose}{\ensuremath{-}1.55}  % ... its Diebold-Mariano statistic (negative: full more accurate)
\providecommand{\rhtaPRidgeRthThreeKClose}{}\renewcommand{\rhtaPRidgeRthThreeKClose}{\ensuremath{p = 0.12}}  % ... its p statement
\providecommand{\rhtaQHarRidgeRthThreeKClose}{}\renewcommand{\rhtaQHarRidgeRthThreeKClose}{0.1418}  % ... mean QLIKE, HAR benchmark
\providecommand{\rhtaQFullRidgeRthThreeKClose}{}\renewcommand{\rhtaQFullRidgeRthThreeKClose}{0.1391}  % ... mean QLIKE, full set
\providecommand{\rhtaPctLassoRthThreeKClose}{}\renewcommand{\rhtaPctLassoRthThreeKClose}{\ensuremath{-}1.9}  % QLIKE % change, full set vs HAR benchmark: Lasso, RTH13_N3k (RTH13, the last 3,000 rows), 16:00 bar
\providecommand{\rhtaDmLassoRthThreeKClose}{}\renewcommand{\rhtaDmLassoRthThreeKClose}{\ensuremath{-}1.29}  % ... its Diebold-Mariano statistic (negative: full more accurate)
\providecommand{\rhtaPLassoRthThreeKClose}{}\renewcommand{\rhtaPLassoRthThreeKClose}{\ensuremath{p = 0.20}}  % ... its p statement
\providecommand{\rhtaQHarLassoRthThreeKClose}{}\renewcommand{\rhtaQHarLassoRthThreeKClose}{0.1418}  % ... mean QLIKE, HAR benchmark
\providecommand{\rhtaQFullLassoRthThreeKClose}{}\renewcommand{\rhtaQFullLassoRthThreeKClose}{0.1392}  % ... mean QLIKE, full set
\providecommand{\rhtaPctEnetRthThreeKClose}{}\renewcommand{\rhtaPctEnetRthThreeKClose}{\ensuremath{-}1.6}  % QLIKE % change, full set vs HAR benchmark: Elastic net, RTH13_N3k (RTH13, the last 3,000 rows), 16:00 bar
\providecommand{\rhtaDmEnetRthThreeKClose}{}\renewcommand{\rhtaDmEnetRthThreeKClose}{\ensuremath{-}0.85}  % ... its Diebold-Mariano statistic (negative: full more accurate)
\providecommand{\rhtaPEnetRthThreeKClose}{}\renewcommand{\rhtaPEnetRthThreeKClose}{\ensuremath{p = 0.39}}  % ... its p statement
\providecommand{\rhtaQHarEnetRthThreeKClose}{}\renewcommand{\rhtaQHarEnetRthThreeKClose}{0.1418}  % ... mean QLIKE, HAR benchmark
\providecommand{\rhtaQFullEnetRthThreeKClose}{}\renewcommand{\rhtaQFullEnetRthThreeKClose}{0.1396}  % ... mean QLIKE, full set
\providecommand{\rhtaPctLgbmRthThreeKClose}{}\renewcommand{\rhtaPctLgbmRthThreeKClose}{\ensuremath{-}3.7}  % QLIKE % change, full set vs HAR benchmark: LightGBM, RTH13_N3k (RTH13, the last 3,000 rows), 16:00 bar
\providecommand{\rhtaDmLgbmRthThreeKClose}{}\renewcommand{\rhtaDmLgbmRthThreeKClose}{\ensuremath{-}2.15}  % ... its Diebold-Mariano statistic (negative: full more accurate)
\providecommand{\rhtaPLgbmRthThreeKClose}{}\renewcommand{\rhtaPLgbmRthThreeKClose}{\ensuremath{p = 0.031}}  % ... its p statement
\providecommand{\rhtaQHarLgbmRthThreeKClose}{}\renewcommand{\rhtaQHarLgbmRthThreeKClose}{0.1401}  % ... mean QLIKE, HAR benchmark
\providecommand{\rhtaQFullLgbmRthThreeKClose}{}\renewcommand{\rhtaQFullLgbmRthThreeKClose}{0.1349}  % ... mean QLIKE, full set
\providecommand{\rhtaPctXgbRthThreeKClose}{}\renewcommand{\rhtaPctXgbRthThreeKClose}{\ensuremath{-}2.3}  % QLIKE % change, full set vs HAR benchmark: XGBoost, RTH13_N3k (RTH13, the last 3,000 rows), 16:00 bar
\providecommand{\rhtaDmXgbRthThreeKClose}{}\renewcommand{\rhtaDmXgbRthThreeKClose}{\ensuremath{-}1.26}  % ... its Diebold-Mariano statistic (negative: full more accurate)
\providecommand{\rhtaPXgbRthThreeKClose}{}\renewcommand{\rhtaPXgbRthThreeKClose}{\ensuremath{p = 0.21}}  % ... its p statement
\providecommand{\rhtaQHarXgbRthThreeKClose}{}\renewcommand{\rhtaQHarXgbRthThreeKClose}{0.1334}  % ... mean QLIKE, HAR benchmark
\providecommand{\rhtaQFullXgbRthThreeKClose}{}\renewcommand{\rhtaQFullXgbRthThreeKClose}{0.1303}  % ... mean QLIKE, full set
\providecommand{\rhtaPctRfRthThreeKClose}{}\renewcommand{\rhtaPctRfRthThreeKClose}{\ensuremath{-}4.5}  % QLIKE % change, full set vs HAR benchmark: Random forest, RTH13_N3k (RTH13, the last 3,000 rows), 16:00 bar
\providecommand{\rhtaDmRfRthThreeKClose}{}\renewcommand{\rhtaDmRfRthThreeKClose}{\ensuremath{-}3.54}  % ... its Diebold-Mariano statistic (negative: full more accurate)
\providecommand{\rhtaPRfRthThreeKClose}{}\renewcommand{\rhtaPRfRthThreeKClose}{\ensuremath{p < 0.001}}  % ... its p statement
\providecommand{\rhtaQHarRfRthThreeKClose}{}\renewcommand{\rhtaQHarRfRthThreeKClose}{0.1400}  % ... mean QLIKE, HAR benchmark
\providecommand{\rhtaQFullRfRthThreeKClose}{}\renewcommand{\rhtaQFullRfRthThreeKClose}{0.1336}  % ... mean QLIKE, full set
\providecommand{\rhtaPctRidgeRthThreeKBar}{}\renewcommand{\rhtaPctRidgeRthThreeKBar}{\ensuremath{-}3.6}  % QLIKE % change, full set vs HAR benchmark: Ridge, RTH13_N3k (RTH13, the last 3,000 rows), every bar
\providecommand{\rhtaDmRidgeRthThreeKBar}{}\renewcommand{\rhtaDmRidgeRthThreeKBar}{\ensuremath{-}9.44}  % ... its Diebold-Mariano statistic (negative: full more accurate)
\providecommand{\rhtaPRidgeRthThreeKBar}{}\renewcommand{\rhtaPRidgeRthThreeKBar}{\ensuremath{p < 0.001}}  % ... its p statement
\providecommand{\rhtaQHarRidgeRthThreeKBar}{}\renewcommand{\rhtaQHarRidgeRthThreeKBar}{0.1910}  % ... mean QLIKE, HAR benchmark
\providecommand{\rhtaQFullRidgeRthThreeKBar}{}\renewcommand{\rhtaQFullRidgeRthThreeKBar}{0.1840}  % ... mean QLIKE, full set
\providecommand{\rhtaPctLassoRthThreeKBar}{}\renewcommand{\rhtaPctLassoRthThreeKBar}{\ensuremath{-}3.8}  % QLIKE % change, full set vs HAR benchmark: Lasso, RTH13_N3k (RTH13, the last 3,000 rows), every bar
\providecommand{\rhtaDmLassoRthThreeKBar}{}\renewcommand{\rhtaDmLassoRthThreeKBar}{\ensuremath{-}7.70}  % ... its Diebold-Mariano statistic (negative: full more accurate)
\providecommand{\rhtaPLassoRthThreeKBar}{}\renewcommand{\rhtaPLassoRthThreeKBar}{\ensuremath{p < 0.001}}  % ... its p statement
\providecommand{\rhtaQHarLassoRthThreeKBar}{}\renewcommand{\rhtaQHarLassoRthThreeKBar}{0.1910}  % ... mean QLIKE, HAR benchmark
\providecommand{\rhtaQFullLassoRthThreeKBar}{}\renewcommand{\rhtaQFullLassoRthThreeKBar}{0.1837}  % ... mean QLIKE, full set
\providecommand{\rhtaPctEnetRthThreeKBar}{}\renewcommand{\rhtaPctEnetRthThreeKBar}{\ensuremath{-}4.2}  % QLIKE % change, full set vs HAR benchmark: Elastic net, RTH13_N3k (RTH13, the last 3,000 rows), every bar
\providecommand{\rhtaDmEnetRthThreeKBar}{}\renewcommand{\rhtaDmEnetRthThreeKBar}{\ensuremath{-}6.30}  % ... its Diebold-Mariano statistic (negative: full more accurate)
\providecommand{\rhtaPEnetRthThreeKBar}{}\renewcommand{\rhtaPEnetRthThreeKBar}{\ensuremath{p < 0.001}}  % ... its p statement
\providecommand{\rhtaQHarEnetRthThreeKBar}{}\renewcommand{\rhtaQHarEnetRthThreeKBar}{0.1910}  % ... mean QLIKE, HAR benchmark
\providecommand{\rhtaQFullEnetRthThreeKBar}{}\renewcommand{\rhtaQFullEnetRthThreeKBar}{0.1830}  % ... mean QLIKE, full set
\providecommand{\rhtaPctLgbmRthThreeKBar}{}\renewcommand{\rhtaPctLgbmRthThreeKBar}{\ensuremath{-}6.0}  % QLIKE % change, full set vs HAR benchmark: LightGBM, RTH13_N3k (RTH13, the last 3,000 rows), every bar
\providecommand{\rhtaDmLgbmRthThreeKBar}{}\renewcommand{\rhtaDmLgbmRthThreeKBar}{\ensuremath{-}7.85}  % ... its Diebold-Mariano statistic (negative: full more accurate)
\providecommand{\rhtaPLgbmRthThreeKBar}{}\renewcommand{\rhtaPLgbmRthThreeKBar}{\ensuremath{p < 0.001}}  % ... its p statement
\providecommand{\rhtaQHarLgbmRthThreeKBar}{}\renewcommand{\rhtaQHarLgbmRthThreeKBar}{0.1972}  % ... mean QLIKE, HAR benchmark
\providecommand{\rhtaQFullLgbmRthThreeKBar}{}\renewcommand{\rhtaQFullLgbmRthThreeKBar}{0.1854}  % ... mean QLIKE, full set
\providecommand{\rhtaPctXgbRthThreeKBar}{}\renewcommand{\rhtaPctXgbRthThreeKBar}{\ensuremath{-}6.8}  % QLIKE % change, full set vs HAR benchmark: XGBoost, RTH13_N3k (RTH13, the last 3,000 rows), every bar
\providecommand{\rhtaDmXgbRthThreeKBar}{}\renewcommand{\rhtaDmXgbRthThreeKBar}{\ensuremath{-}9.13}  % ... its Diebold-Mariano statistic (negative: full more accurate)
\providecommand{\rhtaPXgbRthThreeKBar}{}\renewcommand{\rhtaPXgbRthThreeKBar}{\ensuremath{p < 0.001}}  % ... its p statement
\providecommand{\rhtaQHarXgbRthThreeKBar}{}\renewcommand{\rhtaQHarXgbRthThreeKBar}{0.1931}  % ... mean QLIKE, HAR benchmark
\providecommand{\rhtaQFullXgbRthThreeKBar}{}\renewcommand{\rhtaQFullXgbRthThreeKBar}{0.1801}  % ... mean QLIKE, full set
\providecommand{\rhtaPctRfRthThreeKBar}{}\renewcommand{\rhtaPctRfRthThreeKBar}{\ensuremath{-}4.7}  % QLIKE % change, full set vs HAR benchmark: Random forest, RTH13_N3k (RTH13, the last 3,000 rows), every bar
\providecommand{\rhtaDmRfRthThreeKBar}{}\renewcommand{\rhtaDmRfRthThreeKBar}{\ensuremath{-}8.98}  % ... its Diebold-Mariano statistic (negative: full more accurate)
\providecommand{\rhtaPRfRthThreeKBar}{}\renewcommand{\rhtaPRfRthThreeKBar}{\ensuremath{p < 0.001}}  % ... its p statement
\providecommand{\rhtaQHarRfRthThreeKBar}{}\renewcommand{\rhtaQHarRfRthThreeKBar}{0.1935}  % ... mean QLIKE, HAR benchmark
\providecommand{\rhtaQFullRfRthThreeKBar}{}\renewcommand{\rhtaQFullRfRthThreeKBar}{0.1844}  % ... mean QLIKE, full set
\providecommand{\rhtaPctRidgeAllThirtyNineKClose}{}\renewcommand{\rhtaPctRidgeAllThirtyNineKClose}{\ensuremath{-}0.9}  % QLIKE % change, full set vs HAR benchmark: Ridge, ALL48_N39k (ALL48, the last 39,000 rows), 16:00 bar
\providecommand{\rhtaDmRidgeAllThirtyNineKClose}{}\renewcommand{\rhtaDmRidgeAllThirtyNineKClose}{\ensuremath{-}0.41}  % ... its Diebold-Mariano statistic (negative: full more accurate)
\providecommand{\rhtaPRidgeAllThirtyNineKClose}{}\renewcommand{\rhtaPRidgeAllThirtyNineKClose}{\ensuremath{p = 0.68}}  % ... its p statement
\providecommand{\rhtaQHarRidgeAllThirtyNineKClose}{}\renewcommand{\rhtaQHarRidgeAllThirtyNineKClose}{0.1296}  % ... mean QLIKE, HAR benchmark
\providecommand{\rhtaQFullRidgeAllThirtyNineKClose}{}\renewcommand{\rhtaQFullRidgeAllThirtyNineKClose}{0.1285}  % ... mean QLIKE, full set
\providecommand{\rhtaPctLassoAllThirtyNineKClose}{}\renewcommand{\rhtaPctLassoAllThirtyNineKClose}{\ensuremath{-}3.6}  % QLIKE % change, full set vs HAR benchmark: Lasso, ALL48_N39k (ALL48, the last 39,000 rows), 16:00 bar
\providecommand{\rhtaDmLassoAllThirtyNineKClose}{}\renewcommand{\rhtaDmLassoAllThirtyNineKClose}{\ensuremath{-}2.19}  % ... its Diebold-Mariano statistic (negative: full more accurate)
\providecommand{\rhtaPLassoAllThirtyNineKClose}{}\renewcommand{\rhtaPLassoAllThirtyNineKClose}{\ensuremath{p = 0.028}}  % ... its p statement
\providecommand{\rhtaQHarLassoAllThirtyNineKClose}{}\renewcommand{\rhtaQHarLassoAllThirtyNineKClose}{0.1311}  % ... mean QLIKE, HAR benchmark
\providecommand{\rhtaQFullLassoAllThirtyNineKClose}{}\renewcommand{\rhtaQFullLassoAllThirtyNineKClose}{0.1264}  % ... mean QLIKE, full set
\providecommand{\rhtaPctEnetAllThirtyNineKClose}{}\renewcommand{\rhtaPctEnetAllThirtyNineKClose}{\ensuremath{-}3.7}  % QLIKE % change, full set vs HAR benchmark: Elastic net, ALL48_N39k (ALL48, the last 39,000 rows), 16:00 bar
\providecommand{\rhtaDmEnetAllThirtyNineKClose}{}\renewcommand{\rhtaDmEnetAllThirtyNineKClose}{\ensuremath{-}3.22}  % ... its Diebold-Mariano statistic (negative: full more accurate)
\providecommand{\rhtaPEnetAllThirtyNineKClose}{}\renewcommand{\rhtaPEnetAllThirtyNineKClose}{\ensuremath{p = 0.001}}  % ... its p statement
\providecommand{\rhtaQHarEnetAllThirtyNineKClose}{}\renewcommand{\rhtaQHarEnetAllThirtyNineKClose}{0.1312}  % ... mean QLIKE, HAR benchmark
\providecommand{\rhtaQFullEnetAllThirtyNineKClose}{}\renewcommand{\rhtaQFullEnetAllThirtyNineKClose}{0.1263}  % ... mean QLIKE, full set
\providecommand{\rhtaPctLgbmAllThirtyNineKClose}{}\renewcommand{\rhtaPctLgbmAllThirtyNineKClose}{\ensuremath{-}5.5}  % QLIKE % change, full set vs HAR benchmark: LightGBM, ALL48_N39k (ALL48, the last 39,000 rows), 16:00 bar
\providecommand{\rhtaDmLgbmAllThirtyNineKClose}{}\renewcommand{\rhtaDmLgbmAllThirtyNineKClose}{\ensuremath{-}4.12}  % ... its Diebold-Mariano statistic (negative: full more accurate)
\providecommand{\rhtaPLgbmAllThirtyNineKClose}{}\renewcommand{\rhtaPLgbmAllThirtyNineKClose}{\ensuremath{p < 0.001}}  % ... its p statement
\providecommand{\rhtaQHarLgbmAllThirtyNineKClose}{}\renewcommand{\rhtaQHarLgbmAllThirtyNineKClose}{0.1276}  % ... mean QLIKE, HAR benchmark
\providecommand{\rhtaQFullLgbmAllThirtyNineKClose}{}\renewcommand{\rhtaQFullLgbmAllThirtyNineKClose}{0.1206}  % ... mean QLIKE, full set
\providecommand{\rhtaPctXgbAllThirtyNineKClose}{}\renewcommand{\rhtaPctXgbAllThirtyNineKClose}{\ensuremath{-}4.7}  % QLIKE % change, full set vs HAR benchmark: XGBoost, ALL48_N39k (ALL48, the last 39,000 rows), 16:00 bar
\providecommand{\rhtaDmXgbAllThirtyNineKClose}{}\renewcommand{\rhtaDmXgbAllThirtyNineKClose}{\ensuremath{-}3.66}  % ... its Diebold-Mariano statistic (negative: full more accurate)
\providecommand{\rhtaPXgbAllThirtyNineKClose}{}\renewcommand{\rhtaPXgbAllThirtyNineKClose}{\ensuremath{p < 0.001}}  % ... its p statement
\providecommand{\rhtaQHarXgbAllThirtyNineKClose}{}\renewcommand{\rhtaQHarXgbAllThirtyNineKClose}{0.1267}  % ... mean QLIKE, HAR benchmark
\providecommand{\rhtaQFullXgbAllThirtyNineKClose}{}\renewcommand{\rhtaQFullXgbAllThirtyNineKClose}{0.1208}  % ... mean QLIKE, full set
\providecommand{\rhtaPctRfAllThirtyNineKClose}{}\renewcommand{\rhtaPctRfAllThirtyNineKClose}{\ensuremath{-}4.9}  % QLIKE % change, full set vs HAR benchmark: Random forest, ALL48_N39k (ALL48, the last 39,000 rows), 16:00 bar
\providecommand{\rhtaDmRfAllThirtyNineKClose}{}\renewcommand{\rhtaDmRfAllThirtyNineKClose}{\ensuremath{-}3.89}  % ... its Diebold-Mariano statistic (negative: full more accurate)
\providecommand{\rhtaPRfAllThirtyNineKClose}{}\renewcommand{\rhtaPRfAllThirtyNineKClose}{\ensuremath{p < 0.001}}  % ... its p statement
\providecommand{\rhtaQHarRfAllThirtyNineKClose}{}\renewcommand{\rhtaQHarRfAllThirtyNineKClose}{0.1277}  % ... mean QLIKE, HAR benchmark
\providecommand{\rhtaQFullRfAllThirtyNineKClose}{}\renewcommand{\rhtaQFullRfAllThirtyNineKClose}{0.1214}  % ... mean QLIKE, full set
\providecommand{\rhtaPctRidgeAllThirtyNineKBar}{}\renewcommand{\rhtaPctRidgeAllThirtyNineKBar}{\ensuremath{-}4.4}  % QLIKE % change, full set vs HAR benchmark: Ridge, ALL48_N39k (ALL48, the last 39,000 rows), every bar
\providecommand{\rhtaDmRidgeAllThirtyNineKBar}{}\renewcommand{\rhtaDmRidgeAllThirtyNineKBar}{\ensuremath{-}4.51}  % ... its Diebold-Mariano statistic (negative: full more accurate)
\providecommand{\rhtaPRidgeAllThirtyNineKBar}{}\renewcommand{\rhtaPRidgeAllThirtyNineKBar}{\ensuremath{p < 0.001}}  % ... its p statement
\providecommand{\rhtaQHarRidgeAllThirtyNineKBar}{}\renewcommand{\rhtaQHarRidgeAllThirtyNineKBar}{0.1925}  % ... mean QLIKE, HAR benchmark
\providecommand{\rhtaQFullRidgeAllThirtyNineKBar}{}\renewcommand{\rhtaQFullRidgeAllThirtyNineKBar}{0.1840}  % ... mean QLIKE, full set
\providecommand{\rhtaPctLassoAllThirtyNineKBar}{}\renewcommand{\rhtaPctLassoAllThirtyNineKBar}{\ensuremath{-}3.8}  % QLIKE % change, full set vs HAR benchmark: Lasso, ALL48_N39k (ALL48, the last 39,000 rows), every bar
\providecommand{\rhtaDmLassoAllThirtyNineKBar}{}\renewcommand{\rhtaDmLassoAllThirtyNineKBar}{\ensuremath{-}7.03}  % ... its Diebold-Mariano statistic (negative: full more accurate)
\providecommand{\rhtaPLassoAllThirtyNineKBar}{}\renewcommand{\rhtaPLassoAllThirtyNineKBar}{\ensuremath{p < 0.001}}  % ... its p statement
\providecommand{\rhtaQHarLassoAllThirtyNineKBar}{}\renewcommand{\rhtaQHarLassoAllThirtyNineKBar}{0.1922}  % ... mean QLIKE, HAR benchmark
\providecommand{\rhtaQFullLassoAllThirtyNineKBar}{}\renewcommand{\rhtaQFullLassoAllThirtyNineKBar}{0.1849}  % ... mean QLIKE, full set
\providecommand{\rhtaPctEnetAllThirtyNineKBar}{}\renewcommand{\rhtaPctEnetAllThirtyNineKBar}{\ensuremath{-}3.0}  % QLIKE % change, full set vs HAR benchmark: Elastic net, ALL48_N39k (ALL48, the last 39,000 rows), every bar
\providecommand{\rhtaDmEnetAllThirtyNineKBar}{}\renewcommand{\rhtaDmEnetAllThirtyNineKBar}{\ensuremath{-}8.37}  % ... its Diebold-Mariano statistic (negative: full more accurate)
\providecommand{\rhtaPEnetAllThirtyNineKBar}{}\renewcommand{\rhtaPEnetAllThirtyNineKBar}{\ensuremath{p < 0.001}}  % ... its p statement
\providecommand{\rhtaQHarEnetAllThirtyNineKBar}{}\renewcommand{\rhtaQHarEnetAllThirtyNineKBar}{0.1922}  % ... mean QLIKE, HAR benchmark
\providecommand{\rhtaQFullEnetAllThirtyNineKBar}{}\renewcommand{\rhtaQFullEnetAllThirtyNineKBar}{0.1865}  % ... mean QLIKE, full set
\providecommand{\rhtaPctLgbmAllThirtyNineKBar}{}\renewcommand{\rhtaPctLgbmAllThirtyNineKBar}{\ensuremath{-}7.2}  % QLIKE % change, full set vs HAR benchmark: LightGBM, ALL48_N39k (ALL48, the last 39,000 rows), every bar
\providecommand{\rhtaDmLgbmAllThirtyNineKBar}{}\renewcommand{\rhtaDmLgbmAllThirtyNineKBar}{\ensuremath{-}10.73}  % ... its Diebold-Mariano statistic (negative: full more accurate)
\providecommand{\rhtaPLgbmAllThirtyNineKBar}{}\renewcommand{\rhtaPLgbmAllThirtyNineKBar}{\ensuremath{p < 0.001}}  % ... its p statement
\providecommand{\rhtaQHarLgbmAllThirtyNineKBar}{}\renewcommand{\rhtaQHarLgbmAllThirtyNineKBar}{0.1920}  % ... mean QLIKE, HAR benchmark
\providecommand{\rhtaQFullLgbmAllThirtyNineKBar}{}\renewcommand{\rhtaQFullLgbmAllThirtyNineKBar}{0.1781}  % ... mean QLIKE, full set
\providecommand{\rhtaPctXgbAllThirtyNineKBar}{}\renewcommand{\rhtaPctXgbAllThirtyNineKBar}{\ensuremath{-}10.1}  % QLIKE % change, full set vs HAR benchmark: XGBoost, ALL48_N39k (ALL48, the last 39,000 rows), every bar
\providecommand{\rhtaDmXgbAllThirtyNineKBar}{}\renewcommand{\rhtaDmXgbAllThirtyNineKBar}{\ensuremath{-}10.02}  % ... its Diebold-Mariano statistic (negative: full more accurate)
\providecommand{\rhtaPXgbAllThirtyNineKBar}{}\renewcommand{\rhtaPXgbAllThirtyNineKBar}{\ensuremath{p < 0.001}}  % ... its p statement
\providecommand{\rhtaQHarXgbAllThirtyNineKBar}{}\renewcommand{\rhtaQHarXgbAllThirtyNineKBar}{0.1927}  % ... mean QLIKE, HAR benchmark
\providecommand{\rhtaQFullXgbAllThirtyNineKBar}{}\renewcommand{\rhtaQFullXgbAllThirtyNineKBar}{0.1732}  % ... mean QLIKE, full set
\providecommand{\rhtaPctRfAllThirtyNineKBar}{}\renewcommand{\rhtaPctRfAllThirtyNineKBar}{\ensuremath{-}5.6}  % QLIKE % change, full set vs HAR benchmark: Random forest, ALL48_N39k (ALL48, the last 39,000 rows), every bar
\providecommand{\rhtaDmRfAllThirtyNineKBar}{}\renewcommand{\rhtaDmRfAllThirtyNineKBar}{\ensuremath{-}12.63}  % ... its Diebold-Mariano statistic (negative: full more accurate)
\providecommand{\rhtaPRfAllThirtyNineKBar}{}\renewcommand{\rhtaPRfAllThirtyNineKBar}{\ensuremath{p < 0.001}}  % ... its p statement
\providecommand{\rhtaQHarRfAllThirtyNineKBar}{}\renewcommand{\rhtaQHarRfAllThirtyNineKBar}{0.1926}  % ... mean QLIKE, HAR benchmark
\providecommand{\rhtaQFullRfAllThirtyNineKBar}{}\renewcommand{\rhtaQFullRfAllThirtyNineKBar}{0.1818}  % ... mean QLIKE, full set
\providecommand{\rhtaXPctRidgeAllVsRthClose}{}\renewcommand{\rhtaXPctRidgeAllVsRthClose}{\ensuremath{-}3.5}  % QLIKE % change, ALL48 vs RTH13 (full set; negative: ALL48 more accurate), Ridge, 16:00 bar
\providecommand{\rhtaXDmRidgeAllVsRthClose}{}\renewcommand{\rhtaXDmRidgeAllVsRthClose}{\ensuremath{-}2.72}  % ... its Diebold-Mariano statistic
\providecommand{\rhtaXPRidgeAllVsRthClose}{}\renewcommand{\rhtaXPRidgeAllVsRthClose}{\ensuremath{p = 0.006}}  % ... its p statement
\providecommand{\rhtaXPctLassoAllVsRthClose}{}\renewcommand{\rhtaXPctLassoAllVsRthClose}{\ensuremath{-}6.0}  % QLIKE % change, ALL48 vs RTH13 (full set; negative: ALL48 more accurate), Lasso, 16:00 bar
\providecommand{\rhtaXDmLassoAllVsRthClose}{}\renewcommand{\rhtaXDmLassoAllVsRthClose}{\ensuremath{-}3.31}  % ... its Diebold-Mariano statistic
\providecommand{\rhtaXPLassoAllVsRthClose}{}\renewcommand{\rhtaXPLassoAllVsRthClose}{\ensuremath{p < 0.001}}  % ... its p statement
\providecommand{\rhtaXPctEnetAllVsRthClose}{}\renewcommand{\rhtaXPctEnetAllVsRthClose}{\ensuremath{-}2.1}  % QLIKE % change, ALL48 vs RTH13 (full set; negative: ALL48 more accurate), Elastic net, 16:00 bar
\providecommand{\rhtaXDmEnetAllVsRthClose}{}\renewcommand{\rhtaXDmEnetAllVsRthClose}{\ensuremath{-}1.64}  % ... its Diebold-Mariano statistic
\providecommand{\rhtaXPEnetAllVsRthClose}{}\renewcommand{\rhtaXPEnetAllVsRthClose}{\ensuremath{p = 0.1006}}  % ... its p statement
\providecommand{\rhtaXPctLgbmAllVsRthClose}{}\renewcommand{\rhtaXPctLgbmAllVsRthClose}{\ensuremath{-}2.3}  % QLIKE % change, ALL48 vs RTH13 (full set; negative: ALL48 more accurate), LightGBM, 16:00 bar
\providecommand{\rhtaXDmLgbmAllVsRthClose}{}\renewcommand{\rhtaXDmLgbmAllVsRthClose}{\ensuremath{-}2.16}  % ... its Diebold-Mariano statistic
\providecommand{\rhtaXPLgbmAllVsRthClose}{}\renewcommand{\rhtaXPLgbmAllVsRthClose}{\ensuremath{p = 0.030}}  % ... its p statement
\providecommand{\rhtaXPctXgbAllVsRthClose}{}\renewcommand{\rhtaXPctXgbAllVsRthClose}{\ensuremath{-}1.9}  % QLIKE % change, ALL48 vs RTH13 (full set; negative: ALL48 more accurate), XGBoost, 16:00 bar
\providecommand{\rhtaXDmXgbAllVsRthClose}{}\renewcommand{\rhtaXDmXgbAllVsRthClose}{\ensuremath{-}2.21}  % ... its Diebold-Mariano statistic
\providecommand{\rhtaXPXgbAllVsRthClose}{}\renewcommand{\rhtaXPXgbAllVsRthClose}{\ensuremath{p = 0.027}}  % ... its p statement
\providecommand{\rhtaXPctRfAllVsRthClose}{}\renewcommand{\rhtaXPctRfAllVsRthClose}{\ensuremath{-}4.3}  % QLIKE % change, ALL48 vs RTH13 (full set; negative: ALL48 more accurate), Random forest, 16:00 bar
\providecommand{\rhtaXDmRfAllVsRthClose}{}\renewcommand{\rhtaXDmRfAllVsRthClose}{\ensuremath{-}4.26}  % ... its Diebold-Mariano statistic
\providecommand{\rhtaXPRfAllVsRthClose}{}\renewcommand{\rhtaXPRfAllVsRthClose}{\ensuremath{p < 0.001}}  % ... its p statement
\providecommand{\rhtaXPctRidgeAllVsRthBar}{}\renewcommand{\rhtaXPctRidgeAllVsRthBar}{1.9}  % QLIKE % change, ALL48 vs RTH13 (full set; negative: ALL48 more accurate), Ridge, every bar
\providecommand{\rhtaXDmRidgeAllVsRthBar}{}\renewcommand{\rhtaXDmRidgeAllVsRthBar}{3.53}  % ... its Diebold-Mariano statistic
\providecommand{\rhtaXPRidgeAllVsRthBar}{}\renewcommand{\rhtaXPRidgeAllVsRthBar}{\ensuremath{p < 0.001}}  % ... its p statement
\providecommand{\rhtaXPctLassoAllVsRthBar}{}\renewcommand{\rhtaXPctLassoAllVsRthBar}{2.5}  % QLIKE % change, ALL48 vs RTH13 (full set; negative: ALL48 more accurate), Lasso, every bar
\providecommand{\rhtaXDmLassoAllVsRthBar}{}\renewcommand{\rhtaXDmLassoAllVsRthBar}{2.64}  % ... its Diebold-Mariano statistic
\providecommand{\rhtaXPLassoAllVsRthBar}{}\renewcommand{\rhtaXPLassoAllVsRthBar}{\ensuremath{p = 0.008}}  % ... its p statement
\providecommand{\rhtaXPctEnetAllVsRthBar}{}\renewcommand{\rhtaXPctEnetAllVsRthBar}{2.2}  % QLIKE % change, ALL48 vs RTH13 (full set; negative: ALL48 more accurate), Elastic net, every bar
\providecommand{\rhtaXDmEnetAllVsRthBar}{}\renewcommand{\rhtaXDmEnetAllVsRthBar}{4.68}  % ... its Diebold-Mariano statistic
\providecommand{\rhtaXPEnetAllVsRthBar}{}\renewcommand{\rhtaXPEnetAllVsRthBar}{\ensuremath{p < 0.001}}  % ... its p statement
\providecommand{\rhtaXPctLgbmAllVsRthBar}{}\renewcommand{\rhtaXPctLgbmAllVsRthBar}{4.5}  % QLIKE % change, ALL48 vs RTH13 (full set; negative: ALL48 more accurate), LightGBM, every bar
\providecommand{\rhtaXDmLgbmAllVsRthBar}{}\renewcommand{\rhtaXDmLgbmAllVsRthBar}{7.39}  % ... its Diebold-Mariano statistic
\providecommand{\rhtaXPLgbmAllVsRthBar}{}\renewcommand{\rhtaXPLgbmAllVsRthBar}{\ensuremath{p < 0.001}}  % ... its p statement
\providecommand{\rhtaXPctXgbAllVsRthBar}{}\renewcommand{\rhtaXPctXgbAllVsRthBar}{3.2}  % QLIKE % change, ALL48 vs RTH13 (full set; negative: ALL48 more accurate), XGBoost, every bar
\providecommand{\rhtaXDmXgbAllVsRthBar}{}\renewcommand{\rhtaXDmXgbAllVsRthBar}{5.64}  % ... its Diebold-Mariano statistic
\providecommand{\rhtaXPXgbAllVsRthBar}{}\renewcommand{\rhtaXPXgbAllVsRthBar}{\ensuremath{p < 0.001}}  % ... its p statement
\providecommand{\rhtaXPctRfAllVsRthBar}{}\renewcommand{\rhtaXPctRfAllVsRthBar}{5.2}  % QLIKE % change, ALL48 vs RTH13 (full set; negative: ALL48 more accurate), Random forest, every bar
\providecommand{\rhtaXDmRfAllVsRthBar}{}\renewcommand{\rhtaXDmRfAllVsRthBar}{6.61}  % ... its Diebold-Mariano statistic
\providecommand{\rhtaXPRfAllVsRthBar}{}\renewcommand{\rhtaXPRfAllVsRthBar}{\ensuremath{p < 0.001}}  % ... its p statement
\providecommand{\rhtaXPctRidgeClVsRthClose}{}\renewcommand{\rhtaXPctRidgeClVsRthClose}{\ensuremath{-}10.2}  % QLIKE % change, CLOSE vs RTH13 (full set; negative: CLOSE more accurate), Ridge, 16:00 bar
\providecommand{\rhtaXDmRidgeClVsRthClose}{}\renewcommand{\rhtaXDmRidgeClVsRthClose}{\ensuremath{-}3.98}  % ... its Diebold-Mariano statistic
\providecommand{\rhtaXPRidgeClVsRthClose}{}\renewcommand{\rhtaXPRidgeClVsRthClose}{\ensuremath{p < 0.001}}  % ... its p statement
\providecommand{\rhtaXPctLassoClVsRthClose}{}\renewcommand{\rhtaXPctLassoClVsRthClose}{\ensuremath{-}11.8}  % QLIKE % change, CLOSE vs RTH13 (full set; negative: CLOSE more accurate), Lasso, 16:00 bar
\providecommand{\rhtaXDmLassoClVsRthClose}{}\renewcommand{\rhtaXDmLassoClVsRthClose}{\ensuremath{-}3.79}  % ... its Diebold-Mariano statistic
\providecommand{\rhtaXPLassoClVsRthClose}{}\renewcommand{\rhtaXPLassoClVsRthClose}{\ensuremath{p < 0.001}}  % ... its p statement
\providecommand{\rhtaXPctEnetClVsRthClose}{}\renewcommand{\rhtaXPctEnetClVsRthClose}{\ensuremath{-}9.4}  % QLIKE % change, CLOSE vs RTH13 (full set; negative: CLOSE more accurate), Elastic net, 16:00 bar
\providecommand{\rhtaXDmEnetClVsRthClose}{}\renewcommand{\rhtaXDmEnetClVsRthClose}{\ensuremath{-}3.86}  % ... its Diebold-Mariano statistic
\providecommand{\rhtaXPEnetClVsRthClose}{}\renewcommand{\rhtaXPEnetClVsRthClose}{\ensuremath{p < 0.001}}  % ... its p statement
\providecommand{\rhtaXPctLgbmClVsRthClose}{}\renewcommand{\rhtaXPctLgbmClVsRthClose}{1.6}  % QLIKE % change, CLOSE vs RTH13 (full set; negative: CLOSE more accurate), LightGBM, 16:00 bar
\providecommand{\rhtaXDmLgbmClVsRthClose}{}\renewcommand{\rhtaXDmLgbmClVsRthClose}{0.45}  % ... its Diebold-Mariano statistic
\providecommand{\rhtaXPLgbmClVsRthClose}{}\renewcommand{\rhtaXPLgbmClVsRthClose}{\ensuremath{p = 0.65}}  % ... its p statement
\providecommand{\rhtaXPctXgbClVsRthClose}{}\renewcommand{\rhtaXPctXgbClVsRthClose}{\ensuremath{-}1.1}  % QLIKE % change, CLOSE vs RTH13 (full set; negative: CLOSE more accurate), XGBoost, 16:00 bar
\providecommand{\rhtaXDmXgbClVsRthClose}{}\renewcommand{\rhtaXDmXgbClVsRthClose}{\ensuremath{-}0.40}  % ... its Diebold-Mariano statistic
\providecommand{\rhtaXPXgbClVsRthClose}{}\renewcommand{\rhtaXPXgbClVsRthClose}{\ensuremath{p = 0.69}}  % ... its p statement
\providecommand{\rhtaXPctRfClVsRthClose}{}\renewcommand{\rhtaXPctRfClVsRthClose}{\ensuremath{-}2.2}  % QLIKE % change, CLOSE vs RTH13 (full set; negative: CLOSE more accurate), Random forest, 16:00 bar
\providecommand{\rhtaXDmRfClVsRthClose}{}\renewcommand{\rhtaXDmRfClVsRthClose}{\ensuremath{-}0.81}  % ... its Diebold-Mariano statistic
\providecommand{\rhtaXPRfClVsRthClose}{}\renewcommand{\rhtaXPRfClVsRthClose}{\ensuremath{p = 0.42}}  % ... its p statement
\providecommand{\rhtaXPctRidgeClVsAllClose}{}\renewcommand{\rhtaXPctRidgeClVsAllClose}{\ensuremath{-}7.0}  % QLIKE % change, CLOSE vs ALL48 (full set; negative: CLOSE more accurate), Ridge, 16:00 bar
\providecommand{\rhtaXDmRidgeClVsAllClose}{}\renewcommand{\rhtaXDmRidgeClVsAllClose}{\ensuremath{-}3.31}  % ... its Diebold-Mariano statistic
\providecommand{\rhtaXPRidgeClVsAllClose}{}\renewcommand{\rhtaXPRidgeClVsAllClose}{\ensuremath{p < 0.001}}  % ... its p statement
\providecommand{\rhtaXPctLassoClVsAllClose}{}\renewcommand{\rhtaXPctLassoClVsAllClose}{\ensuremath{-}6.2}  % QLIKE % change, CLOSE vs ALL48 (full set; negative: CLOSE more accurate), Lasso, 16:00 bar
\providecommand{\rhtaXDmLassoClVsAllClose}{}\renewcommand{\rhtaXDmLassoClVsAllClose}{\ensuremath{-}3.16}  % ... its Diebold-Mariano statistic
\providecommand{\rhtaXPLassoClVsAllClose}{}\renewcommand{\rhtaXPLassoClVsAllClose}{\ensuremath{p = 0.002}}  % ... its p statement
\providecommand{\rhtaXPctEnetClVsAllClose}{}\renewcommand{\rhtaXPctEnetClVsAllClose}{\ensuremath{-}7.5}  % QLIKE % change, CLOSE vs ALL48 (full set; negative: CLOSE more accurate), Elastic net, 16:00 bar
\providecommand{\rhtaXDmEnetClVsAllClose}{}\renewcommand{\rhtaXDmEnetClVsAllClose}{\ensuremath{-}3.70}  % ... its Diebold-Mariano statistic
\providecommand{\rhtaXPEnetClVsAllClose}{}\renewcommand{\rhtaXPEnetClVsAllClose}{\ensuremath{p < 0.001}}  % ... its p statement
\providecommand{\rhtaXPctLgbmClVsAllClose}{}\renewcommand{\rhtaXPctLgbmClVsAllClose}{4.0}  % QLIKE % change, CLOSE vs ALL48 (full set; negative: CLOSE more accurate), LightGBM, 16:00 bar
\providecommand{\rhtaXDmLgbmClVsAllClose}{}\renewcommand{\rhtaXDmLgbmClVsAllClose}{0.94}  % ... its Diebold-Mariano statistic
\providecommand{\rhtaXPLgbmClVsAllClose}{}\renewcommand{\rhtaXPLgbmClVsAllClose}{\ensuremath{p = 0.35}}  % ... its p statement
\providecommand{\rhtaXPctXgbClVsAllClose}{}\renewcommand{\rhtaXPctXgbClVsAllClose}{0.8}  % QLIKE % change, CLOSE vs ALL48 (full set; negative: CLOSE more accurate), XGBoost, 16:00 bar
\providecommand{\rhtaXDmXgbClVsAllClose}{}\renewcommand{\rhtaXDmXgbClVsAllClose}{0.28}  % ... its Diebold-Mariano statistic
\providecommand{\rhtaXPXgbClVsAllClose}{}\renewcommand{\rhtaXPXgbClVsAllClose}{\ensuremath{p = 0.78}}  % ... its p statement
\providecommand{\rhtaXPctRfClVsAllClose}{}\renewcommand{\rhtaXPctRfClVsAllClose}{2.2}  % QLIKE % change, CLOSE vs ALL48 (full set; negative: CLOSE more accurate), Random forest, 16:00 bar
\providecommand{\rhtaXDmRfClVsAllClose}{}\renewcommand{\rhtaXDmRfClVsAllClose}{0.74}  % ... its Diebold-Mariano statistic
\providecommand{\rhtaXPRfClVsAllClose}{}\renewcommand{\rhtaXPRfClVsAllClose}{\ensuremath{p = 0.46}}  % ... its p statement
\providecommand{\rhtaXPctRidgeClNatVsRthClose}{}\renewcommand{\rhtaXPctRidgeClNatVsRthClose}{\ensuremath{-}10.4}  % QLIKE % change, CLOSE_NAT vs RTH13 (full set; negative: CLOSE_NAT more accurate), Ridge, 16:00 bar
\providecommand{\rhtaXDmRidgeClNatVsRthClose}{}\renewcommand{\rhtaXDmRidgeClNatVsRthClose}{\ensuremath{-}3.91}  % ... its Diebold-Mariano statistic
\providecommand{\rhtaXPRidgeClNatVsRthClose}{}\renewcommand{\rhtaXPRidgeClNatVsRthClose}{\ensuremath{p < 0.001}}  % ... its p statement
\providecommand{\rhtaXPctLassoClNatVsRthClose}{}\renewcommand{\rhtaXPctLassoClNatVsRthClose}{\ensuremath{-}11.5}  % QLIKE % change, CLOSE_NAT vs RTH13 (full set; negative: CLOSE_NAT more accurate), Lasso, 16:00 bar
\providecommand{\rhtaXDmLassoClNatVsRthClose}{}\renewcommand{\rhtaXDmLassoClNatVsRthClose}{\ensuremath{-}3.57}  % ... its Diebold-Mariano statistic
\providecommand{\rhtaXPLassoClNatVsRthClose}{}\renewcommand{\rhtaXPLassoClNatVsRthClose}{\ensuremath{p < 0.001}}  % ... its p statement
\providecommand{\rhtaXPctEnetClNatVsRthClose}{}\renewcommand{\rhtaXPctEnetClNatVsRthClose}{\ensuremath{-}9.3}  % QLIKE % change, CLOSE_NAT vs RTH13 (full set; negative: CLOSE_NAT more accurate), Elastic net, 16:00 bar
\providecommand{\rhtaXDmEnetClNatVsRthClose}{}\renewcommand{\rhtaXDmEnetClNatVsRthClose}{\ensuremath{-}3.58}  % ... its Diebold-Mariano statistic
\providecommand{\rhtaXPEnetClNatVsRthClose}{}\renewcommand{\rhtaXPEnetClNatVsRthClose}{\ensuremath{p < 0.001}}  % ... its p statement
\providecommand{\rhtaXPctRidgeAllCtVsRthClose}{}\renewcommand{\rhtaXPctRidgeAllCtVsRthClose}{\ensuremath{-}3.2}  % QLIKE % change, ALL48@16 vs RTH13 (full set; negative: ALL48@16 more accurate), Ridge, 16:00 bar
\providecommand{\rhtaXDmRidgeAllCtVsRthClose}{}\renewcommand{\rhtaXDmRidgeAllCtVsRthClose}{\ensuremath{-}2.35}  % ... its Diebold-Mariano statistic
\providecommand{\rhtaXPRidgeAllCtVsRthClose}{}\renewcommand{\rhtaXPRidgeAllCtVsRthClose}{\ensuremath{p = 0.019}}  % ... its p statement
\providecommand{\rhtaXPctLassoAllCtVsRthClose}{}\renewcommand{\rhtaXPctLassoAllCtVsRthClose}{\ensuremath{-}5.4}  % QLIKE % change, ALL48@16 vs RTH13 (full set; negative: ALL48@16 more accurate), Lasso, 16:00 bar
\providecommand{\rhtaXDmLassoAllCtVsRthClose}{}\renewcommand{\rhtaXDmLassoAllCtVsRthClose}{\ensuremath{-}3.04}  % ... its Diebold-Mariano statistic
\providecommand{\rhtaXPLassoAllCtVsRthClose}{}\renewcommand{\rhtaXPLassoAllCtVsRthClose}{\ensuremath{p = 0.002}}  % ... its p statement
\providecommand{\rhtaXPctEnetAllCtVsRthClose}{}\renewcommand{\rhtaXPctEnetAllCtVsRthClose}{\ensuremath{-}2.1}  % QLIKE % change, ALL48@16 vs RTH13 (full set; negative: ALL48@16 more accurate), Elastic net, 16:00 bar
\providecommand{\rhtaXDmEnetAllCtVsRthClose}{}\renewcommand{\rhtaXDmEnetAllCtVsRthClose}{\ensuremath{-}1.64}  % ... its Diebold-Mariano statistic
\providecommand{\rhtaXPEnetAllCtVsRthClose}{}\renewcommand{\rhtaXPEnetAllCtVsRthClose}{\ensuremath{p = 0.1006}}  % ... its p statement
\providecommand{\rhtaXPctRidgeRthCtVsRthClose}{}\renewcommand{\rhtaXPctRidgeRthCtVsRthClose}{0.0}  % QLIKE % change, RTH13@16 vs RTH13 (full set; negative: RTH13@16 more accurate), Ridge, 16:00 bar
\providecommand{\rhtaXDmRidgeRthCtVsRthClose}{}\renewcommand{\rhtaXDmRidgeRthCtVsRthClose}{0.00}  % ... its Diebold-Mariano statistic
\providecommand{\rhtaXPRidgeRthCtVsRthClose}{}\renewcommand{\rhtaXPRidgeRthCtVsRthClose}{\ensuremath{p = 1.00}}  % ... its p statement
\providecommand{\rhtaXPctLassoRthCtVsRthClose}{}\renewcommand{\rhtaXPctLassoRthCtVsRthClose}{\ensuremath{-}5.4}  % QLIKE % change, RTH13@16 vs RTH13 (full set; negative: RTH13@16 more accurate), Lasso, 16:00 bar
\providecommand{\rhtaXDmLassoRthCtVsRthClose}{}\renewcommand{\rhtaXDmLassoRthCtVsRthClose}{\ensuremath{-}3.40}  % ... its Diebold-Mariano statistic
\providecommand{\rhtaXPLassoRthCtVsRthClose}{}\renewcommand{\rhtaXPLassoRthCtVsRthClose}{\ensuremath{p < 0.001}}  % ... its p statement
\providecommand{\rhtaXPctEnetRthCtVsRthClose}{}\renewcommand{\rhtaXPctEnetRthCtVsRthClose}{0.0}  % QLIKE % change, RTH13@16 vs RTH13 (full set; negative: RTH13@16 more accurate), Elastic net, 16:00 bar
\providecommand{\rhtaXDmEnetRthCtVsRthClose}{}\renewcommand{\rhtaXDmEnetRthCtVsRthClose}{0.00}  % ... its Diebold-Mariano statistic
\providecommand{\rhtaXPEnetRthCtVsRthClose}{}\renewcommand{\rhtaXPEnetRthCtVsRthClose}{\ensuremath{p = 1.00}}  % ... its p statement
\providecommand{\rhtaXPctRidgeAllCtVsRthBar}{}\renewcommand{\rhtaXPctRidgeAllCtVsRthBar}{1.6}  % QLIKE % change, ALL48@16 vs RTH13 (full set; negative: ALL48@16 more accurate), Ridge, every bar
\providecommand{\rhtaXDmRidgeAllCtVsRthBar}{}\renewcommand{\rhtaXDmRidgeAllCtVsRthBar}{3.24}  % ... its Diebold-Mariano statistic
\providecommand{\rhtaXPRidgeAllCtVsRthBar}{}\renewcommand{\rhtaXPRidgeAllCtVsRthBar}{\ensuremath{p = 0.001}}  % ... its p statement
\providecommand{\rhtaXPctLassoAllCtVsRthBar}{}\renewcommand{\rhtaXPctLassoAllCtVsRthBar}{2.2}  % QLIKE % change, ALL48@16 vs RTH13 (full set; negative: ALL48@16 more accurate), Lasso, every bar
\providecommand{\rhtaXDmLassoAllCtVsRthBar}{}\renewcommand{\rhtaXDmLassoAllCtVsRthBar}{2.65}  % ... its Diebold-Mariano statistic
\providecommand{\rhtaXPLassoAllCtVsRthBar}{}\renewcommand{\rhtaXPLassoAllCtVsRthBar}{\ensuremath{p = 0.008}}  % ... its p statement
\providecommand{\rhtaXPctEnetAllCtVsRthBar}{}\renewcommand{\rhtaXPctEnetAllCtVsRthBar}{2.2}  % QLIKE % change, ALL48@16 vs RTH13 (full set; negative: ALL48@16 more accurate), Elastic net, every bar
\providecommand{\rhtaXDmEnetAllCtVsRthBar}{}\renewcommand{\rhtaXDmEnetAllCtVsRthBar}{4.68}  % ... its Diebold-Mariano statistic
\providecommand{\rhtaXPEnetAllCtVsRthBar}{}\renewcommand{\rhtaXPEnetAllCtVsRthBar}{\ensuremath{p < 0.001}}  % ... its p statement
\providecommand{\rhtaXPctRidgeRthCtVsRthBar}{}\renewcommand{\rhtaXPctRidgeRthCtVsRthBar}{0.0}  % QLIKE % change, RTH13@16 vs RTH13 (full set; negative: RTH13@16 more accurate), Ridge, every bar
\providecommand{\rhtaXDmRidgeRthCtVsRthBar}{}\renewcommand{\rhtaXDmRidgeRthCtVsRthBar}{0.00}  % ... its Diebold-Mariano statistic
\providecommand{\rhtaXPRidgeRthCtVsRthBar}{}\renewcommand{\rhtaXPRidgeRthCtVsRthBar}{\ensuremath{p = 1.00}}  % ... its p statement
\providecommand{\rhtaXPctLassoRthCtVsRthBar}{}\renewcommand{\rhtaXPctLassoRthCtVsRthBar}{0.2}  % QLIKE % change, RTH13@16 vs RTH13 (full set; negative: RTH13@16 more accurate), Lasso, every bar
\providecommand{\rhtaXDmLassoRthCtVsRthBar}{}\renewcommand{\rhtaXDmLassoRthCtVsRthBar}{0.19}  % ... its Diebold-Mariano statistic
\providecommand{\rhtaXPLassoRthCtVsRthBar}{}\renewcommand{\rhtaXPLassoRthCtVsRthBar}{\ensuremath{p = 0.85}}  % ... its p statement
\providecommand{\rhtaXPctEnetRthCtVsRthBar}{}\renewcommand{\rhtaXPctEnetRthCtVsRthBar}{0.0}  % QLIKE % change, RTH13@16 vs RTH13 (full set; negative: RTH13@16 more accurate), Elastic net, every bar
\providecommand{\rhtaXDmEnetRthCtVsRthBar}{}\renewcommand{\rhtaXDmEnetRthCtVsRthBar}{0.00}  % ... its Diebold-Mariano statistic
\providecommand{\rhtaXPEnetRthCtVsRthBar}{}\renewcommand{\rhtaXPEnetRthCtVsRthBar}{\ensuremath{p = 1.00}}  % ... its p statement
\providecommand{\rhtaXPctRidgeAllThreeKVsRthThreeKClose}{}\renewcommand{\rhtaXPctRidgeAllThreeKVsRthThreeKClose}{\ensuremath{-}8.1}  % QLIKE % change, ALL48_N3k vs RTH13_N3k (full set; negative: ALL48_N3k more accurate), Ridge, 16:00 bar
\providecommand{\rhtaXDmRidgeAllThreeKVsRthThreeKClose}{}\renewcommand{\rhtaXDmRidgeAllThreeKVsRthThreeKClose}{\ensuremath{-}5.43}  % ... its Diebold-Mariano statistic
\providecommand{\rhtaXPRidgeAllThreeKVsRthThreeKClose}{}\renewcommand{\rhtaXPRidgeAllThreeKVsRthThreeKClose}{\ensuremath{p < 0.001}}  % ... its p statement
\providecommand{\rhtaXPctLassoAllThreeKVsRthThreeKClose}{}\renewcommand{\rhtaXPctLassoAllThreeKVsRthThreeKClose}{\ensuremath{-}6.5}  % QLIKE % change, ALL48_N3k vs RTH13_N3k (full set; negative: ALL48_N3k more accurate), Lasso, 16:00 bar
\providecommand{\rhtaXDmLassoAllThreeKVsRthThreeKClose}{}\renewcommand{\rhtaXDmLassoAllThreeKVsRthThreeKClose}{\ensuremath{-}3.99}  % ... its Diebold-Mariano statistic
\providecommand{\rhtaXPLassoAllThreeKVsRthThreeKClose}{}\renewcommand{\rhtaXPLassoAllThreeKVsRthThreeKClose}{\ensuremath{p < 0.001}}  % ... its p statement
\providecommand{\rhtaXPctEnetAllThreeKVsRthThreeKClose}{}\renewcommand{\rhtaXPctEnetAllThreeKVsRthThreeKClose}{\ensuremath{-}4.2}  % QLIKE % change, ALL48_N3k vs RTH13_N3k (full set; negative: ALL48_N3k more accurate), Elastic net, 16:00 bar
\providecommand{\rhtaXDmEnetAllThreeKVsRthThreeKClose}{}\renewcommand{\rhtaXDmEnetAllThreeKVsRthThreeKClose}{\ensuremath{-}2.13}  % ... its Diebold-Mariano statistic
\providecommand{\rhtaXPEnetAllThreeKVsRthThreeKClose}{}\renewcommand{\rhtaXPEnetAllThreeKVsRthThreeKClose}{\ensuremath{p = 0.033}}  % ... its p statement
\providecommand{\rhtaXPctLgbmAllThreeKVsRthThreeKClose}{}\renewcommand{\rhtaXPctLgbmAllThreeKVsRthThreeKClose}{\ensuremath{-}3.9}  % QLIKE % change, ALL48_N3k vs RTH13_N3k (full set; negative: ALL48_N3k more accurate), LightGBM, 16:00 bar
\providecommand{\rhtaXDmLgbmAllThreeKVsRthThreeKClose}{}\renewcommand{\rhtaXDmLgbmAllThreeKVsRthThreeKClose}{\ensuremath{-}2.07}  % ... its Diebold-Mariano statistic
\providecommand{\rhtaXPLgbmAllThreeKVsRthThreeKClose}{}\renewcommand{\rhtaXPLgbmAllThreeKVsRthThreeKClose}{\ensuremath{p = 0.038}}  % ... its p statement
\providecommand{\rhtaXPctXgbAllThreeKVsRthThreeKClose}{}\renewcommand{\rhtaXPctXgbAllThreeKVsRthThreeKClose}{\ensuremath{-}3.3}  % QLIKE % change, ALL48_N3k vs RTH13_N3k (full set; negative: ALL48_N3k more accurate), XGBoost, 16:00 bar
\providecommand{\rhtaXDmXgbAllThreeKVsRthThreeKClose}{}\renewcommand{\rhtaXDmXgbAllThreeKVsRthThreeKClose}{\ensuremath{-}1.80}  % ... its Diebold-Mariano statistic
\providecommand{\rhtaXPXgbAllThreeKVsRthThreeKClose}{}\renewcommand{\rhtaXPXgbAllThreeKVsRthThreeKClose}{\ensuremath{p = 0.072}}  % ... its p statement
\providecommand{\rhtaXPctRfAllThreeKVsRthThreeKClose}{}\renewcommand{\rhtaXPctRfAllThreeKVsRthThreeKClose}{\ensuremath{-}5.7}  % QLIKE % change, ALL48_N3k vs RTH13_N3k (full set; negative: ALL48_N3k more accurate), Random forest, 16:00 bar
\providecommand{\rhtaXDmRfAllThreeKVsRthThreeKClose}{}\renewcommand{\rhtaXDmRfAllThreeKVsRthThreeKClose}{\ensuremath{-}3.60}  % ... its Diebold-Mariano statistic
\providecommand{\rhtaXPRfAllThreeKVsRthThreeKClose}{}\renewcommand{\rhtaXPRfAllThreeKVsRthThreeKClose}{\ensuremath{p < 0.001}}  % ... its p statement
\providecommand{\rhtaXPctRidgeAllThreeKVsRthThreeKBar}{}\renewcommand{\rhtaXPctRidgeAllThreeKVsRthThreeKBar}{4.5}  % QLIKE % change, ALL48_N3k vs RTH13_N3k (full set; negative: ALL48_N3k more accurate), Ridge, every bar
\providecommand{\rhtaXDmRidgeAllThreeKVsRthThreeKBar}{}\renewcommand{\rhtaXDmRidgeAllThreeKVsRthThreeKBar}{7.72}  % ... its Diebold-Mariano statistic
\providecommand{\rhtaXPRidgeAllThreeKVsRthThreeKBar}{}\renewcommand{\rhtaXPRidgeAllThreeKVsRthThreeKBar}{\ensuremath{p < 0.001}}  % ... its p statement
\providecommand{\rhtaXPctLassoAllThreeKVsRthThreeKBar}{}\renewcommand{\rhtaXPctLassoAllThreeKVsRthThreeKBar}{4.9}  % QLIKE % change, ALL48_N3k vs RTH13_N3k (full set; negative: ALL48_N3k more accurate), Lasso, every bar
\providecommand{\rhtaXDmLassoAllThreeKVsRthThreeKBar}{}\renewcommand{\rhtaXDmLassoAllThreeKVsRthThreeKBar}{7.96}  % ... its Diebold-Mariano statistic
\providecommand{\rhtaXPLassoAllThreeKVsRthThreeKBar}{}\renewcommand{\rhtaXPLassoAllThreeKVsRthThreeKBar}{\ensuremath{p < 0.001}}  % ... its p statement
\providecommand{\rhtaXPctEnetAllThreeKVsRthThreeKBar}{}\renewcommand{\rhtaXPctEnetAllThreeKVsRthThreeKBar}{5.8}  % QLIKE % change, ALL48_N3k vs RTH13_N3k (full set; negative: ALL48_N3k more accurate), Elastic net, every bar
\providecommand{\rhtaXDmEnetAllThreeKVsRthThreeKBar}{}\renewcommand{\rhtaXDmEnetAllThreeKVsRthThreeKBar}{7.58}  % ... its Diebold-Mariano statistic
\providecommand{\rhtaXPEnetAllThreeKVsRthThreeKBar}{}\renewcommand{\rhtaXPEnetAllThreeKVsRthThreeKBar}{\ensuremath{p < 0.001}}  % ... its p statement
\providecommand{\rhtaXPctLgbmAllThreeKVsRthThreeKBar}{}\renewcommand{\rhtaXPctLgbmAllThreeKVsRthThreeKBar}{4.7}  % QLIKE % change, ALL48_N3k vs RTH13_N3k (full set; negative: ALL48_N3k more accurate), LightGBM, every bar
\providecommand{\rhtaXDmLgbmAllThreeKVsRthThreeKBar}{}\renewcommand{\rhtaXDmLgbmAllThreeKVsRthThreeKBar}{5.00}  % ... its Diebold-Mariano statistic
\providecommand{\rhtaXPLgbmAllThreeKVsRthThreeKBar}{}\renewcommand{\rhtaXPLgbmAllThreeKVsRthThreeKBar}{\ensuremath{p < 0.001}}  % ... its p statement
\providecommand{\rhtaXPctXgbAllThreeKVsRthThreeKBar}{}\renewcommand{\rhtaXPctXgbAllThreeKVsRthThreeKBar}{5.1}  % QLIKE % change, ALL48_N3k vs RTH13_N3k (full set; negative: ALL48_N3k more accurate), XGBoost, every bar
\providecommand{\rhtaXDmXgbAllThreeKVsRthThreeKBar}{}\renewcommand{\rhtaXDmXgbAllThreeKVsRthThreeKBar}{5.92}  % ... its Diebold-Mariano statistic
\providecommand{\rhtaXPXgbAllThreeKVsRthThreeKBar}{}\renewcommand{\rhtaXPXgbAllThreeKVsRthThreeKBar}{\ensuremath{p < 0.001}}  % ... its p statement
\providecommand{\rhtaXPctRfAllThreeKVsRthThreeKBar}{}\renewcommand{\rhtaXPctRfAllThreeKVsRthThreeKBar}{2.7}  % QLIKE % change, ALL48_N3k vs RTH13_N3k (full set; negative: ALL48_N3k more accurate), Random forest, every bar
\providecommand{\rhtaXDmRfAllThreeKVsRthThreeKBar}{}\renewcommand{\rhtaXDmRfAllThreeKVsRthThreeKBar}{3.94}  % ... its Diebold-Mariano statistic
\providecommand{\rhtaXPRfAllThreeKVsRthThreeKBar}{}\renewcommand{\rhtaXPRfAllThreeKVsRthThreeKBar}{\ensuremath{p < 0.001}}  % ... its p statement
\providecommand{\rhtaXPctRidgeClVsRthThreeKClose}{}\renewcommand{\rhtaXPctRidgeClVsRthThreeKClose}{\ensuremath{-}17.4}  % QLIKE % change, CLOSE vs RTH13_N3k (full set; negative: CLOSE more accurate), Ridge, 16:00 bar
\providecommand{\rhtaXDmRidgeClVsRthThreeKClose}{}\renewcommand{\rhtaXDmRidgeClVsRthThreeKClose}{\ensuremath{-}6.84}  % ... its Diebold-Mariano statistic
\providecommand{\rhtaXPRidgeClVsRthThreeKClose}{}\renewcommand{\rhtaXPRidgeClVsRthThreeKClose}{\ensuremath{p < 0.001}}  % ... its p statement
\providecommand{\rhtaXPctLassoClVsRthThreeKClose}{}\renewcommand{\rhtaXPctLassoClVsRthThreeKClose}{\ensuremath{-}16.9}  % QLIKE % change, CLOSE vs RTH13_N3k (full set; negative: CLOSE more accurate), Lasso, 16:00 bar
\providecommand{\rhtaXDmLassoClVsRthThreeKClose}{}\renewcommand{\rhtaXDmLassoClVsRthThreeKClose}{\ensuremath{-}6.57}  % ... its Diebold-Mariano statistic
\providecommand{\rhtaXPLassoClVsRthThreeKClose}{}\renewcommand{\rhtaXPLassoClVsRthThreeKClose}{\ensuremath{p < 0.001}}  % ... its p statement
\providecommand{\rhtaXPctEnetClVsRthThreeKClose}{}\renewcommand{\rhtaXPctEnetClVsRthThreeKClose}{\ensuremath{-}18.0}  % QLIKE % change, CLOSE vs RTH13_N3k (full set; negative: CLOSE more accurate), Elastic net, 16:00 bar
\providecommand{\rhtaXDmEnetClVsRthThreeKClose}{}\renewcommand{\rhtaXDmEnetClVsRthThreeKClose}{\ensuremath{-}6.66}  % ... its Diebold-Mariano statistic
\providecommand{\rhtaXPEnetClVsRthThreeKClose}{}\renewcommand{\rhtaXPEnetClVsRthThreeKClose}{\ensuremath{p < 0.001}}  % ... its p statement
\providecommand{\rhtaXPctLgbmClVsRthThreeKClose}{}\renewcommand{\rhtaXPctLgbmClVsRthThreeKClose}{\ensuremath{-}9.5}  % QLIKE % change, CLOSE vs RTH13_N3k (full set; negative: CLOSE more accurate), LightGBM, 16:00 bar
\providecommand{\rhtaXDmLgbmClVsRthThreeKClose}{}\renewcommand{\rhtaXDmLgbmClVsRthThreeKClose}{\ensuremath{-}2.88}  % ... its Diebold-Mariano statistic
\providecommand{\rhtaXPLgbmClVsRthThreeKClose}{}\renewcommand{\rhtaXPLgbmClVsRthThreeKClose}{\ensuremath{p = 0.004}}  % ... its p statement
\providecommand{\rhtaXPctXgbClVsRthThreeKClose}{}\renewcommand{\rhtaXPctXgbClVsRthThreeKClose}{\ensuremath{-}9.5}  % QLIKE % change, CLOSE vs RTH13_N3k (full set; negative: CLOSE more accurate), XGBoost, 16:00 bar
\providecommand{\rhtaXDmXgbClVsRthThreeKClose}{}\renewcommand{\rhtaXDmXgbClVsRthThreeKClose}{\ensuremath{-}3.40}  % ... its Diebold-Mariano statistic
\providecommand{\rhtaXPXgbClVsRthThreeKClose}{}\renewcommand{\rhtaXPXgbClVsRthThreeKClose}{\ensuremath{p < 0.001}}  % ... its p statement
\providecommand{\rhtaXPctRfClVsRthThreeKClose}{}\renewcommand{\rhtaXPctRfClVsRthThreeKClose}{\ensuremath{-}8.7}  % QLIKE % change, CLOSE vs RTH13_N3k (full set; negative: CLOSE more accurate), Random forest, 16:00 bar
\providecommand{\rhtaXDmRfClVsRthThreeKClose}{}\renewcommand{\rhtaXDmRfClVsRthThreeKClose}{\ensuremath{-}3.08}  % ... its Diebold-Mariano statistic
\providecommand{\rhtaXPRfClVsRthThreeKClose}{}\renewcommand{\rhtaXPRfClVsRthThreeKClose}{\ensuremath{p = 0.002}}  % ... its p statement
\providecommand{\rhtaXPctRidgeClVsAllThreeKClose}{}\renewcommand{\rhtaXPctRidgeClVsAllThreeKClose}{\ensuremath{-}10.2}  % QLIKE % change, CLOSE vs ALL48_N3k (full set; negative: CLOSE more accurate), Ridge, 16:00 bar
\providecommand{\rhtaXDmRidgeClVsAllThreeKClose}{}\renewcommand{\rhtaXDmRidgeClVsAllThreeKClose}{\ensuremath{-}4.63}  % ... its Diebold-Mariano statistic
\providecommand{\rhtaXPRidgeClVsAllThreeKClose}{}\renewcommand{\rhtaXPRidgeClVsAllThreeKClose}{\ensuremath{p < 0.001}}  % ... its p statement
\providecommand{\rhtaXPctLassoClVsAllThreeKClose}{}\renewcommand{\rhtaXPctLassoClVsAllThreeKClose}{\ensuremath{-}11.1}  % QLIKE % change, CLOSE vs ALL48_N3k (full set; negative: CLOSE more accurate), Lasso, 16:00 bar
\providecommand{\rhtaXDmLassoClVsAllThreeKClose}{}\renewcommand{\rhtaXDmLassoClVsAllThreeKClose}{\ensuremath{-}4.62}  % ... its Diebold-Mariano statistic
\providecommand{\rhtaXPLassoClVsAllThreeKClose}{}\renewcommand{\rhtaXPLassoClVsAllThreeKClose}{\ensuremath{p < 0.001}}  % ... its p statement
\providecommand{\rhtaXPctEnetClVsAllThreeKClose}{}\renewcommand{\rhtaXPctEnetClVsAllThreeKClose}{\ensuremath{-}14.5}  % QLIKE % change, CLOSE vs ALL48_N3k (full set; negative: CLOSE more accurate), Elastic net, 16:00 bar
\providecommand{\rhtaXDmEnetClVsAllThreeKClose}{}\renewcommand{\rhtaXDmEnetClVsAllThreeKClose}{\ensuremath{-}5.41}  % ... its Diebold-Mariano statistic
\providecommand{\rhtaXPEnetClVsAllThreeKClose}{}\renewcommand{\rhtaXPEnetClVsAllThreeKClose}{\ensuremath{p < 0.001}}  % ... its p statement
\providecommand{\rhtaXPctLgbmClVsAllThreeKClose}{}\renewcommand{\rhtaXPctLgbmClVsAllThreeKClose}{\ensuremath{-}5.8}  % QLIKE % change, CLOSE vs ALL48_N3k (full set; negative: CLOSE more accurate), LightGBM, 16:00 bar
\providecommand{\rhtaXDmLgbmClVsAllThreeKClose}{}\renewcommand{\rhtaXDmLgbmClVsAllThreeKClose}{\ensuremath{-}1.68}  % ... its Diebold-Mariano statistic
\providecommand{\rhtaXPLgbmClVsAllThreeKClose}{}\renewcommand{\rhtaXPLgbmClVsAllThreeKClose}{\ensuremath{p = 0.093}}  % ... its p statement
\providecommand{\rhtaXPctXgbClVsAllThreeKClose}{}\renewcommand{\rhtaXPctXgbClVsAllThreeKClose}{\ensuremath{-}6.4}  % QLIKE % change, CLOSE vs ALL48_N3k (full set; negative: CLOSE more accurate), XGBoost, 16:00 bar
\providecommand{\rhtaXDmXgbClVsAllThreeKClose}{}\renewcommand{\rhtaXDmXgbClVsAllThreeKClose}{\ensuremath{-}2.23}  % ... its Diebold-Mariano statistic
\providecommand{\rhtaXPXgbClVsAllThreeKClose}{}\renewcommand{\rhtaXPXgbClVsAllThreeKClose}{\ensuremath{p = 0.026}}  % ... its p statement
\providecommand{\rhtaXPctRfClVsAllThreeKClose}{}\renewcommand{\rhtaXPctRfClVsAllThreeKClose}{\ensuremath{-}3.2}  % QLIKE % change, CLOSE vs ALL48_N3k (full set; negative: CLOSE more accurate), Random forest, 16:00 bar
\providecommand{\rhtaXDmRfClVsAllThreeKClose}{}\renewcommand{\rhtaXDmRfClVsAllThreeKClose}{\ensuremath{-}1.10}  % ... its Diebold-Mariano statistic
\providecommand{\rhtaXPRfClVsAllThreeKClose}{}\renewcommand{\rhtaXPRfClVsAllThreeKClose}{\ensuremath{p = 0.27}}  % ... its p statement
\providecommand{\rhtaXPctRidgeAllThirtyNineKVsRthClose}{}\renewcommand{\rhtaXPctRidgeAllThirtyNineKVsRthClose}{0.5}  % QLIKE % change, ALL48_N39k vs RTH13 (full set; negative: ALL48_N39k more accurate), Ridge, 16:00 bar
\providecommand{\rhtaXDmRidgeAllThirtyNineKVsRthClose}{}\renewcommand{\rhtaXDmRidgeAllThirtyNineKVsRthClose}{0.34}  % ... its Diebold-Mariano statistic
\providecommand{\rhtaXPRidgeAllThirtyNineKVsRthClose}{}\renewcommand{\rhtaXPRidgeAllThirtyNineKVsRthClose}{\ensuremath{p = 0.73}}  % ... its p statement
\providecommand{\rhtaXPctLassoAllThirtyNineKVsRthClose}{}\renewcommand{\rhtaXPctLassoAllThirtyNineKVsRthClose}{\ensuremath{-}3.6}  % QLIKE % change, ALL48_N39k vs RTH13 (full set; negative: ALL48_N39k more accurate), Lasso, 16:00 bar
\providecommand{\rhtaXDmLassoAllThirtyNineKVsRthClose}{}\renewcommand{\rhtaXDmLassoAllThirtyNineKVsRthClose}{\ensuremath{-}2.08}  % ... its Diebold-Mariano statistic
\providecommand{\rhtaXPLassoAllThirtyNineKVsRthClose}{}\renewcommand{\rhtaXPLassoAllThirtyNineKVsRthClose}{\ensuremath{p = 0.038}}  % ... its p statement
\providecommand{\rhtaXPctEnetAllThirtyNineKVsRthClose}{}\renewcommand{\rhtaXPctEnetAllThirtyNineKVsRthClose}{0.1}  % QLIKE % change, ALL48_N39k vs RTH13 (full set; negative: ALL48_N39k more accurate), Elastic net, 16:00 bar
\providecommand{\rhtaXDmEnetAllThirtyNineKVsRthClose}{}\renewcommand{\rhtaXDmEnetAllThirtyNineKVsRthClose}{0.07}  % ... its Diebold-Mariano statistic
\providecommand{\rhtaXPEnetAllThirtyNineKVsRthClose}{}\renewcommand{\rhtaXPEnetAllThirtyNineKVsRthClose}{\ensuremath{p = 0.95}}  % ... its p statement
\providecommand{\rhtaXPctLgbmAllThirtyNineKVsRthClose}{}\renewcommand{\rhtaXPctLgbmAllThirtyNineKVsRthClose}{0.4}  % QLIKE % change, ALL48_N39k vs RTH13 (full set; negative: ALL48_N39k more accurate), LightGBM, 16:00 bar
\providecommand{\rhtaXDmLgbmAllThirtyNineKVsRthClose}{}\renewcommand{\rhtaXDmLgbmAllThirtyNineKVsRthClose}{0.36}  % ... its Diebold-Mariano statistic
\providecommand{\rhtaXPLgbmAllThirtyNineKVsRthClose}{}\renewcommand{\rhtaXPLgbmAllThirtyNineKVsRthClose}{\ensuremath{p = 0.72}}  % ... its p statement
\providecommand{\rhtaXPctXgbAllThirtyNineKVsRthClose}{}\renewcommand{\rhtaXPctXgbAllThirtyNineKVsRthClose}{1.3}  % QLIKE % change, ALL48_N39k vs RTH13 (full set; negative: ALL48_N39k more accurate), XGBoost, 16:00 bar
\providecommand{\rhtaXDmXgbAllThirtyNineKVsRthClose}{}\renewcommand{\rhtaXDmXgbAllThirtyNineKVsRthClose}{1.45}  % ... its Diebold-Mariano statistic
\providecommand{\rhtaXPXgbAllThirtyNineKVsRthClose}{}\renewcommand{\rhtaXPXgbAllThirtyNineKVsRthClose}{\ensuremath{p = 0.15}}  % ... its p statement
\providecommand{\rhtaXPctRfAllThirtyNineKVsRthClose}{}\renewcommand{\rhtaXPctRfAllThirtyNineKVsRthClose}{\ensuremath{-}2.7}  % QLIKE % change, ALL48_N39k vs RTH13 (full set; negative: ALL48_N39k more accurate), Random forest, 16:00 bar
\providecommand{\rhtaXDmRfAllThirtyNineKVsRthClose}{}\renewcommand{\rhtaXDmRfAllThirtyNineKVsRthClose}{\ensuremath{-}2.72}  % ... its Diebold-Mariano statistic
\providecommand{\rhtaXPRfAllThirtyNineKVsRthClose}{}\renewcommand{\rhtaXPRfAllThirtyNineKVsRthClose}{\ensuremath{p = 0.007}}  % ... its p statement
\providecommand{\rhtaXPctRidgeAllThirtyNineKVsRthBar}{}\renewcommand{\rhtaXPctRidgeAllThirtyNineKVsRthBar}{2.7}  % QLIKE % change, ALL48_N39k vs RTH13 (full set; negative: ALL48_N39k more accurate), Ridge, every bar
\providecommand{\rhtaXDmRidgeAllThirtyNineKVsRthBar}{}\renewcommand{\rhtaXDmRidgeAllThirtyNineKVsRthBar}{4.54}  % ... its Diebold-Mariano statistic
\providecommand{\rhtaXPRidgeAllThirtyNineKVsRthBar}{}\renewcommand{\rhtaXPRidgeAllThirtyNineKVsRthBar}{\ensuremath{p < 0.001}}  % ... its p statement
\providecommand{\rhtaXPctLassoAllThirtyNineKVsRthBar}{}\renewcommand{\rhtaXPctLassoAllThirtyNineKVsRthBar}{4.2}  % QLIKE % change, ALL48_N39k vs RTH13 (full set; negative: ALL48_N39k more accurate), Lasso, every bar
\providecommand{\rhtaXDmLassoAllThirtyNineKVsRthBar}{}\renewcommand{\rhtaXDmLassoAllThirtyNineKVsRthBar}{3.44}  % ... its Diebold-Mariano statistic
\providecommand{\rhtaXPLassoAllThirtyNineKVsRthBar}{}\renewcommand{\rhtaXPLassoAllThirtyNineKVsRthBar}{\ensuremath{p < 0.001}}  % ... its p statement
\providecommand{\rhtaXPctEnetAllThirtyNineKVsRthBar}{}\renewcommand{\rhtaXPctEnetAllThirtyNineKVsRthBar}{4.9}  % QLIKE % change, ALL48_N39k vs RTH13 (full set; negative: ALL48_N39k more accurate), Elastic net, every bar
\providecommand{\rhtaXDmEnetAllThirtyNineKVsRthBar}{}\renewcommand{\rhtaXDmEnetAllThirtyNineKVsRthBar}{5.04}  % ... its Diebold-Mariano statistic
\providecommand{\rhtaXPEnetAllThirtyNineKVsRthBar}{}\renewcommand{\rhtaXPEnetAllThirtyNineKVsRthBar}{\ensuremath{p < 0.001}}  % ... its p statement
\providecommand{\rhtaXPctLgbmAllThirtyNineKVsRthBar}{}\renewcommand{\rhtaXPctLgbmAllThirtyNineKVsRthBar}{10.4}  % QLIKE % change, ALL48_N39k vs RTH13 (full set; negative: ALL48_N39k more accurate), LightGBM, every bar
\providecommand{\rhtaXDmLgbmAllThirtyNineKVsRthBar}{}\renewcommand{\rhtaXDmLgbmAllThirtyNineKVsRthBar}{8.35}  % ... its Diebold-Mariano statistic
\providecommand{\rhtaXPLgbmAllThirtyNineKVsRthBar}{}\renewcommand{\rhtaXPLgbmAllThirtyNineKVsRthBar}{\ensuremath{p < 0.001}}  % ... its p statement
\providecommand{\rhtaXPctXgbAllThirtyNineKVsRthBar}{}\renewcommand{\rhtaXPctXgbAllThirtyNineKVsRthBar}{6.4}  % QLIKE % change, ALL48_N39k vs RTH13 (full set; negative: ALL48_N39k more accurate), XGBoost, every bar
\providecommand{\rhtaXDmXgbAllThirtyNineKVsRthBar}{}\renewcommand{\rhtaXDmXgbAllThirtyNineKVsRthBar}{8.06}  % ... its Diebold-Mariano statistic
\providecommand{\rhtaXPXgbAllThirtyNineKVsRthBar}{}\renewcommand{\rhtaXPXgbAllThirtyNineKVsRthBar}{\ensuremath{p < 0.001}}  % ... its p statement
\providecommand{\rhtaXPctRfAllThirtyNineKVsRthBar}{}\renewcommand{\rhtaXPctRfAllThirtyNineKVsRthBar}{7.9}  % QLIKE % change, ALL48_N39k vs RTH13 (full set; negative: ALL48_N39k more accurate), Random forest, every bar
\providecommand{\rhtaXDmRfAllThirtyNineKVsRthBar}{}\renewcommand{\rhtaXDmRfAllThirtyNineKVsRthBar}{6.79}  % ... its Diebold-Mariano statistic
\providecommand{\rhtaXPRfAllThirtyNineKVsRthBar}{}\renewcommand{\rhtaXPRfAllThirtyNineKVsRthBar}{\ensuremath{p < 0.001}}  % ... its p statement
\providecommand{\rhtaXPctRidgeAllThreeKVsAllClose}{}\renewcommand{\rhtaXPctRidgeAllThreeKVsAllClose}{3.5}  % QLIKE % change, ALL48_N3k vs ALL48 (full set; negative: ALL48_N3k more accurate), Ridge, 16:00 bar
\providecommand{\rhtaXDmRidgeAllThreeKVsAllClose}{}\renewcommand{\rhtaXDmRidgeAllThreeKVsAllClose}{2.28}  % ... its Diebold-Mariano statistic
\providecommand{\rhtaXPRidgeAllThreeKVsAllClose}{}\renewcommand{\rhtaXPRidgeAllThreeKVsAllClose}{\ensuremath{p = 0.022}}  % ... its p statement
\providecommand{\rhtaXPctLassoAllThreeKVsAllClose}{}\renewcommand{\rhtaXPctLassoAllThreeKVsAllClose}{5.6}  % QLIKE % change, ALL48_N3k vs ALL48 (full set; negative: ALL48_N3k more accurate), Lasso, 16:00 bar
\providecommand{\rhtaXDmLassoAllThreeKVsAllClose}{}\renewcommand{\rhtaXDmLassoAllThreeKVsAllClose}{3.42}  % ... its Diebold-Mariano statistic
\providecommand{\rhtaXPLassoAllThreeKVsAllClose}{}\renewcommand{\rhtaXPLassoAllThreeKVsAllClose}{\ensuremath{p < 0.001}}  % ... its p statement
\providecommand{\rhtaXPctEnetAllThreeKVsAllClose}{}\renewcommand{\rhtaXPctEnetAllThreeKVsAllClose}{8.2}  % QLIKE % change, ALL48_N3k vs ALL48 (full set; negative: ALL48_N3k more accurate), Elastic net, 16:00 bar
\providecommand{\rhtaXDmEnetAllThreeKVsAllClose}{}\renewcommand{\rhtaXDmEnetAllThreeKVsAllClose}{4.01}  % ... its Diebold-Mariano statistic
\providecommand{\rhtaXPEnetAllThreeKVsAllClose}{}\renewcommand{\rhtaXPEnetAllThreeKVsAllClose}{\ensuremath{p < 0.001}}  % ... its p statement
\providecommand{\rhtaXPctLgbmAllThreeKVsAllClose}{}\renewcommand{\rhtaXPctLgbmAllThreeKVsAllClose}{10.4}  % QLIKE % change, ALL48_N3k vs ALL48 (full set; negative: ALL48_N3k more accurate), LightGBM, 16:00 bar
\providecommand{\rhtaXDmLgbmAllThreeKVsAllClose}{}\renewcommand{\rhtaXDmLgbmAllThreeKVsAllClose}{4.29}  % ... its Diebold-Mariano statistic
\providecommand{\rhtaXPLgbmAllThreeKVsAllClose}{}\renewcommand{\rhtaXPLgbmAllThreeKVsAllClose}{\ensuremath{p < 0.001}}  % ... its p statement
\providecommand{\rhtaXPctXgbAllThreeKVsAllClose}{}\renewcommand{\rhtaXPctXgbAllThreeKVsAllClose}{7.7}  % QLIKE % change, ALL48_N3k vs ALL48 (full set; negative: ALL48_N3k more accurate), XGBoost, 16:00 bar
\providecommand{\rhtaXDmXgbAllThreeKVsAllClose}{}\renewcommand{\rhtaXDmXgbAllThreeKVsAllClose}{4.20}  % ... its Diebold-Mariano statistic
\providecommand{\rhtaXPXgbAllThreeKVsAllClose}{}\renewcommand{\rhtaXPXgbAllThreeKVsAllClose}{\ensuremath{p < 0.001}}  % ... its p statement
\providecommand{\rhtaXPctRfAllThreeKVsAllClose}{}\renewcommand{\rhtaXPctRfAllThreeKVsAllClose}{5.6}  % QLIKE % change, ALL48_N3k vs ALL48 (full set; negative: ALL48_N3k more accurate), Random forest, 16:00 bar
\providecommand{\rhtaXDmRfAllThreeKVsAllClose}{}\renewcommand{\rhtaXDmRfAllThreeKVsAllClose}{3.47}  % ... its Diebold-Mariano statistic
\providecommand{\rhtaXPRfAllThreeKVsAllClose}{}\renewcommand{\rhtaXPRfAllThreeKVsAllClose}{\ensuremath{p < 0.001}}  % ... its p statement
\providecommand{\rhtaXPctRidgeAllThreeKVsAllBar}{}\renewcommand{\rhtaXPctRidgeAllThreeKVsAllBar}{5.4}  % QLIKE % change, ALL48_N3k vs ALL48 (full set; negative: ALL48_N3k more accurate), Ridge, every bar
\providecommand{\rhtaXDmRidgeAllThreeKVsAllBar}{}\renewcommand{\rhtaXDmRidgeAllThreeKVsAllBar}{6.67}  % ... its Diebold-Mariano statistic
\providecommand{\rhtaXPRidgeAllThreeKVsAllBar}{}\renewcommand{\rhtaXPRidgeAllThreeKVsAllBar}{\ensuremath{p < 0.001}}  % ... its p statement
\providecommand{\rhtaXPctLassoAllThreeKVsAllBar}{}\renewcommand{\rhtaXPctLassoAllThreeKVsAllBar}{5.9}  % QLIKE % change, ALL48_N3k vs ALL48 (full set; negative: ALL48_N3k more accurate), Lasso, every bar
\providecommand{\rhtaXDmLassoAllThreeKVsAllBar}{}\renewcommand{\rhtaXDmLassoAllThreeKVsAllBar}{7.47}  % ... its Diebold-Mariano statistic
\providecommand{\rhtaXPLassoAllThreeKVsAllBar}{}\renewcommand{\rhtaXPLassoAllThreeKVsAllBar}{\ensuremath{p < 0.001}}  % ... its p statement
\providecommand{\rhtaXPctEnetAllThreeKVsAllBar}{}\renewcommand{\rhtaXPctEnetAllThreeKVsAllBar}{6.6}  % QLIKE % change, ALL48_N3k vs ALL48 (full set; negative: ALL48_N3k more accurate), Elastic net, every bar
\providecommand{\rhtaXDmEnetAllThreeKVsAllBar}{}\renewcommand{\rhtaXDmEnetAllThreeKVsAllBar}{7.48}  % ... its Diebold-Mariano statistic
\providecommand{\rhtaXPEnetAllThreeKVsAllBar}{}\renewcommand{\rhtaXPEnetAllThreeKVsAllBar}{\ensuremath{p < 0.001}}  % ... its p statement
\providecommand{\rhtaXPctLgbmAllThreeKVsAllBar}{}\renewcommand{\rhtaXPctLgbmAllThreeKVsAllBar}{15.2}  % QLIKE % change, ALL48_N3k vs ALL48 (full set; negative: ALL48_N3k more accurate), LightGBM, every bar
\providecommand{\rhtaXDmLgbmAllThreeKVsAllBar}{}\renewcommand{\rhtaXDmLgbmAllThreeKVsAllBar}{8.62}  % ... its Diebold-Mariano statistic
\providecommand{\rhtaXPLgbmAllThreeKVsAllBar}{}\renewcommand{\rhtaXPLgbmAllThreeKVsAllBar}{\ensuremath{p < 0.001}}  % ... its p statement
\providecommand{\rhtaXPctXgbAllThreeKVsAllBar}{}\renewcommand{\rhtaXPctXgbAllThreeKVsAllBar}{12.7}  % QLIKE % change, ALL48_N3k vs ALL48 (full set; negative: ALL48_N3k more accurate), XGBoost, every bar
\providecommand{\rhtaXDmXgbAllThreeKVsAllBar}{}\renewcommand{\rhtaXDmXgbAllThreeKVsAllBar}{8.10}  % ... its Diebold-Mariano statistic
\providecommand{\rhtaXPXgbAllThreeKVsAllBar}{}\renewcommand{\rhtaXPXgbAllThreeKVsAllBar}{\ensuremath{p < 0.001}}  % ... its p statement
\providecommand{\rhtaXPctRfAllThreeKVsAllBar}{}\renewcommand{\rhtaXPctRfAllThreeKVsAllBar}{6.8}  % QLIKE % change, ALL48_N3k vs ALL48 (full set; negative: ALL48_N3k more accurate), Random forest, every bar
\providecommand{\rhtaXDmRfAllThreeKVsAllBar}{}\renewcommand{\rhtaXDmRfAllThreeKVsAllBar}{7.80}  % ... its Diebold-Mariano statistic
\providecommand{\rhtaXPRfAllThreeKVsAllBar}{}\renewcommand{\rhtaXPRfAllThreeKVsAllBar}{\ensuremath{p < 0.001}}  % ... its p statement
\providecommand{\rhtaXPctRidgeRthThreeKVsRthClose}{}\renewcommand{\rhtaXPctRidgeRthThreeKVsRthClose}{8.7}  % QLIKE % change, RTH13_N3k vs RTH13 (full set; negative: RTH13_N3k more accurate), Ridge, 16:00 bar
\providecommand{\rhtaXDmRidgeRthThreeKVsRthClose}{}\renewcommand{\rhtaXDmRidgeRthThreeKVsRthClose}{4.61}  % ... its Diebold-Mariano statistic
\providecommand{\rhtaXPRidgeRthThreeKVsRthClose}{}\renewcommand{\rhtaXPRidgeRthThreeKVsRthClose}{\ensuremath{p < 0.001}}  % ... its p statement
\providecommand{\rhtaXPctLassoRthThreeKVsRthClose}{}\renewcommand{\rhtaXPctLassoRthThreeKVsRthClose}{6.2}  % QLIKE % change, RTH13_N3k vs RTH13 (full set; negative: RTH13_N3k more accurate), Lasso, 16:00 bar
\providecommand{\rhtaXDmLassoRthThreeKVsRthClose}{}\renewcommand{\rhtaXDmLassoRthThreeKVsRthClose}{2.63}  % ... its Diebold-Mariano statistic
\providecommand{\rhtaXPLassoRthThreeKVsRthClose}{}\renewcommand{\rhtaXPLassoRthThreeKVsRthClose}{\ensuremath{p = 0.009}}  % ... its p statement
\providecommand{\rhtaXPctEnetRthThreeKVsRthClose}{}\renewcommand{\rhtaXPctEnetRthThreeKVsRthClose}{10.6}  % QLIKE % change, RTH13_N3k vs RTH13 (full set; negative: RTH13_N3k more accurate), Elastic net, 16:00 bar
\providecommand{\rhtaXDmEnetRthThreeKVsRthClose}{}\renewcommand{\rhtaXDmEnetRthThreeKVsRthClose}{5.13}  % ... its Diebold-Mariano statistic
\providecommand{\rhtaXPEnetRthThreeKVsRthClose}{}\renewcommand{\rhtaXPEnetRthThreeKVsRthClose}{\ensuremath{p < 0.001}}  % ... its p statement
\providecommand{\rhtaXPctLgbmRthThreeKVsRthClose}{}\renewcommand{\rhtaXPctLgbmRthThreeKVsRthClose}{12.3}  % QLIKE % change, RTH13_N3k vs RTH13 (full set; negative: RTH13_N3k more accurate), LightGBM, 16:00 bar
\providecommand{\rhtaXDmLgbmRthThreeKVsRthClose}{}\renewcommand{\rhtaXDmLgbmRthThreeKVsRthClose}{5.35}  % ... its Diebold-Mariano statistic
\providecommand{\rhtaXPLgbmRthThreeKVsRthClose}{}\renewcommand{\rhtaXPLgbmRthThreeKVsRthClose}{\ensuremath{p < 0.001}}  % ... its p statement
\providecommand{\rhtaXPctXgbRthThreeKVsRthClose}{}\renewcommand{\rhtaXPctXgbRthThreeKVsRthClose}{9.3}  % QLIKE % change, RTH13_N3k vs RTH13 (full set; negative: RTH13_N3k more accurate), XGBoost, 16:00 bar
\providecommand{\rhtaXDmXgbRthThreeKVsRthClose}{}\renewcommand{\rhtaXDmXgbRthThreeKVsRthClose}{5.11}  % ... its Diebold-Mariano statistic
\providecommand{\rhtaXPXgbRthThreeKVsRthClose}{}\renewcommand{\rhtaXPXgbRthThreeKVsRthClose}{\ensuremath{p < 0.001}}  % ... its p statement
\providecommand{\rhtaXPctRfRthThreeKVsRthClose}{}\renewcommand{\rhtaXPctRfRthThreeKVsRthClose}{7.1}  % QLIKE % change, RTH13_N3k vs RTH13 (full set; negative: RTH13_N3k more accurate), Random forest, 16:00 bar
\providecommand{\rhtaXDmRfRthThreeKVsRthClose}{}\renewcommand{\rhtaXDmRfRthThreeKVsRthClose}{5.16}  % ... its Diebold-Mariano statistic
\providecommand{\rhtaXPRfRthThreeKVsRthClose}{}\renewcommand{\rhtaXPRfRthThreeKVsRthClose}{\ensuremath{p < 0.001}}  % ... its p statement
\providecommand{\rhtaXPctRidgeRthThreeKVsRthBar}{}\renewcommand{\rhtaXPctRidgeRthThreeKVsRthBar}{2.7}  % QLIKE % change, RTH13_N3k vs RTH13 (full set; negative: RTH13_N3k more accurate), Ridge, every bar
\providecommand{\rhtaXDmRidgeRthThreeKVsRthBar}{}\renewcommand{\rhtaXDmRidgeRthThreeKVsRthBar}{3.01}  % ... its Diebold-Mariano statistic
\providecommand{\rhtaXPRidgeRthThreeKVsRthBar}{}\renewcommand{\rhtaXPRidgeRthThreeKVsRthBar}{\ensuremath{p = 0.003}}  % ... its p statement
\providecommand{\rhtaXPctLassoRthThreeKVsRthBar}{}\renewcommand{\rhtaXPctLassoRthThreeKVsRthBar}{3.5}  % QLIKE % change, RTH13_N3k vs RTH13 (full set; negative: RTH13_N3k more accurate), Lasso, every bar
\providecommand{\rhtaXDmLassoRthThreeKVsRthBar}{}\renewcommand{\rhtaXDmLassoRthThreeKVsRthBar}{2.71}  % ... its Diebold-Mariano statistic
\providecommand{\rhtaXPLassoRthThreeKVsRthBar}{}\renewcommand{\rhtaXPLassoRthThreeKVsRthBar}{\ensuremath{p = 0.007}}  % ... its p statement
\providecommand{\rhtaXPctEnetRthThreeKVsRthBar}{}\renewcommand{\rhtaXPctEnetRthThreeKVsRthBar}{2.9}  % QLIKE % change, RTH13_N3k vs RTH13 (full set; negative: RTH13_N3k more accurate), Elastic net, every bar
\providecommand{\rhtaXDmEnetRthThreeKVsRthBar}{}\renewcommand{\rhtaXDmEnetRthThreeKVsRthBar}{3.54}  % ... its Diebold-Mariano statistic
\providecommand{\rhtaXPEnetRthThreeKVsRthBar}{}\renewcommand{\rhtaXPEnetRthThreeKVsRthBar}{\ensuremath{p < 0.001}}  % ... its p statement
\providecommand{\rhtaXPctLgbmRthThreeKVsRthBar}{}\renewcommand{\rhtaXPctLgbmRthThreeKVsRthBar}{14.9}  % QLIKE % change, RTH13_N3k vs RTH13 (full set; negative: RTH13_N3k more accurate), LightGBM, every bar
\providecommand{\rhtaXDmLgbmRthThreeKVsRthBar}{}\renewcommand{\rhtaXDmLgbmRthThreeKVsRthBar}{8.05}  % ... its Diebold-Mariano statistic
\providecommand{\rhtaXPLgbmRthThreeKVsRthBar}{}\renewcommand{\rhtaXPLgbmRthThreeKVsRthBar}{\ensuremath{p < 0.001}}  % ... its p statement
\providecommand{\rhtaXPctXgbRthThreeKVsRthBar}{}\renewcommand{\rhtaXPctXgbRthThreeKVsRthBar}{10.7}  % QLIKE % change, RTH13_N3k vs RTH13 (full set; negative: RTH13_N3k more accurate), XGBoost, every bar
\providecommand{\rhtaXDmXgbRthThreeKVsRthBar}{}\renewcommand{\rhtaXDmXgbRthThreeKVsRthBar}{7.31}  % ... its Diebold-Mariano statistic
\providecommand{\rhtaXPXgbRthThreeKVsRthBar}{}\renewcommand{\rhtaXPXgbRthThreeKVsRthBar}{\ensuremath{p < 0.001}}  % ... its p statement
\providecommand{\rhtaXPctRfRthThreeKVsRthBar}{}\renewcommand{\rhtaXPctRfRthThreeKVsRthBar}{9.4}  % QLIKE % change, RTH13_N3k vs RTH13 (full set; negative: RTH13_N3k more accurate), Random forest, every bar
\providecommand{\rhtaXDmRfRthThreeKVsRthBar}{}\renewcommand{\rhtaXDmRfRthThreeKVsRthBar}{7.22}  % ... its Diebold-Mariano statistic
\providecommand{\rhtaXPRfRthThreeKVsRthBar}{}\renewcommand{\rhtaXPRfRthThreeKVsRthBar}{\ensuremath{p < 0.001}}  % ... its p statement
\providecommand{\rhtaXPctRidgeAllThirtyNineKVsAllClose}{}\renewcommand{\rhtaXPctRidgeAllThirtyNineKVsAllClose}{4.1}  % QLIKE % change, ALL48_N39k vs ALL48 (full set; negative: ALL48_N39k more accurate), Ridge, 16:00 bar
\providecommand{\rhtaXDmRidgeAllThirtyNineKVsAllClose}{}\renewcommand{\rhtaXDmRidgeAllThirtyNineKVsAllClose}{4.03}  % ... its Diebold-Mariano statistic
\providecommand{\rhtaXPRidgeAllThirtyNineKVsAllClose}{}\renewcommand{\rhtaXPRidgeAllThirtyNineKVsAllClose}{\ensuremath{p < 0.001}}  % ... its p statement
\providecommand{\rhtaXPctLassoAllThirtyNineKVsAllClose}{}\renewcommand{\rhtaXPctLassoAllThirtyNineKVsAllClose}{2.6}  % QLIKE % change, ALL48_N39k vs ALL48 (full set; negative: ALL48_N39k more accurate), Lasso, 16:00 bar
\providecommand{\rhtaXDmLassoAllThirtyNineKVsAllClose}{}\renewcommand{\rhtaXDmLassoAllThirtyNineKVsAllClose}{2.67}  % ... its Diebold-Mariano statistic
\providecommand{\rhtaXPLassoAllThirtyNineKVsAllClose}{}\renewcommand{\rhtaXPLassoAllThirtyNineKVsAllClose}{\ensuremath{p = 0.008}}  % ... its p statement
\providecommand{\rhtaXPctEnetAllThirtyNineKVsAllClose}{}\renewcommand{\rhtaXPctEnetAllThirtyNineKVsAllClose}{2.2}  % QLIKE % change, ALL48_N39k vs ALL48 (full set; negative: ALL48_N39k more accurate), Elastic net, 16:00 bar
\providecommand{\rhtaXDmEnetAllThirtyNineKVsAllClose}{}\renewcommand{\rhtaXDmEnetAllThirtyNineKVsAllClose}{1.89}  % ... its Diebold-Mariano statistic
\providecommand{\rhtaXPEnetAllThirtyNineKVsAllClose}{}\renewcommand{\rhtaXPEnetAllThirtyNineKVsAllClose}{\ensuremath{p = 0.059}}  % ... its p statement
\providecommand{\rhtaXPctLgbmAllThirtyNineKVsAllClose}{}\renewcommand{\rhtaXPctLgbmAllThirtyNineKVsAllClose}{2.7}  % QLIKE % change, ALL48_N39k vs ALL48 (full set; negative: ALL48_N39k more accurate), LightGBM, 16:00 bar
\providecommand{\rhtaXDmLgbmAllThirtyNineKVsAllClose}{}\renewcommand{\rhtaXDmLgbmAllThirtyNineKVsAllClose}{2.87}  % ... its Diebold-Mariano statistic
\providecommand{\rhtaXPLgbmAllThirtyNineKVsAllClose}{}\renewcommand{\rhtaXPLgbmAllThirtyNineKVsAllClose}{\ensuremath{p = 0.004}}  % ... its p statement
\providecommand{\rhtaXPctXgbAllThirtyNineKVsAllClose}{}\renewcommand{\rhtaXPctXgbAllThirtyNineKVsAllClose}{3.3}  % QLIKE % change, ALL48_N39k vs ALL48 (full set; negative: ALL48_N39k more accurate), XGBoost, 16:00 bar
\providecommand{\rhtaXDmXgbAllThirtyNineKVsAllClose}{}\renewcommand{\rhtaXDmXgbAllThirtyNineKVsAllClose}{4.05}  % ... its Diebold-Mariano statistic
\providecommand{\rhtaXPXgbAllThirtyNineKVsAllClose}{}\renewcommand{\rhtaXPXgbAllThirtyNineKVsAllClose}{\ensuremath{p < 0.001}}  % ... its p statement
\providecommand{\rhtaXPctRfAllThirtyNineKVsAllClose}{}\renewcommand{\rhtaXPctRfAllThirtyNineKVsAllClose}{1.7}  % QLIKE % change, ALL48_N39k vs ALL48 (full set; negative: ALL48_N39k more accurate), Random forest, 16:00 bar
\providecommand{\rhtaXDmRfAllThirtyNineKVsAllClose}{}\renewcommand{\rhtaXDmRfAllThirtyNineKVsAllClose}{2.38}  % ... its Diebold-Mariano statistic
\providecommand{\rhtaXPRfAllThirtyNineKVsAllClose}{}\renewcommand{\rhtaXPRfAllThirtyNineKVsAllClose}{\ensuremath{p = 0.017}}  % ... its p statement
\providecommand{\rhtaXPctRidgeAllThirtyNineKVsAllBar}{}\renewcommand{\rhtaXPctRidgeAllThirtyNineKVsAllBar}{0.8}  % QLIKE % change, ALL48_N39k vs ALL48 (full set; negative: ALL48_N39k more accurate), Ridge, every bar
\providecommand{\rhtaXDmRidgeAllThirtyNineKVsAllBar}{}\renewcommand{\rhtaXDmRidgeAllThirtyNineKVsAllBar}{1.56}  % ... its Diebold-Mariano statistic
\providecommand{\rhtaXPRidgeAllThirtyNineKVsAllBar}{}\renewcommand{\rhtaXPRidgeAllThirtyNineKVsAllBar}{\ensuremath{p = 0.12}}  % ... its p statement
\providecommand{\rhtaXPctLassoAllThirtyNineKVsAllBar}{}\renewcommand{\rhtaXPctLassoAllThirtyNineKVsAllBar}{1.6}  % QLIKE % change, ALL48_N39k vs ALL48 (full set; negative: ALL48_N39k more accurate), Lasso, every bar
\providecommand{\rhtaXDmLassoAllThirtyNineKVsAllBar}{}\renewcommand{\rhtaXDmLassoAllThirtyNineKVsAllBar}{3.26}  % ... its Diebold-Mariano statistic
\providecommand{\rhtaXPLassoAllThirtyNineKVsAllBar}{}\renewcommand{\rhtaXPLassoAllThirtyNineKVsAllBar}{\ensuremath{p = 0.001}}  % ... its p statement
\providecommand{\rhtaXPctEnetAllThirtyNineKVsAllBar}{}\renewcommand{\rhtaXPctEnetAllThirtyNineKVsAllBar}{2.7}  % QLIKE % change, ALL48_N39k vs ALL48 (full set; negative: ALL48_N39k more accurate), Elastic net, every bar
\providecommand{\rhtaXDmEnetAllThirtyNineKVsAllBar}{}\renewcommand{\rhtaXDmEnetAllThirtyNineKVsAllBar}{3.42}  % ... its Diebold-Mariano statistic
\providecommand{\rhtaXPEnetAllThirtyNineKVsAllBar}{}\renewcommand{\rhtaXPEnetAllThirtyNineKVsAllBar}{\ensuremath{p < 0.001}}  % ... its p statement
\providecommand{\rhtaXPctLgbmAllThirtyNineKVsAllBar}{}\renewcommand{\rhtaXPctLgbmAllThirtyNineKVsAllBar}{5.7}  % QLIKE % change, ALL48_N39k vs ALL48 (full set; negative: ALL48_N39k more accurate), LightGBM, every bar
\providecommand{\rhtaXDmLgbmAllThirtyNineKVsAllBar}{}\renewcommand{\rhtaXDmLgbmAllThirtyNineKVsAllBar}{6.76}  % ... its Diebold-Mariano statistic
\providecommand{\rhtaXPLgbmAllThirtyNineKVsAllBar}{}\renewcommand{\rhtaXPLgbmAllThirtyNineKVsAllBar}{\ensuremath{p < 0.001}}  % ... its p statement
\providecommand{\rhtaXPctXgbAllThirtyNineKVsAllBar}{}\renewcommand{\rhtaXPctXgbAllThirtyNineKVsAllBar}{3.2}  % QLIKE % change, ALL48_N39k vs ALL48 (full set; negative: ALL48_N39k more accurate), XGBoost, every bar
\providecommand{\rhtaXDmXgbAllThirtyNineKVsAllBar}{}\renewcommand{\rhtaXDmXgbAllThirtyNineKVsAllBar}{6.86}  % ... its Diebold-Mariano statistic
\providecommand{\rhtaXPXgbAllThirtyNineKVsAllBar}{}\renewcommand{\rhtaXPXgbAllThirtyNineKVsAllBar}{\ensuremath{p < 0.001}}  % ... its p statement
\providecommand{\rhtaXPctRfAllThirtyNineKVsAllBar}{}\renewcommand{\rhtaXPctRfAllThirtyNineKVsAllBar}{2.6}  % QLIKE % change, ALL48_N39k vs ALL48 (full set; negative: ALL48_N39k more accurate), Random forest, every bar
\providecommand{\rhtaXDmRfAllThirtyNineKVsAllBar}{}\renewcommand{\rhtaXDmRfAllThirtyNineKVsAllBar}{4.91}  % ... its Diebold-Mariano statistic
\providecommand{\rhtaXPRfAllThirtyNineKVsAllBar}{}\renewcommand{\rhtaXPRfAllThirtyNineKVsAllBar}{\ensuremath{p < 0.001}}  % ... its p statement
\providecommand{\rhtaShRidgeHarRth}{}\renewcommand{\rhtaShRidgeHarRth}{0.95}  % Sharpe ratio (D19), mid-quotes: Ridge, baseline, RTH13 (main design: the 13 regular-hours bars, 3,000 days)
\providecommand{\rhtaShXRidgeHarRth}{}\renewcommand{\rhtaShXRidgeHarRth}{0.60}  % ... at quoted bid/ask prices
\providecommand{\rhtaShRidgeFullRth}{}\renewcommand{\rhtaShRidgeFullRth}{1.77}  % Sharpe ratio (D19), mid-quotes: Ridge, all_features, RTH13 (main design: the 13 regular-hours bars, 3,000 days)
\providecommand{\rhtaShXRidgeFullRth}{}\renewcommand{\rhtaShXRidgeFullRth}{1.41}  % ... at quoted bid/ask prices
\providecommand{\rhtaPShRidgeRth}{}\renewcommand{\rhtaPShRidgeRth}{\ensuremath{p = 0.11}}  % Ledoit-Wolf p, Sharpe ratio full minus HAR (mid): Ridge, RTH13
\providecommand{\rhtaShLassoHarRth}{}\renewcommand{\rhtaShLassoHarRth}{0.95}  % Sharpe ratio (D19), mid-quotes: Lasso, baseline, RTH13 (main design: the 13 regular-hours bars, 3,000 days)
\providecommand{\rhtaShXLassoHarRth}{}\renewcommand{\rhtaShXLassoHarRth}{0.60}  % ... at quoted bid/ask prices
\providecommand{\rhtaShLassoFullRth}{}\renewcommand{\rhtaShLassoFullRth}{1.72}  % Sharpe ratio (D19), mid-quotes: Lasso, all_features, RTH13 (main design: the 13 regular-hours bars, 3,000 days)
\providecommand{\rhtaShXLassoFullRth}{}\renewcommand{\rhtaShXLassoFullRth}{1.37}  % ... at quoted bid/ask prices
\providecommand{\rhtaPShLassoRth}{}\renewcommand{\rhtaPShLassoRth}{\ensuremath{p = 0.12}}  % Ledoit-Wolf p, Sharpe ratio full minus HAR (mid): Lasso, RTH13
\providecommand{\rhtaShEnetHarRth}{}\renewcommand{\rhtaShEnetHarRth}{0.95}  % Sharpe ratio (D19), mid-quotes: Elastic net, baseline, RTH13 (main design: the 13 regular-hours bars, 3,000 days)
\providecommand{\rhtaShXEnetHarRth}{}\renewcommand{\rhtaShXEnetHarRth}{0.60}  % ... at quoted bid/ask prices
\providecommand{\rhtaShEnetFullRth}{}\renewcommand{\rhtaShEnetFullRth}{1.86}  % Sharpe ratio (D19), mid-quotes: Elastic net, all_features, RTH13 (main design: the 13 regular-hours bars, 3,000 days)
\providecommand{\rhtaShXEnetFullRth}{}\renewcommand{\rhtaShXEnetFullRth}{1.50}  % ... at quoted bid/ask prices
\providecommand{\rhtaPShEnetRth}{}\renewcommand{\rhtaPShEnetRth}{\ensuremath{p = 0.077}}  % Ledoit-Wolf p, Sharpe ratio full minus HAR (mid): Elastic net, RTH13
\providecommand{\rhtaShLgbmHarRth}{}\renewcommand{\rhtaShLgbmHarRth}{0.65}  % Sharpe ratio (D19), mid-quotes: LightGBM, baseline, RTH13 (main design: the 13 regular-hours bars, 3,000 days)
\providecommand{\rhtaShXLgbmHarRth}{}\renewcommand{\rhtaShXLgbmHarRth}{0.31}  % ... at quoted bid/ask prices
\providecommand{\rhtaShLgbmFullRth}{}\renewcommand{\rhtaShLgbmFullRth}{1.59}  % Sharpe ratio (D19), mid-quotes: LightGBM, all_features, RTH13 (main design: the 13 regular-hours bars, 3,000 days)
\providecommand{\rhtaShXLgbmFullRth}{}\renewcommand{\rhtaShXLgbmFullRth}{1.23}  % ... at quoted bid/ask prices
\providecommand{\rhtaPShLgbmRth}{}\renewcommand{\rhtaPShLgbmRth}{\ensuremath{p = 0.040}}  % Ledoit-Wolf p, Sharpe ratio full minus HAR (mid): LightGBM, RTH13
\providecommand{\rhtaShXgbHarRth}{}\renewcommand{\rhtaShXgbHarRth}{0.98}  % Sharpe ratio (D19), mid-quotes: XGBoost, baseline, RTH13 (main design: the 13 regular-hours bars, 3,000 days)
\providecommand{\rhtaShXXgbHarRth}{}\renewcommand{\rhtaShXXgbHarRth}{0.63}  % ... at quoted bid/ask prices
\providecommand{\rhtaShXgbFullRth}{}\renewcommand{\rhtaShXgbFullRth}{1.47}  % Sharpe ratio (D19), mid-quotes: XGBoost, all_features, RTH13 (main design: the 13 regular-hours bars, 3,000 days)
\providecommand{\rhtaShXXgbFullRth}{}\renewcommand{\rhtaShXXgbFullRth}{1.12}  % ... at quoted bid/ask prices
\providecommand{\rhtaPShXgbRth}{}\renewcommand{\rhtaPShXgbRth}{\ensuremath{p = 0.27}}  % Ledoit-Wolf p, Sharpe ratio full minus HAR (mid): XGBoost, RTH13
\providecommand{\rhtaShRfHarRth}{}\renewcommand{\rhtaShRfHarRth}{0.32}  % Sharpe ratio (D19), mid-quotes: Random forest, baseline, RTH13 (main design: the 13 regular-hours bars, 3,000 days)
\providecommand{\rhtaShXRfHarRth}{}\renewcommand{\rhtaShXRfHarRth}{\ensuremath{-}0.03}  % ... at quoted bid/ask prices
\providecommand{\rhtaShRfFullRth}{}\renewcommand{\rhtaShRfFullRth}{1.88}  % Sharpe ratio (D19), mid-quotes: Random forest, all_features, RTH13 (main design: the 13 regular-hours bars, 3,000 days)
\providecommand{\rhtaShXRfFullRth}{}\renewcommand{\rhtaShXRfFullRth}{1.52}  % ... at quoted bid/ask prices
\providecommand{\rhtaPShRfRth}{}\renewcommand{\rhtaPShRfRth}{\ensuremath{p = 0.020}}  % Ledoit-Wolf p, Sharpe ratio full minus HAR (mid): Random forest, RTH13
\providecommand{\rhtaShRidgeHarAll}{}\renewcommand{\rhtaShRidgeHarAll}{0.95}  % Sharpe ratio (D19), mid-quotes: Ridge, baseline, ALL48 (all 48 intervals, 3,000 days)
\providecommand{\rhtaShXRidgeHarAll}{}\renewcommand{\rhtaShXRidgeHarAll}{0.59}  % ... at quoted bid/ask prices
\providecommand{\rhtaShRidgeFullAll}{}\renewcommand{\rhtaShRidgeFullAll}{1.61}  % Sharpe ratio (D19), mid-quotes: Ridge, all_features, ALL48 (all 48 intervals, 3,000 days)
\providecommand{\rhtaShXRidgeFullAll}{}\renewcommand{\rhtaShXRidgeFullAll}{1.26}  % ... at quoted bid/ask prices
\providecommand{\rhtaShLassoHarAll}{}\renewcommand{\rhtaShLassoHarAll}{0.93}  % Sharpe ratio (D19), mid-quotes: Lasso, baseline, ALL48 (all 48 intervals, 3,000 days)
\providecommand{\rhtaShXLassoHarAll}{}\renewcommand{\rhtaShXLassoHarAll}{0.58}  % ... at quoted bid/ask prices
\providecommand{\rhtaShLassoFullAll}{}\renewcommand{\rhtaShLassoFullAll}{1.50}  % Sharpe ratio (D19), mid-quotes: Lasso, all_features, ALL48 (all 48 intervals, 3,000 days)
\providecommand{\rhtaShXLassoFullAll}{}\renewcommand{\rhtaShXLassoFullAll}{1.14}  % ... at quoted bid/ask prices
\providecommand{\rhtaShEnetHarAll}{}\renewcommand{\rhtaShEnetHarAll}{0.89}  % Sharpe ratio (D19), mid-quotes: Elastic net, baseline, ALL48 (all 48 intervals, 3,000 days)
\providecommand{\rhtaShXEnetHarAll}{}\renewcommand{\rhtaShXEnetHarAll}{0.54}  % ... at quoted bid/ask prices
\providecommand{\rhtaShEnetFullAll}{}\renewcommand{\rhtaShEnetFullAll}{1.52}  % Sharpe ratio (D19), mid-quotes: Elastic net, all_features, ALL48 (all 48 intervals, 3,000 days)
\providecommand{\rhtaShXEnetFullAll}{}\renewcommand{\rhtaShXEnetFullAll}{1.17}  % ... at quoted bid/ask prices
\providecommand{\rhtaShLgbmHarAll}{}\renewcommand{\rhtaShLgbmHarAll}{0.21}  % Sharpe ratio (D19), mid-quotes: LightGBM, baseline, ALL48 (all 48 intervals, 3,000 days)
\providecommand{\rhtaShXLgbmHarAll}{}\renewcommand{\rhtaShXLgbmHarAll}{\ensuremath{-}0.14}  % ... at quoted bid/ask prices
\providecommand{\rhtaShLgbmFullAll}{}\renewcommand{\rhtaShLgbmFullAll}{1.25}  % Sharpe ratio (D19), mid-quotes: LightGBM, all_features, ALL48 (all 48 intervals, 3,000 days)
\providecommand{\rhtaShXLgbmFullAll}{}\renewcommand{\rhtaShXLgbmFullAll}{0.90}  % ... at quoted bid/ask prices
\providecommand{\rhtaPShLgbmAll}{}\renewcommand{\rhtaPShLgbmAll}{\ensuremath{p = 0.11}}  % Ledoit-Wolf p, Sharpe ratio full minus HAR (mid): LightGBM, ALL48
\providecommand{\rhtaShXgbHarAll}{}\renewcommand{\rhtaShXgbHarAll}{0.73}  % Sharpe ratio (D19), mid-quotes: XGBoost, baseline, ALL48 (all 48 intervals, 3,000 days)
\providecommand{\rhtaShXXgbHarAll}{}\renewcommand{\rhtaShXXgbHarAll}{0.38}  % ... at quoted bid/ask prices
\providecommand{\rhtaShXgbFullAll}{}\renewcommand{\rhtaShXgbFullAll}{1.50}  % Sharpe ratio (D19), mid-quotes: XGBoost, all_features, ALL48 (all 48 intervals, 3,000 days)
\providecommand{\rhtaShXXgbFullAll}{}\renewcommand{\rhtaShXXgbFullAll}{1.15}  % ... at quoted bid/ask prices
\providecommand{\rhtaPShXgbAll}{}\renewcommand{\rhtaPShXgbAll}{\ensuremath{p = 0.11}}  % Ledoit-Wolf p, Sharpe ratio full minus HAR (mid): XGBoost, ALL48
\providecommand{\rhtaShRfHarAll}{}\renewcommand{\rhtaShRfHarAll}{0.65}  % Sharpe ratio (D19), mid-quotes: Random forest, baseline, ALL48 (all 48 intervals, 3,000 days)
\providecommand{\rhtaShXRfHarAll}{}\renewcommand{\rhtaShXRfHarAll}{0.29}  % ... at quoted bid/ask prices
\providecommand{\rhtaShRfFullAll}{}\renewcommand{\rhtaShRfFullAll}{1.82}  % Sharpe ratio (D19), mid-quotes: Random forest, all_features, ALL48 (all 48 intervals, 3,000 days)
\providecommand{\rhtaShXRfFullAll}{}\renewcommand{\rhtaShXRfFullAll}{1.46}  % ... at quoted bid/ask prices
\providecommand{\rhtaPShRfAll}{}\renewcommand{\rhtaPShRfAll}{\ensuremath{p = 0.066}}  % Ledoit-Wolf p, Sharpe ratio full minus HAR (mid): Random forest, ALL48
\providecommand{\rhtaShRidgeHarCl}{}\renewcommand{\rhtaShRidgeHarCl}{1.03}  % Sharpe ratio (D19), mid-quotes: Ridge, baseline, CLOSE (the closing half hour only, 3,000 closes)
\providecommand{\rhtaShXRidgeHarCl}{}\renewcommand{\rhtaShXRidgeHarCl}{0.68}  % ... at quoted bid/ask prices
\providecommand{\rhtaShRidgeFullCl}{}\renewcommand{\rhtaShRidgeFullCl}{1.26}  % Sharpe ratio (D19), mid-quotes: Ridge, all_features, CLOSE (the closing half hour only, 3,000 closes)
\providecommand{\rhtaShXRidgeFullCl}{}\renewcommand{\rhtaShXRidgeFullCl}{0.91}  % ... at quoted bid/ask prices
\providecommand{\rhtaShLassoHarCl}{}\renewcommand{\rhtaShLassoHarCl}{1.03}  % Sharpe ratio (D19), mid-quotes: Lasso, baseline, CLOSE (the closing half hour only, 3,000 closes)
\providecommand{\rhtaShXLassoHarCl}{}\renewcommand{\rhtaShXLassoHarCl}{0.68}  % ... at quoted bid/ask prices
\providecommand{\rhtaShLassoFullCl}{}\renewcommand{\rhtaShLassoFullCl}{1.58}  % Sharpe ratio (D19), mid-quotes: Lasso, all_features, CLOSE (the closing half hour only, 3,000 closes)
\providecommand{\rhtaShXLassoFullCl}{}\renewcommand{\rhtaShXLassoFullCl}{1.23}  % ... at quoted bid/ask prices
\providecommand{\rhtaShEnetHarCl}{}\renewcommand{\rhtaShEnetHarCl}{1.03}  % Sharpe ratio (D19), mid-quotes: Elastic net, baseline, CLOSE (the closing half hour only, 3,000 closes)
\providecommand{\rhtaShXEnetHarCl}{}\renewcommand{\rhtaShXEnetHarCl}{0.68}  % ... at quoted bid/ask prices
\providecommand{\rhtaShEnetFullCl}{}\renewcommand{\rhtaShEnetFullCl}{1.71}  % Sharpe ratio (D19), mid-quotes: Elastic net, all_features, CLOSE (the closing half hour only, 3,000 closes)
\providecommand{\rhtaShXEnetFullCl}{}\renewcommand{\rhtaShXEnetFullCl}{1.36}  % ... at quoted bid/ask prices
\providecommand{\rhtaShLgbmHarCl}{}\renewcommand{\rhtaShLgbmHarCl}{\ensuremath{-}0.03}  % Sharpe ratio (D19), mid-quotes: LightGBM, baseline, CLOSE (the closing half hour only, 3,000 closes)
\providecommand{\rhtaShXLgbmHarCl}{}\renewcommand{\rhtaShXLgbmHarCl}{\ensuremath{-}0.37}  % ... at quoted bid/ask prices
\providecommand{\rhtaShLgbmFullCl}{}\renewcommand{\rhtaShLgbmFullCl}{1.94}  % Sharpe ratio (D19), mid-quotes: LightGBM, all_features, CLOSE (the closing half hour only, 3,000 closes)
\providecommand{\rhtaShXLgbmFullCl}{}\renewcommand{\rhtaShXLgbmFullCl}{1.60}  % ... at quoted bid/ask prices
\providecommand{\rhtaPShLgbmCl}{}\renewcommand{\rhtaPShLgbmCl}{\ensuremath{p = 0.003}}  % Ledoit-Wolf p, Sharpe ratio full minus HAR (mid): LightGBM, CLOSE
\providecommand{\rhtaShXgbHarCl}{}\renewcommand{\rhtaShXgbHarCl}{0.40}  % Sharpe ratio (D19), mid-quotes: XGBoost, baseline, CLOSE (the closing half hour only, 3,000 closes)
\providecommand{\rhtaShXXgbHarCl}{}\renewcommand{\rhtaShXXgbHarCl}{0.05}  % ... at quoted bid/ask prices
\providecommand{\rhtaShXgbFullCl}{}\renewcommand{\rhtaShXgbFullCl}{1.62}  % Sharpe ratio (D19), mid-quotes: XGBoost, all_features, CLOSE (the closing half hour only, 3,000 closes)
\providecommand{\rhtaShXXgbFullCl}{}\renewcommand{\rhtaShXXgbFullCl}{1.27}  % ... at quoted bid/ask prices
\providecommand{\rhtaPShXgbCl}{}\renewcommand{\rhtaPShXgbCl}{\ensuremath{p = 0.055}}  % Ledoit-Wolf p, Sharpe ratio full minus HAR (mid): XGBoost, CLOSE
\providecommand{\rhtaShRfHarCl}{}\renewcommand{\rhtaShRfHarCl}{0.29}  % Sharpe ratio (D19), mid-quotes: Random forest, baseline, CLOSE (the closing half hour only, 3,000 closes)
\providecommand{\rhtaShXRfHarCl}{}\renewcommand{\rhtaShXRfHarCl}{\ensuremath{-}0.05}  % ... at quoted bid/ask prices
\providecommand{\rhtaShRfFullCl}{}\renewcommand{\rhtaShRfFullCl}{1.29}  % Sharpe ratio (D19), mid-quotes: Random forest, all_features, CLOSE (the closing half hour only, 3,000 closes)
\providecommand{\rhtaShXRfFullCl}{}\renewcommand{\rhtaShXRfFullCl}{0.94}  % ... at quoted bid/ask prices
\providecommand{\rhtaPShRfCl}{}\renewcommand{\rhtaPShRfCl}{\ensuremath{p = 0.14}}  % Ledoit-Wolf p, Sharpe ratio full minus HAR (mid): Random forest, CLOSE
\providecommand{\rhtaShRidgeHarClNat}{}\renewcommand{\rhtaShRidgeHarClNat}{0.94}  % Sharpe ratio (D19), mid-quotes: Ridge, baseline, CLOSE_NAT (closing rows, design built natively on them; linear only)
\providecommand{\rhtaShXRidgeHarClNat}{}\renewcommand{\rhtaShXRidgeHarClNat}{0.59}  % ... at quoted bid/ask prices
\providecommand{\rhtaShRidgeFullClNat}{}\renewcommand{\rhtaShRidgeFullClNat}{1.24}  % Sharpe ratio (D19), mid-quotes: Ridge, all_features, CLOSE_NAT (closing rows, design built natively on them; linear only)
\providecommand{\rhtaShXRidgeFullClNat}{}\renewcommand{\rhtaShXRidgeFullClNat}{0.89}  % ... at quoted bid/ask prices
\providecommand{\rhtaShLassoHarClNat}{}\renewcommand{\rhtaShLassoHarClNat}{0.94}  % Sharpe ratio (D19), mid-quotes: Lasso, baseline, CLOSE_NAT (closing rows, design built natively on them; linear only)
\providecommand{\rhtaShXLassoHarClNat}{}\renewcommand{\rhtaShXLassoHarClNat}{0.59}  % ... at quoted bid/ask prices
\providecommand{\rhtaShLassoFullClNat}{}\renewcommand{\rhtaShLassoFullClNat}{1.23}  % Sharpe ratio (D19), mid-quotes: Lasso, all_features, CLOSE_NAT (closing rows, design built natively on them; linear only)
\providecommand{\rhtaShXLassoFullClNat}{}\renewcommand{\rhtaShXLassoFullClNat}{0.89}  % ... at quoted bid/ask prices
\providecommand{\rhtaShEnetHarClNat}{}\renewcommand{\rhtaShEnetHarClNat}{0.94}  % Sharpe ratio (D19), mid-quotes: Elastic net, baseline, CLOSE_NAT (closing rows, design built natively on them; linear only)
\providecommand{\rhtaShXEnetHarClNat}{}\renewcommand{\rhtaShXEnetHarClNat}{0.59}  % ... at quoted bid/ask prices
\providecommand{\rhtaShEnetFullClNat}{}\renewcommand{\rhtaShEnetFullClNat}{1.71}  % Sharpe ratio (D19), mid-quotes: Elastic net, all_features, CLOSE_NAT (closing rows, design built natively on them; linear only)
\providecommand{\rhtaShXEnetFullClNat}{}\renewcommand{\rhtaShXEnetFullClNat}{1.37}  % ... at quoted bid/ask prices
\providecommand{\rhtaShRidgeHarAllCt}{}\renewcommand{\rhtaShRidgeHarAllCt}{0.95}  % Sharpe ratio (D19), mid-quotes: Ridge, baseline, ALL48@16 (ALL48, penalties chosen on the 16:00 validation rows; linear only)
\providecommand{\rhtaShXRidgeHarAllCt}{}\renewcommand{\rhtaShXRidgeHarAllCt}{0.59}  % ... at quoted bid/ask prices
\providecommand{\rhtaShRidgeFullAllCt}{}\renewcommand{\rhtaShRidgeFullAllCt}{1.40}  % Sharpe ratio (D19), mid-quotes: Ridge, all_features, ALL48@16 (ALL48, penalties chosen on the 16:00 validation rows; linear only)
\providecommand{\rhtaShXRidgeFullAllCt}{}\renewcommand{\rhtaShXRidgeFullAllCt}{1.05}  % ... at quoted bid/ask prices
\providecommand{\rhtaShLassoHarAllCt}{}\renewcommand{\rhtaShLassoHarAllCt}{0.96}  % Sharpe ratio (D19), mid-quotes: Lasso, baseline, ALL48@16 (ALL48, penalties chosen on the 16:00 validation rows; linear only)
\providecommand{\rhtaShXLassoHarAllCt}{}\renewcommand{\rhtaShXLassoHarAllCt}{0.60}  % ... at quoted bid/ask prices
\providecommand{\rhtaShLassoFullAllCt}{}\renewcommand{\rhtaShLassoFullAllCt}{1.35}  % Sharpe ratio (D19), mid-quotes: Lasso, all_features, ALL48@16 (ALL48, penalties chosen on the 16:00 validation rows; linear only)
\providecommand{\rhtaShXLassoFullAllCt}{}\renewcommand{\rhtaShXLassoFullAllCt}{1.00}  % ... at quoted bid/ask prices
\providecommand{\rhtaShEnetHarAllCt}{}\renewcommand{\rhtaShEnetHarAllCt}{0.80}  % Sharpe ratio (D19), mid-quotes: Elastic net, baseline, ALL48@16 (ALL48, penalties chosen on the 16:00 validation rows; linear only)
\providecommand{\rhtaShXEnetHarAllCt}{}\renewcommand{\rhtaShXEnetHarAllCt}{0.45}  % ... at quoted bid/ask prices
\providecommand{\rhtaShEnetFullAllCt}{}\renewcommand{\rhtaShEnetFullAllCt}{1.52}  % Sharpe ratio (D19), mid-quotes: Elastic net, all_features, ALL48@16 (ALL48, penalties chosen on the 16:00 validation rows; linear only)
\providecommand{\rhtaShXEnetFullAllCt}{}\renewcommand{\rhtaShXEnetFullAllCt}{1.17}  % ... at quoted bid/ask prices
\providecommand{\rhtaShRidgeHarRthCt}{}\renewcommand{\rhtaShRidgeHarRthCt}{0.95}  % Sharpe ratio (D19), mid-quotes: Ridge, baseline, RTH13@16 (RTH13, penalties chosen on the 16:00 validation rows; linear only)
\providecommand{\rhtaShXRidgeHarRthCt}{}\renewcommand{\rhtaShXRidgeHarRthCt}{0.60}  % ... at quoted bid/ask prices
\providecommand{\rhtaShRidgeFullRthCt}{}\renewcommand{\rhtaShRidgeFullRthCt}{1.77}  % Sharpe ratio (D19), mid-quotes: Ridge, all_features, RTH13@16 (RTH13, penalties chosen on the 16:00 validation rows; linear only)
\providecommand{\rhtaShXRidgeFullRthCt}{}\renewcommand{\rhtaShXRidgeFullRthCt}{1.41}  % ... at quoted bid/ask prices
\providecommand{\rhtaShLassoHarRthCt}{}\renewcommand{\rhtaShLassoHarRthCt}{1.01}  % Sharpe ratio (D19), mid-quotes: Lasso, baseline, RTH13@16 (RTH13, penalties chosen on the 16:00 validation rows; linear only)
\providecommand{\rhtaShXLassoHarRthCt}{}\renewcommand{\rhtaShXLassoHarRthCt}{0.65}  % ... at quoted bid/ask prices
\providecommand{\rhtaShLassoFullRthCt}{}\renewcommand{\rhtaShLassoFullRthCt}{1.88}  % Sharpe ratio (D19), mid-quotes: Lasso, all_features, RTH13@16 (RTH13, penalties chosen on the 16:00 validation rows; linear only)
\providecommand{\rhtaShXLassoFullRthCt}{}\renewcommand{\rhtaShXLassoFullRthCt}{1.52}  % ... at quoted bid/ask prices
\providecommand{\rhtaShEnetHarRthCt}{}\renewcommand{\rhtaShEnetHarRthCt}{1.01}  % Sharpe ratio (D19), mid-quotes: Elastic net, baseline, RTH13@16 (RTH13, penalties chosen on the 16:00 validation rows; linear only)
\providecommand{\rhtaShXEnetHarRthCt}{}\renewcommand{\rhtaShXEnetHarRthCt}{0.66}  % ... at quoted bid/ask prices
\providecommand{\rhtaShEnetFullRthCt}{}\renewcommand{\rhtaShEnetFullRthCt}{1.86}  % Sharpe ratio (D19), mid-quotes: Elastic net, all_features, RTH13@16 (RTH13, penalties chosen on the 16:00 validation rows; linear only)
\providecommand{\rhtaShXEnetFullRthCt}{}\renewcommand{\rhtaShXEnetFullRthCt}{1.50}  % ... at quoted bid/ask prices
\providecommand{\rhtaShRidgeHarAllThreeK}{}\renewcommand{\rhtaShRidgeHarAllThreeK}{1.56}  % Sharpe ratio (D19), mid-quotes: Ridge, baseline, ALL48_N3k (ALL48, the last 3,000 rows)
\providecommand{\rhtaShXRidgeHarAllThreeK}{}\renewcommand{\rhtaShXRidgeHarAllThreeK}{1.21}  % ... at quoted bid/ask prices
\providecommand{\rhtaShRidgeFullAllThreeK}{}\renewcommand{\rhtaShRidgeFullAllThreeK}{1.45}  % Sharpe ratio (D19), mid-quotes: Ridge, all_features, ALL48_N3k (ALL48, the last 3,000 rows)
\providecommand{\rhtaShXRidgeFullAllThreeK}{}\renewcommand{\rhtaShXRidgeFullAllThreeK}{1.10}  % ... at quoted bid/ask prices
\providecommand{\rhtaShLassoHarAllThreeK}{}\renewcommand{\rhtaShLassoHarAllThreeK}{1.57}  % Sharpe ratio (D19), mid-quotes: Lasso, baseline, ALL48_N3k (ALL48, the last 3,000 rows)
\providecommand{\rhtaShXLassoHarAllThreeK}{}\renewcommand{\rhtaShXLassoHarAllThreeK}{1.21}  % ... at quoted bid/ask prices
\providecommand{\rhtaShLassoFullAllThreeK}{}\renewcommand{\rhtaShLassoFullAllThreeK}{1.41}  % Sharpe ratio (D19), mid-quotes: Lasso, all_features, ALL48_N3k (ALL48, the last 3,000 rows)
\providecommand{\rhtaShXLassoFullAllThreeK}{}\renewcommand{\rhtaShXLassoFullAllThreeK}{1.06}  % ... at quoted bid/ask prices
\providecommand{\rhtaShEnetHarAllThreeK}{}\renewcommand{\rhtaShEnetHarAllThreeK}{1.52}  % Sharpe ratio (D19), mid-quotes: Elastic net, baseline, ALL48_N3k (ALL48, the last 3,000 rows)
\providecommand{\rhtaShXEnetHarAllThreeK}{}\renewcommand{\rhtaShXEnetHarAllThreeK}{1.16}  % ... at quoted bid/ask prices
\providecommand{\rhtaShEnetFullAllThreeK}{}\renewcommand{\rhtaShEnetFullAllThreeK}{1.63}  % Sharpe ratio (D19), mid-quotes: Elastic net, all_features, ALL48_N3k (ALL48, the last 3,000 rows)
\providecommand{\rhtaShXEnetFullAllThreeK}{}\renewcommand{\rhtaShXEnetFullAllThreeK}{1.27}  % ... at quoted bid/ask prices
\providecommand{\rhtaShLgbmHarAllThreeK}{}\renewcommand{\rhtaShLgbmHarAllThreeK}{0.92}  % Sharpe ratio (D19), mid-quotes: LightGBM, baseline, ALL48_N3k (ALL48, the last 3,000 rows)
\providecommand{\rhtaShXLgbmHarAllThreeK}{}\renewcommand{\rhtaShXLgbmHarAllThreeK}{0.57}  % ... at quoted bid/ask prices
\providecommand{\rhtaShLgbmFullAllThreeK}{}\renewcommand{\rhtaShLgbmFullAllThreeK}{1.35}  % Sharpe ratio (D19), mid-quotes: LightGBM, all_features, ALL48_N3k (ALL48, the last 3,000 rows)
\providecommand{\rhtaShXLgbmFullAllThreeK}{}\renewcommand{\rhtaShXLgbmFullAllThreeK}{0.99}  % ... at quoted bid/ask prices
\providecommand{\rhtaPShLgbmAllThreeK}{}\renewcommand{\rhtaPShLgbmAllThreeK}{\ensuremath{p = 0.20}}  % Ledoit-Wolf p, Sharpe ratio full minus HAR (mid): LightGBM, ALL48_N3k
\providecommand{\rhtaShXgbHarAllThreeK}{}\renewcommand{\rhtaShXgbHarAllThreeK}{0.93}  % Sharpe ratio (D19), mid-quotes: XGBoost, baseline, ALL48_N3k (ALL48, the last 3,000 rows)
\providecommand{\rhtaShXXgbHarAllThreeK}{}\renewcommand{\rhtaShXXgbHarAllThreeK}{0.58}  % ... at quoted bid/ask prices
\providecommand{\rhtaShXgbFullAllThreeK}{}\renewcommand{\rhtaShXgbFullAllThreeK}{1.24}  % Sharpe ratio (D19), mid-quotes: XGBoost, all_features, ALL48_N3k (ALL48, the last 3,000 rows)
\providecommand{\rhtaShXXgbFullAllThreeK}{}\renewcommand{\rhtaShXXgbFullAllThreeK}{0.88}  % ... at quoted bid/ask prices
\providecommand{\rhtaPShXgbAllThreeK}{}\renewcommand{\rhtaPShXgbAllThreeK}{\ensuremath{p = 0.41}}  % Ledoit-Wolf p, Sharpe ratio full minus HAR (mid): XGBoost, ALL48_N3k
\providecommand{\rhtaShRfHarAllThreeK}{}\renewcommand{\rhtaShRfHarAllThreeK}{0.87}  % Sharpe ratio (D19), mid-quotes: Random forest, baseline, ALL48_N3k (ALL48, the last 3,000 rows)
\providecommand{\rhtaShXRfHarAllThreeK}{}\renewcommand{\rhtaShXRfHarAllThreeK}{0.52}  % ... at quoted bid/ask prices
\providecommand{\rhtaShRfFullAllThreeK}{}\renewcommand{\rhtaShRfFullAllThreeK}{1.30}  % Sharpe ratio (D19), mid-quotes: Random forest, all_features, ALL48_N3k (ALL48, the last 3,000 rows)
\providecommand{\rhtaShXRfFullAllThreeK}{}\renewcommand{\rhtaShXRfFullAllThreeK}{0.94}  % ... at quoted bid/ask prices
\providecommand{\rhtaPShRfAllThreeK}{}\renewcommand{\rhtaPShRfAllThreeK}{\ensuremath{p = 0.21}}  % Ledoit-Wolf p, Sharpe ratio full minus HAR (mid): Random forest, ALL48_N3k
\providecommand{\rhtaShRidgeHarRthThreeK}{}\renewcommand{\rhtaShRidgeHarRthThreeK}{1.16}  % Sharpe ratio (D19), mid-quotes: Ridge, baseline, RTH13_N3k (RTH13, the last 3,000 rows)
\providecommand{\rhtaShXRidgeHarRthThreeK}{}\renewcommand{\rhtaShXRidgeHarRthThreeK}{0.80}  % ... at quoted bid/ask prices
\providecommand{\rhtaShRidgeFullRthThreeK}{}\renewcommand{\rhtaShRidgeFullRthThreeK}{1.12}  % Sharpe ratio (D19), mid-quotes: Ridge, all_features, RTH13_N3k (RTH13, the last 3,000 rows)
\providecommand{\rhtaShXRidgeFullRthThreeK}{}\renewcommand{\rhtaShXRidgeFullRthThreeK}{0.77}  % ... at quoted bid/ask prices
\providecommand{\rhtaShLassoHarRthThreeK}{}\renewcommand{\rhtaShLassoHarRthThreeK}{1.16}  % Sharpe ratio (D19), mid-quotes: Lasso, baseline, RTH13_N3k (RTH13, the last 3,000 rows)
\providecommand{\rhtaShXLassoHarRthThreeK}{}\renewcommand{\rhtaShXLassoHarRthThreeK}{0.80}  % ... at quoted bid/ask prices
\providecommand{\rhtaShLassoFullRthThreeK}{}\renewcommand{\rhtaShLassoFullRthThreeK}{1.55}  % Sharpe ratio (D19), mid-quotes: Lasso, all_features, RTH13_N3k (RTH13, the last 3,000 rows)
\providecommand{\rhtaShXLassoFullRthThreeK}{}\renewcommand{\rhtaShXLassoFullRthThreeK}{1.19}  % ... at quoted bid/ask prices
\providecommand{\rhtaShEnetHarRthThreeK}{}\renewcommand{\rhtaShEnetHarRthThreeK}{1.16}  % Sharpe ratio (D19), mid-quotes: Elastic net, baseline, RTH13_N3k (RTH13, the last 3,000 rows)
\providecommand{\rhtaShXEnetHarRthThreeK}{}\renewcommand{\rhtaShXEnetHarRthThreeK}{0.80}  % ... at quoted bid/ask prices
\providecommand{\rhtaShEnetFullRthThreeK}{}\renewcommand{\rhtaShEnetFullRthThreeK}{1.25}  % Sharpe ratio (D19), mid-quotes: Elastic net, all_features, RTH13_N3k (RTH13, the last 3,000 rows)
\providecommand{\rhtaShXEnetFullRthThreeK}{}\renewcommand{\rhtaShXEnetFullRthThreeK}{0.89}  % ... at quoted bid/ask prices
\providecommand{\rhtaShLgbmHarRthThreeK}{}\renewcommand{\rhtaShLgbmHarRthThreeK}{0.25}  % Sharpe ratio (D19), mid-quotes: LightGBM, baseline, RTH13_N3k (RTH13, the last 3,000 rows)
\providecommand{\rhtaShXLgbmHarRthThreeK}{}\renewcommand{\rhtaShXLgbmHarRthThreeK}{\ensuremath{-}0.10}  % ... at quoted bid/ask prices
\providecommand{\rhtaShLgbmFullRthThreeK}{}\renewcommand{\rhtaShLgbmFullRthThreeK}{0.83}  % Sharpe ratio (D19), mid-quotes: LightGBM, all_features, RTH13_N3k (RTH13, the last 3,000 rows)
\providecommand{\rhtaShXLgbmFullRthThreeK}{}\renewcommand{\rhtaShXLgbmFullRthThreeK}{0.47}  % ... at quoted bid/ask prices
\providecommand{\rhtaPShLgbmRthThreeK}{}\renewcommand{\rhtaPShLgbmRthThreeK}{\ensuremath{p = 0.22}}  % Ledoit-Wolf p, Sharpe ratio full minus HAR (mid): LightGBM, RTH13_N3k
\providecommand{\rhtaShXgbHarRthThreeK}{}\renewcommand{\rhtaShXgbHarRthThreeK}{0.87}  % Sharpe ratio (D19), mid-quotes: XGBoost, baseline, RTH13_N3k (RTH13, the last 3,000 rows)
\providecommand{\rhtaShXXgbHarRthThreeK}{}\renewcommand{\rhtaShXXgbHarRthThreeK}{0.52}  % ... at quoted bid/ask prices
\providecommand{\rhtaShXgbFullRthThreeK}{}\renewcommand{\rhtaShXgbFullRthThreeK}{0.86}  % Sharpe ratio (D19), mid-quotes: XGBoost, all_features, RTH13_N3k (RTH13, the last 3,000 rows)
\providecommand{\rhtaShXXgbFullRthThreeK}{}\renewcommand{\rhtaShXXgbFullRthThreeK}{0.51}  % ... at quoted bid/ask prices
\providecommand{\rhtaPShXgbRthThreeK}{}\renewcommand{\rhtaPShXgbRthThreeK}{\ensuremath{p = 0.98}}  % Ledoit-Wolf p, Sharpe ratio full minus HAR (mid): XGBoost, RTH13_N3k
\providecommand{\rhtaShRfHarRthThreeK}{}\renewcommand{\rhtaShRfHarRthThreeK}{0.67}  % Sharpe ratio (D19), mid-quotes: Random forest, baseline, RTH13_N3k (RTH13, the last 3,000 rows)
\providecommand{\rhtaShXRfHarRthThreeK}{}\renewcommand{\rhtaShXRfHarRthThreeK}{0.31}  % ... at quoted bid/ask prices
\providecommand{\rhtaShRfFullRthThreeK}{}\renewcommand{\rhtaShRfFullRthThreeK}{0.91}  % Sharpe ratio (D19), mid-quotes: Random forest, all_features, RTH13_N3k (RTH13, the last 3,000 rows)
\providecommand{\rhtaShXRfFullRthThreeK}{}\renewcommand{\rhtaShXRfFullRthThreeK}{0.56}  % ... at quoted bid/ask prices
\providecommand{\rhtaPShRfRthThreeK}{}\renewcommand{\rhtaPShRfRthThreeK}{\ensuremath{p = 0.53}}  % Ledoit-Wolf p, Sharpe ratio full minus HAR (mid): Random forest, RTH13_N3k
\providecommand{\rhtaShRidgeHarAllThirtyNineK}{}\renewcommand{\rhtaShRidgeHarAllThirtyNineK}{1.23}  % Sharpe ratio (D19), mid-quotes: Ridge, baseline, ALL48_N39k (ALL48, the last 39,000 rows)
\providecommand{\rhtaShXRidgeHarAllThirtyNineK}{}\renewcommand{\rhtaShXRidgeHarAllThirtyNineK}{0.87}  % ... at quoted bid/ask prices
\providecommand{\rhtaShRidgeFullAllThirtyNineK}{}\renewcommand{\rhtaShRidgeFullAllThirtyNineK}{0.76}  % Sharpe ratio (D19), mid-quotes: Ridge, all_features, ALL48_N39k (ALL48, the last 39,000 rows)
\providecommand{\rhtaShXRidgeFullAllThirtyNineK}{}\renewcommand{\rhtaShXRidgeFullAllThirtyNineK}{0.40}  % ... at quoted bid/ask prices
\providecommand{\rhtaShLassoHarAllThirtyNineK}{}\renewcommand{\rhtaShLassoHarAllThirtyNineK}{1.20}  % Sharpe ratio (D19), mid-quotes: Lasso, baseline, ALL48_N39k (ALL48, the last 39,000 rows)
\providecommand{\rhtaShXLassoHarAllThirtyNineK}{}\renewcommand{\rhtaShXLassoHarAllThirtyNineK}{0.84}  % ... at quoted bid/ask prices
\providecommand{\rhtaShLassoFullAllThirtyNineK}{}\renewcommand{\rhtaShLassoFullAllThirtyNineK}{1.33}  % Sharpe ratio (D19), mid-quotes: Lasso, all_features, ALL48_N39k (ALL48, the last 39,000 rows)
\providecommand{\rhtaShXLassoFullAllThirtyNineK}{}\renewcommand{\rhtaShXLassoFullAllThirtyNineK}{0.98}  % ... at quoted bid/ask prices
\providecommand{\rhtaShEnetHarAllThirtyNineK}{}\renewcommand{\rhtaShEnetHarAllThirtyNineK}{1.22}  % Sharpe ratio (D19), mid-quotes: Elastic net, baseline, ALL48_N39k (ALL48, the last 39,000 rows)
\providecommand{\rhtaShXEnetHarAllThirtyNineK}{}\renewcommand{\rhtaShXEnetHarAllThirtyNineK}{0.87}  % ... at quoted bid/ask prices
\providecommand{\rhtaShEnetFullAllThirtyNineK}{}\renewcommand{\rhtaShEnetFullAllThirtyNineK}{1.63}  % Sharpe ratio (D19), mid-quotes: Elastic net, all_features, ALL48_N39k (ALL48, the last 39,000 rows)
\providecommand{\rhtaShXEnetFullAllThirtyNineK}{}\renewcommand{\rhtaShXEnetFullAllThirtyNineK}{1.27}  % ... at quoted bid/ask prices
\providecommand{\rhtaShLgbmHarAllThirtyNineK}{}\renewcommand{\rhtaShLgbmHarAllThirtyNineK}{0.48}  % Sharpe ratio (D19), mid-quotes: LightGBM, baseline, ALL48_N39k (ALL48, the last 39,000 rows)
\providecommand{\rhtaShXLgbmHarAllThirtyNineK}{}\renewcommand{\rhtaShXLgbmHarAllThirtyNineK}{0.13}  % ... at quoted bid/ask prices
\providecommand{\rhtaShLgbmFullAllThirtyNineK}{}\renewcommand{\rhtaShLgbmFullAllThirtyNineK}{1.73}  % Sharpe ratio (D19), mid-quotes: LightGBM, all_features, ALL48_N39k (ALL48, the last 39,000 rows)
\providecommand{\rhtaShXLgbmFullAllThirtyNineK}{}\renewcommand{\rhtaShXLgbmFullAllThirtyNineK}{1.37}  % ... at quoted bid/ask prices
\providecommand{\rhtaPShLgbmAllThirtyNineK}{}\renewcommand{\rhtaPShLgbmAllThirtyNineK}{\ensuremath{p = 0.060}}  % Ledoit-Wolf p, Sharpe ratio full minus HAR (mid): LightGBM, ALL48_N39k
\providecommand{\rhtaShXgbHarAllThirtyNineK}{}\renewcommand{\rhtaShXgbHarAllThirtyNineK}{0.62}  % Sharpe ratio (D19), mid-quotes: XGBoost, baseline, ALL48_N39k (ALL48, the last 39,000 rows)
\providecommand{\rhtaShXXgbHarAllThirtyNineK}{}\renewcommand{\rhtaShXXgbHarAllThirtyNineK}{0.27}  % ... at quoted bid/ask prices
\providecommand{\rhtaShXgbFullAllThirtyNineK}{}\renewcommand{\rhtaShXgbFullAllThirtyNineK}{1.64}  % Sharpe ratio (D19), mid-quotes: XGBoost, all_features, ALL48_N39k (ALL48, the last 39,000 rows)
\providecommand{\rhtaShXXgbFullAllThirtyNineK}{}\renewcommand{\rhtaShXXgbFullAllThirtyNineK}{1.28}  % ... at quoted bid/ask prices
\providecommand{\rhtaPShXgbAllThirtyNineK}{}\renewcommand{\rhtaPShXgbAllThirtyNineK}{\ensuremath{p = 0.12}}  % Ledoit-Wolf p, Sharpe ratio full minus HAR (mid): XGBoost, ALL48_N39k
\providecommand{\rhtaShRfHarAllThirtyNineK}{}\renewcommand{\rhtaShRfHarAllThirtyNineK}{0.87}  % Sharpe ratio (D19), mid-quotes: Random forest, baseline, ALL48_N39k (ALL48, the last 39,000 rows)
\providecommand{\rhtaShXRfHarAllThirtyNineK}{}\renewcommand{\rhtaShXRfHarAllThirtyNineK}{0.50}  % ... at quoted bid/ask prices
\providecommand{\rhtaShRfFullAllThirtyNineK}{}\renewcommand{\rhtaShRfFullAllThirtyNineK}{0.95}  % Sharpe ratio (D19), mid-quotes: Random forest, all_features, ALL48_N39k (ALL48, the last 39,000 rows)
\providecommand{\rhtaShXRfFullAllThirtyNineK}{}\renewcommand{\rhtaShXRfFullAllThirtyNineK}{0.58}  % ... at quoted bid/ask prices
\providecommand{\rhtaPShRfAllThirtyNineK}{}\renewcommand{\rhtaPShRfAllThirtyNineK}{\ensuremath{p = 0.83}}  % Ledoit-Wolf p, Sharpe ratio full minus HAR (mid): Random forest, ALL48_N39k
\providecommand{\rhtaDShRidgeAllVsRth}{}\renewcommand{\rhtaDShRidgeAllVsRth}{\ensuremath{-}0.16}  % Sharpe ratio (D19, mid), ALL48 minus RTH13, full set, Ridge
\providecommand{\rhtaPDShRidgeAllVsRth}{}\renewcommand{\rhtaPDShRidgeAllVsRth}{\ensuremath{p = 0.63}}  % ... its Ledoit-Wolf p
\providecommand{\rhtaDShLassoAllVsRth}{}\renewcommand{\rhtaDShLassoAllVsRth}{\ensuremath{-}0.23}  % Sharpe ratio (D19, mid), ALL48 minus RTH13, full set, Lasso
\providecommand{\rhtaPDShLassoAllVsRth}{}\renewcommand{\rhtaPDShLassoAllVsRth}{\ensuremath{p = 0.51}}  % ... its Ledoit-Wolf p
\providecommand{\rhtaDShEnetAllVsRth}{}\renewcommand{\rhtaDShEnetAllVsRth}{\ensuremath{-}0.33}  % Sharpe ratio (D19, mid), ALL48 minus RTH13, full set, Elastic net
\providecommand{\rhtaPDShEnetAllVsRth}{}\renewcommand{\rhtaPDShEnetAllVsRth}{\ensuremath{p = 0.32}}  % ... its Ledoit-Wolf p
\providecommand{\rhtaDShLgbmAllVsRth}{}\renewcommand{\rhtaDShLgbmAllVsRth}{\ensuremath{-}0.34}  % Sharpe ratio (D19, mid), ALL48 minus RTH13, full set, LightGBM
\providecommand{\rhtaPDShLgbmAllVsRth}{}\renewcommand{\rhtaPDShLgbmAllVsRth}{\ensuremath{p = 0.17}}  % ... its Ledoit-Wolf p
\providecommand{\rhtaDShXgbAllVsRth}{}\renewcommand{\rhtaDShXgbAllVsRth}{0.03}  % Sharpe ratio (D19, mid), ALL48 minus RTH13, full set, XGBoost
\providecommand{\rhtaPDShXgbAllVsRth}{}\renewcommand{\rhtaPDShXgbAllVsRth}{\ensuremath{p = 0.90}}  % ... its Ledoit-Wolf p
\providecommand{\rhtaDShRfAllVsRth}{}\renewcommand{\rhtaDShRfAllVsRth}{\ensuremath{-}0.06}  % Sharpe ratio (D19, mid), ALL48 minus RTH13, full set, Random forest
\providecommand{\rhtaPDShRfAllVsRth}{}\renewcommand{\rhtaPDShRfAllVsRth}{\ensuremath{p = 0.85}}  % ... its Ledoit-Wolf p
\providecommand{\rhtaDShRidgeClVsRth}{}\renewcommand{\rhtaDShRidgeClVsRth}{\ensuremath{-}0.51}  % Sharpe ratio (D19, mid), CLOSE minus RTH13, full set, Ridge
\providecommand{\rhtaPDShRidgeClVsRth}{}\renewcommand{\rhtaPDShRidgeClVsRth}{\ensuremath{p = 0.22}}  % ... its Ledoit-Wolf p
\providecommand{\rhtaDShLassoClVsRth}{}\renewcommand{\rhtaDShLassoClVsRth}{\ensuremath{-}0.15}  % Sharpe ratio (D19, mid), CLOSE minus RTH13, full set, Lasso
\providecommand{\rhtaPDShLassoClVsRth}{}\renewcommand{\rhtaPDShLassoClVsRth}{\ensuremath{p = 0.72}}  % ... its Ledoit-Wolf p
\providecommand{\rhtaDShEnetClVsRth}{}\renewcommand{\rhtaDShEnetClVsRth}{\ensuremath{-}0.15}  % Sharpe ratio (D19, mid), CLOSE minus RTH13, full set, Elastic net
\providecommand{\rhtaPDShEnetClVsRth}{}\renewcommand{\rhtaPDShEnetClVsRth}{\ensuremath{p = 0.70}}  % ... its Ledoit-Wolf p
\providecommand{\rhtaDShLgbmClVsRth}{}\renewcommand{\rhtaDShLgbmClVsRth}{0.36}  % Sharpe ratio (D19, mid), CLOSE minus RTH13, full set, LightGBM
\providecommand{\rhtaPDShLgbmClVsRth}{}\renewcommand{\rhtaPDShLgbmClVsRth}{\ensuremath{p = 0.41}}  % ... its Ledoit-Wolf p
\providecommand{\rhtaDShXgbClVsRth}{}\renewcommand{\rhtaDShXgbClVsRth}{0.15}  % Sharpe ratio (D19, mid), CLOSE minus RTH13, full set, XGBoost
\providecommand{\rhtaPDShXgbClVsRth}{}\renewcommand{\rhtaPDShXgbClVsRth}{\ensuremath{p = 0.60}}  % ... its Ledoit-Wolf p
\providecommand{\rhtaDShRfClVsRth}{}\renewcommand{\rhtaDShRfClVsRth}{\ensuremath{-}0.59}  % Sharpe ratio (D19, mid), CLOSE minus RTH13, full set, Random forest
\providecommand{\rhtaPDShRfClVsRth}{}\renewcommand{\rhtaPDShRfClVsRth}{\ensuremath{p = 0.22}}  % ... its Ledoit-Wolf p
\providecommand{\rhtaWsPctRidgeWOne}{}\renewcommand{\rhtaWsPctRidgeWOne}{\ensuremath{-}4.8}  % CLOSE at W = 1,000 closes, 16:00, QLIKE % change full vs HAR, Ridge
\providecommand{\rhtaWsPRidgeWOne}{}\renewcommand{\rhtaWsPRidgeWOne}{\ensuremath{p < 0.001}}  % ... its p statement
\providecommand{\rhtaWsPctLassoWOne}{}\renewcommand{\rhtaWsPctLassoWOne}{\ensuremath{-}2.6}  % CLOSE at W = 1,000 closes, 16:00, QLIKE % change full vs HAR, Lasso
\providecommand{\rhtaWsPLassoWOne}{}\renewcommand{\rhtaWsPLassoWOne}{\ensuremath{p = 0.12}}  % ... its p statement
\providecommand{\rhtaWsPctEnetWOne}{}\renewcommand{\rhtaWsPctEnetWOne}{\ensuremath{-}1.5}  % CLOSE at W = 1,000 closes, 16:00, QLIKE % change full vs HAR, Elastic net
\providecommand{\rhtaWsPEnetWOne}{}\renewcommand{\rhtaWsPEnetWOne}{\ensuremath{p = 0.51}}  % ... its p statement
\providecommand{\rhtaWsPctLgbmWOne}{}\renewcommand{\rhtaWsPctLgbmWOne}{\ensuremath{-}7.8}  % CLOSE at W = 1,000 closes, 16:00, QLIKE % change full vs HAR, LightGBM
\providecommand{\rhtaWsPLgbmWOne}{}\renewcommand{\rhtaWsPLgbmWOne}{\ensuremath{p < 0.001}}  % ... its p statement
\providecommand{\rhtaWsPctXgbWOne}{}\renewcommand{\rhtaWsPctXgbWOne}{\ensuremath{-}6.8}  % CLOSE at W = 1,000 closes, 16:00, QLIKE % change full vs HAR, XGBoost
\providecommand{\rhtaWsPXgbWOne}{}\renewcommand{\rhtaWsPXgbWOne}{\ensuremath{p < 0.001}}  % ... its p statement
\providecommand{\rhtaWsPctRfWOne}{}\renewcommand{\rhtaWsPctRfWOne}{\ensuremath{-}2.1}  % CLOSE at W = 1,000 closes, 16:00, QLIKE % change full vs HAR, Random forest
\providecommand{\rhtaWsPRfWOne}{}\renewcommand{\rhtaWsPRfWOne}{\ensuremath{p = 0.069}}  % ... its p statement
\providecommand{\rhtaWsPctRidgeWOneHalf}{}\renewcommand{\rhtaWsPctRidgeWOneHalf}{\ensuremath{-}4.1}  % CLOSE at W = 1,500 closes, 16:00, QLIKE % change full vs HAR, Ridge
\providecommand{\rhtaWsPRidgeWOneHalf}{}\renewcommand{\rhtaWsPRidgeWOneHalf}{\ensuremath{p = 0.036}}  % ... its p statement
\providecommand{\rhtaWsPctLassoWOneHalf}{}\renewcommand{\rhtaWsPctLassoWOneHalf}{\ensuremath{-}5.0}  % CLOSE at W = 1,500 closes, 16:00, QLIKE % change full vs HAR, Lasso
\providecommand{\rhtaWsPLassoWOneHalf}{}\renewcommand{\rhtaWsPLassoWOneHalf}{\ensuremath{p < 0.001}}  % ... its p statement
\providecommand{\rhtaWsPctEnetWOneHalf}{}\renewcommand{\rhtaWsPctEnetWOneHalf}{\ensuremath{-}4.3}  % CLOSE at W = 1,500 closes, 16:00, QLIKE % change full vs HAR, Elastic net
\providecommand{\rhtaWsPEnetWOneHalf}{}\renewcommand{\rhtaWsPEnetWOneHalf}{\ensuremath{p = 0.035}}  % ... its p statement
\providecommand{\rhtaWsPctLgbmWOneHalf}{}\renewcommand{\rhtaWsPctLgbmWOneHalf}{\ensuremath{-}5.8}  % CLOSE at W = 1,500 closes, 16:00, QLIKE % change full vs HAR, LightGBM
\providecommand{\rhtaWsPLgbmWOneHalf}{}\renewcommand{\rhtaWsPLgbmWOneHalf}{\ensuremath{p = 0.001}}  % ... its p statement
\providecommand{\rhtaWsPctXgbWOneHalf}{}\renewcommand{\rhtaWsPctXgbWOneHalf}{\ensuremath{-}7.2}  % CLOSE at W = 1,500 closes, 16:00, QLIKE % change full vs HAR, XGBoost
\providecommand{\rhtaWsPXgbWOneHalf}{}\renewcommand{\rhtaWsPXgbWOneHalf}{\ensuremath{p < 0.001}}  % ... its p statement
\providecommand{\rhtaWsPctRfWOneHalf}{}\renewcommand{\rhtaWsPctRfWOneHalf}{\ensuremath{-}2.2}  % CLOSE at W = 1,500 closes, 16:00, QLIKE % change full vs HAR, Random forest
\providecommand{\rhtaWsPRfWOneHalf}{}\renewcommand{\rhtaWsPRfWOneHalf}{\ensuremath{p = 0.052}}  % ... its p statement
\providecommand{\rhtaWsPctRidgeWTwo}{}\renewcommand{\rhtaWsPctRidgeWTwo}{\ensuremath{-}3.9}  % CLOSE at W = 2,000 closes, 16:00, QLIKE % change full vs HAR, Ridge
\providecommand{\rhtaWsPRidgeWTwo}{}\renewcommand{\rhtaWsPRidgeWTwo}{\ensuremath{p = 0.049}}  % ... its p statement
\providecommand{\rhtaWsPctLassoWTwo}{}\renewcommand{\rhtaWsPctLassoWTwo}{\ensuremath{-}4.6}  % CLOSE at W = 2,000 closes, 16:00, QLIKE % change full vs HAR, Lasso
\providecommand{\rhtaWsPLassoWTwo}{}\renewcommand{\rhtaWsPLassoWTwo}{\ensuremath{p < 0.001}}  % ... its p statement
\providecommand{\rhtaWsPctEnetWTwo}{}\renewcommand{\rhtaWsPctEnetWTwo}{\ensuremath{-}4.6}  % CLOSE at W = 2,000 closes, 16:00, QLIKE % change full vs HAR, Elastic net
\providecommand{\rhtaWsPEnetWTwo}{}\renewcommand{\rhtaWsPEnetWTwo}{\ensuremath{p = 0.007}}  % ... its p statement
\providecommand{\rhtaWsPctLgbmWTwo}{}\renewcommand{\rhtaWsPctLgbmWTwo}{\ensuremath{-}7.4}  % CLOSE at W = 2,000 closes, 16:00, QLIKE % change full vs HAR, LightGBM
\providecommand{\rhtaWsPLgbmWTwo}{}\renewcommand{\rhtaWsPLgbmWTwo}{\ensuremath{p < 0.001}}  % ... its p statement
\providecommand{\rhtaWsPctXgbWTwo}{}\renewcommand{\rhtaWsPctXgbWTwo}{\ensuremath{-}7.8}  % CLOSE at W = 2,000 closes, 16:00, QLIKE % change full vs HAR, XGBoost
\providecommand{\rhtaWsPXgbWTwo}{}\renewcommand{\rhtaWsPXgbWTwo}{\ensuremath{p < 0.001}}  % ... its p statement
\providecommand{\rhtaWsPctRfWTwo}{}\renewcommand{\rhtaWsPctRfWTwo}{\ensuremath{-}3.2}  % CLOSE at W = 2,000 closes, 16:00, QLIKE % change full vs HAR, Random forest
\providecommand{\rhtaWsPRfWTwo}{}\renewcommand{\rhtaWsPRfWTwo}{\ensuremath{p = 0.002}}  % ... its p statement
\providecommand{\rhtaWsPctRidgeWTwoHalf}{}\renewcommand{\rhtaWsPctRidgeWTwoHalf}{\ensuremath{-}4.8}  % CLOSE at W = 2,500 closes, 16:00, QLIKE % change full vs HAR, Ridge
\providecommand{\rhtaWsPRidgeWTwoHalf}{}\renewcommand{\rhtaWsPRidgeWTwoHalf}{\ensuremath{p = 0.008}}  % ... its p statement
\providecommand{\rhtaWsPctLassoWTwoHalf}{}\renewcommand{\rhtaWsPctLassoWTwoHalf}{\ensuremath{-}5.3}  % CLOSE at W = 2,500 closes, 16:00, QLIKE % change full vs HAR, Lasso
\providecommand{\rhtaWsPLassoWTwoHalf}{}\renewcommand{\rhtaWsPLassoWTwoHalf}{\ensuremath{p < 0.001}}  % ... its p statement
\providecommand{\rhtaWsPctEnetWTwoHalf}{}\renewcommand{\rhtaWsPctEnetWTwoHalf}{\ensuremath{-}5.8}  % CLOSE at W = 2,500 closes, 16:00, QLIKE % change full vs HAR, Elastic net
\providecommand{\rhtaWsPEnetWTwoHalf}{}\renewcommand{\rhtaWsPEnetWTwoHalf}{\ensuremath{p < 0.001}}  % ... its p statement
\providecommand{\rhtaWsPctLgbmWTwoHalf}{}\renewcommand{\rhtaWsPctLgbmWTwoHalf}{\ensuremath{-}7.0}  % CLOSE at W = 2,500 closes, 16:00, QLIKE % change full vs HAR, LightGBM
\providecommand{\rhtaWsPLgbmWTwoHalf}{}\renewcommand{\rhtaWsPLgbmWTwoHalf}{\ensuremath{p < 0.001}}  % ... its p statement
\providecommand{\rhtaWsPctXgbWTwoHalf}{}\renewcommand{\rhtaWsPctXgbWTwoHalf}{\ensuremath{-}7.6}  % CLOSE at W = 2,500 closes, 16:00, QLIKE % change full vs HAR, XGBoost
\providecommand{\rhtaWsPXgbWTwoHalf}{}\renewcommand{\rhtaWsPXgbWTwoHalf}{\ensuremath{p < 0.001}}  % ... its p statement
\providecommand{\rhtaWsPctRfWTwoHalf}{}\renewcommand{\rhtaWsPctRfWTwoHalf}{\ensuremath{-}3.7}  % CLOSE at W = 2,500 closes, 16:00, QLIKE % change full vs HAR, Random forest
\providecommand{\rhtaWsPRfWTwoHalf}{}\renewcommand{\rhtaWsPRfWTwoHalf}{\ensuremath{p < 0.001}}  % ... its p statement
\providecommand{\rhtaWsPctRidgeWThree}{}\renewcommand{\rhtaWsPctRidgeWThree}{\ensuremath{-}6.1}  % CLOSE at W = 3,000 closes, 16:00, QLIKE % change full vs HAR, Ridge
\providecommand{\rhtaWsPRidgeWThree}{}\renewcommand{\rhtaWsPRidgeWThree}{\ensuremath{p < 0.001}}  % ... its p statement
\providecommand{\rhtaWsPctLassoWThree}{}\renewcommand{\rhtaWsPctLassoWThree}{\ensuremath{-}5.4}  % CLOSE at W = 3,000 closes, 16:00, QLIKE % change full vs HAR, Lasso
\providecommand{\rhtaWsPLassoWThree}{}\renewcommand{\rhtaWsPLassoWThree}{\ensuremath{p < 0.001}}  % ... its p statement
\providecommand{\rhtaWsPctEnetWThree}{}\renewcommand{\rhtaWsPctEnetWThree}{\ensuremath{-}6.5}  % CLOSE at W = 3,000 closes, 16:00, QLIKE % change full vs HAR, Elastic net
\providecommand{\rhtaWsPEnetWThree}{}\renewcommand{\rhtaWsPEnetWThree}{\ensuremath{p < 0.001}}  % ... its p statement
\providecommand{\rhtaWsPctLgbmWThree}{}\renewcommand{\rhtaWsPctLgbmWThree}{\ensuremath{-}7.0}  % CLOSE at W = 3,000 closes, 16:00, QLIKE % change full vs HAR, LightGBM
\providecommand{\rhtaWsPLgbmWThree}{}\renewcommand{\rhtaWsPLgbmWThree}{\ensuremath{p < 0.001}}  % ... its p statement
\providecommand{\rhtaWsPctXgbWThree}{}\renewcommand{\rhtaWsPctXgbWThree}{\ensuremath{-}8.5}  % CLOSE at W = 3,000 closes, 16:00, QLIKE % change full vs HAR, XGBoost
\providecommand{\rhtaWsPXgbWThree}{}\renewcommand{\rhtaWsPXgbWThree}{\ensuremath{p < 0.001}}  % ... its p statement
\providecommand{\rhtaWsPctRfWThree}{}\renewcommand{\rhtaWsPctRfWThree}{\ensuremath{-}3.7}  % CLOSE at W = 3,000 closes, 16:00, QLIKE % change full vs HAR, Random forest
\providecommand{\rhtaWsPRfWThree}{}\renewcommand{\rhtaWsPRfWThree}{\ensuremath{p < 0.001}}  % ... its p statement
\providecommand{\rhtaWsPctRidgeWThreeHalf}{}\renewcommand{\rhtaWsPctRidgeWThreeHalf}{\ensuremath{-}6.5}  % CLOSE at W = 3,500 closes, 16:00, QLIKE % change full vs HAR, Ridge
\providecommand{\rhtaWsPRidgeWThreeHalf}{}\renewcommand{\rhtaWsPRidgeWThreeHalf}{\ensuremath{p < 0.001}}  % ... its p statement
\providecommand{\rhtaWsPctLassoWThreeHalf}{}\renewcommand{\rhtaWsPctLassoWThreeHalf}{\ensuremath{-}5.1}  % CLOSE at W = 3,500 closes, 16:00, QLIKE % change full vs HAR, Lasso
\providecommand{\rhtaWsPLassoWThreeHalf}{}\renewcommand{\rhtaWsPLassoWThreeHalf}{\ensuremath{p = 0.011}}  % ... its p statement
\providecommand{\rhtaWsPctEnetWThreeHalf}{}\renewcommand{\rhtaWsPctEnetWThreeHalf}{\ensuremath{-}6.7}  % CLOSE at W = 3,500 closes, 16:00, QLIKE % change full vs HAR, Elastic net
\providecommand{\rhtaWsPEnetWThreeHalf}{}\renewcommand{\rhtaWsPEnetWThreeHalf}{\ensuremath{p < 0.001}}  % ... its p statement
\providecommand{\rhtaWsPctLgbmWThreeHalf}{}\renewcommand{\rhtaWsPctLgbmWThreeHalf}{\ensuremath{-}7.2}  % CLOSE at W = 3,500 closes, 16:00, QLIKE % change full vs HAR, LightGBM
\providecommand{\rhtaWsPLgbmWThreeHalf}{}\renewcommand{\rhtaWsPLgbmWThreeHalf}{\ensuremath{p < 0.001}}  % ... its p statement
\providecommand{\rhtaWsPctXgbWThreeHalf}{}\renewcommand{\rhtaWsPctXgbWThreeHalf}{\ensuremath{-}7.3}  % CLOSE at W = 3,500 closes, 16:00, QLIKE % change full vs HAR, XGBoost
\providecommand{\rhtaWsPXgbWThreeHalf}{}\renewcommand{\rhtaWsPXgbWThreeHalf}{\ensuremath{p < 0.001}}  % ... its p statement
\providecommand{\rhtaWsPctRfWThreeHalf}{}\renewcommand{\rhtaWsPctRfWThreeHalf}{\ensuremath{-}3.4}  % CLOSE at W = 3,500 closes, 16:00, QLIKE % change full vs HAR, Random forest
\providecommand{\rhtaWsPRfWThreeHalf}{}\renewcommand{\rhtaWsPRfWThreeHalf}{\ensuremath{p = 0.001}}  % ... its p statement
\providecommand{\rhtaLinPctAllCloseMin}{}\renewcommand{\rhtaLinPctAllCloseMin}{\ensuremath{-}2.6}  % 16:00, full vs HAR, ALL48, penalized linear models: lowest of the range, %
\providecommand{\rhtaLinPctAllCloseMax}{}\renewcommand{\rhtaLinPctAllCloseMax}{\ensuremath{-}1.8}  % 16:00, full vs HAR, ALL48, penalized linear models: highest of the range, %
\providecommand{\rhtaTreePctAllCloseMin}{}\renewcommand{\rhtaTreePctAllCloseMin}{\ensuremath{-}6.5}  % 16:00, full vs HAR, ALL48, tree ensembles: lowest of the range, %
\providecommand{\rhtaTreePctAllCloseMax}{}\renewcommand{\rhtaTreePctAllCloseMax}{\ensuremath{-}5.6}  % 16:00, full vs HAR, ALL48, tree ensembles: highest of the range, %
\providecommand{\rhtaPctAllCloseMin}{}\renewcommand{\rhtaPctAllCloseMin}{\ensuremath{-}6.5}  % 16:00, full vs HAR, ALL48, all six classes: lowest of the range, %
\providecommand{\rhtaPctAllCloseMax}{}\renewcommand{\rhtaPctAllCloseMax}{\ensuremath{-}1.8}  % 16:00, full vs HAR, ALL48, all six classes: highest of the range, %
\providecommand{\rhtaLinPctClCloseMin}{}\renewcommand{\rhtaLinPctClCloseMin}{\ensuremath{-}6.5}  % 16:00, full vs HAR, CLOSE, penalized linear models: lowest of the range, %
\providecommand{\rhtaLinPctClCloseMax}{}\renewcommand{\rhtaLinPctClCloseMax}{\ensuremath{-}5.4}  % 16:00, full vs HAR, CLOSE, penalized linear models: highest of the range, %
\providecommand{\rhtaTreePctClCloseMin}{}\renewcommand{\rhtaTreePctClCloseMin}{\ensuremath{-}8.5}  % 16:00, full vs HAR, CLOSE, tree ensembles: lowest of the range, %
\providecommand{\rhtaTreePctClCloseMax}{}\renewcommand{\rhtaTreePctClCloseMax}{\ensuremath{-}3.7}  % 16:00, full vs HAR, CLOSE, tree ensembles: highest of the range, %
\providecommand{\rhtaPctClCloseMin}{}\renewcommand{\rhtaPctClCloseMin}{\ensuremath{-}8.5}  % 16:00, full vs HAR, CLOSE, all six classes: lowest of the range, %
\providecommand{\rhtaPctClCloseMax}{}\renewcommand{\rhtaPctClCloseMax}{\ensuremath{-}3.7}  % 16:00, full vs HAR, CLOSE, all six classes: highest of the range, %
\providecommand{\rhtaLinPctRthCloseMin}{}\renewcommand{\rhtaLinPctRthCloseMin}{\ensuremath{-}0.5}  % 16:00, full vs HAR, RTH13, penalized linear models: lowest of the range, %
\providecommand{\rhtaLinPctRthCloseMax}{}\renewcommand{\rhtaLinPctRthCloseMax}{3.3}  % 16:00, full vs HAR, RTH13, penalized linear models: highest of the range, %
\providecommand{\rhtaTreePctRthCloseMin}{}\renewcommand{\rhtaTreePctRthCloseMin}{\ensuremath{-}3.4}  % 16:00, full vs HAR, RTH13, tree ensembles: lowest of the range, %
\providecommand{\rhtaTreePctRthCloseMax}{}\renewcommand{\rhtaTreePctRthCloseMax}{\ensuremath{-}1.4}  % 16:00, full vs HAR, RTH13, tree ensembles: highest of the range, %
\providecommand{\rhtaPctRthCloseMin}{}\renewcommand{\rhtaPctRthCloseMin}{\ensuremath{-}3.4}  % 16:00, full vs HAR, RTH13, all six classes: lowest of the range, %
\providecommand{\rhtaPctRthCloseMax}{}\renewcommand{\rhtaPctRthCloseMax}{3.3}  % 16:00, full vs HAR, RTH13, all six classes: highest of the range, %
\providecommand{\rhtaLinPctClNatCloseMin}{}\renewcommand{\rhtaLinPctClNatCloseMin}{\ensuremath{-}6.6}  % 16:00, full vs HAR, CLOSE_NAT, penalized linear models: lowest of the range, %
\providecommand{\rhtaLinPctClNatCloseMax}{}\renewcommand{\rhtaLinPctClNatCloseMax}{\ensuremath{-}5.4}  % 16:00, full vs HAR, CLOSE_NAT, penalized linear models: highest of the range, %
\providecommand{\rhtaLinAllClosePMin}{}\renewcommand{\rhtaLinAllClosePMin}{\ensuremath{p = 0.15}}  % 16:00, ALL48, linear: smallest p
\providecommand{\rhtaLinAllCtClosePMin}{}\renewcommand{\rhtaLinAllCtClosePMin}{\ensuremath{p = 0.33}}  % 16:00, ALL48@16, linear: smallest p
\providecommand{\rhtaLinRthCtClosePMin}{}\renewcommand{\rhtaLinRthCtClosePMin}{\ensuremath{p = 0.38}}  % 16:00, RTH13@16, linear: smallest p
\providecommand{\rhtaLinPctRthBarMin}{}\renewcommand{\rhtaLinPctRthBarMin}{\ensuremath{-}6.2}  % every bar, full vs HAR, RTH13, penalized linear models: lowest of the range, %
\providecommand{\rhtaLinPctRthBarMax}{}\renewcommand{\rhtaLinPctRthBarMax}{\ensuremath{-}5.4}  % every bar, full vs HAR, RTH13, penalized linear models: highest of the range, %
\providecommand{\rhtaTreePctRthBarMin}{}\renewcommand{\rhtaTreePctRthBarMin}{\ensuremath{-}14.0}  % every bar, full vs HAR, RTH13, tree ensembles: lowest of the range, %
\providecommand{\rhtaTreePctRthBarMax}{}\renewcommand{\rhtaTreePctRthBarMax}{\ensuremath{-}11.1}  % every bar, full vs HAR, RTH13, tree ensembles: highest of the range, %
\providecommand{\rhtaLinPctAllBarMin}{}\renewcommand{\rhtaLinPctAllBarMin}{\ensuremath{-}5.2}  % every bar, full vs HAR, ALL48, penalized linear models: lowest of the range, %
\providecommand{\rhtaLinPctAllBarMax}{}\renewcommand{\rhtaLinPctAllBarMax}{\ensuremath{-}4.9}  % every bar, full vs HAR, ALL48, penalized linear models: highest of the range, %
\providecommand{\rhtaTreePctAllBarMin}{}\renewcommand{\rhtaTreePctAllBarMin}{\ensuremath{-}12.0}  % every bar, full vs HAR, ALL48, tree ensembles: lowest of the range, %
\providecommand{\rhtaTreePctAllBarMax}{}\renewcommand{\rhtaTreePctAllBarMax}{\ensuremath{-}7.0}  % every bar, full vs HAR, ALL48, tree ensembles: highest of the range, %
\providecommand{\rhtaAllVsRthClosePctMin}{}\renewcommand{\rhtaAllVsRthClosePctMin}{\ensuremath{-}6.0}  % 16:00, ALL48 vs RTH13, full set, all six classes: lowest of the range, %
\providecommand{\rhtaAllVsRthClosePctMax}{}\renewcommand{\rhtaAllVsRthClosePctMax}{\ensuremath{-}1.9}  % 16:00, ALL48 vs RTH13, full set, all six classes: highest of the range, %
\providecommand{\rhtaNAllVsRthCloseSig}{}\renewcommand{\rhtaNAllVsRthCloseSig}{five}  % 16:00, ALL48 vs RTH13: classes with p < 0.05 (lower)
\providecommand{\rhtaAllVsRthBarPctMin}{}\renewcommand{\rhtaAllVsRthBarPctMin}{1.9}  % every bar, ALL48 vs RTH13, full set, all six classes (RTH13 more accurate): lowest of the range, %
\providecommand{\rhtaAllVsRthBarPctMax}{}\renewcommand{\rhtaAllVsRthBarPctMax}{5.2}  % every bar, ALL48 vs RTH13, full set, all six classes (RTH13 more accurate): highest of the range, %
\providecommand{\rhtaAllVsRthBarPMax}{}\renewcommand{\rhtaAllVsRthBarPMax}{\ensuremath{p = 0.008}}  % every bar, ALL48 vs RTH13: largest p of the six
\providecommand{\rhtaLinClVsRthClosePctMin}{}\renewcommand{\rhtaLinClVsRthClosePctMin}{\ensuremath{-}11.8}  % 16:00, CLOSE vs RTH13, full set, penalized linear models: lowest of the range, %
\providecommand{\rhtaLinClVsRthClosePctMax}{}\renewcommand{\rhtaLinClVsRthClosePctMax}{\ensuremath{-}9.4}  % 16:00, CLOSE vs RTH13, full set, penalized linear models: highest of the range, %
\providecommand{\rhtaTreeClVsRthClosePctMin}{}\renewcommand{\rhtaTreeClVsRthClosePctMin}{\ensuremath{-}2.2}  % 16:00, CLOSE vs RTH13, full set, tree ensembles: lowest of the range, %
\providecommand{\rhtaTreeClVsRthClosePctMax}{}\renewcommand{\rhtaTreeClVsRthClosePctMax}{1.6}  % 16:00, CLOSE vs RTH13, full set, tree ensembles: highest of the range, %
\providecommand{\rhtaTreeClVsRthClosePMin}{}\renewcommand{\rhtaTreeClVsRthClosePMin}{\ensuremath{p = 0.42}}  % ... tree ensembles: smallest p
\providecommand{\rhtaTreeClVsRthClosePMax}{}\renewcommand{\rhtaTreeClVsRthClosePMax}{\ensuremath{p = 0.69}}  % ... tree ensembles: largest p
\providecommand{\rhtaAllThirtyNineKVsRthClosePctMin}{}\renewcommand{\rhtaAllThirtyNineKVsRthClosePctMin}{\ensuremath{-}3.6}  % 16:00, ALL48_N39k vs RTH13, full set, all six classes: lowest of the range, %
\providecommand{\rhtaAllThirtyNineKVsRthClosePctMax}{}\renewcommand{\rhtaAllThirtyNineKVsRthClosePctMax}{1.3}  % 16:00, ALL48_N39k vs RTH13, full set, all six classes: highest of the range, %
\providecommand{\rhtaNAllThirtyNineKVsRthCloseSig}{}\renewcommand{\rhtaNAllThirtyNineKVsRthCloseSig}{two}  % ... classes with p < 0.05
\providecommand{\rhtaAllThirtyNineKVsRthBarPctMin}{}\renewcommand{\rhtaAllThirtyNineKVsRthBarPctMin}{2.7}  % every bar, ALL48_N39k vs RTH13, full set, all six classes: lowest of the range, %
\providecommand{\rhtaAllThirtyNineKVsRthBarPctMax}{}\renewcommand{\rhtaAllThirtyNineKVsRthBarPctMax}{10.4}  % every bar, ALL48_N39k vs RTH13, full set, all six classes: highest of the range, %
\providecommand{\rhtaAllThirtyNineKVsAllClosePctMin}{}\renewcommand{\rhtaAllThirtyNineKVsAllClosePctMin}{1.7}  % 16:00, ALL48_N39k vs ALL48 (39,000 vs 144,000 rows), full set: lowest of the range, %
\providecommand{\rhtaAllThirtyNineKVsAllClosePctMax}{}\renewcommand{\rhtaAllThirtyNineKVsAllClosePctMax}{4.1}  % 16:00, ALL48_N39k vs ALL48 (39,000 vs 144,000 rows), full set: highest of the range, %
\providecommand{\rhtaNAllThirtyNineKVsAllCloseSig}{}\renewcommand{\rhtaNAllThirtyNineKVsAllCloseSig}{five}  % ... classes with p < 0.05 (higher loss)
\providecommand{\rhtaAllThreeKVsRthThreeKClosePctMin}{}\renewcommand{\rhtaAllThreeKVsRthThreeKClosePctMin}{\ensuremath{-}8.1}  % 16:00, ALL48_N3k vs RTH13_N3k, full set, all six classes: lowest of the range, %
\providecommand{\rhtaAllThreeKVsRthThreeKClosePctMax}{}\renewcommand{\rhtaAllThreeKVsRthThreeKClosePctMax}{\ensuremath{-}3.3}  % 16:00, ALL48_N3k vs RTH13_N3k, full set, all six classes: highest of the range, %
\providecommand{\rhtaNAllThreeKVsRthThreeKCloseSig}{}\renewcommand{\rhtaNAllThreeKVsRthThreeKCloseSig}{five}  % ... classes with p < 0.05 (lower)
\providecommand{\rhtaAllThreeKVsRthThreeKBarPctMin}{}\renewcommand{\rhtaAllThreeKVsRthThreeKBarPctMin}{2.7}  % every bar, ALL48_N3k vs RTH13_N3k, full set, all six classes: lowest of the range, %
\providecommand{\rhtaAllThreeKVsRthThreeKBarPctMax}{}\renewcommand{\rhtaAllThreeKVsRthThreeKBarPctMax}{5.8}  % every bar, ALL48_N3k vs RTH13_N3k, full set, all six classes: highest of the range, %
\providecommand{\rhtaClVsRthThreeKClosePctMin}{}\renewcommand{\rhtaClVsRthThreeKClosePctMin}{\ensuremath{-}18.0}  % 16:00, CLOSE vs RTH13_N3k (3,000 rows each), full set: lowest of the range, %
\providecommand{\rhtaClVsRthThreeKClosePctMax}{}\renewcommand{\rhtaClVsRthThreeKClosePctMax}{\ensuremath{-}8.7}  % 16:00, CLOSE vs RTH13_N3k (3,000 rows each), full set: highest of the range, %
\providecommand{\rhtaClVsRthThreeKClosePMax}{}\renewcommand{\rhtaClVsRthThreeKClosePMax}{\ensuremath{p = 0.004}}  % ... largest p of the six
\providecommand{\rhtaRthThreeKVsRthClosePctMin}{}\renewcommand{\rhtaRthThreeKVsRthClosePctMin}{6.2}  % 16:00, RTH13_N3k vs RTH13 (3,000 vs 39,000 rows), full set: lowest of the range, %
\providecommand{\rhtaRthThreeKVsRthClosePctMax}{}\renewcommand{\rhtaRthThreeKVsRthClosePctMax}{12.3}  % 16:00, RTH13_N3k vs RTH13 (3,000 vs 39,000 rows), full set: highest of the range, %
\providecommand{\rhtaRthThreeKVsRthBarPctMin}{}\renewcommand{\rhtaRthThreeKVsRthBarPctMin}{2.7}  % every bar, RTH13_N3k vs RTH13 (3,000 vs 39,000 rows), full set: lowest of the range, %
\providecommand{\rhtaRthThreeKVsRthBarPctMax}{}\renewcommand{\rhtaRthThreeKVsRthBarPctMax}{14.9}  % every bar, RTH13_N3k vs RTH13 (3,000 vs 39,000 rows), full set: highest of the range, %
\providecommand{\rhtaShFullRthMin}{}\renewcommand{\rhtaShFullRthMin}{1.47}  % Sharpe ratio (D19, mid), full set, RTH13: lowest of the six
\providecommand{\rhtaShFullRthMax}{}\renewcommand{\rhtaShFullRthMax}{1.88}  % ... highest of the six
\providecommand{\rhtaNTreeShSigRth}{}\renewcommand{\rhtaNTreeShSigRth}{two}  % tree classes whose full-vs-HAR Sharpe ratio difference has Ledoit-Wolf p < 0.05, RTH13
\providecommand{\rhtaShFullAllMin}{}\renewcommand{\rhtaShFullAllMin}{1.25}  % Sharpe ratio (D19, mid), full set, ALL48: lowest of the six
\providecommand{\rhtaShFullAllMax}{}\renewcommand{\rhtaShFullAllMax}{1.82}  % ... highest of the six
\providecommand{\rhtaNTreeShSigAll}{}\renewcommand{\rhtaNTreeShSigAll}{none}  % tree classes whose full-vs-HAR Sharpe ratio difference has Ledoit-Wolf p < 0.05, ALL48
\providecommand{\rhtaShFullClMin}{}\renewcommand{\rhtaShFullClMin}{1.26}  % Sharpe ratio (D19, mid), full set, CLOSE: lowest of the six
\providecommand{\rhtaShFullClMax}{}\renewcommand{\rhtaShFullClMax}{1.94}  % ... highest of the six
\providecommand{\rhtaNTreeShSigCl}{}\renewcommand{\rhtaNTreeShSigCl}{one}  % tree classes whose full-vs-HAR Sharpe ratio difference has Ledoit-Wolf p < 0.05, CLOSE
\providecommand{\rhtaShTreeHarClMin}{}\renewcommand{\rhtaShTreeHarClMin}{\ensuremath{-}0.03}  % Sharpe ratio (D19, mid), HAR benchmark, CLOSE, tree ensembles: lowest
\providecommand{\rhtaShTreeHarClMax}{}\renewcommand{\rhtaShTreeHarClMax}{0.40}  % ... highest
\providecommand{\rhtaNDShTests}{}\renewcommand{\rhtaNDShTests}{12}  % Sharpe ratio tests ALL48 - RTH13 and CLOSE - RTH13, full set (six classes each)
\providecommand{\rhtaNDShSig}{}\renewcommand{\rhtaNDShSig}{none}  % ... of them with Ledoit-Wolf p < 0.05
\providecommand{\rhtaDShPMin}{}\renewcommand{\rhtaDShPMin}{\ensuremath{p = 0.17}}  % ... the smallest Ledoit-Wolf p
\providecommand{\rhtaWsNCells}{}\renewcommand{\rhtaWsNCells}{36}  % window x class cells of the close-only sweep (16:00, full vs HAR)
\providecommand{\rhtaWsNLower}{}\renewcommand{\rhtaWsNLower}{36}  % ... with the lower QLIKE on the full set
\providecommand{\rhtaWsNSig}{}\renewcommand{\rhtaWsNSig}{32}  % ... with p < 0.05
\providecommand{\rhtaWsSigFrom}{}\renewcommand{\rhtaWsSigFrom}{2{,}000}  % smallest W from which every cell has p < 0.05
\providecommand{\rhtaWsLinPctMin}{}\renewcommand{\rhtaWsLinPctMin}{\ensuremath{-}6.7}  % close-only sweep, 16:00, full vs HAR, linear, all W: lowest of the range, %
\providecommand{\rhtaWsLinPctMax}{}\renewcommand{\rhtaWsLinPctMax}{\ensuremath{-}1.5}  % close-only sweep, 16:00, full vs HAR, linear, all W: highest of the range, %
\providecommand{\rhtaWsTreePctMin}{}\renewcommand{\rhtaWsTreePctMin}{\ensuremath{-}8.5}  % close-only sweep, 16:00, full vs HAR, trees, all W: lowest of the range, %
\providecommand{\rhtaWsTreePctMax}{}\renewcommand{\rhtaWsTreePctMax}{\ensuremath{-}2.1}  % close-only sweep, 16:00, full vs HAR, trees, all W: highest of the range, %
\fi
\ifdefined\rhcaNdaysSpan\else% AUTO-GENERATED by writeup/make_rehaul_close_accuracy.py -- do not edit; re-run the script.
% \rhca... macros: closing-half-hour accuracy on both spans (tab:rh_close_accuracy) and the calibration
% diagnostics of the 12 pooled forecasts (tab:rh_calibration, fig:rh_calibration). Names: <Stat><Model><Set>
% <Sample><Span>; Model Ridge Lasso Enet Lgbm Xgb Rf Ols (pooled), Ci<Model> (trained on the closing half
% hour), Bm (closing-interval HAR benchmark); Set Har Full; Sample Close (16:00 bar) Bar (every bar);
% Span Span (2014-2024) Pre (2014-2019) Oos (T_oos, 2020-2024).
% --- spans and constants
\providecommand{\rhcaNdaysSpan}{}\renewcommand{\rhcaNdaysSpan}{2{,}579}  % trading days, 2014-01-02 .. 2024-04-30
\providecommand{\rhcaStartSpan}{}\renewcommand{\rhcaStartSpan}{2 January 2014}  % first day
\providecommand{\rhcaEndSpan}{}\renewcommand{\rhcaEndSpan}{30 April 2024}  % last day
\providecommand{\rhcaYearsSpan}{}\renewcommand{\rhcaYearsSpan}{2014--2024}  % years
\providecommand{\rhcaDMLagSpan}{}\renewcommand{\rhcaDMLagSpan}{8}  % Newey-West lag of the daily DM tests (dm_test automatic)
\providecommand{\rhcaHacLagBarSpan}{}\renewcommand{\rhcaHacLagBarSpan}{104}  % Newey-West lag of the every-bar MZ regressions, bars
\providecommand{\rhcaNdaysPre}{}\renewcommand{\rhcaNdaysPre}{1{,}496}  % trading days, 2014-01-02 .. 2019-12-31
\providecommand{\rhcaStartPre}{}\renewcommand{\rhcaStartPre}{2 January 2014}  % first day
\providecommand{\rhcaEndPre}{}\renewcommand{\rhcaEndPre}{31 December 2019}  % last day
\providecommand{\rhcaYearsPre}{}\renewcommand{\rhcaYearsPre}{2014--2019}  % years
\providecommand{\rhcaDMLagPre}{}\renewcommand{\rhcaDMLagPre}{7}  % Newey-West lag of the daily DM tests (dm_test automatic)
\providecommand{\rhcaHacLagBarPre}{}\renewcommand{\rhcaHacLagBarPre}{91}  % Newey-West lag of the every-bar MZ regressions, bars
\providecommand{\rhcaNdaysOos}{}\renewcommand{\rhcaNdaysOos}{1{,}082}  % trading days, 2020-01-03 .. 2024-04-30
\providecommand{\rhcaStartOos}{}\renewcommand{\rhcaStartOos}{3 January 2020}  % first day
\providecommand{\rhcaEndOos}{}\renewcommand{\rhcaEndOos}{30 April 2024}  % last day
\providecommand{\rhcaYearsOos}{}\renewcommand{\rhcaYearsOos}{2020--2024}  % years
\providecommand{\rhcaDMLagOos}{}\renewcommand{\rhcaDMLagOos}{6}  % Newey-West lag of the daily DM tests (dm_test automatic)
\providecommand{\rhcaHacLagBarOos}{}\renewcommand{\rhcaHacLagBarOos}{78}  % Newey-West lag of the every-bar MZ regressions, bars
\providecommand{\rhcaNbars}{}\renewcommand{\rhcaNbars}{33{,}527}  % scored bars of Table 1, 2014-2024
\providecommand{\rhcaCloseFirst}{}\renewcommand{\rhcaCloseFirst}{14 June 2006}  % first forecast of the closing-interval models
\providecommand{\rhcaMseExp}{}\renewcommand{\rhcaMseExp}{10}  % MSE printed x 10^this
\providecommand{\rhcaUnitExp}{}\renewcommand{\rhcaUnitExp}{6}  % MAE and the levels intercept printed x 10^this
\providecommand{\rhcaCrashStart}{}\renewcommand{\rhcaCrashStart}{24 February 2020}  % first day of the excluded window (make_rehaul_results.CRASH)
\providecommand{\rhcaCrashEnd}{}\renewcommand{\rhcaCrashEnd}{30 April 2020}  % last day of the excluded window
\providecommand{\rhcaTopDays}{}\renewcommand{\rhcaTopDays}{5}  % days in the extreme-day share
% --- closing half hour (16:00 bar): pooled forecasts, Span / Pre / Oos; Q QLIKE, Mse (x 10^10), Mae (x 10^6), Rsq R2_OOS %, Pct / Dm / P QLIKE full vs HAR, MsePct / MseDm / MseP, Gap / GapDm / GapP vs the closing-interval HAR benchmark (Oos only)
\providecommand{\rhcaQOlsHarCloseSpan}{}\renewcommand{\rhcaQOlsHarCloseSpan}{0.1269}
\providecommand{\rhcaMseOlsHarCloseSpan}{}\renewcommand{\rhcaMseOlsHarCloseSpan}{1.97}
\providecommand{\rhcaMaeOlsHarCloseSpan}{}\renewcommand{\rhcaMaeOlsHarCloseSpan}{2.92}
\providecommand{\rhcaRsqOlsHarCloseSpan}{}\renewcommand{\rhcaRsqOlsHarCloseSpan}{77.4}
\providecommand{\rhcaQRidgeHarCloseSpan}{}\renewcommand{\rhcaQRidgeHarCloseSpan}{0.1269}
\providecommand{\rhcaMseRidgeHarCloseSpan}{}\renewcommand{\rhcaMseRidgeHarCloseSpan}{1.97}
\providecommand{\rhcaMaeRidgeHarCloseSpan}{}\renewcommand{\rhcaMaeRidgeHarCloseSpan}{2.92}
\providecommand{\rhcaRsqRidgeHarCloseSpan}{}\renewcommand{\rhcaRsqRidgeHarCloseSpan}{77.4}
\providecommand{\rhcaQRidgeFullCloseSpan}{}\renewcommand{\rhcaQRidgeFullCloseSpan}{0.1291}
\providecommand{\rhcaMseRidgeFullCloseSpan}{}\renewcommand{\rhcaMseRidgeFullCloseSpan}{3.49}
\providecommand{\rhcaMaeRidgeFullCloseSpan}{}\renewcommand{\rhcaMaeRidgeFullCloseSpan}{3.22}
\providecommand{\rhcaRsqRidgeFullCloseSpan}{}\renewcommand{\rhcaRsqRidgeFullCloseSpan}{59.9}
\providecommand{\rhcaPctRidgeCloseSpan}{}\renewcommand{\rhcaPctRidgeCloseSpan}{1.8}  % QLIKE %, full vs HAR of the same class
\providecommand{\rhcaDmRidgeCloseSpan}{}\renewcommand{\rhcaDmRidgeCloseSpan}{0.67}
\providecommand{\rhcaPRidgeCloseSpan}{}\renewcommand{\rhcaPRidgeCloseSpan}{\ensuremath{p = 0.50}}
\providecommand{\rhcaMsePctRidgeCloseSpan}{}\renewcommand{\rhcaMsePctRidgeCloseSpan}{77.5}  % MSE %, full vs HAR
\providecommand{\rhcaMseDmRidgeCloseSpan}{}\renewcommand{\rhcaMseDmRidgeCloseSpan}{1.18}
\providecommand{\rhcaMsePRidgeCloseSpan}{}\renewcommand{\rhcaMsePRidgeCloseSpan}{\ensuremath{p = 0.24}}
\providecommand{\rhcaQLassoHarCloseSpan}{}\renewcommand{\rhcaQLassoHarCloseSpan}{0.1269}
\providecommand{\rhcaMseLassoHarCloseSpan}{}\renewcommand{\rhcaMseLassoHarCloseSpan}{1.97}
\providecommand{\rhcaMaeLassoHarCloseSpan}{}\renewcommand{\rhcaMaeLassoHarCloseSpan}{2.92}
\providecommand{\rhcaRsqLassoHarCloseSpan}{}\renewcommand{\rhcaRsqLassoHarCloseSpan}{77.4}
\providecommand{\rhcaQLassoFullCloseSpan}{}\renewcommand{\rhcaQLassoFullCloseSpan}{0.1322}
\providecommand{\rhcaMseLassoFullCloseSpan}{}\renewcommand{\rhcaMseLassoFullCloseSpan}{3.57}
\providecommand{\rhcaMaeLassoFullCloseSpan}{}\renewcommand{\rhcaMaeLassoFullCloseSpan}{3.24}
\providecommand{\rhcaRsqLassoFullCloseSpan}{}\renewcommand{\rhcaRsqLassoFullCloseSpan}{59.0}
\providecommand{\rhcaPctLassoCloseSpan}{}\renewcommand{\rhcaPctLassoCloseSpan}{4.2}  % QLIKE %, full vs HAR of the same class
\providecommand{\rhcaDmLassoCloseSpan}{}\renewcommand{\rhcaDmLassoCloseSpan}{1.35}
\providecommand{\rhcaPLassoCloseSpan}{}\renewcommand{\rhcaPLassoCloseSpan}{\ensuremath{p = 0.18}}
\providecommand{\rhcaMsePctLassoCloseSpan}{}\renewcommand{\rhcaMsePctLassoCloseSpan}{81.5}  % MSE %, full vs HAR
\providecommand{\rhcaMseDmLassoCloseSpan}{}\renewcommand{\rhcaMseDmLassoCloseSpan}{1.21}
\providecommand{\rhcaMsePLassoCloseSpan}{}\renewcommand{\rhcaMsePLassoCloseSpan}{\ensuremath{p = 0.23}}
\providecommand{\rhcaQEnetHarCloseSpan}{}\renewcommand{\rhcaQEnetHarCloseSpan}{0.1269}
\providecommand{\rhcaMseEnetHarCloseSpan}{}\renewcommand{\rhcaMseEnetHarCloseSpan}{1.97}
\providecommand{\rhcaMaeEnetHarCloseSpan}{}\renewcommand{\rhcaMaeEnetHarCloseSpan}{2.92}
\providecommand{\rhcaRsqEnetHarCloseSpan}{}\renewcommand{\rhcaRsqEnetHarCloseSpan}{77.4}
\providecommand{\rhcaQEnetFullCloseSpan}{}\renewcommand{\rhcaQEnetFullCloseSpan}{0.1276}
\providecommand{\rhcaMseEnetFullCloseSpan}{}\renewcommand{\rhcaMseEnetFullCloseSpan}{3.92}
\providecommand{\rhcaMaeEnetFullCloseSpan}{}\renewcommand{\rhcaMaeEnetFullCloseSpan}{3.28}
\providecommand{\rhcaRsqEnetFullCloseSpan}{}\renewcommand{\rhcaRsqEnetFullCloseSpan}{55.0}
\providecommand{\rhcaPctEnetCloseSpan}{}\renewcommand{\rhcaPctEnetCloseSpan}{0.6}  % QLIKE %, full vs HAR of the same class
\providecommand{\rhcaDmEnetCloseSpan}{}\renewcommand{\rhcaDmEnetCloseSpan}{0.24}
\providecommand{\rhcaPEnetCloseSpan}{}\renewcommand{\rhcaPEnetCloseSpan}{\ensuremath{p = 0.81}}
\providecommand{\rhcaMsePctEnetCloseSpan}{}\renewcommand{\rhcaMsePctEnetCloseSpan}{99.1}  % MSE %, full vs HAR
\providecommand{\rhcaMseDmEnetCloseSpan}{}\renewcommand{\rhcaMseDmEnetCloseSpan}{1.16}
\providecommand{\rhcaMsePEnetCloseSpan}{}\renewcommand{\rhcaMsePEnetCloseSpan}{\ensuremath{p = 0.25}}
\providecommand{\rhcaQLgbmHarCloseSpan}{}\renewcommand{\rhcaQLgbmHarCloseSpan}{0.1235}
\providecommand{\rhcaMseLgbmHarCloseSpan}{}\renewcommand{\rhcaMseLgbmHarCloseSpan}{1.77}
\providecommand{\rhcaMaeLgbmHarCloseSpan}{}\renewcommand{\rhcaMaeLgbmHarCloseSpan}{3.01}
\providecommand{\rhcaRsqLgbmHarCloseSpan}{}\renewcommand{\rhcaRsqLgbmHarCloseSpan}{79.6}
\providecommand{\rhcaQLgbmFullCloseSpan}{}\renewcommand{\rhcaQLgbmFullCloseSpan}{0.1202}
\providecommand{\rhcaMseLgbmFullCloseSpan}{}\renewcommand{\rhcaMseLgbmFullCloseSpan}{2.36}
\providecommand{\rhcaMaeLgbmFullCloseSpan}{}\renewcommand{\rhcaMaeLgbmFullCloseSpan}{3.02}
\providecommand{\rhcaRsqLgbmFullCloseSpan}{}\renewcommand{\rhcaRsqLgbmFullCloseSpan}{72.9}
\providecommand{\rhcaPctLgbmCloseSpan}{}\renewcommand{\rhcaPctLgbmCloseSpan}{\ensuremath{-}2.7}  % QLIKE %, full vs HAR of the same class
\providecommand{\rhcaDmLgbmCloseSpan}{}\renewcommand{\rhcaDmLgbmCloseSpan}{\ensuremath{-}1.76}
\providecommand{\rhcaPLgbmCloseSpan}{}\renewcommand{\rhcaPLgbmCloseSpan}{\ensuremath{p = 0.079}}
\providecommand{\rhcaMsePctLgbmCloseSpan}{}\renewcommand{\rhcaMsePctLgbmCloseSpan}{33.0}  % MSE %, full vs HAR
\providecommand{\rhcaMseDmLgbmCloseSpan}{}\renewcommand{\rhcaMseDmLgbmCloseSpan}{1.10}
\providecommand{\rhcaMsePLgbmCloseSpan}{}\renewcommand{\rhcaMsePLgbmCloseSpan}{\ensuremath{p = 0.27}}
\providecommand{\rhcaQXgbHarCloseSpan}{}\renewcommand{\rhcaQXgbHarCloseSpan}{0.1234}
\providecommand{\rhcaMseXgbHarCloseSpan}{}\renewcommand{\rhcaMseXgbHarCloseSpan}{1.73}
\providecommand{\rhcaMaeXgbHarCloseSpan}{}\renewcommand{\rhcaMaeXgbHarCloseSpan}{2.95}
\providecommand{\rhcaRsqXgbHarCloseSpan}{}\renewcommand{\rhcaRsqXgbHarCloseSpan}{80.1}
\providecommand{\rhcaQXgbFullCloseSpan}{}\renewcommand{\rhcaQXgbFullCloseSpan}{0.1195}
\providecommand{\rhcaMseXgbFullCloseSpan}{}\renewcommand{\rhcaMseXgbFullCloseSpan}{2.19}
\providecommand{\rhcaMaeXgbFullCloseSpan}{}\renewcommand{\rhcaMaeXgbFullCloseSpan}{3.01}
\providecommand{\rhcaRsqXgbFullCloseSpan}{}\renewcommand{\rhcaRsqXgbFullCloseSpan}{74.9}
\providecommand{\rhcaPctXgbCloseSpan}{}\renewcommand{\rhcaPctXgbCloseSpan}{\ensuremath{-}3.1}  % QLIKE %, full vs HAR of the same class
\providecommand{\rhcaDmXgbCloseSpan}{}\renewcommand{\rhcaDmXgbCloseSpan}{\ensuremath{-}2.01}
\providecommand{\rhcaPXgbCloseSpan}{}\renewcommand{\rhcaPXgbCloseSpan}{\ensuremath{p = 0.044}}
\providecommand{\rhcaMsePctXgbCloseSpan}{}\renewcommand{\rhcaMsePctXgbCloseSpan}{26.3}  % MSE %, full vs HAR
\providecommand{\rhcaMseDmXgbCloseSpan}{}\renewcommand{\rhcaMseDmXgbCloseSpan}{1.09}
\providecommand{\rhcaMsePXgbCloseSpan}{}\renewcommand{\rhcaMsePXgbCloseSpan}{\ensuremath{p = 0.28}}
\providecommand{\rhcaQRfHarCloseSpan}{}\renewcommand{\rhcaQRfHarCloseSpan}{0.1266}
\providecommand{\rhcaMseRfHarCloseSpan}{}\renewcommand{\rhcaMseRfHarCloseSpan}{1.63}
\providecommand{\rhcaMaeRfHarCloseSpan}{}\renewcommand{\rhcaMaeRfHarCloseSpan}{2.91}
\providecommand{\rhcaRsqRfHarCloseSpan}{}\renewcommand{\rhcaRsqRfHarCloseSpan}{81.3}
\providecommand{\rhcaQRfFullCloseSpan}{}\renewcommand{\rhcaQRfFullCloseSpan}{0.1246}
\providecommand{\rhcaMseRfFullCloseSpan}{}\renewcommand{\rhcaMseRfFullCloseSpan}{1.95}
\providecommand{\rhcaMaeRfFullCloseSpan}{}\renewcommand{\rhcaMaeRfFullCloseSpan}{2.97}
\providecommand{\rhcaRsqRfFullCloseSpan}{}\renewcommand{\rhcaRsqRfFullCloseSpan}{77.6}
\providecommand{\rhcaPctRfCloseSpan}{}\renewcommand{\rhcaPctRfCloseSpan}{\ensuremath{-}1.6}  % QLIKE %, full vs HAR of the same class
\providecommand{\rhcaDmRfCloseSpan}{}\renewcommand{\rhcaDmRfCloseSpan}{\ensuremath{-}1.03}
\providecommand{\rhcaPRfCloseSpan}{}\renewcommand{\rhcaPRfCloseSpan}{\ensuremath{p = 0.30}}
\providecommand{\rhcaMsePctRfCloseSpan}{}\renewcommand{\rhcaMsePctRfCloseSpan}{19.7}  % MSE %, full vs HAR
\providecommand{\rhcaMseDmRfCloseSpan}{}\renewcommand{\rhcaMseDmRfCloseSpan}{1.11}
\providecommand{\rhcaMsePRfCloseSpan}{}\renewcommand{\rhcaMsePRfCloseSpan}{\ensuremath{p = 0.27}}
\providecommand{\rhcaQOlsHarClosePre}{}\renewcommand{\rhcaQOlsHarClosePre}{0.1365}
\providecommand{\rhcaMseOlsHarClosePre}{}\renewcommand{\rhcaMseOlsHarClosePre}{0.32}
\providecommand{\rhcaMaeOlsHarClosePre}{}\renewcommand{\rhcaMaeOlsHarClosePre}{1.74}
\providecommand{\rhcaRsqOlsHarClosePre}{}\renewcommand{\rhcaRsqOlsHarClosePre}{73.9}
\providecommand{\rhcaQRidgeHarClosePre}{}\renewcommand{\rhcaQRidgeHarClosePre}{0.1365}
\providecommand{\rhcaMseRidgeHarClosePre}{}\renewcommand{\rhcaMseRidgeHarClosePre}{0.32}
\providecommand{\rhcaMaeRidgeHarClosePre}{}\renewcommand{\rhcaMaeRidgeHarClosePre}{1.74}
\providecommand{\rhcaRsqRidgeHarClosePre}{}\renewcommand{\rhcaRsqRidgeHarClosePre}{73.9}
\providecommand{\rhcaQRidgeFullClosePre}{}\renewcommand{\rhcaQRidgeFullClosePre}{0.1309}
\providecommand{\rhcaMseRidgeFullClosePre}{}\renewcommand{\rhcaMseRidgeFullClosePre}{0.28}
\providecommand{\rhcaMaeRidgeFullClosePre}{}\renewcommand{\rhcaMaeRidgeFullClosePre}{1.66}
\providecommand{\rhcaRsqRidgeFullClosePre}{}\renewcommand{\rhcaRsqRidgeFullClosePre}{77.3}
\providecommand{\rhcaPctRidgeClosePre}{}\renewcommand{\rhcaPctRidgeClosePre}{\ensuremath{-}4.1}  % QLIKE %, full vs HAR of the same class
\providecommand{\rhcaDmRidgeClosePre}{}\renewcommand{\rhcaDmRidgeClosePre}{\ensuremath{-}1.42}
\providecommand{\rhcaPRidgeClosePre}{}\renewcommand{\rhcaPRidgeClosePre}{\ensuremath{p = 0.16}}
\providecommand{\rhcaMsePctRidgeClosePre}{}\renewcommand{\rhcaMsePctRidgeClosePre}{\ensuremath{-}13.1}  % MSE %, full vs HAR
\providecommand{\rhcaMseDmRidgeClosePre}{}\renewcommand{\rhcaMseDmRidgeClosePre}{\ensuremath{-}1.79}
\providecommand{\rhcaMsePRidgeClosePre}{}\renewcommand{\rhcaMsePRidgeClosePre}{\ensuremath{p = 0.074}}
\providecommand{\rhcaQLassoHarClosePre}{}\renewcommand{\rhcaQLassoHarClosePre}{0.1365}
\providecommand{\rhcaMseLassoHarClosePre}{}\renewcommand{\rhcaMseLassoHarClosePre}{0.32}
\providecommand{\rhcaMaeLassoHarClosePre}{}\renewcommand{\rhcaMaeLassoHarClosePre}{1.74}
\providecommand{\rhcaRsqLassoHarClosePre}{}\renewcommand{\rhcaRsqLassoHarClosePre}{73.9}
\providecommand{\rhcaQLassoFullClosePre}{}\renewcommand{\rhcaQLassoFullClosePre}{0.1333}
\providecommand{\rhcaMseLassoFullClosePre}{}\renewcommand{\rhcaMseLassoFullClosePre}{0.27}
\providecommand{\rhcaMaeLassoFullClosePre}{}\renewcommand{\rhcaMaeLassoFullClosePre}{1.65}
\providecommand{\rhcaRsqLassoFullClosePre}{}\renewcommand{\rhcaRsqLassoFullClosePre}{78.0}
\providecommand{\rhcaPctLassoClosePre}{}\renewcommand{\rhcaPctLassoClosePre}{\ensuremath{-}2.3}  % QLIKE %, full vs HAR of the same class
\providecommand{\rhcaDmLassoClosePre}{}\renewcommand{\rhcaDmLassoClosePre}{\ensuremath{-}0.67}
\providecommand{\rhcaPLassoClosePre}{}\renewcommand{\rhcaPLassoClosePre}{\ensuremath{p = 0.50}}
\providecommand{\rhcaMsePctLassoClosePre}{}\renewcommand{\rhcaMsePctLassoClosePre}{\ensuremath{-}15.8}  % MSE %, full vs HAR
\providecommand{\rhcaMseDmLassoClosePre}{}\renewcommand{\rhcaMseDmLassoClosePre}{\ensuremath{-}1.78}
\providecommand{\rhcaMsePLassoClosePre}{}\renewcommand{\rhcaMsePLassoClosePre}{\ensuremath{p = 0.075}}
\providecommand{\rhcaQEnetHarClosePre}{}\renewcommand{\rhcaQEnetHarClosePre}{0.1365}
\providecommand{\rhcaMseEnetHarClosePre}{}\renewcommand{\rhcaMseEnetHarClosePre}{0.32}
\providecommand{\rhcaMaeEnetHarClosePre}{}\renewcommand{\rhcaMaeEnetHarClosePre}{1.74}
\providecommand{\rhcaRsqEnetHarClosePre}{}\renewcommand{\rhcaRsqEnetHarClosePre}{73.9}
\providecommand{\rhcaQEnetFullClosePre}{}\renewcommand{\rhcaQEnetFullClosePre}{0.1300}
\providecommand{\rhcaMseEnetFullClosePre}{}\renewcommand{\rhcaMseEnetFullClosePre}{0.28}
\providecommand{\rhcaMaeEnetFullClosePre}{}\renewcommand{\rhcaMaeEnetFullClosePre}{1.64}
\providecommand{\rhcaRsqEnetFullClosePre}{}\renewcommand{\rhcaRsqEnetFullClosePre}{77.6}
\providecommand{\rhcaPctEnetClosePre}{}\renewcommand{\rhcaPctEnetClosePre}{\ensuremath{-}4.8}  % QLIKE %, full vs HAR of the same class
\providecommand{\rhcaDmEnetClosePre}{}\renewcommand{\rhcaDmEnetClosePre}{\ensuremath{-}1.80}
\providecommand{\rhcaPEnetClosePre}{}\renewcommand{\rhcaPEnetClosePre}{\ensuremath{p = 0.071}}
\providecommand{\rhcaMsePctEnetClosePre}{}\renewcommand{\rhcaMsePctEnetClosePre}{\ensuremath{-}14.4}  % MSE %, full vs HAR
\providecommand{\rhcaMseDmEnetClosePre}{}\renewcommand{\rhcaMseDmEnetClosePre}{\ensuremath{-}1.90}
\providecommand{\rhcaMsePEnetClosePre}{}\renewcommand{\rhcaMsePEnetClosePre}{\ensuremath{p = 0.058}}
\providecommand{\rhcaQLgbmHarClosePre}{}\renewcommand{\rhcaQLgbmHarClosePre}{0.1333}
\providecommand{\rhcaMseLgbmHarClosePre}{}\renewcommand{\rhcaMseLgbmHarClosePre}{0.40}
\providecommand{\rhcaMaeLgbmHarClosePre}{}\renewcommand{\rhcaMaeLgbmHarClosePre}{1.84}
\providecommand{\rhcaRsqLgbmHarClosePre}{}\renewcommand{\rhcaRsqLgbmHarClosePre}{68.1}
\providecommand{\rhcaQLgbmFullClosePre}{}\renewcommand{\rhcaQLgbmFullClosePre}{0.1287}
\providecommand{\rhcaMseLgbmFullClosePre}{}\renewcommand{\rhcaMseLgbmFullClosePre}{0.37}
\providecommand{\rhcaMaeLgbmFullClosePre}{}\renewcommand{\rhcaMaeLgbmFullClosePre}{1.74}
\providecommand{\rhcaRsqLgbmFullClosePre}{}\renewcommand{\rhcaRsqLgbmFullClosePre}{70.5}
\providecommand{\rhcaPctLgbmClosePre}{}\renewcommand{\rhcaPctLgbmClosePre}{\ensuremath{-}3.4}  % QLIKE %, full vs HAR of the same class
\providecommand{\rhcaDmLgbmClosePre}{}\renewcommand{\rhcaDmLgbmClosePre}{\ensuremath{-}1.99}
\providecommand{\rhcaPLgbmClosePre}{}\renewcommand{\rhcaPLgbmClosePre}{\ensuremath{p = 0.046}}
\providecommand{\rhcaMsePctLgbmClosePre}{}\renewcommand{\rhcaMsePctLgbmClosePre}{\ensuremath{-}7.5}  % MSE %, full vs HAR
\providecommand{\rhcaMseDmLgbmClosePre}{}\renewcommand{\rhcaMseDmLgbmClosePre}{\ensuremath{-}1.78}
\providecommand{\rhcaMsePLgbmClosePre}{}\renewcommand{\rhcaMsePLgbmClosePre}{\ensuremath{p = 0.076}}
\providecommand{\rhcaQXgbHarClosePre}{}\renewcommand{\rhcaQXgbHarClosePre}{0.1331}
\providecommand{\rhcaMseXgbHarClosePre}{}\renewcommand{\rhcaMseXgbHarClosePre}{0.37}
\providecommand{\rhcaMaeXgbHarClosePre}{}\renewcommand{\rhcaMaeXgbHarClosePre}{1.82}
\providecommand{\rhcaRsqXgbHarClosePre}{}\renewcommand{\rhcaRsqXgbHarClosePre}{70.0}
\providecommand{\rhcaQXgbFullClosePre}{}\renewcommand{\rhcaQXgbFullClosePre}{0.1277}
\providecommand{\rhcaMseXgbFullClosePre}{}\renewcommand{\rhcaMseXgbFullClosePre}{0.33}
\providecommand{\rhcaMaeXgbFullClosePre}{}\renewcommand{\rhcaMaeXgbFullClosePre}{1.74}
\providecommand{\rhcaRsqXgbFullClosePre}{}\renewcommand{\rhcaRsqXgbFullClosePre}{73.0}
\providecommand{\rhcaPctXgbClosePre}{}\renewcommand{\rhcaPctXgbClosePre}{\ensuremath{-}4.0}  % QLIKE %, full vs HAR of the same class
\providecommand{\rhcaDmXgbClosePre}{}\renewcommand{\rhcaDmXgbClosePre}{\ensuremath{-}2.13}
\providecommand{\rhcaPXgbClosePre}{}\renewcommand{\rhcaPXgbClosePre}{\ensuremath{p = 0.033}}
\providecommand{\rhcaMsePctXgbClosePre}{}\renewcommand{\rhcaMsePctXgbClosePre}{\ensuremath{-}10.0}  % MSE %, full vs HAR
\providecommand{\rhcaMseDmXgbClosePre}{}\renewcommand{\rhcaMseDmXgbClosePre}{\ensuremath{-}1.66}
\providecommand{\rhcaMsePXgbClosePre}{}\renewcommand{\rhcaMsePXgbClosePre}{\ensuremath{p = 0.097}}
\providecommand{\rhcaQRfHarClosePre}{}\renewcommand{\rhcaQRfHarClosePre}{0.1364}
\providecommand{\rhcaMseRfHarClosePre}{}\renewcommand{\rhcaMseRfHarClosePre}{0.36}
\providecommand{\rhcaMaeRfHarClosePre}{}\renewcommand{\rhcaMaeRfHarClosePre}{1.81}
\providecommand{\rhcaRsqRfHarClosePre}{}\renewcommand{\rhcaRsqRfHarClosePre}{70.7}
\providecommand{\rhcaQRfFullClosePre}{}\renewcommand{\rhcaQRfFullClosePre}{0.1335}
\providecommand{\rhcaMseRfFullClosePre}{}\renewcommand{\rhcaMseRfFullClosePre}{0.34}
\providecommand{\rhcaMaeRfFullClosePre}{}\renewcommand{\rhcaMaeRfFullClosePre}{1.77}
\providecommand{\rhcaRsqRfFullClosePre}{}\renewcommand{\rhcaRsqRfFullClosePre}{72.3}
\providecommand{\rhcaPctRfClosePre}{}\renewcommand{\rhcaPctRfClosePre}{\ensuremath{-}2.1}  % QLIKE %, full vs HAR of the same class
\providecommand{\rhcaDmRfClosePre}{}\renewcommand{\rhcaDmRfClosePre}{\ensuremath{-}1.14}
\providecommand{\rhcaPRfClosePre}{}\renewcommand{\rhcaPRfClosePre}{\ensuremath{p = 0.25}}
\providecommand{\rhcaMsePctRfClosePre}{}\renewcommand{\rhcaMsePctRfClosePre}{\ensuremath{-}5.3}  % MSE %, full vs HAR
\providecommand{\rhcaMseDmRfClosePre}{}\renewcommand{\rhcaMseDmRfClosePre}{\ensuremath{-}1.39}
\providecommand{\rhcaMsePRfClosePre}{}\renewcommand{\rhcaMsePRfClosePre}{\ensuremath{p = 0.17}}
\providecommand{\rhcaQOlsHarCloseOos}{}\renewcommand{\rhcaQOlsHarCloseOos}{0.1137}
\providecommand{\rhcaMseOlsHarCloseOos}{}\renewcommand{\rhcaMseOlsHarCloseOos}{4.24}
\providecommand{\rhcaMaeOlsHarCloseOos}{}\renewcommand{\rhcaMaeOlsHarCloseOos}{4.55}
\providecommand{\rhcaRsqOlsHarCloseOos}{}\renewcommand{\rhcaRsqOlsHarCloseOos}{77.7}
\providecommand{\rhcaGapOlsHarCloseOos}{}\renewcommand{\rhcaGapOlsHarCloseOos}{10.9}  % QLIKE % vs the closing-interval HAR benchmark
\providecommand{\rhcaGapDmOlsHarCloseOos}{}\renewcommand{\rhcaGapDmOlsHarCloseOos}{2.84}
\providecommand{\rhcaGapPOlsHarCloseOos}{}\renewcommand{\rhcaGapPOlsHarCloseOos}{\ensuremath{p = 0.004}}
\providecommand{\rhcaQRidgeHarCloseOos}{}\renewcommand{\rhcaQRidgeHarCloseOos}{0.1137}
\providecommand{\rhcaMseRidgeHarCloseOos}{}\renewcommand{\rhcaMseRidgeHarCloseOos}{4.24}
\providecommand{\rhcaMaeRidgeHarCloseOos}{}\renewcommand{\rhcaMaeRidgeHarCloseOos}{4.55}
\providecommand{\rhcaRsqRidgeHarCloseOos}{}\renewcommand{\rhcaRsqRidgeHarCloseOos}{77.7}
\providecommand{\rhcaGapRidgeHarCloseOos}{}\renewcommand{\rhcaGapRidgeHarCloseOos}{10.9}  % QLIKE % vs the closing-interval HAR benchmark
\providecommand{\rhcaGapDmRidgeHarCloseOos}{}\renewcommand{\rhcaGapDmRidgeHarCloseOos}{2.84}
\providecommand{\rhcaGapPRidgeHarCloseOos}{}\renewcommand{\rhcaGapPRidgeHarCloseOos}{\ensuremath{p = 0.004}}
\providecommand{\rhcaQRidgeFullCloseOos}{}\renewcommand{\rhcaQRidgeFullCloseOos}{0.1269}
\providecommand{\rhcaMseRidgeFullCloseOos}{}\renewcommand{\rhcaMseRidgeFullCloseOos}{7.94}
\providecommand{\rhcaMaeRidgeFullCloseOos}{}\renewcommand{\rhcaMaeRidgeFullCloseOos}{5.39}
\providecommand{\rhcaRsqRidgeFullCloseOos}{}\renewcommand{\rhcaRsqRidgeFullCloseOos}{58.4}
\providecommand{\rhcaPctRidgeCloseOos}{}\renewcommand{\rhcaPctRidgeCloseOos}{11.6}  % QLIKE %, full vs HAR of the same class
\providecommand{\rhcaDmRidgeCloseOos}{}\renewcommand{\rhcaDmRidgeCloseOos}{2.44}
\providecommand{\rhcaPRidgeCloseOos}{}\renewcommand{\rhcaPRidgeCloseOos}{\ensuremath{p = 0.015}}
\providecommand{\rhcaMsePctRidgeCloseOos}{}\renewcommand{\rhcaMsePctRidgeCloseOos}{87.0}  % MSE %, full vs HAR
\providecommand{\rhcaMseDmRidgeCloseOos}{}\renewcommand{\rhcaMseDmRidgeCloseOos}{1.29}
\providecommand{\rhcaMsePRidgeCloseOos}{}\renewcommand{\rhcaMsePRidgeCloseOos}{\ensuremath{p = 0.20}}
\providecommand{\rhcaGapRidgeFullCloseOos}{}\renewcommand{\rhcaGapRidgeFullCloseOos}{23.8}  % QLIKE % vs the closing-interval HAR benchmark
\providecommand{\rhcaGapDmRidgeFullCloseOos}{}\renewcommand{\rhcaGapDmRidgeFullCloseOos}{3.72}
\providecommand{\rhcaGapPRidgeFullCloseOos}{}\renewcommand{\rhcaGapPRidgeFullCloseOos}{\ensuremath{p < 0.001}}
\providecommand{\rhcaQLassoHarCloseOos}{}\renewcommand{\rhcaQLassoHarCloseOos}{0.1137}
\providecommand{\rhcaMseLassoHarCloseOos}{}\renewcommand{\rhcaMseLassoHarCloseOos}{4.24}
\providecommand{\rhcaMaeLassoHarCloseOos}{}\renewcommand{\rhcaMaeLassoHarCloseOos}{4.55}
\providecommand{\rhcaRsqLassoHarCloseOos}{}\renewcommand{\rhcaRsqLassoHarCloseOos}{77.7}
\providecommand{\rhcaGapLassoHarCloseOos}{}\renewcommand{\rhcaGapLassoHarCloseOos}{10.9}  % QLIKE % vs the closing-interval HAR benchmark
\providecommand{\rhcaGapDmLassoHarCloseOos}{}\renewcommand{\rhcaGapDmLassoHarCloseOos}{2.84}
\providecommand{\rhcaGapPLassoHarCloseOos}{}\renewcommand{\rhcaGapPLassoHarCloseOos}{\ensuremath{p = 0.004}}
\providecommand{\rhcaQLassoFullCloseOos}{}\renewcommand{\rhcaQLassoFullCloseOos}{0.1309}
\providecommand{\rhcaMseLassoFullCloseOos}{}\renewcommand{\rhcaMseLassoFullCloseOos}{8.14}
\providecommand{\rhcaMaeLassoFullCloseOos}{}\renewcommand{\rhcaMaeLassoFullCloseOos}{5.45}
\providecommand{\rhcaRsqLassoFullCloseOos}{}\renewcommand{\rhcaRsqLassoFullCloseOos}{57.3}
\providecommand{\rhcaPctLassoCloseOos}{}\renewcommand{\rhcaPctLassoCloseOos}{15.2}  % QLIKE %, full vs HAR of the same class
\providecommand{\rhcaDmLassoCloseOos}{}\renewcommand{\rhcaDmLassoCloseOos}{2.68}
\providecommand{\rhcaPLassoCloseOos}{}\renewcommand{\rhcaPLassoCloseOos}{\ensuremath{p = 0.007}}
\providecommand{\rhcaMsePctLassoCloseOos}{}\renewcommand{\rhcaMsePctLassoCloseOos}{91.7}  % MSE %, full vs HAR
\providecommand{\rhcaMseDmLassoCloseOos}{}\renewcommand{\rhcaMseDmLassoCloseOos}{1.33}
\providecommand{\rhcaMsePLassoCloseOos}{}\renewcommand{\rhcaMsePLassoCloseOos}{\ensuremath{p = 0.18}}
\providecommand{\rhcaGapLassoFullCloseOos}{}\renewcommand{\rhcaGapLassoFullCloseOos}{27.7}  % QLIKE % vs the closing-interval HAR benchmark
\providecommand{\rhcaGapDmLassoFullCloseOos}{}\renewcommand{\rhcaGapDmLassoFullCloseOos}{3.80}
\providecommand{\rhcaGapPLassoFullCloseOos}{}\renewcommand{\rhcaGapPLassoFullCloseOos}{\ensuremath{p < 0.001}}
\providecommand{\rhcaQEnetHarCloseOos}{}\renewcommand{\rhcaQEnetHarCloseOos}{0.1137}
\providecommand{\rhcaMseEnetHarCloseOos}{}\renewcommand{\rhcaMseEnetHarCloseOos}{4.24}
\providecommand{\rhcaMaeEnetHarCloseOos}{}\renewcommand{\rhcaMaeEnetHarCloseOos}{4.55}
\providecommand{\rhcaRsqEnetHarCloseOos}{}\renewcommand{\rhcaRsqEnetHarCloseOos}{77.7}
\providecommand{\rhcaGapEnetHarCloseOos}{}\renewcommand{\rhcaGapEnetHarCloseOos}{10.9}  % QLIKE % vs the closing-interval HAR benchmark
\providecommand{\rhcaGapDmEnetHarCloseOos}{}\renewcommand{\rhcaGapDmEnetHarCloseOos}{2.84}
\providecommand{\rhcaGapPEnetHarCloseOos}{}\renewcommand{\rhcaGapPEnetHarCloseOos}{\ensuremath{p = 0.004}}
\providecommand{\rhcaQEnetFullCloseOos}{}\renewcommand{\rhcaQEnetFullCloseOos}{0.1245}
\providecommand{\rhcaMseEnetFullCloseOos}{}\renewcommand{\rhcaMseEnetFullCloseOos}{8.96}
\providecommand{\rhcaMaeEnetFullCloseOos}{}\renewcommand{\rhcaMaeEnetFullCloseOos}{5.54}
\providecommand{\rhcaRsqEnetFullCloseOos}{}\renewcommand{\rhcaRsqEnetFullCloseOos}{53.0}
\providecommand{\rhcaPctEnetCloseOos}{}\renewcommand{\rhcaPctEnetCloseOos}{9.5}  % QLIKE %, full vs HAR of the same class
\providecommand{\rhcaDmEnetCloseOos}{}\renewcommand{\rhcaDmEnetCloseOos}{2.09}
\providecommand{\rhcaPEnetCloseOos}{}\renewcommand{\rhcaPEnetCloseOos}{\ensuremath{p = 0.036}}
\providecommand{\rhcaMsePctEnetCloseOos}{}\renewcommand{\rhcaMsePctEnetCloseOos}{111.1}  % MSE %, full vs HAR
\providecommand{\rhcaMseDmEnetCloseOos}{}\renewcommand{\rhcaMseDmEnetCloseOos}{1.25}
\providecommand{\rhcaMsePEnetCloseOos}{}\renewcommand{\rhcaMsePEnetCloseOos}{\ensuremath{p = 0.21}}
\providecommand{\rhcaGapEnetFullCloseOos}{}\renewcommand{\rhcaGapEnetFullCloseOos}{21.4}  % QLIKE % vs the closing-interval HAR benchmark
\providecommand{\rhcaGapDmEnetFullCloseOos}{}\renewcommand{\rhcaGapDmEnetFullCloseOos}{3.48}
\providecommand{\rhcaGapPEnetFullCloseOos}{}\renewcommand{\rhcaGapPEnetFullCloseOos}{\ensuremath{p < 0.001}}
\providecommand{\rhcaQLgbmHarCloseOos}{}\renewcommand{\rhcaQLgbmHarCloseOos}{0.1101}
\providecommand{\rhcaMseLgbmHarCloseOos}{}\renewcommand{\rhcaMseLgbmHarCloseOos}{3.68}
\providecommand{\rhcaMaeLgbmHarCloseOos}{}\renewcommand{\rhcaMaeLgbmHarCloseOos}{4.63}
\providecommand{\rhcaRsqLgbmHarCloseOos}{}\renewcommand{\rhcaRsqLgbmHarCloseOos}{80.7}
\providecommand{\rhcaGapLgbmHarCloseOos}{}\renewcommand{\rhcaGapLgbmHarCloseOos}{7.4}  % QLIKE % vs the closing-interval HAR benchmark
\providecommand{\rhcaGapDmLgbmHarCloseOos}{}\renewcommand{\rhcaGapDmLgbmHarCloseOos}{2.60}
\providecommand{\rhcaGapPLgbmHarCloseOos}{}\renewcommand{\rhcaGapPLgbmHarCloseOos}{\ensuremath{p = 0.009}}
\providecommand{\rhcaQLgbmFullCloseOos}{}\renewcommand{\rhcaQLgbmFullCloseOos}{0.1086}
\providecommand{\rhcaMseLgbmFullCloseOos}{}\renewcommand{\rhcaMseLgbmFullCloseOos}{5.12}
\providecommand{\rhcaMaeLgbmFullCloseOos}{}\renewcommand{\rhcaMaeLgbmFullCloseOos}{4.80}
\providecommand{\rhcaRsqLgbmFullCloseOos}{}\renewcommand{\rhcaRsqLgbmFullCloseOos}{73.1}
\providecommand{\rhcaPctLgbmCloseOos}{}\renewcommand{\rhcaPctLgbmCloseOos}{\ensuremath{-}1.4}  % QLIKE %, full vs HAR of the same class
\providecommand{\rhcaDmLgbmCloseOos}{}\renewcommand{\rhcaDmLgbmCloseOos}{\ensuremath{-}0.49}
\providecommand{\rhcaPLgbmCloseOos}{}\renewcommand{\rhcaPLgbmCloseOos}{\ensuremath{p = 0.62}}
\providecommand{\rhcaMsePctLgbmCloseOos}{}\renewcommand{\rhcaMsePctLgbmCloseOos}{38.9}  % MSE %, full vs HAR
\providecommand{\rhcaMseDmLgbmCloseOos}{}\renewcommand{\rhcaMseDmLgbmCloseOos}{1.22}
\providecommand{\rhcaMsePLgbmCloseOos}{}\renewcommand{\rhcaMsePLgbmCloseOos}{\ensuremath{p = 0.22}}
\providecommand{\rhcaGapLgbmFullCloseOos}{}\renewcommand{\rhcaGapLgbmFullCloseOos}{6.0}  % QLIKE % vs the closing-interval HAR benchmark
\providecommand{\rhcaGapDmLgbmFullCloseOos}{}\renewcommand{\rhcaGapDmLgbmFullCloseOos}{1.41}
\providecommand{\rhcaGapPLgbmFullCloseOos}{}\renewcommand{\rhcaGapPLgbmFullCloseOos}{\ensuremath{p = 0.16}}
\providecommand{\rhcaQXgbHarCloseOos}{}\renewcommand{\rhcaQXgbHarCloseOos}{0.1102}
\providecommand{\rhcaMseXgbHarCloseOos}{}\renewcommand{\rhcaMseXgbHarCloseOos}{3.62}
\providecommand{\rhcaMaeXgbHarCloseOos}{}\renewcommand{\rhcaMaeXgbHarCloseOos}{4.51}
\providecommand{\rhcaRsqXgbHarCloseOos}{}\renewcommand{\rhcaRsqXgbHarCloseOos}{81.0}
\providecommand{\rhcaGapXgbHarCloseOos}{}\renewcommand{\rhcaGapXgbHarCloseOos}{7.5}  % QLIKE % vs the closing-interval HAR benchmark
\providecommand{\rhcaGapDmXgbHarCloseOos}{}\renewcommand{\rhcaGapDmXgbHarCloseOos}{2.70}
\providecommand{\rhcaGapPXgbHarCloseOos}{}\renewcommand{\rhcaGapPXgbHarCloseOos}{\ensuremath{p = 0.007}}
\providecommand{\rhcaQXgbFullCloseOos}{}\renewcommand{\rhcaQXgbFullCloseOos}{0.1083}
\providecommand{\rhcaMseXgbFullCloseOos}{}\renewcommand{\rhcaMseXgbFullCloseOos}{4.76}
\providecommand{\rhcaMaeXgbFullCloseOos}{}\renewcommand{\rhcaMaeXgbFullCloseOos}{4.76}
\providecommand{\rhcaRsqXgbFullCloseOos}{}\renewcommand{\rhcaRsqXgbFullCloseOos}{75.0}
\providecommand{\rhcaPctXgbCloseOos}{}\renewcommand{\rhcaPctXgbCloseOos}{\ensuremath{-}1.7}  % QLIKE %, full vs HAR of the same class
\providecommand{\rhcaDmXgbCloseOos}{}\renewcommand{\rhcaDmXgbCloseOos}{\ensuremath{-}0.63}
\providecommand{\rhcaPXgbCloseOos}{}\renewcommand{\rhcaPXgbCloseOos}{\ensuremath{p = 0.53}}
\providecommand{\rhcaMsePctXgbCloseOos}{}\renewcommand{\rhcaMsePctXgbCloseOos}{31.4}  % MSE %, full vs HAR
\providecommand{\rhcaMseDmXgbCloseOos}{}\renewcommand{\rhcaMseDmXgbCloseOos}{1.23}
\providecommand{\rhcaMsePXgbCloseOos}{}\renewcommand{\rhcaMsePXgbCloseOos}{\ensuremath{p = 0.22}}
\providecommand{\rhcaGapXgbFullCloseOos}{}\renewcommand{\rhcaGapXgbFullCloseOos}{5.7}  % QLIKE % vs the closing-interval HAR benchmark
\providecommand{\rhcaGapDmXgbFullCloseOos}{}\renewcommand{\rhcaGapDmXgbFullCloseOos}{1.40}
\providecommand{\rhcaGapPXgbFullCloseOos}{}\renewcommand{\rhcaGapPXgbFullCloseOos}{\ensuremath{p = 0.16}}
\providecommand{\rhcaQRfHarCloseOos}{}\renewcommand{\rhcaQRfHarCloseOos}{0.1131}
\providecommand{\rhcaMseRfHarCloseOos}{}\renewcommand{\rhcaMseRfHarCloseOos}{3.39}
\providecommand{\rhcaMaeRfHarCloseOos}{}\renewcommand{\rhcaMaeRfHarCloseOos}{4.42}
\providecommand{\rhcaRsqRfHarCloseOos}{}\renewcommand{\rhcaRsqRfHarCloseOos}{82.2}
\providecommand{\rhcaGapRfHarCloseOos}{}\renewcommand{\rhcaGapRfHarCloseOos}{10.4}  % QLIKE % vs the closing-interval HAR benchmark
\providecommand{\rhcaGapDmRfHarCloseOos}{}\renewcommand{\rhcaGapDmRfHarCloseOos}{3.12}
\providecommand{\rhcaGapPRfHarCloseOos}{}\renewcommand{\rhcaGapPRfHarCloseOos}{\ensuremath{p = 0.002}}
\providecommand{\rhcaQRfFullCloseOos}{}\renewcommand{\rhcaQRfFullCloseOos}{0.1123}
\providecommand{\rhcaMseRfFullCloseOos}{}\renewcommand{\rhcaMseRfFullCloseOos}{4.18}
\providecommand{\rhcaMaeRfFullCloseOos}{}\renewcommand{\rhcaMaeRfFullCloseOos}{4.63}
\providecommand{\rhcaRsqRfFullCloseOos}{}\renewcommand{\rhcaRsqRfFullCloseOos}{78.0}
\providecommand{\rhcaPctRfCloseOos}{}\renewcommand{\rhcaPctRfCloseOos}{\ensuremath{-}0.7}  % QLIKE %, full vs HAR of the same class
\providecommand{\rhcaDmRfCloseOos}{}\renewcommand{\rhcaDmRfCloseOos}{\ensuremath{-}0.27}
\providecommand{\rhcaPRfCloseOos}{}\renewcommand{\rhcaPRfCloseOos}{\ensuremath{p = 0.79}}
\providecommand{\rhcaMsePctRfCloseOos}{}\renewcommand{\rhcaMsePctRfCloseOos}{23.3}  % MSE %, full vs HAR
\providecommand{\rhcaMseDmRfCloseOos}{}\renewcommand{\rhcaMseDmRfCloseOos}{1.25}
\providecommand{\rhcaMsePRfCloseOos}{}\renewcommand{\rhcaMsePRfCloseOos}{\ensuremath{p = 0.21}}
\providecommand{\rhcaGapRfFullCloseOos}{}\renewcommand{\rhcaGapRfFullCloseOos}{9.6}  % QLIKE % vs the closing-interval HAR benchmark
\providecommand{\rhcaGapDmRfFullCloseOos}{}\renewcommand{\rhcaGapDmRfFullCloseOos}{2.13}
\providecommand{\rhcaGapPRfFullCloseOos}{}\renewcommand{\rhcaGapPRfFullCloseOos}{\ensuremath{p = 0.033}}
\providecommand{\rhcaQBmCloseOos}{}\renewcommand{\rhcaQBmCloseOos}{0.1025}
\providecommand{\rhcaMseBmCloseOos}{}\renewcommand{\rhcaMseBmCloseOos}{4.31}
\providecommand{\rhcaMaeBmCloseOos}{}\renewcommand{\rhcaMaeBmCloseOos}{4.78}
\providecommand{\rhcaRsqBmCloseOos}{}\renewcommand{\rhcaRsqBmCloseOos}{77.4}
\providecommand{\rhcaQCiRidgeHarCloseOos}{}\renewcommand{\rhcaQCiRidgeHarCloseOos}{0.1033}
\providecommand{\rhcaMseCiRidgeHarCloseOos}{}\renewcommand{\rhcaMseCiRidgeHarCloseOos}{4.15}
\providecommand{\rhcaMaeCiRidgeHarCloseOos}{}\renewcommand{\rhcaMaeCiRidgeHarCloseOos}{4.67}
\providecommand{\rhcaRsqCiRidgeHarCloseOos}{}\renewcommand{\rhcaRsqCiRidgeHarCloseOos}{78.2}
\providecommand{\rhcaGapCiRidgeHarCloseOos}{}\renewcommand{\rhcaGapCiRidgeHarCloseOos}{0.7}  % QLIKE % vs the closing-interval HAR benchmark
\providecommand{\rhcaGapDmCiRidgeHarCloseOos}{}\renewcommand{\rhcaGapDmCiRidgeHarCloseOos}{0.86}
\providecommand{\rhcaGapPCiRidgeHarCloseOos}{}\renewcommand{\rhcaGapPCiRidgeHarCloseOos}{\ensuremath{p = 0.39}}
\providecommand{\rhcaQCiRidgeFullCloseOos}{}\renewcommand{\rhcaQCiRidgeFullCloseOos}{0.1031}
\providecommand{\rhcaMseCiRidgeFullCloseOos}{}\renewcommand{\rhcaMseCiRidgeFullCloseOos}{7.30}
\providecommand{\rhcaMaeCiRidgeFullCloseOos}{}\renewcommand{\rhcaMaeCiRidgeFullCloseOos}{5.27}
\providecommand{\rhcaRsqCiRidgeFullCloseOos}{}\renewcommand{\rhcaRsqCiRidgeFullCloseOos}{61.7}
\providecommand{\rhcaPctCiRidgeCloseOos}{}\renewcommand{\rhcaPctCiRidgeCloseOos}{\ensuremath{-}0.1}  % QLIKE %, full vs HAR of the same class
\providecommand{\rhcaDmCiRidgeCloseOos}{}\renewcommand{\rhcaDmCiRidgeCloseOos}{\ensuremath{-}0.04}
\providecommand{\rhcaPCiRidgeCloseOos}{}\renewcommand{\rhcaPCiRidgeCloseOos}{\ensuremath{p = 0.97}}
\providecommand{\rhcaMsePctCiRidgeCloseOos}{}\renewcommand{\rhcaMsePctCiRidgeCloseOos}{75.8}  % MSE %, full vs HAR
\providecommand{\rhcaMseDmCiRidgeCloseOos}{}\renewcommand{\rhcaMseDmCiRidgeCloseOos}{1.45}
\providecommand{\rhcaMsePCiRidgeCloseOos}{}\renewcommand{\rhcaMsePCiRidgeCloseOos}{\ensuremath{p = 0.15}}
\providecommand{\rhcaGapCiRidgeFullCloseOos}{}\renewcommand{\rhcaGapCiRidgeFullCloseOos}{0.6}  % QLIKE % vs the closing-interval HAR benchmark
\providecommand{\rhcaGapDmCiRidgeFullCloseOos}{}\renewcommand{\rhcaGapDmCiRidgeFullCloseOos}{0.15}
\providecommand{\rhcaGapPCiRidgeFullCloseOos}{}\renewcommand{\rhcaGapPCiRidgeFullCloseOos}{\ensuremath{p = 0.88}}
\providecommand{\rhcaQCiLassoHarCloseOos}{}\renewcommand{\rhcaQCiLassoHarCloseOos}{0.1033}
\providecommand{\rhcaMseCiLassoHarCloseOos}{}\renewcommand{\rhcaMseCiLassoHarCloseOos}{4.15}
\providecommand{\rhcaMaeCiLassoHarCloseOos}{}\renewcommand{\rhcaMaeCiLassoHarCloseOos}{4.67}
\providecommand{\rhcaRsqCiLassoHarCloseOos}{}\renewcommand{\rhcaRsqCiLassoHarCloseOos}{78.2}
\providecommand{\rhcaGapCiLassoHarCloseOos}{}\renewcommand{\rhcaGapCiLassoHarCloseOos}{0.7}  % QLIKE % vs the closing-interval HAR benchmark
\providecommand{\rhcaGapDmCiLassoHarCloseOos}{}\renewcommand{\rhcaGapDmCiLassoHarCloseOos}{0.86}
\providecommand{\rhcaGapPCiLassoHarCloseOos}{}\renewcommand{\rhcaGapPCiLassoHarCloseOos}{\ensuremath{p = 0.39}}
\providecommand{\rhcaQCiLassoFullCloseOos}{}\renewcommand{\rhcaQCiLassoFullCloseOos}{0.0992}
\providecommand{\rhcaMseCiLassoFullCloseOos}{}\renewcommand{\rhcaMseCiLassoFullCloseOos}{5.94}
\providecommand{\rhcaMaeCiLassoFullCloseOos}{}\renewcommand{\rhcaMaeCiLassoFullCloseOos}{5.07}
\providecommand{\rhcaRsqCiLassoFullCloseOos}{}\renewcommand{\rhcaRsqCiLassoFullCloseOos}{68.8}
\providecommand{\rhcaPctCiLassoCloseOos}{}\renewcommand{\rhcaPctCiLassoCloseOos}{\ensuremath{-}4.0}  % QLIKE %, full vs HAR of the same class
\providecommand{\rhcaDmCiLassoCloseOos}{}\renewcommand{\rhcaDmCiLassoCloseOos}{\ensuremath{-}1.91}
\providecommand{\rhcaPCiLassoCloseOos}{}\renewcommand{\rhcaPCiLassoCloseOos}{\ensuremath{p = 0.056}}
\providecommand{\rhcaMsePctCiLassoCloseOos}{}\renewcommand{\rhcaMsePctCiLassoCloseOos}{43.1}  % MSE %, full vs HAR
\providecommand{\rhcaMseDmCiLassoCloseOos}{}\renewcommand{\rhcaMseDmCiLassoCloseOos}{1.34}
\providecommand{\rhcaMsePCiLassoCloseOos}{}\renewcommand{\rhcaMsePCiLassoCloseOos}{\ensuremath{p = 0.18}}
\providecommand{\rhcaGapCiLassoFullCloseOos}{}\renewcommand{\rhcaGapCiLassoFullCloseOos}{\ensuremath{-}3.3}  % QLIKE % vs the closing-interval HAR benchmark
\providecommand{\rhcaGapDmCiLassoFullCloseOos}{}\renewcommand{\rhcaGapDmCiLassoFullCloseOos}{\ensuremath{-}1.44}
\providecommand{\rhcaGapPCiLassoFullCloseOos}{}\renewcommand{\rhcaGapPCiLassoFullCloseOos}{\ensuremath{p = 0.15}}
\providecommand{\rhcaQCiEnetHarCloseOos}{}\renewcommand{\rhcaQCiEnetHarCloseOos}{0.1033}
\providecommand{\rhcaMseCiEnetHarCloseOos}{}\renewcommand{\rhcaMseCiEnetHarCloseOos}{4.15}
\providecommand{\rhcaMaeCiEnetHarCloseOos}{}\renewcommand{\rhcaMaeCiEnetHarCloseOos}{4.67}
\providecommand{\rhcaRsqCiEnetHarCloseOos}{}\renewcommand{\rhcaRsqCiEnetHarCloseOos}{78.2}
\providecommand{\rhcaGapCiEnetHarCloseOos}{}\renewcommand{\rhcaGapCiEnetHarCloseOos}{0.7}  % QLIKE % vs the closing-interval HAR benchmark
\providecommand{\rhcaGapDmCiEnetHarCloseOos}{}\renewcommand{\rhcaGapDmCiEnetHarCloseOos}{0.86}
\providecommand{\rhcaGapPCiEnetHarCloseOos}{}\renewcommand{\rhcaGapPCiEnetHarCloseOos}{\ensuremath{p = 0.39}}
\providecommand{\rhcaQCiEnetFullCloseOos}{}\renewcommand{\rhcaQCiEnetFullCloseOos}{0.1014}
\providecommand{\rhcaMseCiEnetFullCloseOos}{}\renewcommand{\rhcaMseCiEnetFullCloseOos}{6.89}
\providecommand{\rhcaMaeCiEnetFullCloseOos}{}\renewcommand{\rhcaMaeCiEnetFullCloseOos}{5.21}
\providecommand{\rhcaRsqCiEnetFullCloseOos}{}\renewcommand{\rhcaRsqCiEnetFullCloseOos}{63.9}
\providecommand{\rhcaPctCiEnetCloseOos}{}\renewcommand{\rhcaPctCiEnetCloseOos}{\ensuremath{-}1.8}  % QLIKE %, full vs HAR of the same class
\providecommand{\rhcaDmCiEnetCloseOos}{}\renewcommand{\rhcaDmCiEnetCloseOos}{\ensuremath{-}0.55}
\providecommand{\rhcaPCiEnetCloseOos}{}\renewcommand{\rhcaPCiEnetCloseOos}{\ensuremath{p = 0.58}}
\providecommand{\rhcaMsePctCiEnetCloseOos}{}\renewcommand{\rhcaMsePctCiEnetCloseOos}{65.9}  % MSE %, full vs HAR
\providecommand{\rhcaMseDmCiEnetCloseOos}{}\renewcommand{\rhcaMseDmCiEnetCloseOos}{1.45}
\providecommand{\rhcaMsePCiEnetCloseOos}{}\renewcommand{\rhcaMsePCiEnetCloseOos}{\ensuremath{p = 0.15}}
\providecommand{\rhcaGapCiEnetFullCloseOos}{}\renewcommand{\rhcaGapCiEnetFullCloseOos}{\ensuremath{-}1.1}  % QLIKE % vs the closing-interval HAR benchmark
\providecommand{\rhcaGapDmCiEnetFullCloseOos}{}\renewcommand{\rhcaGapDmCiEnetFullCloseOos}{\ensuremath{-}0.32}
\providecommand{\rhcaGapPCiEnetFullCloseOos}{}\renewcommand{\rhcaGapPCiEnetFullCloseOos}{\ensuremath{p = 0.75}}
\providecommand{\rhcaQCiLgbmHarCloseOos}{}\renewcommand{\rhcaQCiLgbmHarCloseOos}{0.1094}
\providecommand{\rhcaMseCiLgbmHarCloseOos}{}\renewcommand{\rhcaMseCiLgbmHarCloseOos}{5.75}
\providecommand{\rhcaMaeCiLgbmHarCloseOos}{}\renewcommand{\rhcaMaeCiLgbmHarCloseOos}{5.02}
\providecommand{\rhcaRsqCiLgbmHarCloseOos}{}\renewcommand{\rhcaRsqCiLgbmHarCloseOos}{69.8}
\providecommand{\rhcaGapCiLgbmHarCloseOos}{}\renewcommand{\rhcaGapCiLgbmHarCloseOos}{6.7}  % QLIKE % vs the closing-interval HAR benchmark
\providecommand{\rhcaGapDmCiLgbmHarCloseOos}{}\renewcommand{\rhcaGapDmCiLgbmHarCloseOos}{3.03}
\providecommand{\rhcaGapPCiLgbmHarCloseOos}{}\renewcommand{\rhcaGapPCiLgbmHarCloseOos}{\ensuremath{p = 0.002}}
\providecommand{\rhcaQCiLgbmFullCloseOos}{}\renewcommand{\rhcaQCiLgbmFullCloseOos}{0.1049}
\providecommand{\rhcaMseCiLgbmFullCloseOos}{}\renewcommand{\rhcaMseCiLgbmFullCloseOos}{17.19}
\providecommand{\rhcaMaeCiLgbmFullCloseOos}{}\renewcommand{\rhcaMaeCiLgbmFullCloseOos}{6.84}
\providecommand{\rhcaRsqCiLgbmFullCloseOos}{}\renewcommand{\rhcaRsqCiLgbmFullCloseOos}{9.8}
\providecommand{\rhcaPctCiLgbmCloseOos}{}\renewcommand{\rhcaPctCiLgbmCloseOos}{\ensuremath{-}4.1}  % QLIKE %, full vs HAR of the same class
\providecommand{\rhcaDmCiLgbmCloseOos}{}\renewcommand{\rhcaDmCiLgbmCloseOos}{\ensuremath{-}1.48}
\providecommand{\rhcaPCiLgbmCloseOos}{}\renewcommand{\rhcaPCiLgbmCloseOos}{\ensuremath{p = 0.14}}
\providecommand{\rhcaMsePctCiLgbmCloseOos}{}\renewcommand{\rhcaMsePctCiLgbmCloseOos}{199.1}  % MSE %, full vs HAR
\providecommand{\rhcaMseDmCiLgbmCloseOos}{}\renewcommand{\rhcaMseDmCiLgbmCloseOos}{1.21}
\providecommand{\rhcaMsePCiLgbmCloseOos}{}\renewcommand{\rhcaMsePCiLgbmCloseOos}{\ensuremath{p = 0.23}}
\providecommand{\rhcaGapCiLgbmFullCloseOos}{}\renewcommand{\rhcaGapCiLgbmFullCloseOos}{2.3}  % QLIKE % vs the closing-interval HAR benchmark
\providecommand{\rhcaGapDmCiLgbmFullCloseOos}{}\renewcommand{\rhcaGapDmCiLgbmFullCloseOos}{0.72}
\providecommand{\rhcaGapPCiLgbmFullCloseOos}{}\renewcommand{\rhcaGapPCiLgbmFullCloseOos}{\ensuremath{p = 0.47}}
\providecommand{\rhcaQCiXgbHarCloseOos}{}\renewcommand{\rhcaQCiXgbHarCloseOos}{0.1076}
\providecommand{\rhcaMseCiXgbHarCloseOos}{}\renewcommand{\rhcaMseCiXgbHarCloseOos}{5.27}
\providecommand{\rhcaMaeCiXgbHarCloseOos}{}\renewcommand{\rhcaMaeCiXgbHarCloseOos}{4.94}
\providecommand{\rhcaRsqCiXgbHarCloseOos}{}\renewcommand{\rhcaRsqCiXgbHarCloseOos}{72.4}
\providecommand{\rhcaGapCiXgbHarCloseOos}{}\renewcommand{\rhcaGapCiXgbHarCloseOos}{5.0}  % QLIKE % vs the closing-interval HAR benchmark
\providecommand{\rhcaGapDmCiXgbHarCloseOos}{}\renewcommand{\rhcaGapDmCiXgbHarCloseOos}{2.45}
\providecommand{\rhcaGapPCiXgbHarCloseOos}{}\renewcommand{\rhcaGapPCiXgbHarCloseOos}{\ensuremath{p = 0.014}}
\providecommand{\rhcaQCiXgbFullCloseOos}{}\renewcommand{\rhcaQCiXgbFullCloseOos}{0.1031}
\providecommand{\rhcaMseCiXgbFullCloseOos}{}\renewcommand{\rhcaMseCiXgbFullCloseOos}{20.94}
\providecommand{\rhcaMaeCiXgbFullCloseOos}{}\renewcommand{\rhcaMaeCiXgbFullCloseOos}{6.98}
\providecommand{\rhcaRsqCiXgbFullCloseOos}{}\renewcommand{\rhcaRsqCiXgbFullCloseOos}{\ensuremath{-}9.9}
\providecommand{\rhcaPctCiXgbCloseOos}{}\renewcommand{\rhcaPctCiXgbCloseOos}{\ensuremath{-}4.1}  % QLIKE %, full vs HAR of the same class
\providecommand{\rhcaDmCiXgbCloseOos}{}\renewcommand{\rhcaDmCiXgbCloseOos}{\ensuremath{-}1.63}
\providecommand{\rhcaPCiXgbCloseOos}{}\renewcommand{\rhcaPCiXgbCloseOos}{\ensuremath{p = 0.1039}}
\providecommand{\rhcaMsePctCiXgbCloseOos}{}\renewcommand{\rhcaMsePctCiXgbCloseOos}{297.6}  % MSE %, full vs HAR
\providecommand{\rhcaMseDmCiXgbCloseOos}{}\renewcommand{\rhcaMseDmCiXgbCloseOos}{1.15}
\providecommand{\rhcaMsePCiXgbCloseOos}{}\renewcommand{\rhcaMsePCiXgbCloseOos}{\ensuremath{p = 0.25}}
\providecommand{\rhcaGapCiXgbFullCloseOos}{}\renewcommand{\rhcaGapCiXgbFullCloseOos}{0.6}  % QLIKE % vs the closing-interval HAR benchmark
\providecommand{\rhcaGapDmCiXgbFullCloseOos}{}\renewcommand{\rhcaGapDmCiXgbFullCloseOos}{0.21}
\providecommand{\rhcaGapPCiXgbFullCloseOos}{}\renewcommand{\rhcaGapPCiXgbFullCloseOos}{\ensuremath{p = 0.83}}
\providecommand{\rhcaQCiRfHarCloseOos}{}\renewcommand{\rhcaQCiRfHarCloseOos}{0.1068}
\providecommand{\rhcaMseCiRfHarCloseOos}{}\renewcommand{\rhcaMseCiRfHarCloseOos}{4.62}
\providecommand{\rhcaMaeCiRfHarCloseOos}{}\renewcommand{\rhcaMaeCiRfHarCloseOos}{4.76}
\providecommand{\rhcaRsqCiRfHarCloseOos}{}\renewcommand{\rhcaRsqCiRfHarCloseOos}{75.8}
\providecommand{\rhcaGapCiRfHarCloseOos}{}\renewcommand{\rhcaGapCiRfHarCloseOos}{4.2}  % QLIKE % vs the closing-interval HAR benchmark
\providecommand{\rhcaGapDmCiRfHarCloseOos}{}\renewcommand{\rhcaGapDmCiRfHarCloseOos}{2.28}
\providecommand{\rhcaGapPCiRfHarCloseOos}{}\renewcommand{\rhcaGapPCiRfHarCloseOos}{\ensuremath{p = 0.022}}
\providecommand{\rhcaQCiRfFullCloseOos}{}\renewcommand{\rhcaQCiRfFullCloseOos}{0.1060}
\providecommand{\rhcaMseCiRfFullCloseOos}{}\renewcommand{\rhcaMseCiRfFullCloseOos}{14.15}
\providecommand{\rhcaMaeCiRfFullCloseOos}{}\renewcommand{\rhcaMaeCiRfFullCloseOos}{6.74}
\providecommand{\rhcaRsqCiRfFullCloseOos}{}\renewcommand{\rhcaRsqCiRfFullCloseOos}{25.8}
\providecommand{\rhcaPctCiRfCloseOos}{}\renewcommand{\rhcaPctCiRfCloseOos}{\ensuremath{-}0.8}  % QLIKE %, full vs HAR of the same class
\providecommand{\rhcaDmCiRfCloseOos}{}\renewcommand{\rhcaDmCiRfCloseOos}{\ensuremath{-}0.28}
\providecommand{\rhcaPCiRfCloseOos}{}\renewcommand{\rhcaPCiRfCloseOos}{\ensuremath{p = 0.78}}
\providecommand{\rhcaMsePctCiRfCloseOos}{}\renewcommand{\rhcaMsePctCiRfCloseOos}{206.6}  % MSE %, full vs HAR
\providecommand{\rhcaMseDmCiRfCloseOos}{}\renewcommand{\rhcaMseDmCiRfCloseOos}{1.22}
\providecommand{\rhcaMsePCiRfCloseOos}{}\renewcommand{\rhcaMsePCiRfCloseOos}{\ensuremath{p = 0.22}}
\providecommand{\rhcaGapCiRfFullCloseOos}{}\renewcommand{\rhcaGapCiRfFullCloseOos}{3.4}  % QLIKE % vs the closing-interval HAR benchmark
\providecommand{\rhcaGapDmCiRfFullCloseOos}{}\renewcommand{\rhcaGapDmCiRfFullCloseOos}{1.11}
\providecommand{\rhcaGapPCiRfFullCloseOos}{}\renewcommand{\rhcaGapPCiRfFullCloseOos}{\ensuremath{p = 0.27}}
% --- every bar: pooled forecasts, Span / Pre / Oos (Table 1 and its subperiods)
\providecommand{\rhcaQOlsHarBarSpan}{}\renewcommand{\rhcaQOlsHarBarSpan}{0.1893}
\providecommand{\rhcaMseOlsHarBarSpan}{}\renewcommand{\rhcaMseOlsHarBarSpan}{1.11}
\providecommand{\rhcaMaeOlsHarBarSpan}{}\renewcommand{\rhcaMaeOlsHarBarSpan}{2.25}
\providecommand{\rhcaRsqOlsHarBarSpan}{}\renewcommand{\rhcaRsqOlsHarBarSpan}{65.4}
\providecommand{\rhcaQRidgeHarBarSpan}{}\renewcommand{\rhcaQRidgeHarBarSpan}{0.1893}
\providecommand{\rhcaMseRidgeHarBarSpan}{}\renewcommand{\rhcaMseRidgeHarBarSpan}{1.11}
\providecommand{\rhcaMaeRidgeHarBarSpan}{}\renewcommand{\rhcaMaeRidgeHarBarSpan}{2.25}
\providecommand{\rhcaRsqRidgeHarBarSpan}{}\renewcommand{\rhcaRsqRidgeHarBarSpan}{65.4}
\providecommand{\rhcaQRidgeFullBarSpan}{}\renewcommand{\rhcaQRidgeFullBarSpan}{0.1778}
\providecommand{\rhcaMseRidgeFullBarSpan}{}\renewcommand{\rhcaMseRidgeFullBarSpan}{2.23}
\providecommand{\rhcaMaeRidgeFullBarSpan}{}\renewcommand{\rhcaMaeRidgeFullBarSpan}{2.38}
\providecommand{\rhcaRsqRidgeFullBarSpan}{}\renewcommand{\rhcaRsqRidgeFullBarSpan}{30.7}
\providecommand{\rhcaPctRidgeBarSpan}{}\renewcommand{\rhcaPctRidgeBarSpan}{\ensuremath{-}6.0}
\providecommand{\rhcaDmRidgeBarSpan}{}\renewcommand{\rhcaDmRidgeBarSpan}{\ensuremath{-}6.28}
\providecommand{\rhcaPRidgeBarSpan}{}\renewcommand{\rhcaPRidgeBarSpan}{\ensuremath{p < 0.001}}
\providecommand{\rhcaMsePctRidgeBarSpan}{}\renewcommand{\rhcaMsePctRidgeBarSpan}{100.6}
\providecommand{\rhcaMseDmRidgeBarSpan}{}\renewcommand{\rhcaMseDmRidgeBarSpan}{1.12}
\providecommand{\rhcaMsePRidgeBarSpan}{}\renewcommand{\rhcaMsePRidgeBarSpan}{\ensuremath{p = 0.26}}
\providecommand{\rhcaQLassoHarBarSpan}{}\renewcommand{\rhcaQLassoHarBarSpan}{0.1893}
\providecommand{\rhcaMseLassoHarBarSpan}{}\renewcommand{\rhcaMseLassoHarBarSpan}{1.11}
\providecommand{\rhcaMaeLassoHarBarSpan}{}\renewcommand{\rhcaMaeLassoHarBarSpan}{2.25}
\providecommand{\rhcaRsqLassoHarBarSpan}{}\renewcommand{\rhcaRsqLassoHarBarSpan}{65.4}
\providecommand{\rhcaQLassoFullBarSpan}{}\renewcommand{\rhcaQLassoFullBarSpan}{0.1723}
\providecommand{\rhcaMseLassoFullBarSpan}{}\renewcommand{\rhcaMseLassoFullBarSpan}{2.39}
\providecommand{\rhcaMaeLassoFullBarSpan}{}\renewcommand{\rhcaMaeLassoFullBarSpan}{2.40}
\providecommand{\rhcaRsqLassoFullBarSpan}{}\renewcommand{\rhcaRsqLassoFullBarSpan}{25.8}
\providecommand{\rhcaPctLassoBarSpan}{}\renewcommand{\rhcaPctLassoBarSpan}{\ensuremath{-}9.0}
\providecommand{\rhcaDmLassoBarSpan}{}\renewcommand{\rhcaDmLassoBarSpan}{\ensuremath{-}6.34}
\providecommand{\rhcaPLassoBarSpan}{}\renewcommand{\rhcaPLassoBarSpan}{\ensuremath{p < 0.001}}
\providecommand{\rhcaMsePctLassoBarSpan}{}\renewcommand{\rhcaMsePctLassoBarSpan}{114.7}
\providecommand{\rhcaMseDmLassoBarSpan}{}\renewcommand{\rhcaMseDmLassoBarSpan}{1.16}
\providecommand{\rhcaMsePLassoBarSpan}{}\renewcommand{\rhcaMsePLassoBarSpan}{\ensuremath{p = 0.25}}
\providecommand{\rhcaQEnetHarBarSpan}{}\renewcommand{\rhcaQEnetHarBarSpan}{0.1893}
\providecommand{\rhcaMseEnetHarBarSpan}{}\renewcommand{\rhcaMseEnetHarBarSpan}{1.11}
\providecommand{\rhcaMaeEnetHarBarSpan}{}\renewcommand{\rhcaMaeEnetHarBarSpan}{2.25}
\providecommand{\rhcaRsqEnetHarBarSpan}{}\renewcommand{\rhcaRsqEnetHarBarSpan}{65.4}
\providecommand{\rhcaQEnetFullBarSpan}{}\renewcommand{\rhcaQEnetFullBarSpan}{0.1766}
\providecommand{\rhcaMseEnetFullBarSpan}{}\renewcommand{\rhcaMseEnetFullBarSpan}{2.42}
\providecommand{\rhcaMaeEnetFullBarSpan}{}\renewcommand{\rhcaMaeEnetFullBarSpan}{2.40}
\providecommand{\rhcaRsqEnetFullBarSpan}{}\renewcommand{\rhcaRsqEnetFullBarSpan}{24.8}
\providecommand{\rhcaPctEnetBarSpan}{}\renewcommand{\rhcaPctEnetBarSpan}{\ensuremath{-}6.7}
\providecommand{\rhcaDmEnetBarSpan}{}\renewcommand{\rhcaDmEnetBarSpan}{\ensuremath{-}7.36}
\providecommand{\rhcaPEnetBarSpan}{}\renewcommand{\rhcaPEnetBarSpan}{\ensuremath{p < 0.001}}
\providecommand{\rhcaMsePctEnetBarSpan}{}\renewcommand{\rhcaMsePctEnetBarSpan}{117.7}
\providecommand{\rhcaMseDmEnetBarSpan}{}\renewcommand{\rhcaMseDmEnetBarSpan}{1.12}
\providecommand{\rhcaMsePEnetBarSpan}{}\renewcommand{\rhcaMsePEnetBarSpan}{\ensuremath{p = 0.26}}
\providecommand{\rhcaQLgbmHarBarSpan}{}\renewcommand{\rhcaQLgbmHarBarSpan}{0.1875}
\providecommand{\rhcaMseLgbmHarBarSpan}{}\renewcommand{\rhcaMseLgbmHarBarSpan}{1.18}
\providecommand{\rhcaMaeLgbmHarBarSpan}{}\renewcommand{\rhcaMaeLgbmHarBarSpan}{2.27}
\providecommand{\rhcaRsqLgbmHarBarSpan}{}\renewcommand{\rhcaRsqLgbmHarBarSpan}{63.3}
\providecommand{\rhcaQLgbmFullBarSpan}{}\renewcommand{\rhcaQLgbmFullBarSpan}{0.1611}
\providecommand{\rhcaMseLgbmFullBarSpan}{}\renewcommand{\rhcaMseLgbmFullBarSpan}{1.42}
\providecommand{\rhcaMaeLgbmFullBarSpan}{}\renewcommand{\rhcaMaeLgbmFullBarSpan}{2.28}
\providecommand{\rhcaRsqLgbmFullBarSpan}{}\renewcommand{\rhcaRsqLgbmFullBarSpan}{55.9}
\providecommand{\rhcaPctLgbmBarSpan}{}\renewcommand{\rhcaPctLgbmBarSpan}{\ensuremath{-}14.1}
\providecommand{\rhcaDmLgbmBarSpan}{}\renewcommand{\rhcaDmLgbmBarSpan}{\ensuremath{-}10.26}
\providecommand{\rhcaPLgbmBarSpan}{}\renewcommand{\rhcaPLgbmBarSpan}{\ensuremath{p < 0.001}}
\providecommand{\rhcaMsePctLgbmBarSpan}{}\renewcommand{\rhcaMsePctLgbmBarSpan}{20.0}
\providecommand{\rhcaMseDmLgbmBarSpan}{}\renewcommand{\rhcaMseDmLgbmBarSpan}{1.03}
\providecommand{\rhcaMsePLgbmBarSpan}{}\renewcommand{\rhcaMsePLgbmBarSpan}{\ensuremath{p = 0.30}}
\providecommand{\rhcaQXgbHarBarSpan}{}\renewcommand{\rhcaQXgbHarBarSpan}{0.1888}
\providecommand{\rhcaMseXgbHarBarSpan}{}\renewcommand{\rhcaMseXgbHarBarSpan}{1.16}
\providecommand{\rhcaMaeXgbHarBarSpan}{}\renewcommand{\rhcaMaeXgbHarBarSpan}{2.26}
\providecommand{\rhcaRsqXgbHarBarSpan}{}\renewcommand{\rhcaRsqXgbHarBarSpan}{63.9}
\providecommand{\rhcaQXgbFullBarSpan}{}\renewcommand{\rhcaQXgbFullBarSpan}{0.1612}
\providecommand{\rhcaMseXgbFullBarSpan}{}\renewcommand{\rhcaMseXgbFullBarSpan}{1.35}
\providecommand{\rhcaMaeXgbFullBarSpan}{}\renewcommand{\rhcaMaeXgbFullBarSpan}{2.27}
\providecommand{\rhcaRsqXgbFullBarSpan}{}\renewcommand{\rhcaRsqXgbFullBarSpan}{58.0}
\providecommand{\rhcaPctXgbBarSpan}{}\renewcommand{\rhcaPctXgbBarSpan}{\ensuremath{-}14.6}
\providecommand{\rhcaDmXgbBarSpan}{}\renewcommand{\rhcaDmXgbBarSpan}{\ensuremath{-}9.94}
\providecommand{\rhcaPXgbBarSpan}{}\renewcommand{\rhcaPXgbBarSpan}{\ensuremath{p < 0.001}}
\providecommand{\rhcaMsePctXgbBarSpan}{}\renewcommand{\rhcaMsePctXgbBarSpan}{16.3}
\providecommand{\rhcaMseDmXgbBarSpan}{}\renewcommand{\rhcaMseDmXgbBarSpan}{1.04}
\providecommand{\rhcaMsePXgbBarSpan}{}\renewcommand{\rhcaMsePXgbBarSpan}{\ensuremath{p = 0.30}}
\providecommand{\rhcaQRfHarBarSpan}{}\renewcommand{\rhcaQRfHarBarSpan}{0.1895}
\providecommand{\rhcaMseRfHarBarSpan}{}\renewcommand{\rhcaMseRfHarBarSpan}{1.19}
\providecommand{\rhcaMaeRfHarBarSpan}{}\renewcommand{\rhcaMaeRfHarBarSpan}{2.27}
\providecommand{\rhcaRsqRfHarBarSpan}{}\renewcommand{\rhcaRsqRfHarBarSpan}{63.0}
\providecommand{\rhcaQRfFullBarSpan}{}\renewcommand{\rhcaQRfFullBarSpan}{0.1681}
\providecommand{\rhcaMseRfFullBarSpan}{}\renewcommand{\rhcaMseRfFullBarSpan}{1.32}
\providecommand{\rhcaMaeRfFullBarSpan}{}\renewcommand{\rhcaMaeRfFullBarSpan}{2.27}
\providecommand{\rhcaRsqRfFullBarSpan}{}\renewcommand{\rhcaRsqRfFullBarSpan}{58.9}
\providecommand{\rhcaPctRfBarSpan}{}\renewcommand{\rhcaPctRfBarSpan}{\ensuremath{-}11.3}
\providecommand{\rhcaDmRfBarSpan}{}\renewcommand{\rhcaDmRfBarSpan}{\ensuremath{-}9.74}
\providecommand{\rhcaPRfBarSpan}{}\renewcommand{\rhcaPRfBarSpan}{\ensuremath{p < 0.001}}
\providecommand{\rhcaMsePctRfBarSpan}{}\renewcommand{\rhcaMsePctRfBarSpan}{10.9}
\providecommand{\rhcaMseDmRfBarSpan}{}\renewcommand{\rhcaMseDmRfBarSpan}{1.10}
\providecommand{\rhcaMsePRfBarSpan}{}\renewcommand{\rhcaMsePRfBarSpan}{\ensuremath{p = 0.27}}
\providecommand{\rhcaQOlsHarBarPre}{}\renewcommand{\rhcaQOlsHarBarPre}{0.1864}
\providecommand{\rhcaMseOlsHarBarPre}{}\renewcommand{\rhcaMseOlsHarBarPre}{0.31}
\providecommand{\rhcaMaeOlsHarBarPre}{}\renewcommand{\rhcaMaeOlsHarBarPre}{1.41}
\providecommand{\rhcaRsqOlsHarBarPre}{}\renewcommand{\rhcaRsqOlsHarBarPre}{61.4}
\providecommand{\rhcaQRidgeHarBarPre}{}\renewcommand{\rhcaQRidgeHarBarPre}{0.1864}
\providecommand{\rhcaMseRidgeHarBarPre}{}\renewcommand{\rhcaMseRidgeHarBarPre}{0.31}
\providecommand{\rhcaMaeRidgeHarBarPre}{}\renewcommand{\rhcaMaeRidgeHarBarPre}{1.41}
\providecommand{\rhcaRsqRidgeHarBarPre}{}\renewcommand{\rhcaRsqRidgeHarBarPre}{61.4}
\providecommand{\rhcaQRidgeFullBarPre}{}\renewcommand{\rhcaQRidgeFullBarPre}{0.1726}
\providecommand{\rhcaMseRidgeFullBarPre}{}\renewcommand{\rhcaMseRidgeFullBarPre}{0.28}
\providecommand{\rhcaMaeRidgeFullBarPre}{}\renewcommand{\rhcaMaeRidgeFullBarPre}{1.35}
\providecommand{\rhcaRsqRidgeFullBarPre}{}\renewcommand{\rhcaRsqRidgeFullBarPre}{64.9}
\providecommand{\rhcaPctRidgeBarPre}{}\renewcommand{\rhcaPctRidgeBarPre}{\ensuremath{-}7.4}
\providecommand{\rhcaDmRidgeBarPre}{}\renewcommand{\rhcaDmRidgeBarPre}{\ensuremath{-}6.82}
\providecommand{\rhcaPRidgeBarPre}{}\renewcommand{\rhcaPRidgeBarPre}{\ensuremath{p < 0.001}}
\providecommand{\rhcaMsePctRidgeBarPre}{}\renewcommand{\rhcaMsePctRidgeBarPre}{\ensuremath{-}8.9}
\providecommand{\rhcaMseDmRidgeBarPre}{}\renewcommand{\rhcaMseDmRidgeBarPre}{\ensuremath{-}1.79}
\providecommand{\rhcaMsePRidgeBarPre}{}\renewcommand{\rhcaMsePRidgeBarPre}{\ensuremath{p = 0.074}}
\providecommand{\rhcaQLassoHarBarPre}{}\renewcommand{\rhcaQLassoHarBarPre}{0.1864}
\providecommand{\rhcaMseLassoHarBarPre}{}\renewcommand{\rhcaMseLassoHarBarPre}{0.31}
\providecommand{\rhcaMaeLassoHarBarPre}{}\renewcommand{\rhcaMaeLassoHarBarPre}{1.41}
\providecommand{\rhcaRsqLassoHarBarPre}{}\renewcommand{\rhcaRsqLassoHarBarPre}{61.4}
\providecommand{\rhcaQLassoFullBarPre}{}\renewcommand{\rhcaQLassoFullBarPre}{0.1667}
\providecommand{\rhcaMseLassoFullBarPre}{}\renewcommand{\rhcaMseLassoFullBarPre}{0.28}
\providecommand{\rhcaMaeLassoFullBarPre}{}\renewcommand{\rhcaMaeLassoFullBarPre}{1.35}
\providecommand{\rhcaRsqLassoFullBarPre}{}\renewcommand{\rhcaRsqLassoFullBarPre}{65.6}
\providecommand{\rhcaPctLassoBarPre}{}\renewcommand{\rhcaPctLassoBarPre}{\ensuremath{-}10.6}
\providecommand{\rhcaDmLassoBarPre}{}\renewcommand{\rhcaDmLassoBarPre}{\ensuremath{-}6.88}
\providecommand{\rhcaPLassoBarPre}{}\renewcommand{\rhcaPLassoBarPre}{\ensuremath{p < 0.001}}
\providecommand{\rhcaMsePctLassoBarPre}{}\renewcommand{\rhcaMsePctLassoBarPre}{\ensuremath{-}10.8}
\providecommand{\rhcaMseDmLassoBarPre}{}\renewcommand{\rhcaMseDmLassoBarPre}{\ensuremath{-}1.81}
\providecommand{\rhcaMsePLassoBarPre}{}\renewcommand{\rhcaMsePLassoBarPre}{\ensuremath{p = 0.070}}
\providecommand{\rhcaQEnetHarBarPre}{}\renewcommand{\rhcaQEnetHarBarPre}{0.1864}
\providecommand{\rhcaMseEnetHarBarPre}{}\renewcommand{\rhcaMseEnetHarBarPre}{0.31}
\providecommand{\rhcaMaeEnetHarBarPre}{}\renewcommand{\rhcaMaeEnetHarBarPre}{1.41}
\providecommand{\rhcaRsqEnetHarBarPre}{}\renewcommand{\rhcaRsqEnetHarBarPre}{61.4}
\providecommand{\rhcaQEnetFullBarPre}{}\renewcommand{\rhcaQEnetFullBarPre}{0.1720}
\providecommand{\rhcaMseEnetFullBarPre}{}\renewcommand{\rhcaMseEnetFullBarPre}{0.28}
\providecommand{\rhcaMaeEnetFullBarPre}{}\renewcommand{\rhcaMaeEnetFullBarPre}{1.35}
\providecommand{\rhcaRsqEnetFullBarPre}{}\renewcommand{\rhcaRsqEnetFullBarPre}{65.2}
\providecommand{\rhcaPctEnetBarPre}{}\renewcommand{\rhcaPctEnetBarPre}{\ensuremath{-}7.7}
\providecommand{\rhcaDmEnetBarPre}{}\renewcommand{\rhcaDmEnetBarPre}{\ensuremath{-}6.97}
\providecommand{\rhcaPEnetBarPre}{}\renewcommand{\rhcaPEnetBarPre}{\ensuremath{p < 0.001}}
\providecommand{\rhcaMsePctEnetBarPre}{}\renewcommand{\rhcaMsePctEnetBarPre}{\ensuremath{-}9.7}
\providecommand{\rhcaMseDmEnetBarPre}{}\renewcommand{\rhcaMseDmEnetBarPre}{\ensuremath{-}1.75}
\providecommand{\rhcaMsePEnetBarPre}{}\renewcommand{\rhcaMsePEnetBarPre}{\ensuremath{p = 0.080}}
\providecommand{\rhcaQLgbmHarBarPre}{}\renewcommand{\rhcaQLgbmHarBarPre}{0.1858}
\providecommand{\rhcaMseLgbmHarBarPre}{}\renewcommand{\rhcaMseLgbmHarBarPre}{0.37}
\providecommand{\rhcaMaeLgbmHarBarPre}{}\renewcommand{\rhcaMaeLgbmHarBarPre}{1.43}
\providecommand{\rhcaRsqLgbmHarBarPre}{}\renewcommand{\rhcaRsqLgbmHarBarPre}{53.3}
\providecommand{\rhcaQLgbmFullBarPre}{}\renewcommand{\rhcaQLgbmFullBarPre}{0.1586}
\providecommand{\rhcaMseLgbmFullBarPre}{}\renewcommand{\rhcaMseLgbmFullBarPre}{0.36}
\providecommand{\rhcaMaeLgbmFullBarPre}{}\renewcommand{\rhcaMaeLgbmFullBarPre}{1.36}
\providecommand{\rhcaRsqLgbmFullBarPre}{}\renewcommand{\rhcaRsqLgbmFullBarPre}{55.5}
\providecommand{\rhcaPctLgbmBarPre}{}\renewcommand{\rhcaPctLgbmBarPre}{\ensuremath{-}14.6}
\providecommand{\rhcaDmLgbmBarPre}{}\renewcommand{\rhcaDmLgbmBarPre}{\ensuremath{-}7.91}
\providecommand{\rhcaPLgbmBarPre}{}\renewcommand{\rhcaPLgbmBarPre}{\ensuremath{p < 0.001}}
\providecommand{\rhcaMsePctLgbmBarPre}{}\renewcommand{\rhcaMsePctLgbmBarPre}{\ensuremath{-}4.7}
\providecommand{\rhcaMseDmLgbmBarPre}{}\renewcommand{\rhcaMseDmLgbmBarPre}{\ensuremath{-}2.59}
\providecommand{\rhcaMsePLgbmBarPre}{}\renewcommand{\rhcaMsePLgbmBarPre}{\ensuremath{p = 0.0097}}
\providecommand{\rhcaQXgbHarBarPre}{}\renewcommand{\rhcaQXgbHarBarPre}{0.1864}
\providecommand{\rhcaMseXgbHarBarPre}{}\renewcommand{\rhcaMseXgbHarBarPre}{0.37}
\providecommand{\rhcaMaeXgbHarBarPre}{}\renewcommand{\rhcaMaeXgbHarBarPre}{1.43}
\providecommand{\rhcaRsqXgbHarBarPre}{}\renewcommand{\rhcaRsqXgbHarBarPre}{53.9}
\providecommand{\rhcaQXgbFullBarPre}{}\renewcommand{\rhcaQXgbFullBarPre}{0.1597}
\providecommand{\rhcaMseXgbFullBarPre}{}\renewcommand{\rhcaMseXgbFullBarPre}{0.36}
\providecommand{\rhcaMaeXgbFullBarPre}{}\renewcommand{\rhcaMaeXgbFullBarPre}{1.37}
\providecommand{\rhcaRsqXgbFullBarPre}{}\renewcommand{\rhcaRsqXgbFullBarPre}{55.5}
\providecommand{\rhcaPctXgbBarPre}{}\renewcommand{\rhcaPctXgbBarPre}{\ensuremath{-}14.3}
\providecommand{\rhcaDmXgbBarPre}{}\renewcommand{\rhcaDmXgbBarPre}{\ensuremath{-}7.48}
\providecommand{\rhcaPXgbBarPre}{}\renewcommand{\rhcaPXgbBarPre}{\ensuremath{p < 0.001}}
\providecommand{\rhcaMsePctXgbBarPre}{}\renewcommand{\rhcaMsePctXgbBarPre}{\ensuremath{-}3.4}
\providecommand{\rhcaMseDmXgbBarPre}{}\renewcommand{\rhcaMseDmXgbBarPre}{\ensuremath{-}2.50}
\providecommand{\rhcaMsePXgbBarPre}{}\renewcommand{\rhcaMsePXgbBarPre}{\ensuremath{p = 0.012}}
\providecommand{\rhcaQRfHarBarPre}{}\renewcommand{\rhcaQRfHarBarPre}{0.1865}
\providecommand{\rhcaMseRfHarBarPre}{}\renewcommand{\rhcaMseRfHarBarPre}{0.37}
\providecommand{\rhcaMaeRfHarBarPre}{}\renewcommand{\rhcaMaeRfHarBarPre}{1.43}
\providecommand{\rhcaRsqRfHarBarPre}{}\renewcommand{\rhcaRsqRfHarBarPre}{53.9}
\providecommand{\rhcaQRfFullBarPre}{}\renewcommand{\rhcaQRfFullBarPre}{0.1646}
\providecommand{\rhcaMseRfFullBarPre}{}\renewcommand{\rhcaMseRfFullBarPre}{0.36}
\providecommand{\rhcaMaeRfFullBarPre}{}\renewcommand{\rhcaMaeRfFullBarPre}{1.37}
\providecommand{\rhcaRsqRfFullBarPre}{}\renewcommand{\rhcaRsqRfFullBarPre}{54.6}
\providecommand{\rhcaPctRfBarPre}{}\renewcommand{\rhcaPctRfBarPre}{\ensuremath{-}11.7}
\providecommand{\rhcaDmRfBarPre}{}\renewcommand{\rhcaDmRfBarPre}{\ensuremath{-}7.21}
\providecommand{\rhcaPRfBarPre}{}\renewcommand{\rhcaPRfBarPre}{\ensuremath{p < 0.001}}
\providecommand{\rhcaMsePctRfBarPre}{}\renewcommand{\rhcaMsePctRfBarPre}{\ensuremath{-}1.6}
\providecommand{\rhcaMseDmRfBarPre}{}\renewcommand{\rhcaMseDmRfBarPre}{\ensuremath{-}1.74}
\providecommand{\rhcaMsePRfBarPre}{}\renewcommand{\rhcaMsePRfBarPre}{\ensuremath{p = 0.083}}
\providecommand{\rhcaQOlsHarBarOos}{}\renewcommand{\rhcaQOlsHarBarOos}{0.1933}
\providecommand{\rhcaMseOlsHarBarOos}{}\renewcommand{\rhcaMseOlsHarBarOos}{2.22}
\providecommand{\rhcaMaeOlsHarBarOos}{}\renewcommand{\rhcaMaeOlsHarBarOos}{3.42}
\providecommand{\rhcaRsqOlsHarBarOos}{}\renewcommand{\rhcaRsqOlsHarBarOos}{66.1}
\providecommand{\rhcaQRidgeHarBarOos}{}\renewcommand{\rhcaQRidgeHarBarOos}{0.1933}
\providecommand{\rhcaMseRidgeHarBarOos}{}\renewcommand{\rhcaMseRidgeHarBarOos}{2.22}
\providecommand{\rhcaMaeRidgeHarBarOos}{}\renewcommand{\rhcaMaeRidgeHarBarOos}{3.42}
\providecommand{\rhcaRsqRidgeHarBarOos}{}\renewcommand{\rhcaRsqRidgeHarBarOos}{66.1}
\providecommand{\rhcaQRidgeFullBarOos}{}\renewcommand{\rhcaQRidgeFullBarOos}{0.1852}
\providecommand{\rhcaMseRidgeFullBarOos}{}\renewcommand{\rhcaMseRidgeFullBarOos}{4.93}
\providecommand{\rhcaMaeRidgeFullBarOos}{}\renewcommand{\rhcaMaeRidgeFullBarOos}{3.80}
\providecommand{\rhcaRsqRidgeFullBarOos}{}\renewcommand{\rhcaRsqRidgeFullBarOos}{24.9}
\providecommand{\rhcaPctRidgeBarOos}{}\renewcommand{\rhcaPctRidgeBarOos}{\ensuremath{-}4.2}
\providecommand{\rhcaDmRidgeBarOos}{}\renewcommand{\rhcaDmRidgeBarOos}{\ensuremath{-}2.52}
\providecommand{\rhcaPRidgeBarOos}{}\renewcommand{\rhcaPRidgeBarOos}{\ensuremath{p = 0.012}}
\providecommand{\rhcaMsePctRidgeBarOos}{}\renewcommand{\rhcaMsePctRidgeBarOos}{121.6}
\providecommand{\rhcaMseDmRidgeBarOos}{}\renewcommand{\rhcaMseDmRidgeBarOos}{1.19}
\providecommand{\rhcaMsePRidgeBarOos}{}\renewcommand{\rhcaMsePRidgeBarOos}{\ensuremath{p = 0.24}}
\providecommand{\rhcaQLassoHarBarOos}{}\renewcommand{\rhcaQLassoHarBarOos}{0.1933}
\providecommand{\rhcaMseLassoHarBarOos}{}\renewcommand{\rhcaMseLassoHarBarOos}{2.22}
\providecommand{\rhcaMaeLassoHarBarOos}{}\renewcommand{\rhcaMaeLassoHarBarOos}{3.42}
\providecommand{\rhcaRsqLassoHarBarOos}{}\renewcommand{\rhcaRsqLassoHarBarOos}{66.1}
\providecommand{\rhcaQLassoFullBarOos}{}\renewcommand{\rhcaQLassoFullBarOos}{0.1801}
\providecommand{\rhcaMseLassoFullBarOos}{}\renewcommand{\rhcaMseLassoFullBarOos}{5.31}
\providecommand{\rhcaMaeLassoFullBarOos}{}\renewcommand{\rhcaMaeLassoFullBarOos}{3.86}
\providecommand{\rhcaRsqLassoFullBarOos}{}\renewcommand{\rhcaRsqLassoFullBarOos}{19.1}
\providecommand{\rhcaPctLassoBarOos}{}\renewcommand{\rhcaPctLassoBarOos}{\ensuremath{-}6.8}
\providecommand{\rhcaDmLassoBarOos}{}\renewcommand{\rhcaDmLassoBarOos}{\ensuremath{-}2.70}
\providecommand{\rhcaPLassoBarOos}{}\renewcommand{\rhcaPLassoBarOos}{\ensuremath{p = 0.007}}
\providecommand{\rhcaMsePctLassoBarOos}{}\renewcommand{\rhcaMsePctLassoBarOos}{138.8}
\providecommand{\rhcaMseDmLassoBarOos}{}\renewcommand{\rhcaMseDmLassoBarOos}{1.23}
\providecommand{\rhcaMsePLassoBarOos}{}\renewcommand{\rhcaMsePLassoBarOos}{\ensuremath{p = 0.22}}
\providecommand{\rhcaQEnetHarBarOos}{}\renewcommand{\rhcaQEnetHarBarOos}{0.1933}
\providecommand{\rhcaMseEnetHarBarOos}{}\renewcommand{\rhcaMseEnetHarBarOos}{2.22}
\providecommand{\rhcaMaeEnetHarBarOos}{}\renewcommand{\rhcaMaeEnetHarBarOos}{3.42}
\providecommand{\rhcaRsqEnetHarBarOos}{}\renewcommand{\rhcaRsqEnetHarBarOos}{66.1}
\providecommand{\rhcaQEnetFullBarOos}{}\renewcommand{\rhcaQEnetFullBarOos}{0.1831}
\providecommand{\rhcaMseEnetFullBarOos}{}\renewcommand{\rhcaMseEnetFullBarOos}{5.38}
\providecommand{\rhcaMaeEnetFullBarOos}{}\renewcommand{\rhcaMaeEnetFullBarOos}{3.85}
\providecommand{\rhcaRsqEnetFullBarOos}{}\renewcommand{\rhcaRsqEnetFullBarOos}{18.0}
\providecommand{\rhcaPctEnetBarOos}{}\renewcommand{\rhcaPctEnetBarOos}{\ensuremath{-}5.3}
\providecommand{\rhcaDmEnetBarOos}{}\renewcommand{\rhcaDmEnetBarOos}{\ensuremath{-}3.58}
\providecommand{\rhcaPEnetBarOos}{}\renewcommand{\rhcaPEnetBarOos}{\ensuremath{p < 0.001}}
\providecommand{\rhcaMsePctEnetBarOos}{}\renewcommand{\rhcaMsePctEnetBarOos}{142.1}
\providecommand{\rhcaMseDmEnetBarOos}{}\renewcommand{\rhcaMseDmEnetBarOos}{1.18}
\providecommand{\rhcaMsePEnetBarOos}{}\renewcommand{\rhcaMsePEnetBarOos}{\ensuremath{p = 0.24}}
\providecommand{\rhcaQLgbmHarBarOos}{}\renewcommand{\rhcaQLgbmHarBarOos}{0.1900}
\providecommand{\rhcaMseLgbmHarBarOos}{}\renewcommand{\rhcaMseLgbmHarBarOos}{2.30}
\providecommand{\rhcaMaeLgbmHarBarOos}{}\renewcommand{\rhcaMaeLgbmHarBarOos}{3.43}
\providecommand{\rhcaRsqLgbmHarBarOos}{}\renewcommand{\rhcaRsqLgbmHarBarOos}{65.0}
\providecommand{\rhcaQLgbmFullBarOos}{}\renewcommand{\rhcaQLgbmFullBarOos}{0.1645}
\providecommand{\rhcaMseLgbmFullBarOos}{}\renewcommand{\rhcaMseLgbmFullBarOos}{2.89}
\providecommand{\rhcaMaeLgbmFullBarOos}{}\renewcommand{\rhcaMaeLgbmFullBarOos}{3.55}
\providecommand{\rhcaRsqLgbmFullBarOos}{}\renewcommand{\rhcaRsqLgbmFullBarOos}{56.0}
\providecommand{\rhcaPctLgbmBarOos}{}\renewcommand{\rhcaPctLgbmBarOos}{\ensuremath{-}13.4}
\providecommand{\rhcaDmLgbmBarOos}{}\renewcommand{\rhcaDmLgbmBarOos}{\ensuremath{-}6.42}
\providecommand{\rhcaPLgbmBarOos}{}\renewcommand{\rhcaPLgbmBarOos}{\ensuremath{p < 0.001}}
\providecommand{\rhcaMsePctLgbmBarOos}{}\renewcommand{\rhcaMsePctLgbmBarOos}{25.5}
\providecommand{\rhcaMseDmLgbmBarOos}{}\renewcommand{\rhcaMseDmLgbmBarOos}{1.13}
\providecommand{\rhcaMsePLgbmBarOos}{}\renewcommand{\rhcaMsePLgbmBarOos}{\ensuremath{p = 0.26}}
\providecommand{\rhcaQXgbHarBarOos}{}\renewcommand{\rhcaQXgbHarBarOos}{0.1921}
\providecommand{\rhcaMseXgbHarBarOos}{}\renewcommand{\rhcaMseXgbHarBarOos}{2.26}
\providecommand{\rhcaMaeXgbHarBarOos}{}\renewcommand{\rhcaMaeXgbHarBarOos}{3.42}
\providecommand{\rhcaRsqXgbHarBarOos}{}\renewcommand{\rhcaRsqXgbHarBarOos}{65.6}
\providecommand{\rhcaQXgbFullBarOos}{}\renewcommand{\rhcaQXgbFullBarOos}{0.1634}
\providecommand{\rhcaMseXgbFullBarOos}{}\renewcommand{\rhcaMseXgbFullBarOos}{2.72}
\providecommand{\rhcaMaeXgbFullBarOos}{}\renewcommand{\rhcaMaeXgbFullBarOos}{3.53}
\providecommand{\rhcaRsqXgbFullBarOos}{}\renewcommand{\rhcaRsqXgbFullBarOos}{58.5}
\providecommand{\rhcaPctXgbBarOos}{}\renewcommand{\rhcaPctXgbBarOos}{\ensuremath{-}14.9}
\providecommand{\rhcaDmXgbBarOos}{}\renewcommand{\rhcaDmXgbBarOos}{\ensuremath{-}6.42}
\providecommand{\rhcaPXgbBarOos}{}\renewcommand{\rhcaPXgbBarOos}{\ensuremath{p < 0.001}}
\providecommand{\rhcaMsePctXgbBarOos}{}\renewcommand{\rhcaMsePctXgbBarOos}{20.8}
\providecommand{\rhcaMseDmXgbBarOos}{}\renewcommand{\rhcaMseDmXgbBarOos}{1.15}
\providecommand{\rhcaMsePXgbBarOos}{}\renewcommand{\rhcaMsePXgbBarOos}{\ensuremath{p = 0.25}}
\providecommand{\rhcaQRfHarBarOos}{}\renewcommand{\rhcaQRfHarBarOos}{0.1938}
\providecommand{\rhcaMseRfHarBarOos}{}\renewcommand{\rhcaMseRfHarBarOos}{2.33}
\providecommand{\rhcaMaeRfHarBarOos}{}\renewcommand{\rhcaMaeRfHarBarOos}{3.44}
\providecommand{\rhcaRsqRfHarBarOos}{}\renewcommand{\rhcaRsqRfHarBarOos}{64.5}
\providecommand{\rhcaQRfFullBarOos}{}\renewcommand{\rhcaQRfFullBarOos}{0.1730}
\providecommand{\rhcaMseRfFullBarOos}{}\renewcommand{\rhcaMseRfFullBarOos}{2.65}
\providecommand{\rhcaMaeRfFullBarOos}{}\renewcommand{\rhcaMaeRfFullBarOos}{3.52}
\providecommand{\rhcaRsqRfFullBarOos}{}\renewcommand{\rhcaRsqRfFullBarOos}{59.6}
\providecommand{\rhcaPctRfBarOos}{}\renewcommand{\rhcaPctRfBarOos}{\ensuremath{-}10.7}
\providecommand{\rhcaDmRfBarOos}{}\renewcommand{\rhcaDmRfBarOos}{\ensuremath{-}6.47}
\providecommand{\rhcaPRfBarOos}{}\renewcommand{\rhcaPRfBarOos}{\ensuremath{p < 0.001}}
\providecommand{\rhcaMsePctRfBarOos}{}\renewcommand{\rhcaMsePctRfBarOos}{13.7}
\providecommand{\rhcaMseDmRfBarOos}{}\renewcommand{\rhcaMseDmRfBarOos}{1.19}
\providecommand{\rhcaMsePRfBarOos}{}\renewcommand{\rhcaMsePRfBarOos}{\ensuremath{p = 0.23}}
% --- calibration, Pre / Oos, every bar and 16:00: Mr mean of RV/sigma2, Rm ratio of means, Under %, MzA (x 10^6), MzB, MzBse, MzRsq (levels MZ), MzlA, MzlB, MzlBse, MzlRsq (logs MZ), ShBias / ShSlope / ShDist (Theil shares of the MSE, %), Ql / Qd (QLIKE level / dispersion terms), OverSe (% of SE from over-predictions)
\providecommand{\rhcaMrRidgeHarBarPre}{}\renewcommand{\rhcaMrRidgeHarBarPre}{1.078}
\providecommand{\rhcaRmRidgeHarBarPre}{}\renewcommand{\rhcaRmRidgeHarBarPre}{1.057}
\providecommand{\rhcaUnderRidgeHarBarPre}{}\renewcommand{\rhcaUnderRidgeHarBarPre}{39.9}
\providecommand{\rhcaMzARidgeHarBarPre}{}\renewcommand{\rhcaMzARidgeHarBarPre}{0.29}
\providecommand{\rhcaMzBRidgeHarBarPre}{}\renewcommand{\rhcaMzBRidgeHarBarPre}{0.97}
\providecommand{\rhcaMzBseRidgeHarBarPre}{}\renewcommand{\rhcaMzBseRidgeHarBarPre}{0.11}
\providecommand{\rhcaMzRsqRidgeHarBarPre}{}\renewcommand{\rhcaMzRsqRidgeHarBarPre}{0.48}
\providecommand{\rhcaMzlARidgeHarBarPre}{}\renewcommand{\rhcaMzlARidgeHarBarPre}{0.04}
\providecommand{\rhcaMzlBRidgeHarBarPre}{}\renewcommand{\rhcaMzlBRidgeHarBarPre}{1.01}
\providecommand{\rhcaMzlBseRidgeHarBarPre}{}\renewcommand{\rhcaMzlBseRidgeHarBarPre}{0.01}
\providecommand{\rhcaMzlRsqRidgeHarBarPre}{}\renewcommand{\rhcaMzlRsqRidgeHarBarPre}{0.76}
\providecommand{\rhcaShBiasRidgeHarBarPre}{}\renewcommand{\rhcaShBiasRidgeHarBarPre}{0.1}
\providecommand{\rhcaShSlopeRidgeHarBarPre}{}\renewcommand{\rhcaShSlopeRidgeHarBarPre}{0.1}
\providecommand{\rhcaShDistRidgeHarBarPre}{}\renewcommand{\rhcaShDistRidgeHarBarPre}{99.8}
\providecommand{\rhcaQlRidgeHarBarPre}{}\renewcommand{\rhcaQlRidgeHarBarPre}{0.0029}
\providecommand{\rhcaQdRidgeHarBarPre}{}\renewcommand{\rhcaQdRidgeHarBarPre}{0.1835}
\providecommand{\rhcaOverSeRidgeHarBarPre}{}\renewcommand{\rhcaOverSeRidgeHarBarPre}{20.0}
\providecommand{\rhcaMrRidgeFullBarPre}{}\renewcommand{\rhcaMrRidgeFullBarPre}{1.096}
\providecommand{\rhcaRmRidgeFullBarPre}{}\renewcommand{\rhcaRmRidgeFullBarPre}{1.049}
\providecommand{\rhcaUnderRidgeFullBarPre}{}\renewcommand{\rhcaUnderRidgeFullBarPre}{41.6}
\providecommand{\rhcaMzARidgeFullBarPre}{}\renewcommand{\rhcaMzARidgeFullBarPre}{0.17}
\providecommand{\rhcaMzBRidgeFullBarPre}{}\renewcommand{\rhcaMzBRidgeFullBarPre}{1.00}
\providecommand{\rhcaMzBseRidgeFullBarPre}{}\renewcommand{\rhcaMzBseRidgeFullBarPre}{0.10}
\providecommand{\rhcaMzRsqRidgeFullBarPre}{}\renewcommand{\rhcaMzRsqRidgeFullBarPre}{0.53}
\providecommand{\rhcaMzlARidgeFullBarPre}{}\renewcommand{\rhcaMzlARidgeFullBarPre}{\ensuremath{-}0.22}
\providecommand{\rhcaMzlBRidgeFullBarPre}{}\renewcommand{\rhcaMzlBRidgeFullBarPre}{0.99}
\providecommand{\rhcaMzlBseRidgeFullBarPre}{}\renewcommand{\rhcaMzlBseRidgeFullBarPre}{0.01}
\providecommand{\rhcaMzlRsqRidgeFullBarPre}{}\renewcommand{\rhcaMzlRsqRidgeFullBarPre}{0.78}
\providecommand{\rhcaShBiasRidgeFullBarPre}{}\renewcommand{\rhcaShBiasRidgeFullBarPre}{0.1}
\providecommand{\rhcaShSlopeRidgeFullBarPre}{}\renewcommand{\rhcaShSlopeRidgeFullBarPre}{0.0}
\providecommand{\rhcaShDistRidgeFullBarPre}{}\renewcommand{\rhcaShDistRidgeFullBarPre}{99.9}
\providecommand{\rhcaQlRidgeFullBarPre}{}\renewcommand{\rhcaQlRidgeFullBarPre}{0.0043}
\providecommand{\rhcaQdRidgeFullBarPre}{}\renewcommand{\rhcaQdRidgeFullBarPre}{0.1683}
\providecommand{\rhcaOverSeRidgeFullBarPre}{}\renewcommand{\rhcaOverSeRidgeFullBarPre}{19.6}
\providecommand{\rhcaMrLassoHarBarPre}{}\renewcommand{\rhcaMrLassoHarBarPre}{1.078}
\providecommand{\rhcaRmLassoHarBarPre}{}\renewcommand{\rhcaRmLassoHarBarPre}{1.057}
\providecommand{\rhcaUnderLassoHarBarPre}{}\renewcommand{\rhcaUnderLassoHarBarPre}{39.9}
\providecommand{\rhcaMzALassoHarBarPre}{}\renewcommand{\rhcaMzALassoHarBarPre}{0.29}
\providecommand{\rhcaMzBLassoHarBarPre}{}\renewcommand{\rhcaMzBLassoHarBarPre}{0.97}
\providecommand{\rhcaMzBseLassoHarBarPre}{}\renewcommand{\rhcaMzBseLassoHarBarPre}{0.11}
\providecommand{\rhcaMzRsqLassoHarBarPre}{}\renewcommand{\rhcaMzRsqLassoHarBarPre}{0.48}
\providecommand{\rhcaMzlALassoHarBarPre}{}\renewcommand{\rhcaMzlALassoHarBarPre}{0.04}
\providecommand{\rhcaMzlBLassoHarBarPre}{}\renewcommand{\rhcaMzlBLassoHarBarPre}{1.01}
\providecommand{\rhcaMzlBseLassoHarBarPre}{}\renewcommand{\rhcaMzlBseLassoHarBarPre}{0.01}
\providecommand{\rhcaMzlRsqLassoHarBarPre}{}\renewcommand{\rhcaMzlRsqLassoHarBarPre}{0.76}
\providecommand{\rhcaShBiasLassoHarBarPre}{}\renewcommand{\rhcaShBiasLassoHarBarPre}{0.1}
\providecommand{\rhcaShSlopeLassoHarBarPre}{}\renewcommand{\rhcaShSlopeLassoHarBarPre}{0.1}
\providecommand{\rhcaShDistLassoHarBarPre}{}\renewcommand{\rhcaShDistLassoHarBarPre}{99.8}
\providecommand{\rhcaQlLassoHarBarPre}{}\renewcommand{\rhcaQlLassoHarBarPre}{0.0029}
\providecommand{\rhcaQdLassoHarBarPre}{}\renewcommand{\rhcaQdLassoHarBarPre}{0.1835}
\providecommand{\rhcaOverSeLassoHarBarPre}{}\renewcommand{\rhcaOverSeLassoHarBarPre}{20.0}
\providecommand{\rhcaMrLassoFullBarPre}{}\renewcommand{\rhcaMrLassoFullBarPre}{1.092}
\providecommand{\rhcaRmLassoFullBarPre}{}\renewcommand{\rhcaRmLassoFullBarPre}{1.045}
\providecommand{\rhcaUnderLassoFullBarPre}{}\renewcommand{\rhcaUnderLassoFullBarPre}{42.1}
\providecommand{\rhcaMzALassoFullBarPre}{}\renewcommand{\rhcaMzALassoFullBarPre}{0.16}
\providecommand{\rhcaMzBLassoFullBarPre}{}\renewcommand{\rhcaMzBLassoFullBarPre}{0.99}
\providecommand{\rhcaMzBseLassoFullBarPre}{}\renewcommand{\rhcaMzBseLassoFullBarPre}{0.10}
\providecommand{\rhcaMzRsqLassoFullBarPre}{}\renewcommand{\rhcaMzRsqLassoFullBarPre}{0.54}
\providecommand{\rhcaMzlALassoFullBarPre}{}\renewcommand{\rhcaMzlALassoFullBarPre}{\ensuremath{-}0.25}
\providecommand{\rhcaMzlBLassoFullBarPre}{}\renewcommand{\rhcaMzlBLassoFullBarPre}{0.99}
\providecommand{\rhcaMzlBseLassoFullBarPre}{}\renewcommand{\rhcaMzlBseLassoFullBarPre}{0.01}
\providecommand{\rhcaMzlRsqLassoFullBarPre}{}\renewcommand{\rhcaMzlRsqLassoFullBarPre}{0.78}
\providecommand{\rhcaShBiasLassoFullBarPre}{}\renewcommand{\rhcaShBiasLassoFullBarPre}{0.1}
\providecommand{\rhcaShSlopeLassoFullBarPre}{}\renewcommand{\rhcaShSlopeLassoFullBarPre}{0.0}
\providecommand{\rhcaShDistLassoFullBarPre}{}\renewcommand{\rhcaShDistLassoFullBarPre}{99.9}
\providecommand{\rhcaQlLassoFullBarPre}{}\renewcommand{\rhcaQlLassoFullBarPre}{0.0040}
\providecommand{\rhcaQdLassoFullBarPre}{}\renewcommand{\rhcaQdLassoFullBarPre}{0.1627}
\providecommand{\rhcaOverSeLassoFullBarPre}{}\renewcommand{\rhcaOverSeLassoFullBarPre}{20.1}
\providecommand{\rhcaMrEnetHarBarPre}{}\renewcommand{\rhcaMrEnetHarBarPre}{1.078}
\providecommand{\rhcaRmEnetHarBarPre}{}\renewcommand{\rhcaRmEnetHarBarPre}{1.057}
\providecommand{\rhcaUnderEnetHarBarPre}{}\renewcommand{\rhcaUnderEnetHarBarPre}{39.9}
\providecommand{\rhcaMzAEnetHarBarPre}{}\renewcommand{\rhcaMzAEnetHarBarPre}{0.29}
\providecommand{\rhcaMzBEnetHarBarPre}{}\renewcommand{\rhcaMzBEnetHarBarPre}{0.97}
\providecommand{\rhcaMzBseEnetHarBarPre}{}\renewcommand{\rhcaMzBseEnetHarBarPre}{0.11}
\providecommand{\rhcaMzRsqEnetHarBarPre}{}\renewcommand{\rhcaMzRsqEnetHarBarPre}{0.48}
\providecommand{\rhcaMzlAEnetHarBarPre}{}\renewcommand{\rhcaMzlAEnetHarBarPre}{0.04}
\providecommand{\rhcaMzlBEnetHarBarPre}{}\renewcommand{\rhcaMzlBEnetHarBarPre}{1.01}
\providecommand{\rhcaMzlBseEnetHarBarPre}{}\renewcommand{\rhcaMzlBseEnetHarBarPre}{0.01}
\providecommand{\rhcaMzlRsqEnetHarBarPre}{}\renewcommand{\rhcaMzlRsqEnetHarBarPre}{0.76}
\providecommand{\rhcaShBiasEnetHarBarPre}{}\renewcommand{\rhcaShBiasEnetHarBarPre}{0.1}
\providecommand{\rhcaShSlopeEnetHarBarPre}{}\renewcommand{\rhcaShSlopeEnetHarBarPre}{0.1}
\providecommand{\rhcaShDistEnetHarBarPre}{}\renewcommand{\rhcaShDistEnetHarBarPre}{99.8}
\providecommand{\rhcaQlEnetHarBarPre}{}\renewcommand{\rhcaQlEnetHarBarPre}{0.0029}
\providecommand{\rhcaQdEnetHarBarPre}{}\renewcommand{\rhcaQdEnetHarBarPre}{0.1835}
\providecommand{\rhcaOverSeEnetHarBarPre}{}\renewcommand{\rhcaOverSeEnetHarBarPre}{20.0}
\providecommand{\rhcaMrEnetFullBarPre}{}\renewcommand{\rhcaMrEnetFullBarPre}{1.091}
\providecommand{\rhcaRmEnetFullBarPre}{}\renewcommand{\rhcaRmEnetFullBarPre}{1.049}
\providecommand{\rhcaUnderEnetFullBarPre}{}\renewcommand{\rhcaUnderEnetFullBarPre}{41.5}
\providecommand{\rhcaMzAEnetFullBarPre}{}\renewcommand{\rhcaMzAEnetFullBarPre}{0.16}
\providecommand{\rhcaMzBEnetFullBarPre}{}\renewcommand{\rhcaMzBEnetFullBarPre}{1.00}
\providecommand{\rhcaMzBseEnetFullBarPre}{}\renewcommand{\rhcaMzBseEnetFullBarPre}{0.10}
\providecommand{\rhcaMzRsqEnetFullBarPre}{}\renewcommand{\rhcaMzRsqEnetFullBarPre}{0.53}
\providecommand{\rhcaMzlAEnetFullBarPre}{}\renewcommand{\rhcaMzlAEnetFullBarPre}{\ensuremath{-}0.16}
\providecommand{\rhcaMzlBEnetFullBarPre}{}\renewcommand{\rhcaMzlBEnetFullBarPre}{0.99}
\providecommand{\rhcaMzlBseEnetFullBarPre}{}\renewcommand{\rhcaMzlBseEnetFullBarPre}{0.01}
\providecommand{\rhcaMzlRsqEnetFullBarPre}{}\renewcommand{\rhcaMzlRsqEnetFullBarPre}{0.78}
\providecommand{\rhcaShBiasEnetFullBarPre}{}\renewcommand{\rhcaShBiasEnetFullBarPre}{0.1}
\providecommand{\rhcaShSlopeEnetFullBarPre}{}\renewcommand{\rhcaShSlopeEnetFullBarPre}{0.0}
\providecommand{\rhcaShDistEnetFullBarPre}{}\renewcommand{\rhcaShDistEnetFullBarPre}{99.9}
\providecommand{\rhcaQlEnetFullBarPre}{}\renewcommand{\rhcaQlEnetFullBarPre}{0.0039}
\providecommand{\rhcaQdEnetFullBarPre}{}\renewcommand{\rhcaQdEnetFullBarPre}{0.1681}
\providecommand{\rhcaOverSeEnetFullBarPre}{}\renewcommand{\rhcaOverSeEnetFullBarPre}{19.7}
\providecommand{\rhcaMrLgbmHarBarPre}{}\renewcommand{\rhcaMrLgbmHarBarPre}{1.083}
\providecommand{\rhcaRmLgbmHarBarPre}{}\renewcommand{\rhcaRmLgbmHarBarPre}{1.075}
\providecommand{\rhcaUnderLgbmHarBarPre}{}\renewcommand{\rhcaUnderLgbmHarBarPre}{40.2}
\providecommand{\rhcaMzALgbmHarBarPre}{}\renewcommand{\rhcaMzALgbmHarBarPre}{0.27}
\providecommand{\rhcaMzBLgbmHarBarPre}{}\renewcommand{\rhcaMzBLgbmHarBarPre}{0.99}
\providecommand{\rhcaMzBseLgbmHarBarPre}{}\renewcommand{\rhcaMzBseLgbmHarBarPre}{0.07}
\providecommand{\rhcaMzRsqLgbmHarBarPre}{}\renewcommand{\rhcaMzRsqLgbmHarBarPre}{0.37}
\providecommand{\rhcaMzlALgbmHarBarPre}{}\renewcommand{\rhcaMzlALgbmHarBarPre}{0.04}
\providecommand{\rhcaMzlBLgbmHarBarPre}{}\renewcommand{\rhcaMzlBLgbmHarBarPre}{1.01}
\providecommand{\rhcaMzlBseLgbmHarBarPre}{}\renewcommand{\rhcaMzlBseLgbmHarBarPre}{0.01}
\providecommand{\rhcaMzlRsqLgbmHarBarPre}{}\renewcommand{\rhcaMzlRsqLgbmHarBarPre}{0.76}
\providecommand{\rhcaShBiasLgbmHarBarPre}{}\renewcommand{\rhcaShBiasLgbmHarBarPre}{0.1}
\providecommand{\rhcaShSlopeLgbmHarBarPre}{}\renewcommand{\rhcaShSlopeLgbmHarBarPre}{0.0}
\providecommand{\rhcaShDistLgbmHarBarPre}{}\renewcommand{\rhcaShDistLgbmHarBarPre}{99.8}
\providecommand{\rhcaQlLgbmHarBarPre}{}\renewcommand{\rhcaQlLgbmHarBarPre}{0.0032}
\providecommand{\rhcaQdLgbmHarBarPre}{}\renewcommand{\rhcaQdLgbmHarBarPre}{0.1826}
\providecommand{\rhcaOverSeLgbmHarBarPre}{}\renewcommand{\rhcaOverSeLgbmHarBarPre}{10.3}
\providecommand{\rhcaMrLgbmFullBarPre}{}\renewcommand{\rhcaMrLgbmFullBarPre}{1.069}
\providecommand{\rhcaRmLgbmFullBarPre}{}\renewcommand{\rhcaRmLgbmFullBarPre}{1.060}
\providecommand{\rhcaUnderLgbmFullBarPre}{}\renewcommand{\rhcaUnderLgbmFullBarPre}{41.0}
\providecommand{\rhcaMzALgbmFullBarPre}{}\renewcommand{\rhcaMzALgbmFullBarPre}{0.29}
\providecommand{\rhcaMzBLgbmFullBarPre}{}\renewcommand{\rhcaMzBLgbmFullBarPre}{0.97}
\providecommand{\rhcaMzBseLgbmFullBarPre}{}\renewcommand{\rhcaMzBseLgbmFullBarPre}{0.06}
\providecommand{\rhcaMzRsqLgbmFullBarPre}{}\renewcommand{\rhcaMzRsqLgbmFullBarPre}{0.40}
\providecommand{\rhcaMzlALgbmFullBarPre}{}\renewcommand{\rhcaMzlALgbmFullBarPre}{\ensuremath{-}0.03}
\providecommand{\rhcaMzlBLgbmFullBarPre}{}\renewcommand{\rhcaMzlBLgbmFullBarPre}{1.00}
\providecommand{\rhcaMzlBseLgbmFullBarPre}{}\renewcommand{\rhcaMzlBseLgbmFullBarPre}{0.01}
\providecommand{\rhcaMzlRsqLgbmFullBarPre}{}\renewcommand{\rhcaMzlRsqLgbmFullBarPre}{0.79}
\providecommand{\rhcaShBiasLgbmFullBarPre}{}\renewcommand{\rhcaShBiasLgbmFullBarPre}{0.1}
\providecommand{\rhcaShSlopeLgbmFullBarPre}{}\renewcommand{\rhcaShSlopeLgbmFullBarPre}{0.1}
\providecommand{\rhcaShDistLgbmFullBarPre}{}\renewcommand{\rhcaShDistLgbmFullBarPre}{99.8}
\providecommand{\rhcaQlLgbmFullBarPre}{}\renewcommand{\rhcaQlLgbmFullBarPre}{0.0023}
\providecommand{\rhcaQdLgbmFullBarPre}{}\renewcommand{\rhcaQdLgbmFullBarPre}{0.1563}
\providecommand{\rhcaOverSeLgbmFullBarPre}{}\renewcommand{\rhcaOverSeLgbmFullBarPre}{12.0}
\providecommand{\rhcaMrXgbHarBarPre}{}\renewcommand{\rhcaMrXgbHarBarPre}{1.083}
\providecommand{\rhcaRmXgbHarBarPre}{}\renewcommand{\rhcaRmXgbHarBarPre}{1.073}
\providecommand{\rhcaUnderXgbHarBarPre}{}\renewcommand{\rhcaUnderXgbHarBarPre}{40.2}
\providecommand{\rhcaMzAXgbHarBarPre}{}\renewcommand{\rhcaMzAXgbHarBarPre}{0.31}
\providecommand{\rhcaMzBXgbHarBarPre}{}\renewcommand{\rhcaMzBXgbHarBarPre}{0.97}
\providecommand{\rhcaMzBseXgbHarBarPre}{}\renewcommand{\rhcaMzBseXgbHarBarPre}{0.09}
\providecommand{\rhcaMzRsqXgbHarBarPre}{}\renewcommand{\rhcaMzRsqXgbHarBarPre}{0.38}
\providecommand{\rhcaMzlAXgbHarBarPre}{}\renewcommand{\rhcaMzlAXgbHarBarPre}{0.06}
\providecommand{\rhcaMzlBXgbHarBarPre}{}\renewcommand{\rhcaMzlBXgbHarBarPre}{1.01}
\providecommand{\rhcaMzlBseXgbHarBarPre}{}\renewcommand{\rhcaMzlBseXgbHarBarPre}{0.01}
\providecommand{\rhcaMzlRsqXgbHarBarPre}{}\renewcommand{\rhcaMzlRsqXgbHarBarPre}{0.76}
\providecommand{\rhcaShBiasXgbHarBarPre}{}\renewcommand{\rhcaShBiasXgbHarBarPre}{0.1}
\providecommand{\rhcaShSlopeXgbHarBarPre}{}\renewcommand{\rhcaShSlopeXgbHarBarPre}{0.0}
\providecommand{\rhcaShDistXgbHarBarPre}{}\renewcommand{\rhcaShDistXgbHarBarPre}{99.8}
\providecommand{\rhcaQlXgbHarBarPre}{}\renewcommand{\rhcaQlXgbHarBarPre}{0.0032}
\providecommand{\rhcaQdXgbHarBarPre}{}\renewcommand{\rhcaQdXgbHarBarPre}{0.1832}
\providecommand{\rhcaOverSeXgbHarBarPre}{}\renewcommand{\rhcaOverSeXgbHarBarPre}{12.2}
\providecommand{\rhcaMrXgbFullBarPre}{}\renewcommand{\rhcaMrXgbFullBarPre}{1.067}
\providecommand{\rhcaRmXgbFullBarPre}{}\renewcommand{\rhcaRmXgbFullBarPre}{1.055}
\providecommand{\rhcaUnderXgbFullBarPre}{}\renewcommand{\rhcaUnderXgbFullBarPre}{40.8}
\providecommand{\rhcaMzAXgbFullBarPre}{}\renewcommand{\rhcaMzAXgbFullBarPre}{0.40}
\providecommand{\rhcaMzBXgbFullBarPre}{}\renewcommand{\rhcaMzBXgbFullBarPre}{0.93}
\providecommand{\rhcaMzBseXgbFullBarPre}{}\renewcommand{\rhcaMzBseXgbFullBarPre}{0.09}
\providecommand{\rhcaMzRsqXgbFullBarPre}{}\renewcommand{\rhcaMzRsqXgbFullBarPre}{0.40}
\providecommand{\rhcaMzlAXgbFullBarPre}{}\renewcommand{\rhcaMzlAXgbFullBarPre}{0.00}
\providecommand{\rhcaMzlBXgbFullBarPre}{}\renewcommand{\rhcaMzlBXgbFullBarPre}{1.01}
\providecommand{\rhcaMzlBseXgbFullBarPre}{}\renewcommand{\rhcaMzlBseXgbFullBarPre}{0.01}
\providecommand{\rhcaMzlRsqXgbFullBarPre}{}\renewcommand{\rhcaMzlRsqXgbFullBarPre}{0.79}
\providecommand{\rhcaShBiasXgbFullBarPre}{}\renewcommand{\rhcaShBiasXgbFullBarPre}{0.1}
\providecommand{\rhcaShSlopeXgbFullBarPre}{}\renewcommand{\rhcaShSlopeXgbFullBarPre}{0.4}
\providecommand{\rhcaShDistXgbFullBarPre}{}\renewcommand{\rhcaShDistXgbFullBarPre}{99.5}
\providecommand{\rhcaQlXgbFullBarPre}{}\renewcommand{\rhcaQlXgbFullBarPre}{0.0021}
\providecommand{\rhcaQdXgbFullBarPre}{}\renewcommand{\rhcaQdXgbFullBarPre}{0.1575}
\providecommand{\rhcaOverSeXgbFullBarPre}{}\renewcommand{\rhcaOverSeXgbFullBarPre}{15.1}
\providecommand{\rhcaMrRfHarBarPre}{}\renewcommand{\rhcaMrRfHarBarPre}{1.085}
\providecommand{\rhcaRmRfHarBarPre}{}\renewcommand{\rhcaRmRfHarBarPre}{1.068}
\providecommand{\rhcaUnderRfHarBarPre}{}\renewcommand{\rhcaUnderRfHarBarPre}{40.1}
\providecommand{\rhcaMzARfHarBarPre}{}\renewcommand{\rhcaMzARfHarBarPre}{0.33}
\providecommand{\rhcaMzBRfHarBarPre}{}\renewcommand{\rhcaMzBRfHarBarPre}{0.96}
\providecommand{\rhcaMzBseRfHarBarPre}{}\renewcommand{\rhcaMzBseRfHarBarPre}{0.10}
\providecommand{\rhcaMzRsqRfHarBarPre}{}\renewcommand{\rhcaMzRsqRfHarBarPre}{0.38}
\providecommand{\rhcaMzlARfHarBarPre}{}\renewcommand{\rhcaMzlARfHarBarPre}{\ensuremath{-}0.04}
\providecommand{\rhcaMzlBRfHarBarPre}{}\renewcommand{\rhcaMzlBRfHarBarPre}{1.00}
\providecommand{\rhcaMzlBseRfHarBarPre}{}\renewcommand{\rhcaMzlBseRfHarBarPre}{0.01}
\providecommand{\rhcaMzlRsqRfHarBarPre}{}\renewcommand{\rhcaMzlRsqRfHarBarPre}{0.76}
\providecommand{\rhcaShBiasRfHarBarPre}{}\renewcommand{\rhcaShBiasRfHarBarPre}{0.1}
\providecommand{\rhcaShSlopeRfHarBarPre}{}\renewcommand{\rhcaShSlopeRfHarBarPre}{0.1}
\providecommand{\rhcaShDistRfHarBarPre}{}\renewcommand{\rhcaShDistRfHarBarPre}{99.8}
\providecommand{\rhcaQlRfHarBarPre}{}\renewcommand{\rhcaQlRfHarBarPre}{0.0034}
\providecommand{\rhcaQdRfHarBarPre}{}\renewcommand{\rhcaQdRfHarBarPre}{0.1831}
\providecommand{\rhcaOverSeRfHarBarPre}{}\renewcommand{\rhcaOverSeRfHarBarPre}{12.8}
\providecommand{\rhcaMrRfFullBarPre}{}\renewcommand{\rhcaMrRfFullBarPre}{1.078}
\providecommand{\rhcaRmRfFullBarPre}{}\renewcommand{\rhcaRmRfFullBarPre}{1.057}
\providecommand{\rhcaUnderRfFullBarPre}{}\renewcommand{\rhcaUnderRfFullBarPre}{41.1}
\providecommand{\rhcaMzARfFullBarPre}{}\renewcommand{\rhcaMzARfFullBarPre}{0.43}
\providecommand{\rhcaMzBRfFullBarPre}{}\renewcommand{\rhcaMzBRfFullBarPre}{0.92}
\providecommand{\rhcaMzBseRfFullBarPre}{}\renewcommand{\rhcaMzBseRfFullBarPre}{0.09}
\providecommand{\rhcaMzRsqRfFullBarPre}{}\renewcommand{\rhcaMzRsqRfFullBarPre}{0.39}
\providecommand{\rhcaMzlARfFullBarPre}{}\renewcommand{\rhcaMzlARfFullBarPre}{\ensuremath{-}0.11}
\providecommand{\rhcaMzlBRfFullBarPre}{}\renewcommand{\rhcaMzlBRfFullBarPre}{1.00}
\providecommand{\rhcaMzlBseRfFullBarPre}{}\renewcommand{\rhcaMzlBseRfFullBarPre}{0.01}
\providecommand{\rhcaMzlRsqRfFullBarPre}{}\renewcommand{\rhcaMzlRsqRfFullBarPre}{0.79}
\providecommand{\rhcaShBiasRfFullBarPre}{}\renewcommand{\rhcaShBiasRfFullBarPre}{0.1}
\providecommand{\rhcaShSlopeRfFullBarPre}{}\renewcommand{\rhcaShSlopeRfFullBarPre}{0.5}
\providecommand{\rhcaShDistRfFullBarPre}{}\renewcommand{\rhcaShDistRfFullBarPre}{99.5}
\providecommand{\rhcaQlRfFullBarPre}{}\renewcommand{\rhcaQlRfFullBarPre}{0.0029}
\providecommand{\rhcaQdRfFullBarPre}{}\renewcommand{\rhcaQdRfFullBarPre}{0.1617}
\providecommand{\rhcaOverSeRfFullBarPre}{}\renewcommand{\rhcaOverSeRfFullBarPre}{14.7}
\providecommand{\rhcaMrRidgeHarClosePre}{}\renewcommand{\rhcaMrRidgeHarClosePre}{1.121}
\providecommand{\rhcaRmRidgeHarClosePre}{}\renewcommand{\rhcaRmRidgeHarClosePre}{1.097}
\providecommand{\rhcaUnderRidgeHarClosePre}{}\renewcommand{\rhcaUnderRidgeHarClosePre}{46.7}
\providecommand{\rhcaMzARidgeHarClosePre}{}\renewcommand{\rhcaMzARidgeHarClosePre}{\ensuremath{-}0.16}
\providecommand{\rhcaMzBRidgeHarClosePre}{}\renewcommand{\rhcaMzBRidgeHarClosePre}{1.14}
\providecommand{\rhcaMzBseRidgeHarClosePre}{}\renewcommand{\rhcaMzBseRidgeHarClosePre}{0.10}
\providecommand{\rhcaMzRsqRidgeHarClosePre}{}\renewcommand{\rhcaMzRsqRidgeHarClosePre}{0.64}
\providecommand{\rhcaMzlARidgeHarClosePre}{}\renewcommand{\rhcaMzlARidgeHarClosePre}{\ensuremath{-}0.68}
\providecommand{\rhcaMzlBRidgeHarClosePre}{}\renewcommand{\rhcaMzlBRidgeHarClosePre}{0.95}
\providecommand{\rhcaMzlBseRidgeHarClosePre}{}\renewcommand{\rhcaMzlBseRidgeHarClosePre}{0.02}
\providecommand{\rhcaMzlRsqRidgeHarClosePre}{}\renewcommand{\rhcaMzlRsqRidgeHarClosePre}{0.78}
\providecommand{\rhcaShBiasRidgeHarClosePre}{}\renewcommand{\rhcaShBiasRidgeHarClosePre}{0.5}
\providecommand{\rhcaShSlopeRidgeHarClosePre}{}\renewcommand{\rhcaShSlopeRidgeHarClosePre}{2.5}
\providecommand{\rhcaShDistRidgeHarClosePre}{}\renewcommand{\rhcaShDistRidgeHarClosePre}{97.0}
\providecommand{\rhcaQlRidgeHarClosePre}{}\renewcommand{\rhcaQlRidgeHarClosePre}{0.0068}
\providecommand{\rhcaQdRidgeHarClosePre}{}\renewcommand{\rhcaQdRidgeHarClosePre}{0.1297}
\providecommand{\rhcaOverSeRidgeHarClosePre}{}\renewcommand{\rhcaOverSeRidgeHarClosePre}{13.1}
\providecommand{\rhcaMrRidgeFullClosePre}{}\renewcommand{\rhcaMrRidgeFullClosePre}{1.137}
\providecommand{\rhcaRmRidgeFullClosePre}{}\renewcommand{\rhcaRmRidgeFullClosePre}{1.061}
\providecommand{\rhcaUnderRidgeFullClosePre}{}\renewcommand{\rhcaUnderRidgeFullClosePre}{47.2}
\providecommand{\rhcaMzARidgeFullClosePre}{}\renewcommand{\rhcaMzARidgeFullClosePre}{\ensuremath{-}0.02}
\providecommand{\rhcaMzBRidgeFullClosePre}{}\renewcommand{\rhcaMzBRidgeFullClosePre}{1.07}
\providecommand{\rhcaMzBseRidgeFullClosePre}{}\renewcommand{\rhcaMzBseRidgeFullClosePre}{0.09}
\providecommand{\rhcaMzRsqRidgeFullClosePre}{}\renewcommand{\rhcaMzRsqRidgeFullClosePre}{0.68}
\providecommand{\rhcaMzlARidgeFullClosePre}{}\renewcommand{\rhcaMzlARidgeFullClosePre}{\ensuremath{-}0.93}
\providecommand{\rhcaMzlBRidgeFullClosePre}{}\renewcommand{\rhcaMzlBRidgeFullClosePre}{0.93}
\providecommand{\rhcaMzlBseRidgeFullClosePre}{}\renewcommand{\rhcaMzlBseRidgeFullClosePre}{0.02}
\providecommand{\rhcaMzlRsqRidgeFullClosePre}{}\renewcommand{\rhcaMzlRsqRidgeFullClosePre}{0.80}
\providecommand{\rhcaShBiasRidgeFullClosePre}{}\renewcommand{\rhcaShBiasRidgeFullClosePre}{0.2}
\providecommand{\rhcaShSlopeRidgeFullClosePre}{}\renewcommand{\rhcaShSlopeRidgeFullClosePre}{0.8}
\providecommand{\rhcaShDistRidgeFullClosePre}{}\renewcommand{\rhcaShDistRidgeFullClosePre}{99.0}
\providecommand{\rhcaQlRidgeFullClosePre}{}\renewcommand{\rhcaQlRidgeFullClosePre}{0.0086}
\providecommand{\rhcaQdRidgeFullClosePre}{}\renewcommand{\rhcaQdRidgeFullClosePre}{0.1222}
\providecommand{\rhcaOverSeRidgeFullClosePre}{}\renewcommand{\rhcaOverSeRidgeFullClosePre}{18.0}
\providecommand{\rhcaMrLassoHarClosePre}{}\renewcommand{\rhcaMrLassoHarClosePre}{1.121}
\providecommand{\rhcaRmLassoHarClosePre}{}\renewcommand{\rhcaRmLassoHarClosePre}{1.097}
\providecommand{\rhcaUnderLassoHarClosePre}{}\renewcommand{\rhcaUnderLassoHarClosePre}{46.7}
\providecommand{\rhcaMzALassoHarClosePre}{}\renewcommand{\rhcaMzALassoHarClosePre}{\ensuremath{-}0.16}
\providecommand{\rhcaMzBLassoHarClosePre}{}\renewcommand{\rhcaMzBLassoHarClosePre}{1.14}
\providecommand{\rhcaMzBseLassoHarClosePre}{}\renewcommand{\rhcaMzBseLassoHarClosePre}{0.10}
\providecommand{\rhcaMzRsqLassoHarClosePre}{}\renewcommand{\rhcaMzRsqLassoHarClosePre}{0.64}
\providecommand{\rhcaMzlALassoHarClosePre}{}\renewcommand{\rhcaMzlALassoHarClosePre}{\ensuremath{-}0.68}
\providecommand{\rhcaMzlBLassoHarClosePre}{}\renewcommand{\rhcaMzlBLassoHarClosePre}{0.95}
\providecommand{\rhcaMzlBseLassoHarClosePre}{}\renewcommand{\rhcaMzlBseLassoHarClosePre}{0.02}
\providecommand{\rhcaMzlRsqLassoHarClosePre}{}\renewcommand{\rhcaMzlRsqLassoHarClosePre}{0.78}
\providecommand{\rhcaShBiasLassoHarClosePre}{}\renewcommand{\rhcaShBiasLassoHarClosePre}{0.5}
\providecommand{\rhcaShSlopeLassoHarClosePre}{}\renewcommand{\rhcaShSlopeLassoHarClosePre}{2.5}
\providecommand{\rhcaShDistLassoHarClosePre}{}\renewcommand{\rhcaShDistLassoHarClosePre}{97.0}
\providecommand{\rhcaQlLassoHarClosePre}{}\renewcommand{\rhcaQlLassoHarClosePre}{0.0068}
\providecommand{\rhcaQdLassoHarClosePre}{}\renewcommand{\rhcaQdLassoHarClosePre}{0.1297}
\providecommand{\rhcaOverSeLassoHarClosePre}{}\renewcommand{\rhcaOverSeLassoHarClosePre}{13.1}
\providecommand{\rhcaMrLassoFullClosePre}{}\renewcommand{\rhcaMrLassoFullClosePre}{1.148}
\providecommand{\rhcaRmLassoFullClosePre}{}\renewcommand{\rhcaRmLassoFullClosePre}{1.059}
\providecommand{\rhcaUnderLassoFullClosePre}{}\renewcommand{\rhcaUnderLassoFullClosePre}{48.6}
\providecommand{\rhcaMzALassoFullClosePre}{}\renewcommand{\rhcaMzALassoFullClosePre}{\ensuremath{-}0.01}
\providecommand{\rhcaMzBLassoFullClosePre}{}\renewcommand{\rhcaMzBLassoFullClosePre}{1.06}
\providecommand{\rhcaMzBseLassoFullClosePre}{}\renewcommand{\rhcaMzBseLassoFullClosePre}{0.09}
\providecommand{\rhcaMzRsqLassoFullClosePre}{}\renewcommand{\rhcaMzRsqLassoFullClosePre}{0.69}
\providecommand{\rhcaMzlALassoFullClosePre}{}\renewcommand{\rhcaMzlALassoFullClosePre}{\ensuremath{-}1.01}
\providecommand{\rhcaMzlBLassoFullClosePre}{}\renewcommand{\rhcaMzlBLassoFullClosePre}{0.92}
\providecommand{\rhcaMzlBseLassoFullClosePre}{}\renewcommand{\rhcaMzlBseLassoFullClosePre}{0.02}
\providecommand{\rhcaMzlRsqLassoFullClosePre}{}\renewcommand{\rhcaMzlRsqLassoFullClosePre}{0.80}
\providecommand{\rhcaShBiasLassoFullClosePre}{}\renewcommand{\rhcaShBiasLassoFullClosePre}{0.2}
\providecommand{\rhcaShSlopeLassoFullClosePre}{}\renewcommand{\rhcaShSlopeLassoFullClosePre}{0.7}
\providecommand{\rhcaShDistLassoFullClosePre}{}\renewcommand{\rhcaShDistLassoFullClosePre}{99.1}
\providecommand{\rhcaQlLassoFullClosePre}{}\renewcommand{\rhcaQlLassoFullClosePre}{0.0100}
\providecommand{\rhcaQdLassoFullClosePre}{}\renewcommand{\rhcaQdLassoFullClosePre}{0.1233}
\providecommand{\rhcaOverSeLassoFullClosePre}{}\renewcommand{\rhcaOverSeLassoFullClosePre}{18.4}
\providecommand{\rhcaMrEnetHarClosePre}{}\renewcommand{\rhcaMrEnetHarClosePre}{1.121}
\providecommand{\rhcaRmEnetHarClosePre}{}\renewcommand{\rhcaRmEnetHarClosePre}{1.097}
\providecommand{\rhcaUnderEnetHarClosePre}{}\renewcommand{\rhcaUnderEnetHarClosePre}{46.7}
\providecommand{\rhcaMzAEnetHarClosePre}{}\renewcommand{\rhcaMzAEnetHarClosePre}{\ensuremath{-}0.16}
\providecommand{\rhcaMzBEnetHarClosePre}{}\renewcommand{\rhcaMzBEnetHarClosePre}{1.14}
\providecommand{\rhcaMzBseEnetHarClosePre}{}\renewcommand{\rhcaMzBseEnetHarClosePre}{0.10}
\providecommand{\rhcaMzRsqEnetHarClosePre}{}\renewcommand{\rhcaMzRsqEnetHarClosePre}{0.64}
\providecommand{\rhcaMzlAEnetHarClosePre}{}\renewcommand{\rhcaMzlAEnetHarClosePre}{\ensuremath{-}0.68}
\providecommand{\rhcaMzlBEnetHarClosePre}{}\renewcommand{\rhcaMzlBEnetHarClosePre}{0.95}
\providecommand{\rhcaMzlBseEnetHarClosePre}{}\renewcommand{\rhcaMzlBseEnetHarClosePre}{0.02}
\providecommand{\rhcaMzlRsqEnetHarClosePre}{}\renewcommand{\rhcaMzlRsqEnetHarClosePre}{0.78}
\providecommand{\rhcaShBiasEnetHarClosePre}{}\renewcommand{\rhcaShBiasEnetHarClosePre}{0.5}
\providecommand{\rhcaShSlopeEnetHarClosePre}{}\renewcommand{\rhcaShSlopeEnetHarClosePre}{2.5}
\providecommand{\rhcaShDistEnetHarClosePre}{}\renewcommand{\rhcaShDistEnetHarClosePre}{97.0}
\providecommand{\rhcaQlEnetHarClosePre}{}\renewcommand{\rhcaQlEnetHarClosePre}{0.0068}
\providecommand{\rhcaQdEnetHarClosePre}{}\renewcommand{\rhcaQdEnetHarClosePre}{0.1297}
\providecommand{\rhcaOverSeEnetHarClosePre}{}\renewcommand{\rhcaOverSeEnetHarClosePre}{13.1}
\providecommand{\rhcaMrEnetFullClosePre}{}\renewcommand{\rhcaMrEnetFullClosePre}{1.138}
\providecommand{\rhcaRmEnetFullClosePre}{}\renewcommand{\rhcaRmEnetFullClosePre}{1.067}
\providecommand{\rhcaUnderEnetFullClosePre}{}\renewcommand{\rhcaUnderEnetFullClosePre}{47.3}
\providecommand{\rhcaMzAEnetFullClosePre}{}\renewcommand{\rhcaMzAEnetFullClosePre}{\ensuremath{-}0.05}
\providecommand{\rhcaMzBEnetFullClosePre}{}\renewcommand{\rhcaMzBEnetFullClosePre}{1.08}
\providecommand{\rhcaMzBseEnetFullClosePre}{}\renewcommand{\rhcaMzBseEnetFullClosePre}{0.09}
\providecommand{\rhcaMzRsqEnetFullClosePre}{}\renewcommand{\rhcaMzRsqEnetFullClosePre}{0.69}
\providecommand{\rhcaMzlAEnetFullClosePre}{}\renewcommand{\rhcaMzlAEnetFullClosePre}{\ensuremath{-}0.90}
\providecommand{\rhcaMzlBEnetFullClosePre}{}\renewcommand{\rhcaMzlBEnetFullClosePre}{0.93}
\providecommand{\rhcaMzlBseEnetFullClosePre}{}\renewcommand{\rhcaMzlBseEnetFullClosePre}{0.02}
\providecommand{\rhcaMzlRsqEnetFullClosePre}{}\renewcommand{\rhcaMzlRsqEnetFullClosePre}{0.80}
\providecommand{\rhcaShBiasEnetFullClosePre}{}\renewcommand{\rhcaShBiasEnetFullClosePre}{0.3}
\providecommand{\rhcaShSlopeEnetFullClosePre}{}\renewcommand{\rhcaShSlopeEnetFullClosePre}{1.2}
\providecommand{\rhcaShDistEnetFullClosePre}{}\renewcommand{\rhcaShDistEnetFullClosePre}{98.5}
\providecommand{\rhcaQlEnetFullClosePre}{}\renewcommand{\rhcaQlEnetFullClosePre}{0.0087}
\providecommand{\rhcaQdEnetFullClosePre}{}\renewcommand{\rhcaQdEnetFullClosePre}{0.1213}
\providecommand{\rhcaOverSeEnetFullClosePre}{}\renewcommand{\rhcaOverSeEnetFullClosePre}{17.3}
\providecommand{\rhcaMrLgbmHarClosePre}{}\renewcommand{\rhcaMrLgbmHarClosePre}{1.068}
\providecommand{\rhcaRmLgbmHarClosePre}{}\renewcommand{\rhcaRmLgbmHarClosePre}{1.062}
\providecommand{\rhcaUnderLgbmHarClosePre}{}\renewcommand{\rhcaUnderLgbmHarClosePre}{42.9}
\providecommand{\rhcaMzALgbmHarClosePre}{}\renewcommand{\rhcaMzALgbmHarClosePre}{0.00}
\providecommand{\rhcaMzBLgbmHarClosePre}{}\renewcommand{\rhcaMzBLgbmHarClosePre}{1.06}
\providecommand{\rhcaMzBseLgbmHarClosePre}{}\renewcommand{\rhcaMzBseLgbmHarClosePre}{0.10}
\providecommand{\rhcaMzRsqLgbmHarClosePre}{}\renewcommand{\rhcaMzRsqLgbmHarClosePre}{0.55}
\providecommand{\rhcaMzlALgbmHarClosePre}{}\renewcommand{\rhcaMzlALgbmHarClosePre}{\ensuremath{-}0.59}
\providecommand{\rhcaMzlBLgbmHarClosePre}{}\renewcommand{\rhcaMzlBLgbmHarClosePre}{0.96}
\providecommand{\rhcaMzlBseLgbmHarClosePre}{}\renewcommand{\rhcaMzlBseLgbmHarClosePre}{0.02}
\providecommand{\rhcaMzlRsqLgbmHarClosePre}{}\renewcommand{\rhcaMzlRsqLgbmHarClosePre}{0.78}
\providecommand{\rhcaShBiasLgbmHarClosePre}{}\renewcommand{\rhcaShBiasLgbmHarClosePre}{0.2}
\providecommand{\rhcaShSlopeLgbmHarClosePre}{}\renewcommand{\rhcaShSlopeLgbmHarClosePre}{0.4}
\providecommand{\rhcaShDistLgbmHarClosePre}{}\renewcommand{\rhcaShDistLgbmHarClosePre}{99.4}
\providecommand{\rhcaQlLgbmHarClosePre}{}\renewcommand{\rhcaQlLgbmHarClosePre}{0.0022}
\providecommand{\rhcaQdLgbmHarClosePre}{}\renewcommand{\rhcaQdLgbmHarClosePre}{0.1310}
\providecommand{\rhcaOverSeLgbmHarClosePre}{}\renewcommand{\rhcaOverSeLgbmHarClosePre}{14.4}
\providecommand{\rhcaMrLgbmFullClosePre}{}\renewcommand{\rhcaMrLgbmFullClosePre}{1.101}
\providecommand{\rhcaRmLgbmFullClosePre}{}\renewcommand{\rhcaRmLgbmFullClosePre}{1.042}
\providecommand{\rhcaUnderLgbmFullClosePre}{}\renewcommand{\rhcaUnderLgbmFullClosePre}{45.9}
\providecommand{\rhcaMzALgbmFullClosePre}{}\renewcommand{\rhcaMzALgbmFullClosePre}{0.26}
\providecommand{\rhcaMzBLgbmFullClosePre}{}\renewcommand{\rhcaMzBLgbmFullClosePre}{0.98}
\providecommand{\rhcaMzBseLgbmFullClosePre}{}\renewcommand{\rhcaMzBseLgbmFullClosePre}{0.07}
\providecommand{\rhcaMzRsqLgbmFullClosePre}{}\renewcommand{\rhcaMzRsqLgbmFullClosePre}{0.58}
\providecommand{\rhcaMzlALgbmFullClosePre}{}\renewcommand{\rhcaMzlALgbmFullClosePre}{\ensuremath{-}0.89}
\providecommand{\rhcaMzlBLgbmFullClosePre}{}\renewcommand{\rhcaMzlBLgbmFullClosePre}{0.93}
\providecommand{\rhcaMzlBseLgbmFullClosePre}{}\renewcommand{\rhcaMzlBseLgbmFullClosePre}{0.02}
\providecommand{\rhcaMzlRsqLgbmFullClosePre}{}\renewcommand{\rhcaMzlRsqLgbmFullClosePre}{0.80}
\providecommand{\rhcaShBiasLgbmFullClosePre}{}\renewcommand{\rhcaShBiasLgbmFullClosePre}{0.1}
\providecommand{\rhcaShSlopeLgbmFullClosePre}{}\renewcommand{\rhcaShSlopeLgbmFullClosePre}{0.1}
\providecommand{\rhcaShDistLgbmFullClosePre}{}\renewcommand{\rhcaShDistLgbmFullClosePre}{99.9}
\providecommand{\rhcaQlLgbmFullClosePre}{}\renewcommand{\rhcaQlLgbmFullClosePre}{0.0048}
\providecommand{\rhcaQdLgbmFullClosePre}{}\renewcommand{\rhcaQdLgbmFullClosePre}{0.1239}
\providecommand{\rhcaOverSeLgbmFullClosePre}{}\renewcommand{\rhcaOverSeLgbmFullClosePre}{17.3}
\providecommand{\rhcaMrXgbHarClosePre}{}\renewcommand{\rhcaMrXgbHarClosePre}{1.064}
\providecommand{\rhcaRmXgbHarClosePre}{}\renewcommand{\rhcaRmXgbHarClosePre}{1.058}
\providecommand{\rhcaUnderXgbHarClosePre}{}\renewcommand{\rhcaUnderXgbHarClosePre}{42.6}
\providecommand{\rhcaMzAXgbHarClosePre}{}\renewcommand{\rhcaMzAXgbHarClosePre}{\ensuremath{-}0.12}
\providecommand{\rhcaMzBXgbHarClosePre}{}\renewcommand{\rhcaMzBXgbHarClosePre}{1.09}
\providecommand{\rhcaMzBseXgbHarClosePre}{}\renewcommand{\rhcaMzBseXgbHarClosePre}{0.09}
\providecommand{\rhcaMzRsqXgbHarClosePre}{}\renewcommand{\rhcaMzRsqXgbHarClosePre}{0.58}
\providecommand{\rhcaMzlAXgbHarClosePre}{}\renewcommand{\rhcaMzlAXgbHarClosePre}{\ensuremath{-}0.58}
\providecommand{\rhcaMzlBXgbHarClosePre}{}\renewcommand{\rhcaMzlBXgbHarClosePre}{0.96}
\providecommand{\rhcaMzlBseXgbHarClosePre}{}\renewcommand{\rhcaMzlBseXgbHarClosePre}{0.02}
\providecommand{\rhcaMzlRsqXgbHarClosePre}{}\renewcommand{\rhcaMzlRsqXgbHarClosePre}{0.78}
\providecommand{\rhcaShBiasXgbHarClosePre}{}\renewcommand{\rhcaShBiasXgbHarClosePre}{0.2}
\providecommand{\rhcaShSlopeXgbHarClosePre}{}\renewcommand{\rhcaShSlopeXgbHarClosePre}{0.9}
\providecommand{\rhcaShDistXgbHarClosePre}{}\renewcommand{\rhcaShDistXgbHarClosePre}{99.0}
\providecommand{\rhcaQlXgbHarClosePre}{}\renewcommand{\rhcaQlXgbHarClosePre}{0.0020}
\providecommand{\rhcaQdXgbHarClosePre}{}\renewcommand{\rhcaQdXgbHarClosePre}{0.1311}
\providecommand{\rhcaOverSeXgbHarClosePre}{}\renewcommand{\rhcaOverSeXgbHarClosePre}{14.2}
\providecommand{\rhcaMrXgbFullClosePre}{}\renewcommand{\rhcaMrXgbFullClosePre}{1.098}
\providecommand{\rhcaRmXgbFullClosePre}{}\renewcommand{\rhcaRmXgbFullClosePre}{1.037}
\providecommand{\rhcaUnderXgbFullClosePre}{}\renewcommand{\rhcaUnderXgbFullClosePre}{45.6}
\providecommand{\rhcaMzAXgbFullClosePre}{}\renewcommand{\rhcaMzAXgbFullClosePre}{0.19}
\providecommand{\rhcaMzBXgbFullClosePre}{}\renewcommand{\rhcaMzBXgbFullClosePre}{0.99}
\providecommand{\rhcaMzBseXgbFullClosePre}{}\renewcommand{\rhcaMzBseXgbFullClosePre}{0.08}
\providecommand{\rhcaMzRsqXgbFullClosePre}{}\renewcommand{\rhcaMzRsqXgbFullClosePre}{0.62}
\providecommand{\rhcaMzlAXgbFullClosePre}{}\renewcommand{\rhcaMzlAXgbFullClosePre}{\ensuremath{-}0.87}
\providecommand{\rhcaMzlBXgbFullClosePre}{}\renewcommand{\rhcaMzlBXgbFullClosePre}{0.94}
\providecommand{\rhcaMzlBseXgbFullClosePre}{}\renewcommand{\rhcaMzlBseXgbFullClosePre}{0.02}
\providecommand{\rhcaMzlRsqXgbFullClosePre}{}\renewcommand{\rhcaMzlRsqXgbFullClosePre}{0.80}
\providecommand{\rhcaShBiasXgbFullClosePre}{}\renewcommand{\rhcaShBiasXgbFullClosePre}{0.1}
\providecommand{\rhcaShSlopeXgbFullClosePre}{}\renewcommand{\rhcaShSlopeXgbFullClosePre}{0.0}
\providecommand{\rhcaShDistXgbFullClosePre}{}\renewcommand{\rhcaShDistXgbFullClosePre}{99.9}
\providecommand{\rhcaQlXgbFullClosePre}{}\renewcommand{\rhcaQlXgbFullClosePre}{0.0045}
\providecommand{\rhcaQdXgbFullClosePre}{}\renewcommand{\rhcaQdXgbFullClosePre}{0.1232}
\providecommand{\rhcaOverSeXgbFullClosePre}{}\renewcommand{\rhcaOverSeXgbFullClosePre}{18.8}
\providecommand{\rhcaMrRfHarClosePre}{}\renewcommand{\rhcaMrRfHarClosePre}{1.091}
\providecommand{\rhcaRmRfHarClosePre}{}\renewcommand{\rhcaRmRfHarClosePre}{1.073}
\providecommand{\rhcaUnderRfHarClosePre}{}\renewcommand{\rhcaUnderRfHarClosePre}{43.6}
\providecommand{\rhcaMzARfHarClosePre}{}\renewcommand{\rhcaMzARfHarClosePre}{\ensuremath{-}0.13}
\providecommand{\rhcaMzBRfHarClosePre}{}\renewcommand{\rhcaMzBRfHarClosePre}{1.11}
\providecommand{\rhcaMzBseRfHarClosePre}{}\renewcommand{\rhcaMzBseRfHarClosePre}{0.10}
\providecommand{\rhcaMzRsqRfHarClosePre}{}\renewcommand{\rhcaMzRsqRfHarClosePre}{0.59}
\providecommand{\rhcaMzlARfHarClosePre}{}\renewcommand{\rhcaMzlARfHarClosePre}{\ensuremath{-}0.68}
\providecommand{\rhcaMzlBRfHarClosePre}{}\renewcommand{\rhcaMzlBRfHarClosePre}{0.95}
\providecommand{\rhcaMzlBseRfHarClosePre}{}\renewcommand{\rhcaMzlBseRfHarClosePre}{0.02}
\providecommand{\rhcaMzlRsqRfHarClosePre}{}\renewcommand{\rhcaMzlRsqRfHarClosePre}{0.78}
\providecommand{\rhcaShBiasRfHarClosePre}{}\renewcommand{\rhcaShBiasRfHarClosePre}{0.2}
\providecommand{\rhcaShSlopeRfHarClosePre}{}\renewcommand{\rhcaShSlopeRfHarClosePre}{1.3}
\providecommand{\rhcaShDistRfHarClosePre}{}\renewcommand{\rhcaShDistRfHarClosePre}{98.5}
\providecommand{\rhcaQlRfHarClosePre}{}\renewcommand{\rhcaQlRfHarClosePre}{0.0039}
\providecommand{\rhcaQdRfHarClosePre}{}\renewcommand{\rhcaQdRfHarClosePre}{0.1325}
\providecommand{\rhcaOverSeRfHarClosePre}{}\renewcommand{\rhcaOverSeRfHarClosePre}{14.1}
\providecommand{\rhcaMrRfFullClosePre}{}\renewcommand{\rhcaMrRfFullClosePre}{1.119}
\providecommand{\rhcaRmRfFullClosePre}{}\renewcommand{\rhcaRmRfFullClosePre}{1.037}
\providecommand{\rhcaUnderRfFullClosePre}{}\renewcommand{\rhcaUnderRfFullClosePre}{46.1}
\providecommand{\rhcaMzARfFullClosePre}{}\renewcommand{\rhcaMzARfFullClosePre}{0.25}
\providecommand{\rhcaMzBRfFullClosePre}{}\renewcommand{\rhcaMzBRfFullClosePre}{0.98}
\providecommand{\rhcaMzBseRfFullClosePre}{}\renewcommand{\rhcaMzBseRfFullClosePre}{0.08}
\providecommand{\rhcaMzRsqRfFullClosePre}{}\renewcommand{\rhcaMzRsqRfFullClosePre}{0.61}
\providecommand{\rhcaMzlARfFullClosePre}{}\renewcommand{\rhcaMzlARfFullClosePre}{\ensuremath{-}1.04}
\providecommand{\rhcaMzlBRfFullClosePre}{}\renewcommand{\rhcaMzlBRfFullClosePre}{0.92}
\providecommand{\rhcaMzlBseRfFullClosePre}{}\renewcommand{\rhcaMzlBseRfFullClosePre}{0.02}
\providecommand{\rhcaMzlRsqRfFullClosePre}{}\renewcommand{\rhcaMzlRsqRfFullClosePre}{0.79}
\providecommand{\rhcaShBiasRfFullClosePre}{}\renewcommand{\rhcaShBiasRfFullClosePre}{0.1}
\providecommand{\rhcaShSlopeRfFullClosePre}{}\renewcommand{\rhcaShSlopeRfFullClosePre}{0.1}
\providecommand{\rhcaShDistRfFullClosePre}{}\renewcommand{\rhcaShDistRfFullClosePre}{99.8}
\providecommand{\rhcaQlRfFullClosePre}{}\renewcommand{\rhcaQlRfFullClosePre}{0.0065}
\providecommand{\rhcaQdRfFullClosePre}{}\renewcommand{\rhcaQdRfFullClosePre}{0.1270}
\providecommand{\rhcaOverSeRfFullClosePre}{}\renewcommand{\rhcaOverSeRfFullClosePre}{18.9}
\providecommand{\rhcaMrRidgeHarBarOos}{}\renewcommand{\rhcaMrRidgeHarBarOos}{1.061}
\providecommand{\rhcaRmRidgeHarBarOos}{}\renewcommand{\rhcaRmRidgeHarBarOos}{1.045}
\providecommand{\rhcaUnderRidgeHarBarOos}{}\renewcommand{\rhcaUnderRidgeHarBarOos}{38.0}
\providecommand{\rhcaMzARidgeHarBarOos}{}\renewcommand{\rhcaMzARidgeHarBarOos}{0.41}
\providecommand{\rhcaMzBRidgeHarBarOos}{}\renewcommand{\rhcaMzBRidgeHarBarOos}{0.99}
\providecommand{\rhcaMzBseRidgeHarBarOos}{}\renewcommand{\rhcaMzBseRidgeHarBarOos}{0.08}
\providecommand{\rhcaMzRsqRidgeHarBarOos}{}\renewcommand{\rhcaMzRsqRidgeHarBarOos}{0.67}
\providecommand{\rhcaMzlARidgeHarBarOos}{}\renewcommand{\rhcaMzlARidgeHarBarOos}{0.15}
\providecommand{\rhcaMzlBRidgeHarBarOos}{}\renewcommand{\rhcaMzlBRidgeHarBarOos}{1.02}
\providecommand{\rhcaMzlBseRidgeHarBarOos}{}\renewcommand{\rhcaMzlBseRidgeHarBarOos}{0.01}
\providecommand{\rhcaMzlRsqRidgeHarBarOos}{}\renewcommand{\rhcaMzlRsqRidgeHarBarOos}{0.79}
\providecommand{\rhcaShBiasRidgeHarBarOos}{}\renewcommand{\rhcaShBiasRidgeHarBarOos}{0.1}
\providecommand{\rhcaShSlopeRidgeHarBarOos}{}\renewcommand{\rhcaShSlopeRidgeHarBarOos}{0.0}
\providecommand{\rhcaShDistRidgeHarBarOos}{}\renewcommand{\rhcaShDistRidgeHarBarOos}{99.9}
\providecommand{\rhcaQlRidgeHarBarOos}{}\renewcommand{\rhcaQlRidgeHarBarOos}{0.0018}
\providecommand{\rhcaQdRidgeHarBarOos}{}\renewcommand{\rhcaQdRidgeHarBarOos}{0.1915}
\providecommand{\rhcaOverSeRidgeHarBarOos}{}\renewcommand{\rhcaOverSeRidgeHarBarOos}{21.5}
\providecommand{\rhcaMrRidgeFullBarOos}{}\renewcommand{\rhcaMrRidgeFullBarOos}{1.117}
\providecommand{\rhcaRmRidgeFullBarOos}{}\renewcommand{\rhcaRmRidgeFullBarOos}{0.977}
\providecommand{\rhcaUnderRidgeFullBarOos}{}\renewcommand{\rhcaUnderRidgeFullBarOos}{42.5}
\providecommand{\rhcaMzARidgeFullBarOos}{}\renewcommand{\rhcaMzARidgeFullBarOos}{3.27}
\providecommand{\rhcaMzBRidgeFullBarOos}{}\renewcommand{\rhcaMzBRidgeFullBarOos}{0.58}
\providecommand{\rhcaMzBseRidgeFullBarOos}{}\renewcommand{\rhcaMzBseRidgeFullBarOos}{0.10}
\providecommand{\rhcaMzRsqRidgeFullBarOos}{}\renewcommand{\rhcaMzRsqRidgeFullBarOos}{0.52}
\providecommand{\rhcaMzlARidgeFullBarOos}{}\renewcommand{\rhcaMzlARidgeFullBarOos}{\ensuremath{-}0.25}
\providecommand{\rhcaMzlBRidgeFullBarOos}{}\renewcommand{\rhcaMzlBRidgeFullBarOos}{0.99}
\providecommand{\rhcaMzlBseRidgeFullBarOos}{}\renewcommand{\rhcaMzlBseRidgeFullBarOos}{0.01}
\providecommand{\rhcaMzlRsqRidgeFullBarOos}{}\renewcommand{\rhcaMzlRsqRidgeFullBarOos}{0.81}
\providecommand{\rhcaShBiasRidgeFullBarOos}{}\renewcommand{\rhcaShBiasRidgeFullBarOos}{0.0}
\providecommand{\rhcaShSlopeRidgeFullBarOos}{}\renewcommand{\rhcaShSlopeRidgeFullBarOos}{35.9}
\providecommand{\rhcaShDistRidgeFullBarOos}{}\renewcommand{\rhcaShDistRidgeFullBarOos}{64.1}
\providecommand{\rhcaQlRidgeFullBarOos}{}\renewcommand{\rhcaQlRidgeFullBarOos}{0.0064}
\providecommand{\rhcaQdRidgeFullBarOos}{}\renewcommand{\rhcaQdRidgeFullBarOos}{0.1788}
\providecommand{\rhcaOverSeRidgeFullBarOos}{}\renewcommand{\rhcaOverSeRidgeFullBarOos}{63.1}
\providecommand{\rhcaMrLassoHarBarOos}{}\renewcommand{\rhcaMrLassoHarBarOos}{1.061}
\providecommand{\rhcaRmLassoHarBarOos}{}\renewcommand{\rhcaRmLassoHarBarOos}{1.045}
\providecommand{\rhcaUnderLassoHarBarOos}{}\renewcommand{\rhcaUnderLassoHarBarOos}{38.0}
\providecommand{\rhcaMzALassoHarBarOos}{}\renewcommand{\rhcaMzALassoHarBarOos}{0.41}
\providecommand{\rhcaMzBLassoHarBarOos}{}\renewcommand{\rhcaMzBLassoHarBarOos}{0.99}
\providecommand{\rhcaMzBseLassoHarBarOos}{}\renewcommand{\rhcaMzBseLassoHarBarOos}{0.08}
\providecommand{\rhcaMzRsqLassoHarBarOos}{}\renewcommand{\rhcaMzRsqLassoHarBarOos}{0.67}
\providecommand{\rhcaMzlALassoHarBarOos}{}\renewcommand{\rhcaMzlALassoHarBarOos}{0.15}
\providecommand{\rhcaMzlBLassoHarBarOos}{}\renewcommand{\rhcaMzlBLassoHarBarOos}{1.02}
\providecommand{\rhcaMzlBseLassoHarBarOos}{}\renewcommand{\rhcaMzlBseLassoHarBarOos}{0.01}
\providecommand{\rhcaMzlRsqLassoHarBarOos}{}\renewcommand{\rhcaMzlRsqLassoHarBarOos}{0.79}
\providecommand{\rhcaShBiasLassoHarBarOos}{}\renewcommand{\rhcaShBiasLassoHarBarOos}{0.1}
\providecommand{\rhcaShSlopeLassoHarBarOos}{}\renewcommand{\rhcaShSlopeLassoHarBarOos}{0.0}
\providecommand{\rhcaShDistLassoHarBarOos}{}\renewcommand{\rhcaShDistLassoHarBarOos}{99.9}
\providecommand{\rhcaQlLassoHarBarOos}{}\renewcommand{\rhcaQlLassoHarBarOos}{0.0018}
\providecommand{\rhcaQdLassoHarBarOos}{}\renewcommand{\rhcaQdLassoHarBarOos}{0.1915}
\providecommand{\rhcaOverSeLassoHarBarOos}{}\renewcommand{\rhcaOverSeLassoHarBarOos}{21.5}
\providecommand{\rhcaMrLassoFullBarOos}{}\renewcommand{\rhcaMrLassoFullBarOos}{1.116}
\providecommand{\rhcaRmLassoFullBarOos}{}\renewcommand{\rhcaRmLassoFullBarOos}{0.978}
\providecommand{\rhcaUnderLassoFullBarOos}{}\renewcommand{\rhcaUnderLassoFullBarOos}{43.2}
\providecommand{\rhcaMzALassoFullBarOos}{}\renewcommand{\rhcaMzALassoFullBarOos}{3.46}
\providecommand{\rhcaMzBLassoFullBarOos}{}\renewcommand{\rhcaMzBLassoFullBarOos}{0.56}
\providecommand{\rhcaMzBseLassoFullBarOos}{}\renewcommand{\rhcaMzBseLassoFullBarOos}{0.10}
\providecommand{\rhcaMzRsqLassoFullBarOos}{}\renewcommand{\rhcaMzRsqLassoFullBarOos}{0.51}
\providecommand{\rhcaMzlALassoFullBarOos}{}\renewcommand{\rhcaMzlALassoFullBarOos}{\ensuremath{-}0.26}
\providecommand{\rhcaMzlBLassoFullBarOos}{}\renewcommand{\rhcaMzlBLassoFullBarOos}{0.98}
\providecommand{\rhcaMzlBseLassoFullBarOos}{}\renewcommand{\rhcaMzlBseLassoFullBarOos}{0.01}
\providecommand{\rhcaMzlRsqLassoFullBarOos}{}\renewcommand{\rhcaMzlRsqLassoFullBarOos}{0.81}
\providecommand{\rhcaShBiasLassoFullBarOos}{}\renewcommand{\rhcaShBiasLassoFullBarOos}{0.0}
\providecommand{\rhcaShSlopeLassoFullBarOos}{}\renewcommand{\rhcaShSlopeLassoFullBarOos}{38.2}
\providecommand{\rhcaShDistLassoFullBarOos}{}\renewcommand{\rhcaShDistLassoFullBarOos}{61.8}
\providecommand{\rhcaQlLassoFullBarOos}{}\renewcommand{\rhcaQlLassoFullBarOos}{0.0062}
\providecommand{\rhcaQdLassoFullBarOos}{}\renewcommand{\rhcaQdLassoFullBarOos}{0.1739}
\providecommand{\rhcaOverSeLassoFullBarOos}{}\renewcommand{\rhcaOverSeLassoFullBarOos}{64.5}
\providecommand{\rhcaMrEnetHarBarOos}{}\renewcommand{\rhcaMrEnetHarBarOos}{1.061}
\providecommand{\rhcaRmEnetHarBarOos}{}\renewcommand{\rhcaRmEnetHarBarOos}{1.045}
\providecommand{\rhcaUnderEnetHarBarOos}{}\renewcommand{\rhcaUnderEnetHarBarOos}{38.0}
\providecommand{\rhcaMzAEnetHarBarOos}{}\renewcommand{\rhcaMzAEnetHarBarOos}{0.41}
\providecommand{\rhcaMzBEnetHarBarOos}{}\renewcommand{\rhcaMzBEnetHarBarOos}{0.99}
\providecommand{\rhcaMzBseEnetHarBarOos}{}\renewcommand{\rhcaMzBseEnetHarBarOos}{0.08}
\providecommand{\rhcaMzRsqEnetHarBarOos}{}\renewcommand{\rhcaMzRsqEnetHarBarOos}{0.67}
\providecommand{\rhcaMzlAEnetHarBarOos}{}\renewcommand{\rhcaMzlAEnetHarBarOos}{0.15}
\providecommand{\rhcaMzlBEnetHarBarOos}{}\renewcommand{\rhcaMzlBEnetHarBarOos}{1.02}
\providecommand{\rhcaMzlBseEnetHarBarOos}{}\renewcommand{\rhcaMzlBseEnetHarBarOos}{0.01}
\providecommand{\rhcaMzlRsqEnetHarBarOos}{}\renewcommand{\rhcaMzlRsqEnetHarBarOos}{0.79}
\providecommand{\rhcaShBiasEnetHarBarOos}{}\renewcommand{\rhcaShBiasEnetHarBarOos}{0.1}
\providecommand{\rhcaShSlopeEnetHarBarOos}{}\renewcommand{\rhcaShSlopeEnetHarBarOos}{0.0}
\providecommand{\rhcaShDistEnetHarBarOos}{}\renewcommand{\rhcaShDistEnetHarBarOos}{99.9}
\providecommand{\rhcaQlEnetHarBarOos}{}\renewcommand{\rhcaQlEnetHarBarOos}{0.0018}
\providecommand{\rhcaQdEnetHarBarOos}{}\renewcommand{\rhcaQdEnetHarBarOos}{0.1915}
\providecommand{\rhcaOverSeEnetHarBarOos}{}\renewcommand{\rhcaOverSeEnetHarBarOos}{21.5}
\providecommand{\rhcaMrEnetFullBarOos}{}\renewcommand{\rhcaMrEnetFullBarOos}{1.109}
\providecommand{\rhcaRmEnetFullBarOos}{}\renewcommand{\rhcaRmEnetFullBarOos}{0.966}
\providecommand{\rhcaUnderEnetFullBarOos}{}\renewcommand{\rhcaUnderEnetFullBarOos}{42.2}
\providecommand{\rhcaMzAEnetFullBarOos}{}\renewcommand{\rhcaMzAEnetFullBarOos}{3.45}
\providecommand{\rhcaMzBEnetFullBarOos}{}\renewcommand{\rhcaMzBEnetFullBarOos}{0.56}
\providecommand{\rhcaMzBseEnetFullBarOos}{}\renewcommand{\rhcaMzBseEnetFullBarOos}{0.10}
\providecommand{\rhcaMzRsqEnetFullBarOos}{}\renewcommand{\rhcaMzRsqEnetFullBarOos}{0.52}
\providecommand{\rhcaMzlAEnetFullBarOos}{}\renewcommand{\rhcaMzlAEnetFullBarOos}{\ensuremath{-}0.22}
\providecommand{\rhcaMzlBEnetFullBarOos}{}\renewcommand{\rhcaMzlBEnetFullBarOos}{0.99}
\providecommand{\rhcaMzlBseEnetFullBarOos}{}\renewcommand{\rhcaMzlBseEnetFullBarOos}{0.01}
\providecommand{\rhcaMzlRsqEnetFullBarOos}{}\renewcommand{\rhcaMzlRsqEnetFullBarOos}{0.81}
\providecommand{\rhcaShBiasEnetFullBarOos}{}\renewcommand{\rhcaShBiasEnetFullBarOos}{0.0}
\providecommand{\rhcaShSlopeEnetFullBarOos}{}\renewcommand{\rhcaShSlopeEnetFullBarOos}{40.8}
\providecommand{\rhcaShDistEnetFullBarOos}{}\renewcommand{\rhcaShDistEnetFullBarOos}{59.1}
\providecommand{\rhcaQlEnetFullBarOos}{}\renewcommand{\rhcaQlEnetFullBarOos}{0.0056}
\providecommand{\rhcaQdEnetFullBarOos}{}\renewcommand{\rhcaQdEnetFullBarOos}{0.1775}
\providecommand{\rhcaOverSeEnetFullBarOos}{}\renewcommand{\rhcaOverSeEnetFullBarOos}{67.0}
\providecommand{\rhcaMrLgbmHarBarOos}{}\renewcommand{\rhcaMrLgbmHarBarOos}{1.067}
\providecommand{\rhcaRmLgbmHarBarOos}{}\renewcommand{\rhcaRmLgbmHarBarOos}{1.028}
\providecommand{\rhcaUnderLgbmHarBarOos}{}\renewcommand{\rhcaUnderLgbmHarBarOos}{38.5}
\providecommand{\rhcaMzALgbmHarBarOos}{}\renewcommand{\rhcaMzALgbmHarBarOos}{0.72}
\providecommand{\rhcaMzBLgbmHarBarOos}{}\renewcommand{\rhcaMzBLgbmHarBarOos}{0.94}
\providecommand{\rhcaMzBseLgbmHarBarOos}{}\renewcommand{\rhcaMzBseLgbmHarBarOos}{0.06}
\providecommand{\rhcaMzRsqLgbmHarBarOos}{}\renewcommand{\rhcaMzRsqLgbmHarBarOos}{0.66}
\providecommand{\rhcaMzlALgbmHarBarOos}{}\renewcommand{\rhcaMzlALgbmHarBarOos}{0.03}
\providecommand{\rhcaMzlBLgbmHarBarOos}{}\renewcommand{\rhcaMzlBLgbmHarBarOos}{1.01}
\providecommand{\rhcaMzlBseLgbmHarBarOos}{}\renewcommand{\rhcaMzlBseLgbmHarBarOos}{0.01}
\providecommand{\rhcaMzlRsqLgbmHarBarOos}{}\renewcommand{\rhcaMzlRsqLgbmHarBarOos}{0.80}
\providecommand{\rhcaShBiasLgbmHarBarOos}{}\renewcommand{\rhcaShBiasLgbmHarBarOos}{0.0}
\providecommand{\rhcaShSlopeLgbmHarBarOos}{}\renewcommand{\rhcaShSlopeLgbmHarBarOos}{0.8}
\providecommand{\rhcaShDistLgbmHarBarOos}{}\renewcommand{\rhcaShDistLgbmHarBarOos}{99.1}
\providecommand{\rhcaQlLgbmHarBarOos}{}\renewcommand{\rhcaQlLgbmHarBarOos}{0.0022}
\providecommand{\rhcaQdLgbmHarBarOos}{}\renewcommand{\rhcaQdLgbmHarBarOos}{0.1878}
\providecommand{\rhcaOverSeLgbmHarBarOos}{}\renewcommand{\rhcaOverSeLgbmHarBarOos}{25.9}
\providecommand{\rhcaMrLgbmFullBarOos}{}\renewcommand{\rhcaMrLgbmFullBarOos}{1.072}
\providecommand{\rhcaRmLgbmFullBarOos}{}\renewcommand{\rhcaRmLgbmFullBarOos}{0.963}
\providecommand{\rhcaUnderLgbmFullBarOos}{}\renewcommand{\rhcaUnderLgbmFullBarOos}{40.6}
\providecommand{\rhcaMzALgbmFullBarOos}{}\renewcommand{\rhcaMzALgbmFullBarOos}{1.86}
\providecommand{\rhcaMzBLgbmFullBarOos}{}\renewcommand{\rhcaMzBLgbmFullBarOos}{0.74}
\providecommand{\rhcaMzBseLgbmFullBarOos}{}\renewcommand{\rhcaMzBseLgbmFullBarOos}{0.07}
\providecommand{\rhcaMzRsqLgbmFullBarOos}{}\renewcommand{\rhcaMzRsqLgbmFullBarOos}{0.64}
\providecommand{\rhcaMzlALgbmFullBarOos}{}\renewcommand{\rhcaMzlALgbmFullBarOos}{\ensuremath{-}0.20}
\providecommand{\rhcaMzlBLgbmFullBarOos}{}\renewcommand{\rhcaMzlBLgbmFullBarOos}{0.99}
\providecommand{\rhcaMzlBseLgbmFullBarOos}{}\renewcommand{\rhcaMzlBseLgbmFullBarOos}{0.01}
\providecommand{\rhcaMzlRsqLgbmFullBarOos}{}\renewcommand{\rhcaMzlRsqLgbmFullBarOos}{0.82}
\providecommand{\rhcaShBiasLgbmFullBarOos}{}\renewcommand{\rhcaShBiasLgbmFullBarOos}{0.0}
\providecommand{\rhcaShSlopeLgbmFullBarOos}{}\renewcommand{\rhcaShSlopeLgbmFullBarOos}{17.7}
\providecommand{\rhcaShDistLgbmFullBarOos}{}\renewcommand{\rhcaShDistLgbmFullBarOos}{82.2}
\providecommand{\rhcaQlLgbmFullBarOos}{}\renewcommand{\rhcaQlLgbmFullBarOos}{0.0025}
\providecommand{\rhcaQdLgbmFullBarOos}{}\renewcommand{\rhcaQdLgbmFullBarOos}{0.1621}
\providecommand{\rhcaOverSeLgbmFullBarOos}{}\renewcommand{\rhcaOverSeLgbmFullBarOos}{49.1}
\providecommand{\rhcaMrXgbHarBarOos}{}\renewcommand{\rhcaMrXgbHarBarOos}{1.069}
\providecommand{\rhcaRmXgbHarBarOos}{}\renewcommand{\rhcaRmXgbHarBarOos}{1.028}
\providecommand{\rhcaUnderXgbHarBarOos}{}\renewcommand{\rhcaUnderXgbHarBarOos}{38.4}
\providecommand{\rhcaMzAXgbHarBarOos}{}\renewcommand{\rhcaMzAXgbHarBarOos}{0.69}
\providecommand{\rhcaMzBXgbHarBarOos}{}\renewcommand{\rhcaMzBXgbHarBarOos}{0.94}
\providecommand{\rhcaMzBseXgbHarBarOos}{}\renewcommand{\rhcaMzBseXgbHarBarOos}{0.06}
\providecommand{\rhcaMzRsqXgbHarBarOos}{}\renewcommand{\rhcaMzRsqXgbHarBarOos}{0.66}
\providecommand{\rhcaMzlAXgbHarBarOos}{}\renewcommand{\rhcaMzlAXgbHarBarOos}{0.04}
\providecommand{\rhcaMzlBXgbHarBarOos}{}\renewcommand{\rhcaMzlBXgbHarBarOos}{1.01}
\providecommand{\rhcaMzlBseXgbHarBarOos}{}\renewcommand{\rhcaMzlBseXgbHarBarOos}{0.01}
\providecommand{\rhcaMzlRsqXgbHarBarOos}{}\renewcommand{\rhcaMzlRsqXgbHarBarOos}{0.80}
\providecommand{\rhcaShBiasXgbHarBarOos}{}\renewcommand{\rhcaShBiasXgbHarBarOos}{0.0}
\providecommand{\rhcaShSlopeXgbHarBarOos}{}\renewcommand{\rhcaShSlopeXgbHarBarOos}{0.8}
\providecommand{\rhcaShDistXgbHarBarOos}{}\renewcommand{\rhcaShDistXgbHarBarOos}{99.2}
\providecommand{\rhcaQlXgbHarBarOos}{}\renewcommand{\rhcaQlXgbHarBarOos}{0.0023}
\providecommand{\rhcaQdXgbHarBarOos}{}\renewcommand{\rhcaQdXgbHarBarOos}{0.1899}
\providecommand{\rhcaOverSeXgbHarBarOos}{}\renewcommand{\rhcaOverSeXgbHarBarOos}{25.9}
\providecommand{\rhcaMrXgbFullBarOos}{}\renewcommand{\rhcaMrXgbFullBarOos}{1.068}
\providecommand{\rhcaRmXgbFullBarOos}{}\renewcommand{\rhcaRmXgbFullBarOos}{0.966}
\providecommand{\rhcaUnderXgbFullBarOos}{}\renewcommand{\rhcaUnderXgbFullBarOos}{40.5}
\providecommand{\rhcaMzAXgbFullBarOos}{}\renewcommand{\rhcaMzAXgbFullBarOos}{1.74}
\providecommand{\rhcaMzBXgbFullBarOos}{}\renewcommand{\rhcaMzBXgbFullBarOos}{0.76}
\providecommand{\rhcaMzBseXgbFullBarOos}{}\renewcommand{\rhcaMzBseXgbFullBarOos}{0.06}
\providecommand{\rhcaMzRsqXgbFullBarOos}{}\renewcommand{\rhcaMzRsqXgbFullBarOos}{0.66}
\providecommand{\rhcaMzlAXgbFullBarOos}{}\renewcommand{\rhcaMzlAXgbFullBarOos}{\ensuremath{-}0.18}
\providecommand{\rhcaMzlBXgbFullBarOos}{}\renewcommand{\rhcaMzlBXgbFullBarOos}{0.99}
\providecommand{\rhcaMzlBseXgbFullBarOos}{}\renewcommand{\rhcaMzlBseXgbFullBarOos}{0.01}
\providecommand{\rhcaMzlRsqXgbFullBarOos}{}\renewcommand{\rhcaMzlRsqXgbFullBarOos}{0.82}
\providecommand{\rhcaShBiasXgbFullBarOos}{}\renewcommand{\rhcaShBiasXgbFullBarOos}{0.0}
\providecommand{\rhcaShSlopeXgbFullBarOos}{}\renewcommand{\rhcaShSlopeXgbFullBarOos}{16.1}
\providecommand{\rhcaShDistXgbFullBarOos}{}\renewcommand{\rhcaShDistXgbFullBarOos}{83.9}
\providecommand{\rhcaQlXgbFullBarOos}{}\renewcommand{\rhcaQlXgbFullBarOos}{0.0022}
\providecommand{\rhcaQdXgbFullBarOos}{}\renewcommand{\rhcaQdXgbFullBarOos}{0.1612}
\providecommand{\rhcaOverSeXgbFullBarOos}{}\renewcommand{\rhcaOverSeXgbFullBarOos}{48.3}
\providecommand{\rhcaMrRfHarBarOos}{}\renewcommand{\rhcaMrRfHarBarOos}{1.077}
\providecommand{\rhcaRmRfHarBarOos}{}\renewcommand{\rhcaRmRfHarBarOos}{1.024}
\providecommand{\rhcaUnderRfHarBarOos}{}\renewcommand{\rhcaUnderRfHarBarOos}{38.8}
\providecommand{\rhcaMzARfHarBarOos}{}\renewcommand{\rhcaMzARfHarBarOos}{0.82}
\providecommand{\rhcaMzBRfHarBarOos}{}\renewcommand{\rhcaMzBRfHarBarOos}{0.92}
\providecommand{\rhcaMzBseRfHarBarOos}{}\renewcommand{\rhcaMzBseRfHarBarOos}{0.07}
\providecommand{\rhcaMzRsqRfHarBarOos}{}\renewcommand{\rhcaMzRsqRfHarBarOos}{0.65}
\providecommand{\rhcaMzlARfHarBarOos}{}\renewcommand{\rhcaMzlARfHarBarOos}{\ensuremath{-}0.07}
\providecommand{\rhcaMzlBRfHarBarOos}{}\renewcommand{\rhcaMzlBRfHarBarOos}{1.00}
\providecommand{\rhcaMzlBseRfHarBarOos}{}\renewcommand{\rhcaMzlBseRfHarBarOos}{0.01}
\providecommand{\rhcaMzlRsqRfHarBarOos}{}\renewcommand{\rhcaMzlRsqRfHarBarOos}{0.80}
\providecommand{\rhcaShBiasRfHarBarOos}{}\renewcommand{\rhcaShBiasRfHarBarOos}{0.0}
\providecommand{\rhcaShSlopeRfHarBarOos}{}\renewcommand{\rhcaShSlopeRfHarBarOos}{1.4}
\providecommand{\rhcaShDistRfHarBarOos}{}\renewcommand{\rhcaShDistRfHarBarOos}{98.6}
\providecommand{\rhcaQlRfHarBarOos}{}\renewcommand{\rhcaQlRfHarBarOos}{0.0028}
\providecommand{\rhcaQdRfHarBarOos}{}\renewcommand{\rhcaQdRfHarBarOos}{0.1909}
\providecommand{\rhcaOverSeRfHarBarOos}{}\renewcommand{\rhcaOverSeRfHarBarOos}{26.7}
\providecommand{\rhcaMrRfFullBarOos}{}\renewcommand{\rhcaMrRfFullBarOos}{1.087}
\providecommand{\rhcaRmRfFullBarOos}{}\renewcommand{\rhcaRmRfFullBarOos}{0.974}
\providecommand{\rhcaUnderRfFullBarOos}{}\renewcommand{\rhcaUnderRfFullBarOos}{41.0}
\providecommand{\rhcaMzARfFullBarOos}{}\renewcommand{\rhcaMzARfFullBarOos}{1.61}
\providecommand{\rhcaMzBRfFullBarOos}{}\renewcommand{\rhcaMzBRfFullBarOos}{0.78}
\providecommand{\rhcaMzBseRfFullBarOos}{}\renewcommand{\rhcaMzBseRfFullBarOos}{0.07}
\providecommand{\rhcaMzRsqRfFullBarOos}{}\renewcommand{\rhcaMzRsqRfFullBarOos}{0.65}
\providecommand{\rhcaMzlARfFullBarOos}{}\renewcommand{\rhcaMzlARfFullBarOos}{\ensuremath{-}0.25}
\providecommand{\rhcaMzlBRfFullBarOos}{}\renewcommand{\rhcaMzlBRfFullBarOos}{0.99}
\providecommand{\rhcaMzlBseRfFullBarOos}{}\renewcommand{\rhcaMzlBseRfFullBarOos}{0.01}
\providecommand{\rhcaMzlRsqRfFullBarOos}{}\renewcommand{\rhcaMzlRsqRfFullBarOos}{0.82}
\providecommand{\rhcaShBiasRfFullBarOos}{}\renewcommand{\rhcaShBiasRfFullBarOos}{0.0}
\providecommand{\rhcaShSlopeRfFullBarOos}{}\renewcommand{\rhcaShSlopeRfFullBarOos}{12.7}
\providecommand{\rhcaShDistRfFullBarOos}{}\renewcommand{\rhcaShDistRfFullBarOos}{87.3}
\providecommand{\rhcaQlRfFullBarOos}{}\renewcommand{\rhcaQlRfFullBarOos}{0.0035}
\providecommand{\rhcaQdRfFullBarOos}{}\renewcommand{\rhcaQdRfFullBarOos}{0.1694}
\providecommand{\rhcaOverSeRfFullBarOos}{}\renewcommand{\rhcaOverSeRfFullBarOos}{44.4}
\providecommand{\rhcaMrRidgeHarCloseOos}{}\renewcommand{\rhcaMrRidgeHarCloseOos}{1.092}
\providecommand{\rhcaRmRidgeHarCloseOos}{}\renewcommand{\rhcaRmRidgeHarCloseOos}{1.103}
\providecommand{\rhcaUnderRidgeHarCloseOos}{}\renewcommand{\rhcaUnderRidgeHarCloseOos}{47.7}
\providecommand{\rhcaMzARidgeHarCloseOos}{}\renewcommand{\rhcaMzARidgeHarCloseOos}{\ensuremath{-}1.23}
\providecommand{\rhcaMzBRidgeHarCloseOos}{}\renewcommand{\rhcaMzBRidgeHarCloseOos}{1.21}
\providecommand{\rhcaMzBseRidgeHarCloseOos}{}\renewcommand{\rhcaMzBseRidgeHarCloseOos}{0.09}
\providecommand{\rhcaMzRsqRidgeHarCloseOos}{}\renewcommand{\rhcaMzRsqRidgeHarCloseOos}{0.80}
\providecommand{\rhcaMzlARidgeHarCloseOos}{}\renewcommand{\rhcaMzlARidgeHarCloseOos}{\ensuremath{-}0.38}
\providecommand{\rhcaMzlBRidgeHarCloseOos}{}\renewcommand{\rhcaMzlBRidgeHarCloseOos}{0.97}
\providecommand{\rhcaMzlBseRidgeHarCloseOos}{}\renewcommand{\rhcaMzlBseRidgeHarCloseOos}{0.02}
\providecommand{\rhcaMzlRsqRidgeHarCloseOos}{}\renewcommand{\rhcaMzlRsqRidgeHarCloseOos}{0.84}
\providecommand{\rhcaShBiasRidgeHarCloseOos}{}\renewcommand{\rhcaShBiasRidgeHarCloseOos}{0.3}
\providecommand{\rhcaShSlopeRidgeHarCloseOos}{}\renewcommand{\rhcaShSlopeRidgeHarCloseOos}{10.7}
\providecommand{\rhcaShDistRidgeHarCloseOos}{}\renewcommand{\rhcaShDistRidgeHarCloseOos}{88.9}
\providecommand{\rhcaQlRidgeHarCloseOos}{}\renewcommand{\rhcaQlRidgeHarCloseOos}{0.0040}
\providecommand{\rhcaQdRidgeHarCloseOos}{}\renewcommand{\rhcaQdRidgeHarCloseOos}{0.1097}
\providecommand{\rhcaOverSeRidgeHarCloseOos}{}\renewcommand{\rhcaOverSeRidgeHarCloseOos}{13.8}
\providecommand{\rhcaMrRidgeFullCloseOos}{}\renewcommand{\rhcaMrRidgeFullCloseOos}{1.169}
\providecommand{\rhcaRmRidgeFullCloseOos}{}\renewcommand{\rhcaRmRidgeFullCloseOos}{0.971}
\providecommand{\rhcaUnderRidgeFullCloseOos}{}\renewcommand{\rhcaUnderRidgeFullCloseOos}{52.1}
\providecommand{\rhcaMzARidgeFullCloseOos}{}\renewcommand{\rhcaMzARidgeFullCloseOos}{3.60}
\providecommand{\rhcaMzBRidgeFullCloseOos}{}\renewcommand{\rhcaMzBRidgeFullCloseOos}{0.70}
\providecommand{\rhcaMzBseRidgeFullCloseOos}{}\renewcommand{\rhcaMzBseRidgeFullCloseOos}{0.06}
\providecommand{\rhcaMzRsqRidgeFullCloseOos}{}\renewcommand{\rhcaMzRsqRidgeFullCloseOos}{0.72}
\providecommand{\rhcaMzlARidgeFullCloseOos}{}\renewcommand{\rhcaMzlARidgeFullCloseOos}{\ensuremath{-}0.95}
\providecommand{\rhcaMzlBRidgeFullCloseOos}{}\renewcommand{\rhcaMzlBRidgeFullCloseOos}{0.92}
\providecommand{\rhcaMzlBseRidgeFullCloseOos}{}\renewcommand{\rhcaMzlBseRidgeFullCloseOos}{0.01}
\providecommand{\rhcaMzlRsqRidgeFullCloseOos}{}\renewcommand{\rhcaMzlRsqRidgeFullCloseOos}{0.84}
\providecommand{\rhcaShBiasRidgeFullCloseOos}{}\renewcommand{\rhcaShBiasRidgeFullCloseOos}{0.0}
\providecommand{\rhcaShSlopeRidgeFullCloseOos}{}\renewcommand{\rhcaShSlopeRidgeFullCloseOos}{32.0}
\providecommand{\rhcaShDistRidgeFullCloseOos}{}\renewcommand{\rhcaShDistRidgeFullCloseOos}{68.0}
\providecommand{\rhcaQlRidgeFullCloseOos}{}\renewcommand{\rhcaQlRidgeFullCloseOos}{0.0128}
\providecommand{\rhcaQdRidgeFullCloseOos}{}\renewcommand{\rhcaQdRidgeFullCloseOos}{0.1140}
\providecommand{\rhcaOverSeRidgeFullCloseOos}{}\renewcommand{\rhcaOverSeRidgeFullCloseOos}{52.6}
\providecommand{\rhcaMrLassoHarCloseOos}{}\renewcommand{\rhcaMrLassoHarCloseOos}{1.092}
\providecommand{\rhcaRmLassoHarCloseOos}{}\renewcommand{\rhcaRmLassoHarCloseOos}{1.103}
\providecommand{\rhcaUnderLassoHarCloseOos}{}\renewcommand{\rhcaUnderLassoHarCloseOos}{47.7}
\providecommand{\rhcaMzALassoHarCloseOos}{}\renewcommand{\rhcaMzALassoHarCloseOos}{\ensuremath{-}1.23}
\providecommand{\rhcaMzBLassoHarCloseOos}{}\renewcommand{\rhcaMzBLassoHarCloseOos}{1.21}
\providecommand{\rhcaMzBseLassoHarCloseOos}{}\renewcommand{\rhcaMzBseLassoHarCloseOos}{0.09}
\providecommand{\rhcaMzRsqLassoHarCloseOos}{}\renewcommand{\rhcaMzRsqLassoHarCloseOos}{0.80}
\providecommand{\rhcaMzlALassoHarCloseOos}{}\renewcommand{\rhcaMzlALassoHarCloseOos}{\ensuremath{-}0.38}
\providecommand{\rhcaMzlBLassoHarCloseOos}{}\renewcommand{\rhcaMzlBLassoHarCloseOos}{0.97}
\providecommand{\rhcaMzlBseLassoHarCloseOos}{}\renewcommand{\rhcaMzlBseLassoHarCloseOos}{0.02}
\providecommand{\rhcaMzlRsqLassoHarCloseOos}{}\renewcommand{\rhcaMzlRsqLassoHarCloseOos}{0.84}
\providecommand{\rhcaShBiasLassoHarCloseOos}{}\renewcommand{\rhcaShBiasLassoHarCloseOos}{0.3}
\providecommand{\rhcaShSlopeLassoHarCloseOos}{}\renewcommand{\rhcaShSlopeLassoHarCloseOos}{10.7}
\providecommand{\rhcaShDistLassoHarCloseOos}{}\renewcommand{\rhcaShDistLassoHarCloseOos}{88.9}
\providecommand{\rhcaQlLassoHarCloseOos}{}\renewcommand{\rhcaQlLassoHarCloseOos}{0.0040}
\providecommand{\rhcaQdLassoHarCloseOos}{}\renewcommand{\rhcaQdLassoHarCloseOos}{0.1097}
\providecommand{\rhcaOverSeLassoHarCloseOos}{}\renewcommand{\rhcaOverSeLassoHarCloseOos}{13.8}
\providecommand{\rhcaMrLassoFullCloseOos}{}\renewcommand{\rhcaMrLassoFullCloseOos}{1.183}
\providecommand{\rhcaRmLassoFullCloseOos}{}\renewcommand{\rhcaRmLassoFullCloseOos}{0.982}
\providecommand{\rhcaUnderLassoFullCloseOos}{}\renewcommand{\rhcaUnderLassoFullCloseOos}{52.6}
\providecommand{\rhcaMzALassoFullCloseOos}{}\renewcommand{\rhcaMzALassoFullCloseOos}{3.74}
\providecommand{\rhcaMzBLassoFullCloseOos}{}\renewcommand{\rhcaMzBLassoFullCloseOos}{0.69}
\providecommand{\rhcaMzBseLassoFullCloseOos}{}\renewcommand{\rhcaMzBseLassoFullCloseOos}{0.05}
\providecommand{\rhcaMzRsqLassoFullCloseOos}{}\renewcommand{\rhcaMzRsqLassoFullCloseOos}{0.71}
\providecommand{\rhcaMzlALassoFullCloseOos}{}\renewcommand{\rhcaMzlALassoFullCloseOos}{\ensuremath{-}0.93}
\providecommand{\rhcaMzlBLassoFullCloseOos}{}\renewcommand{\rhcaMzlBLassoFullCloseOos}{0.92}
\providecommand{\rhcaMzlBseLassoFullCloseOos}{}\renewcommand{\rhcaMzlBseLassoFullCloseOos}{0.01}
\providecommand{\rhcaMzlRsqLassoFullCloseOos}{}\renewcommand{\rhcaMzlRsqLassoFullCloseOos}{0.84}
\providecommand{\rhcaShBiasLassoFullCloseOos}{}\renewcommand{\rhcaShBiasLassoFullCloseOos}{0.0}
\providecommand{\rhcaShSlopeLassoFullCloseOos}{}\renewcommand{\rhcaShSlopeLassoFullCloseOos}{32.1}
\providecommand{\rhcaShDistLassoFullCloseOos}{}\renewcommand{\rhcaShDistLassoFullCloseOos}{67.9}
\providecommand{\rhcaQlLassoFullCloseOos}{}\renewcommand{\rhcaQlLassoFullCloseOos}{0.0150}
\providecommand{\rhcaQdLassoFullCloseOos}{}\renewcommand{\rhcaQdLassoFullCloseOos}{0.1159}
\providecommand{\rhcaOverSeLassoFullCloseOos}{}\renewcommand{\rhcaOverSeLassoFullCloseOos}{51.5}
\providecommand{\rhcaMrEnetHarCloseOos}{}\renewcommand{\rhcaMrEnetHarCloseOos}{1.092}
\providecommand{\rhcaRmEnetHarCloseOos}{}\renewcommand{\rhcaRmEnetHarCloseOos}{1.103}
\providecommand{\rhcaUnderEnetHarCloseOos}{}\renewcommand{\rhcaUnderEnetHarCloseOos}{47.7}
\providecommand{\rhcaMzAEnetHarCloseOos}{}\renewcommand{\rhcaMzAEnetHarCloseOos}{\ensuremath{-}1.23}
\providecommand{\rhcaMzBEnetHarCloseOos}{}\renewcommand{\rhcaMzBEnetHarCloseOos}{1.21}
\providecommand{\rhcaMzBseEnetHarCloseOos}{}\renewcommand{\rhcaMzBseEnetHarCloseOos}{0.09}
\providecommand{\rhcaMzRsqEnetHarCloseOos}{}\renewcommand{\rhcaMzRsqEnetHarCloseOos}{0.80}
\providecommand{\rhcaMzlAEnetHarCloseOos}{}\renewcommand{\rhcaMzlAEnetHarCloseOos}{\ensuremath{-}0.38}
\providecommand{\rhcaMzlBEnetHarCloseOos}{}\renewcommand{\rhcaMzlBEnetHarCloseOos}{0.97}
\providecommand{\rhcaMzlBseEnetHarCloseOos}{}\renewcommand{\rhcaMzlBseEnetHarCloseOos}{0.02}
\providecommand{\rhcaMzlRsqEnetHarCloseOos}{}\renewcommand{\rhcaMzlRsqEnetHarCloseOos}{0.84}
\providecommand{\rhcaShBiasEnetHarCloseOos}{}\renewcommand{\rhcaShBiasEnetHarCloseOos}{0.3}
\providecommand{\rhcaShSlopeEnetHarCloseOos}{}\renewcommand{\rhcaShSlopeEnetHarCloseOos}{10.7}
\providecommand{\rhcaShDistEnetHarCloseOos}{}\renewcommand{\rhcaShDistEnetHarCloseOos}{88.9}
\providecommand{\rhcaQlEnetHarCloseOos}{}\renewcommand{\rhcaQlEnetHarCloseOos}{0.0040}
\providecommand{\rhcaQdEnetHarCloseOos}{}\renewcommand{\rhcaQdEnetHarCloseOos}{0.1097}
\providecommand{\rhcaOverSeEnetHarCloseOos}{}\renewcommand{\rhcaOverSeEnetHarCloseOos}{13.8}
\providecommand{\rhcaMrEnetFullCloseOos}{}\renewcommand{\rhcaMrEnetFullCloseOos}{1.165}
\providecommand{\rhcaRmEnetFullCloseOos}{}\renewcommand{\rhcaRmEnetFullCloseOos}{0.955}
\providecommand{\rhcaUnderEnetFullCloseOos}{}\renewcommand{\rhcaUnderEnetFullCloseOos}{52.9}
\providecommand{\rhcaMzAEnetFullCloseOos}{}\renewcommand{\rhcaMzAEnetFullCloseOos}{3.93}
\providecommand{\rhcaMzBEnetFullCloseOos}{}\renewcommand{\rhcaMzBEnetFullCloseOos}{0.66}
\providecommand{\rhcaMzBseEnetFullCloseOos}{}\renewcommand{\rhcaMzBseEnetFullCloseOos}{0.05}
\providecommand{\rhcaMzRsqEnetFullCloseOos}{}\renewcommand{\rhcaMzRsqEnetFullCloseOos}{0.72}
\providecommand{\rhcaMzlAEnetFullCloseOos}{}\renewcommand{\rhcaMzlAEnetFullCloseOos}{\ensuremath{-}0.94}
\providecommand{\rhcaMzlBEnetFullCloseOos}{}\renewcommand{\rhcaMzlBEnetFullCloseOos}{0.92}
\providecommand{\rhcaMzlBseEnetFullCloseOos}{}\renewcommand{\rhcaMzlBseEnetFullCloseOos}{0.01}
\providecommand{\rhcaMzlRsqEnetFullCloseOos}{}\renewcommand{\rhcaMzlRsqEnetFullCloseOos}{0.84}
\providecommand{\rhcaShBiasEnetFullCloseOos}{}\renewcommand{\rhcaShBiasEnetFullCloseOos}{0.0}
\providecommand{\rhcaShSlopeEnetFullCloseOos}{}\renewcommand{\rhcaShSlopeEnetFullCloseOos}{39.7}
\providecommand{\rhcaShDistEnetFullCloseOos}{}\renewcommand{\rhcaShDistEnetFullCloseOos}{60.3}
\providecommand{\rhcaQlEnetFullCloseOos}{}\renewcommand{\rhcaQlEnetFullCloseOos}{0.0123}
\providecommand{\rhcaQdEnetFullCloseOos}{}\renewcommand{\rhcaQdEnetFullCloseOos}{0.1122}
\providecommand{\rhcaOverSeEnetFullCloseOos}{}\renewcommand{\rhcaOverSeEnetFullCloseOos}{59.7}
\providecommand{\rhcaMrLgbmHarCloseOos}{}\renewcommand{\rhcaMrLgbmHarCloseOos}{1.035}
\providecommand{\rhcaRmLgbmHarCloseOos}{}\renewcommand{\rhcaRmLgbmHarCloseOos}{1.028}
\providecommand{\rhcaUnderLgbmHarCloseOos}{}\renewcommand{\rhcaUnderLgbmHarCloseOos}{42.9}
\providecommand{\rhcaMzALgbmHarCloseOos}{}\renewcommand{\rhcaMzALgbmHarCloseOos}{\ensuremath{-}0.68}
\providecommand{\rhcaMzBLgbmHarCloseOos}{}\renewcommand{\rhcaMzBLgbmHarCloseOos}{1.08}
\providecommand{\rhcaMzBseLgbmHarCloseOos}{}\renewcommand{\rhcaMzBseLgbmHarCloseOos}{0.06}
\providecommand{\rhcaMzRsqLgbmHarCloseOos}{}\renewcommand{\rhcaMzRsqLgbmHarCloseOos}{0.81}
\providecommand{\rhcaMzlALgbmHarCloseOos}{}\renewcommand{\rhcaMzlALgbmHarCloseOos}{\ensuremath{-}0.33}
\providecommand{\rhcaMzlBLgbmHarCloseOos}{}\renewcommand{\rhcaMzlBLgbmHarCloseOos}{0.98}
\providecommand{\rhcaMzlBseLgbmHarCloseOos}{}\renewcommand{\rhcaMzlBseLgbmHarCloseOos}{0.02}
\providecommand{\rhcaMzlRsqLgbmHarCloseOos}{}\renewcommand{\rhcaMzlRsqLgbmHarCloseOos}{0.84}
\providecommand{\rhcaShBiasLgbmHarCloseOos}{}\renewcommand{\rhcaShBiasLgbmHarCloseOos}{0.0}
\providecommand{\rhcaShSlopeLgbmHarCloseOos}{}\renewcommand{\rhcaShSlopeLgbmHarCloseOos}{2.4}
\providecommand{\rhcaShDistLgbmHarCloseOos}{}\renewcommand{\rhcaShDistLgbmHarCloseOos}{97.6}
\providecommand{\rhcaQlLgbmHarCloseOos}{}\renewcommand{\rhcaQlLgbmHarCloseOos}{0.0006}
\providecommand{\rhcaQdLgbmHarCloseOos}{}\renewcommand{\rhcaQdLgbmHarCloseOos}{0.1095}
\providecommand{\rhcaOverSeLgbmHarCloseOos}{}\renewcommand{\rhcaOverSeLgbmHarCloseOos}{22.8}
\providecommand{\rhcaMrLgbmFullCloseOos}{}\renewcommand{\rhcaMrLgbmFullCloseOos}{1.093}
\providecommand{\rhcaRmLgbmFullCloseOos}{}\renewcommand{\rhcaRmLgbmFullCloseOos}{0.969}
\providecommand{\rhcaUnderLgbmFullCloseOos}{}\renewcommand{\rhcaUnderLgbmFullCloseOos}{47.1}
\providecommand{\rhcaMzALgbmFullCloseOos}{}\renewcommand{\rhcaMzALgbmFullCloseOos}{2.03}
\providecommand{\rhcaMzBLgbmFullCloseOos}{}\renewcommand{\rhcaMzBLgbmFullCloseOos}{0.81}
\providecommand{\rhcaMzBseLgbmFullCloseOos}{}\renewcommand{\rhcaMzBseLgbmFullCloseOos}{0.05}
\providecommand{\rhcaMzRsqLgbmFullCloseOos}{}\renewcommand{\rhcaMzRsqLgbmFullCloseOos}{0.77}
\providecommand{\rhcaMzlALgbmFullCloseOos}{}\renewcommand{\rhcaMzlALgbmFullCloseOos}{\ensuremath{-}0.69}
\providecommand{\rhcaMzlBLgbmFullCloseOos}{}\renewcommand{\rhcaMzlBLgbmFullCloseOos}{0.95}
\providecommand{\rhcaMzlBseLgbmFullCloseOos}{}\renewcommand{\rhcaMzlBseLgbmFullCloseOos}{0.01}
\providecommand{\rhcaMzlRsqLgbmFullCloseOos}{}\renewcommand{\rhcaMzlRsqLgbmFullCloseOos}{0.85}
\providecommand{\rhcaShBiasLgbmFullCloseOos}{}\renewcommand{\rhcaShBiasLgbmFullCloseOos}{0.0}
\providecommand{\rhcaShSlopeLgbmFullCloseOos}{}\renewcommand{\rhcaShSlopeLgbmFullCloseOos}{14.8}
\providecommand{\rhcaShDistLgbmFullCloseOos}{}\renewcommand{\rhcaShDistLgbmFullCloseOos}{85.2}
\providecommand{\rhcaQlLgbmFullCloseOos}{}\renewcommand{\rhcaQlLgbmFullCloseOos}{0.0041}
\providecommand{\rhcaQdLgbmFullCloseOos}{}\renewcommand{\rhcaQdLgbmFullCloseOos}{0.1045}
\providecommand{\rhcaOverSeLgbmFullCloseOos}{}\renewcommand{\rhcaOverSeLgbmFullCloseOos}{44.8}
\providecommand{\rhcaMrXgbHarCloseOos}{}\renewcommand{\rhcaMrXgbHarCloseOos}{1.033}
\providecommand{\rhcaRmXgbHarCloseOos}{}\renewcommand{\rhcaRmXgbHarCloseOos}{1.029}
\providecommand{\rhcaUnderXgbHarCloseOos}{}\renewcommand{\rhcaUnderXgbHarCloseOos}{43.2}
\providecommand{\rhcaMzAXgbHarCloseOos}{}\renewcommand{\rhcaMzAXgbHarCloseOos}{\ensuremath{-}0.70}
\providecommand{\rhcaMzBXgbHarCloseOos}{}\renewcommand{\rhcaMzBXgbHarCloseOos}{1.08}
\providecommand{\rhcaMzBseXgbHarCloseOos}{}\renewcommand{\rhcaMzBseXgbHarCloseOos}{0.06}
\providecommand{\rhcaMzRsqXgbHarCloseOos}{}\renewcommand{\rhcaMzRsqXgbHarCloseOos}{0.81}
\providecommand{\rhcaMzlAXgbHarCloseOos}{}\renewcommand{\rhcaMzlAXgbHarCloseOos}{\ensuremath{-}0.30}
\providecommand{\rhcaMzlBXgbHarCloseOos}{}\renewcommand{\rhcaMzlBXgbHarCloseOos}{0.98}
\providecommand{\rhcaMzlBseXgbHarCloseOos}{}\renewcommand{\rhcaMzlBseXgbHarCloseOos}{0.02}
\providecommand{\rhcaMzlRsqXgbHarCloseOos}{}\renewcommand{\rhcaMzlRsqXgbHarCloseOos}{0.84}
\providecommand{\rhcaShBiasXgbHarCloseOos}{}\renewcommand{\rhcaShBiasXgbHarCloseOos}{0.0}
\providecommand{\rhcaShSlopeXgbHarCloseOos}{}\renewcommand{\rhcaShSlopeXgbHarCloseOos}{2.6}
\providecommand{\rhcaShDistXgbHarCloseOos}{}\renewcommand{\rhcaShDistXgbHarCloseOos}{97.4}
\providecommand{\rhcaQlXgbHarCloseOos}{}\renewcommand{\rhcaQlXgbHarCloseOos}{0.0005}
\providecommand{\rhcaQdXgbHarCloseOos}{}\renewcommand{\rhcaQdXgbHarCloseOos}{0.1096}
\providecommand{\rhcaOverSeXgbHarCloseOos}{}\renewcommand{\rhcaOverSeXgbHarCloseOos}{22.7}
\providecommand{\rhcaMrXgbFullCloseOos}{}\renewcommand{\rhcaMrXgbFullCloseOos}{1.091}
\providecommand{\rhcaRmXgbFullCloseOos}{}\renewcommand{\rhcaRmXgbFullCloseOos}{0.965}
\providecommand{\rhcaUnderXgbFullCloseOos}{}\renewcommand{\rhcaUnderXgbFullCloseOos}{47.9}
\providecommand{\rhcaMzAXgbFullCloseOos}{}\renewcommand{\rhcaMzAXgbFullCloseOos}{1.80}
\providecommand{\rhcaMzBXgbFullCloseOos}{}\renewcommand{\rhcaMzBXgbFullCloseOos}{0.83}
\providecommand{\rhcaMzBseXgbFullCloseOos}{}\renewcommand{\rhcaMzBseXgbFullCloseOos}{0.04}
\providecommand{\rhcaMzRsqXgbFullCloseOos}{}\renewcommand{\rhcaMzRsqXgbFullCloseOos}{0.78}
\providecommand{\rhcaMzlAXgbFullCloseOos}{}\renewcommand{\rhcaMzlAXgbFullCloseOos}{\ensuremath{-}0.68}
\providecommand{\rhcaMzlBXgbFullCloseOos}{}\renewcommand{\rhcaMzlBXgbFullCloseOos}{0.95}
\providecommand{\rhcaMzlBseXgbFullCloseOos}{}\renewcommand{\rhcaMzlBseXgbFullCloseOos}{0.01}
\providecommand{\rhcaMzlRsqXgbFullCloseOos}{}\renewcommand{\rhcaMzlRsqXgbFullCloseOos}{0.85}
\providecommand{\rhcaShBiasXgbFullCloseOos}{}\renewcommand{\rhcaShBiasXgbFullCloseOos}{0.0}
\providecommand{\rhcaShSlopeXgbFullCloseOos}{}\renewcommand{\rhcaShSlopeXgbFullCloseOos}{13.1}
\providecommand{\rhcaShDistXgbFullCloseOos}{}\renewcommand{\rhcaShDistXgbFullCloseOos}{86.8}
\providecommand{\rhcaQlXgbFullCloseOos}{}\renewcommand{\rhcaQlXgbFullCloseOos}{0.0039}
\providecommand{\rhcaQdXgbFullCloseOos}{}\renewcommand{\rhcaQdXgbFullCloseOos}{0.1044}
\providecommand{\rhcaOverSeXgbFullCloseOos}{}\renewcommand{\rhcaOverSeXgbFullCloseOos}{46.1}
\providecommand{\rhcaMrRfHarCloseOos}{}\renewcommand{\rhcaMrRfHarCloseOos}{1.063}
\providecommand{\rhcaRmRfHarCloseOos}{}\renewcommand{\rhcaRmRfHarCloseOos}{1.033}
\providecommand{\rhcaUnderRfHarCloseOos}{}\renewcommand{\rhcaUnderRfHarCloseOos}{45.5}
\providecommand{\rhcaMzARfHarCloseOos}{}\renewcommand{\rhcaMzARfHarCloseOos}{\ensuremath{-}0.38}
\providecommand{\rhcaMzBRfHarCloseOos}{}\renewcommand{\rhcaMzBRfHarCloseOos}{1.06}
\providecommand{\rhcaMzBseRfHarCloseOos}{}\renewcommand{\rhcaMzBseRfHarCloseOos}{0.07}
\providecommand{\rhcaMzRsqRfHarCloseOos}{}\renewcommand{\rhcaMzRsqRfHarCloseOos}{0.82}
\providecommand{\rhcaMzlARfHarCloseOos}{}\renewcommand{\rhcaMzlARfHarCloseOos}{\ensuremath{-}0.46}
\providecommand{\rhcaMzlBRfHarCloseOos}{}\renewcommand{\rhcaMzlBRfHarCloseOos}{0.97}
\providecommand{\rhcaMzlBseRfHarCloseOos}{}\renewcommand{\rhcaMzlBseRfHarCloseOos}{0.02}
\providecommand{\rhcaMzlRsqRfHarCloseOos}{}\renewcommand{\rhcaMzlRsqRfHarCloseOos}{0.83}
\providecommand{\rhcaShBiasRfHarCloseOos}{}\renewcommand{\rhcaShBiasRfHarCloseOos}{0.1}
\providecommand{\rhcaShSlopeRfHarCloseOos}{}\renewcommand{\rhcaShSlopeRfHarCloseOos}{1.7}
\providecommand{\rhcaShDistRfHarCloseOos}{}\renewcommand{\rhcaShDistRfHarCloseOos}{98.3}
\providecommand{\rhcaQlRfHarCloseOos}{}\renewcommand{\rhcaQlRfHarCloseOos}{0.0019}
\providecommand{\rhcaQdRfHarCloseOos}{}\renewcommand{\rhcaQdRfHarCloseOos}{0.1112}
\providecommand{\rhcaOverSeRfHarCloseOos}{}\renewcommand{\rhcaOverSeRfHarCloseOos}{25.7}
\providecommand{\rhcaMrRfFullCloseOos}{}\renewcommand{\rhcaMrRfFullCloseOos}{1.113}
\providecommand{\rhcaRmRfFullCloseOos}{}\renewcommand{\rhcaRmRfFullCloseOos}{0.978}
\providecommand{\rhcaUnderRfFullCloseOos}{}\renewcommand{\rhcaUnderRfFullCloseOos}{50.1}
\providecommand{\rhcaMzARfFullCloseOos}{}\renewcommand{\rhcaMzARfFullCloseOos}{1.15}
\providecommand{\rhcaMzBRfFullCloseOos}{}\renewcommand{\rhcaMzBRfFullCloseOos}{0.89}
\providecommand{\rhcaMzBseRfFullCloseOos}{}\renewcommand{\rhcaMzBseRfFullCloseOos}{0.06}
\providecommand{\rhcaMzRsqRfFullCloseOos}{}\renewcommand{\rhcaMzRsqRfFullCloseOos}{0.79}
\providecommand{\rhcaMzlARfFullCloseOos}{}\renewcommand{\rhcaMzlARfFullCloseOos}{\ensuremath{-}0.87}
\providecommand{\rhcaMzlBRfFullCloseOos}{}\renewcommand{\rhcaMzlBRfFullCloseOos}{0.93}
\providecommand{\rhcaMzlBseRfFullCloseOos}{}\renewcommand{\rhcaMzlBseRfFullCloseOos}{0.01}
\providecommand{\rhcaMzlRsqRfFullCloseOos}{}\renewcommand{\rhcaMzlRsqRfFullCloseOos}{0.85}
\providecommand{\rhcaShBiasRfFullCloseOos}{}\renewcommand{\rhcaShBiasRfFullCloseOos}{0.0}
\providecommand{\rhcaShSlopeRfFullCloseOos}{}\renewcommand{\rhcaShSlopeRfFullCloseOos}{5.5}
\providecommand{\rhcaShDistRfFullCloseOos}{}\renewcommand{\rhcaShDistRfFullCloseOos}{94.5}
\providecommand{\rhcaQlRfFullCloseOos}{}\renewcommand{\rhcaQlRfFullCloseOos}{0.0059}
\providecommand{\rhcaQdRfFullCloseOos}{}\renewcommand{\rhcaQdRfFullCloseOos}{0.1064}
\providecommand{\rhcaOverSeRfFullCloseOos}{}\renewcommand{\rhcaOverSeRfFullCloseOos}{43.1}
% --- aggregates for the prose
\providecommand{\rhcaNLowerCloseSpan}{}\renewcommand{\rhcaNLowerCloseSpan}{3}  % classes (of 6) with the lower 16:00 QLIKE on the full set, Span
\providecommand{\rhcaNLowerSigCloseSpan}{}\renewcommand{\rhcaNLowerSigCloseSpan}{1}  % ... and DM p < 0.05
\providecommand{\rhcaNHigherCloseSpan}{}\renewcommand{\rhcaNHigherCloseSpan}{3}  % classes with the higher 16:00 QLIKE on the full set
\providecommand{\rhcaNHigherSigCloseSpan}{}\renewcommand{\rhcaNHigherSigCloseSpan}{0}  % ... and DM p < 0.05
\providecommand{\rhcaNLowerCloseWordSpan}{}\renewcommand{\rhcaNLowerCloseWordSpan}{three}  % as a word
\providecommand{\rhcaNLinHigherSigCloseSpan}{}\renewcommand{\rhcaNLinHigherSigCloseSpan}{0}  % linear classes (of 3): full-set 16:00 QLIKE significantly higher
\providecommand{\rhcaNLinLowerCloseSpan}{}\renewcommand{\rhcaNLinLowerCloseSpan}{0}  % linear classes with a lower 16:00 QLIKE
\providecommand{\rhcaNTreeLowerSigCloseSpan}{}\renewcommand{\rhcaNTreeLowerSigCloseSpan}{1}  % tree classes (of 3): full-set 16:00 QLIKE significantly lower
\providecommand{\rhcaLinPctCloseSpanMin}{}\renewcommand{\rhcaLinPctCloseSpanMin}{0.6}  % linear 16:00 QLIKE %, full vs HAR, Span: min
\providecommand{\rhcaLinPctCloseSpanMax}{}\renewcommand{\rhcaLinPctCloseSpanMax}{4.2}  % ... max
\providecommand{\rhcaLinPctCloseSpanRng}{}\renewcommand{\rhcaLinPctCloseSpanRng}{0.6--4.2}  % ... the range as text
\providecommand{\rhcaTreePctCloseSpanMin}{}\renewcommand{\rhcaTreePctCloseSpanMin}{\ensuremath{-}3.1}  % tree 16:00 QLIKE %, full vs HAR, Span: min
\providecommand{\rhcaTreePctCloseSpanMax}{}\renewcommand{\rhcaTreePctCloseSpanMax}{\ensuremath{-}1.6}  % ... max
\providecommand{\rhcaTreePctCloseSpanRng}{}\renewcommand{\rhcaTreePctCloseSpanRng}{\ensuremath{-}3.1 to \ensuremath{-}1.6}  % ... the range as text
\providecommand{\rhcaLinMsePctCloseSpanMin}{}\renewcommand{\rhcaLinMsePctCloseSpanMin}{77.5}  % linear 16:00 MSE %, full vs HAR, Span: min
\providecommand{\rhcaLinMsePctCloseSpanMax}{}\renewcommand{\rhcaLinMsePctCloseSpanMax}{99.1}  % ... max
\providecommand{\rhcaLinMsePctCloseSpanRng}{}\renewcommand{\rhcaLinMsePctCloseSpanRng}{77.5--99.1}  % ... the range as text
\providecommand{\rhcaTreeMsePctCloseSpanMin}{}\renewcommand{\rhcaTreeMsePctCloseSpanMin}{19.7}  % tree 16:00 MSE %, full vs HAR, Span: min
\providecommand{\rhcaTreeMsePctCloseSpanMax}{}\renewcommand{\rhcaTreeMsePctCloseSpanMax}{33.0}  % ... max
\providecommand{\rhcaTreeMsePctCloseSpanRng}{}\renewcommand{\rhcaTreeMsePctCloseSpanRng}{19.7--33.0}  % ... the range as text
\providecommand{\rhcaNMseHigherCloseSpan}{}\renewcommand{\rhcaNMseHigherCloseSpan}{6}  % classes with a higher 16:00 MSE, full
\providecommand{\rhcaNMseHigherSigCloseSpan}{}\renewcommand{\rhcaNMseHigherSigCloseSpan}{0}  % ... significantly
\providecommand{\rhcaNLowerClosePre}{}\renewcommand{\rhcaNLowerClosePre}{6}  % classes (of 6) with the lower 16:00 QLIKE on the full set, Pre
\providecommand{\rhcaNLowerSigClosePre}{}\renewcommand{\rhcaNLowerSigClosePre}{2}  % ... and DM p < 0.05
\providecommand{\rhcaNHigherClosePre}{}\renewcommand{\rhcaNHigherClosePre}{0}  % classes with the higher 16:00 QLIKE on the full set
\providecommand{\rhcaNHigherSigClosePre}{}\renewcommand{\rhcaNHigherSigClosePre}{0}  % ... and DM p < 0.05
\providecommand{\rhcaNLowerCloseWordPre}{}\renewcommand{\rhcaNLowerCloseWordPre}{six}  % as a word
\providecommand{\rhcaNLinHigherSigClosePre}{}\renewcommand{\rhcaNLinHigherSigClosePre}{0}  % linear classes (of 3): full-set 16:00 QLIKE significantly higher
\providecommand{\rhcaNLinLowerClosePre}{}\renewcommand{\rhcaNLinLowerClosePre}{3}  % linear classes with a lower 16:00 QLIKE
\providecommand{\rhcaNTreeLowerSigClosePre}{}\renewcommand{\rhcaNTreeLowerSigClosePre}{2}  % tree classes (of 3): full-set 16:00 QLIKE significantly lower
\providecommand{\rhcaLinPctClosePreMin}{}\renewcommand{\rhcaLinPctClosePreMin}{\ensuremath{-}4.8}  % linear 16:00 QLIKE %, full vs HAR, Pre: min
\providecommand{\rhcaLinPctClosePreMax}{}\renewcommand{\rhcaLinPctClosePreMax}{\ensuremath{-}2.3}  % ... max
\providecommand{\rhcaLinPctClosePreRng}{}\renewcommand{\rhcaLinPctClosePreRng}{\ensuremath{-}4.8 to \ensuremath{-}2.3}  % ... the range as text
\providecommand{\rhcaTreePctClosePreMin}{}\renewcommand{\rhcaTreePctClosePreMin}{\ensuremath{-}4.0}  % tree 16:00 QLIKE %, full vs HAR, Pre: min
\providecommand{\rhcaTreePctClosePreMax}{}\renewcommand{\rhcaTreePctClosePreMax}{\ensuremath{-}2.1}  % ... max
\providecommand{\rhcaTreePctClosePreRng}{}\renewcommand{\rhcaTreePctClosePreRng}{\ensuremath{-}4.0 to \ensuremath{-}2.1}  % ... the range as text
\providecommand{\rhcaLinMsePctClosePreMin}{}\renewcommand{\rhcaLinMsePctClosePreMin}{\ensuremath{-}15.8}  % linear 16:00 MSE %, full vs HAR, Pre: min
\providecommand{\rhcaLinMsePctClosePreMax}{}\renewcommand{\rhcaLinMsePctClosePreMax}{\ensuremath{-}13.1}  % ... max
\providecommand{\rhcaLinMsePctClosePreRng}{}\renewcommand{\rhcaLinMsePctClosePreRng}{\ensuremath{-}15.8 to \ensuremath{-}13.1}  % ... the range as text
\providecommand{\rhcaTreeMsePctClosePreMin}{}\renewcommand{\rhcaTreeMsePctClosePreMin}{\ensuremath{-}10.0}  % tree 16:00 MSE %, full vs HAR, Pre: min
\providecommand{\rhcaTreeMsePctClosePreMax}{}\renewcommand{\rhcaTreeMsePctClosePreMax}{\ensuremath{-}5.3}  % ... max
\providecommand{\rhcaTreeMsePctClosePreRng}{}\renewcommand{\rhcaTreeMsePctClosePreRng}{\ensuremath{-}10.0 to \ensuremath{-}5.3}  % ... the range as text
\providecommand{\rhcaNMseHigherClosePre}{}\renewcommand{\rhcaNMseHigherClosePre}{0}  % classes with a higher 16:00 MSE, full
\providecommand{\rhcaNMseHigherSigClosePre}{}\renewcommand{\rhcaNMseHigherSigClosePre}{0}  % ... significantly
\providecommand{\rhcaNLowerCloseOos}{}\renewcommand{\rhcaNLowerCloseOos}{3}  % classes (of 6) with the lower 16:00 QLIKE on the full set, Oos
\providecommand{\rhcaNLowerSigCloseOos}{}\renewcommand{\rhcaNLowerSigCloseOos}{0}  % ... and DM p < 0.05
\providecommand{\rhcaNHigherCloseOos}{}\renewcommand{\rhcaNHigherCloseOos}{3}  % classes with the higher 16:00 QLIKE on the full set
\providecommand{\rhcaNHigherSigCloseOos}{}\renewcommand{\rhcaNHigherSigCloseOos}{3}  % ... and DM p < 0.05
\providecommand{\rhcaNLowerCloseWordOos}{}\renewcommand{\rhcaNLowerCloseWordOos}{three}  % as a word
\providecommand{\rhcaNLinHigherSigCloseOos}{}\renewcommand{\rhcaNLinHigherSigCloseOos}{3}  % linear classes (of 3): full-set 16:00 QLIKE significantly higher
\providecommand{\rhcaNLinLowerCloseOos}{}\renewcommand{\rhcaNLinLowerCloseOos}{0}  % linear classes with a lower 16:00 QLIKE
\providecommand{\rhcaNTreeLowerSigCloseOos}{}\renewcommand{\rhcaNTreeLowerSigCloseOos}{0}  % tree classes (of 3): full-set 16:00 QLIKE significantly lower
\providecommand{\rhcaLinPctCloseOosMin}{}\renewcommand{\rhcaLinPctCloseOosMin}{9.5}  % linear 16:00 QLIKE %, full vs HAR, Oos: min
\providecommand{\rhcaLinPctCloseOosMax}{}\renewcommand{\rhcaLinPctCloseOosMax}{15.2}  % ... max
\providecommand{\rhcaLinPctCloseOosRng}{}\renewcommand{\rhcaLinPctCloseOosRng}{9.5--15.2}  % ... the range as text
\providecommand{\rhcaTreePctCloseOosMin}{}\renewcommand{\rhcaTreePctCloseOosMin}{\ensuremath{-}1.7}  % tree 16:00 QLIKE %, full vs HAR, Oos: min
\providecommand{\rhcaTreePctCloseOosMax}{}\renewcommand{\rhcaTreePctCloseOosMax}{\ensuremath{-}0.7}  % ... max
\providecommand{\rhcaTreePctCloseOosRng}{}\renewcommand{\rhcaTreePctCloseOosRng}{\ensuremath{-}1.7 to \ensuremath{-}0.7}  % ... the range as text
\providecommand{\rhcaLinMsePctCloseOosMin}{}\renewcommand{\rhcaLinMsePctCloseOosMin}{87.0}  % linear 16:00 MSE %, full vs HAR, Oos: min
\providecommand{\rhcaLinMsePctCloseOosMax}{}\renewcommand{\rhcaLinMsePctCloseOosMax}{111.1}  % ... max
\providecommand{\rhcaLinMsePctCloseOosRng}{}\renewcommand{\rhcaLinMsePctCloseOosRng}{87.0--111.1}  % ... the range as text
\providecommand{\rhcaTreeMsePctCloseOosMin}{}\renewcommand{\rhcaTreeMsePctCloseOosMin}{23.3}  % tree 16:00 MSE %, full vs HAR, Oos: min
\providecommand{\rhcaTreeMsePctCloseOosMax}{}\renewcommand{\rhcaTreeMsePctCloseOosMax}{38.9}  % ... max
\providecommand{\rhcaTreeMsePctCloseOosRng}{}\renewcommand{\rhcaTreeMsePctCloseOosRng}{23.3--38.9}  % ... the range as text
\providecommand{\rhcaNMseHigherCloseOos}{}\renewcommand{\rhcaNMseHigherCloseOos}{6}  % classes with a higher 16:00 MSE, full
\providecommand{\rhcaNMseHigherSigCloseOos}{}\renewcommand{\rhcaNMseHigherSigCloseOos}{0}  % ... significantly
\providecommand{\rhcaNMseHigherBarSpan}{}\renewcommand{\rhcaNMseHigherBarSpan}{6}  % classes with a higher every-bar MSE, full, Span
\providecommand{\rhcaNMseHigherSigBarSpan}{}\renewcommand{\rhcaNMseHigherSigBarSpan}{0}  % ... significantly
\providecommand{\rhcaNMseLowerBarSpan}{}\renewcommand{\rhcaNMseLowerBarSpan}{0}  % classes with a lower every-bar MSE, full
\providecommand{\rhcaNMseLowerSigBarSpan}{}\renewcommand{\rhcaNMseLowerSigBarSpan}{0}  % ... significantly
\providecommand{\rhcaLinMsePctBarSpanMin}{}\renewcommand{\rhcaLinMsePctBarSpanMin}{100.6}  % linear every-bar MSE %, full vs HAR, Span: min
\providecommand{\rhcaLinMsePctBarSpanMax}{}\renewcommand{\rhcaLinMsePctBarSpanMax}{117.7}  % ... max
\providecommand{\rhcaLinMsePctBarSpanRng}{}\renewcommand{\rhcaLinMsePctBarSpanRng}{100.6--117.7}  % ... the range as text
\providecommand{\rhcaTreeMsePctBarSpanMin}{}\renewcommand{\rhcaTreeMsePctBarSpanMin}{10.9}  % tree every-bar MSE %, full vs HAR, Span: min
\providecommand{\rhcaTreeMsePctBarSpanMax}{}\renewcommand{\rhcaTreeMsePctBarSpanMax}{20.0}  % ... max
\providecommand{\rhcaTreeMsePctBarSpanRng}{}\renewcommand{\rhcaTreeMsePctBarSpanRng}{10.9--20.0}  % ... the range as text
\providecommand{\rhcaNMseHigherBarPre}{}\renewcommand{\rhcaNMseHigherBarPre}{0}  % classes with a higher every-bar MSE, full, Pre
\providecommand{\rhcaNMseHigherSigBarPre}{}\renewcommand{\rhcaNMseHigherSigBarPre}{0}  % ... significantly
\providecommand{\rhcaNMseLowerBarPre}{}\renewcommand{\rhcaNMseLowerBarPre}{6}  % classes with a lower every-bar MSE, full
\providecommand{\rhcaNMseLowerSigBarPre}{}\renewcommand{\rhcaNMseLowerSigBarPre}{2}  % ... significantly
\providecommand{\rhcaLinMsePctBarPreMin}{}\renewcommand{\rhcaLinMsePctBarPreMin}{\ensuremath{-}10.8}  % linear every-bar MSE %, full vs HAR, Pre: min
\providecommand{\rhcaLinMsePctBarPreMax}{}\renewcommand{\rhcaLinMsePctBarPreMax}{\ensuremath{-}8.9}  % ... max
\providecommand{\rhcaLinMsePctBarPreRng}{}\renewcommand{\rhcaLinMsePctBarPreRng}{\ensuremath{-}10.8 to \ensuremath{-}8.9}  % ... the range as text
\providecommand{\rhcaTreeMsePctBarPreMin}{}\renewcommand{\rhcaTreeMsePctBarPreMin}{\ensuremath{-}4.7}  % tree every-bar MSE %, full vs HAR, Pre: min
\providecommand{\rhcaTreeMsePctBarPreMax}{}\renewcommand{\rhcaTreeMsePctBarPreMax}{\ensuremath{-}1.6}  % ... max
\providecommand{\rhcaTreeMsePctBarPreRng}{}\renewcommand{\rhcaTreeMsePctBarPreRng}{\ensuremath{-}4.7 to \ensuremath{-}1.6}  % ... the range as text
\providecommand{\rhcaNMseHigherBarOos}{}\renewcommand{\rhcaNMseHigherBarOos}{6}  % classes with a higher every-bar MSE, full, Oos
\providecommand{\rhcaNMseHigherSigBarOos}{}\renewcommand{\rhcaNMseHigherSigBarOos}{0}  % ... significantly
\providecommand{\rhcaNMseLowerBarOos}{}\renewcommand{\rhcaNMseLowerBarOos}{0}  % classes with a lower every-bar MSE, full
\providecommand{\rhcaNMseLowerSigBarOos}{}\renewcommand{\rhcaNMseLowerSigBarOos}{0}  % ... significantly
\providecommand{\rhcaLinMsePctBarOosMin}{}\renewcommand{\rhcaLinMsePctBarOosMin}{121.6}  % linear every-bar MSE %, full vs HAR, Oos: min
\providecommand{\rhcaLinMsePctBarOosMax}{}\renewcommand{\rhcaLinMsePctBarOosMax}{142.1}  % ... max
\providecommand{\rhcaLinMsePctBarOosRng}{}\renewcommand{\rhcaLinMsePctBarOosRng}{121.6--142.1}  % ... the range as text
\providecommand{\rhcaTreeMsePctBarOosMin}{}\renewcommand{\rhcaTreeMsePctBarOosMin}{13.7}  % tree every-bar MSE %, full vs HAR, Oos: min
\providecommand{\rhcaTreeMsePctBarOosMax}{}\renewcommand{\rhcaTreeMsePctBarOosMax}{25.5}  % ... max
\providecommand{\rhcaTreeMsePctBarOosRng}{}\renewcommand{\rhcaTreeMsePctBarOosRng}{13.7--25.5}  % ... the range as text
\providecommand{\rhcaNCiLowerBm}{}\renewcommand{\rhcaNCiLowerBm}{2}  % closing-interval models (of 12) below the benchmark's QLIKE
\providecommand{\rhcaNCiLowerSigBm}{}\renewcommand{\rhcaNCiLowerSigBm}{0}  % ... significantly
\providecommand{\rhcaNCi}{}\renewcommand{\rhcaNCi}{12}  % closing-interval models other than the benchmark
\providecommand{\rhcaCiHarGapBound}{}\renewcommand{\rhcaCiHarGapBound}{0.8}  % |QLIKE %| of the closing-interval HAR-set linear models vs the benchmark is below
\providecommand{\rhcaLinFullMzBBarPreMin}{}\renewcommand{\rhcaLinFullMzBBarPreMin}{0.99}  % levels MZ slope, linear full, BarPre: min
\providecommand{\rhcaLinFullMzBBarPreMax}{}\renewcommand{\rhcaLinFullMzBBarPreMax}{1.00}  % ... max
\providecommand{\rhcaLinFullMzBBarPreRng}{}\renewcommand{\rhcaLinFullMzBBarPreRng}{0.99--1.00}  % ... the range as text
\providecommand{\rhcaTreeFullMzBBarPreMin}{}\renewcommand{\rhcaTreeFullMzBBarPreMin}{0.92}  % levels MZ slope, tree full, BarPre: min
\providecommand{\rhcaTreeFullMzBBarPreMax}{}\renewcommand{\rhcaTreeFullMzBBarPreMax}{0.97}  % ... max
\providecommand{\rhcaTreeFullMzBBarPreRng}{}\renewcommand{\rhcaTreeFullMzBBarPreRng}{0.92--0.97}  % ... the range as text
\providecommand{\rhcaTreeHarMzBBarPreMin}{}\renewcommand{\rhcaTreeHarMzBBarPreMin}{0.96}  % levels MZ slope, tree HAR, BarPre: min
\providecommand{\rhcaTreeHarMzBBarPreMax}{}\renewcommand{\rhcaTreeHarMzBBarPreMax}{0.99}  % ... max
\providecommand{\rhcaTreeHarMzBBarPreRng}{}\renewcommand{\rhcaTreeHarMzBBarPreRng}{0.96--0.99}  % ... the range as text
\providecommand{\rhcaLinHarMzBBarPre}{}\renewcommand{\rhcaLinHarMzBBarPre}{0.97}  % levels MZ slope, linear HAR (= OLS), BarPre
\providecommand{\rhcaLinFullMzlBBarPreMin}{}\renewcommand{\rhcaLinFullMzlBBarPreMin}{0.99}  % logs MZ slope, linear full, BarPre: min
\providecommand{\rhcaLinFullMzlBBarPreMax}{}\renewcommand{\rhcaLinFullMzlBBarPreMax}{0.99}  % ... max
\providecommand{\rhcaLinFullMzlBBarPreRng}{}\renewcommand{\rhcaLinFullMzlBBarPreRng}{0.99}  % ... the range as text
\providecommand{\rhcaTreeFullMzlBBarPreMin}{}\renewcommand{\rhcaTreeFullMzlBBarPreMin}{1.00}  % logs MZ slope, tree full, BarPre: min
\providecommand{\rhcaTreeFullMzlBBarPreMax}{}\renewcommand{\rhcaTreeFullMzlBBarPreMax}{1.01}  % ... max
\providecommand{\rhcaTreeFullMzlBBarPreRng}{}\renewcommand{\rhcaTreeFullMzlBBarPreRng}{1.00--1.01}  % ... the range as text
\providecommand{\rhcaAllHarMzlBBarPreMin}{}\renewcommand{\rhcaAllHarMzlBBarPreMin}{1.00}  % logs MZ slope, HAR forecasts, BarPre: min
\providecommand{\rhcaAllHarMzlBBarPreMax}{}\renewcommand{\rhcaAllHarMzlBBarPreMax}{1.01}  % ... max
\providecommand{\rhcaAllHarMzlBBarPreRng}{}\renewcommand{\rhcaAllHarMzlBBarPreRng}{1.00--1.01}  % ... the range as text
\providecommand{\rhcaLinFullMzlRsqBarPreMin}{}\renewcommand{\rhcaLinFullMzlRsqBarPreMin}{0.78}  % logs MZ R2, linear full, BarPre: min
\providecommand{\rhcaLinFullMzlRsqBarPreMax}{}\renewcommand{\rhcaLinFullMzlRsqBarPreMax}{0.78}  % ... max
\providecommand{\rhcaLinFullMzlRsqBarPreRng}{}\renewcommand{\rhcaLinFullMzlRsqBarPreRng}{0.78}  % ... the range as text
\providecommand{\rhcaTreeFullMzlRsqBarPreMin}{}\renewcommand{\rhcaTreeFullMzlRsqBarPreMin}{0.79}  % logs MZ R2, tree full, BarPre: min
\providecommand{\rhcaTreeFullMzlRsqBarPreMax}{}\renewcommand{\rhcaTreeFullMzlRsqBarPreMax}{0.79}  % ... max
\providecommand{\rhcaTreeFullMzlRsqBarPreRng}{}\renewcommand{\rhcaTreeFullMzlRsqBarPreRng}{0.79}  % ... the range as text
\providecommand{\rhcaAllHarMzlRsqBarPreMin}{}\renewcommand{\rhcaAllHarMzlRsqBarPreMin}{0.76}  % logs MZ R2, HAR forecasts, BarPre: min
\providecommand{\rhcaAllHarMzlRsqBarPreMax}{}\renewcommand{\rhcaAllHarMzlRsqBarPreMax}{0.76}  % ... max
\providecommand{\rhcaAllHarMzlRsqBarPreRng}{}\renewcommand{\rhcaAllHarMzlRsqBarPreRng}{0.76}  % ... the range as text
\providecommand{\rhcaLinFullMzRsqBarPreMin}{}\renewcommand{\rhcaLinFullMzRsqBarPreMin}{0.53}  % levels MZ R2, linear full, BarPre: min
\providecommand{\rhcaLinFullMzRsqBarPreMax}{}\renewcommand{\rhcaLinFullMzRsqBarPreMax}{0.54}  % ... max
\providecommand{\rhcaLinFullMzRsqBarPreRng}{}\renewcommand{\rhcaLinFullMzRsqBarPreRng}{0.53--0.54}  % ... the range as text
\providecommand{\rhcaAllHarMzRsqBarPreMin}{}\renewcommand{\rhcaAllHarMzRsqBarPreMin}{0.37}  % levels MZ R2, HAR forecasts, BarPre: min
\providecommand{\rhcaAllHarMzRsqBarPreMax}{}\renewcommand{\rhcaAllHarMzRsqBarPreMax}{0.48}  % ... max
\providecommand{\rhcaAllHarMzRsqBarPreRng}{}\renewcommand{\rhcaAllHarMzRsqBarPreRng}{0.37--0.48}  % ... the range as text
\providecommand{\rhcaLinFullShSlopeBarPreMin}{}\renewcommand{\rhcaLinFullShSlopeBarPreMin}{0.0}  % Theil slope share %, linear full, BarPre: min
\providecommand{\rhcaLinFullShSlopeBarPreMax}{}\renewcommand{\rhcaLinFullShSlopeBarPreMax}{0.0}  % ... max
\providecommand{\rhcaLinFullShSlopeBarPreRng}{}\renewcommand{\rhcaLinFullShSlopeBarPreRng}{0.0}  % ... the range as text
\providecommand{\rhcaTreeFullShSlopeBarPreMin}{}\renewcommand{\rhcaTreeFullShSlopeBarPreMin}{0.1}  % Theil slope share %, tree full, BarPre: min
\providecommand{\rhcaTreeFullShSlopeBarPreMax}{}\renewcommand{\rhcaTreeFullShSlopeBarPreMax}{0.5}  % ... max
\providecommand{\rhcaTreeFullShSlopeBarPreRng}{}\renewcommand{\rhcaTreeFullShSlopeBarPreRng}{0.1--0.5}  % ... the range as text
\providecommand{\rhcaAllHarShSlopeBarPreMin}{}\renewcommand{\rhcaAllHarShSlopeBarPreMin}{0.0}  % Theil slope share %, HAR forecasts, BarPre: min
\providecommand{\rhcaAllHarShSlopeBarPreMax}{}\renewcommand{\rhcaAllHarShSlopeBarPreMax}{0.1}  % ... max
\providecommand{\rhcaAllHarShSlopeBarPreRng}{}\renewcommand{\rhcaAllHarShSlopeBarPreRng}{0.0--0.1}  % ... the range as text
\providecommand{\rhcaAllShBiasBarPreMin}{}\renewcommand{\rhcaAllShBiasBarPreMin}{0.1}  % Theil bias share %, all 12, BarPre: min
\providecommand{\rhcaAllShBiasBarPreMax}{}\renewcommand{\rhcaAllShBiasBarPreMax}{0.1}  % ... max
\providecommand{\rhcaAllShBiasBarPreRng}{}\renewcommand{\rhcaAllShBiasBarPreRng}{0.1}  % ... the range as text
\providecommand{\rhcaAllRmBarPreMin}{}\renewcommand{\rhcaAllRmBarPreMin}{1.04}  % ratio of means, all 12, BarPre: min
\providecommand{\rhcaAllRmBarPreMax}{}\renewcommand{\rhcaAllRmBarPreMax}{1.07}  % ... max
\providecommand{\rhcaAllRmBarPreRng}{}\renewcommand{\rhcaAllRmBarPreRng}{1.04--1.07}  % ... the range as text
\providecommand{\rhcaFullRmBarPreMin}{}\renewcommand{\rhcaFullRmBarPreMin}{1.04}  % ratio of means, full set, BarPre: min
\providecommand{\rhcaFullRmBarPreMax}{}\renewcommand{\rhcaFullRmBarPreMax}{1.06}  % ... max
\providecommand{\rhcaFullRmBarPreRng}{}\renewcommand{\rhcaFullRmBarPreRng}{1.04--1.06}  % ... the range as text
\providecommand{\rhcaHarRmBarPreMin}{}\renewcommand{\rhcaHarRmBarPreMin}{1.06}  % ratio of means, HAR, BarPre: min
\providecommand{\rhcaHarRmBarPreMax}{}\renewcommand{\rhcaHarRmBarPreMax}{1.07}  % ... max
\providecommand{\rhcaHarRmBarPreRng}{}\renewcommand{\rhcaHarRmBarPreRng}{1.06--1.07}  % ... the range as text
\providecommand{\rhcaAllMrBarPreMin}{}\renewcommand{\rhcaAllMrBarPreMin}{1.07}  % mean of RV/sigma2, all 12, BarPre: min
\providecommand{\rhcaAllMrBarPreMax}{}\renewcommand{\rhcaAllMrBarPreMax}{1.10}  % ... max
\providecommand{\rhcaAllMrBarPreRng}{}\renewcommand{\rhcaAllMrBarPreRng}{1.07--1.10}  % ... the range as text
\providecommand{\rhcaAllUnderBarPreMin}{}\renewcommand{\rhcaAllUnderBarPreMin}{39.9}  % under-prediction %, all 12, BarPre: min
\providecommand{\rhcaAllUnderBarPreMax}{}\renewcommand{\rhcaAllUnderBarPreMax}{42.1}  % ... max
\providecommand{\rhcaAllUnderBarPreRng}{}\renewcommand{\rhcaAllUnderBarPreRng}{39.9--42.1}  % ... the range as text
\providecommand{\rhcaLinFullUnderBarPreMin}{}\renewcommand{\rhcaLinFullUnderBarPreMin}{41.5}  % under-prediction %, linear full, BarPre: min
\providecommand{\rhcaLinFullUnderBarPreMax}{}\renewcommand{\rhcaLinFullUnderBarPreMax}{42.1}  % ... max
\providecommand{\rhcaLinFullUnderBarPreRng}{}\renewcommand{\rhcaLinFullUnderBarPreRng}{41.5--42.1}  % ... the range as text
\providecommand{\rhcaLinHarUnderBarPre}{}\renewcommand{\rhcaLinHarUnderBarPre}{39.9}  % under-prediction %, linear HAR, BarPre
\providecommand{\rhcaLinFullOverSeBarPreMin}{}\renewcommand{\rhcaLinFullOverSeBarPreMin}{19.6}  % % of SE from over-predictions, linear full, BarPre: min
\providecommand{\rhcaLinFullOverSeBarPreMax}{}\renewcommand{\rhcaLinFullOverSeBarPreMax}{20.1}  % ... max
\providecommand{\rhcaLinFullOverSeBarPreRng}{}\renewcommand{\rhcaLinFullOverSeBarPreRng}{19.6--20.1}  % ... the range as text
\providecommand{\rhcaTreeFullOverSeBarPreMin}{}\renewcommand{\rhcaTreeFullOverSeBarPreMin}{12.0}  % % of SE from over-predictions, tree full, BarPre: min
\providecommand{\rhcaTreeFullOverSeBarPreMax}{}\renewcommand{\rhcaTreeFullOverSeBarPreMax}{15.1}  % ... max
\providecommand{\rhcaTreeFullOverSeBarPreRng}{}\renewcommand{\rhcaTreeFullOverSeBarPreRng}{12.0--15.1}  % ... the range as text
\providecommand{\rhcaAllHarOverSeBarPreMin}{}\renewcommand{\rhcaAllHarOverSeBarPreMin}{10.3}  % % of SE from over-predictions, HAR, BarPre: min
\providecommand{\rhcaAllHarOverSeBarPreMax}{}\renewcommand{\rhcaAllHarOverSeBarPreMax}{20.0}  % ... max
\providecommand{\rhcaAllHarOverSeBarPreRng}{}\renewcommand{\rhcaAllHarOverSeBarPreRng}{10.3--20.0}  % ... the range as text
\providecommand{\rhcaNMaeLowerBarPre}{}\renewcommand{\rhcaNMaeLowerBarPre}{6}  % classes (of 6) with the lower MAE on the full set, BarPre
\providecommand{\rhcaLinFullMrBarPreMin}{}\renewcommand{\rhcaLinFullMrBarPreMin}{1.091}  % mean of RV/sigma2, linear full, BarPre: min
\providecommand{\rhcaLinFullMrBarPreMax}{}\renewcommand{\rhcaLinFullMrBarPreMax}{1.096}  % ... max
\providecommand{\rhcaLinFullMrBarPreRng}{}\renewcommand{\rhcaLinFullMrBarPreRng}{1.091--1.096}  % ... the range as text
\providecommand{\rhcaLinHarMrBarPre}{}\renewcommand{\rhcaLinHarMrBarPre}{1.078}  % mean of RV/sigma2, linear HAR, BarPre
\providecommand{\rhcaLinFullQlBarPreMin}{}\renewcommand{\rhcaLinFullQlBarPreMin}{0.0039}  % QLIKE level term, linear full, BarPre: min
\providecommand{\rhcaLinFullQlBarPreMax}{}\renewcommand{\rhcaLinFullQlBarPreMax}{0.0043}  % ... max
\providecommand{\rhcaLinFullQlBarPreRng}{}\renewcommand{\rhcaLinFullQlBarPreRng}{0.0039--0.0043}  % ... the range as text
\providecommand{\rhcaLinHarQlBarPre}{}\renewcommand{\rhcaLinHarQlBarPre}{0.0029}  % QLIKE level term, linear HAR, BarPre
\providecommand{\rhcaLinFullQdBarPreMin}{}\renewcommand{\rhcaLinFullQdBarPreMin}{0.1627}  % QLIKE dispersion term, linear full, BarPre: min
\providecommand{\rhcaLinFullQdBarPreMax}{}\renewcommand{\rhcaLinFullQdBarPreMax}{0.1683}  % ... max
\providecommand{\rhcaLinFullQdBarPreRng}{}\renewcommand{\rhcaLinFullQdBarPreRng}{0.1627--0.1683}  % ... the range as text
\providecommand{\rhcaLinHarQdBarPre}{}\renewcommand{\rhcaLinHarQdBarPre}{0.1835}  % QLIKE dispersion term, linear HAR, BarPre
\providecommand{\rhcaTreeFullQlBarPreMin}{}\renewcommand{\rhcaTreeFullQlBarPreMin}{0.0021}  % QLIKE level term, tree full, BarPre: min
\providecommand{\rhcaTreeFullQlBarPreMax}{}\renewcommand{\rhcaTreeFullQlBarPreMax}{0.0029}  % ... max
\providecommand{\rhcaTreeFullQlBarPreRng}{}\renewcommand{\rhcaTreeFullQlBarPreRng}{0.0021--0.0029}  % ... the range as text
\providecommand{\rhcaTreeHarQlBarPreMin}{}\renewcommand{\rhcaTreeHarQlBarPreMin}{0.0032}  % QLIKE level term, tree HAR, BarPre: min
\providecommand{\rhcaTreeHarQlBarPreMax}{}\renewcommand{\rhcaTreeHarQlBarPreMax}{0.0034}  % ... max
\providecommand{\rhcaTreeHarQlBarPreRng}{}\renewcommand{\rhcaTreeHarQlBarPreRng}{0.0032--0.0034}  % ... the range as text
\providecommand{\rhcaTreeFullQdBarPreMin}{}\renewcommand{\rhcaTreeFullQdBarPreMin}{0.1563}  % QLIKE dispersion term, tree full, BarPre: min
\providecommand{\rhcaTreeFullQdBarPreMax}{}\renewcommand{\rhcaTreeFullQdBarPreMax}{0.1617}  % ... max
\providecommand{\rhcaTreeFullQdBarPreRng}{}\renewcommand{\rhcaTreeFullQdBarPreRng}{0.1563--0.1617}  % ... the range as text
\providecommand{\rhcaTreeHarQdBarPreMin}{}\renewcommand{\rhcaTreeHarQdBarPreMin}{0.1826}  % QLIKE dispersion term, tree HAR, BarPre: min
\providecommand{\rhcaTreeHarQdBarPreMax}{}\renewcommand{\rhcaTreeHarQdBarPreMax}{0.1832}  % ... max
\providecommand{\rhcaTreeHarQdBarPreRng}{}\renewcommand{\rhcaTreeHarQdBarPreRng}{0.1826--0.1832}  % ... the range as text
\providecommand{\rhcaTreeFullUnderBarPreMin}{}\renewcommand{\rhcaTreeFullUnderBarPreMin}{40.8}  % under-prediction %, tree full, BarPre: min
\providecommand{\rhcaTreeFullUnderBarPreMax}{}\renewcommand{\rhcaTreeFullUnderBarPreMax}{41.1}  % ... max
\providecommand{\rhcaTreeFullUnderBarPreRng}{}\renewcommand{\rhcaTreeFullUnderBarPreRng}{40.8--41.1}  % ... the range as text
\providecommand{\rhcaTreeHarUnderBarPreMin}{}\renewcommand{\rhcaTreeHarUnderBarPreMin}{40.1}  % under-prediction %, tree HAR, BarPre: min
\providecommand{\rhcaTreeHarUnderBarPreMax}{}\renewcommand{\rhcaTreeHarUnderBarPreMax}{40.2}  % ... max
\providecommand{\rhcaTreeHarUnderBarPreRng}{}\renewcommand{\rhcaTreeHarUnderBarPreRng}{40.1--40.2}  % ... the range as text
\providecommand{\rhcaLinFullMzBBarOosMin}{}\renewcommand{\rhcaLinFullMzBBarOosMin}{0.56}  % levels MZ slope, linear full, BarOos: min
\providecommand{\rhcaLinFullMzBBarOosMax}{}\renewcommand{\rhcaLinFullMzBBarOosMax}{0.58}  % ... max
\providecommand{\rhcaLinFullMzBBarOosRng}{}\renewcommand{\rhcaLinFullMzBBarOosRng}{0.56--0.58}  % ... the range as text
\providecommand{\rhcaTreeFullMzBBarOosMin}{}\renewcommand{\rhcaTreeFullMzBBarOosMin}{0.74}  % levels MZ slope, tree full, BarOos: min
\providecommand{\rhcaTreeFullMzBBarOosMax}{}\renewcommand{\rhcaTreeFullMzBBarOosMax}{0.78}  % ... max
\providecommand{\rhcaTreeFullMzBBarOosRng}{}\renewcommand{\rhcaTreeFullMzBBarOosRng}{0.74--0.78}  % ... the range as text
\providecommand{\rhcaTreeHarMzBBarOosMin}{}\renewcommand{\rhcaTreeHarMzBBarOosMin}{0.92}  % levels MZ slope, tree HAR, BarOos: min
\providecommand{\rhcaTreeHarMzBBarOosMax}{}\renewcommand{\rhcaTreeHarMzBBarOosMax}{0.94}  % ... max
\providecommand{\rhcaTreeHarMzBBarOosRng}{}\renewcommand{\rhcaTreeHarMzBBarOosRng}{0.92--0.94}  % ... the range as text
\providecommand{\rhcaLinHarMzBBarOos}{}\renewcommand{\rhcaLinHarMzBBarOos}{0.99}  % levels MZ slope, linear HAR (= OLS), BarOos
\providecommand{\rhcaLinFullMzlBBarOosMin}{}\renewcommand{\rhcaLinFullMzlBBarOosMin}{0.98}  % logs MZ slope, linear full, BarOos: min
\providecommand{\rhcaLinFullMzlBBarOosMax}{}\renewcommand{\rhcaLinFullMzlBBarOosMax}{0.99}  % ... max
\providecommand{\rhcaLinFullMzlBBarOosRng}{}\renewcommand{\rhcaLinFullMzlBBarOosRng}{0.98--0.99}  % ... the range as text
\providecommand{\rhcaTreeFullMzlBBarOosMin}{}\renewcommand{\rhcaTreeFullMzlBBarOosMin}{0.99}  % logs MZ slope, tree full, BarOos: min
\providecommand{\rhcaTreeFullMzlBBarOosMax}{}\renewcommand{\rhcaTreeFullMzlBBarOosMax}{0.99}  % ... max
\providecommand{\rhcaTreeFullMzlBBarOosRng}{}\renewcommand{\rhcaTreeFullMzlBBarOosRng}{0.99}  % ... the range as text
\providecommand{\rhcaAllHarMzlBBarOosMin}{}\renewcommand{\rhcaAllHarMzlBBarOosMin}{1.00}  % logs MZ slope, HAR forecasts, BarOos: min
\providecommand{\rhcaAllHarMzlBBarOosMax}{}\renewcommand{\rhcaAllHarMzlBBarOosMax}{1.02}  % ... max
\providecommand{\rhcaAllHarMzlBBarOosRng}{}\renewcommand{\rhcaAllHarMzlBBarOosRng}{1.00--1.02}  % ... the range as text
\providecommand{\rhcaLinFullMzlRsqBarOosMin}{}\renewcommand{\rhcaLinFullMzlRsqBarOosMin}{0.81}  % logs MZ R2, linear full, BarOos: min
\providecommand{\rhcaLinFullMzlRsqBarOosMax}{}\renewcommand{\rhcaLinFullMzlRsqBarOosMax}{0.81}  % ... max
\providecommand{\rhcaLinFullMzlRsqBarOosRng}{}\renewcommand{\rhcaLinFullMzlRsqBarOosRng}{0.81}  % ... the range as text
\providecommand{\rhcaTreeFullMzlRsqBarOosMin}{}\renewcommand{\rhcaTreeFullMzlRsqBarOosMin}{0.82}  % logs MZ R2, tree full, BarOos: min
\providecommand{\rhcaTreeFullMzlRsqBarOosMax}{}\renewcommand{\rhcaTreeFullMzlRsqBarOosMax}{0.82}  % ... max
\providecommand{\rhcaTreeFullMzlRsqBarOosRng}{}\renewcommand{\rhcaTreeFullMzlRsqBarOosRng}{0.82}  % ... the range as text
\providecommand{\rhcaAllHarMzlRsqBarOosMin}{}\renewcommand{\rhcaAllHarMzlRsqBarOosMin}{0.79}  % logs MZ R2, HAR forecasts, BarOos: min
\providecommand{\rhcaAllHarMzlRsqBarOosMax}{}\renewcommand{\rhcaAllHarMzlRsqBarOosMax}{0.80}  % ... max
\providecommand{\rhcaAllHarMzlRsqBarOosRng}{}\renewcommand{\rhcaAllHarMzlRsqBarOosRng}{0.79--0.80}  % ... the range as text
\providecommand{\rhcaLinFullMzRsqBarOosMin}{}\renewcommand{\rhcaLinFullMzRsqBarOosMin}{0.51}  % levels MZ R2, linear full, BarOos: min
\providecommand{\rhcaLinFullMzRsqBarOosMax}{}\renewcommand{\rhcaLinFullMzRsqBarOosMax}{0.52}  % ... max
\providecommand{\rhcaLinFullMzRsqBarOosRng}{}\renewcommand{\rhcaLinFullMzRsqBarOosRng}{0.51--0.52}  % ... the range as text
\providecommand{\rhcaAllHarMzRsqBarOosMin}{}\renewcommand{\rhcaAllHarMzRsqBarOosMin}{0.65}  % levels MZ R2, HAR forecasts, BarOos: min
\providecommand{\rhcaAllHarMzRsqBarOosMax}{}\renewcommand{\rhcaAllHarMzRsqBarOosMax}{0.67}  % ... max
\providecommand{\rhcaAllHarMzRsqBarOosRng}{}\renewcommand{\rhcaAllHarMzRsqBarOosRng}{0.65--0.67}  % ... the range as text
\providecommand{\rhcaLinFullShSlopeBarOosMin}{}\renewcommand{\rhcaLinFullShSlopeBarOosMin}{35.9}  % Theil slope share %, linear full, BarOos: min
\providecommand{\rhcaLinFullShSlopeBarOosMax}{}\renewcommand{\rhcaLinFullShSlopeBarOosMax}{40.8}  % ... max
\providecommand{\rhcaLinFullShSlopeBarOosRng}{}\renewcommand{\rhcaLinFullShSlopeBarOosRng}{35.9--40.8}  % ... the range as text
\providecommand{\rhcaTreeFullShSlopeBarOosMin}{}\renewcommand{\rhcaTreeFullShSlopeBarOosMin}{12.7}  % Theil slope share %, tree full, BarOos: min
\providecommand{\rhcaTreeFullShSlopeBarOosMax}{}\renewcommand{\rhcaTreeFullShSlopeBarOosMax}{17.7}  % ... max
\providecommand{\rhcaTreeFullShSlopeBarOosRng}{}\renewcommand{\rhcaTreeFullShSlopeBarOosRng}{12.7--17.7}  % ... the range as text
\providecommand{\rhcaAllHarShSlopeBarOosMin}{}\renewcommand{\rhcaAllHarShSlopeBarOosMin}{0.0}  % Theil slope share %, HAR forecasts, BarOos: min
\providecommand{\rhcaAllHarShSlopeBarOosMax}{}\renewcommand{\rhcaAllHarShSlopeBarOosMax}{1.4}  % ... max
\providecommand{\rhcaAllHarShSlopeBarOosRng}{}\renewcommand{\rhcaAllHarShSlopeBarOosRng}{0.0--1.4}  % ... the range as text
\providecommand{\rhcaAllShBiasBarOosMin}{}\renewcommand{\rhcaAllShBiasBarOosMin}{0.0}  % Theil bias share %, all 12, BarOos: min
\providecommand{\rhcaAllShBiasBarOosMax}{}\renewcommand{\rhcaAllShBiasBarOosMax}{0.1}  % ... max
\providecommand{\rhcaAllShBiasBarOosRng}{}\renewcommand{\rhcaAllShBiasBarOosRng}{0.0--0.1}  % ... the range as text
\providecommand{\rhcaAllRmBarOosMin}{}\renewcommand{\rhcaAllRmBarOosMin}{0.96}  % ratio of means, all 12, BarOos: min
\providecommand{\rhcaAllRmBarOosMax}{}\renewcommand{\rhcaAllRmBarOosMax}{1.05}  % ... max
\providecommand{\rhcaAllRmBarOosRng}{}\renewcommand{\rhcaAllRmBarOosRng}{0.96--1.05}  % ... the range as text
\providecommand{\rhcaFullRmBarOosMin}{}\renewcommand{\rhcaFullRmBarOosMin}{0.96}  % ratio of means, full set, BarOos: min
\providecommand{\rhcaFullRmBarOosMax}{}\renewcommand{\rhcaFullRmBarOosMax}{0.98}  % ... max
\providecommand{\rhcaFullRmBarOosRng}{}\renewcommand{\rhcaFullRmBarOosRng}{0.96--0.98}  % ... the range as text
\providecommand{\rhcaHarRmBarOosMin}{}\renewcommand{\rhcaHarRmBarOosMin}{1.02}  % ratio of means, HAR, BarOos: min
\providecommand{\rhcaHarRmBarOosMax}{}\renewcommand{\rhcaHarRmBarOosMax}{1.05}  % ... max
\providecommand{\rhcaHarRmBarOosRng}{}\renewcommand{\rhcaHarRmBarOosRng}{1.02--1.05}  % ... the range as text
\providecommand{\rhcaAllMrBarOosMin}{}\renewcommand{\rhcaAllMrBarOosMin}{1.06}  % mean of RV/sigma2, all 12, BarOos: min
\providecommand{\rhcaAllMrBarOosMax}{}\renewcommand{\rhcaAllMrBarOosMax}{1.12}  % ... max
\providecommand{\rhcaAllMrBarOosRng}{}\renewcommand{\rhcaAllMrBarOosRng}{1.06--1.12}  % ... the range as text
\providecommand{\rhcaAllUnderBarOosMin}{}\renewcommand{\rhcaAllUnderBarOosMin}{38.0}  % under-prediction %, all 12, BarOos: min
\providecommand{\rhcaAllUnderBarOosMax}{}\renewcommand{\rhcaAllUnderBarOosMax}{43.2}  % ... max
\providecommand{\rhcaAllUnderBarOosRng}{}\renewcommand{\rhcaAllUnderBarOosRng}{38.0--43.2}  % ... the range as text
\providecommand{\rhcaLinFullUnderBarOosMin}{}\renewcommand{\rhcaLinFullUnderBarOosMin}{42.2}  % under-prediction %, linear full, BarOos: min
\providecommand{\rhcaLinFullUnderBarOosMax}{}\renewcommand{\rhcaLinFullUnderBarOosMax}{43.2}  % ... max
\providecommand{\rhcaLinFullUnderBarOosRng}{}\renewcommand{\rhcaLinFullUnderBarOosRng}{42.2--43.2}  % ... the range as text
\providecommand{\rhcaLinHarUnderBarOos}{}\renewcommand{\rhcaLinHarUnderBarOos}{38.0}  % under-prediction %, linear HAR, BarOos
\providecommand{\rhcaLinFullOverSeBarOosMin}{}\renewcommand{\rhcaLinFullOverSeBarOosMin}{63.1}  % % of SE from over-predictions, linear full, BarOos: min
\providecommand{\rhcaLinFullOverSeBarOosMax}{}\renewcommand{\rhcaLinFullOverSeBarOosMax}{67.0}  % ... max
\providecommand{\rhcaLinFullOverSeBarOosRng}{}\renewcommand{\rhcaLinFullOverSeBarOosRng}{63.1--67.0}  % ... the range as text
\providecommand{\rhcaTreeFullOverSeBarOosMin}{}\renewcommand{\rhcaTreeFullOverSeBarOosMin}{44.4}  % % of SE from over-predictions, tree full, BarOos: min
\providecommand{\rhcaTreeFullOverSeBarOosMax}{}\renewcommand{\rhcaTreeFullOverSeBarOosMax}{49.1}  % ... max
\providecommand{\rhcaTreeFullOverSeBarOosRng}{}\renewcommand{\rhcaTreeFullOverSeBarOosRng}{44.4--49.1}  % ... the range as text
\providecommand{\rhcaAllHarOverSeBarOosMin}{}\renewcommand{\rhcaAllHarOverSeBarOosMin}{21.5}  % % of SE from over-predictions, HAR, BarOos: min
\providecommand{\rhcaAllHarOverSeBarOosMax}{}\renewcommand{\rhcaAllHarOverSeBarOosMax}{26.7}  % ... max
\providecommand{\rhcaAllHarOverSeBarOosRng}{}\renewcommand{\rhcaAllHarOverSeBarOosRng}{21.5--26.7}  % ... the range as text
\providecommand{\rhcaNMaeLowerBarOos}{}\renewcommand{\rhcaNMaeLowerBarOos}{0}  % classes (of 6) with the lower MAE on the full set, BarOos
\providecommand{\rhcaLinFullMrBarOosMin}{}\renewcommand{\rhcaLinFullMrBarOosMin}{1.109}  % mean of RV/sigma2, linear full, BarOos: min
\providecommand{\rhcaLinFullMrBarOosMax}{}\renewcommand{\rhcaLinFullMrBarOosMax}{1.117}  % ... max
\providecommand{\rhcaLinFullMrBarOosRng}{}\renewcommand{\rhcaLinFullMrBarOosRng}{1.109--1.117}  % ... the range as text
\providecommand{\rhcaLinHarMrBarOos}{}\renewcommand{\rhcaLinHarMrBarOos}{1.061}  % mean of RV/sigma2, linear HAR, BarOos
\providecommand{\rhcaLinFullQlBarOosMin}{}\renewcommand{\rhcaLinFullQlBarOosMin}{0.0056}  % QLIKE level term, linear full, BarOos: min
\providecommand{\rhcaLinFullQlBarOosMax}{}\renewcommand{\rhcaLinFullQlBarOosMax}{0.0064}  % ... max
\providecommand{\rhcaLinFullQlBarOosRng}{}\renewcommand{\rhcaLinFullQlBarOosRng}{0.0056--0.0064}  % ... the range as text
\providecommand{\rhcaLinHarQlBarOos}{}\renewcommand{\rhcaLinHarQlBarOos}{0.0018}  % QLIKE level term, linear HAR, BarOos
\providecommand{\rhcaLinFullQdBarOosMin}{}\renewcommand{\rhcaLinFullQdBarOosMin}{0.1739}  % QLIKE dispersion term, linear full, BarOos: min
\providecommand{\rhcaLinFullQdBarOosMax}{}\renewcommand{\rhcaLinFullQdBarOosMax}{0.1788}  % ... max
\providecommand{\rhcaLinFullQdBarOosRng}{}\renewcommand{\rhcaLinFullQdBarOosRng}{0.1739--0.1788}  % ... the range as text
\providecommand{\rhcaLinHarQdBarOos}{}\renewcommand{\rhcaLinHarQdBarOos}{0.1915}  % QLIKE dispersion term, linear HAR, BarOos
\providecommand{\rhcaTreeFullQlBarOosMin}{}\renewcommand{\rhcaTreeFullQlBarOosMin}{0.0022}  % QLIKE level term, tree full, BarOos: min
\providecommand{\rhcaTreeFullQlBarOosMax}{}\renewcommand{\rhcaTreeFullQlBarOosMax}{0.0035}  % ... max
\providecommand{\rhcaTreeFullQlBarOosRng}{}\renewcommand{\rhcaTreeFullQlBarOosRng}{0.0022--0.0035}  % ... the range as text
\providecommand{\rhcaTreeHarQlBarOosMin}{}\renewcommand{\rhcaTreeHarQlBarOosMin}{0.0022}  % QLIKE level term, tree HAR, BarOos: min
\providecommand{\rhcaTreeHarQlBarOosMax}{}\renewcommand{\rhcaTreeHarQlBarOosMax}{0.0028}  % ... max
\providecommand{\rhcaTreeHarQlBarOosRng}{}\renewcommand{\rhcaTreeHarQlBarOosRng}{0.0022--0.0028}  % ... the range as text
\providecommand{\rhcaTreeFullQdBarOosMin}{}\renewcommand{\rhcaTreeFullQdBarOosMin}{0.1612}  % QLIKE dispersion term, tree full, BarOos: min
\providecommand{\rhcaTreeFullQdBarOosMax}{}\renewcommand{\rhcaTreeFullQdBarOosMax}{0.1694}  % ... max
\providecommand{\rhcaTreeFullQdBarOosRng}{}\renewcommand{\rhcaTreeFullQdBarOosRng}{0.1612--0.1694}  % ... the range as text
\providecommand{\rhcaTreeHarQdBarOosMin}{}\renewcommand{\rhcaTreeHarQdBarOosMin}{0.1878}  % QLIKE dispersion term, tree HAR, BarOos: min
\providecommand{\rhcaTreeHarQdBarOosMax}{}\renewcommand{\rhcaTreeHarQdBarOosMax}{0.1909}  % ... max
\providecommand{\rhcaTreeHarQdBarOosRng}{}\renewcommand{\rhcaTreeHarQdBarOosRng}{0.1878--0.1909}  % ... the range as text
\providecommand{\rhcaTreeFullUnderBarOosMin}{}\renewcommand{\rhcaTreeFullUnderBarOosMin}{40.5}  % under-prediction %, tree full, BarOos: min
\providecommand{\rhcaTreeFullUnderBarOosMax}{}\renewcommand{\rhcaTreeFullUnderBarOosMax}{41.0}  % ... max
\providecommand{\rhcaTreeFullUnderBarOosRng}{}\renewcommand{\rhcaTreeFullUnderBarOosRng}{40.5--41.0}  % ... the range as text
\providecommand{\rhcaTreeHarUnderBarOosMin}{}\renewcommand{\rhcaTreeHarUnderBarOosMin}{38.4}  % under-prediction %, tree HAR, BarOos: min
\providecommand{\rhcaTreeHarUnderBarOosMax}{}\renewcommand{\rhcaTreeHarUnderBarOosMax}{38.8}  % ... max
\providecommand{\rhcaTreeHarUnderBarOosRng}{}\renewcommand{\rhcaTreeHarUnderBarOosRng}{38.4--38.8}  % ... the range as text
\providecommand{\rhcaLinFullMzBClosePreMin}{}\renewcommand{\rhcaLinFullMzBClosePreMin}{1.06}  % levels MZ slope, linear full, ClosePre: min
\providecommand{\rhcaLinFullMzBClosePreMax}{}\renewcommand{\rhcaLinFullMzBClosePreMax}{1.08}  % ... max
\providecommand{\rhcaLinFullMzBClosePreRng}{}\renewcommand{\rhcaLinFullMzBClosePreRng}{1.06--1.08}  % ... the range as text
\providecommand{\rhcaTreeFullMzBClosePreMin}{}\renewcommand{\rhcaTreeFullMzBClosePreMin}{0.98}  % levels MZ slope, tree full, ClosePre: min
\providecommand{\rhcaTreeFullMzBClosePreMax}{}\renewcommand{\rhcaTreeFullMzBClosePreMax}{0.99}  % ... max
\providecommand{\rhcaTreeFullMzBClosePreRng}{}\renewcommand{\rhcaTreeFullMzBClosePreRng}{0.98--0.99}  % ... the range as text
\providecommand{\rhcaTreeHarMzBClosePreMin}{}\renewcommand{\rhcaTreeHarMzBClosePreMin}{1.06}  % levels MZ slope, tree HAR, ClosePre: min
\providecommand{\rhcaTreeHarMzBClosePreMax}{}\renewcommand{\rhcaTreeHarMzBClosePreMax}{1.11}  % ... max
\providecommand{\rhcaTreeHarMzBClosePreRng}{}\renewcommand{\rhcaTreeHarMzBClosePreRng}{1.06--1.11}  % ... the range as text
\providecommand{\rhcaLinHarMzBClosePre}{}\renewcommand{\rhcaLinHarMzBClosePre}{1.14}  % levels MZ slope, linear HAR (= OLS), ClosePre
\providecommand{\rhcaLinFullMzlBClosePreMin}{}\renewcommand{\rhcaLinFullMzlBClosePreMin}{0.92}  % logs MZ slope, linear full, ClosePre: min
\providecommand{\rhcaLinFullMzlBClosePreMax}{}\renewcommand{\rhcaLinFullMzlBClosePreMax}{0.93}  % ... max
\providecommand{\rhcaLinFullMzlBClosePreRng}{}\renewcommand{\rhcaLinFullMzlBClosePreRng}{0.92--0.93}  % ... the range as text
\providecommand{\rhcaTreeFullMzlBClosePreMin}{}\renewcommand{\rhcaTreeFullMzlBClosePreMin}{0.92}  % logs MZ slope, tree full, ClosePre: min
\providecommand{\rhcaTreeFullMzlBClosePreMax}{}\renewcommand{\rhcaTreeFullMzlBClosePreMax}{0.94}  % ... max
\providecommand{\rhcaTreeFullMzlBClosePreRng}{}\renewcommand{\rhcaTreeFullMzlBClosePreRng}{0.92--0.94}  % ... the range as text
\providecommand{\rhcaAllHarMzlBClosePreMin}{}\renewcommand{\rhcaAllHarMzlBClosePreMin}{0.95}  % logs MZ slope, HAR forecasts, ClosePre: min
\providecommand{\rhcaAllHarMzlBClosePreMax}{}\renewcommand{\rhcaAllHarMzlBClosePreMax}{0.96}  % ... max
\providecommand{\rhcaAllHarMzlBClosePreRng}{}\renewcommand{\rhcaAllHarMzlBClosePreRng}{0.95--0.96}  % ... the range as text
\providecommand{\rhcaLinFullMzlRsqClosePreMin}{}\renewcommand{\rhcaLinFullMzlRsqClosePreMin}{0.80}  % logs MZ R2, linear full, ClosePre: min
\providecommand{\rhcaLinFullMzlRsqClosePreMax}{}\renewcommand{\rhcaLinFullMzlRsqClosePreMax}{0.80}  % ... max
\providecommand{\rhcaLinFullMzlRsqClosePreRng}{}\renewcommand{\rhcaLinFullMzlRsqClosePreRng}{0.80}  % ... the range as text
\providecommand{\rhcaTreeFullMzlRsqClosePreMin}{}\renewcommand{\rhcaTreeFullMzlRsqClosePreMin}{0.79}  % logs MZ R2, tree full, ClosePre: min
\providecommand{\rhcaTreeFullMzlRsqClosePreMax}{}\renewcommand{\rhcaTreeFullMzlRsqClosePreMax}{0.80}  % ... max
\providecommand{\rhcaTreeFullMzlRsqClosePreRng}{}\renewcommand{\rhcaTreeFullMzlRsqClosePreRng}{0.79--0.80}  % ... the range as text
\providecommand{\rhcaAllHarMzlRsqClosePreMin}{}\renewcommand{\rhcaAllHarMzlRsqClosePreMin}{0.78}  % logs MZ R2, HAR forecasts, ClosePre: min
\providecommand{\rhcaAllHarMzlRsqClosePreMax}{}\renewcommand{\rhcaAllHarMzlRsqClosePreMax}{0.78}  % ... max
\providecommand{\rhcaAllHarMzlRsqClosePreRng}{}\renewcommand{\rhcaAllHarMzlRsqClosePreRng}{0.78}  % ... the range as text
\providecommand{\rhcaLinFullMzRsqClosePreMin}{}\renewcommand{\rhcaLinFullMzRsqClosePreMin}{0.68}  % levels MZ R2, linear full, ClosePre: min
\providecommand{\rhcaLinFullMzRsqClosePreMax}{}\renewcommand{\rhcaLinFullMzRsqClosePreMax}{0.69}  % ... max
\providecommand{\rhcaLinFullMzRsqClosePreRng}{}\renewcommand{\rhcaLinFullMzRsqClosePreRng}{0.68--0.69}  % ... the range as text
\providecommand{\rhcaAllHarMzRsqClosePreMin}{}\renewcommand{\rhcaAllHarMzRsqClosePreMin}{0.55}  % levels MZ R2, HAR forecasts, ClosePre: min
\providecommand{\rhcaAllHarMzRsqClosePreMax}{}\renewcommand{\rhcaAllHarMzRsqClosePreMax}{0.64}  % ... max
\providecommand{\rhcaAllHarMzRsqClosePreRng}{}\renewcommand{\rhcaAllHarMzRsqClosePreRng}{0.55--0.64}  % ... the range as text
\providecommand{\rhcaLinFullShSlopeClosePreMin}{}\renewcommand{\rhcaLinFullShSlopeClosePreMin}{0.7}  % Theil slope share %, linear full, ClosePre: min
\providecommand{\rhcaLinFullShSlopeClosePreMax}{}\renewcommand{\rhcaLinFullShSlopeClosePreMax}{1.2}  % ... max
\providecommand{\rhcaLinFullShSlopeClosePreRng}{}\renewcommand{\rhcaLinFullShSlopeClosePreRng}{0.7--1.2}  % ... the range as text
\providecommand{\rhcaTreeFullShSlopeClosePreMin}{}\renewcommand{\rhcaTreeFullShSlopeClosePreMin}{0.0}  % Theil slope share %, tree full, ClosePre: min
\providecommand{\rhcaTreeFullShSlopeClosePreMax}{}\renewcommand{\rhcaTreeFullShSlopeClosePreMax}{0.1}  % ... max
\providecommand{\rhcaTreeFullShSlopeClosePreRng}{}\renewcommand{\rhcaTreeFullShSlopeClosePreRng}{0.0--0.1}  % ... the range as text
\providecommand{\rhcaAllHarShSlopeClosePreMin}{}\renewcommand{\rhcaAllHarShSlopeClosePreMin}{0.4}  % Theil slope share %, HAR forecasts, ClosePre: min
\providecommand{\rhcaAllHarShSlopeClosePreMax}{}\renewcommand{\rhcaAllHarShSlopeClosePreMax}{2.5}  % ... max
\providecommand{\rhcaAllHarShSlopeClosePreRng}{}\renewcommand{\rhcaAllHarShSlopeClosePreRng}{0.4--2.5}  % ... the range as text
\providecommand{\rhcaAllShBiasClosePreMin}{}\renewcommand{\rhcaAllShBiasClosePreMin}{0.1}  % Theil bias share %, all 12, ClosePre: min
\providecommand{\rhcaAllShBiasClosePreMax}{}\renewcommand{\rhcaAllShBiasClosePreMax}{0.5}  % ... max
\providecommand{\rhcaAllShBiasClosePreRng}{}\renewcommand{\rhcaAllShBiasClosePreRng}{0.1--0.5}  % ... the range as text
\providecommand{\rhcaAllRmClosePreMin}{}\renewcommand{\rhcaAllRmClosePreMin}{1.04}  % ratio of means, all 12, ClosePre: min
\providecommand{\rhcaAllRmClosePreMax}{}\renewcommand{\rhcaAllRmClosePreMax}{1.10}  % ... max
\providecommand{\rhcaAllRmClosePreRng}{}\renewcommand{\rhcaAllRmClosePreRng}{1.04--1.10}  % ... the range as text
\providecommand{\rhcaFullRmClosePreMin}{}\renewcommand{\rhcaFullRmClosePreMin}{1.04}  % ratio of means, full set, ClosePre: min
\providecommand{\rhcaFullRmClosePreMax}{}\renewcommand{\rhcaFullRmClosePreMax}{1.07}  % ... max
\providecommand{\rhcaFullRmClosePreRng}{}\renewcommand{\rhcaFullRmClosePreRng}{1.04--1.07}  % ... the range as text
\providecommand{\rhcaHarRmClosePreMin}{}\renewcommand{\rhcaHarRmClosePreMin}{1.06}  % ratio of means, HAR, ClosePre: min
\providecommand{\rhcaHarRmClosePreMax}{}\renewcommand{\rhcaHarRmClosePreMax}{1.10}  % ... max
\providecommand{\rhcaHarRmClosePreRng}{}\renewcommand{\rhcaHarRmClosePreRng}{1.06--1.10}  % ... the range as text
\providecommand{\rhcaAllMrClosePreMin}{}\renewcommand{\rhcaAllMrClosePreMin}{1.06}  % mean of RV/sigma2, all 12, ClosePre: min
\providecommand{\rhcaAllMrClosePreMax}{}\renewcommand{\rhcaAllMrClosePreMax}{1.15}  % ... max
\providecommand{\rhcaAllMrClosePreRng}{}\renewcommand{\rhcaAllMrClosePreRng}{1.06--1.15}  % ... the range as text
\providecommand{\rhcaAllUnderClosePreMin}{}\renewcommand{\rhcaAllUnderClosePreMin}{42.6}  % under-prediction %, all 12, ClosePre: min
\providecommand{\rhcaAllUnderClosePreMax}{}\renewcommand{\rhcaAllUnderClosePreMax}{48.6}  % ... max
\providecommand{\rhcaAllUnderClosePreRng}{}\renewcommand{\rhcaAllUnderClosePreRng}{42.6--48.6}  % ... the range as text
\providecommand{\rhcaLinFullUnderClosePreMin}{}\renewcommand{\rhcaLinFullUnderClosePreMin}{47.2}  % under-prediction %, linear full, ClosePre: min
\providecommand{\rhcaLinFullUnderClosePreMax}{}\renewcommand{\rhcaLinFullUnderClosePreMax}{48.6}  % ... max
\providecommand{\rhcaLinFullUnderClosePreRng}{}\renewcommand{\rhcaLinFullUnderClosePreRng}{47.2--48.6}  % ... the range as text
\providecommand{\rhcaLinHarUnderClosePre}{}\renewcommand{\rhcaLinHarUnderClosePre}{46.7}  % under-prediction %, linear HAR, ClosePre
\providecommand{\rhcaLinFullOverSeClosePreMin}{}\renewcommand{\rhcaLinFullOverSeClosePreMin}{17.3}  % % of SE from over-predictions, linear full, ClosePre: min
\providecommand{\rhcaLinFullOverSeClosePreMax}{}\renewcommand{\rhcaLinFullOverSeClosePreMax}{18.4}  % ... max
\providecommand{\rhcaLinFullOverSeClosePreRng}{}\renewcommand{\rhcaLinFullOverSeClosePreRng}{17.3--18.4}  % ... the range as text
\providecommand{\rhcaTreeFullOverSeClosePreMin}{}\renewcommand{\rhcaTreeFullOverSeClosePreMin}{17.3}  % % of SE from over-predictions, tree full, ClosePre: min
\providecommand{\rhcaTreeFullOverSeClosePreMax}{}\renewcommand{\rhcaTreeFullOverSeClosePreMax}{18.9}  % ... max
\providecommand{\rhcaTreeFullOverSeClosePreRng}{}\renewcommand{\rhcaTreeFullOverSeClosePreRng}{17.3--18.9}  % ... the range as text
\providecommand{\rhcaAllHarOverSeClosePreMin}{}\renewcommand{\rhcaAllHarOverSeClosePreMin}{13.1}  % % of SE from over-predictions, HAR, ClosePre: min
\providecommand{\rhcaAllHarOverSeClosePreMax}{}\renewcommand{\rhcaAllHarOverSeClosePreMax}{14.4}  % ... max
\providecommand{\rhcaAllHarOverSeClosePreRng}{}\renewcommand{\rhcaAllHarOverSeClosePreRng}{13.1--14.4}  % ... the range as text
\providecommand{\rhcaNMaeLowerClosePre}{}\renewcommand{\rhcaNMaeLowerClosePre}{6}  % classes (of 6) with the lower MAE on the full set, ClosePre
\providecommand{\rhcaLinFullMrClosePreMin}{}\renewcommand{\rhcaLinFullMrClosePreMin}{1.137}  % mean of RV/sigma2, linear full, ClosePre: min
\providecommand{\rhcaLinFullMrClosePreMax}{}\renewcommand{\rhcaLinFullMrClosePreMax}{1.148}  % ... max
\providecommand{\rhcaLinFullMrClosePreRng}{}\renewcommand{\rhcaLinFullMrClosePreRng}{1.137--1.148}  % ... the range as text
\providecommand{\rhcaLinHarMrClosePre}{}\renewcommand{\rhcaLinHarMrClosePre}{1.121}  % mean of RV/sigma2, linear HAR, ClosePre
\providecommand{\rhcaLinFullQlClosePreMin}{}\renewcommand{\rhcaLinFullQlClosePreMin}{0.0086}  % QLIKE level term, linear full, ClosePre: min
\providecommand{\rhcaLinFullQlClosePreMax}{}\renewcommand{\rhcaLinFullQlClosePreMax}{0.0100}  % ... max
\providecommand{\rhcaLinFullQlClosePreRng}{}\renewcommand{\rhcaLinFullQlClosePreRng}{0.0086--0.0100}  % ... the range as text
\providecommand{\rhcaLinHarQlClosePre}{}\renewcommand{\rhcaLinHarQlClosePre}{0.0068}  % QLIKE level term, linear HAR, ClosePre
\providecommand{\rhcaLinFullQdClosePreMin}{}\renewcommand{\rhcaLinFullQdClosePreMin}{0.1213}  % QLIKE dispersion term, linear full, ClosePre: min
\providecommand{\rhcaLinFullQdClosePreMax}{}\renewcommand{\rhcaLinFullQdClosePreMax}{0.1233}  % ... max
\providecommand{\rhcaLinFullQdClosePreRng}{}\renewcommand{\rhcaLinFullQdClosePreRng}{0.1213--0.1233}  % ... the range as text
\providecommand{\rhcaLinHarQdClosePre}{}\renewcommand{\rhcaLinHarQdClosePre}{0.1297}  % QLIKE dispersion term, linear HAR, ClosePre
\providecommand{\rhcaTreeFullQlClosePreMin}{}\renewcommand{\rhcaTreeFullQlClosePreMin}{0.0045}  % QLIKE level term, tree full, ClosePre: min
\providecommand{\rhcaTreeFullQlClosePreMax}{}\renewcommand{\rhcaTreeFullQlClosePreMax}{0.0065}  % ... max
\providecommand{\rhcaTreeFullQlClosePreRng}{}\renewcommand{\rhcaTreeFullQlClosePreRng}{0.0045--0.0065}  % ... the range as text
\providecommand{\rhcaTreeHarQlClosePreMin}{}\renewcommand{\rhcaTreeHarQlClosePreMin}{0.0020}  % QLIKE level term, tree HAR, ClosePre: min
\providecommand{\rhcaTreeHarQlClosePreMax}{}\renewcommand{\rhcaTreeHarQlClosePreMax}{0.0039}  % ... max
\providecommand{\rhcaTreeHarQlClosePreRng}{}\renewcommand{\rhcaTreeHarQlClosePreRng}{0.0020--0.0039}  % ... the range as text
\providecommand{\rhcaTreeFullQdClosePreMin}{}\renewcommand{\rhcaTreeFullQdClosePreMin}{0.1232}  % QLIKE dispersion term, tree full, ClosePre: min
\providecommand{\rhcaTreeFullQdClosePreMax}{}\renewcommand{\rhcaTreeFullQdClosePreMax}{0.1270}  % ... max
\providecommand{\rhcaTreeFullQdClosePreRng}{}\renewcommand{\rhcaTreeFullQdClosePreRng}{0.1232--0.1270}  % ... the range as text
\providecommand{\rhcaTreeHarQdClosePreMin}{}\renewcommand{\rhcaTreeHarQdClosePreMin}{0.1310}  % QLIKE dispersion term, tree HAR, ClosePre: min
\providecommand{\rhcaTreeHarQdClosePreMax}{}\renewcommand{\rhcaTreeHarQdClosePreMax}{0.1325}  % ... max
\providecommand{\rhcaTreeHarQdClosePreRng}{}\renewcommand{\rhcaTreeHarQdClosePreRng}{0.1310--0.1325}  % ... the range as text
\providecommand{\rhcaTreeFullUnderClosePreMin}{}\renewcommand{\rhcaTreeFullUnderClosePreMin}{45.6}  % under-prediction %, tree full, ClosePre: min
\providecommand{\rhcaTreeFullUnderClosePreMax}{}\renewcommand{\rhcaTreeFullUnderClosePreMax}{46.1}  % ... max
\providecommand{\rhcaTreeFullUnderClosePreRng}{}\renewcommand{\rhcaTreeFullUnderClosePreRng}{45.6--46.1}  % ... the range as text
\providecommand{\rhcaTreeHarUnderClosePreMin}{}\renewcommand{\rhcaTreeHarUnderClosePreMin}{42.6}  % under-prediction %, tree HAR, ClosePre: min
\providecommand{\rhcaTreeHarUnderClosePreMax}{}\renewcommand{\rhcaTreeHarUnderClosePreMax}{43.6}  % ... max
\providecommand{\rhcaTreeHarUnderClosePreRng}{}\renewcommand{\rhcaTreeHarUnderClosePreRng}{42.6--43.6}  % ... the range as text
\providecommand{\rhcaLinFullMzBCloseOosMin}{}\renewcommand{\rhcaLinFullMzBCloseOosMin}{0.66}  % levels MZ slope, linear full, CloseOos: min
\providecommand{\rhcaLinFullMzBCloseOosMax}{}\renewcommand{\rhcaLinFullMzBCloseOosMax}{0.70}  % ... max
\providecommand{\rhcaLinFullMzBCloseOosRng}{}\renewcommand{\rhcaLinFullMzBCloseOosRng}{0.66--0.70}  % ... the range as text
\providecommand{\rhcaTreeFullMzBCloseOosMin}{}\renewcommand{\rhcaTreeFullMzBCloseOosMin}{0.81}  % levels MZ slope, tree full, CloseOos: min
\providecommand{\rhcaTreeFullMzBCloseOosMax}{}\renewcommand{\rhcaTreeFullMzBCloseOosMax}{0.89}  % ... max
\providecommand{\rhcaTreeFullMzBCloseOosRng}{}\renewcommand{\rhcaTreeFullMzBCloseOosRng}{0.81--0.89}  % ... the range as text
\providecommand{\rhcaTreeHarMzBCloseOosMin}{}\renewcommand{\rhcaTreeHarMzBCloseOosMin}{1.06}  % levels MZ slope, tree HAR, CloseOos: min
\providecommand{\rhcaTreeHarMzBCloseOosMax}{}\renewcommand{\rhcaTreeHarMzBCloseOosMax}{1.08}  % ... max
\providecommand{\rhcaTreeHarMzBCloseOosRng}{}\renewcommand{\rhcaTreeHarMzBCloseOosRng}{1.06--1.08}  % ... the range as text
\providecommand{\rhcaLinHarMzBCloseOos}{}\renewcommand{\rhcaLinHarMzBCloseOos}{1.21}  % levels MZ slope, linear HAR (= OLS), CloseOos
\providecommand{\rhcaLinFullMzlBCloseOosMin}{}\renewcommand{\rhcaLinFullMzlBCloseOosMin}{0.92}  % logs MZ slope, linear full, CloseOos: min
\providecommand{\rhcaLinFullMzlBCloseOosMax}{}\renewcommand{\rhcaLinFullMzlBCloseOosMax}{0.92}  % ... max
\providecommand{\rhcaLinFullMzlBCloseOosRng}{}\renewcommand{\rhcaLinFullMzlBCloseOosRng}{0.92}  % ... the range as text
\providecommand{\rhcaTreeFullMzlBCloseOosMin}{}\renewcommand{\rhcaTreeFullMzlBCloseOosMin}{0.93}  % logs MZ slope, tree full, CloseOos: min
\providecommand{\rhcaTreeFullMzlBCloseOosMax}{}\renewcommand{\rhcaTreeFullMzlBCloseOosMax}{0.95}  % ... max
\providecommand{\rhcaTreeFullMzlBCloseOosRng}{}\renewcommand{\rhcaTreeFullMzlBCloseOosRng}{0.93--0.95}  % ... the range as text
\providecommand{\rhcaAllHarMzlBCloseOosMin}{}\renewcommand{\rhcaAllHarMzlBCloseOosMin}{0.97}  % logs MZ slope, HAR forecasts, CloseOos: min
\providecommand{\rhcaAllHarMzlBCloseOosMax}{}\renewcommand{\rhcaAllHarMzlBCloseOosMax}{0.98}  % ... max
\providecommand{\rhcaAllHarMzlBCloseOosRng}{}\renewcommand{\rhcaAllHarMzlBCloseOosRng}{0.97--0.98}  % ... the range as text
\providecommand{\rhcaLinFullMzlRsqCloseOosMin}{}\renewcommand{\rhcaLinFullMzlRsqCloseOosMin}{0.84}  % logs MZ R2, linear full, CloseOos: min
\providecommand{\rhcaLinFullMzlRsqCloseOosMax}{}\renewcommand{\rhcaLinFullMzlRsqCloseOosMax}{0.84}  % ... max
\providecommand{\rhcaLinFullMzlRsqCloseOosRng}{}\renewcommand{\rhcaLinFullMzlRsqCloseOosRng}{0.84}  % ... the range as text
\providecommand{\rhcaTreeFullMzlRsqCloseOosMin}{}\renewcommand{\rhcaTreeFullMzlRsqCloseOosMin}{0.85}  % logs MZ R2, tree full, CloseOos: min
\providecommand{\rhcaTreeFullMzlRsqCloseOosMax}{}\renewcommand{\rhcaTreeFullMzlRsqCloseOosMax}{0.85}  % ... max
\providecommand{\rhcaTreeFullMzlRsqCloseOosRng}{}\renewcommand{\rhcaTreeFullMzlRsqCloseOosRng}{0.85}  % ... the range as text
\providecommand{\rhcaAllHarMzlRsqCloseOosMin}{}\renewcommand{\rhcaAllHarMzlRsqCloseOosMin}{0.83}  % logs MZ R2, HAR forecasts, CloseOos: min
\providecommand{\rhcaAllHarMzlRsqCloseOosMax}{}\renewcommand{\rhcaAllHarMzlRsqCloseOosMax}{0.84}  % ... max
\providecommand{\rhcaAllHarMzlRsqCloseOosRng}{}\renewcommand{\rhcaAllHarMzlRsqCloseOosRng}{0.83--0.84}  % ... the range as text
\providecommand{\rhcaLinFullMzRsqCloseOosMin}{}\renewcommand{\rhcaLinFullMzRsqCloseOosMin}{0.71}  % levels MZ R2, linear full, CloseOos: min
\providecommand{\rhcaLinFullMzRsqCloseOosMax}{}\renewcommand{\rhcaLinFullMzRsqCloseOosMax}{0.72}  % ... max
\providecommand{\rhcaLinFullMzRsqCloseOosRng}{}\renewcommand{\rhcaLinFullMzRsqCloseOosRng}{0.71--0.72}  % ... the range as text
\providecommand{\rhcaAllHarMzRsqCloseOosMin}{}\renewcommand{\rhcaAllHarMzRsqCloseOosMin}{0.80}  % levels MZ R2, HAR forecasts, CloseOos: min
\providecommand{\rhcaAllHarMzRsqCloseOosMax}{}\renewcommand{\rhcaAllHarMzRsqCloseOosMax}{0.82}  % ... max
\providecommand{\rhcaAllHarMzRsqCloseOosRng}{}\renewcommand{\rhcaAllHarMzRsqCloseOosRng}{0.80--0.82}  % ... the range as text
\providecommand{\rhcaLinFullShSlopeCloseOosMin}{}\renewcommand{\rhcaLinFullShSlopeCloseOosMin}{32.0}  % Theil slope share %, linear full, CloseOos: min
\providecommand{\rhcaLinFullShSlopeCloseOosMax}{}\renewcommand{\rhcaLinFullShSlopeCloseOosMax}{39.7}  % ... max
\providecommand{\rhcaLinFullShSlopeCloseOosRng}{}\renewcommand{\rhcaLinFullShSlopeCloseOosRng}{32.0--39.7}  % ... the range as text
\providecommand{\rhcaTreeFullShSlopeCloseOosMin}{}\renewcommand{\rhcaTreeFullShSlopeCloseOosMin}{5.5}  % Theil slope share %, tree full, CloseOos: min
\providecommand{\rhcaTreeFullShSlopeCloseOosMax}{}\renewcommand{\rhcaTreeFullShSlopeCloseOosMax}{14.8}  % ... max
\providecommand{\rhcaTreeFullShSlopeCloseOosRng}{}\renewcommand{\rhcaTreeFullShSlopeCloseOosRng}{5.5--14.8}  % ... the range as text
\providecommand{\rhcaAllHarShSlopeCloseOosMin}{}\renewcommand{\rhcaAllHarShSlopeCloseOosMin}{1.7}  % Theil slope share %, HAR forecasts, CloseOos: min
\providecommand{\rhcaAllHarShSlopeCloseOosMax}{}\renewcommand{\rhcaAllHarShSlopeCloseOosMax}{10.7}  % ... max
\providecommand{\rhcaAllHarShSlopeCloseOosRng}{}\renewcommand{\rhcaAllHarShSlopeCloseOosRng}{1.7--10.7}  % ... the range as text
\providecommand{\rhcaAllShBiasCloseOosMin}{}\renewcommand{\rhcaAllShBiasCloseOosMin}{0.0}  % Theil bias share %, all 12, CloseOos: min
\providecommand{\rhcaAllShBiasCloseOosMax}{}\renewcommand{\rhcaAllShBiasCloseOosMax}{0.3}  % ... max
\providecommand{\rhcaAllShBiasCloseOosRng}{}\renewcommand{\rhcaAllShBiasCloseOosRng}{0.0--0.3}  % ... the range as text
\providecommand{\rhcaAllRmCloseOosMin}{}\renewcommand{\rhcaAllRmCloseOosMin}{0.96}  % ratio of means, all 12, CloseOos: min
\providecommand{\rhcaAllRmCloseOosMax}{}\renewcommand{\rhcaAllRmCloseOosMax}{1.10}  % ... max
\providecommand{\rhcaAllRmCloseOosRng}{}\renewcommand{\rhcaAllRmCloseOosRng}{0.96--1.10}  % ... the range as text
\providecommand{\rhcaFullRmCloseOosMin}{}\renewcommand{\rhcaFullRmCloseOosMin}{0.96}  % ratio of means, full set, CloseOos: min
\providecommand{\rhcaFullRmCloseOosMax}{}\renewcommand{\rhcaFullRmCloseOosMax}{0.98}  % ... max
\providecommand{\rhcaFullRmCloseOosRng}{}\renewcommand{\rhcaFullRmCloseOosRng}{0.96--0.98}  % ... the range as text
\providecommand{\rhcaHarRmCloseOosMin}{}\renewcommand{\rhcaHarRmCloseOosMin}{1.03}  % ratio of means, HAR, CloseOos: min
\providecommand{\rhcaHarRmCloseOosMax}{}\renewcommand{\rhcaHarRmCloseOosMax}{1.10}  % ... max
\providecommand{\rhcaHarRmCloseOosRng}{}\renewcommand{\rhcaHarRmCloseOosRng}{1.03--1.10}  % ... the range as text
\providecommand{\rhcaAllMrCloseOosMin}{}\renewcommand{\rhcaAllMrCloseOosMin}{1.03}  % mean of RV/sigma2, all 12, CloseOos: min
\providecommand{\rhcaAllMrCloseOosMax}{}\renewcommand{\rhcaAllMrCloseOosMax}{1.18}  % ... max
\providecommand{\rhcaAllMrCloseOosRng}{}\renewcommand{\rhcaAllMrCloseOosRng}{1.03--1.18}  % ... the range as text
\providecommand{\rhcaAllUnderCloseOosMin}{}\renewcommand{\rhcaAllUnderCloseOosMin}{42.9}  % under-prediction %, all 12, CloseOos: min
\providecommand{\rhcaAllUnderCloseOosMax}{}\renewcommand{\rhcaAllUnderCloseOosMax}{52.9}  % ... max
\providecommand{\rhcaAllUnderCloseOosRng}{}\renewcommand{\rhcaAllUnderCloseOosRng}{42.9--52.9}  % ... the range as text
\providecommand{\rhcaLinFullUnderCloseOosMin}{}\renewcommand{\rhcaLinFullUnderCloseOosMin}{52.1}  % under-prediction %, linear full, CloseOos: min
\providecommand{\rhcaLinFullUnderCloseOosMax}{}\renewcommand{\rhcaLinFullUnderCloseOosMax}{52.9}  % ... max
\providecommand{\rhcaLinFullUnderCloseOosRng}{}\renewcommand{\rhcaLinFullUnderCloseOosRng}{52.1--52.9}  % ... the range as text
\providecommand{\rhcaLinHarUnderCloseOos}{}\renewcommand{\rhcaLinHarUnderCloseOos}{47.7}  % under-prediction %, linear HAR, CloseOos
\providecommand{\rhcaLinFullOverSeCloseOosMin}{}\renewcommand{\rhcaLinFullOverSeCloseOosMin}{51.5}  % % of SE from over-predictions, linear full, CloseOos: min
\providecommand{\rhcaLinFullOverSeCloseOosMax}{}\renewcommand{\rhcaLinFullOverSeCloseOosMax}{59.7}  % ... max
\providecommand{\rhcaLinFullOverSeCloseOosRng}{}\renewcommand{\rhcaLinFullOverSeCloseOosRng}{51.5--59.7}  % ... the range as text
\providecommand{\rhcaTreeFullOverSeCloseOosMin}{}\renewcommand{\rhcaTreeFullOverSeCloseOosMin}{43.1}  % % of SE from over-predictions, tree full, CloseOos: min
\providecommand{\rhcaTreeFullOverSeCloseOosMax}{}\renewcommand{\rhcaTreeFullOverSeCloseOosMax}{46.1}  % ... max
\providecommand{\rhcaTreeFullOverSeCloseOosRng}{}\renewcommand{\rhcaTreeFullOverSeCloseOosRng}{43.1--46.1}  % ... the range as text
\providecommand{\rhcaAllHarOverSeCloseOosMin}{}\renewcommand{\rhcaAllHarOverSeCloseOosMin}{13.8}  % % of SE from over-predictions, HAR, CloseOos: min
\providecommand{\rhcaAllHarOverSeCloseOosMax}{}\renewcommand{\rhcaAllHarOverSeCloseOosMax}{25.7}  % ... max
\providecommand{\rhcaAllHarOverSeCloseOosRng}{}\renewcommand{\rhcaAllHarOverSeCloseOosRng}{13.8--25.7}  % ... the range as text
\providecommand{\rhcaNMaeLowerCloseOos}{}\renewcommand{\rhcaNMaeLowerCloseOos}{0}  % classes (of 6) with the lower MAE on the full set, CloseOos
\providecommand{\rhcaLinQlShareCloseOosMin}{}\renewcommand{\rhcaLinQlShareCloseOosMin}{64}  % % of the linear full-set QLIKE increase from the level term, CloseOos: min
\providecommand{\rhcaLinQlShareCloseOosMax}{}\renewcommand{\rhcaLinQlShareCloseOosMax}{76}  % ... max
\providecommand{\rhcaLinQlShareCloseOosRng}{}\renewcommand{\rhcaLinQlShareCloseOosRng}{64--76}  % ... the range as text
\providecommand{\rhcaLinFullMrCloseOosMin}{}\renewcommand{\rhcaLinFullMrCloseOosMin}{1.165}  % mean of RV/sigma2, linear full, CloseOos: min
\providecommand{\rhcaLinFullMrCloseOosMax}{}\renewcommand{\rhcaLinFullMrCloseOosMax}{1.183}  % ... max
\providecommand{\rhcaLinFullMrCloseOosRng}{}\renewcommand{\rhcaLinFullMrCloseOosRng}{1.165--1.183}  % ... the range as text
\providecommand{\rhcaLinHarMrCloseOos}{}\renewcommand{\rhcaLinHarMrCloseOos}{1.092}  % mean of RV/sigma2, linear HAR, CloseOos
\providecommand{\rhcaLinFullQlCloseOosMin}{}\renewcommand{\rhcaLinFullQlCloseOosMin}{0.0123}  % QLIKE level term, linear full, CloseOos: min
\providecommand{\rhcaLinFullQlCloseOosMax}{}\renewcommand{\rhcaLinFullQlCloseOosMax}{0.0150}  % ... max
\providecommand{\rhcaLinFullQlCloseOosRng}{}\renewcommand{\rhcaLinFullQlCloseOosRng}{0.0123--0.0150}  % ... the range as text
\providecommand{\rhcaLinHarQlCloseOos}{}\renewcommand{\rhcaLinHarQlCloseOos}{0.0040}  % QLIKE level term, linear HAR, CloseOos
\providecommand{\rhcaLinFullQdCloseOosMin}{}\renewcommand{\rhcaLinFullQdCloseOosMin}{0.1122}  % QLIKE dispersion term, linear full, CloseOos: min
\providecommand{\rhcaLinFullQdCloseOosMax}{}\renewcommand{\rhcaLinFullQdCloseOosMax}{0.1159}  % ... max
\providecommand{\rhcaLinFullQdCloseOosRng}{}\renewcommand{\rhcaLinFullQdCloseOosRng}{0.1122--0.1159}  % ... the range as text
\providecommand{\rhcaLinHarQdCloseOos}{}\renewcommand{\rhcaLinHarQdCloseOos}{0.1097}  % QLIKE dispersion term, linear HAR, CloseOos
\providecommand{\rhcaTreeFullQlCloseOosMin}{}\renewcommand{\rhcaTreeFullQlCloseOosMin}{0.0039}  % QLIKE level term, tree full, CloseOos: min
\providecommand{\rhcaTreeFullQlCloseOosMax}{}\renewcommand{\rhcaTreeFullQlCloseOosMax}{0.0059}  % ... max
\providecommand{\rhcaTreeFullQlCloseOosRng}{}\renewcommand{\rhcaTreeFullQlCloseOosRng}{0.0039--0.0059}  % ... the range as text
\providecommand{\rhcaTreeHarQlCloseOosMin}{}\renewcommand{\rhcaTreeHarQlCloseOosMin}{0.0005}  % QLIKE level term, tree HAR, CloseOos: min
\providecommand{\rhcaTreeHarQlCloseOosMax}{}\renewcommand{\rhcaTreeHarQlCloseOosMax}{0.0019}  % ... max
\providecommand{\rhcaTreeHarQlCloseOosRng}{}\renewcommand{\rhcaTreeHarQlCloseOosRng}{0.0005--0.0019}  % ... the range as text
\providecommand{\rhcaTreeFullQdCloseOosMin}{}\renewcommand{\rhcaTreeFullQdCloseOosMin}{0.1044}  % QLIKE dispersion term, tree full, CloseOos: min
\providecommand{\rhcaTreeFullQdCloseOosMax}{}\renewcommand{\rhcaTreeFullQdCloseOosMax}{0.1064}  % ... max
\providecommand{\rhcaTreeFullQdCloseOosRng}{}\renewcommand{\rhcaTreeFullQdCloseOosRng}{0.1044--0.1064}  % ... the range as text
\providecommand{\rhcaTreeHarQdCloseOosMin}{}\renewcommand{\rhcaTreeHarQdCloseOosMin}{0.1095}  % QLIKE dispersion term, tree HAR, CloseOos: min
\providecommand{\rhcaTreeHarQdCloseOosMax}{}\renewcommand{\rhcaTreeHarQdCloseOosMax}{0.1112}  % ... max
\providecommand{\rhcaTreeHarQdCloseOosRng}{}\renewcommand{\rhcaTreeHarQdCloseOosRng}{0.1095--0.1112}  % ... the range as text
\providecommand{\rhcaTreeFullUnderCloseOosMin}{}\renewcommand{\rhcaTreeFullUnderCloseOosMin}{47.1}  % under-prediction %, tree full, CloseOos: min
\providecommand{\rhcaTreeFullUnderCloseOosMax}{}\renewcommand{\rhcaTreeFullUnderCloseOosMax}{50.1}  % ... max
\providecommand{\rhcaTreeFullUnderCloseOosRng}{}\renewcommand{\rhcaTreeFullUnderCloseOosRng}{47.1--50.1}  % ... the range as text
\providecommand{\rhcaTreeHarUnderCloseOosMin}{}\renewcommand{\rhcaTreeHarUnderCloseOosMin}{42.9}  % under-prediction %, tree HAR, CloseOos: min
\providecommand{\rhcaTreeHarUnderCloseOosMax}{}\renewcommand{\rhcaTreeHarUnderCloseOosMax}{45.5}  % ... max
\providecommand{\rhcaTreeHarUnderCloseOosRng}{}\renewcommand{\rhcaTreeHarUnderCloseOosRng}{42.9--45.5}  % ... the range as text
\providecommand{\rhcaLinTopShareBarOosMin}{}\renewcommand{\rhcaLinTopShareBarOosMin}{92}  % % of the full-minus-HAR squared-error sum (every bar, Oos) on the 5 days of largest difference, linear: min
\providecommand{\rhcaLinTopShareBarOosMax}{}\renewcommand{\rhcaLinTopShareBarOosMax}{94}  % ... max
\providecommand{\rhcaLinTopShareBarOosRng}{}\renewcommand{\rhcaLinTopShareBarOosRng}{92--94}  % ... the range as text
\providecommand{\rhcaTreeTopShareBarOosMin}{}\renewcommand{\rhcaTreeTopShareBarOosMin}{92}  % ... tree ensembles: min
\providecommand{\rhcaTreeTopShareBarOosMax}{}\renewcommand{\rhcaTreeTopShareBarOosMax}{101}  % ... max
\providecommand{\rhcaTreeTopShareBarOosRng}{}\renewcommand{\rhcaTreeTopShareBarOosRng}{92--101}  % ... the range as text
\providecommand{\rhcaTopFirstBarOos}{}\renewcommand{\rhcaTopFirstBarOos}{12 March 2020}  % earliest of those days, any class
\providecommand{\rhcaTopLastBarOos}{}\renewcommand{\rhcaTopLastBarOos}{2 April 2020}  % latest of those days, any class
\providecommand{\rhcaLinMsePctBarOosExCrashMin}{}\renewcommand{\rhcaLinMsePctBarOosExCrashMin}{\ensuremath{-}7.3}  % linear every-bar MSE %, full vs HAR, Oos without the crash window: min
\providecommand{\rhcaLinMsePctBarOosExCrashMax}{}\renewcommand{\rhcaLinMsePctBarOosExCrashMax}{\ensuremath{-}5.2}  % ... max
\providecommand{\rhcaLinMsePctBarOosExCrashRng}{}\renewcommand{\rhcaLinMsePctBarOosExCrashRng}{\ensuremath{-}7.3 to \ensuremath{-}5.2}  % ... the range as text
\providecommand{\rhcaTreeMsePctBarOosExCrashMin}{}\renewcommand{\rhcaTreeMsePctBarOosExCrashMin}{\ensuremath{-}9.0}  % ... tree ensembles: min
\providecommand{\rhcaTreeMsePctBarOosExCrashMax}{}\renewcommand{\rhcaTreeMsePctBarOosExCrashMax}{\ensuremath{-}6.0}  % ... max
\providecommand{\rhcaTreeMsePctBarOosExCrashRng}{}\renewcommand{\rhcaTreeMsePctBarOosExCrashRng}{\ensuremath{-}9.0 to \ensuremath{-}6.0}  % ... the range as text
\providecommand{\rhcaNMseLowerBarOosExCrash}{}\renewcommand{\rhcaNMseLowerBarOosExCrash}{6}  % classes (of 6) with the lower every-bar MSE on the full set, Oos without the crash window
\providecommand{\rhcaNMseLowerSigBarOosExCrash}{}\renewcommand{\rhcaNMseLowerSigBarOosExCrash}{6}  % ... significantly (DM on the daily mean squared error)
\providecommand{\rhcaLinMsePctCloseOosExCrashMin}{}\renewcommand{\rhcaLinMsePctCloseOosExCrashMin}{6.6}  % linear 16:00 MSE %, full vs HAR, Oos without the crash window: min
\providecommand{\rhcaLinMsePctCloseOosExCrashMax}{}\renewcommand{\rhcaLinMsePctCloseOosExCrashMax}{10.0}  % ... max
\providecommand{\rhcaLinMsePctCloseOosExCrashRng}{}\renewcommand{\rhcaLinMsePctCloseOosExCrashRng}{6.6--10.0}  % ... the range as text
\providecommand{\rhcaNMseLowerCloseOosExCrash}{}\renewcommand{\rhcaNMseLowerCloseOosExCrash}{1}  % classes with the lower 16:00 MSE on the full set, Oos without the crash window
\providecommand{\rhcaNMseHigherSigCloseOosExCrash}{}\renewcommand{\rhcaNMseHigherSigCloseOosExCrash}{0}  % ... classes with a significantly higher 16:00 MSE there
\providecommand{\rhcaTreeMsePctCloseOosExCrashMin}{}\renewcommand{\rhcaTreeMsePctCloseOosExCrashMin}{\ensuremath{-}0.7}  % tree 16:00 MSE %, full vs HAR, Oos without the crash window: min
\providecommand{\rhcaTreeMsePctCloseOosExCrashMax}{}\renewcommand{\rhcaTreeMsePctCloseOosExCrashMax}{3.4}  % ... max
\providecommand{\rhcaTreeMsePctCloseOosExCrashRng}{}\renewcommand{\rhcaTreeMsePctCloseOosExCrashRng}{\ensuremath{-}0.7 to 3.4}  % ... the range as text
\providecommand{\rhcaAllMaePctBarOosMin}{}\renewcommand{\rhcaAllMaePctBarOosMin}{2.3}  % every-bar MAE %, full vs HAR, Oos, all six classes: min
\providecommand{\rhcaAllMaePctBarOosMax}{}\renewcommand{\rhcaAllMaePctBarOosMax}{12.8}  % ... max
\providecommand{\rhcaAllMaePctBarOosRng}{}\renewcommand{\rhcaAllMaePctBarOosRng}{2.3--12.8}  % ... the range as text
\providecommand{\rhcaAllMaePctBarOosExCrashMin}{}\renewcommand{\rhcaAllMaePctBarOosExCrashMin}{\ensuremath{-}6.3}  % ... without the crash window: min
\providecommand{\rhcaAllMaePctBarOosExCrashMax}{}\renewcommand{\rhcaAllMaePctBarOosExCrashMax}{\ensuremath{-}4.4}  % ... max
\providecommand{\rhcaAllMaePctBarOosExCrashRng}{}\renewcommand{\rhcaAllMaePctBarOosExCrashRng}{\ensuremath{-}6.3 to \ensuremath{-}4.4}  % ... the range as text
\providecommand{\rhcaNMaeLowerBarOosExCrash}{}\renewcommand{\rhcaNMaeLowerBarOosExCrash}{6}  % classes with the lower every-bar MAE on the full set, Oos without the crash window
\providecommand{\rhcaNMaeLowerSigBarOosExCrash}{}\renewcommand{\rhcaNMaeLowerSigBarOosExCrash}{6}  % ... significantly
\providecommand{\rhcaTopShareRidgeBarOos}{}\renewcommand{\rhcaTopShareRidgeBarOos}{94}  % top-day share, %
\providecommand{\rhcaMsePctRidgeBarOosExCrash}{}\renewcommand{\rhcaMsePctRidgeBarOosExCrash}{\ensuremath{-}5.2}  % every-bar MSE %, ex crash
\providecommand{\rhcaMseDmRidgeBarOosExCrash}{}\renewcommand{\rhcaMseDmRidgeBarOosExCrash}{\ensuremath{-}3.10}
\providecommand{\rhcaMsePRidgeBarOosExCrash}{}\renewcommand{\rhcaMsePRidgeBarOosExCrash}{\ensuremath{p = 0.002}}
\providecommand{\rhcaTopShareLassoBarOos}{}\renewcommand{\rhcaTopShareLassoBarOos}{92}  % top-day share, %
\providecommand{\rhcaMsePctLassoBarOosExCrash}{}\renewcommand{\rhcaMsePctLassoBarOosExCrash}{\ensuremath{-}7.3}  % every-bar MSE %, ex crash
\providecommand{\rhcaMseDmLassoBarOosExCrash}{}\renewcommand{\rhcaMseDmLassoBarOosExCrash}{\ensuremath{-}3.29}
\providecommand{\rhcaMsePLassoBarOosExCrash}{}\renewcommand{\rhcaMsePLassoBarOosExCrash}{\ensuremath{p = 0.001}}
\providecommand{\rhcaTopShareEnetBarOos}{}\renewcommand{\rhcaTopShareEnetBarOos}{94}  % top-day share, %
\providecommand{\rhcaMsePctEnetBarOosExCrash}{}\renewcommand{\rhcaMsePctEnetBarOosExCrash}{\ensuremath{-}6.0}  % every-bar MSE %, ex crash
\providecommand{\rhcaMseDmEnetBarOosExCrash}{}\renewcommand{\rhcaMseDmEnetBarOosExCrash}{\ensuremath{-}3.43}
\providecommand{\rhcaMsePEnetBarOosExCrash}{}\renewcommand{\rhcaMsePEnetBarOosExCrash}{\ensuremath{p < 0.001}}
\providecommand{\rhcaTopShareLgbmBarOos}{}\renewcommand{\rhcaTopShareLgbmBarOos}{101}  % top-day share, %
\providecommand{\rhcaMsePctLgbmBarOosExCrash}{}\renewcommand{\rhcaMsePctLgbmBarOosExCrash}{\ensuremath{-}8.8}  % every-bar MSE %, ex crash
\providecommand{\rhcaMseDmLgbmBarOosExCrash}{}\renewcommand{\rhcaMseDmLgbmBarOosExCrash}{\ensuremath{-}4.30}
\providecommand{\rhcaMsePLgbmBarOosExCrash}{}\renewcommand{\rhcaMsePLgbmBarOosExCrash}{\ensuremath{p < 0.001}}
\providecommand{\rhcaTopShareXgbBarOos}{}\renewcommand{\rhcaTopShareXgbBarOos}{99}  % top-day share, %
\providecommand{\rhcaMsePctXgbBarOosExCrash}{}\renewcommand{\rhcaMsePctXgbBarOosExCrash}{\ensuremath{-}9.0}  % every-bar MSE %, ex crash
\providecommand{\rhcaMseDmXgbBarOosExCrash}{}\renewcommand{\rhcaMseDmXgbBarOosExCrash}{\ensuremath{-}3.67}
\providecommand{\rhcaMsePXgbBarOosExCrash}{}\renewcommand{\rhcaMsePXgbBarOosExCrash}{\ensuremath{p < 0.001}}
\providecommand{\rhcaTopShareRfBarOos}{}\renewcommand{\rhcaTopShareRfBarOos}{92}  % top-day share, %
\providecommand{\rhcaMsePctRfBarOosExCrash}{}\renewcommand{\rhcaMsePctRfBarOosExCrash}{\ensuremath{-}6.0}  % every-bar MSE %, ex crash
\providecommand{\rhcaMseDmRfBarOosExCrash}{}\renewcommand{\rhcaMseDmRfBarOosExCrash}{\ensuremath{-}3.97}
\providecommand{\rhcaMsePRfBarOosExCrash}{}\renewcommand{\rhcaMsePRfBarOosExCrash}{\ensuremath{p < 0.001}}
\providecommand{\rhcaMaxFcRidge}{}\renewcommand{\rhcaMaxFcRidge}{1947}  % largest full-set forecast on a bar, Oos, x 10^6
\providecommand{\rhcaMaxFcRvRidge}{}\renewcommand{\rhcaMaxFcRvRidge}{501}  % the realized variance of that bar, x 10^6
\providecommand{\rhcaMaxFcHarRidge}{}\renewcommand{\rhcaMaxFcHarRidge}{627}  % the HAR forecast of that bar, x 10^6
\providecommand{\rhcaMaxFcRatioRidge}{}\renewcommand{\rhcaMaxFcRatioRidge}{3.9}  % forecast / realized variance on that bar
\providecommand{\rhcaMaxFcLasso}{}\renewcommand{\rhcaMaxFcLasso}{2035}  % largest full-set forecast on a bar, Oos, x 10^6
\providecommand{\rhcaMaxFcRvLasso}{}\renewcommand{\rhcaMaxFcRvLasso}{501}  % the realized variance of that bar, x 10^6
\providecommand{\rhcaMaxFcHarLasso}{}\renewcommand{\rhcaMaxFcHarLasso}{627}  % the HAR forecast of that bar, x 10^6
\providecommand{\rhcaMaxFcRatioLasso}{}\renewcommand{\rhcaMaxFcRatioLasso}{4.1}  % forecast / realized variance on that bar
\providecommand{\rhcaMaxFcEnet}{}\renewcommand{\rhcaMaxFcEnet}{2064}  % largest full-set forecast on a bar, Oos, x 10^6
\providecommand{\rhcaMaxFcRvEnet}{}\renewcommand{\rhcaMaxFcRvEnet}{501}  % the realized variance of that bar, x 10^6
\providecommand{\rhcaMaxFcHarEnet}{}\renewcommand{\rhcaMaxFcHarEnet}{627}  % the HAR forecast of that bar, x 10^6
\providecommand{\rhcaMaxFcRatioEnet}{}\renewcommand{\rhcaMaxFcRatioEnet}{4.1}  % forecast / realized variance on that bar
\providecommand{\rhcaMaxFcDate}{}\renewcommand{\rhcaMaxFcDate}{23 March 2020}  % the bar of the largest linear full-set forecast (all three)
\providecommand{\rhcaMaxFcClock}{}\renewcommand{\rhcaMaxFcClock}{10:00}  % its bar end, ET
\providecommand{\rhcaMaxRvOos}{}\renewcommand{\rhcaMaxRvOos}{1150}  % largest bar RV on Oos, x 10^6
\providecommand{\rhcaYrsUpRidgeCloseOos}{}\renewcommand{\rhcaYrsUpRidgeCloseOos}{4}  % calendar years of Oos (of 5) with a higher 16:00 QLIKE on the full set
\providecommand{\rhcaYrsDownRidgeClosePre}{}\renewcommand{\rhcaYrsDownRidgeClosePre}{4}  % calendar years of Pre (of 6) with a lower 16:00 QLIKE
\providecommand{\rhcaWorstYrRidgeCloseOos}{}\renewcommand{\rhcaWorstYrRidgeCloseOos}{2023}  % the Oos year of the largest full-set increase
\providecommand{\rhcaWorstYrPctRidgeCloseOos}{}\renewcommand{\rhcaWorstYrPctRidgeCloseOos}{32.0}  % ... its QLIKE %
\providecommand{\rhcaYrsUpLassoCloseOos}{}\renewcommand{\rhcaYrsUpLassoCloseOos}{4}  % calendar years of Oos (of 5) with a higher 16:00 QLIKE on the full set
\providecommand{\rhcaYrsDownLassoClosePre}{}\renewcommand{\rhcaYrsDownLassoClosePre}{3}  % calendar years of Pre (of 6) with a lower 16:00 QLIKE
\providecommand{\rhcaWorstYrLassoCloseOos}{}\renewcommand{\rhcaWorstYrLassoCloseOos}{2023}  % the Oos year of the largest full-set increase
\providecommand{\rhcaWorstYrPctLassoCloseOos}{}\renewcommand{\rhcaWorstYrPctLassoCloseOos}{33.0}  % ... its QLIKE %
\providecommand{\rhcaYrsUpEnetCloseOos}{}\renewcommand{\rhcaYrsUpEnetCloseOos}{4}  % calendar years of Oos (of 5) with a higher 16:00 QLIKE on the full set
\providecommand{\rhcaYrsDownEnetClosePre}{}\renewcommand{\rhcaYrsDownEnetClosePre}{5}  % calendar years of Pre (of 6) with a lower 16:00 QLIKE
\providecommand{\rhcaWorstYrEnetCloseOos}{}\renewcommand{\rhcaWorstYrEnetCloseOos}{2023}  % the Oos year of the largest full-set increase
\providecommand{\rhcaWorstYrPctEnetCloseOos}{}\renewcommand{\rhcaWorstYrPctEnetCloseOos}{30.7}  % ... its QLIKE %
\providecommand{\rhcaYrsUpLgbmCloseOos}{}\renewcommand{\rhcaYrsUpLgbmCloseOos}{2}  % calendar years of Oos (of 5) with a higher 16:00 QLIKE on the full set
\providecommand{\rhcaYrsDownLgbmClosePre}{}\renewcommand{\rhcaYrsDownLgbmClosePre}{4}  % calendar years of Pre (of 6) with a lower 16:00 QLIKE
\providecommand{\rhcaWorstYrLgbmCloseOos}{}\renewcommand{\rhcaWorstYrLgbmCloseOos}{2024}  % the Oos year of the largest full-set increase
\providecommand{\rhcaWorstYrPctLgbmCloseOos}{}\renewcommand{\rhcaWorstYrPctLgbmCloseOos}{10.5}  % ... its QLIKE %
\providecommand{\rhcaYrsUpXgbCloseOos}{}\renewcommand{\rhcaYrsUpXgbCloseOos}{2}  % calendar years of Oos (of 5) with a higher 16:00 QLIKE on the full set
\providecommand{\rhcaYrsDownXgbClosePre}{}\renewcommand{\rhcaYrsDownXgbClosePre}{4}  % calendar years of Pre (of 6) with a lower 16:00 QLIKE
\providecommand{\rhcaWorstYrXgbCloseOos}{}\renewcommand{\rhcaWorstYrXgbCloseOos}{2024}  % the Oos year of the largest full-set increase
\providecommand{\rhcaWorstYrPctXgbCloseOos}{}\renewcommand{\rhcaWorstYrPctXgbCloseOos}{11.9}  % ... its QLIKE %
\providecommand{\rhcaYrsUpRfCloseOos}{}\renewcommand{\rhcaYrsUpRfCloseOos}{2}  % calendar years of Oos (of 5) with a higher 16:00 QLIKE on the full set
\providecommand{\rhcaYrsDownRfClosePre}{}\renewcommand{\rhcaYrsDownRfClosePre}{4}  % calendar years of Pre (of 6) with a lower 16:00 QLIKE
\providecommand{\rhcaWorstYrRfCloseOos}{}\renewcommand{\rhcaWorstYrRfCloseOos}{2021}  % the Oos year of the largest full-set increase
\providecommand{\rhcaWorstYrPctRfCloseOos}{}\renewcommand{\rhcaWorstYrPctRfCloseOos}{7.3}  % ... its QLIKE %
\providecommand{\rhcaNYrsOos}{}\renewcommand{\rhcaNYrsOos}{5}  % calendar years of Oos
\providecommand{\rhcaNYrsPre}{}\renewcommand{\rhcaNYrsPre}{6}  % calendar years of Pre
\providecommand{\rhcaLinMedLogRatioClosePreMin}{}\renewcommand{\rhcaLinMedLogRatioClosePreMin}{\ensuremath{-}3.1}  % median 100 ln(full / HAR forecast), linear, 16:00, Pre: min
\providecommand{\rhcaLinMedLogRatioClosePreMax}{}\renewcommand{\rhcaLinMedLogRatioClosePreMax}{\ensuremath{-}2.1}  % ... max
\providecommand{\rhcaLinMedLogRatioClosePreRng}{}\renewcommand{\rhcaLinMedLogRatioClosePreRng}{\ensuremath{-}3.1 to \ensuremath{-}2.1}  % ... the range as text
\providecommand{\rhcaTreeMedLogRatioClosePreMin}{}\renewcommand{\rhcaTreeMedLogRatioClosePreMin}{\ensuremath{-}5.1}  % median 100 ln(full / HAR forecast), trees, 16:00, Pre: min
\providecommand{\rhcaTreeMedLogRatioClosePreMax}{}\renewcommand{\rhcaTreeMedLogRatioClosePreMax}{\ensuremath{-}3.6}  % ... max
\providecommand{\rhcaTreeMedLogRatioClosePreRng}{}\renewcommand{\rhcaTreeMedLogRatioClosePreRng}{\ensuremath{-}5.1 to \ensuremath{-}3.6}  % ... the range as text
\providecommand{\rhcaLinMedLogRatioCloseOosMin}{}\renewcommand{\rhcaLinMedLogRatioCloseOosMin}{\ensuremath{-}6.9}  % median 100 ln(full / HAR forecast), linear, 16:00, Oos: min
\providecommand{\rhcaLinMedLogRatioCloseOosMax}{}\renewcommand{\rhcaLinMedLogRatioCloseOosMax}{\ensuremath{-}5.6}  % ... max
\providecommand{\rhcaLinMedLogRatioCloseOosRng}{}\renewcommand{\rhcaLinMedLogRatioCloseOosRng}{\ensuremath{-}6.9 to \ensuremath{-}5.6}  % ... the range as text
\providecommand{\rhcaTreeMedLogRatioCloseOosMin}{}\renewcommand{\rhcaTreeMedLogRatioCloseOosMin}{\ensuremath{-}8.6}  % median 100 ln(full / HAR forecast), trees, 16:00, Oos: min
\providecommand{\rhcaTreeMedLogRatioCloseOosMax}{}\renewcommand{\rhcaTreeMedLogRatioCloseOosMax}{\ensuremath{-}6.1}  % ... max
\providecommand{\rhcaTreeMedLogRatioCloseOosRng}{}\renewcommand{\rhcaTreeMedLogRatioCloseOosRng}{\ensuremath{-}8.6 to \ensuremath{-}6.1}  % ... the range as text
\providecommand{\rhcaTopDecRatioEnetHarBarPre}{}\renewcommand{\rhcaTopDecRatioEnetHarBarPre}{1.04}  % top forecast decile: mean RV / mean forecast
\providecommand{\rhcaTopDecSeEnetHarBarPre}{}\renewcommand{\rhcaTopDecSeEnetHarBarPre}{85}  % top forecast decile: % of the squared error
\providecommand{\rhcaBotDecRatioEnetHarBarPre}{}\renewcommand{\rhcaBotDecRatioEnetHarBarPre}{1.12}  % bottom forecast decile: mean RV / mean forecast
\providecommand{\rhcaTopDecRatioEnetFullBarPre}{}\renewcommand{\rhcaTopDecRatioEnetFullBarPre}{1.01}  % top forecast decile: mean RV / mean forecast
\providecommand{\rhcaTopDecSeEnetFullBarPre}{}\renewcommand{\rhcaTopDecSeEnetFullBarPre}{84}  % top forecast decile: % of the squared error
\providecommand{\rhcaBotDecRatioEnetFullBarPre}{}\renewcommand{\rhcaBotDecRatioEnetFullBarPre}{1.16}  % bottom forecast decile: mean RV / mean forecast
\providecommand{\rhcaTopDecRatioLgbmHarBarPre}{}\renewcommand{\rhcaTopDecRatioLgbmHarBarPre}{1.08}  % top forecast decile: mean RV / mean forecast
\providecommand{\rhcaTopDecSeLgbmHarBarPre}{}\renewcommand{\rhcaTopDecSeLgbmHarBarPre}{88}  % top forecast decile: % of the squared error
\providecommand{\rhcaBotDecRatioLgbmHarBarPre}{}\renewcommand{\rhcaBotDecRatioLgbmHarBarPre}{1.14}  % bottom forecast decile: mean RV / mean forecast
\providecommand{\rhcaTopDecRatioLgbmFullBarPre}{}\renewcommand{\rhcaTopDecRatioLgbmFullBarPre}{1.06}  % top forecast decile: mean RV / mean forecast
\providecommand{\rhcaTopDecSeLgbmFullBarPre}{}\renewcommand{\rhcaTopDecSeLgbmFullBarPre}{91}  % top forecast decile: % of the squared error
\providecommand{\rhcaBotDecRatioLgbmFullBarPre}{}\renewcommand{\rhcaBotDecRatioLgbmFullBarPre}{1.09}  % bottom forecast decile: mean RV / mean forecast
\providecommand{\rhcaTopDecRatioEnetHarBarOos}{}\renewcommand{\rhcaTopDecRatioEnetHarBarOos}{1.03}  % top forecast decile: mean RV / mean forecast
\providecommand{\rhcaTopDecSeEnetHarBarOos}{}\renewcommand{\rhcaTopDecSeEnetHarBarOos}{91}  % top forecast decile: % of the squared error
\providecommand{\rhcaBotDecRatioEnetHarBarOos}{}\renewcommand{\rhcaBotDecRatioEnetHarBarOos}{1.06}  % bottom forecast decile: mean RV / mean forecast
\providecommand{\rhcaTopDecRatioEnetFullBarOos}{}\renewcommand{\rhcaTopDecRatioEnetFullBarOos}{0.88}  % top forecast decile: mean RV / mean forecast
\providecommand{\rhcaTopDecSeEnetFullBarOos}{}\renewcommand{\rhcaTopDecSeEnetFullBarOos}{97}  % top forecast decile: % of the squared error
\providecommand{\rhcaBotDecRatioEnetFullBarOos}{}\renewcommand{\rhcaBotDecRatioEnetFullBarOos}{1.14}  % bottom forecast decile: mean RV / mean forecast
\providecommand{\rhcaTopDecRatioLgbmHarBarOos}{}\renewcommand{\rhcaTopDecRatioLgbmHarBarOos}{1.00}  % top forecast decile: mean RV / mean forecast
\providecommand{\rhcaTopDecSeLgbmHarBarOos}{}\renewcommand{\rhcaTopDecSeLgbmHarBarOos}{92}  % top forecast decile: % of the squared error
\providecommand{\rhcaBotDecRatioLgbmHarBarOos}{}\renewcommand{\rhcaBotDecRatioLgbmHarBarOos}{1.07}  % bottom forecast decile: mean RV / mean forecast
\providecommand{\rhcaTopDecRatioLgbmFullBarOos}{}\renewcommand{\rhcaTopDecRatioLgbmFullBarOos}{0.89}  % top forecast decile: mean RV / mean forecast
\providecommand{\rhcaTopDecSeLgbmFullBarOos}{}\renewcommand{\rhcaTopDecSeLgbmFullBarOos}{95}  % top forecast decile: % of the squared error
\providecommand{\rhcaBotDecRatioLgbmFullBarOos}{}\renewcommand{\rhcaBotDecRatioLgbmFullBarOos}{1.12}  % bottom forecast decile: mean RV / mean forecast
\providecommand{\rhcaTopDecRatioEnetHarClosePre}{}\renewcommand{\rhcaTopDecRatioEnetHarClosePre}{1.11}  % top forecast decile: mean RV / mean forecast
\providecommand{\rhcaTopDecSeEnetHarClosePre}{}\renewcommand{\rhcaTopDecSeEnetHarClosePre}{89}  % top forecast decile: % of the squared error
\providecommand{\rhcaBotDecRatioEnetHarClosePre}{}\renewcommand{\rhcaBotDecRatioEnetHarClosePre}{1.31}  % bottom forecast decile: mean RV / mean forecast
\providecommand{\rhcaTopDecRatioEnetFullClosePre}{}\renewcommand{\rhcaTopDecRatioEnetFullClosePre}{1.03}  % top forecast decile: mean RV / mean forecast
\providecommand{\rhcaTopDecSeEnetFullClosePre}{}\renewcommand{\rhcaTopDecSeEnetFullClosePre}{88}  % top forecast decile: % of the squared error
\providecommand{\rhcaBotDecRatioEnetFullClosePre}{}\renewcommand{\rhcaBotDecRatioEnetFullClosePre}{1.31}  % bottom forecast decile: mean RV / mean forecast
\providecommand{\rhcaTopDecRatioLgbmHarClosePre}{}\renewcommand{\rhcaTopDecRatioLgbmHarClosePre}{1.08}  % top forecast decile: mean RV / mean forecast
\providecommand{\rhcaTopDecSeLgbmHarClosePre}{}\renewcommand{\rhcaTopDecSeLgbmHarClosePre}{91}  % top forecast decile: % of the squared error
\providecommand{\rhcaBotDecRatioLgbmHarClosePre}{}\renewcommand{\rhcaBotDecRatioLgbmHarClosePre}{1.19}  % bottom forecast decile: mean RV / mean forecast
\providecommand{\rhcaTopDecRatioLgbmFullClosePre}{}\renewcommand{\rhcaTopDecRatioLgbmFullClosePre}{1.01}  % top forecast decile: mean RV / mean forecast
\providecommand{\rhcaTopDecSeLgbmFullClosePre}{}\renewcommand{\rhcaTopDecSeLgbmFullClosePre}{91}  % top forecast decile: % of the squared error
\providecommand{\rhcaBotDecRatioLgbmFullClosePre}{}\renewcommand{\rhcaBotDecRatioLgbmFullClosePre}{1.23}  % bottom forecast decile: mean RV / mean forecast
\providecommand{\rhcaTopDecRatioEnetHarCloseOos}{}\renewcommand{\rhcaTopDecRatioEnetHarCloseOos}{1.13}  % top forecast decile: mean RV / mean forecast
\providecommand{\rhcaTopDecSeEnetHarCloseOos}{}\renewcommand{\rhcaTopDecSeEnetHarCloseOos}{98}  % top forecast decile: % of the squared error
\providecommand{\rhcaBotDecRatioEnetHarCloseOos}{}\renewcommand{\rhcaBotDecRatioEnetHarCloseOos}{1.15}  % bottom forecast decile: mean RV / mean forecast
\providecommand{\rhcaTopDecRatioEnetFullCloseOos}{}\renewcommand{\rhcaTopDecRatioEnetFullCloseOos}{0.87}  % top forecast decile: mean RV / mean forecast
\providecommand{\rhcaTopDecSeEnetFullCloseOos}{}\renewcommand{\rhcaTopDecSeEnetFullCloseOos}{99}  % top forecast decile: % of the squared error
\providecommand{\rhcaBotDecRatioEnetFullCloseOos}{}\renewcommand{\rhcaBotDecRatioEnetFullCloseOos}{1.31}  % bottom forecast decile: mean RV / mean forecast
\providecommand{\rhcaTopDecRatioLgbmHarCloseOos}{}\renewcommand{\rhcaTopDecRatioLgbmHarCloseOos}{1.02}  % top forecast decile: mean RV / mean forecast
\providecommand{\rhcaTopDecSeLgbmHarCloseOos}{}\renewcommand{\rhcaTopDecSeLgbmHarCloseOos}{97}  % top forecast decile: % of the squared error
\providecommand{\rhcaBotDecRatioLgbmHarCloseOos}{}\renewcommand{\rhcaBotDecRatioLgbmHarCloseOos}{1.11}  % bottom forecast decile: mean RV / mean forecast
\providecommand{\rhcaTopDecRatioLgbmFullCloseOos}{}\renewcommand{\rhcaTopDecRatioLgbmFullCloseOos}{0.91}  % top forecast decile: mean RV / mean forecast
\providecommand{\rhcaTopDecSeLgbmFullCloseOos}{}\renewcommand{\rhcaTopDecSeLgbmFullCloseOos}{98}  % top forecast decile: % of the squared error
\providecommand{\rhcaBotDecRatioLgbmFullCloseOos}{}\renewcommand{\rhcaBotDecRatioLgbmFullCloseOos}{1.22}  % bottom forecast decile: mean RV / mean forecast
\providecommand{\rhcaHarLinGap}{}\renewcommand{\rhcaHarLinGap}{5 \times 10^{-6}}  % max relative gap, HAR linear vs OLS forecasts (math mode)
\fi
\ifdefined\rhmcsBlock\else% AUTO-GENERATED by writeup/make_rehaul_mcs.py -- do not edit; re-run the script.
% \rhmcs... macros: model confidence set (Hansen, Lunde and Nason 2011), SPA (Hansen 2005) and StepM (Romano and
% Wolf 2005) on the 12 pooled forecasts of the final rerun (the three linear HAR forecasts as one) and the
% closing-interval HAR benchmark. Samples: Bar (every bar, 2014-2024, daily-averaged losses), All (16:00, 2014-2024),
% Oos (16:00, T_oos), Ch (16:00, the benchmark's span). Losses: Qlike, Mse. MCS statistics: Tmax, Trange (T_R).
% SPA / StepM families: <Sample><Q|M>Af (Full vs own HAR), <Sample><Q|M>Bm (Full vs the closing-interval HAR),
% <Mid|Ba><Short|Twin><Mean|Sr> (sign strategies on the 614 post-warmup days, vs the short straddle or the HAR twin,
% mean return or the paper's pooled Sharpe ratio, D19). p-values of StepM and SPA are one-sided (H_k: no better than the benchmark).
\providecommand{\rhmcsB}{}\renewcommand{\rhmcsB}{4{,}999}  % bootstrap replications, every MCS / SPA / StepM (accuracy and strategies)
\providecommand{\rhmcsBlock}{}\renewcommand{\rhmcsBlock}{10}  % block length in trading days (circular moving-block bootstrap), main
\providecommand{\rhmcsBlockLo}{}\renewcommand{\rhmcsBlockLo}{5}  % block length, sensitivity (shorter)
\providecommand{\rhmcsBlockHi}{}\renewcommand{\rhmcsBlockHi}{20}  % block length, sensitivity (longer)
\providecommand{\rhmcsSeed}{}\renewcommand{\rhmcsSeed}{20261025}  % random seed of the accuracy bootstrap (MCS, SPA, StepM on losses)
\providecommand{\rhmcsSeedEcon}{}\renewcommand{\rhmcsSeedEcon}{20261022}  % random seed of the strategy bootstrap (= the paper's D6 bootstrap seed)
\providecommand{\rhmcsNModels}{}\renewcommand{\rhmcsNModels}{10}  % distinct pooled forecasts in the MCS (the three linear HAR forecasts as one)
\providecommand{\rhmcsNModelsWord}{}\renewcommand{\rhmcsNModelsWord}{ten}  % rhmcsNModels as a word
\providecommand{\rhmcsNModelsBench}{}\renewcommand{\rhmcsNModelsBench}{11}  % ... with the closing-interval HAR benchmark
\providecommand{\rhmcsNModelsBenchWord}{}\renewcommand{\rhmcsNModelsBenchWord}{eleven}  % rhmcsNModelsBench as a word
\providecommand{\rhmcsNPool}{}\renewcommand{\rhmcsNPool}{12}  % pooled forecasts (6 classes x HAR, Full)
\providecommand{\rhmcsNAlt}{}\renewcommand{\rhmcsNAlt}{6}  % alternatives in every SPA / StepM family (the six full-set forecasts or strategies)
\providecommand{\rhmcsNAltWord}{}\renewcommand{\rhmcsNAltWord}{six}  % rhmcsNAlt as a word
\providecommand{\rhmcsLinGap}{}\renewcommand{\rhmcsLinGap}{5 \times 10^{-6}}  % largest relative gap between the daily every-bar QLIKE of a linear HAR forecast and the pooled OLS (math mode)
\providecommand{\rhmcsDaysAll}{}\renewcommand{\rhmcsDaysAll}{2{,}579}  % trading days 2014-2024 (Table 1's span)
\providecommand{\rhmcsDaysOos}{}\renewcommand{\rhmcsDaysOos}{1{,}082}  % trading days of T_oos
\providecommand{\rhmcsDaysCh}{}\renewcommand{\rhmcsDaysCh}{2{,}579}  % trading days of the closing-interval benchmark's span within 2014-2024
\providecommand{\rhmcsChStart}{}\renewcommand{\rhmcsChStart}{2 January 2014}  % first day of the closing-interval HAR benchmark's forecasts
\providecommand{\rhmcsNBars}{}\renewcommand{\rhmcsNBars}{33{,}527}  % bars of the every-bar sample (13 a day)
\providecommand{\rhmcsNPw}{}\renewcommand{\rhmcsNPw}{614}  % post-warmup option days of the strategy tests
\providecommand{\rhmcsLinSameTrades}{}\renewcommand{\rhmcsLinSameTrades}{yes}  % the three linear HAR sign strategies have identical daily returns (mid and bid/ask)
\providecommand{\rhmcsBarQlikeTmaxPLinHar}{}\renewcommand{\rhmcsBarQlikeTmaxPLinHar}{\ensuremath{<}0.001}  % MCS p-value, every bar, 2014--2024, Qlike, Tmax: Linear, HAR
\providecommand{\rhmcsBarQlikeTmaxPRidgeFull}{}\renewcommand{\rhmcsBarQlikeTmaxPRidgeFull}{\ensuremath{<}0.001}  % MCS p-value, every bar, 2014--2024, Qlike, Tmax: Ridge, Full
\providecommand{\rhmcsBarQlikeTmaxPLassoFull}{}\renewcommand{\rhmcsBarQlikeTmaxPLassoFull}{\ensuremath{<}0.001}  % MCS p-value, every bar, 2014--2024, Qlike, Tmax: Lasso, Full
\providecommand{\rhmcsBarQlikeTmaxPEnetFull}{}\renewcommand{\rhmcsBarQlikeTmaxPEnetFull}{\ensuremath{<}0.001}  % MCS p-value, every bar, 2014--2024, Qlike, Tmax: Elastic net, Full
\providecommand{\rhmcsBarQlikeTmaxPLgbmHar}{}\renewcommand{\rhmcsBarQlikeTmaxPLgbmHar}{\ensuremath{<}0.001}  % MCS p-value, every bar, 2014--2024, Qlike, Tmax: LightGBM, HAR
\providecommand{\rhmcsBarQlikeTmaxPLgbmFull}{}\renewcommand{\rhmcsBarQlikeTmaxPLgbmFull}{1.00}  % MCS p-value, every bar, 2014--2024, Qlike, Tmax: LightGBM, Full
\providecommand{\rhmcsBarQlikeTmaxPXgbHar}{}\renewcommand{\rhmcsBarQlikeTmaxPXgbHar}{\ensuremath{<}0.001}  % MCS p-value, every bar, 2014--2024, Qlike, Tmax: XGBoost, HAR
\providecommand{\rhmcsBarQlikeTmaxPXgbFull}{}\renewcommand{\rhmcsBarQlikeTmaxPXgbFull}{0.69}  % MCS p-value, every bar, 2014--2024, Qlike, Tmax: XGBoost, Full
\providecommand{\rhmcsBarQlikeTmaxPRfHar}{}\renewcommand{\rhmcsBarQlikeTmaxPRfHar}{\ensuremath{<}0.001}  % MCS p-value, every bar, 2014--2024, Qlike, Tmax: Random forest, HAR
\providecommand{\rhmcsBarQlikeTmaxPRfFull}{}\renewcommand{\rhmcsBarQlikeTmaxPRfFull}{\ensuremath{<}0.001}  % MCS p-value, every bar, 2014--2024, Qlike, Tmax: Random forest, Full
\providecommand{\rhmcsBarQlikeTmaxNNinety}{}\renewcommand{\rhmcsBarQlikeTmaxNNinety}{2}  % models in the 90 % MCS (of 10)
\providecommand{\rhmcsBarQlikeTmaxSetNinety}{}\renewcommand{\rhmcsBarQlikeTmaxSetNinety}{LightGBM on the full predictor set and XGBoost on the full predictor set}  % the 90 % MCS, best (last eliminated) first
\providecommand{\rhmcsBarQlikeTmaxStableNinety}{}\renewcommand{\rhmcsBarQlikeTmaxStableNinety}{yes}  % the 90 % MCS is the same at block lengths (5, 20)
\providecommand{\rhmcsBarQlikeTmaxPMaxOutNinety}{}\renewcommand{\rhmcsBarQlikeTmaxPMaxOutNinety}{\ensuremath{<}0.001}  % largest MCS p-value outside the 90 % set
\providecommand{\rhmcsBarQlikeTmaxNNinetyLo}{}\renewcommand{\rhmcsBarQlikeTmaxNNinetyLo}{2}  % ... at block length 5
\providecommand{\rhmcsBarQlikeTmaxSetNinetyLo}{}\renewcommand{\rhmcsBarQlikeTmaxSetNinetyLo}{LightGBM on the full predictor set and XGBoost on the full predictor set}  % ... at block length 5
\providecommand{\rhmcsBarQlikeTmaxNNinetyHi}{}\renewcommand{\rhmcsBarQlikeTmaxNNinetyHi}{2}  % ... at block length 20
\providecommand{\rhmcsBarQlikeTmaxSetNinetyHi}{}\renewcommand{\rhmcsBarQlikeTmaxSetNinetyHi}{LightGBM on the full predictor set and XGBoost on the full predictor set}  % ... at block length 20
\providecommand{\rhmcsBarQlikeTmaxNSevenFive}{}\renewcommand{\rhmcsBarQlikeTmaxNSevenFive}{2}  % models in the 75 % MCS (of 10)
\providecommand{\rhmcsBarQlikeTmaxSetSevenFive}{}\renewcommand{\rhmcsBarQlikeTmaxSetSevenFive}{LightGBM on the full predictor set and XGBoost on the full predictor set}  % the 75 % MCS, best (last eliminated) first
\providecommand{\rhmcsBarQlikeTmaxStableSevenFive}{}\renewcommand{\rhmcsBarQlikeTmaxStableSevenFive}{yes}  % the 75 % MCS is the same at block lengths (5, 20)
\providecommand{\rhmcsBarQlikeTmaxPMaxOutSevenFive}{}\renewcommand{\rhmcsBarQlikeTmaxPMaxOutSevenFive}{\ensuremath{<}0.001}  % largest MCS p-value outside the 75 % set
\providecommand{\rhmcsBarQlikeTmaxNSevenFiveLo}{}\renewcommand{\rhmcsBarQlikeTmaxNSevenFiveLo}{2}  % ... at block length 5
\providecommand{\rhmcsBarQlikeTmaxSetSevenFiveLo}{}\renewcommand{\rhmcsBarQlikeTmaxSetSevenFiveLo}{LightGBM on the full predictor set and XGBoost on the full predictor set}  % ... at block length 5
\providecommand{\rhmcsBarQlikeTmaxNSevenFiveHi}{}\renewcommand{\rhmcsBarQlikeTmaxNSevenFiveHi}{2}  % ... at block length 20
\providecommand{\rhmcsBarQlikeTmaxSetSevenFiveHi}{}\renewcommand{\rhmcsBarQlikeTmaxSetSevenFiveHi}{LightGBM on the full predictor set and XGBoost on the full predictor set}  % ... at block length 20
\providecommand{\rhmcsBarQlikeTmaxBest}{}\renewcommand{\rhmcsBarQlikeTmaxBest}{LightGBM on the full predictor set}  % the last model left (MCS p-value 1)
\providecommand{\rhmcsBarQlikeTrangePLinHar}{}\renewcommand{\rhmcsBarQlikeTrangePLinHar}{\ensuremath{<}0.001}  % MCS p-value, every bar, 2014--2024, Qlike, Trange: Linear, HAR
\providecommand{\rhmcsBarQlikeTrangePRidgeFull}{}\renewcommand{\rhmcsBarQlikeTrangePRidgeFull}{\ensuremath{<}0.001}  % MCS p-value, every bar, 2014--2024, Qlike, Trange: Ridge, Full
\providecommand{\rhmcsBarQlikeTrangePLassoFull}{}\renewcommand{\rhmcsBarQlikeTrangePLassoFull}{\ensuremath{<}0.001}  % MCS p-value, every bar, 2014--2024, Qlike, Trange: Lasso, Full
\providecommand{\rhmcsBarQlikeTrangePEnetFull}{}\renewcommand{\rhmcsBarQlikeTrangePEnetFull}{\ensuremath{<}0.001}  % MCS p-value, every bar, 2014--2024, Qlike, Trange: Elastic net, Full
\providecommand{\rhmcsBarQlikeTrangePLgbmHar}{}\renewcommand{\rhmcsBarQlikeTrangePLgbmHar}{\ensuremath{<}0.001}  % MCS p-value, every bar, 2014--2024, Qlike, Trange: LightGBM, HAR
\providecommand{\rhmcsBarQlikeTrangePLgbmFull}{}\renewcommand{\rhmcsBarQlikeTrangePLgbmFull}{1.00}  % MCS p-value, every bar, 2014--2024, Qlike, Trange: LightGBM, Full
\providecommand{\rhmcsBarQlikeTrangePXgbHar}{}\renewcommand{\rhmcsBarQlikeTrangePXgbHar}{\ensuremath{<}0.001}  % MCS p-value, every bar, 2014--2024, Qlike, Trange: XGBoost, HAR
\providecommand{\rhmcsBarQlikeTrangePXgbFull}{}\renewcommand{\rhmcsBarQlikeTrangePXgbFull}{0.69}  % MCS p-value, every bar, 2014--2024, Qlike, Trange: XGBoost, Full
\providecommand{\rhmcsBarQlikeTrangePRfHar}{}\renewcommand{\rhmcsBarQlikeTrangePRfHar}{\ensuremath{<}0.001}  % MCS p-value, every bar, 2014--2024, Qlike, Trange: Random forest, HAR
\providecommand{\rhmcsBarQlikeTrangePRfFull}{}\renewcommand{\rhmcsBarQlikeTrangePRfFull}{\ensuremath{<}0.001}  % MCS p-value, every bar, 2014--2024, Qlike, Trange: Random forest, Full
\providecommand{\rhmcsBarQlikeTrangeNNinety}{}\renewcommand{\rhmcsBarQlikeTrangeNNinety}{2}  % models in the 90 % MCS (of 10)
\providecommand{\rhmcsBarQlikeTrangeSetNinety}{}\renewcommand{\rhmcsBarQlikeTrangeSetNinety}{LightGBM on the full predictor set and XGBoost on the full predictor set}  % the 90 % MCS, best (last eliminated) first
\providecommand{\rhmcsBarQlikeTrangeStableNinety}{}\renewcommand{\rhmcsBarQlikeTrangeStableNinety}{yes}  % the 90 % MCS is the same at block lengths (5, 20)
\providecommand{\rhmcsBarQlikeTrangePMaxOutNinety}{}\renewcommand{\rhmcsBarQlikeTrangePMaxOutNinety}{\ensuremath{<}0.001}  % largest MCS p-value outside the 90 % set
\providecommand{\rhmcsBarQlikeTrangeNNinetyLo}{}\renewcommand{\rhmcsBarQlikeTrangeNNinetyLo}{2}  % ... at block length 5
\providecommand{\rhmcsBarQlikeTrangeSetNinetyLo}{}\renewcommand{\rhmcsBarQlikeTrangeSetNinetyLo}{LightGBM on the full predictor set and XGBoost on the full predictor set}  % ... at block length 5
\providecommand{\rhmcsBarQlikeTrangeNNinetyHi}{}\renewcommand{\rhmcsBarQlikeTrangeNNinetyHi}{2}  % ... at block length 20
\providecommand{\rhmcsBarQlikeTrangeSetNinetyHi}{}\renewcommand{\rhmcsBarQlikeTrangeSetNinetyHi}{LightGBM on the full predictor set and XGBoost on the full predictor set}  % ... at block length 20
\providecommand{\rhmcsBarQlikeTrangeNSevenFive}{}\renewcommand{\rhmcsBarQlikeTrangeNSevenFive}{2}  % models in the 75 % MCS (of 10)
\providecommand{\rhmcsBarQlikeTrangeSetSevenFive}{}\renewcommand{\rhmcsBarQlikeTrangeSetSevenFive}{LightGBM on the full predictor set and XGBoost on the full predictor set}  % the 75 % MCS, best (last eliminated) first
\providecommand{\rhmcsBarQlikeTrangeStableSevenFive}{}\renewcommand{\rhmcsBarQlikeTrangeStableSevenFive}{yes}  % the 75 % MCS is the same at block lengths (5, 20)
\providecommand{\rhmcsBarQlikeTrangePMaxOutSevenFive}{}\renewcommand{\rhmcsBarQlikeTrangePMaxOutSevenFive}{\ensuremath{<}0.001}  % largest MCS p-value outside the 75 % set
\providecommand{\rhmcsBarQlikeTrangeNSevenFiveLo}{}\renewcommand{\rhmcsBarQlikeTrangeNSevenFiveLo}{2}  % ... at block length 5
\providecommand{\rhmcsBarQlikeTrangeSetSevenFiveLo}{}\renewcommand{\rhmcsBarQlikeTrangeSetSevenFiveLo}{LightGBM on the full predictor set and XGBoost on the full predictor set}  % ... at block length 5
\providecommand{\rhmcsBarQlikeTrangeNSevenFiveHi}{}\renewcommand{\rhmcsBarQlikeTrangeNSevenFiveHi}{2}  % ... at block length 20
\providecommand{\rhmcsBarQlikeTrangeSetSevenFiveHi}{}\renewcommand{\rhmcsBarQlikeTrangeSetSevenFiveHi}{LightGBM on the full predictor set and XGBoost on the full predictor set}  % ... at block length 20
\providecommand{\rhmcsBarQlikeTrangeBest}{}\renewcommand{\rhmcsBarQlikeTrangeBest}{LightGBM on the full predictor set}  % the last model left (MCS p-value 1)
\providecommand{\rhmcsBarMseTmaxPLinHar}{}\renewcommand{\rhmcsBarMseTmaxPLinHar}{1.00}  % MCS p-value, every bar, 2014--2024, Mse, Tmax: Linear, HAR
\providecommand{\rhmcsBarMseTmaxPRidgeFull}{}\renewcommand{\rhmcsBarMseTmaxPRidgeFull}{0.49}  % MCS p-value, every bar, 2014--2024, Mse, Tmax: Ridge, Full
\providecommand{\rhmcsBarMseTmaxPLassoFull}{}\renewcommand{\rhmcsBarMseTmaxPLassoFull}{0.32}  % MCS p-value, every bar, 2014--2024, Mse, Tmax: Lasso, Full
\providecommand{\rhmcsBarMseTmaxPEnetFull}{}\renewcommand{\rhmcsBarMseTmaxPEnetFull}{0.43}  % MCS p-value, every bar, 2014--2024, Mse, Tmax: Elastic net, Full
\providecommand{\rhmcsBarMseTmaxPLgbmHar}{}\renewcommand{\rhmcsBarMseTmaxPLgbmHar}{0.49}  % MCS p-value, every bar, 2014--2024, Mse, Tmax: LightGBM, HAR
\providecommand{\rhmcsBarMseTmaxPLgbmFull}{}\renewcommand{\rhmcsBarMseTmaxPLgbmFull}{0.49}  % MCS p-value, every bar, 2014--2024, Mse, Tmax: LightGBM, Full
\providecommand{\rhmcsBarMseTmaxPXgbHar}{}\renewcommand{\rhmcsBarMseTmaxPXgbHar}{0.49}  % MCS p-value, every bar, 2014--2024, Mse, Tmax: XGBoost, HAR
\providecommand{\rhmcsBarMseTmaxPXgbFull}{}\renewcommand{\rhmcsBarMseTmaxPXgbFull}{0.49}  % MCS p-value, every bar, 2014--2024, Mse, Tmax: XGBoost, Full
\providecommand{\rhmcsBarMseTmaxPRfHar}{}\renewcommand{\rhmcsBarMseTmaxPRfHar}{0.49}  % MCS p-value, every bar, 2014--2024, Mse, Tmax: Random forest, HAR
\providecommand{\rhmcsBarMseTmaxPRfFull}{}\renewcommand{\rhmcsBarMseTmaxPRfFull}{0.49}  % MCS p-value, every bar, 2014--2024, Mse, Tmax: Random forest, Full
\providecommand{\rhmcsBarMseTmaxNNinety}{}\renewcommand{\rhmcsBarMseTmaxNNinety}{10}  % models in the 90 % MCS (of 10)
\providecommand{\rhmcsBarMseTmaxSetNinety}{}\renewcommand{\rhmcsBarMseTmaxSetNinety}{the linear models on the HAR benchmark, XGBoost on the HAR benchmark, LightGBM on the HAR benchmark, the random forest on the HAR benchmark, XGBoost on the full predictor set, LightGBM on the full predictor set, the random forest on the full predictor set, ridge regression on the full predictor set, the elastic net on the full predictor set, and the lasso on the full predictor set}  % the 90 % MCS, best (last eliminated) first
\providecommand{\rhmcsBarMseTmaxStableNinety}{}\renewcommand{\rhmcsBarMseTmaxStableNinety}{yes}  % the 90 % MCS is the same at block lengths (5, 20)
\providecommand{\rhmcsBarMseTmaxPMaxOutNinety}{}\renewcommand{\rhmcsBarMseTmaxPMaxOutNinety}{--}  % largest MCS p-value outside the 90 % set
\providecommand{\rhmcsBarMseTmaxNNinetyLo}{}\renewcommand{\rhmcsBarMseTmaxNNinetyLo}{10}  % ... at block length 5
\providecommand{\rhmcsBarMseTmaxSetNinetyLo}{}\renewcommand{\rhmcsBarMseTmaxSetNinetyLo}{the linear models on the HAR benchmark, XGBoost on the HAR benchmark, LightGBM on the HAR benchmark, the random forest on the HAR benchmark, XGBoost on the full predictor set, LightGBM on the full predictor set, the random forest on the full predictor set, the elastic net on the full predictor set, ridge regression on the full predictor set, and the lasso on the full predictor set}  % ... at block length 5
\providecommand{\rhmcsBarMseTmaxNNinetyHi}{}\renewcommand{\rhmcsBarMseTmaxNNinetyHi}{10}  % ... at block length 20
\providecommand{\rhmcsBarMseTmaxSetNinetyHi}{}\renewcommand{\rhmcsBarMseTmaxSetNinetyHi}{the linear models on the HAR benchmark, XGBoost on the HAR benchmark, LightGBM on the HAR benchmark, the random forest on the HAR benchmark, XGBoost on the full predictor set, LightGBM on the full predictor set, the random forest on the full predictor set, ridge regression on the full predictor set, the elastic net on the full predictor set, and the lasso on the full predictor set}  % ... at block length 20
\providecommand{\rhmcsBarMseTmaxNSevenFive}{}\renewcommand{\rhmcsBarMseTmaxNSevenFive}{10}  % models in the 75 % MCS (of 10)
\providecommand{\rhmcsBarMseTmaxSetSevenFive}{}\renewcommand{\rhmcsBarMseTmaxSetSevenFive}{the linear models on the HAR benchmark, XGBoost on the HAR benchmark, LightGBM on the HAR benchmark, the random forest on the HAR benchmark, XGBoost on the full predictor set, LightGBM on the full predictor set, the random forest on the full predictor set, ridge regression on the full predictor set, the elastic net on the full predictor set, and the lasso on the full predictor set}  % the 75 % MCS, best (last eliminated) first
\providecommand{\rhmcsBarMseTmaxStableSevenFive}{}\renewcommand{\rhmcsBarMseTmaxStableSevenFive}{no}  % the 75 % MCS is the same at block lengths (5, 20)
\providecommand{\rhmcsBarMseTmaxPMaxOutSevenFive}{}\renewcommand{\rhmcsBarMseTmaxPMaxOutSevenFive}{--}  % largest MCS p-value outside the 75 % set
\providecommand{\rhmcsBarMseTmaxNSevenFiveLo}{}\renewcommand{\rhmcsBarMseTmaxNSevenFiveLo}{9}  % ... at block length 5
\providecommand{\rhmcsBarMseTmaxSetSevenFiveLo}{}\renewcommand{\rhmcsBarMseTmaxSetSevenFiveLo}{the linear models on the HAR benchmark, XGBoost on the HAR benchmark, LightGBM on the HAR benchmark, the random forest on the HAR benchmark, XGBoost on the full predictor set, LightGBM on the full predictor set, the random forest on the full predictor set, the elastic net on the full predictor set, and ridge regression on the full predictor set}  % ... at block length 5
\providecommand{\rhmcsBarMseTmaxNSevenFiveHi}{}\renewcommand{\rhmcsBarMseTmaxNSevenFiveHi}{10}  % ... at block length 20
\providecommand{\rhmcsBarMseTmaxSetSevenFiveHi}{}\renewcommand{\rhmcsBarMseTmaxSetSevenFiveHi}{the linear models on the HAR benchmark, XGBoost on the HAR benchmark, LightGBM on the HAR benchmark, the random forest on the HAR benchmark, XGBoost on the full predictor set, LightGBM on the full predictor set, the random forest on the full predictor set, ridge regression on the full predictor set, the elastic net on the full predictor set, and the lasso on the full predictor set}  % ... at block length 20
\providecommand{\rhmcsBarMseTmaxBest}{}\renewcommand{\rhmcsBarMseTmaxBest}{the linear models on the HAR benchmark}  % the last model left (MCS p-value 1)
\providecommand{\rhmcsBarMseTrangePLinHar}{}\renewcommand{\rhmcsBarMseTrangePLinHar}{1.00}  % MCS p-value, every bar, 2014--2024, Mse, Trange: Linear, HAR
\providecommand{\rhmcsBarMseTrangePRidgeFull}{}\renewcommand{\rhmcsBarMseTrangePRidgeFull}{0.42}  % MCS p-value, every bar, 2014--2024, Mse, Trange: Ridge, Full
\providecommand{\rhmcsBarMseTrangePLassoFull}{}\renewcommand{\rhmcsBarMseTrangePLassoFull}{0.42}  % MCS p-value, every bar, 2014--2024, Mse, Trange: Lasso, Full
\providecommand{\rhmcsBarMseTrangePEnetFull}{}\renewcommand{\rhmcsBarMseTrangePEnetFull}{0.42}  % MCS p-value, every bar, 2014--2024, Mse, Trange: Elastic net, Full
\providecommand{\rhmcsBarMseTrangePLgbmHar}{}\renewcommand{\rhmcsBarMseTrangePLgbmHar}{0.31}  % MCS p-value, every bar, 2014--2024, Mse, Trange: LightGBM, HAR
\providecommand{\rhmcsBarMseTrangePLgbmFull}{}\renewcommand{\rhmcsBarMseTrangePLgbmFull}{0.42}  % MCS p-value, every bar, 2014--2024, Mse, Trange: LightGBM, Full
\providecommand{\rhmcsBarMseTrangePXgbHar}{}\renewcommand{\rhmcsBarMseTrangePXgbHar}{0.42}  % MCS p-value, every bar, 2014--2024, Mse, Trange: XGBoost, HAR
\providecommand{\rhmcsBarMseTrangePXgbFull}{}\renewcommand{\rhmcsBarMseTrangePXgbFull}{0.42}  % MCS p-value, every bar, 2014--2024, Mse, Trange: XGBoost, Full
\providecommand{\rhmcsBarMseTrangePRfHar}{}\renewcommand{\rhmcsBarMseTrangePRfHar}{0.31}  % MCS p-value, every bar, 2014--2024, Mse, Trange: Random forest, HAR
\providecommand{\rhmcsBarMseTrangePRfFull}{}\renewcommand{\rhmcsBarMseTrangePRfFull}{0.42}  % MCS p-value, every bar, 2014--2024, Mse, Trange: Random forest, Full
\providecommand{\rhmcsBarMseTrangeNNinety}{}\renewcommand{\rhmcsBarMseTrangeNNinety}{10}  % models in the 90 % MCS (of 10)
\providecommand{\rhmcsBarMseTrangeSetNinety}{}\renewcommand{\rhmcsBarMseTrangeSetNinety}{the linear models on the HAR benchmark, the elastic net on the full predictor set, ridge regression on the full predictor set, XGBoost on the full predictor set, LightGBM on the full predictor set, the random forest on the full predictor set, the lasso on the full predictor set, XGBoost on the HAR benchmark, LightGBM on the HAR benchmark, and the random forest on the HAR benchmark}  % the 90 % MCS, best (last eliminated) first
\providecommand{\rhmcsBarMseTrangeStableNinety}{}\renewcommand{\rhmcsBarMseTrangeStableNinety}{yes}  % the 90 % MCS is the same at block lengths (5, 20)
\providecommand{\rhmcsBarMseTrangePMaxOutNinety}{}\renewcommand{\rhmcsBarMseTrangePMaxOutNinety}{--}  % largest MCS p-value outside the 90 % set
\providecommand{\rhmcsBarMseTrangeNNinetyLo}{}\renewcommand{\rhmcsBarMseTrangeNNinetyLo}{10}  % ... at block length 5
\providecommand{\rhmcsBarMseTrangeSetNinetyLo}{}\renewcommand{\rhmcsBarMseTrangeSetNinetyLo}{the linear models on the HAR benchmark, the elastic net on the full predictor set, ridge regression on the full predictor set, XGBoost on the HAR benchmark, XGBoost on the full predictor set, LightGBM on the full predictor set, the random forest on the full predictor set, LightGBM on the HAR benchmark, the lasso on the full predictor set, and the random forest on the HAR benchmark}  % ... at block length 5
\providecommand{\rhmcsBarMseTrangeNNinetyHi}{}\renewcommand{\rhmcsBarMseTrangeNNinetyHi}{10}  % ... at block length 20
\providecommand{\rhmcsBarMseTrangeSetNinetyHi}{}\renewcommand{\rhmcsBarMseTrangeSetNinetyHi}{the linear models on the HAR benchmark, the elastic net on the full predictor set, ridge regression on the full predictor set, XGBoost on the full predictor set, the lasso on the full predictor set, LightGBM on the full predictor set, the random forest on the full predictor set, XGBoost on the HAR benchmark, LightGBM on the HAR benchmark, and the random forest on the HAR benchmark}  % ... at block length 20
\providecommand{\rhmcsBarMseTrangeNSevenFive}{}\renewcommand{\rhmcsBarMseTrangeNSevenFive}{10}  % models in the 75 % MCS (of 10)
\providecommand{\rhmcsBarMseTrangeSetSevenFive}{}\renewcommand{\rhmcsBarMseTrangeSetSevenFive}{the linear models on the HAR benchmark, the elastic net on the full predictor set, ridge regression on the full predictor set, XGBoost on the full predictor set, LightGBM on the full predictor set, the random forest on the full predictor set, the lasso on the full predictor set, XGBoost on the HAR benchmark, LightGBM on the HAR benchmark, and the random forest on the HAR benchmark}  % the 75 % MCS, best (last eliminated) first
\providecommand{\rhmcsBarMseTrangeStableSevenFive}{}\renewcommand{\rhmcsBarMseTrangeStableSevenFive}{yes}  % the 75 % MCS is the same at block lengths (5, 20)
\providecommand{\rhmcsBarMseTrangePMaxOutSevenFive}{}\renewcommand{\rhmcsBarMseTrangePMaxOutSevenFive}{--}  % largest MCS p-value outside the 75 % set
\providecommand{\rhmcsBarMseTrangeNSevenFiveLo}{}\renewcommand{\rhmcsBarMseTrangeNSevenFiveLo}{10}  % ... at block length 5
\providecommand{\rhmcsBarMseTrangeSetSevenFiveLo}{}\renewcommand{\rhmcsBarMseTrangeSetSevenFiveLo}{the linear models on the HAR benchmark, the elastic net on the full predictor set, ridge regression on the full predictor set, XGBoost on the HAR benchmark, XGBoost on the full predictor set, LightGBM on the full predictor set, the random forest on the full predictor set, LightGBM on the HAR benchmark, the lasso on the full predictor set, and the random forest on the HAR benchmark}  % ... at block length 5
\providecommand{\rhmcsBarMseTrangeNSevenFiveHi}{}\renewcommand{\rhmcsBarMseTrangeNSevenFiveHi}{10}  % ... at block length 20
\providecommand{\rhmcsBarMseTrangeSetSevenFiveHi}{}\renewcommand{\rhmcsBarMseTrangeSetSevenFiveHi}{the linear models on the HAR benchmark, the elastic net on the full predictor set, ridge regression on the full predictor set, XGBoost on the full predictor set, the lasso on the full predictor set, LightGBM on the full predictor set, the random forest on the full predictor set, XGBoost on the HAR benchmark, LightGBM on the HAR benchmark, and the random forest on the HAR benchmark}  % ... at block length 20
\providecommand{\rhmcsBarMseTrangeBest}{}\renewcommand{\rhmcsBarMseTrangeBest}{the linear models on the HAR benchmark}  % the last model left (MCS p-value 1)
\providecommand{\rhmcsAllQlikeTmaxPLinHar}{}\renewcommand{\rhmcsAllQlikeTmaxPLinHar}{0.19}  % MCS p-value, 16:00 bar, 2014--2024, Qlike, Tmax: Linear, HAR
\providecommand{\rhmcsAllQlikeTmaxPRidgeFull}{}\renewcommand{\rhmcsAllQlikeTmaxPRidgeFull}{0.19}  % MCS p-value, 16:00 bar, 2014--2024, Qlike, Tmax: Ridge, Full
\providecommand{\rhmcsAllQlikeTmaxPLassoFull}{}\renewcommand{\rhmcsAllQlikeTmaxPLassoFull}{0.071}  % MCS p-value, 16:00 bar, 2014--2024, Qlike, Tmax: Lasso, Full
\providecommand{\rhmcsAllQlikeTmaxPEnetFull}{}\renewcommand{\rhmcsAllQlikeTmaxPEnetFull}{0.19}  % MCS p-value, 16:00 bar, 2014--2024, Qlike, Tmax: Elastic net, Full
\providecommand{\rhmcsAllQlikeTmaxPLgbmHar}{}\renewcommand{\rhmcsAllQlikeTmaxPLgbmHar}{0.19}  % MCS p-value, 16:00 bar, 2014--2024, Qlike, Tmax: LightGBM, HAR
\providecommand{\rhmcsAllQlikeTmaxPLgbmFull}{}\renewcommand{\rhmcsAllQlikeTmaxPLgbmFull}{0.29}  % MCS p-value, 16:00 bar, 2014--2024, Qlike, Tmax: LightGBM, Full
\providecommand{\rhmcsAllQlikeTmaxPXgbHar}{}\renewcommand{\rhmcsAllQlikeTmaxPXgbHar}{0.19}  % MCS p-value, 16:00 bar, 2014--2024, Qlike, Tmax: XGBoost, HAR
\providecommand{\rhmcsAllQlikeTmaxPXgbFull}{}\renewcommand{\rhmcsAllQlikeTmaxPXgbFull}{1.00}  % MCS p-value, 16:00 bar, 2014--2024, Qlike, Tmax: XGBoost, Full
\providecommand{\rhmcsAllQlikeTmaxPRfHar}{}\renewcommand{\rhmcsAllQlikeTmaxPRfHar}{0.19}  % MCS p-value, 16:00 bar, 2014--2024, Qlike, Tmax: Random forest, HAR
\providecommand{\rhmcsAllQlikeTmaxPRfFull}{}\renewcommand{\rhmcsAllQlikeTmaxPRfFull}{0.19}  % MCS p-value, 16:00 bar, 2014--2024, Qlike, Tmax: Random forest, Full
\providecommand{\rhmcsAllQlikeTmaxNNinety}{}\renewcommand{\rhmcsAllQlikeTmaxNNinety}{9}  % models in the 90 % MCS (of 10)
\providecommand{\rhmcsAllQlikeTmaxSetNinety}{}\renewcommand{\rhmcsAllQlikeTmaxSetNinety}{XGBoost on the full predictor set, LightGBM on the full predictor set, XGBoost on the HAR benchmark, LightGBM on the HAR benchmark, the random forest on the full predictor set, the elastic net on the full predictor set, the random forest on the HAR benchmark, ridge regression on the full predictor set, and the linear models on the HAR benchmark}  % the 90 % MCS, best (last eliminated) first
\providecommand{\rhmcsAllQlikeTmaxStableNinety}{}\renewcommand{\rhmcsAllQlikeTmaxStableNinety}{yes}  % the 90 % MCS is the same at block lengths (5, 20)
\providecommand{\rhmcsAllQlikeTmaxPMaxOutNinety}{}\renewcommand{\rhmcsAllQlikeTmaxPMaxOutNinety}{0.071}  % largest MCS p-value outside the 90 % set
\providecommand{\rhmcsAllQlikeTmaxNNinetyLo}{}\renewcommand{\rhmcsAllQlikeTmaxNNinetyLo}{9}  % ... at block length 5
\providecommand{\rhmcsAllQlikeTmaxSetNinetyLo}{}\renewcommand{\rhmcsAllQlikeTmaxSetNinetyLo}{XGBoost on the full predictor set, LightGBM on the full predictor set, XGBoost on the HAR benchmark, LightGBM on the HAR benchmark, the random forest on the full predictor set, the elastic net on the full predictor set, the random forest on the HAR benchmark, ridge regression on the full predictor set, and the linear models on the HAR benchmark}  % ... at block length 5
\providecommand{\rhmcsAllQlikeTmaxNNinetyHi}{}\renewcommand{\rhmcsAllQlikeTmaxNNinetyHi}{9}  % ... at block length 20
\providecommand{\rhmcsAllQlikeTmaxSetNinetyHi}{}\renewcommand{\rhmcsAllQlikeTmaxSetNinetyHi}{XGBoost on the full predictor set, LightGBM on the full predictor set, XGBoost on the HAR benchmark, LightGBM on the HAR benchmark, the random forest on the full predictor set, the elastic net on the full predictor set, the random forest on the HAR benchmark, ridge regression on the full predictor set, and the linear models on the HAR benchmark}  % ... at block length 20
\providecommand{\rhmcsAllQlikeTmaxNSevenFive}{}\renewcommand{\rhmcsAllQlikeTmaxNSevenFive}{2}  % models in the 75 % MCS (of 10)
\providecommand{\rhmcsAllQlikeTmaxSetSevenFive}{}\renewcommand{\rhmcsAllQlikeTmaxSetSevenFive}{XGBoost on the full predictor set and LightGBM on the full predictor set}  % the 75 % MCS, best (last eliminated) first
\providecommand{\rhmcsAllQlikeTmaxStableSevenFive}{}\renewcommand{\rhmcsAllQlikeTmaxStableSevenFive}{yes}  % the 75 % MCS is the same at block lengths (5, 20)
\providecommand{\rhmcsAllQlikeTmaxPMaxOutSevenFive}{}\renewcommand{\rhmcsAllQlikeTmaxPMaxOutSevenFive}{0.19}  % largest MCS p-value outside the 75 % set
\providecommand{\rhmcsAllQlikeTmaxNSevenFiveLo}{}\renewcommand{\rhmcsAllQlikeTmaxNSevenFiveLo}{2}  % ... at block length 5
\providecommand{\rhmcsAllQlikeTmaxSetSevenFiveLo}{}\renewcommand{\rhmcsAllQlikeTmaxSetSevenFiveLo}{XGBoost on the full predictor set and LightGBM on the full predictor set}  % ... at block length 5
\providecommand{\rhmcsAllQlikeTmaxNSevenFiveHi}{}\renewcommand{\rhmcsAllQlikeTmaxNSevenFiveHi}{2}  % ... at block length 20
\providecommand{\rhmcsAllQlikeTmaxSetSevenFiveHi}{}\renewcommand{\rhmcsAllQlikeTmaxSetSevenFiveHi}{XGBoost on the full predictor set and LightGBM on the full predictor set}  % ... at block length 20
\providecommand{\rhmcsAllQlikeTmaxBest}{}\renewcommand{\rhmcsAllQlikeTmaxBest}{XGBoost on the full predictor set}  % the last model left (MCS p-value 1)
\providecommand{\rhmcsAllQlikeTrangePLinHar}{}\renewcommand{\rhmcsAllQlikeTrangePLinHar}{\ensuremath{<}0.001}  % MCS p-value, 16:00 bar, 2014--2024, Qlike, Trange: Linear, HAR
\providecommand{\rhmcsAllQlikeTrangePRidgeFull}{}\renewcommand{\rhmcsAllQlikeTrangePRidgeFull}{0.008}  % MCS p-value, 16:00 bar, 2014--2024, Qlike, Trange: Ridge, Full
\providecommand{\rhmcsAllQlikeTrangePLassoFull}{}\renewcommand{\rhmcsAllQlikeTrangePLassoFull}{0.006}  % MCS p-value, 16:00 bar, 2014--2024, Qlike, Trange: Lasso, Full
\providecommand{\rhmcsAllQlikeTrangePEnetFull}{}\renewcommand{\rhmcsAllQlikeTrangePEnetFull}{0.010}  % MCS p-value, 16:00 bar, 2014--2024, Qlike, Trange: Elastic net, Full
\providecommand{\rhmcsAllQlikeTrangePLgbmHar}{}\renewcommand{\rhmcsAllQlikeTrangePLgbmHar}{0.14}  % MCS p-value, 16:00 bar, 2014--2024, Qlike, Trange: LightGBM, HAR
\providecommand{\rhmcsAllQlikeTrangePLgbmFull}{}\renewcommand{\rhmcsAllQlikeTrangePLgbmFull}{0.29}  % MCS p-value, 16:00 bar, 2014--2024, Qlike, Trange: LightGBM, Full
\providecommand{\rhmcsAllQlikeTrangePXgbHar}{}\renewcommand{\rhmcsAllQlikeTrangePXgbHar}{0.14}  % MCS p-value, 16:00 bar, 2014--2024, Qlike, Trange: XGBoost, HAR
\providecommand{\rhmcsAllQlikeTrangePXgbFull}{}\renewcommand{\rhmcsAllQlikeTrangePXgbFull}{1.00}  % MCS p-value, 16:00 bar, 2014--2024, Qlike, Trange: XGBoost, Full
\providecommand{\rhmcsAllQlikeTrangePRfHar}{}\renewcommand{\rhmcsAllQlikeTrangePRfHar}{0.001}  % MCS p-value, 16:00 bar, 2014--2024, Qlike, Trange: Random forest, HAR
\providecommand{\rhmcsAllQlikeTrangePRfFull}{}\renewcommand{\rhmcsAllQlikeTrangePRfFull}{\ensuremath{<}0.001}  % MCS p-value, 16:00 bar, 2014--2024, Qlike, Trange: Random forest, Full
\providecommand{\rhmcsAllQlikeTrangeNNinety}{}\renewcommand{\rhmcsAllQlikeTrangeNNinety}{4}  % models in the 90 % MCS (of 10)
\providecommand{\rhmcsAllQlikeTrangeSetNinety}{}\renewcommand{\rhmcsAllQlikeTrangeSetNinety}{XGBoost on the full predictor set, LightGBM on the full predictor set, XGBoost on the HAR benchmark, and LightGBM on the HAR benchmark}  % the 90 % MCS, best (last eliminated) first
\providecommand{\rhmcsAllQlikeTrangeStableNinety}{}\renewcommand{\rhmcsAllQlikeTrangeStableNinety}{yes}  % the 90 % MCS is the same at block lengths (5, 20)
\providecommand{\rhmcsAllQlikeTrangePMaxOutNinety}{}\renewcommand{\rhmcsAllQlikeTrangePMaxOutNinety}{0.010}  % largest MCS p-value outside the 90 % set
\providecommand{\rhmcsAllQlikeTrangeNNinetyLo}{}\renewcommand{\rhmcsAllQlikeTrangeNNinetyLo}{4}  % ... at block length 5
\providecommand{\rhmcsAllQlikeTrangeSetNinetyLo}{}\renewcommand{\rhmcsAllQlikeTrangeSetNinetyLo}{XGBoost on the full predictor set, LightGBM on the full predictor set, LightGBM on the HAR benchmark, and XGBoost on the HAR benchmark}  % ... at block length 5
\providecommand{\rhmcsAllQlikeTrangeNNinetyHi}{}\renewcommand{\rhmcsAllQlikeTrangeNNinetyHi}{4}  % ... at block length 20
\providecommand{\rhmcsAllQlikeTrangeSetNinetyHi}{}\renewcommand{\rhmcsAllQlikeTrangeSetNinetyHi}{XGBoost on the full predictor set, LightGBM on the full predictor set, LightGBM on the HAR benchmark, and XGBoost on the HAR benchmark}  % ... at block length 20
\providecommand{\rhmcsAllQlikeTrangeNSevenFive}{}\renewcommand{\rhmcsAllQlikeTrangeNSevenFive}{2}  % models in the 75 % MCS (of 10)
\providecommand{\rhmcsAllQlikeTrangeSetSevenFive}{}\renewcommand{\rhmcsAllQlikeTrangeSetSevenFive}{XGBoost on the full predictor set and LightGBM on the full predictor set}  % the 75 % MCS, best (last eliminated) first
\providecommand{\rhmcsAllQlikeTrangeStableSevenFive}{}\renewcommand{\rhmcsAllQlikeTrangeStableSevenFive}{yes}  % the 75 % MCS is the same at block lengths (5, 20)
\providecommand{\rhmcsAllQlikeTrangePMaxOutSevenFive}{}\renewcommand{\rhmcsAllQlikeTrangePMaxOutSevenFive}{0.14}  % largest MCS p-value outside the 75 % set
\providecommand{\rhmcsAllQlikeTrangeNSevenFiveLo}{}\renewcommand{\rhmcsAllQlikeTrangeNSevenFiveLo}{2}  % ... at block length 5
\providecommand{\rhmcsAllQlikeTrangeSetSevenFiveLo}{}\renewcommand{\rhmcsAllQlikeTrangeSetSevenFiveLo}{XGBoost on the full predictor set and LightGBM on the full predictor set}  % ... at block length 5
\providecommand{\rhmcsAllQlikeTrangeNSevenFiveHi}{}\renewcommand{\rhmcsAllQlikeTrangeNSevenFiveHi}{2}  % ... at block length 20
\providecommand{\rhmcsAllQlikeTrangeSetSevenFiveHi}{}\renewcommand{\rhmcsAllQlikeTrangeSetSevenFiveHi}{XGBoost on the full predictor set and LightGBM on the full predictor set}  % ... at block length 20
\providecommand{\rhmcsAllQlikeTrangeBest}{}\renewcommand{\rhmcsAllQlikeTrangeBest}{XGBoost on the full predictor set}  % the last model left (MCS p-value 1)
\providecommand{\rhmcsAllMseTmaxPLinHar}{}\renewcommand{\rhmcsAllMseTmaxPLinHar}{0.59}  % MCS p-value, 16:00 bar, 2014--2024, Mse, Tmax: Linear, HAR
\providecommand{\rhmcsAllMseTmaxPRidgeFull}{}\renewcommand{\rhmcsAllMseTmaxPRidgeFull}{0.43}  % MCS p-value, 16:00 bar, 2014--2024, Mse, Tmax: Ridge, Full
\providecommand{\rhmcsAllMseTmaxPLassoFull}{}\renewcommand{\rhmcsAllMseTmaxPLassoFull}{0.38}  % MCS p-value, 16:00 bar, 2014--2024, Mse, Tmax: Lasso, Full
\providecommand{\rhmcsAllMseTmaxPEnetFull}{}\renewcommand{\rhmcsAllMseTmaxPEnetFull}{0.48}  % MCS p-value, 16:00 bar, 2014--2024, Mse, Tmax: Elastic net, Full
\providecommand{\rhmcsAllMseTmaxPLgbmHar}{}\renewcommand{\rhmcsAllMseTmaxPLgbmHar}{0.59}  % MCS p-value, 16:00 bar, 2014--2024, Mse, Tmax: LightGBM, HAR
\providecommand{\rhmcsAllMseTmaxPLgbmFull}{}\renewcommand{\rhmcsAllMseTmaxPLgbmFull}{0.48}  % MCS p-value, 16:00 bar, 2014--2024, Mse, Tmax: LightGBM, Full
\providecommand{\rhmcsAllMseTmaxPXgbHar}{}\renewcommand{\rhmcsAllMseTmaxPXgbHar}{0.59}  % MCS p-value, 16:00 bar, 2014--2024, Mse, Tmax: XGBoost, HAR
\providecommand{\rhmcsAllMseTmaxPXgbFull}{}\renewcommand{\rhmcsAllMseTmaxPXgbFull}{0.48}  % MCS p-value, 16:00 bar, 2014--2024, Mse, Tmax: XGBoost, Full
\providecommand{\rhmcsAllMseTmaxPRfHar}{}\renewcommand{\rhmcsAllMseTmaxPRfHar}{1.00}  % MCS p-value, 16:00 bar, 2014--2024, Mse, Tmax: Random forest, HAR
\providecommand{\rhmcsAllMseTmaxPRfFull}{}\renewcommand{\rhmcsAllMseTmaxPRfFull}{0.59}  % MCS p-value, 16:00 bar, 2014--2024, Mse, Tmax: Random forest, Full
\providecommand{\rhmcsAllMseTmaxNNinety}{}\renewcommand{\rhmcsAllMseTmaxNNinety}{10}  % models in the 90 % MCS (of 10)
\providecommand{\rhmcsAllMseTmaxSetNinety}{}\renewcommand{\rhmcsAllMseTmaxSetNinety}{the random forest on the HAR benchmark, XGBoost on the HAR benchmark, LightGBM on the HAR benchmark, the linear models on the HAR benchmark, the random forest on the full predictor set, XGBoost on the full predictor set, LightGBM on the full predictor set, the elastic net on the full predictor set, ridge regression on the full predictor set, and the lasso on the full predictor set}  % the 90 % MCS, best (last eliminated) first
\providecommand{\rhmcsAllMseTmaxStableNinety}{}\renewcommand{\rhmcsAllMseTmaxStableNinety}{yes}  % the 90 % MCS is the same at block lengths (5, 20)
\providecommand{\rhmcsAllMseTmaxPMaxOutNinety}{}\renewcommand{\rhmcsAllMseTmaxPMaxOutNinety}{--}  % largest MCS p-value outside the 90 % set
\providecommand{\rhmcsAllMseTmaxNNinetyLo}{}\renewcommand{\rhmcsAllMseTmaxNNinetyLo}{10}  % ... at block length 5
\providecommand{\rhmcsAllMseTmaxSetNinetyLo}{}\renewcommand{\rhmcsAllMseTmaxSetNinetyLo}{the random forest on the HAR benchmark, XGBoost on the HAR benchmark, LightGBM on the HAR benchmark, the linear models on the HAR benchmark, the random forest on the full predictor set, XGBoost on the full predictor set, LightGBM on the full predictor set, the elastic net on the full predictor set, ridge regression on the full predictor set, and the lasso on the full predictor set}  % ... at block length 5
\providecommand{\rhmcsAllMseTmaxNNinetyHi}{}\renewcommand{\rhmcsAllMseTmaxNNinetyHi}{10}  % ... at block length 20
\providecommand{\rhmcsAllMseTmaxSetNinetyHi}{}\renewcommand{\rhmcsAllMseTmaxSetNinetyHi}{the random forest on the HAR benchmark, XGBoost on the HAR benchmark, LightGBM on the HAR benchmark, the linear models on the HAR benchmark, the random forest on the full predictor set, XGBoost on the full predictor set, LightGBM on the full predictor set, the elastic net on the full predictor set, ridge regression on the full predictor set, and the lasso on the full predictor set}  % ... at block length 20
\providecommand{\rhmcsAllMseTmaxNSevenFive}{}\renewcommand{\rhmcsAllMseTmaxNSevenFive}{10}  % models in the 75 % MCS (of 10)
\providecommand{\rhmcsAllMseTmaxSetSevenFive}{}\renewcommand{\rhmcsAllMseTmaxSetSevenFive}{the random forest on the HAR benchmark, XGBoost on the HAR benchmark, LightGBM on the HAR benchmark, the linear models on the HAR benchmark, the random forest on the full predictor set, XGBoost on the full predictor set, LightGBM on the full predictor set, the elastic net on the full predictor set, ridge regression on the full predictor set, and the lasso on the full predictor set}  % the 75 % MCS, best (last eliminated) first
\providecommand{\rhmcsAllMseTmaxStableSevenFive}{}\renewcommand{\rhmcsAllMseTmaxStableSevenFive}{no}  % the 75 % MCS is the same at block lengths (5, 20)
\providecommand{\rhmcsAllMseTmaxPMaxOutSevenFive}{}\renewcommand{\rhmcsAllMseTmaxPMaxOutSevenFive}{--}  % largest MCS p-value outside the 75 % set
\providecommand{\rhmcsAllMseTmaxNSevenFiveLo}{}\renewcommand{\rhmcsAllMseTmaxNSevenFiveLo}{9}  % ... at block length 5
\providecommand{\rhmcsAllMseTmaxSetSevenFiveLo}{}\renewcommand{\rhmcsAllMseTmaxSetSevenFiveLo}{the random forest on the HAR benchmark, XGBoost on the HAR benchmark, LightGBM on the HAR benchmark, the linear models on the HAR benchmark, the random forest on the full predictor set, XGBoost on the full predictor set, LightGBM on the full predictor set, the elastic net on the full predictor set, and ridge regression on the full predictor set}  % ... at block length 5
\providecommand{\rhmcsAllMseTmaxNSevenFiveHi}{}\renewcommand{\rhmcsAllMseTmaxNSevenFiveHi}{10}  % ... at block length 20
\providecommand{\rhmcsAllMseTmaxSetSevenFiveHi}{}\renewcommand{\rhmcsAllMseTmaxSetSevenFiveHi}{the random forest on the HAR benchmark, XGBoost on the HAR benchmark, LightGBM on the HAR benchmark, the linear models on the HAR benchmark, the random forest on the full predictor set, XGBoost on the full predictor set, LightGBM on the full predictor set, the elastic net on the full predictor set, ridge regression on the full predictor set, and the lasso on the full predictor set}  % ... at block length 20
\providecommand{\rhmcsAllMseTmaxBest}{}\renewcommand{\rhmcsAllMseTmaxBest}{the random forest on the HAR benchmark}  % the last model left (MCS p-value 1)
\providecommand{\rhmcsAllMseTrangePLinHar}{}\renewcommand{\rhmcsAllMseTrangePLinHar}{0.60}  % MCS p-value, 16:00 bar, 2014--2024, Mse, Trange: Linear, HAR
\providecommand{\rhmcsAllMseTrangePRidgeFull}{}\renewcommand{\rhmcsAllMseTrangePRidgeFull}{0.60}  % MCS p-value, 16:00 bar, 2014--2024, Mse, Trange: Ridge, Full
\providecommand{\rhmcsAllMseTrangePLassoFull}{}\renewcommand{\rhmcsAllMseTrangePLassoFull}{0.60}  % MCS p-value, 16:00 bar, 2014--2024, Mse, Trange: Lasso, Full
\providecommand{\rhmcsAllMseTrangePEnetFull}{}\renewcommand{\rhmcsAllMseTrangePEnetFull}{0.60}  % MCS p-value, 16:00 bar, 2014--2024, Mse, Trange: Elastic net, Full
\providecommand{\rhmcsAllMseTrangePLgbmHar}{}\renewcommand{\rhmcsAllMseTrangePLgbmHar}{0.35}  % MCS p-value, 16:00 bar, 2014--2024, Mse, Trange: LightGBM, HAR
\providecommand{\rhmcsAllMseTrangePLgbmFull}{}\renewcommand{\rhmcsAllMseTrangePLgbmFull}{0.60}  % MCS p-value, 16:00 bar, 2014--2024, Mse, Trange: LightGBM, Full
\providecommand{\rhmcsAllMseTrangePXgbHar}{}\renewcommand{\rhmcsAllMseTrangePXgbHar}{0.60}  % MCS p-value, 16:00 bar, 2014--2024, Mse, Trange: XGBoost, HAR
\providecommand{\rhmcsAllMseTrangePXgbFull}{}\renewcommand{\rhmcsAllMseTrangePXgbFull}{0.60}  % MCS p-value, 16:00 bar, 2014--2024, Mse, Trange: XGBoost, Full
\providecommand{\rhmcsAllMseTrangePRfHar}{}\renewcommand{\rhmcsAllMseTrangePRfHar}{1.00}  % MCS p-value, 16:00 bar, 2014--2024, Mse, Trange: Random forest, HAR
\providecommand{\rhmcsAllMseTrangePRfFull}{}\renewcommand{\rhmcsAllMseTrangePRfFull}{0.60}  % MCS p-value, 16:00 bar, 2014--2024, Mse, Trange: Random forest, Full
\providecommand{\rhmcsAllMseTrangeNNinety}{}\renewcommand{\rhmcsAllMseTrangeNNinety}{10}  % models in the 90 % MCS (of 10)
\providecommand{\rhmcsAllMseTrangeSetNinety}{}\renewcommand{\rhmcsAllMseTrangeSetNinety}{the random forest on the HAR benchmark, the linear models on the HAR benchmark, XGBoost on the HAR benchmark, the random forest on the full predictor set, XGBoost on the full predictor set, the elastic net on the full predictor set, ridge regression on the full predictor set, the lasso on the full predictor set, LightGBM on the full predictor set, and LightGBM on the HAR benchmark}  % the 90 % MCS, best (last eliminated) first
\providecommand{\rhmcsAllMseTrangeStableNinety}{}\renewcommand{\rhmcsAllMseTrangeStableNinety}{yes}  % the 90 % MCS is the same at block lengths (5, 20)
\providecommand{\rhmcsAllMseTrangePMaxOutNinety}{}\renewcommand{\rhmcsAllMseTrangePMaxOutNinety}{--}  % largest MCS p-value outside the 90 % set
\providecommand{\rhmcsAllMseTrangeNNinetyLo}{}\renewcommand{\rhmcsAllMseTrangeNNinetyLo}{10}  % ... at block length 5
\providecommand{\rhmcsAllMseTrangeSetNinetyLo}{}\renewcommand{\rhmcsAllMseTrangeSetNinetyLo}{the random forest on the HAR benchmark, the linear models on the HAR benchmark, XGBoost on the HAR benchmark, XGBoost on the full predictor set, the random forest on the full predictor set, the elastic net on the full predictor set, ridge regression on the full predictor set, LightGBM on the full predictor set, the lasso on the full predictor set, and LightGBM on the HAR benchmark}  % ... at block length 5
\providecommand{\rhmcsAllMseTrangeNNinetyHi}{}\renewcommand{\rhmcsAllMseTrangeNNinetyHi}{10}  % ... at block length 20
\providecommand{\rhmcsAllMseTrangeSetNinetyHi}{}\renewcommand{\rhmcsAllMseTrangeSetNinetyHi}{the random forest on the HAR benchmark, the linear models on the HAR benchmark, the random forest on the full predictor set, XGBoost on the full predictor set, the elastic net on the full predictor set, ridge regression on the full predictor set, the lasso on the full predictor set, XGBoost on the HAR benchmark, LightGBM on the full predictor set, and LightGBM on the HAR benchmark}  % ... at block length 20
\providecommand{\rhmcsAllMseTrangeNSevenFive}{}\renewcommand{\rhmcsAllMseTrangeNSevenFive}{10}  % models in the 75 % MCS (of 10)
\providecommand{\rhmcsAllMseTrangeSetSevenFive}{}\renewcommand{\rhmcsAllMseTrangeSetSevenFive}{the random forest on the HAR benchmark, the linear models on the HAR benchmark, XGBoost on the HAR benchmark, the random forest on the full predictor set, XGBoost on the full predictor set, the elastic net on the full predictor set, ridge regression on the full predictor set, the lasso on the full predictor set, LightGBM on the full predictor set, and LightGBM on the HAR benchmark}  % the 75 % MCS, best (last eliminated) first
\providecommand{\rhmcsAllMseTrangeStableSevenFive}{}\renewcommand{\rhmcsAllMseTrangeStableSevenFive}{yes}  % the 75 % MCS is the same at block lengths (5, 20)
\providecommand{\rhmcsAllMseTrangePMaxOutSevenFive}{}\renewcommand{\rhmcsAllMseTrangePMaxOutSevenFive}{--}  % largest MCS p-value outside the 75 % set
\providecommand{\rhmcsAllMseTrangeNSevenFiveLo}{}\renewcommand{\rhmcsAllMseTrangeNSevenFiveLo}{10}  % ... at block length 5
\providecommand{\rhmcsAllMseTrangeSetSevenFiveLo}{}\renewcommand{\rhmcsAllMseTrangeSetSevenFiveLo}{the random forest on the HAR benchmark, the linear models on the HAR benchmark, XGBoost on the HAR benchmark, XGBoost on the full predictor set, the random forest on the full predictor set, the elastic net on the full predictor set, ridge regression on the full predictor set, LightGBM on the full predictor set, the lasso on the full predictor set, and LightGBM on the HAR benchmark}  % ... at block length 5
\providecommand{\rhmcsAllMseTrangeNSevenFiveHi}{}\renewcommand{\rhmcsAllMseTrangeNSevenFiveHi}{10}  % ... at block length 20
\providecommand{\rhmcsAllMseTrangeSetSevenFiveHi}{}\renewcommand{\rhmcsAllMseTrangeSetSevenFiveHi}{the random forest on the HAR benchmark, the linear models on the HAR benchmark, the random forest on the full predictor set, XGBoost on the full predictor set, the elastic net on the full predictor set, ridge regression on the full predictor set, the lasso on the full predictor set, XGBoost on the HAR benchmark, LightGBM on the full predictor set, and LightGBM on the HAR benchmark}  % ... at block length 20
\providecommand{\rhmcsAllMseTrangeBest}{}\renewcommand{\rhmcsAllMseTrangeBest}{the random forest on the HAR benchmark}  % the last model left (MCS p-value 1)
\providecommand{\rhmcsOosQlikeTmaxPLinHar}{}\renewcommand{\rhmcsOosQlikeTmaxPLinHar}{0.043}  % MCS p-value, 16:00 bar, T_oos, Qlike, Tmax: Linear, HAR
\providecommand{\rhmcsOosQlikeTmaxPRidgeFull}{}\renewcommand{\rhmcsOosQlikeTmaxPRidgeFull}{0.008}  % MCS p-value, 16:00 bar, T_oos, Qlike, Tmax: Ridge, Full
\providecommand{\rhmcsOosQlikeTmaxPLassoFull}{}\renewcommand{\rhmcsOosQlikeTmaxPLassoFull}{0.008}  % MCS p-value, 16:00 bar, T_oos, Qlike, Tmax: Lasso, Full
\providecommand{\rhmcsOosQlikeTmaxPEnetFull}{}\renewcommand{\rhmcsOosQlikeTmaxPEnetFull}{0.024}  % MCS p-value, 16:00 bar, T_oos, Qlike, Tmax: Elastic net, Full
\providecommand{\rhmcsOosQlikeTmaxPLgbmHar}{}\renewcommand{\rhmcsOosQlikeTmaxPLgbmHar}{0.16}  % MCS p-value, 16:00 bar, T_oos, Qlike, Tmax: LightGBM, HAR
\providecommand{\rhmcsOosQlikeTmaxPLgbmFull}{}\renewcommand{\rhmcsOosQlikeTmaxPLgbmFull}{0.18}  % MCS p-value, 16:00 bar, T_oos, Qlike, Tmax: LightGBM, Full
\providecommand{\rhmcsOosQlikeTmaxPXgbHar}{}\renewcommand{\rhmcsOosQlikeTmaxPXgbHar}{0.16}  % MCS p-value, 16:00 bar, T_oos, Qlike, Tmax: XGBoost, HAR
\providecommand{\rhmcsOosQlikeTmaxPXgbFull}{}\renewcommand{\rhmcsOosQlikeTmaxPXgbFull}{0.18}  % MCS p-value, 16:00 bar, T_oos, Qlike, Tmax: XGBoost, Full
\providecommand{\rhmcsOosQlikeTmaxPRfHar}{}\renewcommand{\rhmcsOosQlikeTmaxPRfHar}{0.043}  % MCS p-value, 16:00 bar, T_oos, Qlike, Tmax: Random forest, HAR
\providecommand{\rhmcsOosQlikeTmaxPRfFull}{}\renewcommand{\rhmcsOosQlikeTmaxPRfFull}{0.14}  % MCS p-value, 16:00 bar, T_oos, Qlike, Tmax: Random forest, Full
\providecommand{\rhmcsOosQlikeTmaxPBench}{}\renewcommand{\rhmcsOosQlikeTmaxPBench}{1.00}  % MCS p-value, 16:00 bar, T_oos, Qlike, Tmax: Closing-interval HAR benchmark
\providecommand{\rhmcsOosQlikeTmaxNNinety}{}\renewcommand{\rhmcsOosQlikeTmaxNNinety}{6}  % models in the 90 % MCS (of 11)
\providecommand{\rhmcsOosQlikeTmaxSetNinety}{}\renewcommand{\rhmcsOosQlikeTmaxSetNinety}{the closing-interval HAR benchmark, XGBoost on the full predictor set, LightGBM on the full predictor set, XGBoost on the HAR benchmark, LightGBM on the HAR benchmark, and the random forest on the full predictor set}  % the 90 % MCS, best (last eliminated) first
\providecommand{\rhmcsOosQlikeTmaxStableNinety}{}\renewcommand{\rhmcsOosQlikeTmaxStableNinety}{yes}  % the 90 % MCS is the same at block lengths (5, 20)
\providecommand{\rhmcsOosQlikeTmaxPMaxOutNinety}{}\renewcommand{\rhmcsOosQlikeTmaxPMaxOutNinety}{0.043}  % largest MCS p-value outside the 90 % set
\providecommand{\rhmcsOosQlikeTmaxNNinetyLo}{}\renewcommand{\rhmcsOosQlikeTmaxNNinetyLo}{6}  % ... at block length 5
\providecommand{\rhmcsOosQlikeTmaxSetNinetyLo}{}\renewcommand{\rhmcsOosQlikeTmaxSetNinetyLo}{the closing-interval HAR benchmark, XGBoost on the full predictor set, LightGBM on the full predictor set, XGBoost on the HAR benchmark, LightGBM on the HAR benchmark, and the random forest on the full predictor set}  % ... at block length 5
\providecommand{\rhmcsOosQlikeTmaxNNinetyHi}{}\renewcommand{\rhmcsOosQlikeTmaxNNinetyHi}{6}  % ... at block length 20
\providecommand{\rhmcsOosQlikeTmaxSetNinetyHi}{}\renewcommand{\rhmcsOosQlikeTmaxSetNinetyHi}{the closing-interval HAR benchmark, XGBoost on the full predictor set, LightGBM on the full predictor set, LightGBM on the HAR benchmark, XGBoost on the HAR benchmark, and the random forest on the full predictor set}  % ... at block length 20
\providecommand{\rhmcsOosQlikeTmaxNSevenFive}{}\renewcommand{\rhmcsOosQlikeTmaxNSevenFive}{1}  % models in the 75 % MCS (of 11)
\providecommand{\rhmcsOosQlikeTmaxSetSevenFive}{}\renewcommand{\rhmcsOosQlikeTmaxSetSevenFive}{the closing-interval HAR benchmark}  % the 75 % MCS, best (last eliminated) first
\providecommand{\rhmcsOosQlikeTmaxStableSevenFive}{}\renewcommand{\rhmcsOosQlikeTmaxStableSevenFive}{yes}  % the 75 % MCS is the same at block lengths (5, 20)
\providecommand{\rhmcsOosQlikeTmaxPMaxOutSevenFive}{}\renewcommand{\rhmcsOosQlikeTmaxPMaxOutSevenFive}{0.18}  % largest MCS p-value outside the 75 % set
\providecommand{\rhmcsOosQlikeTmaxNSevenFiveLo}{}\renewcommand{\rhmcsOosQlikeTmaxNSevenFiveLo}{1}  % ... at block length 5
\providecommand{\rhmcsOosQlikeTmaxSetSevenFiveLo}{}\renewcommand{\rhmcsOosQlikeTmaxSetSevenFiveLo}{the closing-interval HAR benchmark}  % ... at block length 5
\providecommand{\rhmcsOosQlikeTmaxNSevenFiveHi}{}\renewcommand{\rhmcsOosQlikeTmaxNSevenFiveHi}{1}  % ... at block length 20
\providecommand{\rhmcsOosQlikeTmaxSetSevenFiveHi}{}\renewcommand{\rhmcsOosQlikeTmaxSetSevenFiveHi}{the closing-interval HAR benchmark}  % ... at block length 20
\providecommand{\rhmcsOosQlikeTmaxBest}{}\renewcommand{\rhmcsOosQlikeTmaxBest}{the closing-interval HAR benchmark}  % the last model left (MCS p-value 1)
\providecommand{\rhmcsOosQlikeTrangePLinHar}{}\renewcommand{\rhmcsOosQlikeTrangePLinHar}{0.034}  % MCS p-value, 16:00 bar, T_oos, Qlike, Trange: Linear, HAR
\providecommand{\rhmcsOosQlikeTrangePRidgeFull}{}\renewcommand{\rhmcsOosQlikeTrangePRidgeFull}{0.007}  % MCS p-value, 16:00 bar, T_oos, Qlike, Trange: Ridge, Full
\providecommand{\rhmcsOosQlikeTrangePLassoFull}{}\renewcommand{\rhmcsOosQlikeTrangePLassoFull}{0.008}  % MCS p-value, 16:00 bar, T_oos, Qlike, Trange: Lasso, Full
\providecommand{\rhmcsOosQlikeTrangePEnetFull}{}\renewcommand{\rhmcsOosQlikeTrangePEnetFull}{0.013}  % MCS p-value, 16:00 bar, T_oos, Qlike, Trange: Elastic net, Full
\providecommand{\rhmcsOosQlikeTrangePLgbmHar}{}\renewcommand{\rhmcsOosQlikeTrangePLgbmHar}{0.034}  % MCS p-value, 16:00 bar, T_oos, Qlike, Trange: LightGBM, HAR
\providecommand{\rhmcsOosQlikeTrangePLgbmFull}{}\renewcommand{\rhmcsOosQlikeTrangePLgbmFull}{0.29}  % MCS p-value, 16:00 bar, T_oos, Qlike, Trange: LightGBM, Full
\providecommand{\rhmcsOosQlikeTrangePXgbHar}{}\renewcommand{\rhmcsOosQlikeTrangePXgbHar}{0.034}  % MCS p-value, 16:00 bar, T_oos, Qlike, Trange: XGBoost, HAR
\providecommand{\rhmcsOosQlikeTrangePXgbFull}{}\renewcommand{\rhmcsOosQlikeTrangePXgbFull}{0.29}  % MCS p-value, 16:00 bar, T_oos, Qlike, Trange: XGBoost, Full
\providecommand{\rhmcsOosQlikeTrangePRfHar}{}\renewcommand{\rhmcsOosQlikeTrangePRfHar}{0.025}  % MCS p-value, 16:00 bar, T_oos, Qlike, Trange: Random forest, HAR
\providecommand{\rhmcsOosQlikeTrangePRfFull}{}\renewcommand{\rhmcsOosQlikeTrangePRfFull}{0.025}  % MCS p-value, 16:00 bar, T_oos, Qlike, Trange: Random forest, Full
\providecommand{\rhmcsOosQlikeTrangePBench}{}\renewcommand{\rhmcsOosQlikeTrangePBench}{1.00}  % MCS p-value, 16:00 bar, T_oos, Qlike, Trange: Closing-interval HAR benchmark
\providecommand{\rhmcsOosQlikeTrangeNNinety}{}\renewcommand{\rhmcsOosQlikeTrangeNNinety}{3}  % models in the 90 % MCS (of 11)
\providecommand{\rhmcsOosQlikeTrangeSetNinety}{}\renewcommand{\rhmcsOosQlikeTrangeSetNinety}{the closing-interval HAR benchmark, XGBoost on the full predictor set, and LightGBM on the full predictor set}  % the 90 % MCS, best (last eliminated) first
\providecommand{\rhmcsOosQlikeTrangeStableNinety}{}\renewcommand{\rhmcsOosQlikeTrangeStableNinety}{yes}  % the 90 % MCS is the same at block lengths (5, 20)
\providecommand{\rhmcsOosQlikeTrangePMaxOutNinety}{}\renewcommand{\rhmcsOosQlikeTrangePMaxOutNinety}{0.034}  % largest MCS p-value outside the 90 % set
\providecommand{\rhmcsOosQlikeTrangeNNinetyLo}{}\renewcommand{\rhmcsOosQlikeTrangeNNinetyLo}{3}  % ... at block length 5
\providecommand{\rhmcsOosQlikeTrangeSetNinetyLo}{}\renewcommand{\rhmcsOosQlikeTrangeSetNinetyLo}{the closing-interval HAR benchmark, XGBoost on the full predictor set, and LightGBM on the full predictor set}  % ... at block length 5
\providecommand{\rhmcsOosQlikeTrangeNNinetyHi}{}\renewcommand{\rhmcsOosQlikeTrangeNNinetyHi}{3}  % ... at block length 20
\providecommand{\rhmcsOosQlikeTrangeSetNinetyHi}{}\renewcommand{\rhmcsOosQlikeTrangeSetNinetyHi}{the closing-interval HAR benchmark, XGBoost on the full predictor set, and LightGBM on the full predictor set}  % ... at block length 20
\providecommand{\rhmcsOosQlikeTrangeNSevenFive}{}\renewcommand{\rhmcsOosQlikeTrangeNSevenFive}{3}  % models in the 75 % MCS (of 11)
\providecommand{\rhmcsOosQlikeTrangeSetSevenFive}{}\renewcommand{\rhmcsOosQlikeTrangeSetSevenFive}{the closing-interval HAR benchmark, XGBoost on the full predictor set, and LightGBM on the full predictor set}  % the 75 % MCS, best (last eliminated) first
\providecommand{\rhmcsOosQlikeTrangeStableSevenFive}{}\renewcommand{\rhmcsOosQlikeTrangeStableSevenFive}{yes}  % the 75 % MCS is the same at block lengths (5, 20)
\providecommand{\rhmcsOosQlikeTrangePMaxOutSevenFive}{}\renewcommand{\rhmcsOosQlikeTrangePMaxOutSevenFive}{0.034}  % largest MCS p-value outside the 75 % set
\providecommand{\rhmcsOosQlikeTrangeNSevenFiveLo}{}\renewcommand{\rhmcsOosQlikeTrangeNSevenFiveLo}{3}  % ... at block length 5
\providecommand{\rhmcsOosQlikeTrangeSetSevenFiveLo}{}\renewcommand{\rhmcsOosQlikeTrangeSetSevenFiveLo}{the closing-interval HAR benchmark, XGBoost on the full predictor set, and LightGBM on the full predictor set}  % ... at block length 5
\providecommand{\rhmcsOosQlikeTrangeNSevenFiveHi}{}\renewcommand{\rhmcsOosQlikeTrangeNSevenFiveHi}{3}  % ... at block length 20
\providecommand{\rhmcsOosQlikeTrangeSetSevenFiveHi}{}\renewcommand{\rhmcsOosQlikeTrangeSetSevenFiveHi}{the closing-interval HAR benchmark, XGBoost on the full predictor set, and LightGBM on the full predictor set}  % ... at block length 20
\providecommand{\rhmcsOosQlikeTrangeBest}{}\renewcommand{\rhmcsOosQlikeTrangeBest}{the closing-interval HAR benchmark}  % the last model left (MCS p-value 1)
\providecommand{\rhmcsOosMseTmaxPLinHar}{}\renewcommand{\rhmcsOosMseTmaxPLinHar}{0.56}  % MCS p-value, 16:00 bar, T_oos, Mse, Tmax: Linear, HAR
\providecommand{\rhmcsOosMseTmaxPRidgeFull}{}\renewcommand{\rhmcsOosMseTmaxPRidgeFull}{0.44}  % MCS p-value, 16:00 bar, T_oos, Mse, Tmax: Ridge, Full
\providecommand{\rhmcsOosMseTmaxPLassoFull}{}\renewcommand{\rhmcsOosMseTmaxPLassoFull}{0.39}  % MCS p-value, 16:00 bar, T_oos, Mse, Tmax: Lasso, Full
\providecommand{\rhmcsOosMseTmaxPEnetFull}{}\renewcommand{\rhmcsOosMseTmaxPEnetFull}{0.49}  % MCS p-value, 16:00 bar, T_oos, Mse, Tmax: Elastic net, Full
\providecommand{\rhmcsOosMseTmaxPLgbmHar}{}\renewcommand{\rhmcsOosMseTmaxPLgbmHar}{0.56}  % MCS p-value, 16:00 bar, T_oos, Mse, Tmax: LightGBM, HAR
\providecommand{\rhmcsOosMseTmaxPLgbmFull}{}\renewcommand{\rhmcsOosMseTmaxPLgbmFull}{0.54}  % MCS p-value, 16:00 bar, T_oos, Mse, Tmax: LightGBM, Full
\providecommand{\rhmcsOosMseTmaxPXgbHar}{}\renewcommand{\rhmcsOosMseTmaxPXgbHar}{0.56}  % MCS p-value, 16:00 bar, T_oos, Mse, Tmax: XGBoost, HAR
\providecommand{\rhmcsOosMseTmaxPXgbFull}{}\renewcommand{\rhmcsOosMseTmaxPXgbFull}{0.56}  % MCS p-value, 16:00 bar, T_oos, Mse, Tmax: XGBoost, Full
\providecommand{\rhmcsOosMseTmaxPRfHar}{}\renewcommand{\rhmcsOosMseTmaxPRfHar}{1.00}  % MCS p-value, 16:00 bar, T_oos, Mse, Tmax: Random forest, HAR
\providecommand{\rhmcsOosMseTmaxPRfFull}{}\renewcommand{\rhmcsOosMseTmaxPRfFull}{0.56}  % MCS p-value, 16:00 bar, T_oos, Mse, Tmax: Random forest, Full
\providecommand{\rhmcsOosMseTmaxPBench}{}\renewcommand{\rhmcsOosMseTmaxPBench}{0.56}  % MCS p-value, 16:00 bar, T_oos, Mse, Tmax: Closing-interval HAR benchmark
\providecommand{\rhmcsOosMseTmaxNNinety}{}\renewcommand{\rhmcsOosMseTmaxNNinety}{11}  % models in the 90 % MCS (of 11)
\providecommand{\rhmcsOosMseTmaxSetNinety}{}\renewcommand{\rhmcsOosMseTmaxSetNinety}{the random forest on the HAR benchmark, XGBoost on the HAR benchmark, LightGBM on the HAR benchmark, the random forest on the full predictor set, the linear models on the HAR benchmark, the closing-interval HAR benchmark, XGBoost on the full predictor set, LightGBM on the full predictor set, the elastic net on the full predictor set, ridge regression on the full predictor set, and the lasso on the full predictor set}  % the 90 % MCS, best (last eliminated) first
\providecommand{\rhmcsOosMseTmaxStableNinety}{}\renewcommand{\rhmcsOosMseTmaxStableNinety}{yes}  % the 90 % MCS is the same at block lengths (5, 20)
\providecommand{\rhmcsOosMseTmaxPMaxOutNinety}{}\renewcommand{\rhmcsOosMseTmaxPMaxOutNinety}{--}  % largest MCS p-value outside the 90 % set
\providecommand{\rhmcsOosMseTmaxNNinetyLo}{}\renewcommand{\rhmcsOosMseTmaxNNinetyLo}{11}  % ... at block length 5
\providecommand{\rhmcsOosMseTmaxSetNinetyLo}{}\renewcommand{\rhmcsOosMseTmaxSetNinetyLo}{the random forest on the HAR benchmark, XGBoost on the HAR benchmark, LightGBM on the HAR benchmark, the linear models on the HAR benchmark, the random forest on the full predictor set, the closing-interval HAR benchmark, XGBoost on the full predictor set, LightGBM on the full predictor set, the elastic net on the full predictor set, ridge regression on the full predictor set, and the lasso on the full predictor set}  % ... at block length 5
\providecommand{\rhmcsOosMseTmaxNNinetyHi}{}\renewcommand{\rhmcsOosMseTmaxNNinetyHi}{11}  % ... at block length 20
\providecommand{\rhmcsOosMseTmaxSetNinetyHi}{}\renewcommand{\rhmcsOosMseTmaxSetNinetyHi}{the random forest on the HAR benchmark, XGBoost on the HAR benchmark, LightGBM on the HAR benchmark, the linear models on the HAR benchmark, the random forest on the full predictor set, the closing-interval HAR benchmark, XGBoost on the full predictor set, LightGBM on the full predictor set, the elastic net on the full predictor set, ridge regression on the full predictor set, and the lasso on the full predictor set}  % ... at block length 20
\providecommand{\rhmcsOosMseTmaxNSevenFive}{}\renewcommand{\rhmcsOosMseTmaxNSevenFive}{11}  % models in the 75 % MCS (of 11)
\providecommand{\rhmcsOosMseTmaxSetSevenFive}{}\renewcommand{\rhmcsOosMseTmaxSetSevenFive}{the random forest on the HAR benchmark, XGBoost on the HAR benchmark, LightGBM on the HAR benchmark, the random forest on the full predictor set, the linear models on the HAR benchmark, the closing-interval HAR benchmark, XGBoost on the full predictor set, LightGBM on the full predictor set, the elastic net on the full predictor set, ridge regression on the full predictor set, and the lasso on the full predictor set}  % the 75 % MCS, best (last eliminated) first
\providecommand{\rhmcsOosMseTmaxStableSevenFive}{}\renewcommand{\rhmcsOosMseTmaxStableSevenFive}{no}  % the 75 % MCS is the same at block lengths (5, 20)
\providecommand{\rhmcsOosMseTmaxPMaxOutSevenFive}{}\renewcommand{\rhmcsOosMseTmaxPMaxOutSevenFive}{--}  % largest MCS p-value outside the 75 % set
\providecommand{\rhmcsOosMseTmaxNSevenFiveLo}{}\renewcommand{\rhmcsOosMseTmaxNSevenFiveLo}{10}  % ... at block length 5
\providecommand{\rhmcsOosMseTmaxSetSevenFiveLo}{}\renewcommand{\rhmcsOosMseTmaxSetSevenFiveLo}{the random forest on the HAR benchmark, XGBoost on the HAR benchmark, LightGBM on the HAR benchmark, the linear models on the HAR benchmark, the random forest on the full predictor set, the closing-interval HAR benchmark, XGBoost on the full predictor set, LightGBM on the full predictor set, the elastic net on the full predictor set, and ridge regression on the full predictor set}  % ... at block length 5
\providecommand{\rhmcsOosMseTmaxNSevenFiveHi}{}\renewcommand{\rhmcsOosMseTmaxNSevenFiveHi}{11}  % ... at block length 20
\providecommand{\rhmcsOosMseTmaxSetSevenFiveHi}{}\renewcommand{\rhmcsOosMseTmaxSetSevenFiveHi}{the random forest on the HAR benchmark, XGBoost on the HAR benchmark, LightGBM on the HAR benchmark, the linear models on the HAR benchmark, the random forest on the full predictor set, the closing-interval HAR benchmark, XGBoost on the full predictor set, LightGBM on the full predictor set, the elastic net on the full predictor set, ridge regression on the full predictor set, and the lasso on the full predictor set}  % ... at block length 20
\providecommand{\rhmcsOosMseTmaxBest}{}\renewcommand{\rhmcsOosMseTmaxBest}{the random forest on the HAR benchmark}  % the last model left (MCS p-value 1)
\providecommand{\rhmcsOosMseTrangePLinHar}{}\renewcommand{\rhmcsOosMseTrangePLinHar}{0.56}  % MCS p-value, 16:00 bar, T_oos, Mse, Trange: Linear, HAR
\providecommand{\rhmcsOosMseTrangePRidgeFull}{}\renewcommand{\rhmcsOosMseTrangePRidgeFull}{0.56}  % MCS p-value, 16:00 bar, T_oos, Mse, Trange: Ridge, Full
\providecommand{\rhmcsOosMseTrangePLassoFull}{}\renewcommand{\rhmcsOosMseTrangePLassoFull}{0.56}  % MCS p-value, 16:00 bar, T_oos, Mse, Trange: Lasso, Full
\providecommand{\rhmcsOosMseTrangePEnetFull}{}\renewcommand{\rhmcsOosMseTrangePEnetFull}{0.56}  % MCS p-value, 16:00 bar, T_oos, Mse, Trange: Elastic net, Full
\providecommand{\rhmcsOosMseTrangePLgbmHar}{}\renewcommand{\rhmcsOosMseTrangePLgbmHar}{0.56}  % MCS p-value, 16:00 bar, T_oos, Mse, Trange: LightGBM, HAR
\providecommand{\rhmcsOosMseTrangePLgbmFull}{}\renewcommand{\rhmcsOosMseTrangePLgbmFull}{0.56}  % MCS p-value, 16:00 bar, T_oos, Mse, Trange: LightGBM, Full
\providecommand{\rhmcsOosMseTrangePXgbHar}{}\renewcommand{\rhmcsOosMseTrangePXgbHar}{0.56}  % MCS p-value, 16:00 bar, T_oos, Mse, Trange: XGBoost, HAR
\providecommand{\rhmcsOosMseTrangePXgbFull}{}\renewcommand{\rhmcsOosMseTrangePXgbFull}{0.56}  % MCS p-value, 16:00 bar, T_oos, Mse, Trange: XGBoost, Full
\providecommand{\rhmcsOosMseTrangePRfHar}{}\renewcommand{\rhmcsOosMseTrangePRfHar}{1.00}  % MCS p-value, 16:00 bar, T_oos, Mse, Trange: Random forest, HAR
\providecommand{\rhmcsOosMseTrangePRfFull}{}\renewcommand{\rhmcsOosMseTrangePRfFull}{0.56}  % MCS p-value, 16:00 bar, T_oos, Mse, Trange: Random forest, Full
\providecommand{\rhmcsOosMseTrangePBench}{}\renewcommand{\rhmcsOosMseTrangePBench}{0.56}  % MCS p-value, 16:00 bar, T_oos, Mse, Trange: Closing-interval HAR benchmark
\providecommand{\rhmcsOosMseTrangeNNinety}{}\renewcommand{\rhmcsOosMseTrangeNNinety}{11}  % models in the 90 % MCS (of 11)
\providecommand{\rhmcsOosMseTrangeSetNinety}{}\renewcommand{\rhmcsOosMseTrangeSetNinety}{the random forest on the HAR benchmark, XGBoost on the HAR benchmark, the linear models on the HAR benchmark, the random forest on the full predictor set, XGBoost on the full predictor set, LightGBM on the full predictor set, the elastic net on the full predictor set, ridge regression on the full predictor set, the lasso on the full predictor set, the closing-interval HAR benchmark, and LightGBM on the HAR benchmark}  % the 90 % MCS, best (last eliminated) first
\providecommand{\rhmcsOosMseTrangeStableNinety}{}\renewcommand{\rhmcsOosMseTrangeStableNinety}{yes}  % the 90 % MCS is the same at block lengths (5, 20)
\providecommand{\rhmcsOosMseTrangePMaxOutNinety}{}\renewcommand{\rhmcsOosMseTrangePMaxOutNinety}{--}  % largest MCS p-value outside the 90 % set
\providecommand{\rhmcsOosMseTrangeNNinetyLo}{}\renewcommand{\rhmcsOosMseTrangeNNinetyLo}{11}  % ... at block length 5
\providecommand{\rhmcsOosMseTrangeSetNinetyLo}{}\renewcommand{\rhmcsOosMseTrangeSetNinetyLo}{the random forest on the HAR benchmark, XGBoost on the HAR benchmark, the linear models on the HAR benchmark, XGBoost on the full predictor set, the random forest on the full predictor set, the elastic net on the full predictor set, LightGBM on the full predictor set, LightGBM on the HAR benchmark, ridge regression on the full predictor set, the lasso on the full predictor set, and the closing-interval HAR benchmark}  % ... at block length 5
\providecommand{\rhmcsOosMseTrangeNNinetyHi}{}\renewcommand{\rhmcsOosMseTrangeNNinetyHi}{11}  % ... at block length 20
\providecommand{\rhmcsOosMseTrangeSetNinetyHi}{}\renewcommand{\rhmcsOosMseTrangeSetNinetyHi}{the random forest on the HAR benchmark, the linear models on the HAR benchmark, the random forest on the full predictor set, XGBoost on the HAR benchmark, XGBoost on the full predictor set, LightGBM on the full predictor set, the elastic net on the full predictor set, ridge regression on the full predictor set, the lasso on the full predictor set, the closing-interval HAR benchmark, and LightGBM on the HAR benchmark}  % ... at block length 20
\providecommand{\rhmcsOosMseTrangeNSevenFive}{}\renewcommand{\rhmcsOosMseTrangeNSevenFive}{11}  % models in the 75 % MCS (of 11)
\providecommand{\rhmcsOosMseTrangeSetSevenFive}{}\renewcommand{\rhmcsOosMseTrangeSetSevenFive}{the random forest on the HAR benchmark, XGBoost on the HAR benchmark, the linear models on the HAR benchmark, the random forest on the full predictor set, XGBoost on the full predictor set, LightGBM on the full predictor set, the elastic net on the full predictor set, ridge regression on the full predictor set, the lasso on the full predictor set, the closing-interval HAR benchmark, and LightGBM on the HAR benchmark}  % the 75 % MCS, best (last eliminated) first
\providecommand{\rhmcsOosMseTrangeStableSevenFive}{}\renewcommand{\rhmcsOosMseTrangeStableSevenFive}{yes}  % the 75 % MCS is the same at block lengths (5, 20)
\providecommand{\rhmcsOosMseTrangePMaxOutSevenFive}{}\renewcommand{\rhmcsOosMseTrangePMaxOutSevenFive}{--}  % largest MCS p-value outside the 75 % set
\providecommand{\rhmcsOosMseTrangeNSevenFiveLo}{}\renewcommand{\rhmcsOosMseTrangeNSevenFiveLo}{11}  % ... at block length 5
\providecommand{\rhmcsOosMseTrangeSetSevenFiveLo}{}\renewcommand{\rhmcsOosMseTrangeSetSevenFiveLo}{the random forest on the HAR benchmark, XGBoost on the HAR benchmark, the linear models on the HAR benchmark, XGBoost on the full predictor set, the random forest on the full predictor set, the elastic net on the full predictor set, LightGBM on the full predictor set, LightGBM on the HAR benchmark, ridge regression on the full predictor set, the lasso on the full predictor set, and the closing-interval HAR benchmark}  % ... at block length 5
\providecommand{\rhmcsOosMseTrangeNSevenFiveHi}{}\renewcommand{\rhmcsOosMseTrangeNSevenFiveHi}{11}  % ... at block length 20
\providecommand{\rhmcsOosMseTrangeSetSevenFiveHi}{}\renewcommand{\rhmcsOosMseTrangeSetSevenFiveHi}{the random forest on the HAR benchmark, the linear models on the HAR benchmark, the random forest on the full predictor set, XGBoost on the HAR benchmark, XGBoost on the full predictor set, LightGBM on the full predictor set, the elastic net on the full predictor set, ridge regression on the full predictor set, the lasso on the full predictor set, the closing-interval HAR benchmark, and LightGBM on the HAR benchmark}  % ... at block length 20
\providecommand{\rhmcsOosMseTrangeBest}{}\renewcommand{\rhmcsOosMseTrangeBest}{the random forest on the HAR benchmark}  % the last model left (MCS p-value 1)
\providecommand{\rhmcsChQlikeTmaxPLinHar}{}\renewcommand{\rhmcsChQlikeTmaxPLinHar}{0.050}  % MCS p-value, 16:00 bar, the benchmark's span, Qlike, Tmax: Linear, HAR
\providecommand{\rhmcsChQlikeTmaxPRidgeFull}{}\renewcommand{\rhmcsChQlikeTmaxPRidgeFull}{0.050}  % MCS p-value, 16:00 bar, the benchmark's span, Qlike, Tmax: Ridge, Full
\providecommand{\rhmcsChQlikeTmaxPLassoFull}{}\renewcommand{\rhmcsChQlikeTmaxPLassoFull}{0.050}  % MCS p-value, 16:00 bar, the benchmark's span, Qlike, Tmax: Lasso, Full
\providecommand{\rhmcsChQlikeTmaxPEnetFull}{}\renewcommand{\rhmcsChQlikeTmaxPEnetFull}{0.102}  % MCS p-value, 16:00 bar, the benchmark's span, Qlike, Tmax: Elastic net, Full
\providecommand{\rhmcsChQlikeTmaxPLgbmHar}{}\renewcommand{\rhmcsChQlikeTmaxPLgbmHar}{0.102}  % MCS p-value, 16:00 bar, the benchmark's span, Qlike, Tmax: LightGBM, HAR
\providecommand{\rhmcsChQlikeTmaxPLgbmFull}{}\renewcommand{\rhmcsChQlikeTmaxPLgbmFull}{0.16}  % MCS p-value, 16:00 bar, the benchmark's span, Qlike, Tmax: LightGBM, Full
\providecommand{\rhmcsChQlikeTmaxPXgbHar}{}\renewcommand{\rhmcsChQlikeTmaxPXgbHar}{0.102}  % MCS p-value, 16:00 bar, the benchmark's span, Qlike, Tmax: XGBoost, HAR
\providecommand{\rhmcsChQlikeTmaxPXgbFull}{}\renewcommand{\rhmcsChQlikeTmaxPXgbFull}{0.22}  % MCS p-value, 16:00 bar, the benchmark's span, Qlike, Tmax: XGBoost, Full
\providecommand{\rhmcsChQlikeTmaxPRfHar}{}\renewcommand{\rhmcsChQlikeTmaxPRfHar}{0.050}  % MCS p-value, 16:00 bar, the benchmark's span, Qlike, Tmax: Random forest, HAR
\providecommand{\rhmcsChQlikeTmaxPRfFull}{}\renewcommand{\rhmcsChQlikeTmaxPRfFull}{0.102}  % MCS p-value, 16:00 bar, the benchmark's span, Qlike, Tmax: Random forest, Full
\providecommand{\rhmcsChQlikeTmaxPBench}{}\renewcommand{\rhmcsChQlikeTmaxPBench}{1.00}  % MCS p-value, 16:00 bar, the benchmark's span, Qlike, Tmax: Closing-interval HAR benchmark
\providecommand{\rhmcsChQlikeTmaxNNinety}{}\renewcommand{\rhmcsChQlikeTmaxNNinety}{7}  % models in the 90 % MCS (of 11)
\providecommand{\rhmcsChQlikeTmaxSetNinety}{}\renewcommand{\rhmcsChQlikeTmaxSetNinety}{the closing-interval HAR benchmark, XGBoost on the full predictor set, LightGBM on the full predictor set, LightGBM on the HAR benchmark, XGBoost on the HAR benchmark, the random forest on the full predictor set, and the elastic net on the full predictor set}  % the 90 % MCS, best (last eliminated) first
\providecommand{\rhmcsChQlikeTmaxStableNinety}{}\renewcommand{\rhmcsChQlikeTmaxStableNinety}{no}  % the 90 % MCS is the same at block lengths (5, 20)
\providecommand{\rhmcsChQlikeTmaxPMaxOutNinety}{}\renewcommand{\rhmcsChQlikeTmaxPMaxOutNinety}{0.050}  % largest MCS p-value outside the 90 % set
\providecommand{\rhmcsChQlikeTmaxNNinetyLo}{}\renewcommand{\rhmcsChQlikeTmaxNNinetyLo}{3}  % ... at block length 5
\providecommand{\rhmcsChQlikeTmaxSetNinetyLo}{}\renewcommand{\rhmcsChQlikeTmaxSetNinetyLo}{the closing-interval HAR benchmark, XGBoost on the full predictor set, and LightGBM on the full predictor set}  % ... at block length 5
\providecommand{\rhmcsChQlikeTmaxNNinetyHi}{}\renewcommand{\rhmcsChQlikeTmaxNNinetyHi}{7}  % ... at block length 20
\providecommand{\rhmcsChQlikeTmaxSetNinetyHi}{}\renewcommand{\rhmcsChQlikeTmaxSetNinetyHi}{the closing-interval HAR benchmark, XGBoost on the full predictor set, LightGBM on the full predictor set, LightGBM on the HAR benchmark, XGBoost on the HAR benchmark, the random forest on the full predictor set, and the elastic net on the full predictor set}  % ... at block length 20
\providecommand{\rhmcsChQlikeTmaxNSevenFive}{}\renewcommand{\rhmcsChQlikeTmaxNSevenFive}{1}  % models in the 75 % MCS (of 11)
\providecommand{\rhmcsChQlikeTmaxSetSevenFive}{}\renewcommand{\rhmcsChQlikeTmaxSetSevenFive}{the closing-interval HAR benchmark}  % the 75 % MCS, best (last eliminated) first
\providecommand{\rhmcsChQlikeTmaxStableSevenFive}{}\renewcommand{\rhmcsChQlikeTmaxStableSevenFive}{yes}  % the 75 % MCS is the same at block lengths (5, 20)
\providecommand{\rhmcsChQlikeTmaxPMaxOutSevenFive}{}\renewcommand{\rhmcsChQlikeTmaxPMaxOutSevenFive}{0.22}  % largest MCS p-value outside the 75 % set
\providecommand{\rhmcsChQlikeTmaxNSevenFiveLo}{}\renewcommand{\rhmcsChQlikeTmaxNSevenFiveLo}{1}  % ... at block length 5
\providecommand{\rhmcsChQlikeTmaxSetSevenFiveLo}{}\renewcommand{\rhmcsChQlikeTmaxSetSevenFiveLo}{the closing-interval HAR benchmark}  % ... at block length 5
\providecommand{\rhmcsChQlikeTmaxNSevenFiveHi}{}\renewcommand{\rhmcsChQlikeTmaxNSevenFiveHi}{1}  % ... at block length 20
\providecommand{\rhmcsChQlikeTmaxSetSevenFiveHi}{}\renewcommand{\rhmcsChQlikeTmaxSetSevenFiveHi}{the closing-interval HAR benchmark}  % ... at block length 20
\providecommand{\rhmcsChQlikeTmaxBest}{}\renewcommand{\rhmcsChQlikeTmaxBest}{the closing-interval HAR benchmark}  % the last model left (MCS p-value 1)
\providecommand{\rhmcsChQlikeTrangePLinHar}{}\renewcommand{\rhmcsChQlikeTrangePLinHar}{\ensuremath{<}0.001}  % MCS p-value, 16:00 bar, the benchmark's span, Qlike, Trange: Linear, HAR
\providecommand{\rhmcsChQlikeTrangePRidgeFull}{}\renewcommand{\rhmcsChQlikeTrangePRidgeFull}{0.008}  % MCS p-value, 16:00 bar, the benchmark's span, Qlike, Trange: Ridge, Full
\providecommand{\rhmcsChQlikeTrangePLassoFull}{}\renewcommand{\rhmcsChQlikeTrangePLassoFull}{0.007}  % MCS p-value, 16:00 bar, the benchmark's span, Qlike, Trange: Lasso, Full
\providecommand{\rhmcsChQlikeTrangePEnetFull}{}\renewcommand{\rhmcsChQlikeTrangePEnetFull}{0.008}  % MCS p-value, 16:00 bar, the benchmark's span, Qlike, Trange: Elastic net, Full
\providecommand{\rhmcsChQlikeTrangePLgbmHar}{}\renewcommand{\rhmcsChQlikeTrangePLgbmHar}{0.008}  % MCS p-value, 16:00 bar, the benchmark's span, Qlike, Trange: LightGBM, HAR
\providecommand{\rhmcsChQlikeTrangePLgbmFull}{}\renewcommand{\rhmcsChQlikeTrangePLgbmFull}{0.31}  % MCS p-value, 16:00 bar, the benchmark's span, Qlike, Trange: LightGBM, Full
\providecommand{\rhmcsChQlikeTrangePXgbHar}{}\renewcommand{\rhmcsChQlikeTrangePXgbHar}{0.007}  % MCS p-value, 16:00 bar, the benchmark's span, Qlike, Trange: XGBoost, HAR
\providecommand{\rhmcsChQlikeTrangePXgbFull}{}\renewcommand{\rhmcsChQlikeTrangePXgbFull}{0.31}  % MCS p-value, 16:00 bar, the benchmark's span, Qlike, Trange: XGBoost, Full
\providecommand{\rhmcsChQlikeTrangePRfHar}{}\renewcommand{\rhmcsChQlikeTrangePRfHar}{\ensuremath{<}0.001}  % MCS p-value, 16:00 bar, the benchmark's span, Qlike, Trange: Random forest, HAR
\providecommand{\rhmcsChQlikeTrangePRfFull}{}\renewcommand{\rhmcsChQlikeTrangePRfFull}{\ensuremath{<}0.001}  % MCS p-value, 16:00 bar, the benchmark's span, Qlike, Trange: Random forest, Full
\providecommand{\rhmcsChQlikeTrangePBench}{}\renewcommand{\rhmcsChQlikeTrangePBench}{1.00}  % MCS p-value, 16:00 bar, the benchmark's span, Qlike, Trange: Closing-interval HAR benchmark
\providecommand{\rhmcsChQlikeTrangeNNinety}{}\renewcommand{\rhmcsChQlikeTrangeNNinety}{3}  % models in the 90 % MCS (of 11)
\providecommand{\rhmcsChQlikeTrangeSetNinety}{}\renewcommand{\rhmcsChQlikeTrangeSetNinety}{the closing-interval HAR benchmark, XGBoost on the full predictor set, and LightGBM on the full predictor set}  % the 90 % MCS, best (last eliminated) first
\providecommand{\rhmcsChQlikeTrangeStableNinety}{}\renewcommand{\rhmcsChQlikeTrangeStableNinety}{yes}  % the 90 % MCS is the same at block lengths (5, 20)
\providecommand{\rhmcsChQlikeTrangePMaxOutNinety}{}\renewcommand{\rhmcsChQlikeTrangePMaxOutNinety}{0.008}  % largest MCS p-value outside the 90 % set
\providecommand{\rhmcsChQlikeTrangeNNinetyLo}{}\renewcommand{\rhmcsChQlikeTrangeNNinetyLo}{3}  % ... at block length 5
\providecommand{\rhmcsChQlikeTrangeSetNinetyLo}{}\renewcommand{\rhmcsChQlikeTrangeSetNinetyLo}{the closing-interval HAR benchmark, XGBoost on the full predictor set, and LightGBM on the full predictor set}  % ... at block length 5
\providecommand{\rhmcsChQlikeTrangeNNinetyHi}{}\renewcommand{\rhmcsChQlikeTrangeNNinetyHi}{3}  % ... at block length 20
\providecommand{\rhmcsChQlikeTrangeSetNinetyHi}{}\renewcommand{\rhmcsChQlikeTrangeSetNinetyHi}{the closing-interval HAR benchmark, XGBoost on the full predictor set, and LightGBM on the full predictor set}  % ... at block length 20
\providecommand{\rhmcsChQlikeTrangeNSevenFive}{}\renewcommand{\rhmcsChQlikeTrangeNSevenFive}{3}  % models in the 75 % MCS (of 11)
\providecommand{\rhmcsChQlikeTrangeSetSevenFive}{}\renewcommand{\rhmcsChQlikeTrangeSetSevenFive}{the closing-interval HAR benchmark, XGBoost on the full predictor set, and LightGBM on the full predictor set}  % the 75 % MCS, best (last eliminated) first
\providecommand{\rhmcsChQlikeTrangeStableSevenFive}{}\renewcommand{\rhmcsChQlikeTrangeStableSevenFive}{yes}  % the 75 % MCS is the same at block lengths (5, 20)
\providecommand{\rhmcsChQlikeTrangePMaxOutSevenFive}{}\renewcommand{\rhmcsChQlikeTrangePMaxOutSevenFive}{0.008}  % largest MCS p-value outside the 75 % set
\providecommand{\rhmcsChQlikeTrangeNSevenFiveLo}{}\renewcommand{\rhmcsChQlikeTrangeNSevenFiveLo}{3}  % ... at block length 5
\providecommand{\rhmcsChQlikeTrangeSetSevenFiveLo}{}\renewcommand{\rhmcsChQlikeTrangeSetSevenFiveLo}{the closing-interval HAR benchmark, XGBoost on the full predictor set, and LightGBM on the full predictor set}  % ... at block length 5
\providecommand{\rhmcsChQlikeTrangeNSevenFiveHi}{}\renewcommand{\rhmcsChQlikeTrangeNSevenFiveHi}{3}  % ... at block length 20
\providecommand{\rhmcsChQlikeTrangeSetSevenFiveHi}{}\renewcommand{\rhmcsChQlikeTrangeSetSevenFiveHi}{the closing-interval HAR benchmark, XGBoost on the full predictor set, and LightGBM on the full predictor set}  % ... at block length 20
\providecommand{\rhmcsChQlikeTrangeBest}{}\renewcommand{\rhmcsChQlikeTrangeBest}{the closing-interval HAR benchmark}  % the last model left (MCS p-value 1)
\providecommand{\rhmcsChMseTmaxPLinHar}{}\renewcommand{\rhmcsChMseTmaxPLinHar}{0.59}  % MCS p-value, 16:00 bar, the benchmark's span, Mse, Tmax: Linear, HAR
\providecommand{\rhmcsChMseTmaxPRidgeFull}{}\renewcommand{\rhmcsChMseTmaxPRidgeFull}{0.47}  % MCS p-value, 16:00 bar, the benchmark's span, Mse, Tmax: Ridge, Full
\providecommand{\rhmcsChMseTmaxPLassoFull}{}\renewcommand{\rhmcsChMseTmaxPLassoFull}{0.42}  % MCS p-value, 16:00 bar, the benchmark's span, Mse, Tmax: Lasso, Full
\providecommand{\rhmcsChMseTmaxPEnetFull}{}\renewcommand{\rhmcsChMseTmaxPEnetFull}{0.53}  % MCS p-value, 16:00 bar, the benchmark's span, Mse, Tmax: Elastic net, Full
\providecommand{\rhmcsChMseTmaxPLgbmHar}{}\renewcommand{\rhmcsChMseTmaxPLgbmHar}{0.59}  % MCS p-value, 16:00 bar, the benchmark's span, Mse, Tmax: LightGBM, HAR
\providecommand{\rhmcsChMseTmaxPLgbmFull}{}\renewcommand{\rhmcsChMseTmaxPLgbmFull}{0.54}  % MCS p-value, 16:00 bar, the benchmark's span, Mse, Tmax: LightGBM, Full
\providecommand{\rhmcsChMseTmaxPXgbHar}{}\renewcommand{\rhmcsChMseTmaxPXgbHar}{0.59}  % MCS p-value, 16:00 bar, the benchmark's span, Mse, Tmax: XGBoost, HAR
\providecommand{\rhmcsChMseTmaxPXgbFull}{}\renewcommand{\rhmcsChMseTmaxPXgbFull}{0.58}  % MCS p-value, 16:00 bar, the benchmark's span, Mse, Tmax: XGBoost, Full
\providecommand{\rhmcsChMseTmaxPRfHar}{}\renewcommand{\rhmcsChMseTmaxPRfHar}{1.00}  % MCS p-value, 16:00 bar, the benchmark's span, Mse, Tmax: Random forest, HAR
\providecommand{\rhmcsChMseTmaxPRfFull}{}\renewcommand{\rhmcsChMseTmaxPRfFull}{0.59}  % MCS p-value, 16:00 bar, the benchmark's span, Mse, Tmax: Random forest, Full
\providecommand{\rhmcsChMseTmaxPBench}{}\renewcommand{\rhmcsChMseTmaxPBench}{0.58}  % MCS p-value, 16:00 bar, the benchmark's span, Mse, Tmax: Closing-interval HAR benchmark
\providecommand{\rhmcsChMseTmaxNNinety}{}\renewcommand{\rhmcsChMseTmaxNNinety}{11}  % models in the 90 % MCS (of 11)
\providecommand{\rhmcsChMseTmaxSetNinety}{}\renewcommand{\rhmcsChMseTmaxSetNinety}{the random forest on the HAR benchmark, XGBoost on the HAR benchmark, LightGBM on the HAR benchmark, the linear models on the HAR benchmark, the random forest on the full predictor set, the closing-interval HAR benchmark, XGBoost on the full predictor set, LightGBM on the full predictor set, the elastic net on the full predictor set, ridge regression on the full predictor set, and the lasso on the full predictor set}  % the 90 % MCS, best (last eliminated) first
\providecommand{\rhmcsChMseTmaxStableNinety}{}\renewcommand{\rhmcsChMseTmaxStableNinety}{yes}  % the 90 % MCS is the same at block lengths (5, 20)
\providecommand{\rhmcsChMseTmaxPMaxOutNinety}{}\renewcommand{\rhmcsChMseTmaxPMaxOutNinety}{--}  % largest MCS p-value outside the 90 % set
\providecommand{\rhmcsChMseTmaxNNinetyLo}{}\renewcommand{\rhmcsChMseTmaxNNinetyLo}{11}  % ... at block length 5
\providecommand{\rhmcsChMseTmaxSetNinetyLo}{}\renewcommand{\rhmcsChMseTmaxSetNinetyLo}{the random forest on the HAR benchmark, XGBoost on the HAR benchmark, LightGBM on the HAR benchmark, the linear models on the HAR benchmark, the random forest on the full predictor set, the closing-interval HAR benchmark, XGBoost on the full predictor set, LightGBM on the full predictor set, the elastic net on the full predictor set, ridge regression on the full predictor set, and the lasso on the full predictor set}  % ... at block length 5
\providecommand{\rhmcsChMseTmaxNNinetyHi}{}\renewcommand{\rhmcsChMseTmaxNNinetyHi}{11}  % ... at block length 20
\providecommand{\rhmcsChMseTmaxSetNinetyHi}{}\renewcommand{\rhmcsChMseTmaxSetNinetyHi}{the random forest on the HAR benchmark, XGBoost on the HAR benchmark, LightGBM on the HAR benchmark, the linear models on the HAR benchmark, the random forest on the full predictor set, the closing-interval HAR benchmark, XGBoost on the full predictor set, LightGBM on the full predictor set, the elastic net on the full predictor set, ridge regression on the full predictor set, and the lasso on the full predictor set}  % ... at block length 20
\providecommand{\rhmcsChMseTmaxNSevenFive}{}\renewcommand{\rhmcsChMseTmaxNSevenFive}{11}  % models in the 75 % MCS (of 11)
\providecommand{\rhmcsChMseTmaxSetSevenFive}{}\renewcommand{\rhmcsChMseTmaxSetSevenFive}{the random forest on the HAR benchmark, XGBoost on the HAR benchmark, LightGBM on the HAR benchmark, the linear models on the HAR benchmark, the random forest on the full predictor set, the closing-interval HAR benchmark, XGBoost on the full predictor set, LightGBM on the full predictor set, the elastic net on the full predictor set, ridge regression on the full predictor set, and the lasso on the full predictor set}  % the 75 % MCS, best (last eliminated) first
\providecommand{\rhmcsChMseTmaxStableSevenFive}{}\renewcommand{\rhmcsChMseTmaxStableSevenFive}{yes}  % the 75 % MCS is the same at block lengths (5, 20)
\providecommand{\rhmcsChMseTmaxPMaxOutSevenFive}{}\renewcommand{\rhmcsChMseTmaxPMaxOutSevenFive}{--}  % largest MCS p-value outside the 75 % set
\providecommand{\rhmcsChMseTmaxNSevenFiveLo}{}\renewcommand{\rhmcsChMseTmaxNSevenFiveLo}{11}  % ... at block length 5
\providecommand{\rhmcsChMseTmaxSetSevenFiveLo}{}\renewcommand{\rhmcsChMseTmaxSetSevenFiveLo}{the random forest on the HAR benchmark, XGBoost on the HAR benchmark, LightGBM on the HAR benchmark, the linear models on the HAR benchmark, the random forest on the full predictor set, the closing-interval HAR benchmark, XGBoost on the full predictor set, LightGBM on the full predictor set, the elastic net on the full predictor set, ridge regression on the full predictor set, and the lasso on the full predictor set}  % ... at block length 5
\providecommand{\rhmcsChMseTmaxNSevenFiveHi}{}\renewcommand{\rhmcsChMseTmaxNSevenFiveHi}{11}  % ... at block length 20
\providecommand{\rhmcsChMseTmaxSetSevenFiveHi}{}\renewcommand{\rhmcsChMseTmaxSetSevenFiveHi}{the random forest on the HAR benchmark, XGBoost on the HAR benchmark, LightGBM on the HAR benchmark, the linear models on the HAR benchmark, the random forest on the full predictor set, the closing-interval HAR benchmark, XGBoost on the full predictor set, LightGBM on the full predictor set, the elastic net on the full predictor set, ridge regression on the full predictor set, and the lasso on the full predictor set}  % ... at block length 20
\providecommand{\rhmcsChMseTmaxBest}{}\renewcommand{\rhmcsChMseTmaxBest}{the random forest on the HAR benchmark}  % the last model left (MCS p-value 1)
\providecommand{\rhmcsChMseTrangePLinHar}{}\renewcommand{\rhmcsChMseTrangePLinHar}{0.60}  % MCS p-value, 16:00 bar, the benchmark's span, Mse, Trange: Linear, HAR
\providecommand{\rhmcsChMseTrangePRidgeFull}{}\renewcommand{\rhmcsChMseTrangePRidgeFull}{0.60}  % MCS p-value, 16:00 bar, the benchmark's span, Mse, Trange: Ridge, Full
\providecommand{\rhmcsChMseTrangePLassoFull}{}\renewcommand{\rhmcsChMseTrangePLassoFull}{0.60}  % MCS p-value, 16:00 bar, the benchmark's span, Mse, Trange: Lasso, Full
\providecommand{\rhmcsChMseTrangePEnetFull}{}\renewcommand{\rhmcsChMseTrangePEnetFull}{0.60}  % MCS p-value, 16:00 bar, the benchmark's span, Mse, Trange: Elastic net, Full
\providecommand{\rhmcsChMseTrangePLgbmHar}{}\renewcommand{\rhmcsChMseTrangePLgbmHar}{0.41}  % MCS p-value, 16:00 bar, the benchmark's span, Mse, Trange: LightGBM, HAR
\providecommand{\rhmcsChMseTrangePLgbmFull}{}\renewcommand{\rhmcsChMseTrangePLgbmFull}{0.60}  % MCS p-value, 16:00 bar, the benchmark's span, Mse, Trange: LightGBM, Full
\providecommand{\rhmcsChMseTrangePXgbHar}{}\renewcommand{\rhmcsChMseTrangePXgbHar}{0.60}  % MCS p-value, 16:00 bar, the benchmark's span, Mse, Trange: XGBoost, HAR
\providecommand{\rhmcsChMseTrangePXgbFull}{}\renewcommand{\rhmcsChMseTrangePXgbFull}{0.60}  % MCS p-value, 16:00 bar, the benchmark's span, Mse, Trange: XGBoost, Full
\providecommand{\rhmcsChMseTrangePRfHar}{}\renewcommand{\rhmcsChMseTrangePRfHar}{1.00}  % MCS p-value, 16:00 bar, the benchmark's span, Mse, Trange: Random forest, HAR
\providecommand{\rhmcsChMseTrangePRfFull}{}\renewcommand{\rhmcsChMseTrangePRfFull}{0.60}  % MCS p-value, 16:00 bar, the benchmark's span, Mse, Trange: Random forest, Full
\providecommand{\rhmcsChMseTrangePBench}{}\renewcommand{\rhmcsChMseTrangePBench}{0.60}  % MCS p-value, 16:00 bar, the benchmark's span, Mse, Trange: Closing-interval HAR benchmark
\providecommand{\rhmcsChMseTrangeNNinety}{}\renewcommand{\rhmcsChMseTrangeNNinety}{11}  % models in the 90 % MCS (of 11)
\providecommand{\rhmcsChMseTrangeSetNinety}{}\renewcommand{\rhmcsChMseTrangeSetNinety}{the random forest on the HAR benchmark, the linear models on the HAR benchmark, XGBoost on the HAR benchmark, the random forest on the full predictor set, XGBoost on the full predictor set, the elastic net on the full predictor set, ridge regression on the full predictor set, the lasso on the full predictor set, LightGBM on the full predictor set, the closing-interval HAR benchmark, and LightGBM on the HAR benchmark}  % the 90 % MCS, best (last eliminated) first
\providecommand{\rhmcsChMseTrangeStableNinety}{}\renewcommand{\rhmcsChMseTrangeStableNinety}{yes}  % the 90 % MCS is the same at block lengths (5, 20)
\providecommand{\rhmcsChMseTrangePMaxOutNinety}{}\renewcommand{\rhmcsChMseTrangePMaxOutNinety}{--}  % largest MCS p-value outside the 90 % set
\providecommand{\rhmcsChMseTrangeNNinetyLo}{}\renewcommand{\rhmcsChMseTrangeNNinetyLo}{11}  % ... at block length 5
\providecommand{\rhmcsChMseTrangeSetNinetyLo}{}\renewcommand{\rhmcsChMseTrangeSetNinetyLo}{the random forest on the HAR benchmark, the linear models on the HAR benchmark, XGBoost on the HAR benchmark, XGBoost on the full predictor set, the random forest on the full predictor set, the elastic net on the full predictor set, ridge regression on the full predictor set, LightGBM on the full predictor set, the lasso on the full predictor set, the closing-interval HAR benchmark, and LightGBM on the HAR benchmark}  % ... at block length 5
\providecommand{\rhmcsChMseTrangeNNinetyHi}{}\renewcommand{\rhmcsChMseTrangeNNinetyHi}{11}  % ... at block length 20
\providecommand{\rhmcsChMseTrangeSetNinetyHi}{}\renewcommand{\rhmcsChMseTrangeSetNinetyHi}{the random forest on the HAR benchmark, the linear models on the HAR benchmark, the random forest on the full predictor set, XGBoost on the full predictor set, the elastic net on the full predictor set, ridge regression on the full predictor set, the lasso on the full predictor set, XGBoost on the HAR benchmark, the closing-interval HAR benchmark, LightGBM on the full predictor set, and LightGBM on the HAR benchmark}  % ... at block length 20
\providecommand{\rhmcsChMseTrangeNSevenFive}{}\renewcommand{\rhmcsChMseTrangeNSevenFive}{11}  % models in the 75 % MCS (of 11)
\providecommand{\rhmcsChMseTrangeSetSevenFive}{}\renewcommand{\rhmcsChMseTrangeSetSevenFive}{the random forest on the HAR benchmark, the linear models on the HAR benchmark, XGBoost on the HAR benchmark, the random forest on the full predictor set, XGBoost on the full predictor set, the elastic net on the full predictor set, ridge regression on the full predictor set, the lasso on the full predictor set, LightGBM on the full predictor set, the closing-interval HAR benchmark, and LightGBM on the HAR benchmark}  % the 75 % MCS, best (last eliminated) first
\providecommand{\rhmcsChMseTrangeStableSevenFive}{}\renewcommand{\rhmcsChMseTrangeStableSevenFive}{yes}  % the 75 % MCS is the same at block lengths (5, 20)
\providecommand{\rhmcsChMseTrangePMaxOutSevenFive}{}\renewcommand{\rhmcsChMseTrangePMaxOutSevenFive}{--}  % largest MCS p-value outside the 75 % set
\providecommand{\rhmcsChMseTrangeNSevenFiveLo}{}\renewcommand{\rhmcsChMseTrangeNSevenFiveLo}{11}  % ... at block length 5
\providecommand{\rhmcsChMseTrangeSetSevenFiveLo}{}\renewcommand{\rhmcsChMseTrangeSetSevenFiveLo}{the random forest on the HAR benchmark, the linear models on the HAR benchmark, XGBoost on the HAR benchmark, XGBoost on the full predictor set, the random forest on the full predictor set, the elastic net on the full predictor set, ridge regression on the full predictor set, LightGBM on the full predictor set, the lasso on the full predictor set, the closing-interval HAR benchmark, and LightGBM on the HAR benchmark}  % ... at block length 5
\providecommand{\rhmcsChMseTrangeNSevenFiveHi}{}\renewcommand{\rhmcsChMseTrangeNSevenFiveHi}{11}  % ... at block length 20
\providecommand{\rhmcsChMseTrangeSetSevenFiveHi}{}\renewcommand{\rhmcsChMseTrangeSetSevenFiveHi}{the random forest on the HAR benchmark, the linear models on the HAR benchmark, the random forest on the full predictor set, XGBoost on the full predictor set, the elastic net on the full predictor set, ridge regression on the full predictor set, the lasso on the full predictor set, XGBoost on the HAR benchmark, the closing-interval HAR benchmark, LightGBM on the full predictor set, and LightGBM on the HAR benchmark}  % ... at block length 20
\providecommand{\rhmcsChMseTrangeBest}{}\renewcommand{\rhmcsChMseTrangeBest}{the random forest on the HAR benchmark}  % the last model left (MCS p-value 1)
\providecommand{\rhmcsBarQAfSpaC}{}\renewcommand{\rhmcsBarQAfSpaC}{\ensuremath{<}0.001}  % SPA p-value (consistent), family BarQAf
\providecommand{\rhmcsBarQAfSpaL}{}\renewcommand{\rhmcsBarQAfSpaL}{\ensuremath{<}0.001}  % ... lower
\providecommand{\rhmcsBarQAfSpaU}{}\renewcommand{\rhmcsBarQAfSpaU}{\ensuremath{<}0.001}  % ... upper
\providecommand{\rhmcsBarQAfNRej}{}\renewcommand{\rhmcsBarQAfNRej}{6}  % StepM rejections at FWER 5 % (of 6)
\providecommand{\rhmcsBarQAfNRejWord}{}\renewcommand{\rhmcsBarQAfNRejWord}{six}  % ... as a word
\providecommand{\rhmcsBarQAfNRaw}{}\renewcommand{\rhmcsBarQAfNRaw}{6}  % single-test bootstrap p < 0.05 (no adjustment)
\providecommand{\rhmcsBarQAfNRawWord}{}\renewcommand{\rhmcsBarQAfNRawWord}{six}  % ... as a word
\providecommand{\rhmcsBarQAfNPos}{}\renewcommand{\rhmcsBarQAfNPos}{6}  % alternatives with a positive estimate (better than the benchmark)
\providecommand{\rhmcsBarQAfPMin}{}\renewcommand{\rhmcsBarQAfPMin}{\ensuremath{<}0.001}  % smallest StepM-adjusted p-value
\providecommand{\rhmcsBarQAfPMax}{}\renewcommand{\rhmcsBarQAfPMax}{\ensuremath{<}0.001}  % largest StepM-adjusted p-value
\providecommand{\rhmcsBarQAfPTwoMin}{}\renewcommand{\rhmcsBarQAfPTwoMin}{\ensuremath{<}0.001}  % smallest two-sided StepM-adjusted p-value
\providecommand{\rhmcsBarQAfPTwoMax}{}\renewcommand{\rhmcsBarQAfPTwoMax}{\ensuremath{<}0.001}  % largest two-sided StepM-adjusted p-value
\providecommand{\rhmcsBarQAfRawMin}{}\renewcommand{\rhmcsBarQAfRawMin}{\ensuremath{<}0.001}  % smallest single-test bootstrap p-value
\providecommand{\rhmcsBarQAfRej}{}\renewcommand{\rhmcsBarQAfRej}{ridge regression on the full predictor set, the lasso on the full predictor set, the elastic net on the full predictor set, LightGBM on the full predictor set, XGBoost on the full predictor set, and the random forest on the full predictor set}  % the alternatives StepM rejects at 5 % (none: none)
\providecommand{\rhmcsBarQAfNRejTwo}{}\renewcommand{\rhmcsBarQAfNRejTwo}{6}  % StepM rejections at FWER 5 %, two-sided
\providecommand{\rhmcsBarQAfNRejTwoWord}{}\renewcommand{\rhmcsBarQAfNRejTwoWord}{six}  % ... as a word
\providecommand{\rhmcsBarQAfNRawTwo}{}\renewcommand{\rhmcsBarQAfNRawTwo}{6}  % single-test bootstrap p < 0.05, two-sided
\providecommand{\rhmcsBarQAfRejTwo}{}\renewcommand{\rhmcsBarQAfRejTwo}{ridge regression on the full predictor set, the lasso on the full predictor set, the elastic net on the full predictor set, LightGBM on the full predictor set, XGBoost on the full predictor set, and the random forest on the full predictor set}  % the alternatives the two-sided StepM rejects at 5 % (none: none)
\providecommand{\rhmcsBarQAfPRidge}{}\renewcommand{\rhmcsBarQAfPRidge}{\ensuremath{<}0.001}  % StepM-adjusted p-value (one-sided), Ridge
\providecommand{\rhmcsBarQAfRawRidge}{}\renewcommand{\rhmcsBarQAfRawRidge}{\ensuremath{<}0.001}  % single-test bootstrap p-value (one-sided), Ridge
\providecommand{\rhmcsBarQAfPTwoRidge}{}\renewcommand{\rhmcsBarQAfPTwoRidge}{\ensuremath{<}0.001}  % StepM-adjusted p-value (two-sided), Ridge
\providecommand{\rhmcsBarQAfRawTwoRidge}{}\renewcommand{\rhmcsBarQAfRawTwoRidge}{\ensuremath{<}0.001}  % single-test bootstrap p-value (two-sided), Ridge
\providecommand{\rhmcsBarQAfPLasso}{}\renewcommand{\rhmcsBarQAfPLasso}{\ensuremath{<}0.001}  % StepM-adjusted p-value (one-sided), Lasso
\providecommand{\rhmcsBarQAfRawLasso}{}\renewcommand{\rhmcsBarQAfRawLasso}{\ensuremath{<}0.001}  % single-test bootstrap p-value (one-sided), Lasso
\providecommand{\rhmcsBarQAfPTwoLasso}{}\renewcommand{\rhmcsBarQAfPTwoLasso}{\ensuremath{<}0.001}  % StepM-adjusted p-value (two-sided), Lasso
\providecommand{\rhmcsBarQAfRawTwoLasso}{}\renewcommand{\rhmcsBarQAfRawTwoLasso}{\ensuremath{<}0.001}  % single-test bootstrap p-value (two-sided), Lasso
\providecommand{\rhmcsBarQAfPEnet}{}\renewcommand{\rhmcsBarQAfPEnet}{\ensuremath{<}0.001}  % StepM-adjusted p-value (one-sided), Elastic net
\providecommand{\rhmcsBarQAfRawEnet}{}\renewcommand{\rhmcsBarQAfRawEnet}{\ensuremath{<}0.001}  % single-test bootstrap p-value (one-sided), Elastic net
\providecommand{\rhmcsBarQAfPTwoEnet}{}\renewcommand{\rhmcsBarQAfPTwoEnet}{\ensuremath{<}0.001}  % StepM-adjusted p-value (two-sided), Elastic net
\providecommand{\rhmcsBarQAfRawTwoEnet}{}\renewcommand{\rhmcsBarQAfRawTwoEnet}{\ensuremath{<}0.001}  % single-test bootstrap p-value (two-sided), Elastic net
\providecommand{\rhmcsBarQAfPLgbm}{}\renewcommand{\rhmcsBarQAfPLgbm}{\ensuremath{<}0.001}  % StepM-adjusted p-value (one-sided), LightGBM
\providecommand{\rhmcsBarQAfRawLgbm}{}\renewcommand{\rhmcsBarQAfRawLgbm}{\ensuremath{<}0.001}  % single-test bootstrap p-value (one-sided), LightGBM
\providecommand{\rhmcsBarQAfPTwoLgbm}{}\renewcommand{\rhmcsBarQAfPTwoLgbm}{\ensuremath{<}0.001}  % StepM-adjusted p-value (two-sided), LightGBM
\providecommand{\rhmcsBarQAfRawTwoLgbm}{}\renewcommand{\rhmcsBarQAfRawTwoLgbm}{\ensuremath{<}0.001}  % single-test bootstrap p-value (two-sided), LightGBM
\providecommand{\rhmcsBarQAfPXgb}{}\renewcommand{\rhmcsBarQAfPXgb}{\ensuremath{<}0.001}  % StepM-adjusted p-value (one-sided), XGBoost
\providecommand{\rhmcsBarQAfRawXgb}{}\renewcommand{\rhmcsBarQAfRawXgb}{\ensuremath{<}0.001}  % single-test bootstrap p-value (one-sided), XGBoost
\providecommand{\rhmcsBarQAfPTwoXgb}{}\renewcommand{\rhmcsBarQAfPTwoXgb}{\ensuremath{<}0.001}  % StepM-adjusted p-value (two-sided), XGBoost
\providecommand{\rhmcsBarQAfRawTwoXgb}{}\renewcommand{\rhmcsBarQAfRawTwoXgb}{\ensuremath{<}0.001}  % single-test bootstrap p-value (two-sided), XGBoost
\providecommand{\rhmcsBarQAfPRf}{}\renewcommand{\rhmcsBarQAfPRf}{\ensuremath{<}0.001}  % StepM-adjusted p-value (one-sided), Random forest
\providecommand{\rhmcsBarQAfRawRf}{}\renewcommand{\rhmcsBarQAfRawRf}{\ensuremath{<}0.001}  % single-test bootstrap p-value (one-sided), Random forest
\providecommand{\rhmcsBarQAfPTwoRf}{}\renewcommand{\rhmcsBarQAfPTwoRf}{\ensuremath{<}0.001}  % StepM-adjusted p-value (two-sided), Random forest
\providecommand{\rhmcsBarQAfRawTwoRf}{}\renewcommand{\rhmcsBarQAfRawTwoRf}{\ensuremath{<}0.001}  % single-test bootstrap p-value (two-sided), Random forest
\providecommand{\rhmcsBarQAfNRejLo}{}\renewcommand{\rhmcsBarQAfNRejLo}{6}  % StepM rejections at 5 %, block length 5
\providecommand{\rhmcsBarQAfSpaCLo}{}\renewcommand{\rhmcsBarQAfSpaCLo}{\ensuremath{<}0.001}  % SPA p-value (consistent), block length 5
\providecommand{\rhmcsBarQAfNRejTwoLo}{}\renewcommand{\rhmcsBarQAfNRejTwoLo}{6}  % ... two-sided, block length 5
\providecommand{\rhmcsBarQAfPMinLo}{}\renewcommand{\rhmcsBarQAfPMinLo}{\ensuremath{<}0.001}  % smallest adjusted p, block length 5
\providecommand{\rhmcsBarQAfPTwoMinLo}{}\renewcommand{\rhmcsBarQAfPTwoMinLo}{\ensuremath{<}0.001}  % ... two-sided, block length 5
\providecommand{\rhmcsBarQAfNRejHi}{}\renewcommand{\rhmcsBarQAfNRejHi}{6}  % StepM rejections at 5 %, block length 20
\providecommand{\rhmcsBarQAfSpaCHi}{}\renewcommand{\rhmcsBarQAfSpaCHi}{\ensuremath{<}0.001}  % SPA p-value (consistent), block length 20
\providecommand{\rhmcsBarQAfNRejTwoHi}{}\renewcommand{\rhmcsBarQAfNRejTwoHi}{6}  % ... two-sided, block length 20
\providecommand{\rhmcsBarQAfPMinHi}{}\renewcommand{\rhmcsBarQAfPMinHi}{\ensuremath{<}0.001}  % smallest adjusted p, block length 20
\providecommand{\rhmcsBarQAfPTwoMinHi}{}\renewcommand{\rhmcsBarQAfPTwoMinHi}{\ensuremath{<}0.001}  % ... two-sided, block length 20
\providecommand{\rhmcsBarMAfSpaC}{}\renewcommand{\rhmcsBarMAfSpaC}{0.64}  % SPA p-value (consistent), family BarMAf
\providecommand{\rhmcsBarMAfSpaL}{}\renewcommand{\rhmcsBarMAfSpaL}{0.22}  % ... lower
\providecommand{\rhmcsBarMAfSpaU}{}\renewcommand{\rhmcsBarMAfSpaU}{0.64}  % ... upper
\providecommand{\rhmcsBarMAfNRej}{}\renewcommand{\rhmcsBarMAfNRej}{0}  % StepM rejections at FWER 5 % (of 6)
\providecommand{\rhmcsBarMAfNRejWord}{}\renewcommand{\rhmcsBarMAfNRejWord}{none}  % ... as a word
\providecommand{\rhmcsBarMAfNRaw}{}\renewcommand{\rhmcsBarMAfNRaw}{0}  % single-test bootstrap p < 0.05 (no adjustment)
\providecommand{\rhmcsBarMAfNRawWord}{}\renewcommand{\rhmcsBarMAfNRawWord}{none}  % ... as a word
\providecommand{\rhmcsBarMAfNPos}{}\renewcommand{\rhmcsBarMAfNPos}{0}  % alternatives with a positive estimate (better than the benchmark)
\providecommand{\rhmcsBarMAfPMin}{}\renewcommand{\rhmcsBarMAfPMin}{0.88}  % smallest StepM-adjusted p-value
\providecommand{\rhmcsBarMAfPMax}{}\renewcommand{\rhmcsBarMAfPMax}{0.89}  % largest StepM-adjusted p-value
\providecommand{\rhmcsBarMAfPTwoMin}{}\renewcommand{\rhmcsBarMAfPTwoMin}{0.28}  % smallest two-sided StepM-adjusted p-value
\providecommand{\rhmcsBarMAfPTwoMax}{}\renewcommand{\rhmcsBarMAfPTwoMax}{0.38}  % largest two-sided StepM-adjusted p-value
\providecommand{\rhmcsBarMAfRawMin}{}\renewcommand{\rhmcsBarMAfRawMin}{0.84}  % smallest single-test bootstrap p-value
\providecommand{\rhmcsBarMAfRej}{}\renewcommand{\rhmcsBarMAfRej}{none}  % the alternatives StepM rejects at 5 % (none: none)
\providecommand{\rhmcsBarMAfNRejTwo}{}\renewcommand{\rhmcsBarMAfNRejTwo}{0}  % StepM rejections at FWER 5 %, two-sided
\providecommand{\rhmcsBarMAfNRejTwoWord}{}\renewcommand{\rhmcsBarMAfNRejTwoWord}{none}  % ... as a word
\providecommand{\rhmcsBarMAfNRawTwo}{}\renewcommand{\rhmcsBarMAfNRawTwo}{0}  % single-test bootstrap p < 0.05, two-sided
\providecommand{\rhmcsBarMAfRejTwo}{}\renewcommand{\rhmcsBarMAfRejTwo}{none}  % the alternatives the two-sided StepM rejects at 5 % (none: none)
\providecommand{\rhmcsBarMAfPRidge}{}\renewcommand{\rhmcsBarMAfPRidge}{0.89}  % StepM-adjusted p-value (one-sided), Ridge
\providecommand{\rhmcsBarMAfRawRidge}{}\renewcommand{\rhmcsBarMAfRawRidge}{0.87}  % single-test bootstrap p-value (one-sided), Ridge
\providecommand{\rhmcsBarMAfPTwoRidge}{}\renewcommand{\rhmcsBarMAfPTwoRidge}{0.32}  % StepM-adjusted p-value (two-sided), Ridge
\providecommand{\rhmcsBarMAfRawTwoRidge}{}\renewcommand{\rhmcsBarMAfRawTwoRidge}{0.23}  % single-test bootstrap p-value (two-sided), Ridge
\providecommand{\rhmcsBarMAfPLasso}{}\renewcommand{\rhmcsBarMAfPLasso}{0.89}  % StepM-adjusted p-value (one-sided), Lasso
\providecommand{\rhmcsBarMAfRawLasso}{}\renewcommand{\rhmcsBarMAfRawLasso}{0.87}  % single-test bootstrap p-value (one-sided), Lasso
\providecommand{\rhmcsBarMAfPTwoLasso}{}\renewcommand{\rhmcsBarMAfPTwoLasso}{0.28}  % StepM-adjusted p-value (two-sided), Lasso
\providecommand{\rhmcsBarMAfRawTwoLasso}{}\renewcommand{\rhmcsBarMAfRawTwoLasso}{0.20}  % single-test bootstrap p-value (two-sided), Lasso
\providecommand{\rhmcsBarMAfPEnet}{}\renewcommand{\rhmcsBarMAfPEnet}{0.89}  % StepM-adjusted p-value (one-sided), Elastic net
\providecommand{\rhmcsBarMAfRawEnet}{}\renewcommand{\rhmcsBarMAfRawEnet}{0.87}  % single-test bootstrap p-value (one-sided), Elastic net
\providecommand{\rhmcsBarMAfPTwoEnet}{}\renewcommand{\rhmcsBarMAfPTwoEnet}{0.32}  % StepM-adjusted p-value (two-sided), Elastic net
\providecommand{\rhmcsBarMAfRawTwoEnet}{}\renewcommand{\rhmcsBarMAfRawTwoEnet}{0.23}  % single-test bootstrap p-value (two-sided), Elastic net
\providecommand{\rhmcsBarMAfPLgbm}{}\renewcommand{\rhmcsBarMAfPLgbm}{0.88}  % StepM-adjusted p-value (one-sided), LightGBM
\providecommand{\rhmcsBarMAfRawLgbm}{}\renewcommand{\rhmcsBarMAfRawLgbm}{0.84}  % single-test bootstrap p-value (one-sided), LightGBM
\providecommand{\rhmcsBarMAfPTwoLgbm}{}\renewcommand{\rhmcsBarMAfPTwoLgbm}{0.38}  % StepM-adjusted p-value (two-sided), LightGBM
\providecommand{\rhmcsBarMAfRawTwoLgbm}{}\renewcommand{\rhmcsBarMAfRawTwoLgbm}{0.36}  % single-test bootstrap p-value (two-sided), LightGBM
\providecommand{\rhmcsBarMAfPXgb}{}\renewcommand{\rhmcsBarMAfPXgb}{0.88}  % StepM-adjusted p-value (one-sided), XGBoost
\providecommand{\rhmcsBarMAfRawXgb}{}\renewcommand{\rhmcsBarMAfRawXgb}{0.84}  % single-test bootstrap p-value (one-sided), XGBoost
\providecommand{\rhmcsBarMAfPTwoXgb}{}\renewcommand{\rhmcsBarMAfPTwoXgb}{0.38}  % StepM-adjusted p-value (two-sided), XGBoost
\providecommand{\rhmcsBarMAfRawTwoXgb}{}\renewcommand{\rhmcsBarMAfRawTwoXgb}{0.35}  % single-test bootstrap p-value (two-sided), XGBoost
\providecommand{\rhmcsBarMAfPRf}{}\renewcommand{\rhmcsBarMAfPRf}{0.89}  % StepM-adjusted p-value (one-sided), Random forest
\providecommand{\rhmcsBarMAfRawRf}{}\renewcommand{\rhmcsBarMAfRawRf}{0.86}  % single-test bootstrap p-value (one-sided), Random forest
\providecommand{\rhmcsBarMAfPTwoRf}{}\renewcommand{\rhmcsBarMAfPTwoRf}{0.35}  % StepM-adjusted p-value (two-sided), Random forest
\providecommand{\rhmcsBarMAfRawTwoRf}{}\renewcommand{\rhmcsBarMAfRawTwoRf}{0.29}  % single-test bootstrap p-value (two-sided), Random forest
\providecommand{\rhmcsBarMAfNRejLo}{}\renewcommand{\rhmcsBarMAfNRejLo}{0}  % StepM rejections at 5 %, block length 5
\providecommand{\rhmcsBarMAfSpaCLo}{}\renewcommand{\rhmcsBarMAfSpaCLo}{0.64}  % SPA p-value (consistent), block length 5
\providecommand{\rhmcsBarMAfNRejTwoLo}{}\renewcommand{\rhmcsBarMAfNRejTwoLo}{0}  % ... two-sided, block length 5
\providecommand{\rhmcsBarMAfPMinLo}{}\renewcommand{\rhmcsBarMAfPMinLo}{0.91}  % smallest adjusted p, block length 5
\providecommand{\rhmcsBarMAfPTwoMinLo}{}\renewcommand{\rhmcsBarMAfPTwoMinLo}{0.21}  % ... two-sided, block length 5
\providecommand{\rhmcsBarMAfNRejHi}{}\renewcommand{\rhmcsBarMAfNRejHi}{0}  % StepM rejections at 5 %, block length 20
\providecommand{\rhmcsBarMAfSpaCHi}{}\renewcommand{\rhmcsBarMAfSpaCHi}{0.63}  % SPA p-value (consistent), block length 20
\providecommand{\rhmcsBarMAfNRejTwoHi}{}\renewcommand{\rhmcsBarMAfNRejTwoHi}{0}  % ... two-sided, block length 20
\providecommand{\rhmcsBarMAfPMinHi}{}\renewcommand{\rhmcsBarMAfPMinHi}{0.84}  % smallest adjusted p, block length 20
\providecommand{\rhmcsBarMAfPTwoMinHi}{}\renewcommand{\rhmcsBarMAfPTwoMinHi}{0.39}  % ... two-sided, block length 20
\providecommand{\rhmcsAllQAfSpaC}{}\renewcommand{\rhmcsAllQAfSpaC}{0.045}  % SPA p-value (consistent), family AllQAf
\providecommand{\rhmcsAllQAfSpaL}{}\renewcommand{\rhmcsAllQAfSpaL}{0.037}  % ... lower
\providecommand{\rhmcsAllQAfSpaU}{}\renewcommand{\rhmcsAllQAfSpaU}{0.045}  % ... upper
\providecommand{\rhmcsAllQAfNRej}{}\renewcommand{\rhmcsAllQAfNRej}{1}  % StepM rejections at FWER 5 % (of 6)
\providecommand{\rhmcsAllQAfNRejWord}{}\renewcommand{\rhmcsAllQAfNRejWord}{one}  % ... as a word
\providecommand{\rhmcsAllQAfNRaw}{}\renewcommand{\rhmcsAllQAfNRaw}{2}  % single-test bootstrap p < 0.05 (no adjustment)
\providecommand{\rhmcsAllQAfNRawWord}{}\renewcommand{\rhmcsAllQAfNRawWord}{two}  % ... as a word
\providecommand{\rhmcsAllQAfNPos}{}\renewcommand{\rhmcsAllQAfNPos}{3}  % alternatives with a positive estimate (better than the benchmark)
\providecommand{\rhmcsAllQAfPMin}{}\renewcommand{\rhmcsAllQAfPMin}{0.045}  % smallest StepM-adjusted p-value
\providecommand{\rhmcsAllQAfPMax}{}\renewcommand{\rhmcsAllQAfPMax}{0.90}  % largest StepM-adjusted p-value
\providecommand{\rhmcsAllQAfPTwoMin}{}\renewcommand{\rhmcsAllQAfPTwoMin}{0.11}  % smallest two-sided StepM-adjusted p-value
\providecommand{\rhmcsAllQAfPTwoMax}{}\renewcommand{\rhmcsAllQAfPTwoMax}{0.81}  % largest two-sided StepM-adjusted p-value
\providecommand{\rhmcsAllQAfRawMin}{}\renewcommand{\rhmcsAllQAfRawMin}{0.020}  % smallest single-test bootstrap p-value
\providecommand{\rhmcsAllQAfRej}{}\renewcommand{\rhmcsAllQAfRej}{XGBoost on the full predictor set}  % the alternatives StepM rejects at 5 % (none: none)
\providecommand{\rhmcsAllQAfNRejTwo}{}\renewcommand{\rhmcsAllQAfNRejTwo}{0}  % StepM rejections at FWER 5 %, two-sided
\providecommand{\rhmcsAllQAfNRejTwoWord}{}\renewcommand{\rhmcsAllQAfNRejTwoWord}{none}  % ... as a word
\providecommand{\rhmcsAllQAfNRawTwo}{}\renewcommand{\rhmcsAllQAfNRawTwo}{0}  % single-test bootstrap p < 0.05, two-sided
\providecommand{\rhmcsAllQAfRejTwo}{}\renewcommand{\rhmcsAllQAfRejTwo}{none}  % the alternatives the two-sided StepM rejects at 5 % (none: none)
\providecommand{\rhmcsAllQAfPRidge}{}\renewcommand{\rhmcsAllQAfPRidge}{0.78}  % StepM-adjusted p-value (one-sided), Ridge
\providecommand{\rhmcsAllQAfRawRidge}{}\renewcommand{\rhmcsAllQAfRawRidge}{0.76}  % single-test bootstrap p-value (one-sided), Ridge
\providecommand{\rhmcsAllQAfPTwoRidge}{}\renewcommand{\rhmcsAllQAfPTwoRidge}{0.55}  % StepM-adjusted p-value (two-sided), Ridge
\providecommand{\rhmcsAllQAfRawTwoRidge}{}\renewcommand{\rhmcsAllQAfRawTwoRidge}{0.50}  % single-test bootstrap p-value (two-sided), Ridge
\providecommand{\rhmcsAllQAfPLasso}{}\renewcommand{\rhmcsAllQAfPLasso}{0.90}  % StepM-adjusted p-value (one-sided), Lasso
\providecommand{\rhmcsAllQAfRawLasso}{}\renewcommand{\rhmcsAllQAfRawLasso}{0.90}  % single-test bootstrap p-value (one-sided), Lasso
\providecommand{\rhmcsAllQAfPTwoLasso}{}\renewcommand{\rhmcsAllQAfPTwoLasso}{0.32}  % StepM-adjusted p-value (two-sided), Lasso
\providecommand{\rhmcsAllQAfRawTwoLasso}{}\renewcommand{\rhmcsAllQAfRawTwoLasso}{0.18}  % single-test bootstrap p-value (two-sided), Lasso
\providecommand{\rhmcsAllQAfPEnet}{}\renewcommand{\rhmcsAllQAfPEnet}{0.67}  % StepM-adjusted p-value (one-sided), Elastic net
\providecommand{\rhmcsAllQAfRawEnet}{}\renewcommand{\rhmcsAllQAfRawEnet}{0.62}  % single-test bootstrap p-value (one-sided), Elastic net
\providecommand{\rhmcsAllQAfPTwoEnet}{}\renewcommand{\rhmcsAllQAfPTwoEnet}{0.81}  % StepM-adjusted p-value (two-sided), Elastic net
\providecommand{\rhmcsAllQAfRawTwoEnet}{}\renewcommand{\rhmcsAllQAfRawTwoEnet}{0.81}  % single-test bootstrap p-value (two-sided), Elastic net
\providecommand{\rhmcsAllQAfPLgbm}{}\renewcommand{\rhmcsAllQAfPLgbm}{0.081}  % StepM-adjusted p-value (one-sided), LightGBM
\providecommand{\rhmcsAllQAfRawLgbm}{}\renewcommand{\rhmcsAllQAfRawLgbm}{0.037}  % single-test bootstrap p-value (one-sided), LightGBM
\providecommand{\rhmcsAllQAfPTwoLgbm}{}\renewcommand{\rhmcsAllQAfPTwoLgbm}{0.18}  % StepM-adjusted p-value (two-sided), LightGBM
\providecommand{\rhmcsAllQAfRawTwoLgbm}{}\renewcommand{\rhmcsAllQAfRawTwoLgbm}{0.081}  % single-test bootstrap p-value (two-sided), LightGBM
\providecommand{\rhmcsAllQAfPXgb}{}\renewcommand{\rhmcsAllQAfPXgb}{0.045}  % StepM-adjusted p-value (one-sided), XGBoost
\providecommand{\rhmcsAllQAfRawXgb}{}\renewcommand{\rhmcsAllQAfRawXgb}{0.020}  % single-test bootstrap p-value (one-sided), XGBoost
\providecommand{\rhmcsAllQAfPTwoXgb}{}\renewcommand{\rhmcsAllQAfPTwoXgb}{0.11}  % StepM-adjusted p-value (two-sided), XGBoost
\providecommand{\rhmcsAllQAfRawTwoXgb}{}\renewcommand{\rhmcsAllQAfRawTwoXgb}{0.050}  % single-test bootstrap p-value (two-sided), XGBoost
\providecommand{\rhmcsAllQAfPRf}{}\renewcommand{\rhmcsAllQAfPRf}{0.26}  % StepM-adjusted p-value (one-sided), Random forest
\providecommand{\rhmcsAllQAfRawRf}{}\renewcommand{\rhmcsAllQAfRawRf}{0.15}  % single-test bootstrap p-value (one-sided), Random forest
\providecommand{\rhmcsAllQAfPTwoRf}{}\renewcommand{\rhmcsAllQAfPTwoRf}{0.48}  % StepM-adjusted p-value (two-sided), Random forest
\providecommand{\rhmcsAllQAfRawTwoRf}{}\renewcommand{\rhmcsAllQAfRawTwoRf}{0.29}  % single-test bootstrap p-value (two-sided), Random forest
\providecommand{\rhmcsAllQAfNRejLo}{}\renewcommand{\rhmcsAllQAfNRejLo}{1}  % StepM rejections at 5 %, block length 5
\providecommand{\rhmcsAllQAfSpaCLo}{}\renewcommand{\rhmcsAllQAfSpaCLo}{0.036}  % SPA p-value (consistent), block length 5
\providecommand{\rhmcsAllQAfNRejTwoLo}{}\renewcommand{\rhmcsAllQAfNRejTwoLo}{0}  % ... two-sided, block length 5
\providecommand{\rhmcsAllQAfPMinLo}{}\renewcommand{\rhmcsAllQAfPMinLo}{0.036}  % smallest adjusted p, block length 5
\providecommand{\rhmcsAllQAfPTwoMinLo}{}\renewcommand{\rhmcsAllQAfPTwoMinLo}{0.091}  % ... two-sided, block length 5
\providecommand{\rhmcsAllQAfNRejHi}{}\renewcommand{\rhmcsAllQAfNRejHi}{0}  % StepM rejections at 5 %, block length 20
\providecommand{\rhmcsAllQAfSpaCHi}{}\renewcommand{\rhmcsAllQAfSpaCHi}{0.051}  % SPA p-value (consistent), block length 20
\providecommand{\rhmcsAllQAfNRejTwoHi}{}\renewcommand{\rhmcsAllQAfNRejTwoHi}{0}  % ... two-sided, block length 20
\providecommand{\rhmcsAllQAfPMinHi}{}\renewcommand{\rhmcsAllQAfPMinHi}{0.051}  % smallest adjusted p, block length 20
\providecommand{\rhmcsAllQAfPTwoMinHi}{}\renewcommand{\rhmcsAllQAfPTwoMinHi}{0.12}  % ... two-sided, block length 20
\providecommand{\rhmcsAllMAfSpaC}{}\renewcommand{\rhmcsAllMAfSpaC}{0.66}  % SPA p-value (consistent), family AllMAf
\providecommand{\rhmcsAllMAfSpaL}{}\renewcommand{\rhmcsAllMAfSpaL}{0.20}  % ... lower
\providecommand{\rhmcsAllMAfSpaU}{}\renewcommand{\rhmcsAllMAfSpaU}{0.66}  % ... upper
\providecommand{\rhmcsAllMAfNRej}{}\renewcommand{\rhmcsAllMAfNRej}{0}  % StepM rejections at FWER 5 % (of 6)
\providecommand{\rhmcsAllMAfNRejWord}{}\renewcommand{\rhmcsAllMAfNRejWord}{none}  % ... as a word
\providecommand{\rhmcsAllMAfNRaw}{}\renewcommand{\rhmcsAllMAfNRaw}{0}  % single-test bootstrap p < 0.05 (no adjustment)
\providecommand{\rhmcsAllMAfNRawWord}{}\renewcommand{\rhmcsAllMAfNRawWord}{none}  % ... as a word
\providecommand{\rhmcsAllMAfNPos}{}\renewcommand{\rhmcsAllMAfNPos}{0}  % alternatives with a positive estimate (better than the benchmark)
\providecommand{\rhmcsAllMAfPMin}{}\renewcommand{\rhmcsAllMAfPMin}{0.91}  % smallest StepM-adjusted p-value
\providecommand{\rhmcsAllMAfPMax}{}\renewcommand{\rhmcsAllMAfPMax}{0.91}  % largest StepM-adjusted p-value
\providecommand{\rhmcsAllMAfPTwoMin}{}\renewcommand{\rhmcsAllMAfPTwoMin}{0.36}  % smallest two-sided StepM-adjusted p-value
\providecommand{\rhmcsAllMAfPTwoMax}{}\renewcommand{\rhmcsAllMAfPTwoMax}{0.37}  % largest two-sided StepM-adjusted p-value
\providecommand{\rhmcsAllMAfRawMin}{}\renewcommand{\rhmcsAllMAfRawMin}{0.86}  % smallest single-test bootstrap p-value
\providecommand{\rhmcsAllMAfRej}{}\renewcommand{\rhmcsAllMAfRej}{none}  % the alternatives StepM rejects at 5 % (none: none)
\providecommand{\rhmcsAllMAfNRejTwo}{}\renewcommand{\rhmcsAllMAfNRejTwo}{0}  % StepM rejections at FWER 5 %, two-sided
\providecommand{\rhmcsAllMAfNRejTwoWord}{}\renewcommand{\rhmcsAllMAfNRejTwoWord}{none}  % ... as a word
\providecommand{\rhmcsAllMAfNRawTwo}{}\renewcommand{\rhmcsAllMAfNRawTwo}{0}  % single-test bootstrap p < 0.05, two-sided
\providecommand{\rhmcsAllMAfRejTwo}{}\renewcommand{\rhmcsAllMAfRejTwo}{none}  % the alternatives the two-sided StepM rejects at 5 % (none: none)
\providecommand{\rhmcsAllMAfPRidge}{}\renewcommand{\rhmcsAllMAfPRidge}{0.91}  % StepM-adjusted p-value (one-sided), Ridge
\providecommand{\rhmcsAllMAfRawRidge}{}\renewcommand{\rhmcsAllMAfRawRidge}{0.87}  % single-test bootstrap p-value (one-sided), Ridge
\providecommand{\rhmcsAllMAfPTwoRidge}{}\renewcommand{\rhmcsAllMAfPTwoRidge}{0.36}  % StepM-adjusted p-value (two-sided), Ridge
\providecommand{\rhmcsAllMAfRawTwoRidge}{}\renewcommand{\rhmcsAllMAfRawTwoRidge}{0.26}  % single-test bootstrap p-value (two-sided), Ridge
\providecommand{\rhmcsAllMAfPLasso}{}\renewcommand{\rhmcsAllMAfPLasso}{0.91}  % StepM-adjusted p-value (one-sided), Lasso
\providecommand{\rhmcsAllMAfRawLasso}{}\renewcommand{\rhmcsAllMAfRawLasso}{0.87}  % single-test bootstrap p-value (one-sided), Lasso
\providecommand{\rhmcsAllMAfPTwoLasso}{}\renewcommand{\rhmcsAllMAfPTwoLasso}{0.36}  % StepM-adjusted p-value (two-sided), Lasso
\providecommand{\rhmcsAllMAfRawTwoLasso}{}\renewcommand{\rhmcsAllMAfRawTwoLasso}{0.26}  % single-test bootstrap p-value (two-sided), Lasso
\providecommand{\rhmcsAllMAfPEnet}{}\renewcommand{\rhmcsAllMAfPEnet}{0.91}  % StepM-adjusted p-value (one-sided), Elastic net
\providecommand{\rhmcsAllMAfRawEnet}{}\renewcommand{\rhmcsAllMAfRawEnet}{0.87}  % single-test bootstrap p-value (one-sided), Elastic net
\providecommand{\rhmcsAllMAfPTwoEnet}{}\renewcommand{\rhmcsAllMAfPTwoEnet}{0.37}  % StepM-adjusted p-value (two-sided), Elastic net
\providecommand{\rhmcsAllMAfRawTwoEnet}{}\renewcommand{\rhmcsAllMAfRawTwoEnet}{0.26}  % single-test bootstrap p-value (two-sided), Elastic net
\providecommand{\rhmcsAllMAfPLgbm}{}\renewcommand{\rhmcsAllMAfPLgbm}{0.91}  % StepM-adjusted p-value (one-sided), LightGBM
\providecommand{\rhmcsAllMAfRawLgbm}{}\renewcommand{\rhmcsAllMAfRawLgbm}{0.86}  % single-test bootstrap p-value (one-sided), LightGBM
\providecommand{\rhmcsAllMAfPTwoLgbm}{}\renewcommand{\rhmcsAllMAfPTwoLgbm}{0.37}  % StepM-adjusted p-value (two-sided), LightGBM
\providecommand{\rhmcsAllMAfRawTwoLgbm}{}\renewcommand{\rhmcsAllMAfRawTwoLgbm}{0.33}  % single-test bootstrap p-value (two-sided), LightGBM
\providecommand{\rhmcsAllMAfPXgb}{}\renewcommand{\rhmcsAllMAfPXgb}{0.91}  % StepM-adjusted p-value (one-sided), XGBoost
\providecommand{\rhmcsAllMAfRawXgb}{}\renewcommand{\rhmcsAllMAfRawXgb}{0.86}  % single-test bootstrap p-value (one-sided), XGBoost
\providecommand{\rhmcsAllMAfPTwoXgb}{}\renewcommand{\rhmcsAllMAfPTwoXgb}{0.37}  % StepM-adjusted p-value (two-sided), XGBoost
\providecommand{\rhmcsAllMAfRawTwoXgb}{}\renewcommand{\rhmcsAllMAfRawTwoXgb}{0.32}  % single-test bootstrap p-value (two-sided), XGBoost
\providecommand{\rhmcsAllMAfPRf}{}\renewcommand{\rhmcsAllMAfPRf}{0.91}  % StepM-adjusted p-value (one-sided), Random forest
\providecommand{\rhmcsAllMAfRawRf}{}\renewcommand{\rhmcsAllMAfRawRf}{0.86}  % single-test bootstrap p-value (one-sided), Random forest
\providecommand{\rhmcsAllMAfPTwoRf}{}\renewcommand{\rhmcsAllMAfPTwoRf}{0.37}  % StepM-adjusted p-value (two-sided), Random forest
\providecommand{\rhmcsAllMAfRawTwoRf}{}\renewcommand{\rhmcsAllMAfRawTwoRf}{0.30}  % single-test bootstrap p-value (two-sided), Random forest
\providecommand{\rhmcsAllMAfNRejLo}{}\renewcommand{\rhmcsAllMAfNRejLo}{0}  % StepM rejections at 5 %, block length 5
\providecommand{\rhmcsAllMAfSpaCLo}{}\renewcommand{\rhmcsAllMAfSpaCLo}{0.73}  % SPA p-value (consistent), block length 5
\providecommand{\rhmcsAllMAfNRejTwoLo}{}\renewcommand{\rhmcsAllMAfNRejTwoLo}{0}  % ... two-sided, block length 5
\providecommand{\rhmcsAllMAfPMinLo}{}\renewcommand{\rhmcsAllMAfPMinLo}{0.97}  % smallest adjusted p, block length 5
\providecommand{\rhmcsAllMAfPTwoMinLo}{}\renewcommand{\rhmcsAllMAfPTwoMinLo}{0.24}  % ... two-sided, block length 5
\providecommand{\rhmcsAllMAfNRejHi}{}\renewcommand{\rhmcsAllMAfNRejHi}{0}  % StepM rejections at 5 %, block length 20
\providecommand{\rhmcsAllMAfSpaCHi}{}\renewcommand{\rhmcsAllMAfSpaCHi}{0.62}  % SPA p-value (consistent), block length 20
\providecommand{\rhmcsAllMAfNRejTwoHi}{}\renewcommand{\rhmcsAllMAfNRejTwoHi}{0}  % ... two-sided, block length 20
\providecommand{\rhmcsAllMAfPMinHi}{}\renewcommand{\rhmcsAllMAfPMinHi}{0.87}  % smallest adjusted p, block length 20
\providecommand{\rhmcsAllMAfPTwoMinHi}{}\renewcommand{\rhmcsAllMAfPTwoMinHi}{0.45}  % ... two-sided, block length 20
\providecommand{\rhmcsOosQAfSpaC}{}\renewcommand{\rhmcsOosQAfSpaC}{0.34}  % SPA p-value (consistent), family OosQAf
\providecommand{\rhmcsOosQAfSpaL}{}\renewcommand{\rhmcsOosQAfSpaL}{0.34}  % ... lower
\providecommand{\rhmcsOosQAfSpaU}{}\renewcommand{\rhmcsOosQAfSpaU}{0.48}  % ... upper
\providecommand{\rhmcsOosQAfNRej}{}\renewcommand{\rhmcsOosQAfNRej}{0}  % StepM rejections at FWER 5 % (of 6)
\providecommand{\rhmcsOosQAfNRejWord}{}\renewcommand{\rhmcsOosQAfNRejWord}{none}  % ... as a word
\providecommand{\rhmcsOosQAfNRaw}{}\renewcommand{\rhmcsOosQAfNRaw}{0}  % single-test bootstrap p < 0.05 (no adjustment)
\providecommand{\rhmcsOosQAfNRawWord}{}\renewcommand{\rhmcsOosQAfNRawWord}{none}  % ... as a word
\providecommand{\rhmcsOosQAfNPos}{}\renewcommand{\rhmcsOosQAfNPos}{3}  % alternatives with a positive estimate (better than the benchmark)
\providecommand{\rhmcsOosQAfPMin}{}\renewcommand{\rhmcsOosQAfPMin}{0.48}  % smallest StepM-adjusted p-value
\providecommand{\rhmcsOosQAfPMax}{}\renewcommand{\rhmcsOosQAfPMax}{0.99}  % largest StepM-adjusted p-value
\providecommand{\rhmcsOosQAfPTwoMin}{}\renewcommand{\rhmcsOosQAfPTwoMin}{0.032}  % smallest two-sided StepM-adjusted p-value
\providecommand{\rhmcsOosQAfPTwoMax}{}\renewcommand{\rhmcsOosQAfPTwoMax}{0.78}  % largest two-sided StepM-adjusted p-value
\providecommand{\rhmcsOosQAfRawMin}{}\renewcommand{\rhmcsOosQAfRawMin}{0.26}  % smallest single-test bootstrap p-value
\providecommand{\rhmcsOosQAfRej}{}\renewcommand{\rhmcsOosQAfRej}{none}  % the alternatives StepM rejects at 5 % (none: none)
\providecommand{\rhmcsOosQAfNRejTwo}{}\renewcommand{\rhmcsOosQAfNRejTwo}{1}  % StepM rejections at FWER 5 %, two-sided
\providecommand{\rhmcsOosQAfNRejTwoWord}{}\renewcommand{\rhmcsOosQAfNRejTwoWord}{one}  % ... as a word
\providecommand{\rhmcsOosQAfNRawTwo}{}\renewcommand{\rhmcsOosQAfNRawTwo}{3}  % single-test bootstrap p < 0.05, two-sided
\providecommand{\rhmcsOosQAfRejTwo}{}\renewcommand{\rhmcsOosQAfRejTwo}{the lasso on the full predictor set}  % the alternatives the two-sided StepM rejects at 5 % (none: none)
\providecommand{\rhmcsOosQAfPRidge}{}\renewcommand{\rhmcsOosQAfPRidge}{0.99}  % StepM-adjusted p-value (one-sided), Ridge
\providecommand{\rhmcsOosQAfRawRidge}{}\renewcommand{\rhmcsOosQAfRawRidge}{0.99}  % single-test bootstrap p-value (one-sided), Ridge
\providecommand{\rhmcsOosQAfPTwoRidge}{}\renewcommand{\rhmcsOosQAfPTwoRidge}{0.052}  % StepM-adjusted p-value (two-sided), Ridge
\providecommand{\rhmcsOosQAfRawTwoRidge}{}\renewcommand{\rhmcsOosQAfRawTwoRidge}{0.019}  % single-test bootstrap p-value (two-sided), Ridge
\providecommand{\rhmcsOosQAfPLasso}{}\renewcommand{\rhmcsOosQAfPLasso}{0.99}  % StepM-adjusted p-value (one-sided), Lasso
\providecommand{\rhmcsOosQAfRawLasso}{}\renewcommand{\rhmcsOosQAfRawLasso}{0.99}  % single-test bootstrap p-value (one-sided), Lasso
\providecommand{\rhmcsOosQAfPTwoLasso}{}\renewcommand{\rhmcsOosQAfPTwoLasso}{0.032}  % StepM-adjusted p-value (two-sided), Lasso
\providecommand{\rhmcsOosQAfRawTwoLasso}{}\renewcommand{\rhmcsOosQAfRawTwoLasso}{0.012}  % single-test bootstrap p-value (two-sided), Lasso
\providecommand{\rhmcsOosQAfPEnet}{}\renewcommand{\rhmcsOosQAfPEnet}{0.98}  % StepM-adjusted p-value (one-sided), Elastic net
\providecommand{\rhmcsOosQAfRawEnet}{}\renewcommand{\rhmcsOosQAfRawEnet}{0.97}  % single-test bootstrap p-value (one-sided), Elastic net
\providecommand{\rhmcsOosQAfPTwoEnet}{}\renewcommand{\rhmcsOosQAfPTwoEnet}{0.11}  % StepM-adjusted p-value (two-sided), Elastic net
\providecommand{\rhmcsOosQAfRawTwoEnet}{}\renewcommand{\rhmcsOosQAfRawTwoEnet}{0.049}  % single-test bootstrap p-value (two-sided), Elastic net
\providecommand{\rhmcsOosQAfPLgbm}{}\renewcommand{\rhmcsOosQAfPLgbm}{0.53}  % StepM-adjusted p-value (one-sided), LightGBM
\providecommand{\rhmcsOosQAfRawLgbm}{}\renewcommand{\rhmcsOosQAfRawLgbm}{0.31}  % single-test bootstrap p-value (one-sided), LightGBM
\providecommand{\rhmcsOosQAfPTwoLgbm}{}\renewcommand{\rhmcsOosQAfPTwoLgbm}{0.76}  % StepM-adjusted p-value (two-sided), LightGBM
\providecommand{\rhmcsOosQAfRawTwoLgbm}{}\renewcommand{\rhmcsOosQAfRawTwoLgbm}{0.62}  % single-test bootstrap p-value (two-sided), LightGBM
\providecommand{\rhmcsOosQAfPXgb}{}\renewcommand{\rhmcsOosQAfPXgb}{0.48}  % StepM-adjusted p-value (one-sided), XGBoost
\providecommand{\rhmcsOosQAfRawXgb}{}\renewcommand{\rhmcsOosQAfRawXgb}{0.26}  % single-test bootstrap p-value (one-sided), XGBoost
\providecommand{\rhmcsOosQAfPTwoXgb}{}\renewcommand{\rhmcsOosQAfPTwoXgb}{0.68}  % StepM-adjusted p-value (two-sided), XGBoost
\providecommand{\rhmcsOosQAfRawTwoXgb}{}\renewcommand{\rhmcsOosQAfRawTwoXgb}{0.53}  % single-test bootstrap p-value (two-sided), XGBoost
\providecommand{\rhmcsOosQAfPRf}{}\renewcommand{\rhmcsOosQAfPRf}{0.59}  % StepM-adjusted p-value (one-sided), Random forest
\providecommand{\rhmcsOosQAfRawRf}{}\renewcommand{\rhmcsOosQAfRawRf}{0.40}  % single-test bootstrap p-value (one-sided), Random forest
\providecommand{\rhmcsOosQAfPTwoRf}{}\renewcommand{\rhmcsOosQAfPTwoRf}{0.78}  % StepM-adjusted p-value (two-sided), Random forest
\providecommand{\rhmcsOosQAfRawTwoRf}{}\renewcommand{\rhmcsOosQAfRawTwoRf}{0.78}  % single-test bootstrap p-value (two-sided), Random forest
\providecommand{\rhmcsOosQAfNRejLo}{}\renewcommand{\rhmcsOosQAfNRejLo}{0}  % StepM rejections at 5 %, block length 5
\providecommand{\rhmcsOosQAfSpaCLo}{}\renewcommand{\rhmcsOosQAfSpaCLo}{0.34}  % SPA p-value (consistent), block length 5
\providecommand{\rhmcsOosQAfNRejTwoLo}{}\renewcommand{\rhmcsOosQAfNRejTwoLo}{2}  % ... two-sided, block length 5
\providecommand{\rhmcsOosQAfPMinLo}{}\renewcommand{\rhmcsOosQAfPMinLo}{0.47}  % smallest adjusted p, block length 5
\providecommand{\rhmcsOosQAfPTwoMinLo}{}\renewcommand{\rhmcsOosQAfPTwoMinLo}{0.017}  % ... two-sided, block length 5
\providecommand{\rhmcsOosQAfNRejHi}{}\renewcommand{\rhmcsOosQAfNRejHi}{0}  % StepM rejections at 5 %, block length 20
\providecommand{\rhmcsOosQAfSpaCHi}{}\renewcommand{\rhmcsOosQAfSpaCHi}{0.43}  % SPA p-value (consistent), block length 20
\providecommand{\rhmcsOosQAfNRejTwoHi}{}\renewcommand{\rhmcsOosQAfNRejTwoHi}{0}  % ... two-sided, block length 20
\providecommand{\rhmcsOosQAfPMinHi}{}\renewcommand{\rhmcsOosQAfPMinHi}{0.46}  % smallest adjusted p, block length 20
\providecommand{\rhmcsOosQAfPTwoMinHi}{}\renewcommand{\rhmcsOosQAfPTwoMinHi}{0.051}  % ... two-sided, block length 20
\providecommand{\rhmcsOosQBmSpaC}{}\renewcommand{\rhmcsOosQBmSpaC}{0.53}  % SPA p-value (consistent), family OosQBm
\providecommand{\rhmcsOosQBmSpaL}{}\renewcommand{\rhmcsOosQBmSpaL}{0.083}  % ... lower
\providecommand{\rhmcsOosQBmSpaU}{}\renewcommand{\rhmcsOosQBmSpaU}{0.69}  % ... upper
\providecommand{\rhmcsOosQBmNRej}{}\renewcommand{\rhmcsOosQBmNRej}{0}  % StepM rejections at FWER 5 % (of 6)
\providecommand{\rhmcsOosQBmNRejWord}{}\renewcommand{\rhmcsOosQBmNRejWord}{none}  % ... as a word
\providecommand{\rhmcsOosQBmNRaw}{}\renewcommand{\rhmcsOosQBmNRaw}{0}  % single-test bootstrap p < 0.05 (no adjustment)
\providecommand{\rhmcsOosQBmNRawWord}{}\renewcommand{\rhmcsOosQBmNRawWord}{none}  % ... as a word
\providecommand{\rhmcsOosQBmNPos}{}\renewcommand{\rhmcsOosQBmNPos}{0}  % alternatives with a positive estimate (better than the benchmark)
\providecommand{\rhmcsOosQBmPMin}{}\renewcommand{\rhmcsOosQBmPMin}{0.98}  % smallest StepM-adjusted p-value
\providecommand{\rhmcsOosQBmPMax}{}\renewcommand{\rhmcsOosQBmPMax}{1.00}  % largest StepM-adjusted p-value
\providecommand{\rhmcsOosQBmPTwoMin}{}\renewcommand{\rhmcsOosQBmPTwoMin}{0.001}  % smallest two-sided StepM-adjusted p-value
\providecommand{\rhmcsOosQBmPTwoMax}{}\renewcommand{\rhmcsOosQBmPTwoMax}{0.17}  % largest two-sided StepM-adjusted p-value
\providecommand{\rhmcsOosQBmRawMin}{}\renewcommand{\rhmcsOosQBmRawMin}{0.92}  % smallest single-test bootstrap p-value
\providecommand{\rhmcsOosQBmRej}{}\renewcommand{\rhmcsOosQBmRej}{none}  % the alternatives StepM rejects at 5 % (none: none)
\providecommand{\rhmcsOosQBmNRejTwo}{}\renewcommand{\rhmcsOosQBmNRejTwo}{4}  % StepM rejections at FWER 5 %, two-sided
\providecommand{\rhmcsOosQBmNRejTwoWord}{}\renewcommand{\rhmcsOosQBmNRejTwoWord}{four}  % ... as a word
\providecommand{\rhmcsOosQBmNRawTwo}{}\renewcommand{\rhmcsOosQBmNRawTwo}{4}  % single-test bootstrap p < 0.05, two-sided
\providecommand{\rhmcsOosQBmRejTwo}{}\renewcommand{\rhmcsOosQBmRejTwo}{ridge regression on the full predictor set, the lasso on the full predictor set, the elastic net on the full predictor set, and the random forest on the full predictor set}  % the alternatives the two-sided StepM rejects at 5 % (none: none)
\providecommand{\rhmcsOosQBmPRidge}{}\renewcommand{\rhmcsOosQBmPRidge}{1.00}  % StepM-adjusted p-value (one-sided), Ridge
\providecommand{\rhmcsOosQBmRawRidge}{}\renewcommand{\rhmcsOosQBmRawRidge}{1.00}  % single-test bootstrap p-value (one-sided), Ridge
\providecommand{\rhmcsOosQBmPTwoRidge}{}\renewcommand{\rhmcsOosQBmPTwoRidge}{0.001}  % StepM-adjusted p-value (two-sided), Ridge
\providecommand{\rhmcsOosQBmRawTwoRidge}{}\renewcommand{\rhmcsOosQBmRawTwoRidge}{\ensuremath{<}0.001}  % single-test bootstrap p-value (two-sided), Ridge
\providecommand{\rhmcsOosQBmPLasso}{}\renewcommand{\rhmcsOosQBmPLasso}{1.00}  % StepM-adjusted p-value (one-sided), Lasso
\providecommand{\rhmcsOosQBmRawLasso}{}\renewcommand{\rhmcsOosQBmRawLasso}{1.00}  % single-test bootstrap p-value (one-sided), Lasso
\providecommand{\rhmcsOosQBmPTwoLasso}{}\renewcommand{\rhmcsOosQBmPTwoLasso}{0.001}  % StepM-adjusted p-value (two-sided), Lasso
\providecommand{\rhmcsOosQBmRawTwoLasso}{}\renewcommand{\rhmcsOosQBmRawTwoLasso}{\ensuremath{<}0.001}  % single-test bootstrap p-value (two-sided), Lasso
\providecommand{\rhmcsOosQBmPEnet}{}\renewcommand{\rhmcsOosQBmPEnet}{1.00}  % StepM-adjusted p-value (one-sided), Elastic net
\providecommand{\rhmcsOosQBmRawEnet}{}\renewcommand{\rhmcsOosQBmRawEnet}{1.00}  % single-test bootstrap p-value (one-sided), Elastic net
\providecommand{\rhmcsOosQBmPTwoEnet}{}\renewcommand{\rhmcsOosQBmPTwoEnet}{0.002}  % StepM-adjusted p-value (two-sided), Elastic net
\providecommand{\rhmcsOosQBmRawTwoEnet}{}\renewcommand{\rhmcsOosQBmRawTwoEnet}{\ensuremath{<}0.001}  % single-test bootstrap p-value (two-sided), Elastic net
\providecommand{\rhmcsOosQBmPLgbm}{}\renewcommand{\rhmcsOosQBmPLgbm}{0.98}  % StepM-adjusted p-value (one-sided), LightGBM
\providecommand{\rhmcsOosQBmRawLgbm}{}\renewcommand{\rhmcsOosQBmRawLgbm}{0.92}  % single-test bootstrap p-value (one-sided), LightGBM
\providecommand{\rhmcsOosQBmPTwoLgbm}{}\renewcommand{\rhmcsOosQBmPTwoLgbm}{0.17}  % StepM-adjusted p-value (two-sided), LightGBM
\providecommand{\rhmcsOosQBmRawTwoLgbm}{}\renewcommand{\rhmcsOosQBmRawTwoLgbm}{0.15}  % single-test bootstrap p-value (two-sided), LightGBM
\providecommand{\rhmcsOosQBmPXgb}{}\renewcommand{\rhmcsOosQBmPXgb}{0.98}  % StepM-adjusted p-value (one-sided), XGBoost
\providecommand{\rhmcsOosQBmRawXgb}{}\renewcommand{\rhmcsOosQBmRawXgb}{0.92}  % single-test bootstrap p-value (one-sided), XGBoost
\providecommand{\rhmcsOosQBmPTwoXgb}{}\renewcommand{\rhmcsOosQBmPTwoXgb}{0.17}  % StepM-adjusted p-value (two-sided), XGBoost
\providecommand{\rhmcsOosQBmRawTwoXgb}{}\renewcommand{\rhmcsOosQBmRawTwoXgb}{0.16}  % single-test bootstrap p-value (two-sided), XGBoost
\providecommand{\rhmcsOosQBmPRf}{}\renewcommand{\rhmcsOosQBmPRf}{1.00}  % StepM-adjusted p-value (one-sided), Random forest
\providecommand{\rhmcsOosQBmRawRf}{}\renewcommand{\rhmcsOosQBmRawRf}{0.98}  % single-test bootstrap p-value (one-sided), Random forest
\providecommand{\rhmcsOosQBmPTwoRf}{}\renewcommand{\rhmcsOosQBmPTwoRf}{0.039}  % StepM-adjusted p-value (two-sided), Random forest
\providecommand{\rhmcsOosQBmRawTwoRf}{}\renewcommand{\rhmcsOosQBmRawTwoRf}{0.030}  % single-test bootstrap p-value (two-sided), Random forest
\providecommand{\rhmcsOosQBmNRejLo}{}\renewcommand{\rhmcsOosQBmNRejLo}{0}  % StepM rejections at 5 %, block length 5
\providecommand{\rhmcsOosQBmSpaCLo}{}\renewcommand{\rhmcsOosQBmSpaCLo}{0.53}  % SPA p-value (consistent), block length 5
\providecommand{\rhmcsOosQBmNRejTwoLo}{}\renewcommand{\rhmcsOosQBmNRejTwoLo}{4}  % ... two-sided, block length 5
\providecommand{\rhmcsOosQBmPMinLo}{}\renewcommand{\rhmcsOosQBmPMinLo}{0.97}  % smallest adjusted p, block length 5
\providecommand{\rhmcsOosQBmPTwoMinLo}{}\renewcommand{\rhmcsOosQBmPTwoMinLo}{\ensuremath{<}0.001}  % ... two-sided, block length 5
\providecommand{\rhmcsOosQBmNRejHi}{}\renewcommand{\rhmcsOosQBmNRejHi}{0}  % StepM rejections at 5 %, block length 20
\providecommand{\rhmcsOosQBmSpaCHi}{}\renewcommand{\rhmcsOosQBmSpaCHi}{0.53}  % SPA p-value (consistent), block length 20
\providecommand{\rhmcsOosQBmNRejTwoHi}{}\renewcommand{\rhmcsOosQBmNRejTwoHi}{4}  % ... two-sided, block length 20
\providecommand{\rhmcsOosQBmPMinHi}{}\renewcommand{\rhmcsOosQBmPMinHi}{0.98}  % smallest adjusted p, block length 20
\providecommand{\rhmcsOosQBmPTwoMinHi}{}\renewcommand{\rhmcsOosQBmPTwoMinHi}{0.003}  % ... two-sided, block length 20
\providecommand{\rhmcsOosQBmAllSpaC}{}\renewcommand{\rhmcsOosQBmAllSpaC}{0.53}  % SPA p-value (consistent), family OosQBmAll
\providecommand{\rhmcsOosQBmAllSpaL}{}\renewcommand{\rhmcsOosQBmAllSpaL}{0.084}  % ... lower
\providecommand{\rhmcsOosQBmAllSpaU}{}\renewcommand{\rhmcsOosQBmAllSpaU}{0.79}  % ... upper
\providecommand{\rhmcsOosQBmAllNRej}{}\renewcommand{\rhmcsOosQBmAllNRej}{0}  % StepM rejections at FWER 5 % (of 10)
\providecommand{\rhmcsOosQBmAllNRejWord}{}\renewcommand{\rhmcsOosQBmAllNRejWord}{none}  % ... as a word
\providecommand{\rhmcsOosQBmAllNRaw}{}\renewcommand{\rhmcsOosQBmAllNRaw}{0}  % single-test bootstrap p < 0.05 (no adjustment)
\providecommand{\rhmcsOosQBmAllNRawWord}{}\renewcommand{\rhmcsOosQBmAllNRawWord}{none}  % ... as a word
\providecommand{\rhmcsOosQBmAllNPos}{}\renewcommand{\rhmcsOosQBmAllNPos}{0}  % alternatives with a positive estimate (better than the benchmark)
\providecommand{\rhmcsOosQBmAllPMin}{}\renewcommand{\rhmcsOosQBmAllPMin}{0.99}  % smallest StepM-adjusted p-value
\providecommand{\rhmcsOosQBmAllPMax}{}\renewcommand{\rhmcsOosQBmAllPMax}{1.00}  % largest StepM-adjusted p-value
\providecommand{\rhmcsOosQBmAllPTwoMin}{}\renewcommand{\rhmcsOosQBmAllPTwoMin}{0.003}  % smallest two-sided StepM-adjusted p-value
\providecommand{\rhmcsOosQBmAllPTwoMax}{}\renewcommand{\rhmcsOosQBmAllPTwoMax}{0.17}  % largest two-sided StepM-adjusted p-value
\providecommand{\rhmcsOosQBmAllRawMin}{}\renewcommand{\rhmcsOosQBmAllRawMin}{0.92}  % smallest single-test bootstrap p-value
\providecommand{\rhmcsOosQBmAllRej}{}\renewcommand{\rhmcsOosQBmAllRej}{none}  % the alternatives StepM rejects at 5 % (none: none)
\providecommand{\rhmcsOosQBmAllNRejTwo}{}\renewcommand{\rhmcsOosQBmAllNRejTwo}{8}  % StepM rejections at FWER 5 %, two-sided
\providecommand{\rhmcsOosQBmAllNRejTwoWord}{}\renewcommand{\rhmcsOosQBmAllNRejTwoWord}{eight}  % ... as a word
\providecommand{\rhmcsOosQBmAllNRawTwo}{}\renewcommand{\rhmcsOosQBmAllNRawTwo}{8}  % single-test bootstrap p < 0.05, two-sided
\providecommand{\rhmcsOosQBmAllRejTwo}{}\renewcommand{\rhmcsOosQBmAllRejTwo}{the linear models on the HAR benchmark, ridge regression on the full predictor set, the lasso on the full predictor set, the elastic net on the full predictor set, LightGBM on the HAR benchmark, XGBoost on the HAR benchmark, the random forest on the HAR benchmark, and the random forest on the full predictor set}  % the alternatives the two-sided StepM rejects at 5 % (none: none)
\providecommand{\rhmcsOosQBmAllPLinHar}{}\renewcommand{\rhmcsOosQBmAllPLinHar}{1.00}  % StepM-adjusted p-value (one-sided), Linear, HAR
\providecommand{\rhmcsOosQBmAllRawLinHar}{}\renewcommand{\rhmcsOosQBmAllRawLinHar}{0.99}  % single-test bootstrap p-value (one-sided), Linear, HAR
\providecommand{\rhmcsOosQBmAllPTwoLinHar}{}\renewcommand{\rhmcsOosQBmAllPTwoLinHar}{0.014}  % StepM-adjusted p-value (two-sided), Linear, HAR
\providecommand{\rhmcsOosQBmAllRawTwoLinHar}{}\renewcommand{\rhmcsOosQBmAllRawTwoLinHar}{0.007}  % single-test bootstrap p-value (two-sided), Linear, HAR
\providecommand{\rhmcsOosQBmAllPRidgeFull}{}\renewcommand{\rhmcsOosQBmAllPRidgeFull}{1.00}  % StepM-adjusted p-value (one-sided), Ridge, Full
\providecommand{\rhmcsOosQBmAllRawRidgeFull}{}\renewcommand{\rhmcsOosQBmAllRawRidgeFull}{1.00}  % single-test bootstrap p-value (one-sided), Ridge, Full
\providecommand{\rhmcsOosQBmAllPTwoRidgeFull}{}\renewcommand{\rhmcsOosQBmAllPTwoRidgeFull}{0.003}  % StepM-adjusted p-value (two-sided), Ridge, Full
\providecommand{\rhmcsOosQBmAllRawTwoRidgeFull}{}\renewcommand{\rhmcsOosQBmAllRawTwoRidgeFull}{\ensuremath{<}0.001}  % single-test bootstrap p-value (two-sided), Ridge, Full
\providecommand{\rhmcsOosQBmAllPLassoFull}{}\renewcommand{\rhmcsOosQBmAllPLassoFull}{1.00}  % StepM-adjusted p-value (one-sided), Lasso, Full
\providecommand{\rhmcsOosQBmAllRawLassoFull}{}\renewcommand{\rhmcsOosQBmAllRawLassoFull}{1.00}  % single-test bootstrap p-value (one-sided), Lasso, Full
\providecommand{\rhmcsOosQBmAllPTwoLassoFull}{}\renewcommand{\rhmcsOosQBmAllPTwoLassoFull}{0.003}  % StepM-adjusted p-value (two-sided), Lasso, Full
\providecommand{\rhmcsOosQBmAllRawTwoLassoFull}{}\renewcommand{\rhmcsOosQBmAllRawTwoLassoFull}{\ensuremath{<}0.001}  % single-test bootstrap p-value (two-sided), Lasso, Full
\providecommand{\rhmcsOosQBmAllPEnetFull}{}\renewcommand{\rhmcsOosQBmAllPEnetFull}{1.00}  % StepM-adjusted p-value (one-sided), Elastic net, Full
\providecommand{\rhmcsOosQBmAllRawEnetFull}{}\renewcommand{\rhmcsOosQBmAllRawEnetFull}{1.00}  % single-test bootstrap p-value (one-sided), Elastic net, Full
\providecommand{\rhmcsOosQBmAllPTwoEnetFull}{}\renewcommand{\rhmcsOosQBmAllPTwoEnetFull}{0.004}  % StepM-adjusted p-value (two-sided), Elastic net, Full
\providecommand{\rhmcsOosQBmAllRawTwoEnetFull}{}\renewcommand{\rhmcsOosQBmAllRawTwoEnetFull}{\ensuremath{<}0.001}  % single-test bootstrap p-value (two-sided), Elastic net, Full
\providecommand{\rhmcsOosQBmAllPLgbmHar}{}\renewcommand{\rhmcsOosQBmAllPLgbmHar}{1.00}  % StepM-adjusted p-value (one-sided), LightGBM, HAR
\providecommand{\rhmcsOosQBmAllRawLgbmHar}{}\renewcommand{\rhmcsOosQBmAllRawLgbmHar}{0.99}  % single-test bootstrap p-value (one-sided), LightGBM, HAR
\providecommand{\rhmcsOosQBmAllPTwoLgbmHar}{}\renewcommand{\rhmcsOosQBmAllPTwoLgbmHar}{0.021}  % StepM-adjusted p-value (two-sided), LightGBM, HAR
\providecommand{\rhmcsOosQBmAllRawTwoLgbmHar}{}\renewcommand{\rhmcsOosQBmAllRawTwoLgbmHar}{0.009}  % single-test bootstrap p-value (two-sided), LightGBM, HAR
\providecommand{\rhmcsOosQBmAllPLgbmFull}{}\renewcommand{\rhmcsOosQBmAllPLgbmFull}{0.99}  % StepM-adjusted p-value (one-sided), LightGBM, Full
\providecommand{\rhmcsOosQBmAllRawLgbmFull}{}\renewcommand{\rhmcsOosQBmAllRawLgbmFull}{0.92}  % single-test bootstrap p-value (one-sided), LightGBM, Full
\providecommand{\rhmcsOosQBmAllPTwoLgbmFull}{}\renewcommand{\rhmcsOosQBmAllPTwoLgbmFull}{0.17}  % StepM-adjusted p-value (two-sided), LightGBM, Full
\providecommand{\rhmcsOosQBmAllRawTwoLgbmFull}{}\renewcommand{\rhmcsOosQBmAllRawTwoLgbmFull}{0.15}  % single-test bootstrap p-value (two-sided), LightGBM, Full
\providecommand{\rhmcsOosQBmAllPXgbHar}{}\renewcommand{\rhmcsOosQBmAllPXgbHar}{1.00}  % StepM-adjusted p-value (one-sided), XGBoost, HAR
\providecommand{\rhmcsOosQBmAllRawXgbHar}{}\renewcommand{\rhmcsOosQBmAllRawXgbHar}{0.99}  % single-test bootstrap p-value (one-sided), XGBoost, HAR
\providecommand{\rhmcsOosQBmAllPTwoXgbHar}{}\renewcommand{\rhmcsOosQBmAllPTwoXgbHar}{0.016}  % StepM-adjusted p-value (two-sided), XGBoost, HAR
\providecommand{\rhmcsOosQBmAllRawTwoXgbHar}{}\renewcommand{\rhmcsOosQBmAllRawTwoXgbHar}{0.006}  % single-test bootstrap p-value (two-sided), XGBoost, HAR
\providecommand{\rhmcsOosQBmAllPXgbFull}{}\renewcommand{\rhmcsOosQBmAllPXgbFull}{0.99}  % StepM-adjusted p-value (one-sided), XGBoost, Full
\providecommand{\rhmcsOosQBmAllRawXgbFull}{}\renewcommand{\rhmcsOosQBmAllRawXgbFull}{0.92}  % single-test bootstrap p-value (one-sided), XGBoost, Full
\providecommand{\rhmcsOosQBmAllPTwoXgbFull}{}\renewcommand{\rhmcsOosQBmAllPTwoXgbFull}{0.17}  % StepM-adjusted p-value (two-sided), XGBoost, Full
\providecommand{\rhmcsOosQBmAllRawTwoXgbFull}{}\renewcommand{\rhmcsOosQBmAllRawTwoXgbFull}{0.16}  % single-test bootstrap p-value (two-sided), XGBoost, Full
\providecommand{\rhmcsOosQBmAllPRfHar}{}\renewcommand{\rhmcsOosQBmAllPRfHar}{1.00}  % StepM-adjusted p-value (one-sided), Random forest, HAR
\providecommand{\rhmcsOosQBmAllRawRfHar}{}\renewcommand{\rhmcsOosQBmAllRawRfHar}{1.00}  % single-test bootstrap p-value (one-sided), Random forest, HAR
\providecommand{\rhmcsOosQBmAllPTwoRfHar}{}\renewcommand{\rhmcsOosQBmAllPTwoRfHar}{0.006}  % StepM-adjusted p-value (two-sided), Random forest, HAR
\providecommand{\rhmcsOosQBmAllRawTwoRfHar}{}\renewcommand{\rhmcsOosQBmAllRawTwoRfHar}{0.002}  % single-test bootstrap p-value (two-sided), Random forest, HAR
\providecommand{\rhmcsOosQBmAllPRfFull}{}\renewcommand{\rhmcsOosQBmAllPRfFull}{1.00}  % StepM-adjusted p-value (one-sided), Random forest, Full
\providecommand{\rhmcsOosQBmAllRawRfFull}{}\renewcommand{\rhmcsOosQBmAllRawRfFull}{0.98}  % single-test bootstrap p-value (one-sided), Random forest, Full
\providecommand{\rhmcsOosQBmAllPTwoRfFull}{}\renewcommand{\rhmcsOosQBmAllPTwoRfFull}{0.039}  % StepM-adjusted p-value (two-sided), Random forest, Full
\providecommand{\rhmcsOosQBmAllRawTwoRfFull}{}\renewcommand{\rhmcsOosQBmAllRawTwoRfFull}{0.030}  % single-test bootstrap p-value (two-sided), Random forest, Full
\providecommand{\rhmcsOosQBmAllNRejLo}{}\renewcommand{\rhmcsOosQBmAllNRejLo}{0}  % StepM rejections at 5 %, block length 5
\providecommand{\rhmcsOosQBmAllSpaCLo}{}\renewcommand{\rhmcsOosQBmAllSpaCLo}{0.53}  % SPA p-value (consistent), block length 5
\providecommand{\rhmcsOosQBmAllNRejTwoLo}{}\renewcommand{\rhmcsOosQBmAllNRejTwoLo}{8}  % ... two-sided, block length 5
\providecommand{\rhmcsOosQBmAllPMinLo}{}\renewcommand{\rhmcsOosQBmAllPMinLo}{0.99}  % smallest adjusted p, block length 5
\providecommand{\rhmcsOosQBmAllPTwoMinLo}{}\renewcommand{\rhmcsOosQBmAllPTwoMinLo}{0.002}  % ... two-sided, block length 5
\providecommand{\rhmcsOosQBmAllNRejHi}{}\renewcommand{\rhmcsOosQBmAllNRejHi}{0}  % StepM rejections at 5 %, block length 20
\providecommand{\rhmcsOosQBmAllSpaCHi}{}\renewcommand{\rhmcsOosQBmAllSpaCHi}{0.53}  % SPA p-value (consistent), block length 20
\providecommand{\rhmcsOosQBmAllNRejTwoHi}{}\renewcommand{\rhmcsOosQBmAllNRejTwoHi}{8}  % ... two-sided, block length 20
\providecommand{\rhmcsOosQBmAllPMinHi}{}\renewcommand{\rhmcsOosQBmAllPMinHi}{0.99}  % smallest adjusted p, block length 20
\providecommand{\rhmcsOosQBmAllPTwoMinHi}{}\renewcommand{\rhmcsOosQBmAllPTwoMinHi}{0.005}  % ... two-sided, block length 20
\providecommand{\rhmcsOosMAfSpaC}{}\renewcommand{\rhmcsOosMAfSpaC}{0.67}  % SPA p-value (consistent), family OosMAf
\providecommand{\rhmcsOosMAfSpaL}{}\renewcommand{\rhmcsOosMAfSpaL}{0.19}  % ... lower
\providecommand{\rhmcsOosMAfSpaU}{}\renewcommand{\rhmcsOosMAfSpaU}{0.67}  % ... upper
\providecommand{\rhmcsOosMAfNRej}{}\renewcommand{\rhmcsOosMAfNRej}{0}  % StepM rejections at FWER 5 % (of 6)
\providecommand{\rhmcsOosMAfNRejWord}{}\renewcommand{\rhmcsOosMAfNRejWord}{none}  % ... as a word
\providecommand{\rhmcsOosMAfNRaw}{}\renewcommand{\rhmcsOosMAfNRaw}{0}  % single-test bootstrap p < 0.05 (no adjustment)
\providecommand{\rhmcsOosMAfNRawWord}{}\renewcommand{\rhmcsOosMAfNRawWord}{none}  % ... as a word
\providecommand{\rhmcsOosMAfNPos}{}\renewcommand{\rhmcsOosMAfNPos}{0}  % alternatives with a positive estimate (better than the benchmark)
\providecommand{\rhmcsOosMAfPMin}{}\renewcommand{\rhmcsOosMAfPMin}{0.92}  % smallest StepM-adjusted p-value
\providecommand{\rhmcsOosMAfPMax}{}\renewcommand{\rhmcsOosMAfPMax}{0.92}  % largest StepM-adjusted p-value
\providecommand{\rhmcsOosMAfPTwoMin}{}\renewcommand{\rhmcsOosMAfPTwoMin}{0.32}  % smallest two-sided StepM-adjusted p-value
\providecommand{\rhmcsOosMAfPTwoMax}{}\renewcommand{\rhmcsOosMAfPTwoMax}{0.34}  % largest two-sided StepM-adjusted p-value
\providecommand{\rhmcsOosMAfRawMin}{}\renewcommand{\rhmcsOosMAfRawMin}{0.86}  % smallest single-test bootstrap p-value
\providecommand{\rhmcsOosMAfRej}{}\renewcommand{\rhmcsOosMAfRej}{none}  % the alternatives StepM rejects at 5 % (none: none)
\providecommand{\rhmcsOosMAfNRejTwo}{}\renewcommand{\rhmcsOosMAfNRejTwo}{0}  % StepM rejections at FWER 5 %, two-sided
\providecommand{\rhmcsOosMAfNRejTwoWord}{}\renewcommand{\rhmcsOosMAfNRejTwoWord}{none}  % ... as a word
\providecommand{\rhmcsOosMAfNRawTwo}{}\renewcommand{\rhmcsOosMAfNRawTwo}{0}  % single-test bootstrap p < 0.05, two-sided
\providecommand{\rhmcsOosMAfRejTwo}{}\renewcommand{\rhmcsOosMAfRejTwo}{none}  % the alternatives the two-sided StepM rejects at 5 % (none: none)
\providecommand{\rhmcsOosMAfPRidge}{}\renewcommand{\rhmcsOosMAfPRidge}{0.92}  % StepM-adjusted p-value (one-sided), Ridge
\providecommand{\rhmcsOosMAfRawRidge}{}\renewcommand{\rhmcsOosMAfRawRidge}{0.88}  % single-test bootstrap p-value (one-sided), Ridge
\providecommand{\rhmcsOosMAfPTwoRidge}{}\renewcommand{\rhmcsOosMAfPTwoRidge}{0.33}  % StepM-adjusted p-value (two-sided), Ridge
\providecommand{\rhmcsOosMAfRawTwoRidge}{}\renewcommand{\rhmcsOosMAfRawTwoRidge}{0.23}  % single-test bootstrap p-value (two-sided), Ridge
\providecommand{\rhmcsOosMAfPLasso}{}\renewcommand{\rhmcsOosMAfPLasso}{0.92}  % StepM-adjusted p-value (one-sided), Lasso
\providecommand{\rhmcsOosMAfRawLasso}{}\renewcommand{\rhmcsOosMAfRawLasso}{0.88}  % single-test bootstrap p-value (one-sided), Lasso
\providecommand{\rhmcsOosMAfPTwoLasso}{}\renewcommand{\rhmcsOosMAfPTwoLasso}{0.32}  % StepM-adjusted p-value (two-sided), Lasso
\providecommand{\rhmcsOosMAfRawTwoLasso}{}\renewcommand{\rhmcsOosMAfRawTwoLasso}{0.22}  % single-test bootstrap p-value (two-sided), Lasso
\providecommand{\rhmcsOosMAfPEnet}{}\renewcommand{\rhmcsOosMAfPEnet}{0.92}  % StepM-adjusted p-value (one-sided), Elastic net
\providecommand{\rhmcsOosMAfRawEnet}{}\renewcommand{\rhmcsOosMAfRawEnet}{0.88}  % single-test bootstrap p-value (one-sided), Elastic net
\providecommand{\rhmcsOosMAfPTwoEnet}{}\renewcommand{\rhmcsOosMAfPTwoEnet}{0.34}  % StepM-adjusted p-value (two-sided), Elastic net
\providecommand{\rhmcsOosMAfRawTwoEnet}{}\renewcommand{\rhmcsOosMAfRawTwoEnet}{0.23}  % single-test bootstrap p-value (two-sided), Elastic net
\providecommand{\rhmcsOosMAfPLgbm}{}\renewcommand{\rhmcsOosMAfPLgbm}{0.92}  % StepM-adjusted p-value (one-sided), LightGBM
\providecommand{\rhmcsOosMAfRawLgbm}{}\renewcommand{\rhmcsOosMAfRawLgbm}{0.86}  % single-test bootstrap p-value (one-sided), LightGBM
\providecommand{\rhmcsOosMAfPTwoLgbm}{}\renewcommand{\rhmcsOosMAfPTwoLgbm}{0.34}  % StepM-adjusted p-value (two-sided), LightGBM
\providecommand{\rhmcsOosMAfRawTwoLgbm}{}\renewcommand{\rhmcsOosMAfRawTwoLgbm}{0.31}  % single-test bootstrap p-value (two-sided), LightGBM
\providecommand{\rhmcsOosMAfPXgb}{}\renewcommand{\rhmcsOosMAfPXgb}{0.92}  % StepM-adjusted p-value (one-sided), XGBoost
\providecommand{\rhmcsOosMAfRawXgb}{}\renewcommand{\rhmcsOosMAfRawXgb}{0.86}  % single-test bootstrap p-value (one-sided), XGBoost
\providecommand{\rhmcsOosMAfPTwoXgb}{}\renewcommand{\rhmcsOosMAfPTwoXgb}{0.34}  % StepM-adjusted p-value (two-sided), XGBoost
\providecommand{\rhmcsOosMAfRawTwoXgb}{}\renewcommand{\rhmcsOosMAfRawTwoXgb}{0.29}  % single-test bootstrap p-value (two-sided), XGBoost
\providecommand{\rhmcsOosMAfPRf}{}\renewcommand{\rhmcsOosMAfPRf}{0.92}  % StepM-adjusted p-value (one-sided), Random forest
\providecommand{\rhmcsOosMAfRawRf}{}\renewcommand{\rhmcsOosMAfRawRf}{0.86}  % single-test bootstrap p-value (one-sided), Random forest
\providecommand{\rhmcsOosMAfPTwoRf}{}\renewcommand{\rhmcsOosMAfPTwoRf}{0.34}  % StepM-adjusted p-value (two-sided), Random forest
\providecommand{\rhmcsOosMAfRawTwoRf}{}\renewcommand{\rhmcsOosMAfRawTwoRf}{0.28}  % single-test bootstrap p-value (two-sided), Random forest
\providecommand{\rhmcsOosMAfNRejLo}{}\renewcommand{\rhmcsOosMAfNRejLo}{0}  % StepM rejections at 5 %, block length 5
\providecommand{\rhmcsOosMAfSpaCLo}{}\renewcommand{\rhmcsOosMAfSpaCLo}{0.73}  % SPA p-value (consistent), block length 5
\providecommand{\rhmcsOosMAfNRejTwoLo}{}\renewcommand{\rhmcsOosMAfNRejTwoLo}{0}  % ... two-sided, block length 5
\providecommand{\rhmcsOosMAfPMinLo}{}\renewcommand{\rhmcsOosMAfPMinLo}{0.97}  % smallest adjusted p, block length 5
\providecommand{\rhmcsOosMAfPTwoMinLo}{}\renewcommand{\rhmcsOosMAfPTwoMinLo}{0.21}  % ... two-sided, block length 5
\providecommand{\rhmcsOosMAfNRejHi}{}\renewcommand{\rhmcsOosMAfNRejHi}{0}  % StepM rejections at 5 %, block length 20
\providecommand{\rhmcsOosMAfSpaCHi}{}\renewcommand{\rhmcsOosMAfSpaCHi}{0.62}  % SPA p-value (consistent), block length 20
\providecommand{\rhmcsOosMAfNRejTwoHi}{}\renewcommand{\rhmcsOosMAfNRejTwoHi}{0}  % ... two-sided, block length 20
\providecommand{\rhmcsOosMAfPMinHi}{}\renewcommand{\rhmcsOosMAfPMinHi}{0.88}  % smallest adjusted p, block length 20
\providecommand{\rhmcsOosMAfPTwoMinHi}{}\renewcommand{\rhmcsOosMAfPTwoMinHi}{0.38}  % ... two-sided, block length 20
\providecommand{\rhmcsOosMBmSpaC}{}\renewcommand{\rhmcsOosMBmSpaC}{0.63}  % SPA p-value (consistent), family OosMBm
\providecommand{\rhmcsOosMBmSpaL}{}\renewcommand{\rhmcsOosMBmSpaL}{0.33}  % ... lower
\providecommand{\rhmcsOosMBmSpaU}{}\renewcommand{\rhmcsOosMBmSpaU}{0.63}  % ... upper
\providecommand{\rhmcsOosMBmNRej}{}\renewcommand{\rhmcsOosMBmNRej}{0}  % StepM rejections at FWER 5 % (of 6)
\providecommand{\rhmcsOosMBmNRejWord}{}\renewcommand{\rhmcsOosMBmNRejWord}{none}  % ... as a word
\providecommand{\rhmcsOosMBmNRaw}{}\renewcommand{\rhmcsOosMBmNRaw}{0}  % single-test bootstrap p < 0.05 (no adjustment)
\providecommand{\rhmcsOosMBmNRawWord}{}\renewcommand{\rhmcsOosMBmNRawWord}{none}  % ... as a word
\providecommand{\rhmcsOosMBmNPos}{}\renewcommand{\rhmcsOosMBmNPos}{1}  % alternatives with a positive estimate (better than the benchmark)
\providecommand{\rhmcsOosMBmPMin}{}\renewcommand{\rhmcsOosMBmPMin}{0.63}  % smallest StepM-adjusted p-value
\providecommand{\rhmcsOosMBmPMax}{}\renewcommand{\rhmcsOosMBmPMax}{0.88}  % largest StepM-adjusted p-value
\providecommand{\rhmcsOosMBmPTwoMin}{}\renewcommand{\rhmcsOosMBmPTwoMin}{0.47}  % smallest two-sided StepM-adjusted p-value
\providecommand{\rhmcsOosMBmPTwoMax}{}\renewcommand{\rhmcsOosMBmPTwoMax}{0.74}  % largest two-sided StepM-adjusted p-value
\providecommand{\rhmcsOosMBmRawMin}{}\renewcommand{\rhmcsOosMBmRawMin}{0.33}  % smallest single-test bootstrap p-value
\providecommand{\rhmcsOosMBmRej}{}\renewcommand{\rhmcsOosMBmRej}{none}  % the alternatives StepM rejects at 5 % (none: none)
\providecommand{\rhmcsOosMBmNRejTwo}{}\renewcommand{\rhmcsOosMBmNRejTwo}{0}  % StepM rejections at FWER 5 %, two-sided
\providecommand{\rhmcsOosMBmNRejTwoWord}{}\renewcommand{\rhmcsOosMBmNRejTwoWord}{none}  % ... as a word
\providecommand{\rhmcsOosMBmNRawTwo}{}\renewcommand{\rhmcsOosMBmNRawTwo}{0}  % single-test bootstrap p < 0.05, two-sided
\providecommand{\rhmcsOosMBmRejTwo}{}\renewcommand{\rhmcsOosMBmRejTwo}{none}  % the alternatives the two-sided StepM rejects at 5 % (none: none)
\providecommand{\rhmcsOosMBmPRidge}{}\renewcommand{\rhmcsOosMBmPRidge}{0.88}  % StepM-adjusted p-value (one-sided), Ridge
\providecommand{\rhmcsOosMBmRawRidge}{}\renewcommand{\rhmcsOosMBmRawRidge}{0.88}  % single-test bootstrap p-value (one-sided), Ridge
\providecommand{\rhmcsOosMBmPTwoRidge}{}\renewcommand{\rhmcsOosMBmPTwoRidge}{0.48}  % StepM-adjusted p-value (two-sided), Ridge
\providecommand{\rhmcsOosMBmRawTwoRidge}{}\renewcommand{\rhmcsOosMBmRawTwoRidge}{0.26}  % single-test bootstrap p-value (two-sided), Ridge
\providecommand{\rhmcsOosMBmPLasso}{}\renewcommand{\rhmcsOosMBmPLasso}{0.88}  % StepM-adjusted p-value (one-sided), Lasso
\providecommand{\rhmcsOosMBmRawLasso}{}\renewcommand{\rhmcsOosMBmRawLasso}{0.88}  % single-test bootstrap p-value (one-sided), Lasso
\providecommand{\rhmcsOosMBmPTwoLasso}{}\renewcommand{\rhmcsOosMBmPTwoLasso}{0.47}  % StepM-adjusted p-value (two-sided), Lasso
\providecommand{\rhmcsOosMBmRawTwoLasso}{}\renewcommand{\rhmcsOosMBmRawTwoLasso}{0.253}  % single-test bootstrap p-value (two-sided), Lasso
\providecommand{\rhmcsOosMBmPEnet}{}\renewcommand{\rhmcsOosMBmPEnet}{0.88}  % StepM-adjusted p-value (one-sided), Elastic net
\providecommand{\rhmcsOosMBmRawEnet}{}\renewcommand{\rhmcsOosMBmRawEnet}{0.88}  % single-test bootstrap p-value (one-sided), Elastic net
\providecommand{\rhmcsOosMBmPTwoEnet}{}\renewcommand{\rhmcsOosMBmPTwoEnet}{0.49}  % StepM-adjusted p-value (two-sided), Elastic net
\providecommand{\rhmcsOosMBmRawTwoEnet}{}\renewcommand{\rhmcsOosMBmRawTwoEnet}{0.26}  % single-test bootstrap p-value (two-sided), Elastic net
\providecommand{\rhmcsOosMBmPLgbm}{}\renewcommand{\rhmcsOosMBmPLgbm}{0.87}  % StepM-adjusted p-value (one-sided), LightGBM
\providecommand{\rhmcsOosMBmRawLgbm}{}\renewcommand{\rhmcsOosMBmRawLgbm}{0.81}  % single-test bootstrap p-value (one-sided), LightGBM
\providecommand{\rhmcsOosMBmPTwoLgbm}{}\renewcommand{\rhmcsOosMBmPTwoLgbm}{0.62}  % StepM-adjusted p-value (two-sided), LightGBM
\providecommand{\rhmcsOosMBmRawTwoLgbm}{}\renewcommand{\rhmcsOosMBmRawTwoLgbm}{0.40}  % single-test bootstrap p-value (two-sided), LightGBM
\providecommand{\rhmcsOosMBmPXgb}{}\renewcommand{\rhmcsOosMBmPXgb}{0.86}  % StepM-adjusted p-value (one-sided), XGBoost
\providecommand{\rhmcsOosMBmRawXgb}{}\renewcommand{\rhmcsOosMBmRawXgb}{0.76}  % single-test bootstrap p-value (one-sided), XGBoost
\providecommand{\rhmcsOosMBmPTwoXgb}{}\renewcommand{\rhmcsOosMBmPTwoXgb}{0.74}  % StepM-adjusted p-value (two-sided), XGBoost
\providecommand{\rhmcsOosMBmRawTwoXgb}{}\renewcommand{\rhmcsOosMBmRawTwoXgb}{0.55}  % single-test bootstrap p-value (two-sided), XGBoost
\providecommand{\rhmcsOosMBmPRf}{}\renewcommand{\rhmcsOosMBmPRf}{0.63}  % StepM-adjusted p-value (one-sided), Random forest
\providecommand{\rhmcsOosMBmRawRf}{}\renewcommand{\rhmcsOosMBmRawRf}{0.33}  % single-test bootstrap p-value (one-sided), Random forest
\providecommand{\rhmcsOosMBmPTwoRf}{}\renewcommand{\rhmcsOosMBmPTwoRf}{0.74}  % StepM-adjusted p-value (two-sided), Random forest
\providecommand{\rhmcsOosMBmRawTwoRf}{}\renewcommand{\rhmcsOosMBmRawTwoRf}{0.72}  % single-test bootstrap p-value (two-sided), Random forest
\providecommand{\rhmcsOosMBmNRejLo}{}\renewcommand{\rhmcsOosMBmNRejLo}{0}  % StepM rejections at 5 %, block length 5
\providecommand{\rhmcsOosMBmSpaCLo}{}\renewcommand{\rhmcsOosMBmSpaCLo}{0.72}  % SPA p-value (consistent), block length 5
\providecommand{\rhmcsOosMBmNRejTwoLo}{}\renewcommand{\rhmcsOosMBmNRejTwoLo}{0}  % ... two-sided, block length 5
\providecommand{\rhmcsOosMBmPMinLo}{}\renewcommand{\rhmcsOosMBmPMinLo}{0.72}  % smallest adjusted p, block length 5
\providecommand{\rhmcsOosMBmPTwoMinLo}{}\renewcommand{\rhmcsOosMBmPTwoMinLo}{0.32}  % ... two-sided, block length 5
\providecommand{\rhmcsOosMBmNRejHi}{}\renewcommand{\rhmcsOosMBmNRejHi}{0}  % StepM rejections at 5 %, block length 20
\providecommand{\rhmcsOosMBmSpaCHi}{}\renewcommand{\rhmcsOosMBmSpaCHi}{0.61}  % SPA p-value (consistent), block length 20
\providecommand{\rhmcsOosMBmNRejTwoHi}{}\renewcommand{\rhmcsOosMBmNRejTwoHi}{0}  % ... two-sided, block length 20
\providecommand{\rhmcsOosMBmPMinHi}{}\renewcommand{\rhmcsOosMBmPMinHi}{0.61}  % smallest adjusted p, block length 20
\providecommand{\rhmcsOosMBmPTwoMinHi}{}\renewcommand{\rhmcsOosMBmPTwoMinHi}{0.58}  % ... two-sided, block length 20
\providecommand{\rhmcsOosMBmAllSpaC}{}\renewcommand{\rhmcsOosMBmAllSpaC}{0.21}  % SPA p-value (consistent), family OosMBmAll
\providecommand{\rhmcsOosMBmAllSpaL}{}\renewcommand{\rhmcsOosMBmAllSpaL}{0.20}  % ... lower
\providecommand{\rhmcsOosMBmAllSpaU}{}\renewcommand{\rhmcsOosMBmAllSpaU}{0.21}  % ... upper
\providecommand{\rhmcsOosMBmAllNRej}{}\renewcommand{\rhmcsOosMBmAllNRej}{0}  % StepM rejections at FWER 5 % (of 10)
\providecommand{\rhmcsOosMBmAllNRejWord}{}\renewcommand{\rhmcsOosMBmAllNRejWord}{none}  % ... as a word
\providecommand{\rhmcsOosMBmAllNRaw}{}\renewcommand{\rhmcsOosMBmAllNRaw}{0}  % single-test bootstrap p < 0.05 (no adjustment)
\providecommand{\rhmcsOosMBmAllNRawWord}{}\renewcommand{\rhmcsOosMBmAllNRawWord}{none}  % ... as a word
\providecommand{\rhmcsOosMBmAllNPos}{}\renewcommand{\rhmcsOosMBmAllNPos}{5}  % alternatives with a positive estimate (better than the benchmark)
\providecommand{\rhmcsOosMBmAllPMin}{}\renewcommand{\rhmcsOosMBmAllPMin}{0.21}  % smallest StepM-adjusted p-value
\providecommand{\rhmcsOosMBmAllPMax}{}\renewcommand{\rhmcsOosMBmAllPMax}{0.88}  % largest StepM-adjusted p-value
\providecommand{\rhmcsOosMBmAllPTwoMin}{}\renewcommand{\rhmcsOosMBmAllPTwoMin}{0.40}  % smallest two-sided StepM-adjusted p-value
\providecommand{\rhmcsOosMBmAllPTwoMax}{}\renewcommand{\rhmcsOosMBmAllPTwoMax}{0.88}  % largest two-sided StepM-adjusted p-value
\providecommand{\rhmcsOosMBmAllRawMin}{}\renewcommand{\rhmcsOosMBmAllRawMin}{0.094}  % smallest single-test bootstrap p-value
\providecommand{\rhmcsOosMBmAllRej}{}\renewcommand{\rhmcsOosMBmAllRej}{none}  % the alternatives StepM rejects at 5 % (none: none)
\providecommand{\rhmcsOosMBmAllNRejTwo}{}\renewcommand{\rhmcsOosMBmAllNRejTwo}{0}  % StepM rejections at FWER 5 %, two-sided
\providecommand{\rhmcsOosMBmAllNRejTwoWord}{}\renewcommand{\rhmcsOosMBmAllNRejTwoWord}{none}  % ... as a word
\providecommand{\rhmcsOosMBmAllNRawTwo}{}\renewcommand{\rhmcsOosMBmAllNRawTwo}{0}  % single-test bootstrap p < 0.05, two-sided
\providecommand{\rhmcsOosMBmAllRejTwo}{}\renewcommand{\rhmcsOosMBmAllRejTwo}{none}  % the alternatives the two-sided StepM rejects at 5 % (none: none)
\providecommand{\rhmcsOosMBmAllPLinHar}{}\renewcommand{\rhmcsOosMBmAllPLinHar}{0.73}  % StepM-adjusted p-value (one-sided), Linear, HAR
\providecommand{\rhmcsOosMBmAllRawLinHar}{}\renewcommand{\rhmcsOosMBmAllRawLinHar}{0.41}  % single-test bootstrap p-value (one-sided), Linear, HAR
\providecommand{\rhmcsOosMBmAllPTwoLinHar}{}\renewcommand{\rhmcsOosMBmAllPTwoLinHar}{0.88}  % StepM-adjusted p-value (two-sided), Linear, HAR
\providecommand{\rhmcsOosMBmAllRawTwoLinHar}{}\renewcommand{\rhmcsOosMBmAllRawTwoLinHar}{0.88}  % single-test bootstrap p-value (two-sided), Linear, HAR
\providecommand{\rhmcsOosMBmAllPRidgeFull}{}\renewcommand{\rhmcsOosMBmAllPRidgeFull}{0.88}  % StepM-adjusted p-value (one-sided), Ridge, Full
\providecommand{\rhmcsOosMBmAllRawRidgeFull}{}\renewcommand{\rhmcsOosMBmAllRawRidgeFull}{0.88}  % single-test bootstrap p-value (one-sided), Ridge, Full
\providecommand{\rhmcsOosMBmAllPTwoRidgeFull}{}\renewcommand{\rhmcsOosMBmAllPTwoRidgeFull}{0.56}  % StepM-adjusted p-value (two-sided), Ridge, Full
\providecommand{\rhmcsOosMBmAllRawTwoRidgeFull}{}\renewcommand{\rhmcsOosMBmAllRawTwoRidgeFull}{0.26}  % single-test bootstrap p-value (two-sided), Ridge, Full
\providecommand{\rhmcsOosMBmAllPLassoFull}{}\renewcommand{\rhmcsOosMBmAllPLassoFull}{0.88}  % StepM-adjusted p-value (one-sided), Lasso, Full
\providecommand{\rhmcsOosMBmAllRawLassoFull}{}\renewcommand{\rhmcsOosMBmAllRawLassoFull}{0.88}  % single-test bootstrap p-value (one-sided), Lasso, Full
\providecommand{\rhmcsOosMBmAllPTwoLassoFull}{}\renewcommand{\rhmcsOosMBmAllPTwoLassoFull}{0.54}  % StepM-adjusted p-value (two-sided), Lasso, Full
\providecommand{\rhmcsOosMBmAllRawTwoLassoFull}{}\renewcommand{\rhmcsOosMBmAllRawTwoLassoFull}{0.253}  % single-test bootstrap p-value (two-sided), Lasso, Full
\providecommand{\rhmcsOosMBmAllPEnetFull}{}\renewcommand{\rhmcsOosMBmAllPEnetFull}{0.88}  % StepM-adjusted p-value (one-sided), Elastic net, Full
\providecommand{\rhmcsOosMBmAllRawEnetFull}{}\renewcommand{\rhmcsOosMBmAllRawEnetFull}{0.88}  % single-test bootstrap p-value (one-sided), Elastic net, Full
\providecommand{\rhmcsOosMBmAllPTwoEnetFull}{}\renewcommand{\rhmcsOosMBmAllPTwoEnetFull}{0.57}  % StepM-adjusted p-value (two-sided), Elastic net, Full
\providecommand{\rhmcsOosMBmAllRawTwoEnetFull}{}\renewcommand{\rhmcsOosMBmAllRawTwoEnetFull}{0.26}  % single-test bootstrap p-value (two-sided), Elastic net, Full
\providecommand{\rhmcsOosMBmAllPLgbmHar}{}\renewcommand{\rhmcsOosMBmAllPLgbmHar}{0.30}  % StepM-adjusted p-value (one-sided), LightGBM, HAR
\providecommand{\rhmcsOosMBmAllRawLgbmHar}{}\renewcommand{\rhmcsOosMBmAllRawLgbmHar}{0.12}  % single-test bootstrap p-value (one-sided), LightGBM, HAR
\providecommand{\rhmcsOosMBmAllPTwoLgbmHar}{}\renewcommand{\rhmcsOosMBmAllPTwoLgbmHar}{0.54}  % StepM-adjusted p-value (two-sided), LightGBM, HAR
\providecommand{\rhmcsOosMBmAllRawTwoLgbmHar}{}\renewcommand{\rhmcsOosMBmAllRawTwoLgbmHar}{0.24}  % single-test bootstrap p-value (two-sided), LightGBM, HAR
\providecommand{\rhmcsOosMBmAllPLgbmFull}{}\renewcommand{\rhmcsOosMBmAllPLgbmFull}{0.87}  % StepM-adjusted p-value (one-sided), LightGBM, Full
\providecommand{\rhmcsOosMBmAllRawLgbmFull}{}\renewcommand{\rhmcsOosMBmAllRawLgbmFull}{0.81}  % single-test bootstrap p-value (one-sided), LightGBM, Full
\providecommand{\rhmcsOosMBmAllPTwoLgbmFull}{}\renewcommand{\rhmcsOosMBmAllPTwoLgbmFull}{0.70}  % StepM-adjusted p-value (two-sided), LightGBM, Full
\providecommand{\rhmcsOosMBmAllRawTwoLgbmFull}{}\renewcommand{\rhmcsOosMBmAllRawTwoLgbmFull}{0.40}  % single-test bootstrap p-value (two-sided), LightGBM, Full
\providecommand{\rhmcsOosMBmAllPXgbHar}{}\renewcommand{\rhmcsOosMBmAllPXgbHar}{0.21}  % StepM-adjusted p-value (one-sided), XGBoost, HAR
\providecommand{\rhmcsOosMBmAllRawXgbHar}{}\renewcommand{\rhmcsOosMBmAllRawXgbHar}{0.094}  % single-test bootstrap p-value (one-sided), XGBoost, HAR
\providecommand{\rhmcsOosMBmAllPTwoXgbHar}{}\renewcommand{\rhmcsOosMBmAllPTwoXgbHar}{0.40}  % StepM-adjusted p-value (two-sided), XGBoost, HAR
\providecommand{\rhmcsOosMBmAllRawTwoXgbHar}{}\renewcommand{\rhmcsOosMBmAllRawTwoXgbHar}{0.16}  % single-test bootstrap p-value (two-sided), XGBoost, HAR
\providecommand{\rhmcsOosMBmAllPXgbFull}{}\renewcommand{\rhmcsOosMBmAllPXgbFull}{0.86}  % StepM-adjusted p-value (one-sided), XGBoost, Full
\providecommand{\rhmcsOosMBmAllRawXgbFull}{}\renewcommand{\rhmcsOosMBmAllRawXgbFull}{0.76}  % single-test bootstrap p-value (one-sided), XGBoost, Full
\providecommand{\rhmcsOosMBmAllPTwoXgbFull}{}\renewcommand{\rhmcsOosMBmAllPTwoXgbFull}{0.82}  % StepM-adjusted p-value (two-sided), XGBoost, Full
\providecommand{\rhmcsOosMBmAllRawTwoXgbFull}{}\renewcommand{\rhmcsOosMBmAllRawTwoXgbFull}{0.55}  % single-test bootstrap p-value (two-sided), XGBoost, Full
\providecommand{\rhmcsOosMBmAllPRfHar}{}\renewcommand{\rhmcsOosMBmAllPRfHar}{0.21}  % StepM-adjusted p-value (one-sided), Random forest, HAR
\providecommand{\rhmcsOosMBmAllRawRfHar}{}\renewcommand{\rhmcsOosMBmAllRawRfHar}{0.100}  % single-test bootstrap p-value (one-sided), Random forest, HAR
\providecommand{\rhmcsOosMBmAllPTwoRfHar}{}\renewcommand{\rhmcsOosMBmAllPTwoRfHar}{0.42}  % StepM-adjusted p-value (two-sided), Random forest, HAR
\providecommand{\rhmcsOosMBmAllRawTwoRfHar}{}\renewcommand{\rhmcsOosMBmAllRawTwoRfHar}{0.18}  % single-test bootstrap p-value (two-sided), Random forest, HAR
\providecommand{\rhmcsOosMBmAllPRfFull}{}\renewcommand{\rhmcsOosMBmAllPRfFull}{0.69}  % StepM-adjusted p-value (one-sided), Random forest, Full
\providecommand{\rhmcsOosMBmAllRawRfFull}{}\renewcommand{\rhmcsOosMBmAllRawRfFull}{0.33}  % single-test bootstrap p-value (one-sided), Random forest, Full
\providecommand{\rhmcsOosMBmAllPTwoRfFull}{}\renewcommand{\rhmcsOosMBmAllPTwoRfFull}{0.85}  % StepM-adjusted p-value (two-sided), Random forest, Full
\providecommand{\rhmcsOosMBmAllRawTwoRfFull}{}\renewcommand{\rhmcsOosMBmAllRawTwoRfFull}{0.72}  % single-test bootstrap p-value (two-sided), Random forest, Full
\providecommand{\rhmcsOosMBmAllNRejLo}{}\renewcommand{\rhmcsOosMBmAllNRejLo}{0}  % StepM rejections at 5 %, block length 5
\providecommand{\rhmcsOosMBmAllSpaCLo}{}\renewcommand{\rhmcsOosMBmAllSpaCLo}{0.14}  % SPA p-value (consistent), block length 5
\providecommand{\rhmcsOosMBmAllNRejTwoLo}{}\renewcommand{\rhmcsOosMBmAllNRejTwoLo}{0}  % ... two-sided, block length 5
\providecommand{\rhmcsOosMBmAllPMinLo}{}\renewcommand{\rhmcsOosMBmAllPMinLo}{0.14}  % smallest adjusted p, block length 5
\providecommand{\rhmcsOosMBmAllPTwoMinLo}{}\renewcommand{\rhmcsOosMBmAllPTwoMinLo}{0.30}  % ... two-sided, block length 5
\providecommand{\rhmcsOosMBmAllNRejHi}{}\renewcommand{\rhmcsOosMBmAllNRejHi}{0}  % StepM rejections at 5 %, block length 20
\providecommand{\rhmcsOosMBmAllSpaCHi}{}\renewcommand{\rhmcsOosMBmAllSpaCHi}{0.31}  % SPA p-value (consistent), block length 20
\providecommand{\rhmcsOosMBmAllNRejTwoHi}{}\renewcommand{\rhmcsOosMBmAllNRejTwoHi}{0}  % ... two-sided, block length 20
\providecommand{\rhmcsOosMBmAllPMinHi}{}\renewcommand{\rhmcsOosMBmAllPMinHi}{0.31}  % smallest adjusted p, block length 20
\providecommand{\rhmcsOosMBmAllPTwoMinHi}{}\renewcommand{\rhmcsOosMBmAllPTwoMinHi}{0.63}  % ... two-sided, block length 20
\providecommand{\rhmcsChQBmSpaC}{}\renewcommand{\rhmcsChQBmSpaC}{0.54}  % SPA p-value (consistent), family ChQBm
\providecommand{\rhmcsChQBmSpaL}{}\renewcommand{\rhmcsChQBmSpaL}{0.11}  % ... lower
\providecommand{\rhmcsChQBmSpaU}{}\renewcommand{\rhmcsChQBmSpaU}{0.67}  % ... upper
\providecommand{\rhmcsChQBmNRej}{}\renewcommand{\rhmcsChQBmNRej}{0}  % StepM rejections at FWER 5 % (of 6)
\providecommand{\rhmcsChQBmNRejWord}{}\renewcommand{\rhmcsChQBmNRejWord}{none}  % ... as a word
\providecommand{\rhmcsChQBmNRaw}{}\renewcommand{\rhmcsChQBmNRaw}{0}  % single-test bootstrap p < 0.05 (no adjustment)
\providecommand{\rhmcsChQBmNRawWord}{}\renewcommand{\rhmcsChQBmNRawWord}{none}  % ... as a word
\providecommand{\rhmcsChQBmNPos}{}\renewcommand{\rhmcsChQBmNPos}{0}  % alternatives with a positive estimate (better than the benchmark)
\providecommand{\rhmcsChQBmPMin}{}\renewcommand{\rhmcsChQBmPMin}{0.95}  % smallest StepM-adjusted p-value
\providecommand{\rhmcsChQBmPMax}{}\renewcommand{\rhmcsChQBmPMax}{1.00}  % largest StepM-adjusted p-value
\providecommand{\rhmcsChQBmPTwoMin}{}\renewcommand{\rhmcsChQBmPTwoMin}{0.004}  % smallest two-sided StepM-adjusted p-value
\providecommand{\rhmcsChQBmPTwoMax}{}\renewcommand{\rhmcsChQBmPTwoMax}{0.22}  % largest two-sided StepM-adjusted p-value
\providecommand{\rhmcsChQBmRawMin}{}\renewcommand{\rhmcsChQBmRawMin}{0.89}  % smallest single-test bootstrap p-value
\providecommand{\rhmcsChQBmRej}{}\renewcommand{\rhmcsChQBmRej}{none}  % the alternatives StepM rejects at 5 % (none: none)
\providecommand{\rhmcsChQBmNRejTwo}{}\renewcommand{\rhmcsChQBmNRejTwo}{4}  % StepM rejections at FWER 5 %, two-sided
\providecommand{\rhmcsChQBmNRejTwoWord}{}\renewcommand{\rhmcsChQBmNRejTwoWord}{four}  % ... as a word
\providecommand{\rhmcsChQBmNRawTwo}{}\renewcommand{\rhmcsChQBmNRawTwo}{4}  % single-test bootstrap p < 0.05, two-sided
\providecommand{\rhmcsChQBmRejTwo}{}\renewcommand{\rhmcsChQBmRejTwo}{ridge regression on the full predictor set, the lasso on the full predictor set, the elastic net on the full predictor set, and the random forest on the full predictor set}  % the alternatives the two-sided StepM rejects at 5 % (none: none)
\providecommand{\rhmcsChQBmPRidge}{}\renewcommand{\rhmcsChQBmPRidge}{1.00}  % StepM-adjusted p-value (one-sided), Ridge
\providecommand{\rhmcsChQBmRawRidge}{}\renewcommand{\rhmcsChQBmRawRidge}{1.00}  % single-test bootstrap p-value (one-sided), Ridge
\providecommand{\rhmcsChQBmPTwoRidge}{}\renewcommand{\rhmcsChQBmPTwoRidge}{0.008}  % StepM-adjusted p-value (two-sided), Ridge
\providecommand{\rhmcsChQBmRawTwoRidge}{}\renewcommand{\rhmcsChQBmRawTwoRidge}{0.003}  % single-test bootstrap p-value (two-sided), Ridge
\providecommand{\rhmcsChQBmPLasso}{}\renewcommand{\rhmcsChQBmPLasso}{1.00}  % StepM-adjusted p-value (one-sided), Lasso
\providecommand{\rhmcsChQBmRawLasso}{}\renewcommand{\rhmcsChQBmRawLasso}{1.00}  % single-test bootstrap p-value (one-sided), Lasso
\providecommand{\rhmcsChQBmPTwoLasso}{}\renewcommand{\rhmcsChQBmPTwoLasso}{0.004}  % StepM-adjusted p-value (two-sided), Lasso
\providecommand{\rhmcsChQBmRawTwoLasso}{}\renewcommand{\rhmcsChQBmRawTwoLasso}{0.001}  % single-test bootstrap p-value (two-sided), Lasso
\providecommand{\rhmcsChQBmPEnet}{}\renewcommand{\rhmcsChQBmPEnet}{1.00}  % StepM-adjusted p-value (one-sided), Elastic net
\providecommand{\rhmcsChQBmRawEnet}{}\renewcommand{\rhmcsChQBmRawEnet}{1.00}  % single-test bootstrap p-value (one-sided), Elastic net
\providecommand{\rhmcsChQBmPTwoEnet}{}\renewcommand{\rhmcsChQBmPTwoEnet}{0.013}  % StepM-adjusted p-value (two-sided), Elastic net
\providecommand{\rhmcsChQBmRawTwoEnet}{}\renewcommand{\rhmcsChQBmRawTwoEnet}{0.005}  % single-test bootstrap p-value (two-sided), Elastic net
\providecommand{\rhmcsChQBmPLgbm}{}\renewcommand{\rhmcsChQBmPLgbm}{0.97}  % StepM-adjusted p-value (one-sided), LightGBM
\providecommand{\rhmcsChQBmRawLgbm}{}\renewcommand{\rhmcsChQBmRawLgbm}{0.91}  % single-test bootstrap p-value (one-sided), LightGBM
\providecommand{\rhmcsChQBmPTwoLgbm}{}\renewcommand{\rhmcsChQBmPTwoLgbm}{0.19}  % StepM-adjusted p-value (two-sided), LightGBM
\providecommand{\rhmcsChQBmRawTwoLgbm}{}\renewcommand{\rhmcsChQBmRawTwoLgbm}{0.17}  % single-test bootstrap p-value (two-sided), LightGBM
\providecommand{\rhmcsChQBmPXgb}{}\renewcommand{\rhmcsChQBmPXgb}{0.95}  % StepM-adjusted p-value (one-sided), XGBoost
\providecommand{\rhmcsChQBmRawXgb}{}\renewcommand{\rhmcsChQBmRawXgb}{0.89}  % single-test bootstrap p-value (one-sided), XGBoost
\providecommand{\rhmcsChQBmPTwoXgb}{}\renewcommand{\rhmcsChQBmPTwoXgb}{0.22}  % StepM-adjusted p-value (two-sided), XGBoost
\providecommand{\rhmcsChQBmRawTwoXgb}{}\renewcommand{\rhmcsChQBmRawTwoXgb}{0.22}  % single-test bootstrap p-value (two-sided), XGBoost
\providecommand{\rhmcsChQBmPRf}{}\renewcommand{\rhmcsChQBmPRf}{1.00}  % StepM-adjusted p-value (one-sided), Random forest
\providecommand{\rhmcsChQBmRawRf}{}\renewcommand{\rhmcsChQBmRawRf}{0.99}  % single-test bootstrap p-value (one-sided), Random forest
\providecommand{\rhmcsChQBmPTwoRf}{}\renewcommand{\rhmcsChQBmPTwoRf}{0.013}  % StepM-adjusted p-value (two-sided), Random forest
\providecommand{\rhmcsChQBmRawTwoRf}{}\renewcommand{\rhmcsChQBmRawTwoRf}{0.008}  % single-test bootstrap p-value (two-sided), Random forest
\providecommand{\rhmcsChQBmNRejLo}{}\renewcommand{\rhmcsChQBmNRejLo}{0}  % StepM rejections at 5 %, block length 5
\providecommand{\rhmcsChQBmSpaCLo}{}\renewcommand{\rhmcsChQBmSpaCLo}{0.54}  % SPA p-value (consistent), block length 5
\providecommand{\rhmcsChQBmNRejTwoLo}{}\renewcommand{\rhmcsChQBmNRejTwoLo}{4}  % ... two-sided, block length 5
\providecommand{\rhmcsChQBmPMinLo}{}\renewcommand{\rhmcsChQBmPMinLo}{0.96}  % smallest adjusted p, block length 5
\providecommand{\rhmcsChQBmPTwoMinLo}{}\renewcommand{\rhmcsChQBmPTwoMinLo}{0.003}  % ... two-sided, block length 5
\providecommand{\rhmcsChQBmNRejHi}{}\renewcommand{\rhmcsChQBmNRejHi}{0}  % StepM rejections at 5 %, block length 20
\providecommand{\rhmcsChQBmSpaCHi}{}\renewcommand{\rhmcsChQBmSpaCHi}{0.54}  % SPA p-value (consistent), block length 20
\providecommand{\rhmcsChQBmNRejTwoHi}{}\renewcommand{\rhmcsChQBmNRejTwoHi}{4}  % ... two-sided, block length 20
\providecommand{\rhmcsChQBmPMinHi}{}\renewcommand{\rhmcsChQBmPMinHi}{0.96}  % smallest adjusted p, block length 20
\providecommand{\rhmcsChQBmPTwoMinHi}{}\renewcommand{\rhmcsChQBmPTwoMinHi}{0.004}  % ... two-sided, block length 20
\providecommand{\rhmcsChQBmAllSpaC}{}\renewcommand{\rhmcsChQBmAllSpaC}{0.54}  % SPA p-value (consistent), family ChQBmAll
\providecommand{\rhmcsChQBmAllSpaL}{}\renewcommand{\rhmcsChQBmAllSpaL}{0.11}  % ... lower
\providecommand{\rhmcsChQBmAllSpaU}{}\renewcommand{\rhmcsChQBmAllSpaU}{0.78}  % ... upper
\providecommand{\rhmcsChQBmAllNRej}{}\renewcommand{\rhmcsChQBmAllNRej}{0}  % StepM rejections at FWER 5 % (of 10)
\providecommand{\rhmcsChQBmAllNRejWord}{}\renewcommand{\rhmcsChQBmAllNRejWord}{none}  % ... as a word
\providecommand{\rhmcsChQBmAllNRaw}{}\renewcommand{\rhmcsChQBmAllNRaw}{0}  % single-test bootstrap p < 0.05 (no adjustment)
\providecommand{\rhmcsChQBmAllNRawWord}{}\renewcommand{\rhmcsChQBmAllNRawWord}{none}  % ... as a word
\providecommand{\rhmcsChQBmAllNPos}{}\renewcommand{\rhmcsChQBmAllNPos}{0}  % alternatives with a positive estimate (better than the benchmark)
\providecommand{\rhmcsChQBmAllPMin}{}\renewcommand{\rhmcsChQBmAllPMin}{0.98}  % smallest StepM-adjusted p-value
\providecommand{\rhmcsChQBmAllPMax}{}\renewcommand{\rhmcsChQBmAllPMax}{1.00}  % largest StepM-adjusted p-value
\providecommand{\rhmcsChQBmAllPTwoMin}{}\renewcommand{\rhmcsChQBmAllPTwoMin}{\ensuremath{<}0.001}  % smallest two-sided StepM-adjusted p-value
\providecommand{\rhmcsChQBmAllPTwoMax}{}\renewcommand{\rhmcsChQBmAllPTwoMax}{0.22}  % largest two-sided StepM-adjusted p-value
\providecommand{\rhmcsChQBmAllRawMin}{}\renewcommand{\rhmcsChQBmAllRawMin}{0.89}  % smallest single-test bootstrap p-value
\providecommand{\rhmcsChQBmAllRej}{}\renewcommand{\rhmcsChQBmAllRej}{none}  % the alternatives StepM rejects at 5 % (none: none)
\providecommand{\rhmcsChQBmAllNRejTwo}{}\renewcommand{\rhmcsChQBmAllNRejTwo}{8}  % StepM rejections at FWER 5 %, two-sided
\providecommand{\rhmcsChQBmAllNRejTwoWord}{}\renewcommand{\rhmcsChQBmAllNRejTwoWord}{eight}  % ... as a word
\providecommand{\rhmcsChQBmAllNRawTwo}{}\renewcommand{\rhmcsChQBmAllNRawTwo}{8}  % single-test bootstrap p < 0.05, two-sided
\providecommand{\rhmcsChQBmAllRejTwo}{}\renewcommand{\rhmcsChQBmAllRejTwo}{the linear models on the HAR benchmark, ridge regression on the full predictor set, the lasso on the full predictor set, the elastic net on the full predictor set, LightGBM on the HAR benchmark, XGBoost on the HAR benchmark, the random forest on the HAR benchmark, and the random forest on the full predictor set}  % the alternatives the two-sided StepM rejects at 5 % (none: none)
\providecommand{\rhmcsChQBmAllPLinHar}{}\renewcommand{\rhmcsChQBmAllPLinHar}{1.00}  % StepM-adjusted p-value (one-sided), Linear, HAR
\providecommand{\rhmcsChQBmAllRawLinHar}{}\renewcommand{\rhmcsChQBmAllRawLinHar}{1.00}  % single-test bootstrap p-value (one-sided), Linear, HAR
\providecommand{\rhmcsChQBmAllPTwoLinHar}{}\renewcommand{\rhmcsChQBmAllPTwoLinHar}{\ensuremath{<}0.001}  % StepM-adjusted p-value (two-sided), Linear, HAR
\providecommand{\rhmcsChQBmAllRawTwoLinHar}{}\renewcommand{\rhmcsChQBmAllRawTwoLinHar}{\ensuremath{<}0.001}  % single-test bootstrap p-value (two-sided), Linear, HAR
\providecommand{\rhmcsChQBmAllPRidgeFull}{}\renewcommand{\rhmcsChQBmAllPRidgeFull}{1.00}  % StepM-adjusted p-value (one-sided), Ridge, Full
\providecommand{\rhmcsChQBmAllRawRidgeFull}{}\renewcommand{\rhmcsChQBmAllRawRidgeFull}{1.00}  % single-test bootstrap p-value (one-sided), Ridge, Full
\providecommand{\rhmcsChQBmAllPTwoRidgeFull}{}\renewcommand{\rhmcsChQBmAllPTwoRidgeFull}{0.008}  % StepM-adjusted p-value (two-sided), Ridge, Full
\providecommand{\rhmcsChQBmAllRawTwoRidgeFull}{}\renewcommand{\rhmcsChQBmAllRawTwoRidgeFull}{0.003}  % single-test bootstrap p-value (two-sided), Ridge, Full
\providecommand{\rhmcsChQBmAllPLassoFull}{}\renewcommand{\rhmcsChQBmAllPLassoFull}{1.00}  % StepM-adjusted p-value (one-sided), Lasso, Full
\providecommand{\rhmcsChQBmAllRawLassoFull}{}\renewcommand{\rhmcsChQBmAllRawLassoFull}{1.00}  % single-test bootstrap p-value (one-sided), Lasso, Full
\providecommand{\rhmcsChQBmAllPTwoLassoFull}{}\renewcommand{\rhmcsChQBmAllPTwoLassoFull}{0.004}  % StepM-adjusted p-value (two-sided), Lasso, Full
\providecommand{\rhmcsChQBmAllRawTwoLassoFull}{}\renewcommand{\rhmcsChQBmAllRawTwoLassoFull}{0.001}  % single-test bootstrap p-value (two-sided), Lasso, Full
\providecommand{\rhmcsChQBmAllPEnetFull}{}\renewcommand{\rhmcsChQBmAllPEnetFull}{1.00}  % StepM-adjusted p-value (one-sided), Elastic net, Full
\providecommand{\rhmcsChQBmAllRawEnetFull}{}\renewcommand{\rhmcsChQBmAllRawEnetFull}{1.00}  % single-test bootstrap p-value (one-sided), Elastic net, Full
\providecommand{\rhmcsChQBmAllPTwoEnetFull}{}\renewcommand{\rhmcsChQBmAllPTwoEnetFull}{0.013}  % StepM-adjusted p-value (two-sided), Elastic net, Full
\providecommand{\rhmcsChQBmAllRawTwoEnetFull}{}\renewcommand{\rhmcsChQBmAllRawTwoEnetFull}{0.005}  % single-test bootstrap p-value (two-sided), Elastic net, Full
\providecommand{\rhmcsChQBmAllPLgbmHar}{}\renewcommand{\rhmcsChQBmAllPLgbmHar}{1.00}  % StepM-adjusted p-value (one-sided), LightGBM, HAR
\providecommand{\rhmcsChQBmAllRawLgbmHar}{}\renewcommand{\rhmcsChQBmAllRawLgbmHar}{1.00}  % single-test bootstrap p-value (one-sided), LightGBM, HAR
\providecommand{\rhmcsChQBmAllPTwoLgbmHar}{}\renewcommand{\rhmcsChQBmAllPTwoLgbmHar}{0.003}  % StepM-adjusted p-value (two-sided), LightGBM, HAR
\providecommand{\rhmcsChQBmAllRawTwoLgbmHar}{}\renewcommand{\rhmcsChQBmAllRawTwoLgbmHar}{\ensuremath{<}0.001}  % single-test bootstrap p-value (two-sided), LightGBM, HAR
\providecommand{\rhmcsChQBmAllPLgbmFull}{}\renewcommand{\rhmcsChQBmAllPLgbmFull}{0.99}  % StepM-adjusted p-value (one-sided), LightGBM, Full
\providecommand{\rhmcsChQBmAllRawLgbmFull}{}\renewcommand{\rhmcsChQBmAllRawLgbmFull}{0.91}  % single-test bootstrap p-value (one-sided), LightGBM, Full
\providecommand{\rhmcsChQBmAllPTwoLgbmFull}{}\renewcommand{\rhmcsChQBmAllPTwoLgbmFull}{0.19}  % StepM-adjusted p-value (two-sided), LightGBM, Full
\providecommand{\rhmcsChQBmAllRawTwoLgbmFull}{}\renewcommand{\rhmcsChQBmAllRawTwoLgbmFull}{0.17}  % single-test bootstrap p-value (two-sided), LightGBM, Full
\providecommand{\rhmcsChQBmAllPXgbHar}{}\renewcommand{\rhmcsChQBmAllPXgbHar}{1.00}  % StepM-adjusted p-value (one-sided), XGBoost, HAR
\providecommand{\rhmcsChQBmAllRawXgbHar}{}\renewcommand{\rhmcsChQBmAllRawXgbHar}{1.00}  % single-test bootstrap p-value (one-sided), XGBoost, HAR
\providecommand{\rhmcsChQBmAllPTwoXgbHar}{}\renewcommand{\rhmcsChQBmAllPTwoXgbHar}{0.002}  % StepM-adjusted p-value (two-sided), XGBoost, HAR
\providecommand{\rhmcsChQBmAllRawTwoXgbHar}{}\renewcommand{\rhmcsChQBmAllRawTwoXgbHar}{\ensuremath{<}0.001}  % single-test bootstrap p-value (two-sided), XGBoost, HAR
\providecommand{\rhmcsChQBmAllPXgbFull}{}\renewcommand{\rhmcsChQBmAllPXgbFull}{0.98}  % StepM-adjusted p-value (one-sided), XGBoost, Full
\providecommand{\rhmcsChQBmAllRawXgbFull}{}\renewcommand{\rhmcsChQBmAllRawXgbFull}{0.89}  % single-test bootstrap p-value (one-sided), XGBoost, Full
\providecommand{\rhmcsChQBmAllPTwoXgbFull}{}\renewcommand{\rhmcsChQBmAllPTwoXgbFull}{0.22}  % StepM-adjusted p-value (two-sided), XGBoost, Full
\providecommand{\rhmcsChQBmAllRawTwoXgbFull}{}\renewcommand{\rhmcsChQBmAllRawTwoXgbFull}{0.22}  % single-test bootstrap p-value (two-sided), XGBoost, Full
\providecommand{\rhmcsChQBmAllPRfHar}{}\renewcommand{\rhmcsChQBmAllPRfHar}{1.00}  % StepM-adjusted p-value (one-sided), Random forest, HAR
\providecommand{\rhmcsChQBmAllRawRfHar}{}\renewcommand{\rhmcsChQBmAllRawRfHar}{1.00}  % single-test bootstrap p-value (one-sided), Random forest, HAR
\providecommand{\rhmcsChQBmAllPTwoRfHar}{}\renewcommand{\rhmcsChQBmAllPTwoRfHar}{\ensuremath{<}0.001}  % StepM-adjusted p-value (two-sided), Random forest, HAR
\providecommand{\rhmcsChQBmAllRawTwoRfHar}{}\renewcommand{\rhmcsChQBmAllRawTwoRfHar}{\ensuremath{<}0.001}  % single-test bootstrap p-value (two-sided), Random forest, HAR
\providecommand{\rhmcsChQBmAllPRfFull}{}\renewcommand{\rhmcsChQBmAllPRfFull}{1.00}  % StepM-adjusted p-value (one-sided), Random forest, Full
\providecommand{\rhmcsChQBmAllRawRfFull}{}\renewcommand{\rhmcsChQBmAllRawRfFull}{0.99}  % single-test bootstrap p-value (one-sided), Random forest, Full
\providecommand{\rhmcsChQBmAllPTwoRfFull}{}\renewcommand{\rhmcsChQBmAllPTwoRfFull}{0.013}  % StepM-adjusted p-value (two-sided), Random forest, Full
\providecommand{\rhmcsChQBmAllRawTwoRfFull}{}\renewcommand{\rhmcsChQBmAllRawTwoRfFull}{0.008}  % single-test bootstrap p-value (two-sided), Random forest, Full
\providecommand{\rhmcsChQBmAllNRejLo}{}\renewcommand{\rhmcsChQBmAllNRejLo}{0}  % StepM rejections at 5 %, block length 5
\providecommand{\rhmcsChQBmAllSpaCLo}{}\renewcommand{\rhmcsChQBmAllSpaCLo}{0.54}  % SPA p-value (consistent), block length 5
\providecommand{\rhmcsChQBmAllNRejTwoLo}{}\renewcommand{\rhmcsChQBmAllNRejTwoLo}{8}  % ... two-sided, block length 5
\providecommand{\rhmcsChQBmAllPMinLo}{}\renewcommand{\rhmcsChQBmAllPMinLo}{0.98}  % smallest adjusted p, block length 5
\providecommand{\rhmcsChQBmAllPTwoMinLo}{}\renewcommand{\rhmcsChQBmAllPTwoMinLo}{\ensuremath{<}0.001}  % ... two-sided, block length 5
\providecommand{\rhmcsChQBmAllNRejHi}{}\renewcommand{\rhmcsChQBmAllNRejHi}{0}  % StepM rejections at 5 %, block length 20
\providecommand{\rhmcsChQBmAllSpaCHi}{}\renewcommand{\rhmcsChQBmAllSpaCHi}{0.54}  % SPA p-value (consistent), block length 20
\providecommand{\rhmcsChQBmAllNRejTwoHi}{}\renewcommand{\rhmcsChQBmAllNRejTwoHi}{8}  % ... two-sided, block length 20
\providecommand{\rhmcsChQBmAllPMinHi}{}\renewcommand{\rhmcsChQBmAllPMinHi}{0.99}  % smallest adjusted p, block length 20
\providecommand{\rhmcsChQBmAllPTwoMinHi}{}\renewcommand{\rhmcsChQBmAllPTwoMinHi}{\ensuremath{<}0.001}  % ... two-sided, block length 20
\providecommand{\rhmcsChMBmSpaC}{}\renewcommand{\rhmcsChMBmSpaC}{0.66}  % SPA p-value (consistent), family ChMBm
\providecommand{\rhmcsChMBmSpaL}{}\renewcommand{\rhmcsChMBmSpaL}{0.36}  % ... lower
\providecommand{\rhmcsChMBmSpaU}{}\renewcommand{\rhmcsChMBmSpaU}{0.66}  % ... upper
\providecommand{\rhmcsChMBmNRej}{}\renewcommand{\rhmcsChMBmNRej}{0}  % StepM rejections at FWER 5 % (of 6)
\providecommand{\rhmcsChMBmNRejWord}{}\renewcommand{\rhmcsChMBmNRejWord}{none}  % ... as a word
\providecommand{\rhmcsChMBmNRaw}{}\renewcommand{\rhmcsChMBmNRaw}{0}  % single-test bootstrap p < 0.05 (no adjustment)
\providecommand{\rhmcsChMBmNRawWord}{}\renewcommand{\rhmcsChMBmNRawWord}{none}  % ... as a word
\providecommand{\rhmcsChMBmNPos}{}\renewcommand{\rhmcsChMBmNPos}{1}  % alternatives with a positive estimate (better than the benchmark)
\providecommand{\rhmcsChMBmPMin}{}\renewcommand{\rhmcsChMBmPMin}{0.66}  % smallest StepM-adjusted p-value
\providecommand{\rhmcsChMBmPMax}{}\renewcommand{\rhmcsChMBmPMax}{0.88}  % largest StepM-adjusted p-value
\providecommand{\rhmcsChMBmPTwoMin}{}\renewcommand{\rhmcsChMBmPTwoMin}{0.51}  % smallest two-sided StepM-adjusted p-value
\providecommand{\rhmcsChMBmPTwoMax}{}\renewcommand{\rhmcsChMBmPTwoMax}{0.77}  % largest two-sided StepM-adjusted p-value
\providecommand{\rhmcsChMBmRawMin}{}\renewcommand{\rhmcsChMBmRawMin}{0.36}  % smallest single-test bootstrap p-value
\providecommand{\rhmcsChMBmRej}{}\renewcommand{\rhmcsChMBmRej}{none}  % the alternatives StepM rejects at 5 % (none: none)
\providecommand{\rhmcsChMBmNRejTwo}{}\renewcommand{\rhmcsChMBmNRejTwo}{0}  % StepM rejections at FWER 5 %, two-sided
\providecommand{\rhmcsChMBmNRejTwoWord}{}\renewcommand{\rhmcsChMBmNRejTwoWord}{none}  % ... as a word
\providecommand{\rhmcsChMBmNRawTwo}{}\renewcommand{\rhmcsChMBmNRawTwo}{0}  % single-test bootstrap p < 0.05, two-sided
\providecommand{\rhmcsChMBmRejTwo}{}\renewcommand{\rhmcsChMBmRejTwo}{none}  % the alternatives the two-sided StepM rejects at 5 % (none: none)
\providecommand{\rhmcsChMBmPRidge}{}\renewcommand{\rhmcsChMBmPRidge}{0.88}  % StepM-adjusted p-value (one-sided), Ridge
\providecommand{\rhmcsChMBmRawRidge}{}\renewcommand{\rhmcsChMBmRawRidge}{0.87}  % single-test bootstrap p-value (one-sided), Ridge
\providecommand{\rhmcsChMBmPTwoRidge}{}\renewcommand{\rhmcsChMBmPTwoRidge}{0.53}  % StepM-adjusted p-value (two-sided), Ridge
\providecommand{\rhmcsChMBmRawTwoRidge}{}\renewcommand{\rhmcsChMBmRawTwoRidge}{0.29}  % single-test bootstrap p-value (two-sided), Ridge
\providecommand{\rhmcsChMBmPLasso}{}\renewcommand{\rhmcsChMBmPLasso}{0.88}  % StepM-adjusted p-value (one-sided), Lasso
\providecommand{\rhmcsChMBmRawLasso}{}\renewcommand{\rhmcsChMBmRawLasso}{0.87}  % single-test bootstrap p-value (one-sided), Lasso
\providecommand{\rhmcsChMBmPTwoLasso}{}\renewcommand{\rhmcsChMBmPTwoLasso}{0.51}  % StepM-adjusted p-value (two-sided), Lasso
\providecommand{\rhmcsChMBmRawTwoLasso}{}\renewcommand{\rhmcsChMBmRawTwoLasso}{0.28}  % single-test bootstrap p-value (two-sided), Lasso
\providecommand{\rhmcsChMBmPEnet}{}\renewcommand{\rhmcsChMBmPEnet}{0.88}  % StepM-adjusted p-value (one-sided), Elastic net
\providecommand{\rhmcsChMBmRawEnet}{}\renewcommand{\rhmcsChMBmRawEnet}{0.87}  % single-test bootstrap p-value (one-sided), Elastic net
\providecommand{\rhmcsChMBmPTwoEnet}{}\renewcommand{\rhmcsChMBmPTwoEnet}{0.53}  % StepM-adjusted p-value (two-sided), Elastic net
\providecommand{\rhmcsChMBmRawTwoEnet}{}\renewcommand{\rhmcsChMBmRawTwoEnet}{0.28}  % single-test bootstrap p-value (two-sided), Elastic net
\providecommand{\rhmcsChMBmPLgbm}{}\renewcommand{\rhmcsChMBmPLgbm}{0.88}  % StepM-adjusted p-value (one-sided), LightGBM
\providecommand{\rhmcsChMBmRawLgbm}{}\renewcommand{\rhmcsChMBmRawLgbm}{0.82}  % single-test bootstrap p-value (one-sided), LightGBM
\providecommand{\rhmcsChMBmPTwoLgbm}{}\renewcommand{\rhmcsChMBmPTwoLgbm}{0.57}  % StepM-adjusted p-value (two-sided), LightGBM
\providecommand{\rhmcsChMBmRawTwoLgbm}{}\renewcommand{\rhmcsChMBmRawTwoLgbm}{0.36}  % single-test bootstrap p-value (two-sided), LightGBM
\providecommand{\rhmcsChMBmPXgb}{}\renewcommand{\rhmcsChMBmPXgb}{0.87}  % StepM-adjusted p-value (one-sided), XGBoost
\providecommand{\rhmcsChMBmRawXgb}{}\renewcommand{\rhmcsChMBmRawXgb}{0.76}  % single-test bootstrap p-value (one-sided), XGBoost
\providecommand{\rhmcsChMBmPTwoXgb}{}\renewcommand{\rhmcsChMBmPTwoXgb}{0.75}  % StepM-adjusted p-value (two-sided), XGBoost
\providecommand{\rhmcsChMBmRawTwoXgb}{}\renewcommand{\rhmcsChMBmRawTwoXgb}{0.55}  % single-test bootstrap p-value (two-sided), XGBoost
\providecommand{\rhmcsChMBmPRf}{}\renewcommand{\rhmcsChMBmPRf}{0.66}  % StepM-adjusted p-value (one-sided), Random forest
\providecommand{\rhmcsChMBmRawRf}{}\renewcommand{\rhmcsChMBmRawRf}{0.36}  % single-test bootstrap p-value (one-sided), Random forest
\providecommand{\rhmcsChMBmPTwoRf}{}\renewcommand{\rhmcsChMBmPTwoRf}{0.77}  % StepM-adjusted p-value (two-sided), Random forest
\providecommand{\rhmcsChMBmRawTwoRf}{}\renewcommand{\rhmcsChMBmRawTwoRf}{0.77}  % single-test bootstrap p-value (two-sided), Random forest
\providecommand{\rhmcsChMBmNRejLo}{}\renewcommand{\rhmcsChMBmNRejLo}{0}  % StepM rejections at 5 %, block length 5
\providecommand{\rhmcsChMBmSpaCLo}{}\renewcommand{\rhmcsChMBmSpaCLo}{0.74}  % SPA p-value (consistent), block length 5
\providecommand{\rhmcsChMBmNRejTwoLo}{}\renewcommand{\rhmcsChMBmNRejTwoLo}{0}  % ... two-sided, block length 5
\providecommand{\rhmcsChMBmPMinLo}{}\renewcommand{\rhmcsChMBmPMinLo}{0.74}  % smallest adjusted p, block length 5
\providecommand{\rhmcsChMBmPTwoMinLo}{}\renewcommand{\rhmcsChMBmPTwoMinLo}{0.34}  % ... two-sided, block length 5
\providecommand{\rhmcsChMBmNRejHi}{}\renewcommand{\rhmcsChMBmNRejHi}{0}  % StepM rejections at 5 %, block length 20
\providecommand{\rhmcsChMBmSpaCHi}{}\renewcommand{\rhmcsChMBmSpaCHi}{0.62}  % SPA p-value (consistent), block length 20
\providecommand{\rhmcsChMBmNRejTwoHi}{}\renewcommand{\rhmcsChMBmNRejTwoHi}{0}  % ... two-sided, block length 20
\providecommand{\rhmcsChMBmPMinHi}{}\renewcommand{\rhmcsChMBmPMinHi}{0.62}  % smallest adjusted p, block length 20
\providecommand{\rhmcsChMBmPTwoMinHi}{}\renewcommand{\rhmcsChMBmPTwoMinHi}{0.65}  % ... two-sided, block length 20
\providecommand{\rhmcsChMBmAllSpaC}{}\renewcommand{\rhmcsChMBmAllSpaC}{0.253}  % SPA p-value (consistent), family ChMBmAll
\providecommand{\rhmcsChMBmAllSpaL}{}\renewcommand{\rhmcsChMBmAllSpaL}{0.23}  % ... lower
\providecommand{\rhmcsChMBmAllSpaU}{}\renewcommand{\rhmcsChMBmAllSpaU}{0.253}  % ... upper
\providecommand{\rhmcsChMBmAllNRej}{}\renewcommand{\rhmcsChMBmAllNRej}{0}  % StepM rejections at FWER 5 % (of 10)
\providecommand{\rhmcsChMBmAllNRejWord}{}\renewcommand{\rhmcsChMBmAllNRejWord}{none}  % ... as a word
\providecommand{\rhmcsChMBmAllNRaw}{}\renewcommand{\rhmcsChMBmAllNRaw}{0}  % single-test bootstrap p < 0.05 (no adjustment)
\providecommand{\rhmcsChMBmAllNRawWord}{}\renewcommand{\rhmcsChMBmAllNRawWord}{none}  % ... as a word
\providecommand{\rhmcsChMBmAllNPos}{}\renewcommand{\rhmcsChMBmAllNPos}{5}  % alternatives with a positive estimate (better than the benchmark)
\providecommand{\rhmcsChMBmAllPMin}{}\renewcommand{\rhmcsChMBmAllPMin}{0.253}  % smallest StepM-adjusted p-value
\providecommand{\rhmcsChMBmAllPMax}{}\renewcommand{\rhmcsChMBmAllPMax}{0.88}  % largest StepM-adjusted p-value
\providecommand{\rhmcsChMBmAllPTwoMin}{}\renewcommand{\rhmcsChMBmAllPTwoMin}{0.48}  % smallest two-sided StepM-adjusted p-value
\providecommand{\rhmcsChMBmAllPTwoMax}{}\renewcommand{\rhmcsChMBmAllPTwoMax}{0.90}  % largest two-sided StepM-adjusted p-value
\providecommand{\rhmcsChMBmAllRawMin}{}\renewcommand{\rhmcsChMBmAllRawMin}{0.11}  % smallest single-test bootstrap p-value
\providecommand{\rhmcsChMBmAllRej}{}\renewcommand{\rhmcsChMBmAllRej}{none}  % the alternatives StepM rejects at 5 % (none: none)
\providecommand{\rhmcsChMBmAllNRejTwo}{}\renewcommand{\rhmcsChMBmAllNRejTwo}{0}  % StepM rejections at FWER 5 %, two-sided
\providecommand{\rhmcsChMBmAllNRejTwoWord}{}\renewcommand{\rhmcsChMBmAllNRejTwoWord}{none}  % ... as a word
\providecommand{\rhmcsChMBmAllNRawTwo}{}\renewcommand{\rhmcsChMBmAllNRawTwo}{0}  % single-test bootstrap p < 0.05, two-sided
\providecommand{\rhmcsChMBmAllRejTwo}{}\renewcommand{\rhmcsChMBmAllRejTwo}{none}  % the alternatives the two-sided StepM rejects at 5 % (none: none)
\providecommand{\rhmcsChMBmAllPLinHar}{}\renewcommand{\rhmcsChMBmAllPLinHar}{0.72}  % StepM-adjusted p-value (one-sided), Linear, HAR
\providecommand{\rhmcsChMBmAllRawLinHar}{}\renewcommand{\rhmcsChMBmAllRawLinHar}{0.41}  % single-test bootstrap p-value (one-sided), Linear, HAR
\providecommand{\rhmcsChMBmAllPTwoLinHar}{}\renewcommand{\rhmcsChMBmAllPTwoLinHar}{0.90}  % StepM-adjusted p-value (two-sided), Linear, HAR
\providecommand{\rhmcsChMBmAllRawTwoLinHar}{}\renewcommand{\rhmcsChMBmAllRawTwoLinHar}{0.86}  % single-test bootstrap p-value (two-sided), Linear, HAR
\providecommand{\rhmcsChMBmAllPRidgeFull}{}\renewcommand{\rhmcsChMBmAllPRidgeFull}{0.88}  % StepM-adjusted p-value (one-sided), Ridge, Full
\providecommand{\rhmcsChMBmAllRawRidgeFull}{}\renewcommand{\rhmcsChMBmAllRawRidgeFull}{0.87}  % single-test bootstrap p-value (one-sided), Ridge, Full
\providecommand{\rhmcsChMBmAllPTwoRidgeFull}{}\renewcommand{\rhmcsChMBmAllPTwoRidgeFull}{0.63}  % StepM-adjusted p-value (two-sided), Ridge, Full
\providecommand{\rhmcsChMBmAllRawTwoRidgeFull}{}\renewcommand{\rhmcsChMBmAllRawTwoRidgeFull}{0.29}  % single-test bootstrap p-value (two-sided), Ridge, Full
\providecommand{\rhmcsChMBmAllPLassoFull}{}\renewcommand{\rhmcsChMBmAllPLassoFull}{0.88}  % StepM-adjusted p-value (one-sided), Lasso, Full
\providecommand{\rhmcsChMBmAllRawLassoFull}{}\renewcommand{\rhmcsChMBmAllRawLassoFull}{0.87}  % single-test bootstrap p-value (one-sided), Lasso, Full
\providecommand{\rhmcsChMBmAllPTwoLassoFull}{}\renewcommand{\rhmcsChMBmAllPTwoLassoFull}{0.61}  % StepM-adjusted p-value (two-sided), Lasso, Full
\providecommand{\rhmcsChMBmAllRawTwoLassoFull}{}\renewcommand{\rhmcsChMBmAllRawTwoLassoFull}{0.28}  % single-test bootstrap p-value (two-sided), Lasso, Full
\providecommand{\rhmcsChMBmAllPEnetFull}{}\renewcommand{\rhmcsChMBmAllPEnetFull}{0.88}  % StepM-adjusted p-value (one-sided), Elastic net, Full
\providecommand{\rhmcsChMBmAllRawEnetFull}{}\renewcommand{\rhmcsChMBmAllRawEnetFull}{0.87}  % single-test bootstrap p-value (one-sided), Elastic net, Full
\providecommand{\rhmcsChMBmAllPTwoEnetFull}{}\renewcommand{\rhmcsChMBmAllPTwoEnetFull}{0.64}  % StepM-adjusted p-value (two-sided), Elastic net, Full
\providecommand{\rhmcsChMBmAllRawTwoEnetFull}{}\renewcommand{\rhmcsChMBmAllRawTwoEnetFull}{0.28}  % single-test bootstrap p-value (two-sided), Elastic net, Full
\providecommand{\rhmcsChMBmAllPLgbmHar}{}\renewcommand{\rhmcsChMBmAllPLgbmHar}{0.48}  % StepM-adjusted p-value (one-sided), LightGBM, HAR
\providecommand{\rhmcsChMBmAllRawLgbmHar}{}\renewcommand{\rhmcsChMBmAllRawLgbmHar}{0.15}  % single-test bootstrap p-value (one-sided), LightGBM, HAR
\providecommand{\rhmcsChMBmAllPTwoLgbmHar}{}\renewcommand{\rhmcsChMBmAllPTwoLgbmHar}{0.68}  % StepM-adjusted p-value (two-sided), LightGBM, HAR
\providecommand{\rhmcsChMBmAllRawTwoLgbmHar}{}\renewcommand{\rhmcsChMBmAllRawTwoLgbmHar}{0.30}  % single-test bootstrap p-value (two-sided), LightGBM, HAR
\providecommand{\rhmcsChMBmAllPLgbmFull}{}\renewcommand{\rhmcsChMBmAllPLgbmFull}{0.88}  % StepM-adjusted p-value (one-sided), LightGBM, Full
\providecommand{\rhmcsChMBmAllRawLgbmFull}{}\renewcommand{\rhmcsChMBmAllRawLgbmFull}{0.82}  % single-test bootstrap p-value (one-sided), LightGBM, Full
\providecommand{\rhmcsChMBmAllPTwoLgbmFull}{}\renewcommand{\rhmcsChMBmAllPTwoLgbmFull}{0.68}  % StepM-adjusted p-value (two-sided), LightGBM, Full
\providecommand{\rhmcsChMBmAllRawTwoLgbmFull}{}\renewcommand{\rhmcsChMBmAllRawTwoLgbmFull}{0.36}  % single-test bootstrap p-value (two-sided), LightGBM, Full
\providecommand{\rhmcsChMBmAllPXgbHar}{}\renewcommand{\rhmcsChMBmAllPXgbHar}{0.255}  % StepM-adjusted p-value (one-sided), XGBoost, HAR
\providecommand{\rhmcsChMBmAllRawXgbHar}{}\renewcommand{\rhmcsChMBmAllRawXgbHar}{0.11}  % single-test bootstrap p-value (one-sided), XGBoost, HAR
\providecommand{\rhmcsChMBmAllPTwoXgbHar}{}\renewcommand{\rhmcsChMBmAllPTwoXgbHar}{0.49}  % StepM-adjusted p-value (two-sided), XGBoost, HAR
\providecommand{\rhmcsChMBmAllRawTwoXgbHar}{}\renewcommand{\rhmcsChMBmAllRawTwoXgbHar}{0.21}  % single-test bootstrap p-value (two-sided), XGBoost, HAR
\providecommand{\rhmcsChMBmAllPXgbFull}{}\renewcommand{\rhmcsChMBmAllPXgbFull}{0.87}  % StepM-adjusted p-value (one-sided), XGBoost, Full
\providecommand{\rhmcsChMBmAllRawXgbFull}{}\renewcommand{\rhmcsChMBmAllRawXgbFull}{0.76}  % single-test bootstrap p-value (one-sided), XGBoost, Full
\providecommand{\rhmcsChMBmAllPTwoXgbFull}{}\renewcommand{\rhmcsChMBmAllPTwoXgbFull}{0.82}  % StepM-adjusted p-value (two-sided), XGBoost, Full
\providecommand{\rhmcsChMBmAllRawTwoXgbFull}{}\renewcommand{\rhmcsChMBmAllRawTwoXgbFull}{0.55}  % single-test bootstrap p-value (two-sided), XGBoost, Full
\providecommand{\rhmcsChMBmAllPRfHar}{}\renewcommand{\rhmcsChMBmAllPRfHar}{0.253}  % StepM-adjusted p-value (one-sided), Random forest, HAR
\providecommand{\rhmcsChMBmAllRawRfHar}{}\renewcommand{\rhmcsChMBmAllRawRfHar}{0.11}  % single-test bootstrap p-value (one-sided), Random forest, HAR
\providecommand{\rhmcsChMBmAllPTwoRfHar}{}\renewcommand{\rhmcsChMBmAllPTwoRfHar}{0.48}  % StepM-adjusted p-value (two-sided), Random forest, HAR
\providecommand{\rhmcsChMBmAllRawTwoRfHar}{}\renewcommand{\rhmcsChMBmAllRawTwoRfHar}{0.21}  % single-test bootstrap p-value (two-sided), Random forest, HAR
\providecommand{\rhmcsChMBmAllPRfFull}{}\renewcommand{\rhmcsChMBmAllPRfFull}{0.71}  % StepM-adjusted p-value (one-sided), Random forest, Full
\providecommand{\rhmcsChMBmAllRawRfFull}{}\renewcommand{\rhmcsChMBmAllRawRfFull}{0.36}  % single-test bootstrap p-value (one-sided), Random forest, Full
\providecommand{\rhmcsChMBmAllPTwoRfFull}{}\renewcommand{\rhmcsChMBmAllPTwoRfFull}{0.90}  % StepM-adjusted p-value (two-sided), Random forest, Full
\providecommand{\rhmcsChMBmAllRawTwoRfFull}{}\renewcommand{\rhmcsChMBmAllRawTwoRfFull}{0.77}  % single-test bootstrap p-value (two-sided), Random forest, Full
\providecommand{\rhmcsChMBmAllNRejLo}{}\renewcommand{\rhmcsChMBmAllNRejLo}{0}  % StepM rejections at 5 %, block length 5
\providecommand{\rhmcsChMBmAllSpaCLo}{}\renewcommand{\rhmcsChMBmAllSpaCLo}{0.21}  % SPA p-value (consistent), block length 5
\providecommand{\rhmcsChMBmAllNRejTwoLo}{}\renewcommand{\rhmcsChMBmAllNRejTwoLo}{0}  % ... two-sided, block length 5
\providecommand{\rhmcsChMBmAllPMinLo}{}\renewcommand{\rhmcsChMBmAllPMinLo}{0.21}  % smallest adjusted p, block length 5
\providecommand{\rhmcsChMBmAllPTwoMinLo}{}\renewcommand{\rhmcsChMBmAllPTwoMinLo}{0.40}  % ... two-sided, block length 5
\providecommand{\rhmcsChMBmAllNRejHi}{}\renewcommand{\rhmcsChMBmAllNRejHi}{0}  % StepM rejections at 5 %, block length 20
\providecommand{\rhmcsChMBmAllSpaCHi}{}\renewcommand{\rhmcsChMBmAllSpaCHi}{0.39}  % SPA p-value (consistent), block length 20
\providecommand{\rhmcsChMBmAllNRejTwoHi}{}\renewcommand{\rhmcsChMBmAllNRejTwoHi}{0}  % ... two-sided, block length 20
\providecommand{\rhmcsChMBmAllPMinHi}{}\renewcommand{\rhmcsChMBmAllPMinHi}{0.39}  % smallest adjusted p, block length 20
\providecommand{\rhmcsChMBmAllPTwoMinHi}{}\renewcommand{\rhmcsChMBmAllPTwoMinHi}{0.69}  % ... two-sided, block length 20
\providecommand{\rhmcsMidShortMeanSpaC}{}\renewcommand{\rhmcsMidShortMeanSpaC}{0.026}  % SPA p-value (consistent), family MidShortMean
\providecommand{\rhmcsMidShortMeanSpaL}{}\renewcommand{\rhmcsMidShortMeanSpaL}{0.026}  % ... lower
\providecommand{\rhmcsMidShortMeanSpaU}{}\renewcommand{\rhmcsMidShortMeanSpaU}{0.026}  % ... upper
\providecommand{\rhmcsMidShortMeanNRej}{}\renewcommand{\rhmcsMidShortMeanNRej}{3}  % StepM rejections at FWER 5 % (of 6)
\providecommand{\rhmcsMidShortMeanNRejWord}{}\renewcommand{\rhmcsMidShortMeanNRejWord}{three}  % ... as a word
\providecommand{\rhmcsMidShortMeanNRaw}{}\renewcommand{\rhmcsMidShortMeanNRaw}{6}  % single-test bootstrap p < 0.05 (no adjustment)
\providecommand{\rhmcsMidShortMeanNRawWord}{}\renewcommand{\rhmcsMidShortMeanNRawWord}{six}  % ... as a word
\providecommand{\rhmcsMidShortMeanNPos}{}\renewcommand{\rhmcsMidShortMeanNPos}{6}  % alternatives with a positive estimate (better than the benchmark)
\providecommand{\rhmcsMidShortMeanPMin}{}\renewcommand{\rhmcsMidShortMeanPMin}{0.026}  % smallest StepM-adjusted p-value
\providecommand{\rhmcsMidShortMeanPMax}{}\renewcommand{\rhmcsMidShortMeanPMax}{0.051}  % largest StepM-adjusted p-value
\providecommand{\rhmcsMidShortMeanPTwoMin}{}\renewcommand{\rhmcsMidShortMeanPTwoMin}{0.034}  % smallest two-sided StepM-adjusted p-value
\providecommand{\rhmcsMidShortMeanPTwoMax}{}\renewcommand{\rhmcsMidShortMeanPTwoMax}{0.083}  % largest two-sided StepM-adjusted p-value
\providecommand{\rhmcsMidShortMeanRawMin}{}\renewcommand{\rhmcsMidShortMeanRawMin}{0.015}  % smallest single-test bootstrap p-value
\providecommand{\rhmcsMidShortMeanRej}{}\renewcommand{\rhmcsMidShortMeanRej}{the lasso on the full predictor set, the elastic net on the full predictor set, and the random forest on the full predictor set}  % the alternatives StepM rejects at 5 % (none: none)
\providecommand{\rhmcsMidShortMeanNRejTwo}{}\renewcommand{\rhmcsMidShortMeanNRejTwo}{2}  % StepM rejections at FWER 5 %, two-sided
\providecommand{\rhmcsMidShortMeanNRejTwoWord}{}\renewcommand{\rhmcsMidShortMeanNRejTwoWord}{two}  % ... as a word
\providecommand{\rhmcsMidShortMeanNRawTwo}{}\renewcommand{\rhmcsMidShortMeanNRawTwo}{3}  % single-test bootstrap p < 0.05, two-sided
\providecommand{\rhmcsMidShortMeanRejTwo}{}\renewcommand{\rhmcsMidShortMeanRejTwo}{the elastic net on the full predictor set and the random forest on the full predictor set}  % the alternatives the two-sided StepM rejects at 5 % (none: none)
\providecommand{\rhmcsMidShortMeanPRidge}{}\renewcommand{\rhmcsMidShortMeanPRidge}{0.051}  % StepM-adjusted p-value (one-sided), Ridge
\providecommand{\rhmcsMidShortMeanRawRidge}{}\renewcommand{\rhmcsMidShortMeanRawRidge}{0.037}  % single-test bootstrap p-value (one-sided), Ridge
\providecommand{\rhmcsMidShortMeanPTwoRidge}{}\renewcommand{\rhmcsMidShortMeanPTwoRidge}{0.083}  % StepM-adjusted p-value (two-sided), Ridge
\providecommand{\rhmcsMidShortMeanRawTwoRidge}{}\renewcommand{\rhmcsMidShortMeanRawTwoRidge}{0.057}  % single-test bootstrap p-value (two-sided), Ridge
\providecommand{\rhmcsMidShortMeanPLasso}{}\renewcommand{\rhmcsMidShortMeanPLasso}{0.038}  % StepM-adjusted p-value (one-sided), Lasso
\providecommand{\rhmcsMidShortMeanRawLasso}{}\renewcommand{\rhmcsMidShortMeanRawLasso}{0.025}  % single-test bootstrap p-value (one-sided), Lasso
\providecommand{\rhmcsMidShortMeanPTwoLasso}{}\renewcommand{\rhmcsMidShortMeanPTwoLasso}{0.059}  % StepM-adjusted p-value (two-sided), Lasso
\providecommand{\rhmcsMidShortMeanRawTwoLasso}{}\renewcommand{\rhmcsMidShortMeanRawTwoLasso}{0.036}  % single-test bootstrap p-value (two-sided), Lasso
\providecommand{\rhmcsMidShortMeanPEnet}{}\renewcommand{\rhmcsMidShortMeanPEnet}{0.031}  % StepM-adjusted p-value (one-sided), Elastic net
\providecommand{\rhmcsMidShortMeanRawEnet}{}\renewcommand{\rhmcsMidShortMeanRawEnet}{0.020}  % single-test bootstrap p-value (one-sided), Elastic net
\providecommand{\rhmcsMidShortMeanPTwoEnet}{}\renewcommand{\rhmcsMidShortMeanPTwoEnet}{0.046}  % StepM-adjusted p-value (two-sided), Elastic net
\providecommand{\rhmcsMidShortMeanRawTwoEnet}{}\renewcommand{\rhmcsMidShortMeanRawTwoEnet}{0.025}  % single-test bootstrap p-value (two-sided), Elastic net
\providecommand{\rhmcsMidShortMeanPLgbm}{}\renewcommand{\rhmcsMidShortMeanPLgbm}{0.051}  % StepM-adjusted p-value (one-sided), LightGBM
\providecommand{\rhmcsMidShortMeanRawLgbm}{}\renewcommand{\rhmcsMidShortMeanRawLgbm}{0.041}  % single-test bootstrap p-value (one-sided), LightGBM
\providecommand{\rhmcsMidShortMeanPTwoLgbm}{}\renewcommand{\rhmcsMidShortMeanPTwoLgbm}{0.083}  % StepM-adjusted p-value (two-sided), LightGBM
\providecommand{\rhmcsMidShortMeanRawTwoLgbm}{}\renewcommand{\rhmcsMidShortMeanRawTwoLgbm}{0.067}  % single-test bootstrap p-value (two-sided), LightGBM
\providecommand{\rhmcsMidShortMeanPXgb}{}\renewcommand{\rhmcsMidShortMeanPXgb}{0.051}  % StepM-adjusted p-value (one-sided), XGBoost
\providecommand{\rhmcsMidShortMeanRawXgb}{}\renewcommand{\rhmcsMidShortMeanRawXgb}{0.039}  % single-test bootstrap p-value (one-sided), XGBoost
\providecommand{\rhmcsMidShortMeanPTwoXgb}{}\renewcommand{\rhmcsMidShortMeanPTwoXgb}{0.083}  % StepM-adjusted p-value (two-sided), XGBoost
\providecommand{\rhmcsMidShortMeanRawTwoXgb}{}\renewcommand{\rhmcsMidShortMeanRawTwoXgb}{0.065}  % single-test bootstrap p-value (two-sided), XGBoost
\providecommand{\rhmcsMidShortMeanPRf}{}\renewcommand{\rhmcsMidShortMeanPRf}{0.026}  % StepM-adjusted p-value (one-sided), Random forest
\providecommand{\rhmcsMidShortMeanRawRf}{}\renewcommand{\rhmcsMidShortMeanRawRf}{0.015}  % single-test bootstrap p-value (one-sided), Random forest
\providecommand{\rhmcsMidShortMeanPTwoRf}{}\renewcommand{\rhmcsMidShortMeanPTwoRf}{0.034}  % StepM-adjusted p-value (two-sided), Random forest
\providecommand{\rhmcsMidShortMeanRawTwoRf}{}\renewcommand{\rhmcsMidShortMeanRawTwoRf}{0.020}  % single-test bootstrap p-value (two-sided), Random forest
\providecommand{\rhmcsMidShortMeanNRejLo}{}\renewcommand{\rhmcsMidShortMeanNRejLo}{3}  % StepM rejections at 5 %, block length 5
\providecommand{\rhmcsMidShortMeanSpaCLo}{}\renewcommand{\rhmcsMidShortMeanSpaCLo}{0.037}  % SPA p-value (consistent), block length 5
\providecommand{\rhmcsMidShortMeanNRejTwoLo}{}\renewcommand{\rhmcsMidShortMeanNRejTwoLo}{0}  % ... two-sided, block length 5
\providecommand{\rhmcsMidShortMeanPMinLo}{}\renewcommand{\rhmcsMidShortMeanPMinLo}{0.037}  % smallest adjusted p, block length 5
\providecommand{\rhmcsMidShortMeanPTwoMinLo}{}\renewcommand{\rhmcsMidShortMeanPTwoMinLo}{0.055}  % ... two-sided, block length 5
\providecommand{\rhmcsMidShortMeanNRejHi}{}\renewcommand{\rhmcsMidShortMeanNRejHi}{3}  % StepM rejections at 5 %, block length 20
\providecommand{\rhmcsMidShortMeanSpaCHi}{}\renewcommand{\rhmcsMidShortMeanSpaCHi}{0.028}  % SPA p-value (consistent), block length 20
\providecommand{\rhmcsMidShortMeanNRejTwoHi}{}\renewcommand{\rhmcsMidShortMeanNRejTwoHi}{1}  % ... two-sided, block length 20
\providecommand{\rhmcsMidShortMeanPMinHi}{}\renewcommand{\rhmcsMidShortMeanPMinHi}{0.028}  % smallest adjusted p, block length 20
\providecommand{\rhmcsMidShortMeanPTwoMinHi}{}\renewcommand{\rhmcsMidShortMeanPTwoMinHi}{0.036}  % ... two-sided, block length 20
\providecommand{\rhmcsMidShortSrSpaC}{}\renewcommand{\rhmcsMidShortSrSpaC}{0.008}  % SPA p-value (consistent), family MidShortSr
\providecommand{\rhmcsMidShortSrSpaL}{}\renewcommand{\rhmcsMidShortSrSpaL}{0.008}  % ... lower
\providecommand{\rhmcsMidShortSrSpaU}{}\renewcommand{\rhmcsMidShortSrSpaU}{0.008}  % ... upper
\providecommand{\rhmcsMidShortSrNRej}{}\renewcommand{\rhmcsMidShortSrNRej}{6}  % StepM rejections at FWER 5 % (of 6)
\providecommand{\rhmcsMidShortSrNRejWord}{}\renewcommand{\rhmcsMidShortSrNRejWord}{six}  % ... as a word
\providecommand{\rhmcsMidShortSrNRaw}{}\renewcommand{\rhmcsMidShortSrNRaw}{6}  % single-test bootstrap p < 0.05 (no adjustment)
\providecommand{\rhmcsMidShortSrNRawWord}{}\renewcommand{\rhmcsMidShortSrNRawWord}{six}  % ... as a word
\providecommand{\rhmcsMidShortSrNPos}{}\renewcommand{\rhmcsMidShortSrNPos}{6}  % alternatives with a positive estimate (better than the benchmark)
\providecommand{\rhmcsMidShortSrPMin}{}\renewcommand{\rhmcsMidShortSrPMin}{0.008}  % smallest StepM-adjusted p-value
\providecommand{\rhmcsMidShortSrPMax}{}\renewcommand{\rhmcsMidShortSrPMax}{0.028}  % largest StepM-adjusted p-value
\providecommand{\rhmcsMidShortSrPTwoMin}{}\renewcommand{\rhmcsMidShortSrPTwoMin}{0.016}  % smallest two-sided StepM-adjusted p-value
\providecommand{\rhmcsMidShortSrPTwoMax}{}\renewcommand{\rhmcsMidShortSrPTwoMax}{0.060}  % largest two-sided StepM-adjusted p-value
\providecommand{\rhmcsMidShortSrRawMin}{}\renewcommand{\rhmcsMidShortSrRawMin}{0.003}  % smallest single-test bootstrap p-value
\providecommand{\rhmcsMidShortSrRej}{}\renewcommand{\rhmcsMidShortSrRej}{ridge regression on the full predictor set, the lasso on the full predictor set, the elastic net on the full predictor set, LightGBM on the full predictor set, XGBoost on the full predictor set, and the random forest on the full predictor set}  % the alternatives StepM rejects at 5 % (none: none)
\providecommand{\rhmcsMidShortSrNRejTwo}{}\renewcommand{\rhmcsMidShortSrNRejTwo}{3}  % StepM rejections at FWER 5 %, two-sided
\providecommand{\rhmcsMidShortSrNRejTwoWord}{}\renewcommand{\rhmcsMidShortSrNRejTwoWord}{three}  % ... as a word
\providecommand{\rhmcsMidShortSrNRawTwo}{}\renewcommand{\rhmcsMidShortSrNRawTwo}{6}  % single-test bootstrap p < 0.05, two-sided
\providecommand{\rhmcsMidShortSrRejTwo}{}\renewcommand{\rhmcsMidShortSrRejTwo}{the lasso on the full predictor set, the elastic net on the full predictor set, and the random forest on the full predictor set}  % the alternatives the two-sided StepM rejects at 5 % (none: none)
\providecommand{\rhmcsMidShortSrPRidge}{}\renewcommand{\rhmcsMidShortSrPRidge}{0.028}  % StepM-adjusted p-value (one-sided), Ridge
\providecommand{\rhmcsMidShortSrRawRidge}{}\renewcommand{\rhmcsMidShortSrRawRidge}{0.019}  % single-test bootstrap p-value (one-sided), Ridge
\providecommand{\rhmcsMidShortSrPTwoRidge}{}\renewcommand{\rhmcsMidShortSrPTwoRidge}{0.060}  % StepM-adjusted p-value (two-sided), Ridge
\providecommand{\rhmcsMidShortSrRawTwoRidge}{}\renewcommand{\rhmcsMidShortSrRawTwoRidge}{0.039}  % single-test bootstrap p-value (two-sided), Ridge
\providecommand{\rhmcsMidShortSrPLasso}{}\renewcommand{\rhmcsMidShortSrPLasso}{0.017}  % StepM-adjusted p-value (one-sided), Lasso
\providecommand{\rhmcsMidShortSrRawLasso}{}\renewcommand{\rhmcsMidShortSrRawLasso}{0.010}  % single-test bootstrap p-value (one-sided), Lasso
\providecommand{\rhmcsMidShortSrPTwoLasso}{}\renewcommand{\rhmcsMidShortSrPTwoLasso}{0.040}  % StepM-adjusted p-value (two-sided), Lasso
\providecommand{\rhmcsMidShortSrRawTwoLasso}{}\renewcommand{\rhmcsMidShortSrRawTwoLasso}{0.021}  % single-test bootstrap p-value (two-sided), Lasso
\providecommand{\rhmcsMidShortSrPEnet}{}\renewcommand{\rhmcsMidShortSrPEnet}{0.012}  % StepM-adjusted p-value (one-sided), Elastic net
\providecommand{\rhmcsMidShortSrRawEnet}{}\renewcommand{\rhmcsMidShortSrRawEnet}{0.007}  % single-test bootstrap p-value (one-sided), Elastic net
\providecommand{\rhmcsMidShortSrPTwoEnet}{}\renewcommand{\rhmcsMidShortSrPTwoEnet}{0.027}  % StepM-adjusted p-value (two-sided), Elastic net
\providecommand{\rhmcsMidShortSrRawTwoEnet}{}\renewcommand{\rhmcsMidShortSrRawTwoEnet}{0.012}  % single-test bootstrap p-value (two-sided), Elastic net
\providecommand{\rhmcsMidShortSrPLgbm}{}\renewcommand{\rhmcsMidShortSrPLgbm}{0.028}  % StepM-adjusted p-value (one-sided), LightGBM
\providecommand{\rhmcsMidShortSrRawLgbm}{}\renewcommand{\rhmcsMidShortSrRawLgbm}{0.021}  % single-test bootstrap p-value (one-sided), LightGBM
\providecommand{\rhmcsMidShortSrPTwoLgbm}{}\renewcommand{\rhmcsMidShortSrPTwoLgbm}{0.060}  % StepM-adjusted p-value (two-sided), LightGBM
\providecommand{\rhmcsMidShortSrRawTwoLgbm}{}\renewcommand{\rhmcsMidShortSrRawTwoLgbm}{0.047}  % single-test bootstrap p-value (two-sided), LightGBM
\providecommand{\rhmcsMidShortSrPXgb}{}\renewcommand{\rhmcsMidShortSrPXgb}{0.028}  % StepM-adjusted p-value (one-sided), XGBoost
\providecommand{\rhmcsMidShortSrRawXgb}{}\renewcommand{\rhmcsMidShortSrRawXgb}{0.020}  % single-test bootstrap p-value (one-sided), XGBoost
\providecommand{\rhmcsMidShortSrPTwoXgb}{}\renewcommand{\rhmcsMidShortSrPTwoXgb}{0.060}  % StepM-adjusted p-value (two-sided), XGBoost
\providecommand{\rhmcsMidShortSrRawTwoXgb}{}\renewcommand{\rhmcsMidShortSrRawTwoXgb}{0.046}  % single-test bootstrap p-value (two-sided), XGBoost
\providecommand{\rhmcsMidShortSrPRf}{}\renewcommand{\rhmcsMidShortSrPRf}{0.008}  % StepM-adjusted p-value (one-sided), Random forest
\providecommand{\rhmcsMidShortSrRawRf}{}\renewcommand{\rhmcsMidShortSrRawRf}{0.003}  % single-test bootstrap p-value (one-sided), Random forest
\providecommand{\rhmcsMidShortSrPTwoRf}{}\renewcommand{\rhmcsMidShortSrPTwoRf}{0.016}  % StepM-adjusted p-value (two-sided), Random forest
\providecommand{\rhmcsMidShortSrRawTwoRf}{}\renewcommand{\rhmcsMidShortSrRawTwoRf}{0.007}  % single-test bootstrap p-value (two-sided), Random forest
\providecommand{\rhmcsMidShortSrNRejLo}{}\renewcommand{\rhmcsMidShortSrNRejLo}{6}  % StepM rejections at 5 %, block length 5
\providecommand{\rhmcsMidShortSrSpaCLo}{}\renewcommand{\rhmcsMidShortSrSpaCLo}{0.010}  % SPA p-value (consistent), block length 5
\providecommand{\rhmcsMidShortSrNRejTwoLo}{}\renewcommand{\rhmcsMidShortSrNRejTwoLo}{3}  % ... two-sided, block length 5
\providecommand{\rhmcsMidShortSrPMinLo}{}\renewcommand{\rhmcsMidShortSrPMinLo}{0.010}  % smallest adjusted p, block length 5
\providecommand{\rhmcsMidShortSrPTwoMinLo}{}\renewcommand{\rhmcsMidShortSrPTwoMinLo}{0.027}  % ... two-sided, block length 5
\providecommand{\rhmcsMidShortSrNRejHi}{}\renewcommand{\rhmcsMidShortSrNRejHi}{6}  % StepM rejections at 5 %, block length 20
\providecommand{\rhmcsMidShortSrSpaCHi}{}\renewcommand{\rhmcsMidShortSrSpaCHi}{0.008}  % SPA p-value (consistent), block length 20
\providecommand{\rhmcsMidShortSrNRejTwoHi}{}\renewcommand{\rhmcsMidShortSrNRejTwoHi}{3}  % ... two-sided, block length 20
\providecommand{\rhmcsMidShortSrPMinHi}{}\renewcommand{\rhmcsMidShortSrPMinHi}{0.008}  % smallest adjusted p, block length 20
\providecommand{\rhmcsMidShortSrPTwoMinHi}{}\renewcommand{\rhmcsMidShortSrPTwoMinHi}{0.015}  % ... two-sided, block length 20
\providecommand{\rhmcsMidTwinMeanSpaC}{}\renewcommand{\rhmcsMidTwinMeanSpaC}{0.027}  % SPA p-value (consistent), family MidTwinMean
\providecommand{\rhmcsMidTwinMeanSpaL}{}\renewcommand{\rhmcsMidTwinMeanSpaL}{0.027}  % ... lower
\providecommand{\rhmcsMidTwinMeanSpaU}{}\renewcommand{\rhmcsMidTwinMeanSpaU}{0.027}  % ... upper
\providecommand{\rhmcsMidTwinMeanNRej}{}\renewcommand{\rhmcsMidTwinMeanNRej}{1}  % StepM rejections at FWER 5 % (of 6)
\providecommand{\rhmcsMidTwinMeanNRejWord}{}\renewcommand{\rhmcsMidTwinMeanNRejWord}{one}  % ... as a word
\providecommand{\rhmcsMidTwinMeanNRaw}{}\renewcommand{\rhmcsMidTwinMeanNRaw}{3}  % single-test bootstrap p < 0.05 (no adjustment)
\providecommand{\rhmcsMidTwinMeanNRawWord}{}\renewcommand{\rhmcsMidTwinMeanNRawWord}{three}  % ... as a word
\providecommand{\rhmcsMidTwinMeanNPos}{}\renewcommand{\rhmcsMidTwinMeanNPos}{6}  % alternatives with a positive estimate (better than the benchmark)
\providecommand{\rhmcsMidTwinMeanPMin}{}\renewcommand{\rhmcsMidTwinMeanPMin}{0.027}  % smallest StepM-adjusted p-value
\providecommand{\rhmcsMidTwinMeanPMax}{}\renewcommand{\rhmcsMidTwinMeanPMax}{0.13}  % largest StepM-adjusted p-value
\providecommand{\rhmcsMidTwinMeanPTwoMin}{}\renewcommand{\rhmcsMidTwinMeanPTwoMin}{0.042}  % smallest two-sided StepM-adjusted p-value
\providecommand{\rhmcsMidTwinMeanPTwoMax}{}\renewcommand{\rhmcsMidTwinMeanPTwoMax}{0.26}  % largest two-sided StepM-adjusted p-value
\providecommand{\rhmcsMidTwinMeanRawMin}{}\renewcommand{\rhmcsMidTwinMeanRawMin}{0.013}  % smallest single-test bootstrap p-value
\providecommand{\rhmcsMidTwinMeanRej}{}\renewcommand{\rhmcsMidTwinMeanRej}{the random forest on the full predictor set}  % the alternatives StepM rejects at 5 % (none: none)
\providecommand{\rhmcsMidTwinMeanNRejTwo}{}\renewcommand{\rhmcsMidTwinMeanNRejTwo}{1}  % StepM rejections at FWER 5 %, two-sided
\providecommand{\rhmcsMidTwinMeanNRejTwoWord}{}\renewcommand{\rhmcsMidTwinMeanNRejTwoWord}{one}  % ... as a word
\providecommand{\rhmcsMidTwinMeanNRawTwo}{}\renewcommand{\rhmcsMidTwinMeanNRawTwo}{2}  % single-test bootstrap p < 0.05, two-sided
\providecommand{\rhmcsMidTwinMeanRejTwo}{}\renewcommand{\rhmcsMidTwinMeanRejTwo}{the random forest on the full predictor set}  % the alternatives the two-sided StepM rejects at 5 % (none: none)
\providecommand{\rhmcsMidTwinMeanPRidge}{}\renewcommand{\rhmcsMidTwinMeanPRidge}{0.13}  % StepM-adjusted p-value (one-sided), Ridge
\providecommand{\rhmcsMidTwinMeanRawRidge}{}\renewcommand{\rhmcsMidTwinMeanRawRidge}{0.082}  % single-test bootstrap p-value (one-sided), Ridge
\providecommand{\rhmcsMidTwinMeanPTwoRidge}{}\renewcommand{\rhmcsMidTwinMeanPTwoRidge}{0.26}  % StepM-adjusted p-value (two-sided), Ridge
\providecommand{\rhmcsMidTwinMeanRawTwoRidge}{}\renewcommand{\rhmcsMidTwinMeanRawTwoRidge}{0.16}  % single-test bootstrap p-value (two-sided), Ridge
\providecommand{\rhmcsMidTwinMeanPLasso}{}\renewcommand{\rhmcsMidTwinMeanPLasso}{0.11}  % StepM-adjusted p-value (one-sided), Lasso
\providecommand{\rhmcsMidTwinMeanRawLasso}{}\renewcommand{\rhmcsMidTwinMeanRawLasso}{0.063}  % single-test bootstrap p-value (one-sided), Lasso
\providecommand{\rhmcsMidTwinMeanPTwoLasso}{}\renewcommand{\rhmcsMidTwinMeanPTwoLasso}{0.21}  % StepM-adjusted p-value (two-sided), Lasso
\providecommand{\rhmcsMidTwinMeanRawTwoLasso}{}\renewcommand{\rhmcsMidTwinMeanRawTwoLasso}{0.12}  % single-test bootstrap p-value (two-sided), Lasso
\providecommand{\rhmcsMidTwinMeanPEnet}{}\renewcommand{\rhmcsMidTwinMeanPEnet}{0.069}  % StepM-adjusted p-value (one-sided), Elastic net
\providecommand{\rhmcsMidTwinMeanRawEnet}{}\renewcommand{\rhmcsMidTwinMeanRawEnet}{0.036}  % single-test bootstrap p-value (one-sided), Elastic net
\providecommand{\rhmcsMidTwinMeanPTwoEnet}{}\renewcommand{\rhmcsMidTwinMeanPTwoEnet}{0.13}  % StepM-adjusted p-value (two-sided), Elastic net
\providecommand{\rhmcsMidTwinMeanRawTwoEnet}{}\renewcommand{\rhmcsMidTwinMeanRawTwoEnet}{0.063}  % single-test bootstrap p-value (two-sided), Elastic net
\providecommand{\rhmcsMidTwinMeanPLgbm}{}\renewcommand{\rhmcsMidTwinMeanPLgbm}{0.060}  % StepM-adjusted p-value (one-sided), LightGBM
\providecommand{\rhmcsMidTwinMeanRawLgbm}{}\renewcommand{\rhmcsMidTwinMeanRawLgbm}{0.025}  % single-test bootstrap p-value (one-sided), LightGBM
\providecommand{\rhmcsMidTwinMeanPTwoLgbm}{}\renewcommand{\rhmcsMidTwinMeanPTwoLgbm}{0.11}  % StepM-adjusted p-value (two-sided), LightGBM
\providecommand{\rhmcsMidTwinMeanRawTwoLgbm}{}\renewcommand{\rhmcsMidTwinMeanRawTwoLgbm}{0.047}  % single-test bootstrap p-value (two-sided), LightGBM
\providecommand{\rhmcsMidTwinMeanPXgb}{}\renewcommand{\rhmcsMidTwinMeanPXgb}{0.13}  % StepM-adjusted p-value (one-sided), XGBoost
\providecommand{\rhmcsMidTwinMeanRawXgb}{}\renewcommand{\rhmcsMidTwinMeanRawXgb}{0.101}  % single-test bootstrap p-value (one-sided), XGBoost
\providecommand{\rhmcsMidTwinMeanPTwoXgb}{}\renewcommand{\rhmcsMidTwinMeanPTwoXgb}{0.26}  % StepM-adjusted p-value (two-sided), XGBoost
\providecommand{\rhmcsMidTwinMeanRawTwoXgb}{}\renewcommand{\rhmcsMidTwinMeanRawTwoXgb}{0.20}  % single-test bootstrap p-value (two-sided), XGBoost
\providecommand{\rhmcsMidTwinMeanPRf}{}\renewcommand{\rhmcsMidTwinMeanPRf}{0.027}  % StepM-adjusted p-value (one-sided), Random forest
\providecommand{\rhmcsMidTwinMeanRawRf}{}\renewcommand{\rhmcsMidTwinMeanRawRf}{0.013}  % single-test bootstrap p-value (one-sided), Random forest
\providecommand{\rhmcsMidTwinMeanPTwoRf}{}\renewcommand{\rhmcsMidTwinMeanPTwoRf}{0.042}  % StepM-adjusted p-value (two-sided), Random forest
\providecommand{\rhmcsMidTwinMeanRawTwoRf}{}\renewcommand{\rhmcsMidTwinMeanRawTwoRf}{0.013}  % single-test bootstrap p-value (two-sided), Random forest
\providecommand{\rhmcsMidTwinMeanNRejLo}{}\renewcommand{\rhmcsMidTwinMeanNRejLo}{1}  % StepM rejections at 5 %, block length 5
\providecommand{\rhmcsMidTwinMeanSpaCLo}{}\renewcommand{\rhmcsMidTwinMeanSpaCLo}{0.036}  % SPA p-value (consistent), block length 5
\providecommand{\rhmcsMidTwinMeanNRejTwoLo}{}\renewcommand{\rhmcsMidTwinMeanNRejTwoLo}{0}  % ... two-sided, block length 5
\providecommand{\rhmcsMidTwinMeanPMinLo}{}\renewcommand{\rhmcsMidTwinMeanPMinLo}{0.036}  % smallest adjusted p, block length 5
\providecommand{\rhmcsMidTwinMeanPTwoMinLo}{}\renewcommand{\rhmcsMidTwinMeanPTwoMinLo}{0.053}  % ... two-sided, block length 5
\providecommand{\rhmcsMidTwinMeanNRejHi}{}\renewcommand{\rhmcsMidTwinMeanNRejHi}{1}  % StepM rejections at 5 %, block length 20
\providecommand{\rhmcsMidTwinMeanSpaCHi}{}\renewcommand{\rhmcsMidTwinMeanSpaCHi}{0.033}  % SPA p-value (consistent), block length 20
\providecommand{\rhmcsMidTwinMeanNRejTwoHi}{}\renewcommand{\rhmcsMidTwinMeanNRejTwoHi}{1}  % ... two-sided, block length 20
\providecommand{\rhmcsMidTwinMeanPMinHi}{}\renewcommand{\rhmcsMidTwinMeanPMinHi}{0.033}  % smallest adjusted p, block length 20
\providecommand{\rhmcsMidTwinMeanPTwoMinHi}{}\renewcommand{\rhmcsMidTwinMeanPTwoMinHi}{0.047}  % ... two-sided, block length 20
\providecommand{\rhmcsMidTwinSrSpaC}{}\renewcommand{\rhmcsMidTwinSrSpaC}{0.007}  % SPA p-value (consistent), family MidTwinSr
\providecommand{\rhmcsMidTwinSrSpaL}{}\renewcommand{\rhmcsMidTwinSrSpaL}{0.007}  % ... lower
\providecommand{\rhmcsMidTwinSrSpaU}{}\renewcommand{\rhmcsMidTwinSrSpaU}{0.007}  % ... upper
\providecommand{\rhmcsMidTwinSrNRej}{}\renewcommand{\rhmcsMidTwinSrNRej}{1}  % StepM rejections at FWER 5 % (of 6)
\providecommand{\rhmcsMidTwinSrNRejWord}{}\renewcommand{\rhmcsMidTwinSrNRejWord}{one}  % ... as a word
\providecommand{\rhmcsMidTwinSrNRaw}{}\renewcommand{\rhmcsMidTwinSrNRaw}{3}  % single-test bootstrap p < 0.05 (no adjustment)
\providecommand{\rhmcsMidTwinSrNRawWord}{}\renewcommand{\rhmcsMidTwinSrNRawWord}{three}  % ... as a word
\providecommand{\rhmcsMidTwinSrNPos}{}\renewcommand{\rhmcsMidTwinSrNPos}{6}  % alternatives with a positive estimate (better than the benchmark)
\providecommand{\rhmcsMidTwinSrPMin}{}\renewcommand{\rhmcsMidTwinSrPMin}{0.007}  % smallest StepM-adjusted p-value
\providecommand{\rhmcsMidTwinSrPMax}{}\renewcommand{\rhmcsMidTwinSrPMax}{0.14}  % largest StepM-adjusted p-value
\providecommand{\rhmcsMidTwinSrPTwoMin}{}\renewcommand{\rhmcsMidTwinSrPTwoMin}{0.012}  % smallest two-sided StepM-adjusted p-value
\providecommand{\rhmcsMidTwinSrPTwoMax}{}\renewcommand{\rhmcsMidTwinSrPTwoMax}{0.27}  % largest two-sided StepM-adjusted p-value
\providecommand{\rhmcsMidTwinSrRawMin}{}\renewcommand{\rhmcsMidTwinSrRawMin}{0.002}  % smallest single-test bootstrap p-value
\providecommand{\rhmcsMidTwinSrRej}{}\renewcommand{\rhmcsMidTwinSrRej}{the random forest on the full predictor set}  % the alternatives StepM rejects at 5 % (none: none)
\providecommand{\rhmcsMidTwinSrNRejTwo}{}\renewcommand{\rhmcsMidTwinSrNRejTwo}{1}  % StepM rejections at FWER 5 %, two-sided
\providecommand{\rhmcsMidTwinSrNRejTwoWord}{}\renewcommand{\rhmcsMidTwinSrNRejTwoWord}{one}  % ... as a word
\providecommand{\rhmcsMidTwinSrNRawTwo}{}\renewcommand{\rhmcsMidTwinSrNRawTwo}{2}  % single-test bootstrap p < 0.05, two-sided
\providecommand{\rhmcsMidTwinSrRejTwo}{}\renewcommand{\rhmcsMidTwinSrRejTwo}{the random forest on the full predictor set}  % the alternatives the two-sided StepM rejects at 5 % (none: none)
\providecommand{\rhmcsMidTwinSrPRidge}{}\renewcommand{\rhmcsMidTwinSrPRidge}{0.14}  % StepM-adjusted p-value (one-sided), Ridge
\providecommand{\rhmcsMidTwinSrRawRidge}{}\renewcommand{\rhmcsMidTwinSrRawRidge}{0.089}  % single-test bootstrap p-value (one-sided), Ridge
\providecommand{\rhmcsMidTwinSrPTwoRidge}{}\renewcommand{\rhmcsMidTwinSrPTwoRidge}{0.27}  % StepM-adjusted p-value (two-sided), Ridge
\providecommand{\rhmcsMidTwinSrRawTwoRidge}{}\renewcommand{\rhmcsMidTwinSrRawTwoRidge}{0.17}  % single-test bootstrap p-value (two-sided), Ridge
\providecommand{\rhmcsMidTwinSrPLasso}{}\renewcommand{\rhmcsMidTwinSrPLasso}{0.12}  % StepM-adjusted p-value (one-sided), Lasso
\providecommand{\rhmcsMidTwinSrRawLasso}{}\renewcommand{\rhmcsMidTwinSrRawLasso}{0.071}  % single-test bootstrap p-value (one-sided), Lasso
\providecommand{\rhmcsMidTwinSrPTwoLasso}{}\renewcommand{\rhmcsMidTwinSrPTwoLasso}{0.22}  % StepM-adjusted p-value (two-sided), Lasso
\providecommand{\rhmcsMidTwinSrRawTwoLasso}{}\renewcommand{\rhmcsMidTwinSrRawTwoLasso}{0.12}  % single-test bootstrap p-value (two-sided), Lasso
\providecommand{\rhmcsMidTwinSrPEnet}{}\renewcommand{\rhmcsMidTwinSrPEnet}{0.082}  % StepM-adjusted p-value (one-sided), Elastic net
\providecommand{\rhmcsMidTwinSrRawEnet}{}\renewcommand{\rhmcsMidTwinSrRawEnet}{0.045}  % single-test bootstrap p-value (one-sided), Elastic net
\providecommand{\rhmcsMidTwinSrPTwoEnet}{}\renewcommand{\rhmcsMidTwinSrPTwoEnet}{0.15}  % StepM-adjusted p-value (two-sided), Elastic net
\providecommand{\rhmcsMidTwinSrRawTwoEnet}{}\renewcommand{\rhmcsMidTwinSrRawTwoEnet}{0.072}  % single-test bootstrap p-value (two-sided), Elastic net
\providecommand{\rhmcsMidTwinSrPLgbm}{}\renewcommand{\rhmcsMidTwinSrPLgbm}{0.065}  % StepM-adjusted p-value (one-sided), LightGBM
\providecommand{\rhmcsMidTwinSrRawLgbm}{}\renewcommand{\rhmcsMidTwinSrRawLgbm}{0.027}  % single-test bootstrap p-value (one-sided), LightGBM
\providecommand{\rhmcsMidTwinSrPTwoLgbm}{}\renewcommand{\rhmcsMidTwinSrPTwoLgbm}{0.12}  % StepM-adjusted p-value (two-sided), LightGBM
\providecommand{\rhmcsMidTwinSrRawTwoLgbm}{}\renewcommand{\rhmcsMidTwinSrRawTwoLgbm}{0.050}  % single-test bootstrap p-value (two-sided), LightGBM
\providecommand{\rhmcsMidTwinSrPXgb}{}\renewcommand{\rhmcsMidTwinSrPXgb}{0.14}  % StepM-adjusted p-value (one-sided), XGBoost
\providecommand{\rhmcsMidTwinSrRawXgb}{}\renewcommand{\rhmcsMidTwinSrRawXgb}{0.11}  % single-test bootstrap p-value (one-sided), XGBoost
\providecommand{\rhmcsMidTwinSrPTwoXgb}{}\renewcommand{\rhmcsMidTwinSrPTwoXgb}{0.27}  % StepM-adjusted p-value (two-sided), XGBoost
\providecommand{\rhmcsMidTwinSrRawTwoXgb}{}\renewcommand{\rhmcsMidTwinSrRawTwoXgb}{0.20}  % single-test bootstrap p-value (two-sided), XGBoost
\providecommand{\rhmcsMidTwinSrPRf}{}\renewcommand{\rhmcsMidTwinSrPRf}{0.007}  % StepM-adjusted p-value (one-sided), Random forest
\providecommand{\rhmcsMidTwinSrRawRf}{}\renewcommand{\rhmcsMidTwinSrRawRf}{0.002}  % single-test bootstrap p-value (one-sided), Random forest
\providecommand{\rhmcsMidTwinSrPTwoRf}{}\renewcommand{\rhmcsMidTwinSrPTwoRf}{0.012}  % StepM-adjusted p-value (two-sided), Random forest
\providecommand{\rhmcsMidTwinSrRawTwoRf}{}\renewcommand{\rhmcsMidTwinSrRawTwoRf}{0.003}  % single-test bootstrap p-value (two-sided), Random forest
\providecommand{\rhmcsMidTwinSrNRejLo}{}\renewcommand{\rhmcsMidTwinSrNRejLo}{1}  % StepM rejections at 5 %, block length 5
\providecommand{\rhmcsMidTwinSrSpaCLo}{}\renewcommand{\rhmcsMidTwinSrSpaCLo}{0.011}  % SPA p-value (consistent), block length 5
\providecommand{\rhmcsMidTwinSrNRejTwoLo}{}\renewcommand{\rhmcsMidTwinSrNRejTwoLo}{1}  % ... two-sided, block length 5
\providecommand{\rhmcsMidTwinSrPMinLo}{}\renewcommand{\rhmcsMidTwinSrPMinLo}{0.011}  % smallest adjusted p, block length 5
\providecommand{\rhmcsMidTwinSrPTwoMinLo}{}\renewcommand{\rhmcsMidTwinSrPTwoMinLo}{0.017}  % ... two-sided, block length 5
\providecommand{\rhmcsMidTwinSrNRejHi}{}\renewcommand{\rhmcsMidTwinSrNRejHi}{1}  % StepM rejections at 5 %, block length 20
\providecommand{\rhmcsMidTwinSrSpaCHi}{}\renewcommand{\rhmcsMidTwinSrSpaCHi}{0.010}  % SPA p-value (consistent), block length 20
\providecommand{\rhmcsMidTwinSrNRejTwoHi}{}\renewcommand{\rhmcsMidTwinSrNRejTwoHi}{1}  % ... two-sided, block length 20
\providecommand{\rhmcsMidTwinSrPMinHi}{}\renewcommand{\rhmcsMidTwinSrPMinHi}{0.010}  % smallest adjusted p, block length 20
\providecommand{\rhmcsMidTwinSrPTwoMinHi}{}\renewcommand{\rhmcsMidTwinSrPTwoMinHi}{0.016}  % ... two-sided, block length 20
\providecommand{\rhmcsBaShortMeanSpaC}{}\renewcommand{\rhmcsBaShortMeanSpaC}{0.025}  % SPA p-value (consistent), family BaShortMean
\providecommand{\rhmcsBaShortMeanSpaL}{}\renewcommand{\rhmcsBaShortMeanSpaL}{0.025}  % ... lower
\providecommand{\rhmcsBaShortMeanSpaU}{}\renewcommand{\rhmcsBaShortMeanSpaU}{0.025}  % ... upper
\providecommand{\rhmcsBaShortMeanNRej}{}\renewcommand{\rhmcsBaShortMeanNRej}{3}  % StepM rejections at FWER 5 % (of 6)
\providecommand{\rhmcsBaShortMeanNRejWord}{}\renewcommand{\rhmcsBaShortMeanNRejWord}{three}  % ... as a word
\providecommand{\rhmcsBaShortMeanNRaw}{}\renewcommand{\rhmcsBaShortMeanNRaw}{6}  % single-test bootstrap p < 0.05 (no adjustment)
\providecommand{\rhmcsBaShortMeanNRawWord}{}\renewcommand{\rhmcsBaShortMeanNRawWord}{six}  % ... as a word
\providecommand{\rhmcsBaShortMeanNPos}{}\renewcommand{\rhmcsBaShortMeanNPos}{6}  % alternatives with a positive estimate (better than the benchmark)
\providecommand{\rhmcsBaShortMeanPMin}{}\renewcommand{\rhmcsBaShortMeanPMin}{0.025}  % smallest StepM-adjusted p-value
\providecommand{\rhmcsBaShortMeanPMax}{}\renewcommand{\rhmcsBaShortMeanPMax}{0.050}  % largest StepM-adjusted p-value
\providecommand{\rhmcsBaShortMeanPTwoMin}{}\renewcommand{\rhmcsBaShortMeanPTwoMin}{0.034}  % smallest two-sided StepM-adjusted p-value
\providecommand{\rhmcsBaShortMeanPTwoMax}{}\renewcommand{\rhmcsBaShortMeanPTwoMax}{0.082}  % largest two-sided StepM-adjusted p-value
\providecommand{\rhmcsBaShortMeanRawMin}{}\renewcommand{\rhmcsBaShortMeanRawMin}{0.015}  % smallest single-test bootstrap p-value
\providecommand{\rhmcsBaShortMeanRej}{}\renewcommand{\rhmcsBaShortMeanRej}{the lasso on the full predictor set, the elastic net on the full predictor set, and the random forest on the full predictor set}  % the alternatives StepM rejects at 5 % (none: none)
\providecommand{\rhmcsBaShortMeanNRejTwo}{}\renewcommand{\rhmcsBaShortMeanNRejTwo}{2}  % StepM rejections at FWER 5 %, two-sided
\providecommand{\rhmcsBaShortMeanNRejTwoWord}{}\renewcommand{\rhmcsBaShortMeanNRejTwoWord}{two}  % ... as a word
\providecommand{\rhmcsBaShortMeanNRawTwo}{}\renewcommand{\rhmcsBaShortMeanNRawTwo}{3}  % single-test bootstrap p < 0.05, two-sided
\providecommand{\rhmcsBaShortMeanRejTwo}{}\renewcommand{\rhmcsBaShortMeanRejTwo}{the elastic net on the full predictor set and the random forest on the full predictor set}  % the alternatives the two-sided StepM rejects at 5 % (none: none)
\providecommand{\rhmcsBaShortMeanPRidge}{}\renewcommand{\rhmcsBaShortMeanPRidge}{0.050}  % StepM-adjusted p-value (one-sided), Ridge
\providecommand{\rhmcsBaShortMeanRawRidge}{}\renewcommand{\rhmcsBaShortMeanRawRidge}{0.036}  % single-test bootstrap p-value (one-sided), Ridge
\providecommand{\rhmcsBaShortMeanPTwoRidge}{}\renewcommand{\rhmcsBaShortMeanPTwoRidge}{0.082}  % StepM-adjusted p-value (two-sided), Ridge
\providecommand{\rhmcsBaShortMeanRawTwoRidge}{}\renewcommand{\rhmcsBaShortMeanRawTwoRidge}{0.056}  % single-test bootstrap p-value (two-sided), Ridge
\providecommand{\rhmcsBaShortMeanPLasso}{}\renewcommand{\rhmcsBaShortMeanPLasso}{0.037}  % StepM-adjusted p-value (one-sided), Lasso
\providecommand{\rhmcsBaShortMeanRawLasso}{}\renewcommand{\rhmcsBaShortMeanRawLasso}{0.025}  % single-test bootstrap p-value (one-sided), Lasso
\providecommand{\rhmcsBaShortMeanPTwoLasso}{}\renewcommand{\rhmcsBaShortMeanPTwoLasso}{0.058}  % StepM-adjusted p-value (two-sided), Lasso
\providecommand{\rhmcsBaShortMeanRawTwoLasso}{}\renewcommand{\rhmcsBaShortMeanRawTwoLasso}{0.035}  % single-test bootstrap p-value (two-sided), Lasso
\providecommand{\rhmcsBaShortMeanPEnet}{}\renewcommand{\rhmcsBaShortMeanPEnet}{0.031}  % StepM-adjusted p-value (one-sided), Elastic net
\providecommand{\rhmcsBaShortMeanRawEnet}{}\renewcommand{\rhmcsBaShortMeanRawEnet}{0.020}  % single-test bootstrap p-value (one-sided), Elastic net
\providecommand{\rhmcsBaShortMeanPTwoEnet}{}\renewcommand{\rhmcsBaShortMeanPTwoEnet}{0.045}  % StepM-adjusted p-value (two-sided), Elastic net
\providecommand{\rhmcsBaShortMeanRawTwoEnet}{}\renewcommand{\rhmcsBaShortMeanRawTwoEnet}{0.025}  % single-test bootstrap p-value (two-sided), Elastic net
\providecommand{\rhmcsBaShortMeanPLgbm}{}\renewcommand{\rhmcsBaShortMeanPLgbm}{0.050}  % StepM-adjusted p-value (one-sided), LightGBM
\providecommand{\rhmcsBaShortMeanRawLgbm}{}\renewcommand{\rhmcsBaShortMeanRawLgbm}{0.041}  % single-test bootstrap p-value (one-sided), LightGBM
\providecommand{\rhmcsBaShortMeanPTwoLgbm}{}\renewcommand{\rhmcsBaShortMeanPTwoLgbm}{0.082}  % StepM-adjusted p-value (two-sided), LightGBM
\providecommand{\rhmcsBaShortMeanRawTwoLgbm}{}\renewcommand{\rhmcsBaShortMeanRawTwoLgbm}{0.067}  % single-test bootstrap p-value (two-sided), LightGBM
\providecommand{\rhmcsBaShortMeanPXgb}{}\renewcommand{\rhmcsBaShortMeanPXgb}{0.050}  % StepM-adjusted p-value (one-sided), XGBoost
\providecommand{\rhmcsBaShortMeanRawXgb}{}\renewcommand{\rhmcsBaShortMeanRawXgb}{0.038}  % single-test bootstrap p-value (one-sided), XGBoost
\providecommand{\rhmcsBaShortMeanPTwoXgb}{}\renewcommand{\rhmcsBaShortMeanPTwoXgb}{0.082}  % StepM-adjusted p-value (two-sided), XGBoost
\providecommand{\rhmcsBaShortMeanRawTwoXgb}{}\renewcommand{\rhmcsBaShortMeanRawTwoXgb}{0.063}  % single-test bootstrap p-value (two-sided), XGBoost
\providecommand{\rhmcsBaShortMeanPRf}{}\renewcommand{\rhmcsBaShortMeanPRf}{0.025}  % StepM-adjusted p-value (one-sided), Random forest
\providecommand{\rhmcsBaShortMeanRawRf}{}\renewcommand{\rhmcsBaShortMeanRawRf}{0.015}  % single-test bootstrap p-value (one-sided), Random forest
\providecommand{\rhmcsBaShortMeanPTwoRf}{}\renewcommand{\rhmcsBaShortMeanPTwoRf}{0.034}  % StepM-adjusted p-value (two-sided), Random forest
\providecommand{\rhmcsBaShortMeanRawTwoRf}{}\renewcommand{\rhmcsBaShortMeanRawTwoRf}{0.019}  % single-test bootstrap p-value (two-sided), Random forest
\providecommand{\rhmcsBaShortMeanNRejLo}{}\renewcommand{\rhmcsBaShortMeanNRejLo}{3}  % StepM rejections at 5 %, block length 5
\providecommand{\rhmcsBaShortMeanSpaCLo}{}\renewcommand{\rhmcsBaShortMeanSpaCLo}{0.037}  % SPA p-value (consistent), block length 5
\providecommand{\rhmcsBaShortMeanNRejTwoLo}{}\renewcommand{\rhmcsBaShortMeanNRejTwoLo}{0}  % ... two-sided, block length 5
\providecommand{\rhmcsBaShortMeanPMinLo}{}\renewcommand{\rhmcsBaShortMeanPMinLo}{0.037}  % smallest adjusted p, block length 5
\providecommand{\rhmcsBaShortMeanPTwoMinLo}{}\renewcommand{\rhmcsBaShortMeanPTwoMinLo}{0.054}  % ... two-sided, block length 5
\providecommand{\rhmcsBaShortMeanNRejHi}{}\renewcommand{\rhmcsBaShortMeanNRejHi}{3}  % StepM rejections at 5 %, block length 20
\providecommand{\rhmcsBaShortMeanSpaCHi}{}\renewcommand{\rhmcsBaShortMeanSpaCHi}{0.027}  % SPA p-value (consistent), block length 20
\providecommand{\rhmcsBaShortMeanNRejTwoHi}{}\renewcommand{\rhmcsBaShortMeanNRejTwoHi}{1}  % ... two-sided, block length 20
\providecommand{\rhmcsBaShortMeanPMinHi}{}\renewcommand{\rhmcsBaShortMeanPMinHi}{0.027}  % smallest adjusted p, block length 20
\providecommand{\rhmcsBaShortMeanPTwoMinHi}{}\renewcommand{\rhmcsBaShortMeanPTwoMinHi}{0.035}  % ... two-sided, block length 20
\providecommand{\rhmcsBaShortSrSpaC}{}\renewcommand{\rhmcsBaShortSrSpaC}{0.007}  % SPA p-value (consistent), family BaShortSr
\providecommand{\rhmcsBaShortSrSpaL}{}\renewcommand{\rhmcsBaShortSrSpaL}{0.007}  % ... lower
\providecommand{\rhmcsBaShortSrSpaU}{}\renewcommand{\rhmcsBaShortSrSpaU}{0.007}  % ... upper
\providecommand{\rhmcsBaShortSrNRej}{}\renewcommand{\rhmcsBaShortSrNRej}{6}  % StepM rejections at FWER 5 % (of 6)
\providecommand{\rhmcsBaShortSrNRejWord}{}\renewcommand{\rhmcsBaShortSrNRejWord}{six}  % ... as a word
\providecommand{\rhmcsBaShortSrNRaw}{}\renewcommand{\rhmcsBaShortSrNRaw}{6}  % single-test bootstrap p < 0.05 (no adjustment)
\providecommand{\rhmcsBaShortSrNRawWord}{}\renewcommand{\rhmcsBaShortSrNRawWord}{six}  % ... as a word
\providecommand{\rhmcsBaShortSrNPos}{}\renewcommand{\rhmcsBaShortSrNPos}{6}  % alternatives with a positive estimate (better than the benchmark)
\providecommand{\rhmcsBaShortSrPMin}{}\renewcommand{\rhmcsBaShortSrPMin}{0.007}  % smallest StepM-adjusted p-value
\providecommand{\rhmcsBaShortSrPMax}{}\renewcommand{\rhmcsBaShortSrPMax}{0.027}  % largest StepM-adjusted p-value
\providecommand{\rhmcsBaShortSrPTwoMin}{}\renewcommand{\rhmcsBaShortSrPTwoMin}{0.015}  % smallest two-sided StepM-adjusted p-value
\providecommand{\rhmcsBaShortSrPTwoMax}{}\renewcommand{\rhmcsBaShortSrPTwoMax}{0.058}  % largest two-sided StepM-adjusted p-value
\providecommand{\rhmcsBaShortSrRawMin}{}\renewcommand{\rhmcsBaShortSrRawMin}{0.002}  % smallest single-test bootstrap p-value
\providecommand{\rhmcsBaShortSrRej}{}\renewcommand{\rhmcsBaShortSrRej}{ridge regression on the full predictor set, the lasso on the full predictor set, the elastic net on the full predictor set, LightGBM on the full predictor set, XGBoost on the full predictor set, and the random forest on the full predictor set}  % the alternatives StepM rejects at 5 % (none: none)
\providecommand{\rhmcsBaShortSrNRejTwo}{}\renewcommand{\rhmcsBaShortSrNRejTwo}{3}  % StepM rejections at FWER 5 %, two-sided
\providecommand{\rhmcsBaShortSrNRejTwoWord}{}\renewcommand{\rhmcsBaShortSrNRejTwoWord}{three}  % ... as a word
\providecommand{\rhmcsBaShortSrNRawTwo}{}\renewcommand{\rhmcsBaShortSrNRawTwo}{6}  % single-test bootstrap p < 0.05, two-sided
\providecommand{\rhmcsBaShortSrRejTwo}{}\renewcommand{\rhmcsBaShortSrRejTwo}{the lasso on the full predictor set, the elastic net on the full predictor set, and the random forest on the full predictor set}  % the alternatives the two-sided StepM rejects at 5 % (none: none)
\providecommand{\rhmcsBaShortSrPRidge}{}\renewcommand{\rhmcsBaShortSrPRidge}{0.027}  % StepM-adjusted p-value (one-sided), Ridge
\providecommand{\rhmcsBaShortSrRawRidge}{}\renewcommand{\rhmcsBaShortSrRawRidge}{0.018}  % single-test bootstrap p-value (one-sided), Ridge
\providecommand{\rhmcsBaShortSrPTwoRidge}{}\renewcommand{\rhmcsBaShortSrPTwoRidge}{0.058}  % StepM-adjusted p-value (two-sided), Ridge
\providecommand{\rhmcsBaShortSrRawTwoRidge}{}\renewcommand{\rhmcsBaShortSrRawTwoRidge}{0.037}  % single-test bootstrap p-value (two-sided), Ridge
\providecommand{\rhmcsBaShortSrPLasso}{}\renewcommand{\rhmcsBaShortSrPLasso}{0.017}  % StepM-adjusted p-value (one-sided), Lasso
\providecommand{\rhmcsBaShortSrRawLasso}{}\renewcommand{\rhmcsBaShortSrRawLasso}{0.010}  % single-test bootstrap p-value (one-sided), Lasso
\providecommand{\rhmcsBaShortSrPTwoLasso}{}\renewcommand{\rhmcsBaShortSrPTwoLasso}{0.039}  % StepM-adjusted p-value (two-sided), Lasso
\providecommand{\rhmcsBaShortSrRawTwoLasso}{}\renewcommand{\rhmcsBaShortSrRawTwoLasso}{0.021}  % single-test bootstrap p-value (two-sided), Lasso
\providecommand{\rhmcsBaShortSrPEnet}{}\renewcommand{\rhmcsBaShortSrPEnet}{0.011}  % StepM-adjusted p-value (one-sided), Elastic net
\providecommand{\rhmcsBaShortSrRawEnet}{}\renewcommand{\rhmcsBaShortSrRawEnet}{0.007}  % single-test bootstrap p-value (one-sided), Elastic net
\providecommand{\rhmcsBaShortSrPTwoEnet}{}\renewcommand{\rhmcsBaShortSrPTwoEnet}{0.026}  % StepM-adjusted p-value (two-sided), Elastic net
\providecommand{\rhmcsBaShortSrRawTwoEnet}{}\renewcommand{\rhmcsBaShortSrRawTwoEnet}{0.012}  % single-test bootstrap p-value (two-sided), Elastic net
\providecommand{\rhmcsBaShortSrPLgbm}{}\renewcommand{\rhmcsBaShortSrPLgbm}{0.027}  % StepM-adjusted p-value (one-sided), LightGBM
\providecommand{\rhmcsBaShortSrRawLgbm}{}\renewcommand{\rhmcsBaShortSrRawLgbm}{0.020}  % single-test bootstrap p-value (one-sided), LightGBM
\providecommand{\rhmcsBaShortSrPTwoLgbm}{}\renewcommand{\rhmcsBaShortSrPTwoLgbm}{0.058}  % StepM-adjusted p-value (two-sided), LightGBM
\providecommand{\rhmcsBaShortSrRawTwoLgbm}{}\renewcommand{\rhmcsBaShortSrRawTwoLgbm}{0.045}  % single-test bootstrap p-value (two-sided), LightGBM
\providecommand{\rhmcsBaShortSrPXgb}{}\renewcommand{\rhmcsBaShortSrPXgb}{0.027}  % StepM-adjusted p-value (one-sided), XGBoost
\providecommand{\rhmcsBaShortSrRawXgb}{}\renewcommand{\rhmcsBaShortSrRawXgb}{0.019}  % single-test bootstrap p-value (one-sided), XGBoost
\providecommand{\rhmcsBaShortSrPTwoXgb}{}\renewcommand{\rhmcsBaShortSrPTwoXgb}{0.058}  % StepM-adjusted p-value (two-sided), XGBoost
\providecommand{\rhmcsBaShortSrRawTwoXgb}{}\renewcommand{\rhmcsBaShortSrRawTwoXgb}{0.043}  % single-test bootstrap p-value (two-sided), XGBoost
\providecommand{\rhmcsBaShortSrPRf}{}\renewcommand{\rhmcsBaShortSrPRf}{0.007}  % StepM-adjusted p-value (one-sided), Random forest
\providecommand{\rhmcsBaShortSrRawRf}{}\renewcommand{\rhmcsBaShortSrRawRf}{0.002}  % single-test bootstrap p-value (one-sided), Random forest
\providecommand{\rhmcsBaShortSrPTwoRf}{}\renewcommand{\rhmcsBaShortSrPTwoRf}{0.015}  % StepM-adjusted p-value (two-sided), Random forest
\providecommand{\rhmcsBaShortSrRawTwoRf}{}\renewcommand{\rhmcsBaShortSrRawTwoRf}{0.006}  % single-test bootstrap p-value (two-sided), Random forest
\providecommand{\rhmcsBaShortSrNRejLo}{}\renewcommand{\rhmcsBaShortSrNRejLo}{6}  % StepM rejections at 5 %, block length 5
\providecommand{\rhmcsBaShortSrSpaCLo}{}\renewcommand{\rhmcsBaShortSrSpaCLo}{0.010}  % SPA p-value (consistent), block length 5
\providecommand{\rhmcsBaShortSrNRejTwoLo}{}\renewcommand{\rhmcsBaShortSrNRejTwoLo}{3}  % ... two-sided, block length 5
\providecommand{\rhmcsBaShortSrPMinLo}{}\renewcommand{\rhmcsBaShortSrPMinLo}{0.010}  % smallest adjusted p, block length 5
\providecommand{\rhmcsBaShortSrPTwoMinLo}{}\renewcommand{\rhmcsBaShortSrPTwoMinLo}{0.027}  % ... two-sided, block length 5
\providecommand{\rhmcsBaShortSrNRejHi}{}\renewcommand{\rhmcsBaShortSrNRejHi}{6}  % StepM rejections at 5 %, block length 20
\providecommand{\rhmcsBaShortSrSpaCHi}{}\renewcommand{\rhmcsBaShortSrSpaCHi}{0.008}  % SPA p-value (consistent), block length 20
\providecommand{\rhmcsBaShortSrNRejTwoHi}{}\renewcommand{\rhmcsBaShortSrNRejTwoHi}{3}  % ... two-sided, block length 20
\providecommand{\rhmcsBaShortSrPMinHi}{}\renewcommand{\rhmcsBaShortSrPMinHi}{0.008}  % smallest adjusted p, block length 20
\providecommand{\rhmcsBaShortSrPTwoMinHi}{}\renewcommand{\rhmcsBaShortSrPTwoMinHi}{0.014}  % ... two-sided, block length 20
\providecommand{\rhmcsBaTwinMeanSpaC}{}\renewcommand{\rhmcsBaTwinMeanSpaC}{0.028}  % SPA p-value (consistent), family BaTwinMean
\providecommand{\rhmcsBaTwinMeanSpaL}{}\renewcommand{\rhmcsBaTwinMeanSpaL}{0.028}  % ... lower
\providecommand{\rhmcsBaTwinMeanSpaU}{}\renewcommand{\rhmcsBaTwinMeanSpaU}{0.028}  % ... upper
\providecommand{\rhmcsBaTwinMeanNRej}{}\renewcommand{\rhmcsBaTwinMeanNRej}{1}  % StepM rejections at FWER 5 % (of 6)
\providecommand{\rhmcsBaTwinMeanNRejWord}{}\renewcommand{\rhmcsBaTwinMeanNRejWord}{one}  % ... as a word
\providecommand{\rhmcsBaTwinMeanNRaw}{}\renewcommand{\rhmcsBaTwinMeanNRaw}{3}  % single-test bootstrap p < 0.05 (no adjustment)
\providecommand{\rhmcsBaTwinMeanNRawWord}{}\renewcommand{\rhmcsBaTwinMeanNRawWord}{three}  % ... as a word
\providecommand{\rhmcsBaTwinMeanNPos}{}\renewcommand{\rhmcsBaTwinMeanNPos}{6}  % alternatives with a positive estimate (better than the benchmark)
\providecommand{\rhmcsBaTwinMeanPMin}{}\renewcommand{\rhmcsBaTwinMeanPMin}{0.028}  % smallest StepM-adjusted p-value
\providecommand{\rhmcsBaTwinMeanPMax}{}\renewcommand{\rhmcsBaTwinMeanPMax}{0.13}  % largest StepM-adjusted p-value
\providecommand{\rhmcsBaTwinMeanPTwoMin}{}\renewcommand{\rhmcsBaTwinMeanPTwoMin}{0.042}  % smallest two-sided StepM-adjusted p-value
\providecommand{\rhmcsBaTwinMeanPTwoMax}{}\renewcommand{\rhmcsBaTwinMeanPTwoMax}{0.26}  % largest two-sided StepM-adjusted p-value
\providecommand{\rhmcsBaTwinMeanRawMin}{}\renewcommand{\rhmcsBaTwinMeanRawMin}{0.013}  % smallest single-test bootstrap p-value
\providecommand{\rhmcsBaTwinMeanRej}{}\renewcommand{\rhmcsBaTwinMeanRej}{the random forest on the full predictor set}  % the alternatives StepM rejects at 5 % (none: none)
\providecommand{\rhmcsBaTwinMeanNRejTwo}{}\renewcommand{\rhmcsBaTwinMeanNRejTwo}{1}  % StepM rejections at FWER 5 %, two-sided
\providecommand{\rhmcsBaTwinMeanNRejTwoWord}{}\renewcommand{\rhmcsBaTwinMeanNRejTwoWord}{one}  % ... as a word
\providecommand{\rhmcsBaTwinMeanNRawTwo}{}\renewcommand{\rhmcsBaTwinMeanNRawTwo}{2}  % single-test bootstrap p < 0.05, two-sided
\providecommand{\rhmcsBaTwinMeanRejTwo}{}\renewcommand{\rhmcsBaTwinMeanRejTwo}{the random forest on the full predictor set}  % the alternatives the two-sided StepM rejects at 5 % (none: none)
\providecommand{\rhmcsBaTwinMeanPRidge}{}\renewcommand{\rhmcsBaTwinMeanPRidge}{0.13}  % StepM-adjusted p-value (one-sided), Ridge
\providecommand{\rhmcsBaTwinMeanRawRidge}{}\renewcommand{\rhmcsBaTwinMeanRawRidge}{0.083}  % single-test bootstrap p-value (one-sided), Ridge
\providecommand{\rhmcsBaTwinMeanPTwoRidge}{}\renewcommand{\rhmcsBaTwinMeanPTwoRidge}{0.26}  % StepM-adjusted p-value (two-sided), Ridge
\providecommand{\rhmcsBaTwinMeanRawTwoRidge}{}\renewcommand{\rhmcsBaTwinMeanRawTwoRidge}{0.16}  % single-test bootstrap p-value (two-sided), Ridge
\providecommand{\rhmcsBaTwinMeanPLasso}{}\renewcommand{\rhmcsBaTwinMeanPLasso}{0.11}  % StepM-adjusted p-value (one-sided), Lasso
\providecommand{\rhmcsBaTwinMeanRawLasso}{}\renewcommand{\rhmcsBaTwinMeanRawLasso}{0.063}  % single-test bootstrap p-value (one-sided), Lasso
\providecommand{\rhmcsBaTwinMeanPTwoLasso}{}\renewcommand{\rhmcsBaTwinMeanPTwoLasso}{0.21}  % StepM-adjusted p-value (two-sided), Lasso
\providecommand{\rhmcsBaTwinMeanRawTwoLasso}{}\renewcommand{\rhmcsBaTwinMeanRawTwoLasso}{0.12}  % single-test bootstrap p-value (two-sided), Lasso
\providecommand{\rhmcsBaTwinMeanPEnet}{}\renewcommand{\rhmcsBaTwinMeanPEnet}{0.069}  % StepM-adjusted p-value (one-sided), Elastic net
\providecommand{\rhmcsBaTwinMeanRawEnet}{}\renewcommand{\rhmcsBaTwinMeanRawEnet}{0.036}  % single-test bootstrap p-value (one-sided), Elastic net
\providecommand{\rhmcsBaTwinMeanPTwoEnet}{}\renewcommand{\rhmcsBaTwinMeanPTwoEnet}{0.13}  % StepM-adjusted p-value (two-sided), Elastic net
\providecommand{\rhmcsBaTwinMeanRawTwoEnet}{}\renewcommand{\rhmcsBaTwinMeanRawTwoEnet}{0.063}  % single-test bootstrap p-value (two-sided), Elastic net
\providecommand{\rhmcsBaTwinMeanPLgbm}{}\renewcommand{\rhmcsBaTwinMeanPLgbm}{0.061}  % StepM-adjusted p-value (one-sided), LightGBM
\providecommand{\rhmcsBaTwinMeanRawLgbm}{}\renewcommand{\rhmcsBaTwinMeanRawLgbm}{0.025}  % single-test bootstrap p-value (one-sided), LightGBM
\providecommand{\rhmcsBaTwinMeanPTwoLgbm}{}\renewcommand{\rhmcsBaTwinMeanPTwoLgbm}{0.12}  % StepM-adjusted p-value (two-sided), LightGBM
\providecommand{\rhmcsBaTwinMeanRawTwoLgbm}{}\renewcommand{\rhmcsBaTwinMeanRawTwoLgbm}{0.048}  % single-test bootstrap p-value (two-sided), LightGBM
\providecommand{\rhmcsBaTwinMeanPXgb}{}\renewcommand{\rhmcsBaTwinMeanPXgb}{0.13}  % StepM-adjusted p-value (one-sided), XGBoost
\providecommand{\rhmcsBaTwinMeanRawXgb}{}\renewcommand{\rhmcsBaTwinMeanRawXgb}{0.103}  % single-test bootstrap p-value (one-sided), XGBoost
\providecommand{\rhmcsBaTwinMeanPTwoXgb}{}\renewcommand{\rhmcsBaTwinMeanPTwoXgb}{0.26}  % StepM-adjusted p-value (two-sided), XGBoost
\providecommand{\rhmcsBaTwinMeanRawTwoXgb}{}\renewcommand{\rhmcsBaTwinMeanRawTwoXgb}{0.20}  % single-test bootstrap p-value (two-sided), XGBoost
\providecommand{\rhmcsBaTwinMeanPRf}{}\renewcommand{\rhmcsBaTwinMeanPRf}{0.028}  % StepM-adjusted p-value (one-sided), Random forest
\providecommand{\rhmcsBaTwinMeanRawRf}{}\renewcommand{\rhmcsBaTwinMeanRawRf}{0.013}  % single-test bootstrap p-value (one-sided), Random forest
\providecommand{\rhmcsBaTwinMeanPTwoRf}{}\renewcommand{\rhmcsBaTwinMeanPTwoRf}{0.042}  % StepM-adjusted p-value (two-sided), Random forest
\providecommand{\rhmcsBaTwinMeanRawTwoRf}{}\renewcommand{\rhmcsBaTwinMeanRawTwoRf}{0.014}  % single-test bootstrap p-value (two-sided), Random forest
\providecommand{\rhmcsBaTwinMeanNRejLo}{}\renewcommand{\rhmcsBaTwinMeanNRejLo}{1}  % StepM rejections at 5 %, block length 5
\providecommand{\rhmcsBaTwinMeanSpaCLo}{}\renewcommand{\rhmcsBaTwinMeanSpaCLo}{0.036}  % SPA p-value (consistent), block length 5
\providecommand{\rhmcsBaTwinMeanNRejTwoLo}{}\renewcommand{\rhmcsBaTwinMeanNRejTwoLo}{0}  % ... two-sided, block length 5
\providecommand{\rhmcsBaTwinMeanPMinLo}{}\renewcommand{\rhmcsBaTwinMeanPMinLo}{0.036}  % smallest adjusted p, block length 5
\providecommand{\rhmcsBaTwinMeanPTwoMinLo}{}\renewcommand{\rhmcsBaTwinMeanPTwoMinLo}{0.054}  % ... two-sided, block length 5
\providecommand{\rhmcsBaTwinMeanNRejHi}{}\renewcommand{\rhmcsBaTwinMeanNRejHi}{1}  % StepM rejections at 5 %, block length 20
\providecommand{\rhmcsBaTwinMeanSpaCHi}{}\renewcommand{\rhmcsBaTwinMeanSpaCHi}{0.034}  % SPA p-value (consistent), block length 20
\providecommand{\rhmcsBaTwinMeanNRejTwoHi}{}\renewcommand{\rhmcsBaTwinMeanNRejTwoHi}{1}  % ... two-sided, block length 20
\providecommand{\rhmcsBaTwinMeanPMinHi}{}\renewcommand{\rhmcsBaTwinMeanPMinHi}{0.034}  % smallest adjusted p, block length 20
\providecommand{\rhmcsBaTwinMeanPTwoMinHi}{}\renewcommand{\rhmcsBaTwinMeanPTwoMinHi}{0.048}  % ... two-sided, block length 20
\providecommand{\rhmcsBaTwinSrSpaC}{}\renewcommand{\rhmcsBaTwinSrSpaC}{0.008}  % SPA p-value (consistent), family BaTwinSr
\providecommand{\rhmcsBaTwinSrSpaL}{}\renewcommand{\rhmcsBaTwinSrSpaL}{0.008}  % ... lower
\providecommand{\rhmcsBaTwinSrSpaU}{}\renewcommand{\rhmcsBaTwinSrSpaU}{0.008}  % ... upper
\providecommand{\rhmcsBaTwinSrNRej}{}\renewcommand{\rhmcsBaTwinSrNRej}{1}  % StepM rejections at FWER 5 % (of 6)
\providecommand{\rhmcsBaTwinSrNRejWord}{}\renewcommand{\rhmcsBaTwinSrNRejWord}{one}  % ... as a word
\providecommand{\rhmcsBaTwinSrNRaw}{}\renewcommand{\rhmcsBaTwinSrNRaw}{3}  % single-test bootstrap p < 0.05 (no adjustment)
\providecommand{\rhmcsBaTwinSrNRawWord}{}\renewcommand{\rhmcsBaTwinSrNRawWord}{three}  % ... as a word
\providecommand{\rhmcsBaTwinSrNPos}{}\renewcommand{\rhmcsBaTwinSrNPos}{6}  % alternatives with a positive estimate (better than the benchmark)
\providecommand{\rhmcsBaTwinSrPMin}{}\renewcommand{\rhmcsBaTwinSrPMin}{0.008}  % smallest StepM-adjusted p-value
\providecommand{\rhmcsBaTwinSrPMax}{}\renewcommand{\rhmcsBaTwinSrPMax}{0.14}  % largest StepM-adjusted p-value
\providecommand{\rhmcsBaTwinSrPTwoMin}{}\renewcommand{\rhmcsBaTwinSrPTwoMin}{0.013}  % smallest two-sided StepM-adjusted p-value
\providecommand{\rhmcsBaTwinSrPTwoMax}{}\renewcommand{\rhmcsBaTwinSrPTwoMax}{0.27}  % largest two-sided StepM-adjusted p-value
\providecommand{\rhmcsBaTwinSrRawMin}{}\renewcommand{\rhmcsBaTwinSrRawMin}{0.002}  % smallest single-test bootstrap p-value
\providecommand{\rhmcsBaTwinSrRej}{}\renewcommand{\rhmcsBaTwinSrRej}{the random forest on the full predictor set}  % the alternatives StepM rejects at 5 % (none: none)
\providecommand{\rhmcsBaTwinSrNRejTwo}{}\renewcommand{\rhmcsBaTwinSrNRejTwo}{1}  % StepM rejections at FWER 5 %, two-sided
\providecommand{\rhmcsBaTwinSrNRejTwoWord}{}\renewcommand{\rhmcsBaTwinSrNRejTwoWord}{one}  % ... as a word
\providecommand{\rhmcsBaTwinSrNRawTwo}{}\renewcommand{\rhmcsBaTwinSrNRawTwo}{2}  % single-test bootstrap p < 0.05, two-sided
\providecommand{\rhmcsBaTwinSrRejTwo}{}\renewcommand{\rhmcsBaTwinSrRejTwo}{the random forest on the full predictor set}  % the alternatives the two-sided StepM rejects at 5 % (none: none)
\providecommand{\rhmcsBaTwinSrPRidge}{}\renewcommand{\rhmcsBaTwinSrPRidge}{0.14}  % StepM-adjusted p-value (one-sided), Ridge
\providecommand{\rhmcsBaTwinSrRawRidge}{}\renewcommand{\rhmcsBaTwinSrRawRidge}{0.089}  % single-test bootstrap p-value (one-sided), Ridge
\providecommand{\rhmcsBaTwinSrPTwoRidge}{}\renewcommand{\rhmcsBaTwinSrPTwoRidge}{0.27}  % StepM-adjusted p-value (two-sided), Ridge
\providecommand{\rhmcsBaTwinSrRawTwoRidge}{}\renewcommand{\rhmcsBaTwinSrRawTwoRidge}{0.17}  % single-test bootstrap p-value (two-sided), Ridge
\providecommand{\rhmcsBaTwinSrPLasso}{}\renewcommand{\rhmcsBaTwinSrPLasso}{0.12}  % StepM-adjusted p-value (one-sided), Lasso
\providecommand{\rhmcsBaTwinSrRawLasso}{}\renewcommand{\rhmcsBaTwinSrRawLasso}{0.071}  % single-test bootstrap p-value (one-sided), Lasso
\providecommand{\rhmcsBaTwinSrPTwoLasso}{}\renewcommand{\rhmcsBaTwinSrPTwoLasso}{0.22}  % StepM-adjusted p-value (two-sided), Lasso
\providecommand{\rhmcsBaTwinSrRawTwoLasso}{}\renewcommand{\rhmcsBaTwinSrRawTwoLasso}{0.12}  % single-test bootstrap p-value (two-sided), Lasso
\providecommand{\rhmcsBaTwinSrPEnet}{}\renewcommand{\rhmcsBaTwinSrPEnet}{0.080}  % StepM-adjusted p-value (one-sided), Elastic net
\providecommand{\rhmcsBaTwinSrRawEnet}{}\renewcommand{\rhmcsBaTwinSrRawEnet}{0.043}  % single-test bootstrap p-value (one-sided), Elastic net
\providecommand{\rhmcsBaTwinSrPTwoEnet}{}\renewcommand{\rhmcsBaTwinSrPTwoEnet}{0.14}  % StepM-adjusted p-value (two-sided), Elastic net
\providecommand{\rhmcsBaTwinSrRawTwoEnet}{}\renewcommand{\rhmcsBaTwinSrRawTwoEnet}{0.070}  % single-test bootstrap p-value (two-sided), Elastic net
\providecommand{\rhmcsBaTwinSrPLgbm}{}\renewcommand{\rhmcsBaTwinSrPLgbm}{0.066}  % StepM-adjusted p-value (one-sided), LightGBM
\providecommand{\rhmcsBaTwinSrRawLgbm}{}\renewcommand{\rhmcsBaTwinSrRawLgbm}{0.027}  % single-test bootstrap p-value (one-sided), LightGBM
\providecommand{\rhmcsBaTwinSrPTwoLgbm}{}\renewcommand{\rhmcsBaTwinSrPTwoLgbm}{0.12}  % StepM-adjusted p-value (two-sided), LightGBM
\providecommand{\rhmcsBaTwinSrRawTwoLgbm}{}\renewcommand{\rhmcsBaTwinSrRawTwoLgbm}{0.050}  % single-test bootstrap p-value (two-sided), LightGBM
\providecommand{\rhmcsBaTwinSrPXgb}{}\renewcommand{\rhmcsBaTwinSrPXgb}{0.14}  % StepM-adjusted p-value (one-sided), XGBoost
\providecommand{\rhmcsBaTwinSrRawXgb}{}\renewcommand{\rhmcsBaTwinSrRawXgb}{0.11}  % single-test bootstrap p-value (one-sided), XGBoost
\providecommand{\rhmcsBaTwinSrPTwoXgb}{}\renewcommand{\rhmcsBaTwinSrPTwoXgb}{0.27}  % StepM-adjusted p-value (two-sided), XGBoost
\providecommand{\rhmcsBaTwinSrRawTwoXgb}{}\renewcommand{\rhmcsBaTwinSrRawTwoXgb}{0.21}  % single-test bootstrap p-value (two-sided), XGBoost
\providecommand{\rhmcsBaTwinSrPRf}{}\renewcommand{\rhmcsBaTwinSrPRf}{0.008}  % StepM-adjusted p-value (one-sided), Random forest
\providecommand{\rhmcsBaTwinSrRawRf}{}\renewcommand{\rhmcsBaTwinSrRawRf}{0.002}  % single-test bootstrap p-value (one-sided), Random forest
\providecommand{\rhmcsBaTwinSrPTwoRf}{}\renewcommand{\rhmcsBaTwinSrPTwoRf}{0.013}  % StepM-adjusted p-value (two-sided), Random forest
\providecommand{\rhmcsBaTwinSrRawTwoRf}{}\renewcommand{\rhmcsBaTwinSrRawTwoRf}{0.003}  % single-test bootstrap p-value (two-sided), Random forest
\providecommand{\rhmcsBaTwinSrNRejLo}{}\renewcommand{\rhmcsBaTwinSrNRejLo}{1}  % StepM rejections at 5 %, block length 5
\providecommand{\rhmcsBaTwinSrSpaCLo}{}\renewcommand{\rhmcsBaTwinSrSpaCLo}{0.011}  % SPA p-value (consistent), block length 5
\providecommand{\rhmcsBaTwinSrNRejTwoLo}{}\renewcommand{\rhmcsBaTwinSrNRejTwoLo}{1}  % ... two-sided, block length 5
\providecommand{\rhmcsBaTwinSrPMinLo}{}\renewcommand{\rhmcsBaTwinSrPMinLo}{0.011}  % smallest adjusted p, block length 5
\providecommand{\rhmcsBaTwinSrPTwoMinLo}{}\renewcommand{\rhmcsBaTwinSrPTwoMinLo}{0.017}  % ... two-sided, block length 5
\providecommand{\rhmcsBaTwinSrNRejHi}{}\renewcommand{\rhmcsBaTwinSrNRejHi}{1}  % StepM rejections at 5 %, block length 20
\providecommand{\rhmcsBaTwinSrSpaCHi}{}\renewcommand{\rhmcsBaTwinSrSpaCHi}{0.012}  % SPA p-value (consistent), block length 20
\providecommand{\rhmcsBaTwinSrNRejTwoHi}{}\renewcommand{\rhmcsBaTwinSrNRejTwoHi}{1}  % ... two-sided, block length 20
\providecommand{\rhmcsBaTwinSrPMinHi}{}\renewcommand{\rhmcsBaTwinSrPMinHi}{0.012}  % smallest adjusted p, block length 20
\providecommand{\rhmcsBaTwinSrPTwoMinHi}{}\renewcommand{\rhmcsBaTwinSrPTwoMinHi}{0.018}  % ... two-sided, block length 20
\providecommand{\rhmcsTwinTwoSidedPRidge}{}\renewcommand{\rhmcsTwinTwoSidedPRidge}{0.17}  % two-sided pooled-SR bootstrap p of the Full - HAR Sharpe difference, Ridge (D19: the check p_bt; the paper's test is Ledoit-Wolf)
\providecommand{\rhmcsTwinTwoSidedPLasso}{}\renewcommand{\rhmcsTwinTwoSidedPLasso}{0.12}  % two-sided pooled-SR bootstrap p of the Full - HAR Sharpe difference, Lasso (D19: the check p_bt; the paper's test is Ledoit-Wolf)
\providecommand{\rhmcsTwinTwoSidedPEnet}{}\renewcommand{\rhmcsTwinTwoSidedPEnet}{0.072}  % two-sided pooled-SR bootstrap p of the Full - HAR Sharpe difference, Elastic net (D19: the check p_bt; the paper's test is Ledoit-Wolf)
\providecommand{\rhmcsTwinTwoSidedPLgbm}{}\renewcommand{\rhmcsTwinTwoSidedPLgbm}{0.050}  % two-sided pooled-SR bootstrap p of the Full - HAR Sharpe difference, LightGBM (D19: the check p_bt; the paper's test is Ledoit-Wolf)
\providecommand{\rhmcsTwinTwoSidedPXgb}{}\renewcommand{\rhmcsTwinTwoSidedPXgb}{0.20}  % two-sided pooled-SR bootstrap p of the Full - HAR Sharpe difference, XGBoost (D19: the check p_bt; the paper's test is Ledoit-Wolf)
\providecommand{\rhmcsTwinTwoSidedPRf}{}\renewcommand{\rhmcsTwinTwoSidedPRf}{0.003}  % two-sided pooled-SR bootstrap p of the Full - HAR Sharpe difference, Random forest (D19: the check p_bt; the paper's test is Ledoit-Wolf)
\fi
\ifdefined\rhemShCiHar\else% AUTO-GENERATED by writeup/make_rehaul_econ_main.py -- do not edit; re-run the script.
% \rhem... macros for Section 7's main economics table (tab:rh_econ_main), the by-year figure
% (fig:rh_econ_byyear) and the break-even spreads. Sample: the 614 post-warmup days of T_opt; the sign strategy.
% Row stems: <Est><Set> (Est = Ridge, Lasso, Enet, Lgbm, Xgb, Rf; Set = Full, Har: the pooled intraday
% forecasts of the final rerun), CiHar (the closing-interval OLS on the HAR benchmark set), Short (q = -1).
% ---- per row: Sharpe ratio (D19: pooled, sqrt(q)), mean daily return (% of premium), at mid and at quoted bid/ask
\providecommand{\rhemShRidgeFull}{}\renewcommand{\rhemShRidgeFull}{1.64}  % poolpt13_ridge_all_features: SR, mid
\providecommand{\rhemShXRidgeFull}{}\renewcommand{\rhemShXRidgeFull}{1.29}  % ... at quoted bid/ask
\providecommand{\rhemMeanRidgeFull}{}\renewcommand{\rhemMeanRidgeFull}{12.3}  % ... mean daily return, %, mid
\providecommand{\rhemMeanXRidgeFull}{}\renewcommand{\rhemMeanXRidgeFull}{9.7}  % ... mean daily return, %, bid/ask
\providecommand{\rhemShDOneRidgeFull}{}\renewcommand{\rhemShDOneRidgeFull}{1.44}  % ... SR mid, per-year sqrt(N_y), N_y-weighted (tab:rh_econ_sharpe_defs)
\providecommand{\rhemShDTwoRidgeFull}{}\renewcommand{\rhemShDTwoRidgeFull}{1.74}  % ... SR mid, pooled daily returns, sqrt(252)
\providecommand{\rhemShDSixRidgeFull}{}\renewcommand{\rhemShDSixRidgeFull}{1.75}  % ... SR mid, D6: equal-year mean of SR_y (robustness, D19)
\providecommand{\rhemShLassoFull}{}\renewcommand{\rhemShLassoFull}{1.73}  % poolpt13_lasso_all_features: SR, mid
\providecommand{\rhemShXLassoFull}{}\renewcommand{\rhemShXLassoFull}{1.37}  % ... at quoted bid/ask
\providecommand{\rhemMeanLassoFull}{}\renewcommand{\rhemMeanLassoFull}{13.0}  % ... mean daily return, %, mid
\providecommand{\rhemMeanXLassoFull}{}\renewcommand{\rhemMeanXLassoFull}{10.3}  % ... mean daily return, %, bid/ask
\providecommand{\rhemShDOneLassoFull}{}\renewcommand{\rhemShDOneLassoFull}{1.62}  % ... SR mid, per-year sqrt(N_y), N_y-weighted (tab:rh_econ_sharpe_defs)
\providecommand{\rhemShDTwoLassoFull}{}\renewcommand{\rhemShDTwoLassoFull}{1.83}  % ... SR mid, pooled daily returns, sqrt(252)
\providecommand{\rhemShDSixLassoFull}{}\renewcommand{\rhemShDSixLassoFull}{1.75}  % ... SR mid, D6: equal-year mean of SR_y (robustness, D19)
\providecommand{\rhemShEnetFull}{}\renewcommand{\rhemShEnetFull}{1.85}  % poolpt13_enet_all_features: SR, mid
\providecommand{\rhemShXEnetFull}{}\renewcommand{\rhemShXEnetFull}{1.50}  % ... at quoted bid/ask
\providecommand{\rhemMeanEnetFull}{}\renewcommand{\rhemMeanEnetFull}{13.9}  % ... mean daily return, %, mid
\providecommand{\rhemMeanXEnetFull}{}\renewcommand{\rhemMeanXEnetFull}{11.3}  % ... mean daily return, %, bid/ask
\providecommand{\rhemShDOneEnetFull}{}\renewcommand{\rhemShDOneEnetFull}{1.67}  % ... SR mid, per-year sqrt(N_y), N_y-weighted (tab:rh_econ_sharpe_defs)
\providecommand{\rhemShDTwoEnetFull}{}\renewcommand{\rhemShDTwoEnetFull}{1.96}  % ... SR mid, pooled daily returns, sqrt(252)
\providecommand{\rhemShDSixEnetFull}{}\renewcommand{\rhemShDSixEnetFull}{1.95}  % ... SR mid, D6: equal-year mean of SR_y (robustness, D19)
\providecommand{\rhemShLgbmFull}{}\renewcommand{\rhemShLgbmFull}{1.54}  % pooltree_lgbm_all_features: SR, mid
\providecommand{\rhemShXLgbmFull}{}\renewcommand{\rhemShXLgbmFull}{1.18}  % ... at quoted bid/ask
\providecommand{\rhemMeanLgbmFull}{}\renewcommand{\rhemMeanLgbmFull}{11.5}  % ... mean daily return, %, mid
\providecommand{\rhemMeanXLgbmFull}{}\renewcommand{\rhemMeanXLgbmFull}{8.9}  % ... mean daily return, %, bid/ask
\providecommand{\rhemShDOneLgbmFull}{}\renewcommand{\rhemShDOneLgbmFull}{1.32}  % ... SR mid, per-year sqrt(N_y), N_y-weighted (tab:rh_econ_sharpe_defs)
\providecommand{\rhemShDTwoLgbmFull}{}\renewcommand{\rhemShDTwoLgbmFull}{1.62}  % ... SR mid, pooled daily returns, sqrt(252)
\providecommand{\rhemShDSixLgbmFull}{}\renewcommand{\rhemShDSixLgbmFull}{1.67}  % ... SR mid, D6: equal-year mean of SR_y (robustness, D19)
\providecommand{\rhemShXgbFull}{}\renewcommand{\rhemShXgbFull}{1.54}  % pooltree_xgb_all_features: SR, mid
\providecommand{\rhemShXXgbFull}{}\renewcommand{\rhemShXXgbFull}{1.19}  % ... at quoted bid/ask
\providecommand{\rhemMeanXgbFull}{}\renewcommand{\rhemMeanXgbFull}{11.6}  % ... mean daily return, %, mid
\providecommand{\rhemMeanXXgbFull}{}\renewcommand{\rhemMeanXXgbFull}{9.0}  % ... mean daily return, %, bid/ask
\providecommand{\rhemShDOneXgbFull}{}\renewcommand{\rhemShDOneXgbFull}{1.29}  % ... SR mid, per-year sqrt(N_y), N_y-weighted (tab:rh_econ_sharpe_defs)
\providecommand{\rhemShDTwoXgbFull}{}\renewcommand{\rhemShDTwoXgbFull}{1.63}  % ... SR mid, pooled daily returns, sqrt(252)
\providecommand{\rhemShDSixXgbFull}{}\renewcommand{\rhemShDSixXgbFull}{1.71}  % ... SR mid, D6: equal-year mean of SR_y (robustness, D19)
\providecommand{\rhemShRfFull}{}\renewcommand{\rhemShRfFull}{1.88}  % pooltree_rf_all_features: SR, mid
\providecommand{\rhemShXRfFull}{}\renewcommand{\rhemShXRfFull}{1.52}  % ... at quoted bid/ask
\providecommand{\rhemMeanRfFull}{}\renewcommand{\rhemMeanRfFull}{14.1}  % ... mean daily return, %, mid
\providecommand{\rhemMeanXRfFull}{}\renewcommand{\rhemMeanXRfFull}{11.4}  % ... mean daily return, %, bid/ask
\providecommand{\rhemShDOneRfFull}{}\renewcommand{\rhemShDOneRfFull}{1.63}  % ... SR mid, per-year sqrt(N_y), N_y-weighted (tab:rh_econ_sharpe_defs)
\providecommand{\rhemShDTwoRfFull}{}\renewcommand{\rhemShDTwoRfFull}{1.99}  % ... SR mid, pooled daily returns, sqrt(252)
\providecommand{\rhemShDSixRfFull}{}\renewcommand{\rhemShDSixRfFull}{1.99}  % ... SR mid, D6: equal-year mean of SR_y (robustness, D19)
\providecommand{\rhemShRidgeHar}{}\renewcommand{\rhemShRidgeHar}{0.95}  % poolpt13_ridge_baseline: SR, mid
\providecommand{\rhemShXRidgeHar}{}\renewcommand{\rhemShXRidgeHar}{0.60}  % ... at quoted bid/ask
\providecommand{\rhemMeanRidgeHar}{}\renewcommand{\rhemMeanRidgeHar}{7.2}  % ... mean daily return, %, mid
\providecommand{\rhemMeanXRidgeHar}{}\renewcommand{\rhemMeanXRidgeHar}{4.6}  % ... mean daily return, %, bid/ask
\providecommand{\rhemShDOneRidgeHar}{}\renewcommand{\rhemShDOneRidgeHar}{0.99}  % ... SR mid, per-year sqrt(N_y), N_y-weighted (tab:rh_econ_sharpe_defs)
\providecommand{\rhemShDTwoRidgeHar}{}\renewcommand{\rhemShDTwoRidgeHar}{1.01}  % ... SR mid, pooled daily returns, sqrt(252)
\providecommand{\rhemShDSixRidgeHar}{}\renewcommand{\rhemShDSixRidgeHar}{0.98}  % ... SR mid, D6: equal-year mean of SR_y (robustness, D19)
\providecommand{\rhemShLassoHar}{}\renewcommand{\rhemShLassoHar}{0.95}  % poolpt13_lasso_baseline: SR, mid
\providecommand{\rhemShXLassoHar}{}\renewcommand{\rhemShXLassoHar}{0.60}  % ... at quoted bid/ask
\providecommand{\rhemMeanLassoHar}{}\renewcommand{\rhemMeanLassoHar}{7.2}  % ... mean daily return, %, mid
\providecommand{\rhemMeanXLassoHar}{}\renewcommand{\rhemMeanXLassoHar}{4.6}  % ... mean daily return, %, bid/ask
\providecommand{\rhemShDOneLassoHar}{}\renewcommand{\rhemShDOneLassoHar}{0.99}  % ... SR mid, per-year sqrt(N_y), N_y-weighted (tab:rh_econ_sharpe_defs)
\providecommand{\rhemShDTwoLassoHar}{}\renewcommand{\rhemShDTwoLassoHar}{1.01}  % ... SR mid, pooled daily returns, sqrt(252)
\providecommand{\rhemShDSixLassoHar}{}\renewcommand{\rhemShDSixLassoHar}{0.98}  % ... SR mid, D6: equal-year mean of SR_y (robustness, D19)
\providecommand{\rhemShEnetHar}{}\renewcommand{\rhemShEnetHar}{0.95}  % poolpt13_enet_baseline: SR, mid
\providecommand{\rhemShXEnetHar}{}\renewcommand{\rhemShXEnetHar}{0.60}  % ... at quoted bid/ask
\providecommand{\rhemMeanEnetHar}{}\renewcommand{\rhemMeanEnetHar}{7.2}  % ... mean daily return, %, mid
\providecommand{\rhemMeanXEnetHar}{}\renewcommand{\rhemMeanXEnetHar}{4.6}  % ... mean daily return, %, bid/ask
\providecommand{\rhemShDOneEnetHar}{}\renewcommand{\rhemShDOneEnetHar}{0.99}  % ... SR mid, per-year sqrt(N_y), N_y-weighted (tab:rh_econ_sharpe_defs)
\providecommand{\rhemShDTwoEnetHar}{}\renewcommand{\rhemShDTwoEnetHar}{1.01}  % ... SR mid, pooled daily returns, sqrt(252)
\providecommand{\rhemShDSixEnetHar}{}\renewcommand{\rhemShDSixEnetHar}{0.98}  % ... SR mid, D6: equal-year mean of SR_y (robustness, D19)
\providecommand{\rhemShLgbmHar}{}\renewcommand{\rhemShLgbmHar}{0.65}  % pooltree_lgbm_baseline: SR, mid
\providecommand{\rhemShXLgbmHar}{}\renewcommand{\rhemShXLgbmHar}{0.31}  % ... at quoted bid/ask
\providecommand{\rhemMeanLgbmHar}{}\renewcommand{\rhemMeanLgbmHar}{4.9}  % ... mean daily return, %, mid
\providecommand{\rhemMeanXLgbmHar}{}\renewcommand{\rhemMeanXLgbmHar}{2.3}  % ... mean daily return, %, bid/ask
\providecommand{\rhemShDOneLgbmHar}{}\renewcommand{\rhemShDOneLgbmHar}{0.70}  % ... SR mid, per-year sqrt(N_y), N_y-weighted (tab:rh_econ_sharpe_defs)
\providecommand{\rhemShDTwoLgbmHar}{}\renewcommand{\rhemShDTwoLgbmHar}{0.69}  % ... SR mid, pooled daily returns, sqrt(252)
\providecommand{\rhemShDSixLgbmHar}{}\renewcommand{\rhemShDSixLgbmHar}{0.61}  % ... SR mid, D6: equal-year mean of SR_y (robustness, D19)
\providecommand{\rhemShXgbHar}{}\renewcommand{\rhemShXgbHar}{0.98}  % pooltree_xgb_baseline: SR, mid
\providecommand{\rhemShXXgbHar}{}\renewcommand{\rhemShXXgbHar}{0.63}  % ... at quoted bid/ask
\providecommand{\rhemMeanXgbHar}{}\renewcommand{\rhemMeanXgbHar}{7.4}  % ... mean daily return, %, mid
\providecommand{\rhemMeanXXgbHar}{}\renewcommand{\rhemMeanXXgbHar}{4.8}  % ... mean daily return, %, bid/ask
\providecommand{\rhemShDOneXgbHar}{}\renewcommand{\rhemShDOneXgbHar}{0.99}  % ... SR mid, per-year sqrt(N_y), N_y-weighted (tab:rh_econ_sharpe_defs)
\providecommand{\rhemShDTwoXgbHar}{}\renewcommand{\rhemShDTwoXgbHar}{1.04}  % ... SR mid, pooled daily returns, sqrt(252)
\providecommand{\rhemShDSixXgbHar}{}\renewcommand{\rhemShDSixXgbHar}{0.91}  % ... SR mid, D6: equal-year mean of SR_y (robustness, D19)
\providecommand{\rhemShRfHar}{}\renewcommand{\rhemShRfHar}{0.32}  % pooltree_rf_baseline: SR, mid
\providecommand{\rhemShXRfHar}{}\renewcommand{\rhemShXRfHar}{\ensuremath{-}0.03}  % ... at quoted bid/ask
\providecommand{\rhemMeanRfHar}{}\renewcommand{\rhemMeanRfHar}{2.4}  % ... mean daily return, %, mid
\providecommand{\rhemMeanXRfHar}{}\renewcommand{\rhemMeanXRfHar}{\ensuremath{-}0.2}  % ... mean daily return, %, bid/ask
\providecommand{\rhemShDOneRfHar}{}\renewcommand{\rhemShDOneRfHar}{0.66}  % ... SR mid, per-year sqrt(N_y), N_y-weighted (tab:rh_econ_sharpe_defs)
\providecommand{\rhemShDTwoRfHar}{}\renewcommand{\rhemShDTwoRfHar}{0.34}  % ... SR mid, pooled daily returns, sqrt(252)
\providecommand{\rhemShDSixRfHar}{}\renewcommand{\rhemShDSixRfHar}{0.01}  % ... SR mid, D6: equal-year mean of SR_y (robustness, D19)
\providecommand{\rhemShCiHar}{}\renewcommand{\rhemShCiHar}{1.11}  % sub_ols_baseline: SR, mid
\providecommand{\rhemShXCiHar}{}\renewcommand{\rhemShXCiHar}{0.76}  % ... at quoted bid/ask
\providecommand{\rhemMeanCiHar}{}\renewcommand{\rhemMeanCiHar}{8.3}  % ... mean daily return, %, mid
\providecommand{\rhemMeanXCiHar}{}\renewcommand{\rhemMeanXCiHar}{5.7}  % ... mean daily return, %, bid/ask
\providecommand{\rhemShDOneCiHar}{}\renewcommand{\rhemShDOneCiHar}{1.20}  % ... SR mid, per-year sqrt(N_y), N_y-weighted (tab:rh_econ_sharpe_defs)
\providecommand{\rhemShDTwoCiHar}{}\renewcommand{\rhemShDTwoCiHar}{1.17}  % ... SR mid, pooled daily returns, sqrt(252)
\providecommand{\rhemShDSixCiHar}{}\renewcommand{\rhemShDSixCiHar}{0.99}  % ... SR mid, D6: equal-year mean of SR_y (robustness, D19)
\providecommand{\rhemShShort}{}\renewcommand{\rhemShShort}{0.32}  % always_short: SR, mid
\providecommand{\rhemShXShort}{}\renewcommand{\rhemShXShort}{\ensuremath{-}0.03}  % ... at quoted bid/ask
\providecommand{\rhemMeanShort}{}\renewcommand{\rhemMeanShort}{2.4}  % ... mean daily return, %, mid
\providecommand{\rhemMeanXShort}{}\renewcommand{\rhemMeanXShort}{\ensuremath{-}0.2}  % ... mean daily return, %, bid/ask
\providecommand{\rhemShDOneShort}{}\renewcommand{\rhemShDOneShort}{0.77}  % ... SR mid, per-year sqrt(N_y), N_y-weighted (tab:rh_econ_sharpe_defs)
\providecommand{\rhemShDTwoShort}{}\renewcommand{\rhemShDTwoShort}{0.34}  % ... SR mid, pooled daily returns, sqrt(252)
\providecommand{\rhemShDSixShort}{}\renewcommand{\rhemShDSixShort}{\ensuremath{-}0.12}  % ... SR mid, D6: equal-year mean of SR_y (robustness, D19)
% ---- against the short straddle: Delta SR with its Ledoit-Wolf p; fixed-b t (and p) of the mean differential
\providecommand{\rhemDShRidgeFullVsShort}{}\renewcommand{\rhemDShRidgeFullVsShort}{1.32}  % poolpt13_ridge_all_features minus the short straddle: Delta SR, mid
\providecommand{\rhemPDShRidgeFullVsShort}{}\renewcommand{\rhemPDShRidgeFullVsShort}{0.059}  % ... its Ledoit-Wolf p (D19)
\providecommand{\rhemDShXRidgeFullVsShort}{}\renewcommand{\rhemDShXRidgeFullVsShort}{1.32}  % ... Delta SR at bid/ask
\providecommand{\rhemPDShXRidgeFullVsShort}{}\renewcommand{\rhemPDShXRidgeFullVsShort}{0.059}  % ... its Ledoit-Wolf p (D19)
\providecommand{\rhemTRidgeFullVsShort}{}\renewcommand{\rhemTRidgeFullVsShort}{1.70}  % ... fixed-b t of the mean mid return differential (S = 33)
\providecommand{\rhemPTRidgeFullVsShort}{}\renewcommand{\rhemPTRidgeFullVsShort}{0.110}  % ... its fixed-b p
\providecommand{\rhemTXRidgeFullVsShort}{}\renewcommand{\rhemTXRidgeFullVsShort}{1.70}  % ... fixed-b t of the mean bid/ask return differential
\providecommand{\rhemPTXRidgeFullVsShort}{}\renewcommand{\rhemPTXRidgeFullVsShort}{0.109}  % ... its fixed-b p
\providecommand{\rhemDShLassoFullVsShort}{}\renewcommand{\rhemDShLassoFullVsShort}{1.40}  % poolpt13_lasso_all_features minus the short straddle: Delta SR, mid
\providecommand{\rhemPDShLassoFullVsShort}{}\renewcommand{\rhemPDShLassoFullVsShort}{0.039}  % ... its Ledoit-Wolf p (D19)
\providecommand{\rhemDShXLassoFullVsShort}{}\renewcommand{\rhemDShXLassoFullVsShort}{1.40}  % ... Delta SR at bid/ask
\providecommand{\rhemPDShXLassoFullVsShort}{}\renewcommand{\rhemPDShXLassoFullVsShort}{0.039}  % ... its Ledoit-Wolf p (D19)
\providecommand{\rhemTLassoFullVsShort}{}\renewcommand{\rhemTLassoFullVsShort}{1.93}  % ... fixed-b t of the mean mid return differential (S = 33)
\providecommand{\rhemPTLassoFullVsShort}{}\renewcommand{\rhemPTLassoFullVsShort}{0.071}  % ... its fixed-b p
\providecommand{\rhemTXLassoFullVsShort}{}\renewcommand{\rhemTXLassoFullVsShort}{1.93}  % ... fixed-b t of the mean bid/ask return differential
\providecommand{\rhemPTXLassoFullVsShort}{}\renewcommand{\rhemPTXLassoFullVsShort}{0.070}  % ... its fixed-b p
\providecommand{\rhemDShEnetFullVsShort}{}\renewcommand{\rhemDShEnetFullVsShort}{1.53}  % poolpt13_enet_all_features minus the short straddle: Delta SR, mid
\providecommand{\rhemPDShEnetFullVsShort}{}\renewcommand{\rhemPDShEnetFullVsShort}{0.026}  % ... its Ledoit-Wolf p (D19)
\providecommand{\rhemDShXEnetFullVsShort}{}\renewcommand{\rhemDShXEnetFullVsShort}{1.53}  % ... Delta SR at bid/ask
\providecommand{\rhemPDShXEnetFullVsShort}{}\renewcommand{\rhemPDShXEnetFullVsShort}{0.027}  % ... its Ledoit-Wolf p (D19)
\providecommand{\rhemTEnetFullVsShort}{}\renewcommand{\rhemTEnetFullVsShort}{1.95}  % ... fixed-b t of the mean mid return differential (S = 33)
\providecommand{\rhemPTEnetFullVsShort}{}\renewcommand{\rhemPTEnetFullVsShort}{0.068}  % ... its fixed-b p
\providecommand{\rhemTXEnetFullVsShort}{}\renewcommand{\rhemTXEnetFullVsShort}{1.95}  % ... fixed-b t of the mean bid/ask return differential
\providecommand{\rhemPTXEnetFullVsShort}{}\renewcommand{\rhemPTXEnetFullVsShort}{0.068}  % ... its fixed-b p
\providecommand{\rhemDShLgbmFullVsShort}{}\renewcommand{\rhemDShLgbmFullVsShort}{1.21}  % pooltree_lgbm_all_features minus the short straddle: Delta SR, mid
\providecommand{\rhemPDShLgbmFullVsShort}{}\renewcommand{\rhemPDShLgbmFullVsShort}{0.070}  % ... its Ledoit-Wolf p (D19)
\providecommand{\rhemDShXLgbmFullVsShort}{}\renewcommand{\rhemDShXLgbmFullVsShort}{1.21}  % ... Delta SR at bid/ask
\providecommand{\rhemPDShXLgbmFullVsShort}{}\renewcommand{\rhemPDShXLgbmFullVsShort}{0.069}  % ... its Ledoit-Wolf p (D19)
\providecommand{\rhemTLgbmFullVsShort}{}\renewcommand{\rhemTLgbmFullVsShort}{1.68}  % ... fixed-b t of the mean mid return differential (S = 33)
\providecommand{\rhemPTLgbmFullVsShort}{}\renewcommand{\rhemPTLgbmFullVsShort}{0.113}  % ... its fixed-b p
\providecommand{\rhemTXLgbmFullVsShort}{}\renewcommand{\rhemTXLgbmFullVsShort}{1.69}  % ... fixed-b t of the mean bid/ask return differential
\providecommand{\rhemPTXLgbmFullVsShort}{}\renewcommand{\rhemPTXLgbmFullVsShort}{0.112}  % ... its fixed-b p
\providecommand{\rhemDShXgbFullVsShort}{}\renewcommand{\rhemDShXgbFullVsShort}{1.22}  % pooltree_xgb_all_features minus the short straddle: Delta SR, mid
\providecommand{\rhemPDShXgbFullVsShort}{}\renewcommand{\rhemPDShXgbFullVsShort}{0.064}  % ... its Ledoit-Wolf p (D19)
\providecommand{\rhemDShXXgbFullVsShort}{}\renewcommand{\rhemDShXXgbFullVsShort}{1.22}  % ... Delta SR at bid/ask
\providecommand{\rhemPDShXXgbFullVsShort}{}\renewcommand{\rhemPDShXXgbFullVsShort}{0.063}  % ... its Ledoit-Wolf p (D19)
\providecommand{\rhemTXgbFullVsShort}{}\renewcommand{\rhemTXgbFullVsShort}{1.61}  % ... fixed-b t of the mean mid return differential (S = 33)
\providecommand{\rhemPTXgbFullVsShort}{}\renewcommand{\rhemPTXgbFullVsShort}{0.127}  % ... its fixed-b p
\providecommand{\rhemTXXgbFullVsShort}{}\renewcommand{\rhemTXXgbFullVsShort}{1.62}  % ... fixed-b t of the mean bid/ask return differential
\providecommand{\rhemPTXXgbFullVsShort}{}\renewcommand{\rhemPTXXgbFullVsShort}{0.126}  % ... its fixed-b p
\providecommand{\rhemDShRfFullVsShort}{}\renewcommand{\rhemDShRfFullVsShort}{1.56}  % pooltree_rf_all_features minus the short straddle: Delta SR, mid
\providecommand{\rhemPDShRfFullVsShort}{}\renewcommand{\rhemPDShRfFullVsShort}{0.018}  % ... its Ledoit-Wolf p (D19)
\providecommand{\rhemDShXRfFullVsShort}{}\renewcommand{\rhemDShXRfFullVsShort}{1.55}  % ... Delta SR at bid/ask
\providecommand{\rhemPDShXRfFullVsShort}{}\renewcommand{\rhemPDShXRfFullVsShort}{0.019}  % ... its Ledoit-Wolf p (D19)
\providecommand{\rhemTRfFullVsShort}{}\renewcommand{\rhemTRfFullVsShort}{2.14}  % ... fixed-b t of the mean mid return differential (S = 33)
\providecommand{\rhemPTRfFullVsShort}{}\renewcommand{\rhemPTRfFullVsShort}{0.046}  % ... its fixed-b p
\providecommand{\rhemTXRfFullVsShort}{}\renewcommand{\rhemTXRfFullVsShort}{2.15}  % ... fixed-b t of the mean bid/ask return differential
\providecommand{\rhemPTXRfFullVsShort}{}\renewcommand{\rhemPTXRfFullVsShort}{0.046}  % ... its fixed-b p
\providecommand{\rhemDShRidgeHarVsShort}{}\renewcommand{\rhemDShRidgeHarVsShort}{0.63}  % poolpt13_ridge_baseline minus the short straddle: Delta SR, mid
\providecommand{\rhemPDShRidgeHarVsShort}{}\renewcommand{\rhemPDShRidgeHarVsShort}{0.350}  % ... its Ledoit-Wolf p (D19)
\providecommand{\rhemDShXRidgeHarVsShort}{}\renewcommand{\rhemDShXRidgeHarVsShort}{0.63}  % ... Delta SR at bid/ask
\providecommand{\rhemPDShXRidgeHarVsShort}{}\renewcommand{\rhemPDShXRidgeHarVsShort}{0.345}  % ... its Ledoit-Wolf p (D19)
\providecommand{\rhemTRidgeHarVsShort}{}\renewcommand{\rhemTRidgeHarVsShort}{0.90}  % ... fixed-b t of the mean mid return differential (S = 33)
\providecommand{\rhemPTRidgeHarVsShort}{}\renewcommand{\rhemPTRidgeHarVsShort}{0.388}  % ... its fixed-b p
\providecommand{\rhemTXRidgeHarVsShort}{}\renewcommand{\rhemTXRidgeHarVsShort}{0.91}  % ... fixed-b t of the mean bid/ask return differential
\providecommand{\rhemPTXRidgeHarVsShort}{}\renewcommand{\rhemPTXRidgeHarVsShort}{0.383}  % ... its fixed-b p
\providecommand{\rhemDShLassoHarVsShort}{}\renewcommand{\rhemDShLassoHarVsShort}{0.63}  % poolpt13_lasso_baseline minus the short straddle: Delta SR, mid
\providecommand{\rhemPDShLassoHarVsShort}{}\renewcommand{\rhemPDShLassoHarVsShort}{0.350}  % ... its Ledoit-Wolf p (D19)
\providecommand{\rhemDShXLassoHarVsShort}{}\renewcommand{\rhemDShXLassoHarVsShort}{0.63}  % ... Delta SR at bid/ask
\providecommand{\rhemPDShXLassoHarVsShort}{}\renewcommand{\rhemPDShXLassoHarVsShort}{0.345}  % ... its Ledoit-Wolf p (D19)
\providecommand{\rhemTLassoHarVsShort}{}\renewcommand{\rhemTLassoHarVsShort}{0.90}  % ... fixed-b t of the mean mid return differential (S = 33)
\providecommand{\rhemPTLassoHarVsShort}{}\renewcommand{\rhemPTLassoHarVsShort}{0.388}  % ... its fixed-b p
\providecommand{\rhemTXLassoHarVsShort}{}\renewcommand{\rhemTXLassoHarVsShort}{0.91}  % ... fixed-b t of the mean bid/ask return differential
\providecommand{\rhemPTXLassoHarVsShort}{}\renewcommand{\rhemPTXLassoHarVsShort}{0.383}  % ... its fixed-b p
\providecommand{\rhemDShEnetHarVsShort}{}\renewcommand{\rhemDShEnetHarVsShort}{0.63}  % poolpt13_enet_baseline minus the short straddle: Delta SR, mid
\providecommand{\rhemPDShEnetHarVsShort}{}\renewcommand{\rhemPDShEnetHarVsShort}{0.350}  % ... its Ledoit-Wolf p (D19)
\providecommand{\rhemDShXEnetHarVsShort}{}\renewcommand{\rhemDShXEnetHarVsShort}{0.63}  % ... Delta SR at bid/ask
\providecommand{\rhemPDShXEnetHarVsShort}{}\renewcommand{\rhemPDShXEnetHarVsShort}{0.345}  % ... its Ledoit-Wolf p (D19)
\providecommand{\rhemTEnetHarVsShort}{}\renewcommand{\rhemTEnetHarVsShort}{0.90}  % ... fixed-b t of the mean mid return differential (S = 33)
\providecommand{\rhemPTEnetHarVsShort}{}\renewcommand{\rhemPTEnetHarVsShort}{0.388}  % ... its fixed-b p
\providecommand{\rhemTXEnetHarVsShort}{}\renewcommand{\rhemTXEnetHarVsShort}{0.91}  % ... fixed-b t of the mean bid/ask return differential
\providecommand{\rhemPTXEnetHarVsShort}{}\renewcommand{\rhemPTXEnetHarVsShort}{0.383}  % ... its fixed-b p
\providecommand{\rhemDShLgbmHarVsShort}{}\renewcommand{\rhemDShLgbmHarVsShort}{0.33}  % pooltree_lgbm_baseline minus the short straddle: Delta SR, mid
\providecommand{\rhemPDShLgbmHarVsShort}{}\renewcommand{\rhemPDShLgbmHarVsShort}{0.614}  % ... its Ledoit-Wolf p (D19)
\providecommand{\rhemDShXLgbmHarVsShort}{}\renewcommand{\rhemDShXLgbmHarVsShort}{0.34}  % ... Delta SR at bid/ask
\providecommand{\rhemPDShXLgbmHarVsShort}{}\renewcommand{\rhemPDShXLgbmHarVsShort}{0.603}  % ... its Ledoit-Wolf p (D19)
\providecommand{\rhemTLgbmHarVsShort}{}\renewcommand{\rhemTLgbmHarVsShort}{0.53}  % ... fixed-b t of the mean mid return differential (S = 33)
\providecommand{\rhemPTLgbmHarVsShort}{}\renewcommand{\rhemPTLgbmHarVsShort}{0.610}  % ... its fixed-b p
\providecommand{\rhemTXLgbmHarVsShort}{}\renewcommand{\rhemTXLgbmHarVsShort}{0.54}  % ... fixed-b t of the mean bid/ask return differential
\providecommand{\rhemPTXLgbmHarVsShort}{}\renewcommand{\rhemPTXLgbmHarVsShort}{0.602}  % ... its fixed-b p
\providecommand{\rhemDShXgbHarVsShort}{}\renewcommand{\rhemDShXgbHarVsShort}{0.66}  % pooltree_xgb_baseline minus the short straddle: Delta SR, mid
\providecommand{\rhemPDShXgbHarVsShort}{}\renewcommand{\rhemPDShXgbHarVsShort}{0.315}  % ... its Ledoit-Wolf p (D19)
\providecommand{\rhemDShXXgbHarVsShort}{}\renewcommand{\rhemDShXXgbHarVsShort}{0.67}  % ... Delta SR at bid/ask
\providecommand{\rhemPDShXXgbHarVsShort}{}\renewcommand{\rhemPDShXXgbHarVsShort}{0.305}  % ... its Ledoit-Wolf p (D19)
\providecommand{\rhemTXgbHarVsShort}{}\renewcommand{\rhemTXgbHarVsShort}{1.08}  % ... fixed-b t of the mean mid return differential (S = 33)
\providecommand{\rhemPTXgbHarVsShort}{}\renewcommand{\rhemPTXgbHarVsShort}{0.304}  % ... its fixed-b p
\providecommand{\rhemTXXgbHarVsShort}{}\renewcommand{\rhemTXXgbHarVsShort}{1.09}  % ... fixed-b t of the mean bid/ask return differential
\providecommand{\rhemPTXXgbHarVsShort}{}\renewcommand{\rhemPTXXgbHarVsShort}{0.299}  % ... its fixed-b p
\providecommand{\rhemDShRfHarVsShort}{}\renewcommand{\rhemDShRfHarVsShort}{0.00}  % pooltree_rf_baseline minus the short straddle: Delta SR, mid
\providecommand{\rhemPDShRfHarVsShort}{}\renewcommand{\rhemPDShRfHarVsShort}{0.992}  % ... its Ledoit-Wolf p (D19)
\providecommand{\rhemDShXRfHarVsShort}{}\renewcommand{\rhemDShXRfHarVsShort}{0.00}  % ... Delta SR at bid/ask
\providecommand{\rhemPDShXRfHarVsShort}{}\renewcommand{\rhemPDShXRfHarVsShort}{0.998}  % ... its Ledoit-Wolf p (D19)
\providecommand{\rhemTRfHarVsShort}{}\renewcommand{\rhemTRfHarVsShort}{\ensuremath{-}0.01}  % ... fixed-b t of the mean mid return differential (S = 33)
\providecommand{\rhemPTRfHarVsShort}{}\renewcommand{\rhemPTRfHarVsShort}{0.993}  % ... its fixed-b p
\providecommand{\rhemTXRfHarVsShort}{}\renewcommand{\rhemTXRfHarVsShort}{0.00}  % ... fixed-b t of the mean bid/ask return differential
\providecommand{\rhemPTXRfHarVsShort}{}\renewcommand{\rhemPTXRfHarVsShort}{0.997}  % ... its fixed-b p
\providecommand{\rhemDShCiHarVsShort}{}\renewcommand{\rhemDShCiHarVsShort}{0.79}  % sub_ols_baseline minus the short straddle: Delta SR, mid
\providecommand{\rhemPDShCiHarVsShort}{}\renewcommand{\rhemPDShCiHarVsShort}{0.254}  % ... its Ledoit-Wolf p (D19)
\providecommand{\rhemDShXCiHarVsShort}{}\renewcommand{\rhemDShXCiHarVsShort}{0.79}  % ... Delta SR at bid/ask
\providecommand{\rhemPDShXCiHarVsShort}{}\renewcommand{\rhemPDShXCiHarVsShort}{0.247}  % ... its Ledoit-Wolf p (D19)
\providecommand{\rhemTCiHarVsShort}{}\renewcommand{\rhemTCiHarVsShort}{1.14}  % ... fixed-b t of the mean mid return differential (S = 33)
\providecommand{\rhemPTCiHarVsShort}{}\renewcommand{\rhemPTCiHarVsShort}{0.278}  % ... its fixed-b p
\providecommand{\rhemTXCiHarVsShort}{}\renewcommand{\rhemTXCiHarVsShort}{1.15}  % ... fixed-b t of the mean bid/ask return differential
\providecommand{\rhemPTXCiHarVsShort}{}\renewcommand{\rhemPTXCiHarVsShort}{0.273}  % ... its fixed-b p
% ---- Full minus its own HAR row (same model class); every pooled row minus CI-HAR
\providecommand{\rhemDShRidgeFullVsHar}{}\renewcommand{\rhemDShRidgeFullVsHar}{0.69}  % ridge: Full minus HAR, Delta SR, mid
\providecommand{\rhemPDShRidgeFullVsHar}{}\renewcommand{\rhemPDShRidgeFullVsHar}{0.171}  % ... its Ledoit-Wolf p (D19)
\providecommand{\rhemDShXRidgeFullVsHar}{}\renewcommand{\rhemDShXRidgeFullVsHar}{0.68}  % ... Delta SR at bid/ask
\providecommand{\rhemPDShXRidgeFullVsHar}{}\renewcommand{\rhemPDShXRidgeFullVsHar}{0.171}  % ... its Ledoit-Wolf p (D19)
\providecommand{\rhemDShLassoFullVsHar}{}\renewcommand{\rhemDShLassoFullVsHar}{0.77}  % lasso: Full minus HAR, Delta SR, mid
\providecommand{\rhemPDShLassoFullVsHar}{}\renewcommand{\rhemPDShLassoFullVsHar}{0.120}  % ... its Ledoit-Wolf p (D19)
\providecommand{\rhemDShXLassoFullVsHar}{}\renewcommand{\rhemDShXLassoFullVsHar}{0.77}  % ... Delta SR at bid/ask
\providecommand{\rhemPDShXLassoFullVsHar}{}\renewcommand{\rhemPDShXLassoFullVsHar}{0.120}  % ... its Ledoit-Wolf p (D19)
\providecommand{\rhemDShEnetFullVsHar}{}\renewcommand{\rhemDShEnetFullVsHar}{0.90}  % enet: Full minus HAR, Delta SR, mid
\providecommand{\rhemPDShEnetFullVsHar}{}\renewcommand{\rhemPDShEnetFullVsHar}{0.074}  % ... its Ledoit-Wolf p (D19)
\providecommand{\rhemDShXEnetFullVsHar}{}\renewcommand{\rhemDShXEnetFullVsHar}{0.89}  % ... Delta SR at bid/ask
\providecommand{\rhemPDShXEnetFullVsHar}{}\renewcommand{\rhemPDShXEnetFullVsHar}{0.075}  % ... its Ledoit-Wolf p (D19)
\providecommand{\rhemDShLgbmFullVsHar}{}\renewcommand{\rhemDShLgbmFullVsHar}{0.88}  % lgbm: Full minus HAR, Delta SR, mid
\providecommand{\rhemPDShLgbmFullVsHar}{}\renewcommand{\rhemPDShLgbmFullVsHar}{0.050}  % ... its Ledoit-Wolf p (D19)
\providecommand{\rhemDShXLgbmFullVsHar}{}\renewcommand{\rhemDShXLgbmFullVsHar}{0.88}  % ... Delta SR at bid/ask
\providecommand{\rhemPDShXLgbmFullVsHar}{}\renewcommand{\rhemPDShXLgbmFullVsHar}{0.051}  % ... its Ledoit-Wolf p (D19)
\providecommand{\rhemDShXgbFullVsHar}{}\renewcommand{\rhemDShXgbFullVsHar}{0.56}  % xgb: Full minus HAR, Delta SR, mid
\providecommand{\rhemPDShXgbFullVsHar}{}\renewcommand{\rhemPDShXgbFullVsHar}{0.204}  % ... its Ledoit-Wolf p (D19)
\providecommand{\rhemDShXXgbFullVsHar}{}\renewcommand{\rhemDShXXgbFullVsHar}{0.56}  % ... Delta SR at bid/ask
\providecommand{\rhemPDShXXgbFullVsHar}{}\renewcommand{\rhemPDShXXgbFullVsHar}{0.210}  % ... its Ledoit-Wolf p (D19)
\providecommand{\rhemDShRfFullVsHar}{}\renewcommand{\rhemDShRfFullVsHar}{1.56}  % rf: Full minus HAR, Delta SR, mid
\providecommand{\rhemPDShRfFullVsHar}{}\renewcommand{\rhemPDShRfFullVsHar}{0.020}  % ... its Ledoit-Wolf p (D19)
\providecommand{\rhemDShXRfFullVsHar}{}\renewcommand{\rhemDShXRfFullVsHar}{1.55}  % ... Delta SR at bid/ask
\providecommand{\rhemPDShXRfFullVsHar}{}\renewcommand{\rhemPDShXRfFullVsHar}{0.024}  % ... its Ledoit-Wolf p (D19)
\providecommand{\rhemDShRidgeFullVsCiHar}{}\renewcommand{\rhemDShRidgeFullVsCiHar}{0.53}  % poolpt13_ridge_all_features minus CI-HAR: Delta SR, mid
\providecommand{\rhemPDShRidgeFullVsCiHar}{}\renewcommand{\rhemPDShRidgeFullVsCiHar}{0.330}  % ... its Ledoit-Wolf p (D19)
\providecommand{\rhemDShXRidgeFullVsCiHar}{}\renewcommand{\rhemDShXRidgeFullVsCiHar}{0.53}  % ... Delta SR at bid/ask
\providecommand{\rhemPDShXRidgeFullVsCiHar}{}\renewcommand{\rhemPDShXRidgeFullVsCiHar}{0.339}  % ... its Ledoit-Wolf p (D19)
\providecommand{\rhemTRidgeFullVsCiHar}{}\renewcommand{\rhemTRidgeFullVsCiHar}{0.98}  % ... fixed-b t of the mean mid return differential
\providecommand{\rhemPTRidgeFullVsCiHar}{}\renewcommand{\rhemPTRidgeFullVsCiHar}{0.347}  % ... its fixed-b p
\providecommand{\rhemDShLassoFullVsCiHar}{}\renewcommand{\rhemDShLassoFullVsCiHar}{0.62}  % poolpt13_lasso_all_features minus CI-HAR: Delta SR, mid
\providecommand{\rhemPDShLassoFullVsCiHar}{}\renewcommand{\rhemPDShLassoFullVsCiHar}{0.255}  % ... its Ledoit-Wolf p (D19)
\providecommand{\rhemDShXLassoFullVsCiHar}{}\renewcommand{\rhemDShXLassoFullVsCiHar}{0.61}  % ... Delta SR at bid/ask
\providecommand{\rhemPDShXLassoFullVsCiHar}{}\renewcommand{\rhemPDShXLassoFullVsCiHar}{0.264}  % ... its Ledoit-Wolf p (D19)
\providecommand{\rhemTLassoFullVsCiHar}{}\renewcommand{\rhemTLassoFullVsCiHar}{1.21}  % ... fixed-b t of the mean mid return differential
\providecommand{\rhemPTLassoFullVsCiHar}{}\renewcommand{\rhemPTLassoFullVsCiHar}{0.250}  % ... its fixed-b p
\providecommand{\rhemDShEnetFullVsCiHar}{}\renewcommand{\rhemDShEnetFullVsCiHar}{0.75}  % poolpt13_enet_all_features minus CI-HAR: Delta SR, mid
\providecommand{\rhemPDShEnetFullVsCiHar}{}\renewcommand{\rhemPDShEnetFullVsCiHar}{0.154}  % ... its Ledoit-Wolf p (D19)
\providecommand{\rhemDShXEnetFullVsCiHar}{}\renewcommand{\rhemDShXEnetFullVsCiHar}{0.74}  % ... Delta SR at bid/ask
\providecommand{\rhemPDShXEnetFullVsCiHar}{}\renewcommand{\rhemPDShXEnetFullVsCiHar}{0.160}  % ... its Ledoit-Wolf p (D19)
\providecommand{\rhemTEnetFullVsCiHar}{}\renewcommand{\rhemTEnetFullVsCiHar}{1.42}  % ... fixed-b t of the mean mid return differential
\providecommand{\rhemPTEnetFullVsCiHar}{}\renewcommand{\rhemPTEnetFullVsCiHar}{0.177}  % ... its fixed-b p
\providecommand{\rhemDShLgbmFullVsCiHar}{}\renewcommand{\rhemDShLgbmFullVsCiHar}{0.43}  % pooltree_lgbm_all_features minus CI-HAR: Delta SR, mid
\providecommand{\rhemPDShLgbmFullVsCiHar}{}\renewcommand{\rhemPDShLgbmFullVsCiHar}{0.358}  % ... its Ledoit-Wolf p (D19)
\providecommand{\rhemDShXLgbmFullVsCiHar}{}\renewcommand{\rhemDShXLgbmFullVsCiHar}{0.42}  % ... Delta SR at bid/ask
\providecommand{\rhemPDShXLgbmFullVsCiHar}{}\renewcommand{\rhemPDShXLgbmFullVsCiHar}{0.365}  % ... its Ledoit-Wolf p (D19)
\providecommand{\rhemTLgbmFullVsCiHar}{}\renewcommand{\rhemTLgbmFullVsCiHar}{0.90}  % ... fixed-b t of the mean mid return differential
\providecommand{\rhemPTLgbmFullVsCiHar}{}\renewcommand{\rhemPTLgbmFullVsCiHar}{0.388}  % ... its fixed-b p
\providecommand{\rhemDShXgbFullVsCiHar}{}\renewcommand{\rhemDShXgbFullVsCiHar}{0.44}  % pooltree_xgb_all_features minus CI-HAR: Delta SR, mid
\providecommand{\rhemPDShXgbFullVsCiHar}{}\renewcommand{\rhemPDShXgbFullVsCiHar}{0.332}  % ... its Ledoit-Wolf p (D19)
\providecommand{\rhemDShXXgbFullVsCiHar}{}\renewcommand{\rhemDShXXgbFullVsCiHar}{0.43}  % ... Delta SR at bid/ask
\providecommand{\rhemPDShXXgbFullVsCiHar}{}\renewcommand{\rhemPDShXXgbFullVsCiHar}{0.341}  % ... its Ledoit-Wolf p (D19)
\providecommand{\rhemTXgbFullVsCiHar}{}\renewcommand{\rhemTXgbFullVsCiHar}{0.88}  % ... fixed-b t of the mean mid return differential
\providecommand{\rhemPTXgbFullVsCiHar}{}\renewcommand{\rhemPTXgbFullVsCiHar}{0.398}  % ... its fixed-b p
\providecommand{\rhemDShRfFullVsCiHar}{}\renewcommand{\rhemDShRfFullVsCiHar}{0.77}  % pooltree_rf_all_features minus CI-HAR: Delta SR, mid
\providecommand{\rhemPDShRfFullVsCiHar}{}\renewcommand{\rhemPDShRfFullVsCiHar}{0.093}  % ... its Ledoit-Wolf p (D19)
\providecommand{\rhemDShXRfFullVsCiHar}{}\renewcommand{\rhemDShXRfFullVsCiHar}{0.76}  % ... Delta SR at bid/ask
\providecommand{\rhemPDShXRfFullVsCiHar}{}\renewcommand{\rhemPDShXRfFullVsCiHar}{0.098}  % ... its Ledoit-Wolf p (D19)
\providecommand{\rhemTRfFullVsCiHar}{}\renewcommand{\rhemTRfFullVsCiHar}{1.69}  % ... fixed-b t of the mean mid return differential
\providecommand{\rhemPTRfFullVsCiHar}{}\renewcommand{\rhemPTRfFullVsCiHar}{0.111}  % ... its fixed-b p
\providecommand{\rhemDShRidgeHarVsCiHar}{}\renewcommand{\rhemDShRidgeHarVsCiHar}{\ensuremath{-}0.15}  % poolpt13_ridge_baseline minus CI-HAR: Delta SR, mid
\providecommand{\rhemPDShRidgeHarVsCiHar}{}\renewcommand{\rhemPDShRidgeHarVsCiHar}{0.673}  % ... its Ledoit-Wolf p (D19)
\providecommand{\rhemDShXRidgeHarVsCiHar}{}\renewcommand{\rhemDShXRidgeHarVsCiHar}{\ensuremath{-}0.16}  % ... Delta SR at bid/ask
\providecommand{\rhemPDShXRidgeHarVsCiHar}{}\renewcommand{\rhemPDShXRidgeHarVsCiHar}{0.660}  % ... its Ledoit-Wolf p (D19)
\providecommand{\rhemTRidgeHarVsCiHar}{}\renewcommand{\rhemTRidgeHarVsCiHar}{\ensuremath{-}0.40}  % ... fixed-b t of the mean mid return differential
\providecommand{\rhemPTRidgeHarVsCiHar}{}\renewcommand{\rhemPTRidgeHarVsCiHar}{0.704}  % ... its fixed-b p
\providecommand{\rhemDShLassoHarVsCiHar}{}\renewcommand{\rhemDShLassoHarVsCiHar}{\ensuremath{-}0.15}  % poolpt13_lasso_baseline minus CI-HAR: Delta SR, mid
\providecommand{\rhemPDShLassoHarVsCiHar}{}\renewcommand{\rhemPDShLassoHarVsCiHar}{0.673}  % ... its Ledoit-Wolf p (D19)
\providecommand{\rhemDShXLassoHarVsCiHar}{}\renewcommand{\rhemDShXLassoHarVsCiHar}{\ensuremath{-}0.16}  % ... Delta SR at bid/ask
\providecommand{\rhemPDShXLassoHarVsCiHar}{}\renewcommand{\rhemPDShXLassoHarVsCiHar}{0.660}  % ... its Ledoit-Wolf p (D19)
\providecommand{\rhemTLassoHarVsCiHar}{}\renewcommand{\rhemTLassoHarVsCiHar}{\ensuremath{-}0.40}  % ... fixed-b t of the mean mid return differential
\providecommand{\rhemPTLassoHarVsCiHar}{}\renewcommand{\rhemPTLassoHarVsCiHar}{0.704}  % ... its fixed-b p
\providecommand{\rhemDShEnetHarVsCiHar}{}\renewcommand{\rhemDShEnetHarVsCiHar}{\ensuremath{-}0.15}  % poolpt13_enet_baseline minus CI-HAR: Delta SR, mid
\providecommand{\rhemPDShEnetHarVsCiHar}{}\renewcommand{\rhemPDShEnetHarVsCiHar}{0.673}  % ... its Ledoit-Wolf p (D19)
\providecommand{\rhemDShXEnetHarVsCiHar}{}\renewcommand{\rhemDShXEnetHarVsCiHar}{\ensuremath{-}0.16}  % ... Delta SR at bid/ask
\providecommand{\rhemPDShXEnetHarVsCiHar}{}\renewcommand{\rhemPDShXEnetHarVsCiHar}{0.660}  % ... its Ledoit-Wolf p (D19)
\providecommand{\rhemTEnetHarVsCiHar}{}\renewcommand{\rhemTEnetHarVsCiHar}{\ensuremath{-}0.40}  % ... fixed-b t of the mean mid return differential
\providecommand{\rhemPTEnetHarVsCiHar}{}\renewcommand{\rhemPTEnetHarVsCiHar}{0.704}  % ... its fixed-b p
\providecommand{\rhemDShLgbmHarVsCiHar}{}\renewcommand{\rhemDShLgbmHarVsCiHar}{\ensuremath{-}0.45}  % pooltree_lgbm_baseline minus CI-HAR: Delta SR, mid
\providecommand{\rhemPDShLgbmHarVsCiHar}{}\renewcommand{\rhemPDShLgbmHarVsCiHar}{0.242}  % ... its Ledoit-Wolf p (D19)
\providecommand{\rhemDShXLgbmHarVsCiHar}{}\renewcommand{\rhemDShXLgbmHarVsCiHar}{\ensuremath{-}0.45}  % ... Delta SR at bid/ask
\providecommand{\rhemPDShXLgbmHarVsCiHar}{}\renewcommand{\rhemPDShXLgbmHarVsCiHar}{0.240}  % ... its Ledoit-Wolf p (D19)
\providecommand{\rhemTLgbmHarVsCiHar}{}\renewcommand{\rhemTLgbmHarVsCiHar}{\ensuremath{-}1.16}  % ... fixed-b t of the mean mid return differential
\providecommand{\rhemPTLgbmHarVsCiHar}{}\renewcommand{\rhemPTLgbmHarVsCiHar}{0.269}  % ... its fixed-b p
\providecommand{\rhemDShXgbHarVsCiHar}{}\renewcommand{\rhemDShXgbHarVsCiHar}{\ensuremath{-}0.13}  % pooltree_xgb_baseline minus CI-HAR: Delta SR, mid
\providecommand{\rhemPDShXgbHarVsCiHar}{}\renewcommand{\rhemPDShXgbHarVsCiHar}{0.765}  % ... its Ledoit-Wolf p (D19)
\providecommand{\rhemDShXXgbHarVsCiHar}{}\renewcommand{\rhemDShXXgbHarVsCiHar}{\ensuremath{-}0.13}  % ... Delta SR at bid/ask
\providecommand{\rhemPDShXXgbHarVsCiHar}{}\renewcommand{\rhemPDShXXgbHarVsCiHar}{0.765}  % ... its Ledoit-Wolf p (D19)
\providecommand{\rhemTXgbHarVsCiHar}{}\renewcommand{\rhemTXgbHarVsCiHar}{\ensuremath{-}0.32}  % ... fixed-b t of the mean mid return differential
\providecommand{\rhemPTXgbHarVsCiHar}{}\renewcommand{\rhemPTXgbHarVsCiHar}{0.757}  % ... its fixed-b p
\providecommand{\rhemDShRfHarVsCiHar}{}\renewcommand{\rhemDShRfHarVsCiHar}{\ensuremath{-}0.79}  % pooltree_rf_baseline minus CI-HAR: Delta SR, mid
\providecommand{\rhemPDShRfHarVsCiHar}{}\renewcommand{\rhemPDShRfHarVsCiHar}{0.190}  % ... its Ledoit-Wolf p (D19)
\providecommand{\rhemDShXRfHarVsCiHar}{}\renewcommand{\rhemDShXRfHarVsCiHar}{\ensuremath{-}0.79}  % ... Delta SR at bid/ask
\providecommand{\rhemPDShXRfHarVsCiHar}{}\renewcommand{\rhemPDShXRfHarVsCiHar}{0.191}  % ... its Ledoit-Wolf p (D19)
\providecommand{\rhemTRfHarVsCiHar}{}\renewcommand{\rhemTRfHarVsCiHar}{\ensuremath{-}1.35}  % ... fixed-b t of the mean mid return differential
\providecommand{\rhemPTRfHarVsCiHar}{}\renewcommand{\rhemPTRfHarVsCiHar}{0.199}  % ... its fixed-b p
% ---- 2026-10-26 (D2): every Full forecast against the RTH13 pooled OLS-HAR strategy (primary benchmark)
\providecommand{\rhemShOlsHar}{}\renewcommand{\rhemShOlsHar}{0.95}  % poolpt13_ols_baseline: SR, mid (the OLS-HAR strategy)
\providecommand{\rhemShXOlsHar}{}\renewcommand{\rhemShXOlsHar}{0.60}  % ... at quoted bid/ask
\providecommand{\rhemMeanOlsHar}{}\renewcommand{\rhemMeanOlsHar}{7.2}  % ... mean daily return, %, mid
\providecommand{\rhemDShOlsHarVsShort}{}\renewcommand{\rhemDShOlsHarVsShort}{0.63}  % ... minus the short straddle: Delta SR, mid
\providecommand{\rhemPDShOlsHarVsShort}{}\renewcommand{\rhemPDShOlsHarVsShort}{0.350}  % ... its Ledoit-Wolf p (D19)
\providecommand{\rhemDShRidgeFullVsOls}{}\renewcommand{\rhemDShRidgeFullVsOls}{0.69}  % ridge: Full minus OLS-HAR, Delta SR, mid
\providecommand{\rhemPDShRidgeFullVsOls}{}\renewcommand{\rhemPDShRidgeFullVsOls}{0.171}  % ... its Ledoit-Wolf p (D19)
\providecommand{\rhemDShXRidgeFullVsOls}{}\renewcommand{\rhemDShXRidgeFullVsOls}{0.68}  % ... Delta SR at bid/ask
\providecommand{\rhemPDShXRidgeFullVsOls}{}\renewcommand{\rhemPDShXRidgeFullVsOls}{0.171}  % ... its Ledoit-Wolf p (D19)
\providecommand{\rhemDShLassoFullVsOls}{}\renewcommand{\rhemDShLassoFullVsOls}{0.77}  % lasso: Full minus OLS-HAR, Delta SR, mid
\providecommand{\rhemPDShLassoFullVsOls}{}\renewcommand{\rhemPDShLassoFullVsOls}{0.120}  % ... its Ledoit-Wolf p (D19)
\providecommand{\rhemDShXLassoFullVsOls}{}\renewcommand{\rhemDShXLassoFullVsOls}{0.77}  % ... Delta SR at bid/ask
\providecommand{\rhemPDShXLassoFullVsOls}{}\renewcommand{\rhemPDShXLassoFullVsOls}{0.120}  % ... its Ledoit-Wolf p (D19)
\providecommand{\rhemDShEnetFullVsOls}{}\renewcommand{\rhemDShEnetFullVsOls}{0.90}  % enet: Full minus OLS-HAR, Delta SR, mid
\providecommand{\rhemPDShEnetFullVsOls}{}\renewcommand{\rhemPDShEnetFullVsOls}{0.074}  % ... its Ledoit-Wolf p (D19)
\providecommand{\rhemDShXEnetFullVsOls}{}\renewcommand{\rhemDShXEnetFullVsOls}{0.89}  % ... Delta SR at bid/ask
\providecommand{\rhemPDShXEnetFullVsOls}{}\renewcommand{\rhemPDShXEnetFullVsOls}{0.075}  % ... its Ledoit-Wolf p (D19)
\providecommand{\rhemDShLgbmFullVsOls}{}\renewcommand{\rhemDShLgbmFullVsOls}{0.58}  % lgbm: Full minus OLS-HAR, Delta SR, mid
\providecommand{\rhemPDShLgbmFullVsOls}{}\renewcommand{\rhemPDShLgbmFullVsOls}{0.177}  % ... its Ledoit-Wolf p (D19)
\providecommand{\rhemDShXLgbmFullVsOls}{}\renewcommand{\rhemDShXLgbmFullVsOls}{0.58}  % ... Delta SR at bid/ask
\providecommand{\rhemPDShXLgbmFullVsOls}{}\renewcommand{\rhemPDShXLgbmFullVsOls}{0.177}  % ... its Ledoit-Wolf p (D19)
\providecommand{\rhemDShXgbFullVsOls}{}\renewcommand{\rhemDShXgbFullVsOls}{0.59}  % xgb: Full minus OLS-HAR, Delta SR, mid
\providecommand{\rhemPDShXgbFullVsOls}{}\renewcommand{\rhemPDShXgbFullVsOls}{0.168}  % ... its Ledoit-Wolf p (D19)
\providecommand{\rhemDShXXgbFullVsOls}{}\renewcommand{\rhemDShXXgbFullVsOls}{0.59}  % ... Delta SR at bid/ask
\providecommand{\rhemPDShXXgbFullVsOls}{}\renewcommand{\rhemPDShXXgbFullVsOls}{0.168}  % ... its Ledoit-Wolf p (D19)
\providecommand{\rhemDShRfFullVsOls}{}\renewcommand{\rhemDShRfFullVsOls}{0.92}  % rf: Full minus OLS-HAR, Delta SR, mid
\providecommand{\rhemPDShRfFullVsOls}{}\renewcommand{\rhemPDShRfFullVsOls}{0.039}  % ... its Ledoit-Wolf p (D19)
\providecommand{\rhemDShXRfFullVsOls}{}\renewcommand{\rhemDShXRfFullVsOls}{0.92}  % ... Delta SR at bid/ask
\providecommand{\rhemPDShXRfFullVsOls}{}\renewcommand{\rhemPDShXRfFullVsOls}{0.040}  % ... its Ledoit-Wolf p (D19)
\providecommand{\rhemFullVsOlsMin}{}\renewcommand{\rhemFullVsOlsMin}{0.58}  % Full minus OLS-HAR, Delta SR mid: min over the six
\providecommand{\rhemFullVsOlsMax}{}\renewcommand{\rhemFullVsOlsMax}{0.92}  % ... max
\providecommand{\rhemFullVsOlsMed}{}\renewcommand{\rhemFullVsOlsMed}{0.73}  % ... median
\providecommand{\rhemFullVsOlsPMin}{}\renewcommand{\rhemFullVsOlsPMin}{0.039}  % ... smallest Ledoit-Wolf p
\providecommand{\rhemFullVsOlsPMax}{}\renewcommand{\rhemFullVsOlsPMax}{0.177}  % ... largest Ledoit-Wolf p
\providecommand{\rhemFullVsOlsPos}{}\renewcommand{\rhemFullVsOlsPos}{6}  % ... Delta SR > 0, mid
\providecommand{\rhemFullVsOlsSig}{}\renewcommand{\rhemFullVsOlsSig}{1}  % ... significantly positive (Ledoit-Wolf, 5 %), mid
\providecommand{\rhemFullVsOlsSigWord}{}\renewcommand{\rhemFullVsOlsSigWord}{one}  % ... in words
\providecommand{\rhemFullVsOlsSigX}{}\renewcommand{\rhemFullVsOlsSigX}{1}  % ... significantly positive at bid/ask
\providecommand{\rhemFullVsOlsSigXWord}{}\renewcommand{\rhemFullVsOlsSigXWord}{one}  % ... in words
\providecommand{\rhemFullVsOlsHarSig}{}\renewcommand{\rhemFullVsOlsHarSig}{1}  % = rhemFullVsOlsSig (FRONT's name)
% ---- counts over the six Full forecasts (5 %; Ledoit-Wolf for Delta SR, fixed-b for mean differentials)
\providecommand{\rhemFullVsShortSig}{}\renewcommand{\rhemFullVsShortSig}{3}  % Full forecasts: Delta SR vs short sig., mid
\providecommand{\rhemFullVsShortSigWord}{}\renewcommand{\rhemFullVsShortSigWord}{three}  % ... in words
\providecommand{\rhemFullVsShortSigX}{}\renewcommand{\rhemFullVsShortSigX}{3}  % ... at bid/ask
\providecommand{\rhemFullVsShortSigXWord}{}\renewcommand{\rhemFullVsShortSigXWord}{three}  % ... in words
\providecommand{\rhemFullTVsShortSig}{}\renewcommand{\rhemFullTVsShortSig}{1}  % Full forecasts: mean mid gain over short sig. (fixed-b)
\providecommand{\rhemFullTVsShortSigWord}{}\renewcommand{\rhemFullTVsShortSigWord}{one}  % ... in words
\providecommand{\rhemFullTXVsShortSig}{}\renewcommand{\rhemFullTXVsShortSig}{1}  % ... mean bid/ask gain
\providecommand{\rhemFullTXVsShortSigWord}{}\renewcommand{\rhemFullTXVsShortSigWord}{one}  % ... in words
\providecommand{\rhemFullVsHarPos}{}\renewcommand{\rhemFullVsHarPos}{6}  % Full minus own HAR: Delta SR > 0, mid
\providecommand{\rhemFullVsHarPosWord}{}\renewcommand{\rhemFullVsHarPosWord}{six}  % ... in words
\providecommand{\rhemFullVsHarSig}{}\renewcommand{\rhemFullVsHarSig}{2}  % ... significant (Ledoit-Wolf)
\providecommand{\rhemFullVsHarSigWord}{}\renewcommand{\rhemFullVsHarSigWord}{two}  % ... in words
\providecommand{\rhemFullVsCiHarPos}{}\renewcommand{\rhemFullVsCiHarPos}{6}  % Full minus CI-HAR: Delta SR > 0, mid
\providecommand{\rhemFullVsCiHarPosWord}{}\renewcommand{\rhemFullVsCiHarPosWord}{six}  % ... in words
\providecommand{\rhemFullVsCiHarSig}{}\renewcommand{\rhemFullVsCiHarSig}{0}  % ... significant (Ledoit-Wolf)
\providecommand{\rhemFullVsCiHarSigWord}{}\renewcommand{\rhemFullVsCiHarSigWord}{zero}  % ... in words
\providecommand{\rhemFullVsCiHarPosX}{}\renewcommand{\rhemFullVsCiHarPosX}{6}  % ... Delta SR > 0 at bid/ask
\providecommand{\rhemFullVsCiHarPosXWord}{}\renewcommand{\rhemFullVsCiHarPosXWord}{six}  % ... in words
\providecommand{\rhemFullVsCiHarSigX}{}\renewcommand{\rhemFullVsCiHarSigX}{0}  % ... significant at bid/ask
\providecommand{\rhemFullVsCiHarSigXWord}{}\renewcommand{\rhemFullVsCiHarSigXWord}{zero}  % ... in words
\providecommand{\rhemFullAboveCiHarDOne}{}\renewcommand{\rhemFullAboveCiHarDOne}{6}  % Full forecasts whose D1 SR (mid) exceeds CI-HAR's
\providecommand{\rhemFullAboveCiHarDOneWord}{}\renewcommand{\rhemFullAboveCiHarDOneWord}{six}  % ... in words
\providecommand{\rhemFullAboveCiHarDTwo}{}\renewcommand{\rhemFullAboveCiHarDTwo}{6}  % Full forecasts whose D2 SR (mid) exceeds CI-HAR's
\providecommand{\rhemFullAboveCiHarDTwoWord}{}\renewcommand{\rhemFullAboveCiHarDTwoWord}{six}  % ... in words
\providecommand{\rhemHarVsShortSig}{}\renewcommand{\rhemHarVsShortSig}{0}  % distinct HAR forecasts: Delta SR vs short sig.
\providecommand{\rhemHarVsShortSigWord}{}\renewcommand{\rhemHarVsShortSigWord}{zero}  % ... in words
\providecommand{\rhemHarVsCiHarNeg}{}\renewcommand{\rhemHarVsCiHarNeg}{4}  % distinct HAR forecasts below CI-HAR, mid
\providecommand{\rhemHarVsCiHarNegWord}{}\renewcommand{\rhemHarVsCiHarNegWord}{four}  % ... in words
\providecommand{\rhemHarVsCiHarSig}{}\renewcommand{\rhemHarVsCiHarSig}{0}  % ... significant (Ledoit-Wolf)
\providecommand{\rhemHarVsCiHarSigWord}{}\renewcommand{\rhemHarVsCiHarSigWord}{zero}  % ... in words
\providecommand{\rhemFullVsCiHarMin}{}\renewcommand{\rhemFullVsCiHarMin}{0.43}  % Full minus CI-HAR, Delta SR mid: min over the six
\providecommand{\rhemFullVsCiHarMax}{}\renewcommand{\rhemFullVsCiHarMax}{0.77}  % ... max
\providecommand{\rhemFullVsCiHarMed}{}\renewcommand{\rhemFullVsCiHarMed}{0.58}  % ... median
\providecommand{\rhemFullVsCiHarPMin}{}\renewcommand{\rhemFullVsCiHarPMin}{0.093}  % ... smallest Ledoit-Wolf p
\providecommand{\rhemFullVsCiHarPMax}{}\renewcommand{\rhemFullVsCiHarPMax}{0.358}  % ... largest Ledoit-Wolf p
\providecommand{\rhemFullShMed}{}\renewcommand{\rhemFullShMed}{1.68}  % median SR mid over the 6 distinct Full forecasts
\providecommand{\rhemFullShXMed}{}\renewcommand{\rhemFullShXMed}{1.33}  % ... at bid/ask
\providecommand{\rhemHarShMed}{}\renewcommand{\rhemHarShMed}{0.80}  % median SR mid over the 4 distinct Har forecasts
\providecommand{\rhemHarShXMed}{}\renewcommand{\rhemHarShXMed}{0.45}  % ... at bid/ask
% ---- the pooled linear HAR forecasts coincide (penalties near zero; they are the pooled OLS)
\providecommand{\rhemHarLinMaxAbsDiff}{}\renewcommand{\rhemHarLinMaxAbsDiff}{\ensuremath{1.1\times 10^{-10}}}  % largest |difference| of the 16:00 forecasts, ridge/lasso/enet/OLS
\providecommand{\rhemHarLinMaxRelDiff}{}\renewcommand{\rhemHarLinMaxRelDiff}{\ensuremath{4.6\times 10^{-6}}}  % ... largest relative difference
\providecommand{\rhemHarLinNRows}{}\renewcommand{\rhemHarLinNRows}{2{,}582}  % ... over these 16:00 rows of the run
\providecommand{\rhemHarDistinct}{}\renewcommand{\rhemHarDistinct}{4}  % distinct HAR-set forecasts among the six pooled HAR rows
\providecommand{\rhemHarDistinctWord}{}\renewcommand{\rhemHarDistinctWord}{four}  % ... in words
% ---- break-even half-spread: multiple k* of the quoted half-spread at which the post-warmup mean
% ---- (BeMean) or the paper's Sharpe ratio (BeSh; D19: pooled) reaches zero; ...Pct: k* x the median half-spread, % of premium
\providecommand{\rhemHalfSpreadMed}{}\renewcommand{\rhemHalfSpreadMed}{2.5}  % median quoted half-spread, % of premium, 614 post-warmup days
\providecommand{\rhemHalfSpreadMedAll}{}\renewcommand{\rhemHalfSpreadMedAll}{2.9}  % ... on all 866 option days (= \cmSvCostStraddle)
\providecommand{\rhemHalfSpreadMeanPw}{}\renewcommand{\rhemHalfSpreadMeanPw}{2.7}  % ... mean, 614 post-warmup days
\providecommand{\rhemBeMeanRidgeFull}{}\renewcommand{\rhemBeMeanRidgeFull}{4.43}  % poolpt13_ridge_all_features: k* at which the mean daily return is zero
\providecommand{\rhemBeMeanPctRidgeFull}{}\renewcommand{\rhemBeMeanPctRidgeFull}{11.2}  % ... in % of premium
\providecommand{\rhemBeShRidgeFull}{}\renewcommand{\rhemBeShRidgeFull}{4.43}  % poolpt13_ridge_all_features: k* at which SR is zero
\providecommand{\rhemBeShPctRidgeFull}{}\renewcommand{\rhemBeShPctRidgeFull}{11.2}  % ... in % of premium
\providecommand{\rhemBeMeanLassoFull}{}\renewcommand{\rhemBeMeanLassoFull}{4.63}  % poolpt13_lasso_all_features: k* at which the mean daily return is zero
\providecommand{\rhemBeMeanPctLassoFull}{}\renewcommand{\rhemBeMeanPctLassoFull}{11.7}  % ... in % of premium
\providecommand{\rhemBeShLassoFull}{}\renewcommand{\rhemBeShLassoFull}{4.63}  % poolpt13_lasso_all_features: k* at which SR is zero
\providecommand{\rhemBeShPctLassoFull}{}\renewcommand{\rhemBeShPctLassoFull}{11.7}  % ... in % of premium
\providecommand{\rhemBeMeanEnetFull}{}\renewcommand{\rhemBeMeanEnetFull}{4.94}  % poolpt13_enet_all_features: k* at which the mean daily return is zero
\providecommand{\rhemBeMeanPctEnetFull}{}\renewcommand{\rhemBeMeanPctEnetFull}{12.4}  % ... in % of premium
\providecommand{\rhemBeShEnetFull}{}\renewcommand{\rhemBeShEnetFull}{4.94}  % poolpt13_enet_all_features: k* at which SR is zero
\providecommand{\rhemBeShPctEnetFull}{}\renewcommand{\rhemBeShPctEnetFull}{12.4}  % ... in % of premium
\providecommand{\rhemBeMeanLgbmFull}{}\renewcommand{\rhemBeMeanLgbmFull}{4.19}  % pooltree_lgbm_all_features: k* at which the mean daily return is zero
\providecommand{\rhemBeMeanPctLgbmFull}{}\renewcommand{\rhemBeMeanPctLgbmFull}{10.6}  % ... in % of premium
\providecommand{\rhemBeShLgbmFull}{}\renewcommand{\rhemBeShLgbmFull}{4.19}  % pooltree_lgbm_all_features: k* at which SR is zero
\providecommand{\rhemBeShPctLgbmFull}{}\renewcommand{\rhemBeShPctLgbmFull}{10.6}  % ... in % of premium
\providecommand{\rhemBeMeanXgbFull}{}\renewcommand{\rhemBeMeanXgbFull}{4.21}  % pooltree_xgb_all_features: k* at which the mean daily return is zero
\providecommand{\rhemBeMeanPctXgbFull}{}\renewcommand{\rhemBeMeanPctXgbFull}{10.6}  % ... in % of premium
\providecommand{\rhemBeShXgbFull}{}\renewcommand{\rhemBeShXgbFull}{4.21}  % pooltree_xgb_all_features: k* at which SR is zero
\providecommand{\rhemBeShPctXgbFull}{}\renewcommand{\rhemBeShPctXgbFull}{10.6}  % ... in % of premium
\providecommand{\rhemBeMeanRfFull}{}\renewcommand{\rhemBeMeanRfFull}{5.00}  % pooltree_rf_all_features: k* at which the mean daily return is zero
\providecommand{\rhemBeMeanPctRfFull}{}\renewcommand{\rhemBeMeanPctRfFull}{12.6}  % ... in % of premium
\providecommand{\rhemBeShRfFull}{}\renewcommand{\rhemBeShRfFull}{5.00}  % pooltree_rf_all_features: k* at which SR is zero
\providecommand{\rhemBeShPctRfFull}{}\renewcommand{\rhemBeShPctRfFull}{12.6}  % ... in % of premium
\providecommand{\rhemBeMeanRidgeHar}{}\renewcommand{\rhemBeMeanRidgeHar}{2.67}  % poolpt13_ridge_baseline: k* at which the mean daily return is zero
\providecommand{\rhemBeMeanPctRidgeHar}{}\renewcommand{\rhemBeMeanPctRidgeHar}{6.7}  % ... in % of premium
\providecommand{\rhemBeShRidgeHar}{}\renewcommand{\rhemBeShRidgeHar}{2.67}  % poolpt13_ridge_baseline: k* at which SR is zero
\providecommand{\rhemBeShPctRidgeHar}{}\renewcommand{\rhemBeShPctRidgeHar}{6.7}  % ... in % of premium
\providecommand{\rhemBeMeanLassoHar}{}\renewcommand{\rhemBeMeanLassoHar}{2.67}  % poolpt13_lasso_baseline: k* at which the mean daily return is zero
\providecommand{\rhemBeMeanPctLassoHar}{}\renewcommand{\rhemBeMeanPctLassoHar}{6.7}  % ... in % of premium
\providecommand{\rhemBeShLassoHar}{}\renewcommand{\rhemBeShLassoHar}{2.67}  % poolpt13_lasso_baseline: k* at which SR is zero
\providecommand{\rhemBeShPctLassoHar}{}\renewcommand{\rhemBeShPctLassoHar}{6.7}  % ... in % of premium
\providecommand{\rhemBeMeanEnetHar}{}\renewcommand{\rhemBeMeanEnetHar}{2.67}  % poolpt13_enet_baseline: k* at which the mean daily return is zero
\providecommand{\rhemBeMeanPctEnetHar}{}\renewcommand{\rhemBeMeanPctEnetHar}{6.7}  % ... in % of premium
\providecommand{\rhemBeShEnetHar}{}\renewcommand{\rhemBeShEnetHar}{2.67}  % poolpt13_enet_baseline: k* at which SR is zero
\providecommand{\rhemBeShPctEnetHar}{}\renewcommand{\rhemBeShPctEnetHar}{6.7}  % ... in % of premium
\providecommand{\rhemBeMeanLgbmHar}{}\renewcommand{\rhemBeMeanLgbmHar}{1.87}  % pooltree_lgbm_baseline: k* at which the mean daily return is zero
\providecommand{\rhemBeMeanPctLgbmHar}{}\renewcommand{\rhemBeMeanPctLgbmHar}{4.7}  % ... in % of premium
\providecommand{\rhemBeShLgbmHar}{}\renewcommand{\rhemBeShLgbmHar}{1.87}  % pooltree_lgbm_baseline: k* at which SR is zero
\providecommand{\rhemBeShPctLgbmHar}{}\renewcommand{\rhemBeShPctLgbmHar}{4.7}  % ... in % of premium
\providecommand{\rhemBeMeanXgbHar}{}\renewcommand{\rhemBeMeanXgbHar}{2.78}  % pooltree_xgb_baseline: k* at which the mean daily return is zero
\providecommand{\rhemBeMeanPctXgbHar}{}\renewcommand{\rhemBeMeanPctXgbHar}{7.0}  % ... in % of premium
\providecommand{\rhemBeShXgbHar}{}\renewcommand{\rhemBeShXgbHar}{2.78}  % pooltree_xgb_baseline: k* at which SR is zero
\providecommand{\rhemBeShPctXgbHar}{}\renewcommand{\rhemBeShPctXgbHar}{7.0}  % ... in % of premium
\providecommand{\rhemBeMeanRfHar}{}\renewcommand{\rhemBeMeanRfHar}{0.91}  % pooltree_rf_baseline: k* at which the mean daily return is zero
\providecommand{\rhemBeMeanPctRfHar}{}\renewcommand{\rhemBeMeanPctRfHar}{2.3}  % ... in % of premium
\providecommand{\rhemBeShRfHar}{}\renewcommand{\rhemBeShRfHar}{0.91}  % pooltree_rf_baseline: k* at which SR is zero
\providecommand{\rhemBeShPctRfHar}{}\renewcommand{\rhemBeShPctRfHar}{2.3}  % ... in % of premium
\providecommand{\rhemBeMeanCiHar}{}\renewcommand{\rhemBeMeanCiHar}{3.13}  % sub_ols_baseline: k* at which the mean daily return is zero
\providecommand{\rhemBeMeanPctCiHar}{}\renewcommand{\rhemBeMeanPctCiHar}{7.9}  % ... in % of premium
\providecommand{\rhemBeShCiHar}{}\renewcommand{\rhemBeShCiHar}{3.13}  % sub_ols_baseline: k* at which SR is zero
\providecommand{\rhemBeShPctCiHar}{}\renewcommand{\rhemBeShPctCiHar}{7.9}  % ... in % of premium
\providecommand{\rhemBeMeanShort}{}\renewcommand{\rhemBeMeanShort}{0.91}  % always_short: k* at which the mean daily return is zero
\providecommand{\rhemBeMeanPctShort}{}\renewcommand{\rhemBeMeanPctShort}{2.3}  % ... in % of premium
\providecommand{\rhemBeShShort}{}\renewcommand{\rhemBeShShort}{0.91}  % always_short: k* at which SR is zero
\providecommand{\rhemBeShPctShort}{}\renewcommand{\rhemBeShPctShort}{2.3}  % ... in % of premium
\providecommand{\rhemBeMeanFullMin}{}\renewcommand{\rhemBeMeanFullMin}{4.19}  % k* (mean): min over the distinct Full forecasts with one
\providecommand{\rhemBeMeanFullMax}{}\renewcommand{\rhemBeMeanFullMax}{5.00}  % ... max
\providecommand{\rhemBeMeanFullMed}{}\renewcommand{\rhemBeMeanFullMed}{4.53}  % ... median
\providecommand{\rhemBeMeanFullMinPct}{}\renewcommand{\rhemBeMeanFullMinPct}{10.6}  % ... min, % of premium
\providecommand{\rhemBeMeanFullMaxPct}{}\renewcommand{\rhemBeMeanFullMaxPct}{12.6}  % ... max, % of premium
\providecommand{\rhemBeShFullMin}{}\renewcommand{\rhemBeShFullMin}{4.19}  % k* (sh): min over the distinct Full forecasts with one
\providecommand{\rhemBeShFullMax}{}\renewcommand{\rhemBeShFullMax}{5.00}  % ... max
\providecommand{\rhemBeShFullMed}{}\renewcommand{\rhemBeShFullMed}{4.53}  % ... median
\providecommand{\rhemBeShFullMinPct}{}\renewcommand{\rhemBeShFullMinPct}{10.6}  % ... min, % of premium
\providecommand{\rhemBeShFullMaxPct}{}\renewcommand{\rhemBeShFullMaxPct}{12.6}  % ... max, % of premium
\providecommand{\rhemBeMeanHarMin}{}\renewcommand{\rhemBeMeanHarMin}{0.91}  % k* (mean): min over the distinct Har forecasts with one
\providecommand{\rhemBeMeanHarMax}{}\renewcommand{\rhemBeMeanHarMax}{2.78}  % ... max
\providecommand{\rhemBeMeanHarMed}{}\renewcommand{\rhemBeMeanHarMed}{2.27}  % ... median
\providecommand{\rhemBeMeanHarMinPct}{}\renewcommand{\rhemBeMeanHarMinPct}{2.3}  % ... min, % of premium
\providecommand{\rhemBeMeanHarMaxPct}{}\renewcommand{\rhemBeMeanHarMaxPct}{7.0}  % ... max, % of premium
\providecommand{\rhemBeShHarMin}{}\renewcommand{\rhemBeShHarMin}{0.91}  % k* (sh): min over the distinct Har forecasts with one
\providecommand{\rhemBeShHarMax}{}\renewcommand{\rhemBeShHarMax}{2.78}  % ... max
\providecommand{\rhemBeShHarMed}{}\renewcommand{\rhemBeShHarMed}{2.27}  % ... median
\providecommand{\rhemBeShHarMinPct}{}\renewcommand{\rhemBeShHarMinPct}{2.3}  % ... min, % of premium
\providecommand{\rhemBeShHarMaxPct}{}\renewcommand{\rhemBeShHarMaxPct}{7.0}  % ... max, % of premium
% ---- by calendar year (mid): SR_y, each year annualized by sqrt(A_y) (descriptive; their mean is D6); years One.. as \rheDSixYearOne..
\providecommand{\rhemYearOne}{}\renewcommand{\rhemYearOne}{2021}  % calendar year One of the post-warmup days
\providecommand{\rhemNyOne}{}\renewcommand{\rhemNyOne}{64}  % ... N_y
\providecommand{\rhemByYearFullMedOne}{}\renewcommand{\rhemByYearFullMedOne}{1.47}  % 2021: median SR_y of the six Full forecasts
\providecommand{\rhemByYearFullMinOne}{}\renewcommand{\rhemByYearFullMinOne}{1.24}  % ... min
\providecommand{\rhemByYearFullMaxOne}{}\renewcommand{\rhemByYearFullMaxOne}{2.00}  % ... max
\providecommand{\rhemByYearHarMedOne}{}\renewcommand{\rhemByYearHarMedOne}{0.12}  % 2021: median SR_y of the four distinct HAR forecasts
\providecommand{\rhemByYearHarMinOne}{}\renewcommand{\rhemByYearHarMinOne}{\ensuremath{-}1.91}  % ... min
\providecommand{\rhemByYearHarMaxOne}{}\renewcommand{\rhemByYearHarMaxOne}{0.27}  % ... max
\providecommand{\rhemByYearCiHarOne}{}\renewcommand{\rhemByYearCiHarOne}{0.12}  % 2021: SR_y of CI-HAR
\providecommand{\rhemByYearShortOne}{}\renewcommand{\rhemByYearShortOne}{\ensuremath{-}1.10}  % 2021: SR_y of the short straddle
\providecommand{\rhemYearTwo}{}\renewcommand{\rhemYearTwo}{2022}  % calendar year Two of the post-warmup days
\providecommand{\rhemNyTwo}{}\renewcommand{\rhemNyTwo}{219}  % ... N_y
\providecommand{\rhemByYearFullMedTwo}{}\renewcommand{\rhemByYearFullMedTwo}{1.45}  % 2022: median SR_y of the six Full forecasts
\providecommand{\rhemByYearFullMinTwo}{}\renewcommand{\rhemByYearFullMinTwo}{0.96}  % ... min
\providecommand{\rhemByYearFullMaxTwo}{}\renewcommand{\rhemByYearFullMaxTwo}{1.86}  % ... max
\providecommand{\rhemByYearHarMedTwo}{}\renewcommand{\rhemByYearHarMedTwo}{1.16}  % 2022: median SR_y of the four distinct HAR forecasts
\providecommand{\rhemByYearHarMinTwo}{}\renewcommand{\rhemByYearHarMinTwo}{1.04}  % ... min
\providecommand{\rhemByYearHarMaxTwo}{}\renewcommand{\rhemByYearHarMaxTwo}{1.74}  % ... max
\providecommand{\rhemByYearCiHarTwo}{}\renewcommand{\rhemByYearCiHarTwo}{1.67}  % 2022: SR_y of CI-HAR
\providecommand{\rhemByYearShortTwo}{}\renewcommand{\rhemByYearShortTwo}{1.81}  % 2022: SR_y of the short straddle
\providecommand{\rhemYearThree}{}\renewcommand{\rhemYearThree}{2023}  % calendar year Three of the post-warmup days
\providecommand{\rhemNyThree}{}\renewcommand{\rhemNyThree}{248}  % ... N_y
\providecommand{\rhemByYearFullMedThree}{}\renewcommand{\rhemByYearFullMedThree}{1.66}  % 2023: median SR_y of the six Full forecasts
\providecommand{\rhemByYearFullMinThree}{}\renewcommand{\rhemByYearFullMinThree}{1.55}  % ... min
\providecommand{\rhemByYearFullMaxThree}{}\renewcommand{\rhemByYearFullMaxThree}{1.94}  % ... max
\providecommand{\rhemByYearHarMedThree}{}\renewcommand{\rhemByYearHarMedThree}{0.77}  % 2023: median SR_y of the four distinct HAR forecasts
\providecommand{\rhemByYearHarMinThree}{}\renewcommand{\rhemByYearHarMinThree}{0.64}  % ... min
\providecommand{\rhemByYearHarMaxThree}{}\renewcommand{\rhemByYearHarMaxThree}{1.30}  % ... max
\providecommand{\rhemByYearCiHarThree}{}\renewcommand{\rhemByYearCiHarThree}{1.33}  % 2023: SR_y of CI-HAR
\providecommand{\rhemByYearShortThree}{}\renewcommand{\rhemByYearShortThree}{0.88}  % 2023: SR_y of the short straddle
\providecommand{\rhemYearFour}{}\renewcommand{\rhemYearFour}{2024}  % calendar year Four of the post-warmup days
\providecommand{\rhemNyFour}{}\renewcommand{\rhemNyFour}{83}  % ... N_y
\providecommand{\rhemByYearFullMedFour}{}\renewcommand{\rhemByYearFullMedFour}{2.63}  % 2024: median SR_y of the six Full forecasts
\providecommand{\rhemByYearFullMinFour}{}\renewcommand{\rhemByYearFullMinFour}{2.00}  % ... min
\providecommand{\rhemByYearFullMaxFour}{}\renewcommand{\rhemByYearFullMaxFour}{2.99}  % ... max
\providecommand{\rhemByYearHarMedFour}{}\renewcommand{\rhemByYearHarMedFour}{0.87}  % 2024: median SR_y of the four distinct HAR forecasts
\providecommand{\rhemByYearHarMinFour}{}\renewcommand{\rhemByYearHarMinFour}{\ensuremath{-}0.18}  % ... min
\providecommand{\rhemByYearHarMaxFour}{}\renewcommand{\rhemByYearHarMaxFour}{1.29}  % ... max
\providecommand{\rhemByYearCiHarFour}{}\renewcommand{\rhemByYearCiHarFour}{0.82}  % 2024: SR_y of CI-HAR
\providecommand{\rhemByYearShortFour}{}\renewcommand{\rhemByYearShortFour}{\ensuremath{-}2.05}  % 2024: SR_y of the short straddle
\providecommand{\rhemNYears}{}\renewcommand{\rhemNYears}{4}  % calendar years of the post-warmup days
\providecommand{\rhemNYearsWord}{}\renewcommand{\rhemNYearsWord}{four}  % ... in words
\providecommand{\rhemYearsFullAllPos}{}\renewcommand{\rhemYearsFullAllPos}{4}  % years in which all six Full forecasts have SR_y > 0
\providecommand{\rhemYearsFullAllPosWord}{}\renewcommand{\rhemYearsFullAllPosWord}{four}  % ... in words
\providecommand{\rhemYearsFullAllAboveShort}{}\renewcommand{\rhemYearsFullAllAboveShort}{3}  % years in which all six Full SR_y exceed the short's
\providecommand{\rhemYearsFullAllAboveShortWord}{}\renewcommand{\rhemYearsFullAllAboveShortWord}{three}  % ... in words
\providecommand{\rhemYearsFullMedAboveCiHar}{}\renewcommand{\rhemYearsFullMedAboveCiHar}{3}  % years in which the Full median SR_y exceeds CI-HAR's
\providecommand{\rhemYearsFullMedAboveCiHarWord}{}\renewcommand{\rhemYearsFullMedAboveCiHarWord}{three}  % ... in words
\providecommand{\rhemYearsFullMedAboveHarMed}{}\renewcommand{\rhemYearsFullMedAboveHarMed}{4}  % years in which the Full median exceeds the HAR median
\providecommand{\rhemYearsFullMedAboveHarMedWord}{}\renewcommand{\rhemYearsFullMedAboveHarMedWord}{four}  % ... in words
\providecommand{\rhemYearsFullMedAboveShort}{}\renewcommand{\rhemYearsFullMedAboveShort}{3}  % years in which the Full median exceeds the short
\providecommand{\rhemYearsFullMedAboveShortWord}{}\renewcommand{\rhemYearsFullMedAboveShortWord}{three}  % ... in words
\providecommand{\rhemYearsCiHarAboveAllFull}{}\renewcommand{\rhemYearsCiHarAboveAllFull}{0}  % years in which CI-HAR's SR_y exceeds every Full one
\providecommand{\rhemYearsCiHarAboveAllFullWord}{}\renewcommand{\rhemYearsCiHarAboveAllFullWord}{zero}  % ... in words
\providecommand{\rhemYearsCiHarPos}{}\renewcommand{\rhemYearsCiHarPos}{4}  % years in which CI-HAR has SR_y > 0
\providecommand{\rhemYearsCiHarPosWord}{}\renewcommand{\rhemYearsCiHarPosWord}{four}  % ... in words
\providecommand{\rhemNPartialDays}{}\renewcommand{\rhemNPartialDays}{147}  % post-warmup days in the partial first and last calendar years
\providecommand{\rhemPartialDaysPct}{}\renewcommand{\rhemPartialDaysPct}{24}  % ... in % of the post-warmup days
\providecommand{\rhemYearsShortNeg}{}\renewcommand{\rhemYearsShortNeg}{2}  % years in which the short straddle has SR_y < 0
\providecommand{\rhemYearsShortNegWord}{}\renewcommand{\rhemYearsShortNegWord}{two}  % ... in words
% ---- 2026-10-29 (X8, Q-2026-10-29-21): the sign strategy on the HAR + RM + FOMC forecasts (ablation arm
% add_rm_fomc); proposed after the author's 27 Oct remark; not pre-registered. Rule: results/econ_rmfomc_2026-10-29/RULE.txt.
% Row stems <Est>RmFomc (Est = Ols, Ridge, Lasso, Enet, Lgbm, Xgb, Rf). Delta SR at mid (X: at quoted bid/ask)
% with its Ledoit-Wolf p (D19); T: fixed-b t of the mean daily mid-return differential (S = 33).
\providecommand{\rhemSigPct}{}\renewcommand{\rhemSigPct}{5}  % significance level of the X8 Ledoit-Wolf counts, %
\providecommand{\rhemShOlsRmFomc}{}\renewcommand{\rhemShOlsRmFomc}{0.96}  % rmfomc_ols: SR, mid
\providecommand{\rhemShXOlsRmFomc}{}\renewcommand{\rhemShXOlsRmFomc}{0.60}  % ... at quoted bid/ask
\providecommand{\rhemMeanOlsRmFomc}{}\renewcommand{\rhemMeanOlsRmFomc}{7.3}  % ... mean daily return, %, mid
\providecommand{\rhemMeanXOlsRmFomc}{}\renewcommand{\rhemMeanXOlsRmFomc}{4.6}  % ... mean daily return, %, bid/ask
\providecommand{\rhemDShOlsRmFomcVsShort}{}\renewcommand{\rhemDShOlsRmFomcVsShort}{0.64}  % ... minus the short straddle: Delta SR, mid
\providecommand{\rhemPDShOlsRmFomcVsShort}{}\renewcommand{\rhemPDShOlsRmFomcVsShort}{0.213}  % ... its Ledoit-Wolf p
\providecommand{\rhemDShXOlsRmFomcVsShort}{}\renewcommand{\rhemDShXOlsRmFomcVsShort}{0.63}  % ... Delta SR at bid/ask
\providecommand{\rhemPDShXOlsRmFomcVsShort}{}\renewcommand{\rhemPDShXOlsRmFomcVsShort}{0.212}  % ... its Ledoit-Wolf p
\providecommand{\rhemTOlsRmFomcVsShort}{}\renewcommand{\rhemTOlsRmFomcVsShort}{1.19}  % ... fixed-b t of the mean mid return differential
\providecommand{\rhemPTOlsRmFomcVsShort}{}\renewcommand{\rhemPTOlsRmFomcVsShort}{0.257}  % ... its fixed-b p
\providecommand{\rhemDShOlsRmFomcVsHar}{}\renewcommand{\rhemDShOlsRmFomcVsHar}{0.01}  % ... minus the class's HAR set: Delta SR, mid
\providecommand{\rhemPDShOlsRmFomcVsHar}{}\renewcommand{\rhemPDShOlsRmFomcVsHar}{0.992}  % ... its Ledoit-Wolf p
\providecommand{\rhemDShXOlsRmFomcVsHar}{}\renewcommand{\rhemDShXOlsRmFomcVsHar}{0.00}  % ... Delta SR at bid/ask
\providecommand{\rhemPDShXOlsRmFomcVsHar}{}\renewcommand{\rhemPDShXOlsRmFomcVsHar}{0.997}  % ... its Ledoit-Wolf p
\providecommand{\rhemDShOlsRmFomcVsOls}{}\renewcommand{\rhemDShOlsRmFomcVsOls}{0.01}  % ... minus the OLS-HAR: Delta SR, mid
\providecommand{\rhemPDShOlsRmFomcVsOls}{}\renewcommand{\rhemPDShOlsRmFomcVsOls}{0.992}  % ... its Ledoit-Wolf p
\providecommand{\rhemDShXOlsRmFomcVsOls}{}\renewcommand{\rhemDShXOlsRmFomcVsOls}{0.00}  % ... Delta SR at bid/ask
\providecommand{\rhemPDShXOlsRmFomcVsOls}{}\renewcommand{\rhemPDShXOlsRmFomcVsOls}{0.997}  % ... its Ledoit-Wolf p
\providecommand{\rhemShRidgeRmFomc}{}\renewcommand{\rhemShRidgeRmFomc}{1.58}  % rmfomc_ridge: SR, mid
\providecommand{\rhemShXRidgeRmFomc}{}\renewcommand{\rhemShXRidgeRmFomc}{1.21}  % ... at quoted bid/ask
\providecommand{\rhemMeanRidgeRmFomc}{}\renewcommand{\rhemMeanRidgeRmFomc}{11.9}  % ... mean daily return, %, mid
\providecommand{\rhemMeanXRidgeRmFomc}{}\renewcommand{\rhemMeanXRidgeRmFomc}{9.2}  % ... mean daily return, %, bid/ask
\providecommand{\rhemDShRidgeRmFomcVsShort}{}\renewcommand{\rhemDShRidgeRmFomcVsShort}{1.26}  % ... minus the short straddle: Delta SR, mid
\providecommand{\rhemPDShRidgeRmFomcVsShort}{}\renewcommand{\rhemPDShRidgeRmFomcVsShort}{0.034}  % ... its Ledoit-Wolf p
\providecommand{\rhemDShXRidgeRmFomcVsShort}{}\renewcommand{\rhemDShXRidgeRmFomcVsShort}{1.24}  % ... Delta SR at bid/ask
\providecommand{\rhemPDShXRidgeRmFomcVsShort}{}\renewcommand{\rhemPDShXRidgeRmFomcVsShort}{0.034}  % ... its Ledoit-Wolf p
\providecommand{\rhemTRidgeRmFomcVsShort}{}\renewcommand{\rhemTRidgeRmFomcVsShort}{2.07}  % ... fixed-b t of the mean mid return differential
\providecommand{\rhemPTRidgeRmFomcVsShort}{}\renewcommand{\rhemPTRidgeRmFomcVsShort}{0.054}  % ... its fixed-b p
\providecommand{\rhemDShRidgeRmFomcVsHar}{}\renewcommand{\rhemDShRidgeRmFomcVsHar}{0.63}  % ... minus the class's HAR set: Delta SR, mid
\providecommand{\rhemPDShRidgeRmFomcVsHar}{}\renewcommand{\rhemPDShRidgeRmFomcVsHar}{0.382}  % ... its Ledoit-Wolf p
\providecommand{\rhemDShXRidgeRmFomcVsHar}{}\renewcommand{\rhemDShXRidgeRmFomcVsHar}{0.61}  % ... Delta SR at bid/ask
\providecommand{\rhemPDShXRidgeRmFomcVsHar}{}\renewcommand{\rhemPDShXRidgeRmFomcVsHar}{0.401}  % ... its Ledoit-Wolf p
\providecommand{\rhemDShRidgeRmFomcVsOls}{}\renewcommand{\rhemDShRidgeRmFomcVsOls}{0.63}  % ... minus the OLS-HAR: Delta SR, mid
\providecommand{\rhemPDShRidgeRmFomcVsOls}{}\renewcommand{\rhemPDShRidgeRmFomcVsOls}{0.382}  % ... its Ledoit-Wolf p
\providecommand{\rhemDShXRidgeRmFomcVsOls}{}\renewcommand{\rhemDShXRidgeRmFomcVsOls}{0.61}  % ... Delta SR at bid/ask
\providecommand{\rhemPDShXRidgeRmFomcVsOls}{}\renewcommand{\rhemPDShXRidgeRmFomcVsOls}{0.401}  % ... its Ledoit-Wolf p
\providecommand{\rhemShLassoRmFomc}{}\renewcommand{\rhemShLassoRmFomc}{1.51}  % rmfomc_lasso: SR, mid
\providecommand{\rhemShXLassoRmFomc}{}\renewcommand{\rhemShXLassoRmFomc}{1.14}  % ... at quoted bid/ask
\providecommand{\rhemMeanLassoRmFomc}{}\renewcommand{\rhemMeanLassoRmFomc}{11.4}  % ... mean daily return, %, mid
\providecommand{\rhemMeanXLassoRmFomc}{}\renewcommand{\rhemMeanXLassoRmFomc}{8.7}  % ... mean daily return, %, bid/ask
\providecommand{\rhemDShLassoRmFomcVsShort}{}\renewcommand{\rhemDShLassoRmFomcVsShort}{1.19}  % ... minus the short straddle: Delta SR, mid
\providecommand{\rhemPDShLassoRmFomcVsShort}{}\renewcommand{\rhemPDShLassoRmFomcVsShort}{0.050}  % ... its Ledoit-Wolf p
\providecommand{\rhemDShXLassoRmFomcVsShort}{}\renewcommand{\rhemDShXLassoRmFomcVsShort}{1.17}  % ... Delta SR at bid/ask
\providecommand{\rhemPDShXLassoRmFomcVsShort}{}\renewcommand{\rhemPDShXLassoRmFomcVsShort}{0.051}  % ... its Ledoit-Wolf p
\providecommand{\rhemTLassoRmFomcVsShort}{}\renewcommand{\rhemTLassoRmFomcVsShort}{1.83}  % ... fixed-b t of the mean mid return differential
\providecommand{\rhemPTLassoRmFomcVsShort}{}\renewcommand{\rhemPTLassoRmFomcVsShort}{0.085}  % ... its fixed-b p
\providecommand{\rhemDShLassoRmFomcVsHar}{}\renewcommand{\rhemDShLassoRmFomcVsHar}{0.56}  % ... minus the class's HAR set: Delta SR, mid
\providecommand{\rhemPDShLassoRmFomcVsHar}{}\renewcommand{\rhemPDShLassoRmFomcVsHar}{0.434}  % ... its Ledoit-Wolf p
\providecommand{\rhemDShXLassoRmFomcVsHar}{}\renewcommand{\rhemDShXLassoRmFomcVsHar}{0.54}  % ... Delta SR at bid/ask
\providecommand{\rhemPDShXLassoRmFomcVsHar}{}\renewcommand{\rhemPDShXLassoRmFomcVsHar}{0.455}  % ... its Ledoit-Wolf p
\providecommand{\rhemDShLassoRmFomcVsOls}{}\renewcommand{\rhemDShLassoRmFomcVsOls}{0.56}  % ... minus the OLS-HAR: Delta SR, mid
\providecommand{\rhemPDShLassoRmFomcVsOls}{}\renewcommand{\rhemPDShLassoRmFomcVsOls}{0.434}  % ... its Ledoit-Wolf p
\providecommand{\rhemDShXLassoRmFomcVsOls}{}\renewcommand{\rhemDShXLassoRmFomcVsOls}{0.54}  % ... Delta SR at bid/ask
\providecommand{\rhemPDShXLassoRmFomcVsOls}{}\renewcommand{\rhemPDShXLassoRmFomcVsOls}{0.455}  % ... its Ledoit-Wolf p
\providecommand{\rhemShEnetRmFomc}{}\renewcommand{\rhemShEnetRmFomc}{2.03}  % rmfomc_enet: SR, mid
\providecommand{\rhemShXEnetRmFomc}{}\renewcommand{\rhemShXEnetRmFomc}{1.68}  % ... at quoted bid/ask
\providecommand{\rhemMeanEnetRmFomc}{}\renewcommand{\rhemMeanEnetRmFomc}{15.2}  % ... mean daily return, %, mid
\providecommand{\rhemMeanXEnetRmFomc}{}\renewcommand{\rhemMeanXEnetRmFomc}{12.6}  % ... mean daily return, %, bid/ask
\providecommand{\rhemDShEnetRmFomcVsShort}{}\renewcommand{\rhemDShEnetRmFomcVsShort}{1.71}  % ... minus the short straddle: Delta SR, mid
\providecommand{\rhemPDShEnetRmFomcVsShort}{}\renewcommand{\rhemPDShEnetRmFomcVsShort}{0.013}  % ... its Ledoit-Wolf p
\providecommand{\rhemDShXEnetRmFomcVsShort}{}\renewcommand{\rhemDShXEnetRmFomcVsShort}{1.71}  % ... Delta SR at bid/ask
\providecommand{\rhemPDShXEnetRmFomcVsShort}{}\renewcommand{\rhemPDShXEnetRmFomcVsShort}{0.014}  % ... its Ledoit-Wolf p
\providecommand{\rhemTEnetRmFomcVsShort}{}\renewcommand{\rhemTEnetRmFomcVsShort}{2.29}  % ... fixed-b t of the mean mid return differential
\providecommand{\rhemPTEnetRmFomcVsShort}{}\renewcommand{\rhemPTEnetRmFomcVsShort}{0.034}  % ... its fixed-b p
\providecommand{\rhemDShEnetRmFomcVsHar}{}\renewcommand{\rhemDShEnetRmFomcVsHar}{1.08}  % ... minus the class's HAR set: Delta SR, mid
\providecommand{\rhemPDShEnetRmFomcVsHar}{}\renewcommand{\rhemPDShEnetRmFomcVsHar}{0.014}  % ... its Ledoit-Wolf p
\providecommand{\rhemDShXEnetRmFomcVsHar}{}\renewcommand{\rhemDShXEnetRmFomcVsHar}{1.08}  % ... Delta SR at bid/ask
\providecommand{\rhemPDShXEnetRmFomcVsHar}{}\renewcommand{\rhemPDShXEnetRmFomcVsHar}{0.014}  % ... its Ledoit-Wolf p
\providecommand{\rhemDShEnetRmFomcVsOls}{}\renewcommand{\rhemDShEnetRmFomcVsOls}{1.08}  % ... minus the OLS-HAR: Delta SR, mid
\providecommand{\rhemPDShEnetRmFomcVsOls}{}\renewcommand{\rhemPDShEnetRmFomcVsOls}{0.014}  % ... its Ledoit-Wolf p
\providecommand{\rhemDShXEnetRmFomcVsOls}{}\renewcommand{\rhemDShXEnetRmFomcVsOls}{1.08}  % ... Delta SR at bid/ask
\providecommand{\rhemPDShXEnetRmFomcVsOls}{}\renewcommand{\rhemPDShXEnetRmFomcVsOls}{0.014}  % ... its Ledoit-Wolf p
\providecommand{\rhemShLgbmRmFomc}{}\renewcommand{\rhemShLgbmRmFomc}{1.68}  % rmfomc_lgbm: SR, mid
\providecommand{\rhemShXLgbmRmFomc}{}\renewcommand{\rhemShXLgbmRmFomc}{1.33}  % ... at quoted bid/ask
\providecommand{\rhemMeanLgbmRmFomc}{}\renewcommand{\rhemMeanLgbmRmFomc}{12.6}  % ... mean daily return, %, mid
\providecommand{\rhemMeanXLgbmRmFomc}{}\renewcommand{\rhemMeanXLgbmRmFomc}{10.0}  % ... mean daily return, %, bid/ask
\providecommand{\rhemDShLgbmRmFomcVsShort}{}\renewcommand{\rhemDShLgbmRmFomcVsShort}{1.36}  % ... minus the short straddle: Delta SR, mid
\providecommand{\rhemPDShLgbmRmFomcVsShort}{}\renewcommand{\rhemPDShLgbmRmFomcVsShort}{0.042}  % ... its Ledoit-Wolf p
\providecommand{\rhemDShXLgbmRmFomcVsShort}{}\renewcommand{\rhemDShXLgbmRmFomcVsShort}{1.36}  % ... Delta SR at bid/ask
\providecommand{\rhemPDShXLgbmRmFomcVsShort}{}\renewcommand{\rhemPDShXLgbmRmFomcVsShort}{0.040}  % ... its Ledoit-Wolf p
\providecommand{\rhemTLgbmRmFomcVsShort}{}\renewcommand{\rhemTLgbmRmFomcVsShort}{1.76}  % ... fixed-b t of the mean mid return differential
\providecommand{\rhemPTLgbmRmFomcVsShort}{}\renewcommand{\rhemPTLgbmRmFomcVsShort}{0.098}  % ... its fixed-b p
\providecommand{\rhemDShLgbmRmFomcVsHar}{}\renewcommand{\rhemDShLgbmRmFomcVsHar}{1.03}  % ... minus the class's HAR set: Delta SR, mid
\providecommand{\rhemPDShLgbmRmFomcVsHar}{}\renewcommand{\rhemPDShLgbmRmFomcVsHar}{0.021}  % ... its Ledoit-Wolf p
\providecommand{\rhemDShXLgbmRmFomcVsHar}{}\renewcommand{\rhemDShXLgbmRmFomcVsHar}{1.02}  % ... Delta SR at bid/ask
\providecommand{\rhemPDShXLgbmRmFomcVsHar}{}\renewcommand{\rhemPDShXLgbmRmFomcVsHar}{0.021}  % ... its Ledoit-Wolf p
\providecommand{\rhemDShLgbmRmFomcVsOls}{}\renewcommand{\rhemDShLgbmRmFomcVsOls}{0.73}  % ... minus the OLS-HAR: Delta SR, mid
\providecommand{\rhemPDShLgbmRmFomcVsOls}{}\renewcommand{\rhemPDShLgbmRmFomcVsOls}{0.086}  % ... its Ledoit-Wolf p
\providecommand{\rhemDShXLgbmRmFomcVsOls}{}\renewcommand{\rhemDShXLgbmRmFomcVsOls}{0.73}  % ... Delta SR at bid/ask
\providecommand{\rhemPDShXLgbmRmFomcVsOls}{}\renewcommand{\rhemPDShXLgbmRmFomcVsOls}{0.085}  % ... its Ledoit-Wolf p
\providecommand{\rhemShXgbRmFomc}{}\renewcommand{\rhemShXgbRmFomc}{1.46}  % rmfomc_xgb: SR, mid
\providecommand{\rhemShXXgbRmFomc}{}\renewcommand{\rhemShXXgbRmFomc}{1.10}  % ... at quoted bid/ask
\providecommand{\rhemMeanXgbRmFomc}{}\renewcommand{\rhemMeanXgbRmFomc}{10.9}  % ... mean daily return, %, mid
\providecommand{\rhemMeanXXgbRmFomc}{}\renewcommand{\rhemMeanXXgbRmFomc}{8.3}  % ... mean daily return, %, bid/ask
\providecommand{\rhemDShXgbRmFomcVsShort}{}\renewcommand{\rhemDShXgbRmFomcVsShort}{1.13}  % ... minus the short straddle: Delta SR, mid
\providecommand{\rhemPDShXgbRmFomcVsShort}{}\renewcommand{\rhemPDShXgbRmFomcVsShort}{0.100}  % ... its Ledoit-Wolf p
\providecommand{\rhemDShXXgbRmFomcVsShort}{}\renewcommand{\rhemDShXXgbRmFomcVsShort}{1.14}  % ... Delta SR at bid/ask
\providecommand{\rhemPDShXXgbRmFomcVsShort}{}\renewcommand{\rhemPDShXXgbRmFomcVsShort}{0.099}  % ... its Ledoit-Wolf p
\providecommand{\rhemTXgbRmFomcVsShort}{}\renewcommand{\rhemTXgbRmFomcVsShort}{1.48}  % ... fixed-b t of the mean mid return differential
\providecommand{\rhemPTXgbRmFomcVsShort}{}\renewcommand{\rhemPTXgbRmFomcVsShort}{0.162}  % ... its fixed-b p
\providecommand{\rhemDShXgbRmFomcVsHar}{}\renewcommand{\rhemDShXgbRmFomcVsHar}{0.47}  % ... minus the class's HAR set: Delta SR, mid
\providecommand{\rhemPDShXgbRmFomcVsHar}{}\renewcommand{\rhemPDShXgbRmFomcVsHar}{0.234}  % ... its Ledoit-Wolf p
\providecommand{\rhemDShXXgbRmFomcVsHar}{}\renewcommand{\rhemDShXXgbRmFomcVsHar}{0.47}  % ... Delta SR at bid/ask
\providecommand{\rhemPDShXXgbRmFomcVsHar}{}\renewcommand{\rhemPDShXXgbRmFomcVsHar}{0.238}  % ... its Ledoit-Wolf p
\providecommand{\rhemDShXgbRmFomcVsOls}{}\renewcommand{\rhemDShXgbRmFomcVsOls}{0.50}  % ... minus the OLS-HAR: Delta SR, mid
\providecommand{\rhemPDShXgbRmFomcVsOls}{}\renewcommand{\rhemPDShXgbRmFomcVsOls}{0.212}  % ... its Ledoit-Wolf p
\providecommand{\rhemDShXXgbRmFomcVsOls}{}\renewcommand{\rhemDShXXgbRmFomcVsOls}{0.50}  % ... Delta SR at bid/ask
\providecommand{\rhemPDShXXgbRmFomcVsOls}{}\renewcommand{\rhemPDShXXgbRmFomcVsOls}{0.209}  % ... its Ledoit-Wolf p
\providecommand{\rhemShRfRmFomc}{}\renewcommand{\rhemShRfRmFomc}{1.50}  % rmfomc_rf: SR, mid
\providecommand{\rhemShXRfRmFomc}{}\renewcommand{\rhemShXRfRmFomc}{1.15}  % ... at quoted bid/ask
\providecommand{\rhemMeanRfRmFomc}{}\renewcommand{\rhemMeanRfRmFomc}{11.3}  % ... mean daily return, %, mid
\providecommand{\rhemMeanXRfRmFomc}{}\renewcommand{\rhemMeanXRfRmFomc}{8.6}  % ... mean daily return, %, bid/ask
\providecommand{\rhemDShRfRmFomcVsShort}{}\renewcommand{\rhemDShRfRmFomcVsShort}{1.18}  % ... minus the short straddle: Delta SR, mid
\providecommand{\rhemPDShRfRmFomcVsShort}{}\renewcommand{\rhemPDShRfRmFomcVsShort}{0.082}  % ... its Ledoit-Wolf p
\providecommand{\rhemDShXRfRmFomcVsShort}{}\renewcommand{\rhemDShXRfRmFomcVsShort}{1.18}  % ... Delta SR at bid/ask
\providecommand{\rhemPDShXRfRmFomcVsShort}{}\renewcommand{\rhemPDShXRfRmFomcVsShort}{0.082}  % ... its Ledoit-Wolf p
\providecommand{\rhemTRfRmFomcVsShort}{}\renewcommand{\rhemTRfRmFomcVsShort}{1.59}  % ... fixed-b t of the mean mid return differential
\providecommand{\rhemPTRfRmFomcVsShort}{}\renewcommand{\rhemPTRfRmFomcVsShort}{0.133}  % ... its fixed-b p
\providecommand{\rhemDShRfRmFomcVsHar}{}\renewcommand{\rhemDShRfRmFomcVsHar}{1.18}  % ... minus the class's HAR set: Delta SR, mid
\providecommand{\rhemPDShRfRmFomcVsHar}{}\renewcommand{\rhemPDShRfRmFomcVsHar}{0.064}  % ... its Ledoit-Wolf p
\providecommand{\rhemDShXRfRmFomcVsHar}{}\renewcommand{\rhemDShXRfRmFomcVsHar}{1.18}  % ... Delta SR at bid/ask
\providecommand{\rhemPDShXRfRmFomcVsHar}{}\renewcommand{\rhemPDShXRfRmFomcVsHar}{0.066}  % ... its Ledoit-Wolf p
\providecommand{\rhemDShRfRmFomcVsOls}{}\renewcommand{\rhemDShRfRmFomcVsOls}{0.55}  % ... minus the OLS-HAR: Delta SR, mid
\providecommand{\rhemPDShRfRmFomcVsOls}{}\renewcommand{\rhemPDShRfRmFomcVsOls}{0.239}  % ... its Ledoit-Wolf p
\providecommand{\rhemDShXRfRmFomcVsOls}{}\renewcommand{\rhemDShXRfRmFomcVsOls}{0.55}  % ... Delta SR at bid/ask
\providecommand{\rhemPDShXRfRmFomcVsOls}{}\renewcommand{\rhemPDShXRfRmFomcVsOls}{0.239}  % ... its Ledoit-Wolf p
% ---- the primary comparison: each class's full set minus its HAR + RM + FOMC (RULE.txt 5 (P))
\providecommand{\rhemDShRidgeFullVsRmFomc}{}\renewcommand{\rhemDShRidgeFullVsRmFomc}{0.06}  % ridge: full set minus HAR + RM + FOMC, Delta SR, mid
\providecommand{\rhemPDShRidgeFullVsRmFomc}{}\renewcommand{\rhemPDShRidgeFullVsRmFomc}{0.920}  % ... its Ledoit-Wolf p
\providecommand{\rhemDShXRidgeFullVsRmFomc}{}\renewcommand{\rhemDShXRidgeFullVsRmFomc}{0.08}  % ... Delta SR at bid/ask
\providecommand{\rhemPDShXRidgeFullVsRmFomc}{}\renewcommand{\rhemPDShXRidgeFullVsRmFomc}{0.902}  % ... its Ledoit-Wolf p
\providecommand{\rhemTRidgeFullVsRmFomc}{}\renewcommand{\rhemTRidgeFullVsRmFomc}{0.10}  % ... fixed-b t of the mean mid return differential
\providecommand{\rhemPTRidgeFullVsRmFomc}{}\renewcommand{\rhemPTRidgeFullVsRmFomc}{0.927}  % ... its fixed-b p
\providecommand{\rhemDShLassoFullVsRmFomc}{}\renewcommand{\rhemDShLassoFullVsRmFomc}{0.21}  % lasso: full set minus HAR + RM + FOMC, Delta SR, mid
\providecommand{\rhemPDShLassoFullVsRmFomc}{}\renewcommand{\rhemPDShLassoFullVsRmFomc}{0.709}  % ... its Ledoit-Wolf p
\providecommand{\rhemDShXLassoFullVsRmFomc}{}\renewcommand{\rhemDShXLassoFullVsRmFomc}{0.23}  % ... Delta SR at bid/ask
\providecommand{\rhemPDShXLassoFullVsRmFomc}{}\renewcommand{\rhemPDShXLassoFullVsRmFomc}{0.696}  % ... its Ledoit-Wolf p
\providecommand{\rhemTLassoFullVsRmFomc}{}\renewcommand{\rhemTLassoFullVsRmFomc}{0.39}  % ... fixed-b t of the mean mid return differential
\providecommand{\rhemPTLassoFullVsRmFomc}{}\renewcommand{\rhemPTLassoFullVsRmFomc}{0.709}  % ... its fixed-b p
\providecommand{\rhemDShEnetFullVsRmFomc}{}\renewcommand{\rhemDShEnetFullVsRmFomc}{\ensuremath{-}0.18}  % enet: full set minus HAR + RM + FOMC, Delta SR, mid
\providecommand{\rhemPDShEnetFullVsRmFomc}{}\renewcommand{\rhemPDShEnetFullVsRmFomc}{0.541}  % ... its Ledoit-Wolf p
\providecommand{\rhemDShXEnetFullVsRmFomc}{}\renewcommand{\rhemDShXEnetFullVsRmFomc}{\ensuremath{-}0.18}  % ... Delta SR at bid/ask
\providecommand{\rhemPDShXEnetFullVsRmFomc}{}\renewcommand{\rhemPDShXEnetFullVsRmFomc}{0.536}  % ... its Ledoit-Wolf p
\providecommand{\rhemTEnetFullVsRmFomc}{}\renewcommand{\rhemTEnetFullVsRmFomc}{\ensuremath{-}0.56}  % ... fixed-b t of the mean mid return differential
\providecommand{\rhemPTEnetFullVsRmFomc}{}\renewcommand{\rhemPTEnetFullVsRmFomc}{0.588}  % ... its fixed-b p
\providecommand{\rhemDShLgbmFullVsRmFomc}{}\renewcommand{\rhemDShLgbmFullVsRmFomc}{\ensuremath{-}0.15}  % lgbm: full set minus HAR + RM + FOMC, Delta SR, mid
\providecommand{\rhemPDShLgbmFullVsRmFomc}{}\renewcommand{\rhemPDShLgbmFullVsRmFomc}{0.568}  % ... its Ledoit-Wolf p
\providecommand{\rhemDShXLgbmFullVsRmFomc}{}\renewcommand{\rhemDShXLgbmFullVsRmFomc}{\ensuremath{-}0.15}  % ... Delta SR at bid/ask
\providecommand{\rhemPDShXLgbmFullVsRmFomc}{}\renewcommand{\rhemPDShXLgbmFullVsRmFomc}{0.562}  % ... its Ledoit-Wolf p
\providecommand{\rhemTLgbmFullVsRmFomc}{}\renewcommand{\rhemTLgbmFullVsRmFomc}{\ensuremath{-}0.58}  % ... fixed-b t of the mean mid return differential
\providecommand{\rhemPTLgbmFullVsRmFomc}{}\renewcommand{\rhemPTLgbmFullVsRmFomc}{0.581}  % ... its fixed-b p
\providecommand{\rhemDShXgbFullVsRmFomc}{}\renewcommand{\rhemDShXgbFullVsRmFomc}{0.09}  % xgb: full set minus HAR + RM + FOMC, Delta SR, mid
\providecommand{\rhemPDShXgbFullVsRmFomc}{}\renewcommand{\rhemPDShXgbFullVsRmFomc}{0.686}  % ... its Ledoit-Wolf p
\providecommand{\rhemDShXXgbFullVsRmFomc}{}\renewcommand{\rhemDShXXgbFullVsRmFomc}{0.09}  % ... Delta SR at bid/ask
\providecommand{\rhemPDShXXgbFullVsRmFomc}{}\renewcommand{\rhemPDShXXgbFullVsRmFomc}{0.691}  % ... its Ledoit-Wolf p
\providecommand{\rhemTXgbFullVsRmFomc}{}\renewcommand{\rhemTXgbFullVsRmFomc}{0.44}  % ... fixed-b t of the mean mid return differential
\providecommand{\rhemPTXgbFullVsRmFomc}{}\renewcommand{\rhemPTXgbFullVsRmFomc}{0.670}  % ... its fixed-b p
\providecommand{\rhemDShRfFullVsRmFomc}{}\renewcommand{\rhemDShRfFullVsRmFomc}{0.38}  % rf: full set minus HAR + RM + FOMC, Delta SR, mid
\providecommand{\rhemPDShRfFullVsRmFomc}{}\renewcommand{\rhemPDShRfFullVsRmFomc}{0.254}  % ... its Ledoit-Wolf p
\providecommand{\rhemDShXRfFullVsRmFomc}{}\renewcommand{\rhemDShXRfFullVsRmFomc}{0.37}  % ... Delta SR at bid/ask
\providecommand{\rhemPDShXRfFullVsRmFomc}{}\renewcommand{\rhemPDShXRfFullVsRmFomc}{0.261}  % ... its Ledoit-Wolf p
\providecommand{\rhemTRfFullVsRmFomc}{}\renewcommand{\rhemTRfFullVsRmFomc}{1.26}  % ... fixed-b t of the mean mid return differential
\providecommand{\rhemPTRfFullVsRmFomc}{}\renewcommand{\rhemPTRfFullVsRmFomc}{0.229}  % ... its fixed-b p
% ---- over the six classes (Ols excluded); 5 % level; Sig = Delta SR > 0 with p < 0.05, SigNeg = < 0
\providecommand{\rhemRmFomcNCls}{}\renewcommand{\rhemRmFomcNCls}{6}  % model classes with a full-set and a HAR + RM + FOMC forecast
\providecommand{\rhemRmFomcNClsWord}{}\renewcommand{\rhemRmFomcNClsWord}{six}  % ... in words
\providecommand{\rhemFullVsRmFomcMin}{}\renewcommand{\rhemFullVsRmFomcMin}{\ensuremath{-}0.18}  % FullVsRmFomc: Delta SR mid, min over the six
\providecommand{\rhemFullVsRmFomcMax}{}\renewcommand{\rhemFullVsRmFomcMax}{0.38}  % ... max
\providecommand{\rhemFullVsRmFomcPMin}{}\renewcommand{\rhemFullVsRmFomcPMin}{0.254}  % ... smallest Ledoit-Wolf p
\providecommand{\rhemFullVsRmFomcPMax}{}\renewcommand{\rhemFullVsRmFomcPMax}{0.920}  % ... largest Ledoit-Wolf p
\providecommand{\rhemFullVsRmFomcPos}{}\renewcommand{\rhemFullVsRmFomcPos}{4}  % ... Delta SR > 0, mid
\providecommand{\rhemFullVsRmFomcPosWord}{}\renewcommand{\rhemFullVsRmFomcPosWord}{four}  % ... in words
\providecommand{\rhemFullVsRmFomcNeg}{}\renewcommand{\rhemFullVsRmFomcNeg}{2}  % ... Delta SR < 0, mid
\providecommand{\rhemFullVsRmFomcNegWord}{}\renewcommand{\rhemFullVsRmFomcNegWord}{two}  % ... in words
\providecommand{\rhemFullVsRmFomcSig}{}\renewcommand{\rhemFullVsRmFomcSig}{0}  % ... significantly positive, mid
\providecommand{\rhemFullVsRmFomcSigWord}{}\renewcommand{\rhemFullVsRmFomcSigWord}{zero}  % ... in words
\providecommand{\rhemFullVsRmFomcSigNeg}{}\renewcommand{\rhemFullVsRmFomcSigNeg}{0}  % ... significantly negative, mid
\providecommand{\rhemFullVsRmFomcSigNegWord}{}\renewcommand{\rhemFullVsRmFomcSigNegWord}{zero}  % ... in words
\providecommand{\rhemFullVsRmFomcPosX}{}\renewcommand{\rhemFullVsRmFomcPosX}{4}  % ... Delta SR > 0 at bid/ask
\providecommand{\rhemFullVsRmFomcPosXWord}{}\renewcommand{\rhemFullVsRmFomcPosXWord}{four}  % ... in words
\providecommand{\rhemFullVsRmFomcSigX}{}\renewcommand{\rhemFullVsRmFomcSigX}{0}  % ... significantly positive, bid/ask
\providecommand{\rhemFullVsRmFomcSigXWord}{}\renewcommand{\rhemFullVsRmFomcSigXWord}{zero}  % ... in words
\providecommand{\rhemRmFomcVsHarMin}{}\renewcommand{\rhemRmFomcVsHarMin}{0.47}  % RmFomcVsHar: Delta SR mid, min over the six
\providecommand{\rhemRmFomcVsHarMax}{}\renewcommand{\rhemRmFomcVsHarMax}{1.18}  % ... max
\providecommand{\rhemRmFomcVsHarPMin}{}\renewcommand{\rhemRmFomcVsHarPMin}{0.014}  % ... smallest Ledoit-Wolf p
\providecommand{\rhemRmFomcVsHarPMax}{}\renewcommand{\rhemRmFomcVsHarPMax}{0.434}  % ... largest Ledoit-Wolf p
\providecommand{\rhemRmFomcVsHarPos}{}\renewcommand{\rhemRmFomcVsHarPos}{6}  % ... Delta SR > 0, mid
\providecommand{\rhemRmFomcVsHarPosWord}{}\renewcommand{\rhemRmFomcVsHarPosWord}{six}  % ... in words
\providecommand{\rhemRmFomcVsHarNeg}{}\renewcommand{\rhemRmFomcVsHarNeg}{0}  % ... Delta SR < 0, mid
\providecommand{\rhemRmFomcVsHarNegWord}{}\renewcommand{\rhemRmFomcVsHarNegWord}{zero}  % ... in words
\providecommand{\rhemRmFomcVsHarSig}{}\renewcommand{\rhemRmFomcVsHarSig}{2}  % ... significantly positive, mid
\providecommand{\rhemRmFomcVsHarSigWord}{}\renewcommand{\rhemRmFomcVsHarSigWord}{two}  % ... in words
\providecommand{\rhemRmFomcVsHarSigNeg}{}\renewcommand{\rhemRmFomcVsHarSigNeg}{0}  % ... significantly negative, mid
\providecommand{\rhemRmFomcVsHarSigNegWord}{}\renewcommand{\rhemRmFomcVsHarSigNegWord}{zero}  % ... in words
\providecommand{\rhemRmFomcVsHarPosX}{}\renewcommand{\rhemRmFomcVsHarPosX}{6}  % ... Delta SR > 0 at bid/ask
\providecommand{\rhemRmFomcVsHarPosXWord}{}\renewcommand{\rhemRmFomcVsHarPosXWord}{six}  % ... in words
\providecommand{\rhemRmFomcVsHarSigX}{}\renewcommand{\rhemRmFomcVsHarSigX}{2}  % ... significantly positive, bid/ask
\providecommand{\rhemRmFomcVsHarSigXWord}{}\renewcommand{\rhemRmFomcVsHarSigXWord}{two}  % ... in words
\providecommand{\rhemRmFomcVsOlsMin}{}\renewcommand{\rhemRmFomcVsOlsMin}{0.50}  % RmFomcVsOls: Delta SR mid, min over the six
\providecommand{\rhemRmFomcVsOlsMax}{}\renewcommand{\rhemRmFomcVsOlsMax}{1.08}  % ... max
\providecommand{\rhemRmFomcVsOlsPMin}{}\renewcommand{\rhemRmFomcVsOlsPMin}{0.014}  % ... smallest Ledoit-Wolf p
\providecommand{\rhemRmFomcVsOlsPMax}{}\renewcommand{\rhemRmFomcVsOlsPMax}{0.434}  % ... largest Ledoit-Wolf p
\providecommand{\rhemRmFomcVsOlsPos}{}\renewcommand{\rhemRmFomcVsOlsPos}{6}  % ... Delta SR > 0, mid
\providecommand{\rhemRmFomcVsOlsPosWord}{}\renewcommand{\rhemRmFomcVsOlsPosWord}{six}  % ... in words
\providecommand{\rhemRmFomcVsOlsNeg}{}\renewcommand{\rhemRmFomcVsOlsNeg}{0}  % ... Delta SR < 0, mid
\providecommand{\rhemRmFomcVsOlsNegWord}{}\renewcommand{\rhemRmFomcVsOlsNegWord}{zero}  % ... in words
\providecommand{\rhemRmFomcVsOlsSig}{}\renewcommand{\rhemRmFomcVsOlsSig}{1}  % ... significantly positive, mid
\providecommand{\rhemRmFomcVsOlsSigWord}{}\renewcommand{\rhemRmFomcVsOlsSigWord}{one}  % ... in words
\providecommand{\rhemRmFomcVsOlsSigNeg}{}\renewcommand{\rhemRmFomcVsOlsSigNeg}{0}  % ... significantly negative, mid
\providecommand{\rhemRmFomcVsOlsSigNegWord}{}\renewcommand{\rhemRmFomcVsOlsSigNegWord}{zero}  % ... in words
\providecommand{\rhemRmFomcVsOlsPosX}{}\renewcommand{\rhemRmFomcVsOlsPosX}{6}  % ... Delta SR > 0 at bid/ask
\providecommand{\rhemRmFomcVsOlsPosXWord}{}\renewcommand{\rhemRmFomcVsOlsPosXWord}{six}  % ... in words
\providecommand{\rhemRmFomcVsOlsSigX}{}\renewcommand{\rhemRmFomcVsOlsSigX}{1}  % ... significantly positive, bid/ask
\providecommand{\rhemRmFomcVsOlsSigXWord}{}\renewcommand{\rhemRmFomcVsOlsSigXWord}{one}  % ... in words
\providecommand{\rhemRmFomcVsShortMin}{}\renewcommand{\rhemRmFomcVsShortMin}{1.13}  % RmFomcVsShort: Delta SR mid, min over the six
\providecommand{\rhemRmFomcVsShortMax}{}\renewcommand{\rhemRmFomcVsShortMax}{1.71}  % ... max
\providecommand{\rhemRmFomcVsShortPMin}{}\renewcommand{\rhemRmFomcVsShortPMin}{0.013}  % ... smallest Ledoit-Wolf p
\providecommand{\rhemRmFomcVsShortPMax}{}\renewcommand{\rhemRmFomcVsShortPMax}{0.100}  % ... largest Ledoit-Wolf p
\providecommand{\rhemRmFomcVsShortPos}{}\renewcommand{\rhemRmFomcVsShortPos}{6}  % ... Delta SR > 0, mid
\providecommand{\rhemRmFomcVsShortPosWord}{}\renewcommand{\rhemRmFomcVsShortPosWord}{six}  % ... in words
\providecommand{\rhemRmFomcVsShortNeg}{}\renewcommand{\rhemRmFomcVsShortNeg}{0}  % ... Delta SR < 0, mid
\providecommand{\rhemRmFomcVsShortNegWord}{}\renewcommand{\rhemRmFomcVsShortNegWord}{zero}  % ... in words
\providecommand{\rhemRmFomcVsShortSig}{}\renewcommand{\rhemRmFomcVsShortSig}{3}  % ... significantly positive, mid
\providecommand{\rhemRmFomcVsShortSigWord}{}\renewcommand{\rhemRmFomcVsShortSigWord}{three}  % ... in words
\providecommand{\rhemRmFomcVsShortSigNeg}{}\renewcommand{\rhemRmFomcVsShortSigNeg}{0}  % ... significantly negative, mid
\providecommand{\rhemRmFomcVsShortSigNegWord}{}\renewcommand{\rhemRmFomcVsShortSigNegWord}{zero}  % ... in words
\providecommand{\rhemRmFomcVsShortPosX}{}\renewcommand{\rhemRmFomcVsShortPosX}{6}  % ... Delta SR > 0 at bid/ask
\providecommand{\rhemRmFomcVsShortPosXWord}{}\renewcommand{\rhemRmFomcVsShortPosXWord}{six}  % ... in words
\providecommand{\rhemRmFomcVsShortSigX}{}\renewcommand{\rhemRmFomcVsShortSigX}{3}  % ... significantly positive, bid/ask
\providecommand{\rhemRmFomcVsShortSigXWord}{}\renewcommand{\rhemRmFomcVsShortSigXWord}{three}  % ... in words
\providecommand{\rhemRmFomcShMin}{}\renewcommand{\rhemRmFomcShMin}{1.46}  % SR mid of the six HAR + RM + FOMC forecasts: min
\providecommand{\rhemRmFomcShMax}{}\renewcommand{\rhemRmFomcShMax}{2.03}  % ... max
\providecommand{\rhemRmFomcShXMin}{}\renewcommand{\rhemRmFomcShXMin}{1.10}  % SR at bid/ask: min
\providecommand{\rhemRmFomcShXMax}{}\renewcommand{\rhemRmFomcShXMax}{1.68}  % ... max
\fi
% 2026-10-26 (coordinator): parked stop-gap, the macro is now generated -- \providecommand{\rhPending}{\textit{pending}}
% 2026-10-26 (coordinator): parked stop-gap, the macro is now generated -- \providecommand{\rhciFirst}{\rhPending}  % first forecast of the rebuilt CI-HAR (14 June 2006)
% 2026-10-26 (coordinator): parked stop-gap, the macro is now generated -- \providecommand{\rhciQOldClose}{\rhPending}  % QLIKE, 16:00 bar, T_close, the CI-HAR forecasts before the rebuild
% 2026-10-26 (coordinator): parked stop-gap, the macro is now generated -- \providecommand{\rhciQNewClose}{\rhPending}  % ... the rebuilt CI-HAR
% 2026-10-26 (coordinator): parked stop-gap, the macro is now generated -- \providecommand{\rhciDmOldNewClose}{\rhPending}  % DM, rebuilt minus earlier, daily QLIKE on T_close
% 2026-10-26 (coordinator): parked stop-gap, the macro is now generated -- \providecommand{\rhciPOldNewClose}{\rhPending}  % ... its p statement (\ensuremath{p = ...})
% 2026-10-26 (coordinator): parked stop-gap, the macro is now generated -- \providecommand{\rhciFlipDays}{\rhPending}  % post-warmup option days on which the CI-HAR sign position changes

% Alias: the label of the parked Table tab:rh_oos_last30 (old closing-interval model set) points at Table
% tab:rh_train_arms, which holds the models trained on the closing half hour (CLOSE, Panels A, C-F). The aux entry is
% copied from the previous run, so the alias resolves one LaTeX run after tab:rh_train_arms does. Remove it once no
% generated note cites tab:rh_oos_last30 (the close-accuracy table note does, 2026-10-26).
\makeatletter
\expandafter\ifx\csname r@tab:rh_train_arms\endcsname\relax\else
  \immediate\write\@auxout{\string\newlabel{tab:rh_oos_last30}{%
    \unexpanded\expandafter\expandafter\expandafter{\csname r@tab:rh_train_arms\endcsname}}}%
\fi
\makeatother

\clearpage
\section{Closing-interval models}
\label{app:rh_closing}
\begin{itemize}
\item \textbf{Scope.} The models of the main text are estimated on every regular-hours bar and forecast every bar.
  This section collects the models estimated on the closing half hour alone ($\mathcal I = \{c\}$, one observation
  a trading day): the closing-interval HAR benchmark, rebuilt under the current data chain, and the pooled models of
  Table~\ref{tab:rh_oos} re-estimated on the closing half hour (the CLOSE sample). It also reports the accuracy of
  the pooled models' closing half hour on both evaluation spans and the multiple-comparison procedures, whose sets
  of the closing half hour include the closing-interval HAR benchmark.
\item \textbf{Not reported.} The closing-interval linear models and tree ensembles of earlier drafts, a separate
  model set fitted under the earlier data rules with its own window, predictor set and tuning, are not reported;
  the CLOSE sample replaces them.
\end{itemize}

\subsection{The closing-interval models}
\label{app:rh_closing_sample}
\label{app:rh_closing_bench}
\label{app:rh_closing_predictors}
\begin{itemize}
% Moved from the parked subsection "Closing-interval HAR benchmark and estimation" (10_appendix.tex), lightly edited.
\item \textbf{The closing-interval HAR benchmark.} A HAR-type model in the spirit of \citet{Corsi2009}, not Corsi's
  daily HAR: OLS on the HAR components of the target at the horizons $\mathcal K$ and \rhdmNCalendar{} calendar
  columns (\cmPNewBase{} columns in all), estimated on closing-interval observations only, over a rolling window of
  $W = \cmPtFitWin$ trading days, and re-estimated every trading day. It has no exogenous series and no FOMC column.
  It is not the paper's benchmark, which is the pooled OLS on the HAR benchmark (Section~\ref{sec:rh_methods}).
\item \textbf{Rebuilt under the current data chain.} The benchmark's forecasts in earlier drafts were built under the
  data rules in force before the data fixes of Section~\ref{sec:rh_data}. They are rebuilt with the same model and
  code, every data rule at its current setting.
  \begin{itemize}
  \item Three rules change the benchmark's inputs. The diurnal factor of the target averages regular in-slot bars
    only, so the bars after the 13:00 close of an early-close session leave it, which moves the target and the HAR
    components for $n_{\mathrm{d}}$ trading days after each early close (Appendix~\ref{sec:rh_exp_samples}). The
    16:00 rows of the early-close sessions leave the design, as they leave the pooled design. The burn-in of the HAR
    components is dropped on the half-hour grid rather than on the series of closes, so the benchmark's estimation
    sample starts earlier and its forecasts cover $\mathcal T_{\mathrm{oos}}$ (first forecast \rhciFirst).
  \item The other data rules act on exogenous series only and do not reach a design without them.
  \item With the earlier rules switched back on, the code reproduces the earlier forecasts up to floating-point
    rounding.
  \item On the closing half hour of $\mathcal T_{\mathrm{close}}$ the rebuilt benchmark's average QLIKE loss is
    \rhciQNewClose, against \rhciQOldClose{} before the rebuild ($\DM = \rhciDmOldNewClose$,
    \rhciPOldNewClose), and its sign strategy takes the other position on \rhciFlipDays{} of the \rheSzDays{}
    post-warmup days.
  \end{itemize}
\item \textbf{Its strategy.} The sign strategy on the benchmark's forecast (Section~\ref{sec:rh_econ_strategy}) has
  $\SR = \rhemShCiHar$ at mid-quotes and \rhemShXCiHar{} at quoted bid and ask prices, against the short straddle
  $\Delta\SR = \rhemDShCiHarVsShort$ (bootstrap $p = \rhemPDShCiHarVsShort$). The full-set strategies of the pooled
  models have the higher Sharpe ratio in \rhemFullVsCiHarPos{} of the \rhorAfPairs{} model classes
  ($\Delta\SR = \rhemFullVsCiHarMin$ to $\rhemFullVsCiHarMax$); \rhemFullVsCiHarSig{} of these differences are
  significant at the 5\% level under the bootstrap test (the smallest $p$, \rhemFullVsCiHarPMin).
\item \textbf{Models trained on the closing half hour (CLOSE).} A pre-registered comparison of training samples
  (Appendix~\ref{app:rh_robust_train}) re-estimated every pooled model of Table~\ref{tab:rh_oos} on the closing half
  hour alone.
  \begin{itemize}
  \item Sample: the 16:00 rows of the main design, over a rolling window of \rhtaRowsCl{} closes, refit at every
    close, on the design before the FOMC fixes of Appendix~\ref{app:rh_dp_fomc}.
  \item Predictors: those of the main design at the 16:00 bar (Section~\ref{sec:rh_predictors}), computed on the full
    half-hour grid.
  \item Tuning: the frozen tree configurations of the main design, and the rule that chooses the linear penalties
    applied to the sample's own pre-sample rows (validation year \rhtaValYear). A natively built CLOSE design
    (CLOSE-nat) checks the linear models.
  \item Scoring: the 16:00 bar of the \rhtaNdays{} trading days from \rhtaStart{} to \rhtaEnd, the days of
    Table~\ref{tab:rh_oos}.
  \end{itemize}
\end{itemize}

% The CLOSE results of drafts/2026-10-26_train_arms.tex (Panels A, C, D, E and F of tab:rh_train_arms, input in
% rehaul/A_robustness.tex), lightly edited. The labels of the parked subsection "Training on the last 30 minutes" are
% kept on it (sec:rh_econ_last30 is cited from 08_feature_importance.tex and the close-accuracy table).
\subsection{Training on the closing half hour}
\label{sec:rh_econ_last30}
\label{sec:rh_results_pooled}
\begin{itemize}
\item \textbf{Against the main design.} Trained on the close alone, the full-set forecasts of the penalized linear
  models have a lower 16:00 loss than those of the main design ($\Delta\% = \rhtaLinClVsRthClosePctMin$ to
  $\rhtaLinClVsRthClosePctMax$, each with $p < 0.001$); those of the tree ensembles do not differ significantly
  ($\Delta\% = \rhtaTreeClVsRthClosePctMin$ to $\rhtaTreeClVsRthClosePctMax$; the smallest $p$,
  \rhtaTreeClVsRthClosePMin; Table~\ref{tab:rh_train_arms}, Panel~C).
  \begin{itemize}
  \item For the trees, training on the close alone raises the 16:00 loss of the HAR benchmark (for LightGBM from
    \rhtaQHarLgbmRthClose{} to \rhtaQHarLgbmClClose).
  \end{itemize}
\item \textbf{Full predictor set against the HAR benchmark.} Within CLOSE the full predictor set lowers the 16:00
  loss in all six classes (Panel~A).
  \begin{itemize}
  \item $\Delta\% = \rhtaLinPctClCloseMin$ to $\rhtaLinPctClCloseMax$ for the penalized linear models and
    $\rhtaTreePctClCloseMin$ to $\rhtaTreePctClCloseMax$ for the tree ensembles, each with $p < 0.001$; CLOSE-nat
    gives the same result for the linear models ($\Delta\% = \rhtaLinPctClNatCloseMin$ to
    $\rhtaLinPctClNatCloseMax$, each with $p < 0.001$).
  \item In the main design the same comparison gives $\Delta\% = \rhtaPctRthCloseMin$ to $\rhtaPctRthCloseMax$,
    significant at the 5\% level for XGBoost only (\rhtaPXgbRthClose).
  \item A full-versus-HAR change compares two forecasts trained on the same sample, so it also moves with the HAR
    benchmark: for the trees, the larger change under CLOSE reflects their weaker HAR benchmark there (above).
  \end{itemize}
\item \textbf{The window.} A follow-up run, specified in its own written rule before it was estimated,
  re-estimated CLOSE on windows of \rhtaWsMin{} to \rhtaWsMax{} closes (Panel~F). The full predictor set has the
  lower 16:00 loss in all \rhtaWsNCells{} window and class pairs ($\Delta\% = \rhtaWsLinPctMin$ to
  $\rhtaWsLinPctMax$ for the penalized linear models, $\rhtaWsTreePctMin$ to $\rhtaWsTreePctMax$ for the tree
  ensembles), significantly at the 5\% level in \rhtaWsNSig{} of them and in every class from \rhtaWsSigFrom{}
  closes on.
\item \textbf{Equal numbers of rows.} CLOSE, with \rhtaRowsCl{} rows, has a lower 16:00 loss than RTH13-3k, the
  main sample cut to the same number of rows, in all six classes ($\Delta\% = \rhtaClVsRthThreeKClosePctMin$ to
  $\rhtaClVsRthThreeKClosePctMax$; the largest $p$, \rhtaClVsRthThreeKClosePMax; Panel~D). Its window spans about
  \rhtaSpanClSmall{} years against about \rhtaSpanRthSmall{} trading days, so this comparison does not separate the
  choice of intervals from the span of the window.
\item \textbf{The trade} (Panel~E). The full-set sign strategies trained on CLOSE have $\SR$ from \rhtaShFullClMin{}
  to \rhtaShFullClMax, against \rhtaShFullRthMin{} to \rhtaShFullRthMax{} in the main design.
  \begin{itemize}
  \item None of the \rhtaNDShTests{} differences in $\SR$ between the full-set strategies of ALL48 or CLOSE and
    those of the main design is significant at the 5\% level (the smallest $p$, \rhtaDShPMin).
  \item For the tree ensembles, the full-set strategy's gain over the HAR benchmark's is significant at the 5\% level
    in \rhtaNTreeShSigCl{} of the three classes under CLOSE, where the trees' HAR benchmark strategies have $\SR$
    from \rhtaShTreeHarClMin{} to \rhtaShTreeHarClMax.
  \end{itemize}
\item \textbf{Caveats.} The tree configurations were chosen on an earlier design and, by the rule, are not
  re-tuned; with \rhtaRowsCl{} rows their sample-size parameters may not suit the window
  (Appendix~\ref{app:rh_tuning}).
\item \textbf{Summary.} Training on the closing half hour alone lowers the 16:00 loss of the penalized linear models
  but not of the tree ensembles; within CLOSE the full predictor set beats the HAR benchmark in all six classes; and
  no Sharpe ratio differs significantly from the main design's.
\end{itemize}

% The first item of drafts/2026-10-25_close_accuracy.tex (RESULTS agent: "belongs with tab:rh_close_accuracy in your
% A_closing"), lightly edited. Its comparison with the closing-interval HAR benchmark is in A_results.tex
% (app:rh_results_classes); its sentence on the old closing-interval lasso is dropped (D3). Panel B, continued, of the
% generated table still holds the old closing-interval model set (NUMBERS asked to drop it or replace it by CLOSE).
\subsection{Accuracy of the closing half hour}
\label{app:rh_closing_accuracy}
\begin{itemize}
\item \textbf{The table.} Table~\ref{tab:rh_close_accuracy} reports the 15:30 forecast of the closing half hour from
  the pooled models of Table~\ref{tab:rh_oos} on the \rhcaNdaysSpan{} trading days of $\mathcal T_{\mathrm{oos}}$,
  from \rhcaStartSpan{} to \rhcaEndSpan{} (Panel~A), and on $\mathcal T_{\mathrm{close}}$ (Panel~B).
\item \textbf{On $\mathcal T_{\mathrm{oos}}$.} The full predictor set lowers the average QLIKE loss of the closing
  half hour in \rhcaNLowerCloseSpan{} of the \rhorAfPairs{} model classes.
  \begin{itemize}
  % 2026-10-26 (coordinator, D13): on the FOMC-patched design the elastic net's change is positive (NUMBERS report).
  % Was: "... for the three tree ensembles (...) and the elastic net ($\Delta\% = \rhcaPctEnetCloseSpan$); ..." and
  % "It raises the loss of ridge regression (...) and the lasso (...), neither significantly."
  \item It lowers it for the three tree ensembles ($\Delta\% = \rhcaTreePctCloseSpanMin$ to
    $\rhcaTreePctCloseSpanMax$); the reduction is statistically significant at the 5\% level for XGBoost only
    ($\Delta\% = \rhcaPctXgbCloseSpan$, $\DM = \rhcaDmXgbCloseSpan$, \rhcaPXgbCloseSpan).
  \item It raises the loss of ridge regression ($\Delta\% = \rhcaPctRidgeCloseSpan$, $\DM = \rhcaDmRidgeCloseSpan$),
    the lasso ($\Delta\% = \rhcaPctLassoCloseSpan$, $\DM = \rhcaDmLassoCloseSpan$) and the elastic net
    ($\Delta\% = \rhcaPctEnetCloseSpan$), none significantly.
  \end{itemize}
\item \textbf{On $\mathcal T_{\mathrm{close}}$.} The same comparison gives lower losses for the tree ensembles
  ($\Delta\% = \rhcaTreePctCloseOosMin$ to $\rhcaTreePctCloseOosMax$, none significant) and higher losses for the
  penalized linear models ($\Delta\% = \rhcaLinPctCloseOosMin$ to $\rhcaLinPctCloseOosMax$, significant at the 5\%
  level for all three). % 2026-10-26 (coordinator, D13): was "for ridge regression and the lasso" (elastic net now p = 0.036)
\item \textbf{Why the linear models differ between the spans.} Their full-set change has opposite signs on the two
  halves of $\mathcal T_{\mathrm{oos}}$.
  \begin{itemize}
  \item On \rhcaYearsPre{} the full predictor set lowers their closing-half-hour loss ($\Delta\% =
    \rhcaPctRidgeClosePre$, \rhcaPctLassoClosePre{} and \rhcaPctEnetClosePre{} for ridge regression, the lasso and
    the elastic net; $\DM = \rhcaDmRidgeClosePre$, \rhcaDmLassoClosePre{} and \rhcaDmEnetClosePre), and in
    \rhcaYrsDownRidgeClosePre, \rhcaYrsDownLassoClosePre{} and \rhcaYrsDownEnetClosePre{} of the \rhcaNYrsPre{}
    calendar years.
  \item On $\mathcal T_{\mathrm{close}}$ it raises their loss in \rhcaYrsUpRidgeCloseOos{} of the \rhcaNYrsOos{}
    calendar years, most in \rhcaWorstYrRidgeCloseOos{} ($\Delta\% = \rhcaWorstYrPctRidgeCloseOos$ for ridge
    regression).
  \item Over the whole span the two halves offset, and no linear change is significant.
  \end{itemize}
\item \textbf{Against the closing-interval HAR benchmark.} The pooled forecasts are compared with it on
  $\mathcal T_{\mathrm{close}}$ in Appendix~\ref{app:rh_results_classes} (Table~\ref{tab:rh_close_accuracy},
  Panel~B), and jointly in Appendix~\ref{app:rh_closing_mcs}.
\end{itemize}
% AUTO-GENERATED by writeup/make_rehaul_close_accuracy.py -- do not edit; re-run the script.
% tab:rh_close_accuracy: the forecast of the closing half hour (16:00 bar) of the 12 pooled forecasts on Table 1's
% 2014-2024 span and on T_oos, and of the closing-interval models on T_oos; DM on the daily 16:00 loss.
\ifdefined\rhcaNdaysSpan\else% AUTO-GENERATED by writeup/make_rehaul_close_accuracy.py -- do not edit; re-run the script.
% \rhca... macros: closing-half-hour accuracy on both spans (tab:rh_close_accuracy) and the calibration
% diagnostics of the 12 pooled forecasts (tab:rh_calibration, fig:rh_calibration). Names: <Stat><Model><Set>
% <Sample><Span>; Model Ridge Lasso Enet Lgbm Xgb Rf Ols (pooled), Ci<Model> (trained on the closing half
% hour), Bm (closing-interval HAR benchmark); Set Har Full; Sample Close (16:00 bar) Bar (every bar);
% Span Span (2014-2024) Pre (2014-2019) Oos (T_oos, 2020-2024).
% --- spans and constants
\providecommand{\rhcaNdaysSpan}{}\renewcommand{\rhcaNdaysSpan}{2{,}579}  % trading days, 2014-01-02 .. 2024-04-30
\providecommand{\rhcaStartSpan}{}\renewcommand{\rhcaStartSpan}{2 January 2014}  % first day
\providecommand{\rhcaEndSpan}{}\renewcommand{\rhcaEndSpan}{30 April 2024}  % last day
\providecommand{\rhcaYearsSpan}{}\renewcommand{\rhcaYearsSpan}{2014--2024}  % years
\providecommand{\rhcaDMLagSpan}{}\renewcommand{\rhcaDMLagSpan}{8}  % Newey-West lag of the daily DM tests (dm_test automatic)
\providecommand{\rhcaHacLagBarSpan}{}\renewcommand{\rhcaHacLagBarSpan}{104}  % Newey-West lag of the every-bar MZ regressions, bars
\providecommand{\rhcaNdaysPre}{}\renewcommand{\rhcaNdaysPre}{1{,}496}  % trading days, 2014-01-02 .. 2019-12-31
\providecommand{\rhcaStartPre}{}\renewcommand{\rhcaStartPre}{2 January 2014}  % first day
\providecommand{\rhcaEndPre}{}\renewcommand{\rhcaEndPre}{31 December 2019}  % last day
\providecommand{\rhcaYearsPre}{}\renewcommand{\rhcaYearsPre}{2014--2019}  % years
\providecommand{\rhcaDMLagPre}{}\renewcommand{\rhcaDMLagPre}{7}  % Newey-West lag of the daily DM tests (dm_test automatic)
\providecommand{\rhcaHacLagBarPre}{}\renewcommand{\rhcaHacLagBarPre}{91}  % Newey-West lag of the every-bar MZ regressions, bars
\providecommand{\rhcaNdaysOos}{}\renewcommand{\rhcaNdaysOos}{1{,}082}  % trading days, 2020-01-03 .. 2024-04-30
\providecommand{\rhcaStartOos}{}\renewcommand{\rhcaStartOos}{3 January 2020}  % first day
\providecommand{\rhcaEndOos}{}\renewcommand{\rhcaEndOos}{30 April 2024}  % last day
\providecommand{\rhcaYearsOos}{}\renewcommand{\rhcaYearsOos}{2020--2024}  % years
\providecommand{\rhcaDMLagOos}{}\renewcommand{\rhcaDMLagOos}{6}  % Newey-West lag of the daily DM tests (dm_test automatic)
\providecommand{\rhcaHacLagBarOos}{}\renewcommand{\rhcaHacLagBarOos}{78}  % Newey-West lag of the every-bar MZ regressions, bars
\providecommand{\rhcaNbars}{}\renewcommand{\rhcaNbars}{33{,}527}  % scored bars of Table 1, 2014-2024
\providecommand{\rhcaCloseFirst}{}\renewcommand{\rhcaCloseFirst}{14 June 2006}  % first forecast of the closing-interval models
\providecommand{\rhcaMseExp}{}\renewcommand{\rhcaMseExp}{10}  % MSE printed x 10^this
\providecommand{\rhcaUnitExp}{}\renewcommand{\rhcaUnitExp}{6}  % MAE and the levels intercept printed x 10^this
\providecommand{\rhcaCrashStart}{}\renewcommand{\rhcaCrashStart}{24 February 2020}  % first day of the excluded window (make_rehaul_results.CRASH)
\providecommand{\rhcaCrashEnd}{}\renewcommand{\rhcaCrashEnd}{30 April 2020}  % last day of the excluded window
\providecommand{\rhcaTopDays}{}\renewcommand{\rhcaTopDays}{5}  % days in the extreme-day share
% --- closing half hour (16:00 bar): pooled forecasts, Span / Pre / Oos; Q QLIKE, Mse (x 10^10), Mae (x 10^6), Rsq R2_OOS %, Pct / Dm / P QLIKE full vs HAR, MsePct / MseDm / MseP, Gap / GapDm / GapP vs the closing-interval HAR benchmark (Oos only)
\providecommand{\rhcaQOlsHarCloseSpan}{}\renewcommand{\rhcaQOlsHarCloseSpan}{0.1269}
\providecommand{\rhcaMseOlsHarCloseSpan}{}\renewcommand{\rhcaMseOlsHarCloseSpan}{1.97}
\providecommand{\rhcaMaeOlsHarCloseSpan}{}\renewcommand{\rhcaMaeOlsHarCloseSpan}{2.92}
\providecommand{\rhcaRsqOlsHarCloseSpan}{}\renewcommand{\rhcaRsqOlsHarCloseSpan}{77.4}
\providecommand{\rhcaQRidgeHarCloseSpan}{}\renewcommand{\rhcaQRidgeHarCloseSpan}{0.1269}
\providecommand{\rhcaMseRidgeHarCloseSpan}{}\renewcommand{\rhcaMseRidgeHarCloseSpan}{1.97}
\providecommand{\rhcaMaeRidgeHarCloseSpan}{}\renewcommand{\rhcaMaeRidgeHarCloseSpan}{2.92}
\providecommand{\rhcaRsqRidgeHarCloseSpan}{}\renewcommand{\rhcaRsqRidgeHarCloseSpan}{77.4}
\providecommand{\rhcaQRidgeFullCloseSpan}{}\renewcommand{\rhcaQRidgeFullCloseSpan}{0.1291}
\providecommand{\rhcaMseRidgeFullCloseSpan}{}\renewcommand{\rhcaMseRidgeFullCloseSpan}{3.49}
\providecommand{\rhcaMaeRidgeFullCloseSpan}{}\renewcommand{\rhcaMaeRidgeFullCloseSpan}{3.22}
\providecommand{\rhcaRsqRidgeFullCloseSpan}{}\renewcommand{\rhcaRsqRidgeFullCloseSpan}{59.9}
\providecommand{\rhcaPctRidgeCloseSpan}{}\renewcommand{\rhcaPctRidgeCloseSpan}{1.8}  % QLIKE %, full vs HAR of the same class
\providecommand{\rhcaDmRidgeCloseSpan}{}\renewcommand{\rhcaDmRidgeCloseSpan}{0.67}
\providecommand{\rhcaPRidgeCloseSpan}{}\renewcommand{\rhcaPRidgeCloseSpan}{\ensuremath{p = 0.50}}
\providecommand{\rhcaMsePctRidgeCloseSpan}{}\renewcommand{\rhcaMsePctRidgeCloseSpan}{77.5}  % MSE %, full vs HAR
\providecommand{\rhcaMseDmRidgeCloseSpan}{}\renewcommand{\rhcaMseDmRidgeCloseSpan}{1.18}
\providecommand{\rhcaMsePRidgeCloseSpan}{}\renewcommand{\rhcaMsePRidgeCloseSpan}{\ensuremath{p = 0.24}}
\providecommand{\rhcaQLassoHarCloseSpan}{}\renewcommand{\rhcaQLassoHarCloseSpan}{0.1269}
\providecommand{\rhcaMseLassoHarCloseSpan}{}\renewcommand{\rhcaMseLassoHarCloseSpan}{1.97}
\providecommand{\rhcaMaeLassoHarCloseSpan}{}\renewcommand{\rhcaMaeLassoHarCloseSpan}{2.92}
\providecommand{\rhcaRsqLassoHarCloseSpan}{}\renewcommand{\rhcaRsqLassoHarCloseSpan}{77.4}
\providecommand{\rhcaQLassoFullCloseSpan}{}\renewcommand{\rhcaQLassoFullCloseSpan}{0.1322}
\providecommand{\rhcaMseLassoFullCloseSpan}{}\renewcommand{\rhcaMseLassoFullCloseSpan}{3.57}
\providecommand{\rhcaMaeLassoFullCloseSpan}{}\renewcommand{\rhcaMaeLassoFullCloseSpan}{3.24}
\providecommand{\rhcaRsqLassoFullCloseSpan}{}\renewcommand{\rhcaRsqLassoFullCloseSpan}{59.0}
\providecommand{\rhcaPctLassoCloseSpan}{}\renewcommand{\rhcaPctLassoCloseSpan}{4.2}  % QLIKE %, full vs HAR of the same class
\providecommand{\rhcaDmLassoCloseSpan}{}\renewcommand{\rhcaDmLassoCloseSpan}{1.35}
\providecommand{\rhcaPLassoCloseSpan}{}\renewcommand{\rhcaPLassoCloseSpan}{\ensuremath{p = 0.18}}
\providecommand{\rhcaMsePctLassoCloseSpan}{}\renewcommand{\rhcaMsePctLassoCloseSpan}{81.5}  % MSE %, full vs HAR
\providecommand{\rhcaMseDmLassoCloseSpan}{}\renewcommand{\rhcaMseDmLassoCloseSpan}{1.21}
\providecommand{\rhcaMsePLassoCloseSpan}{}\renewcommand{\rhcaMsePLassoCloseSpan}{\ensuremath{p = 0.23}}
\providecommand{\rhcaQEnetHarCloseSpan}{}\renewcommand{\rhcaQEnetHarCloseSpan}{0.1269}
\providecommand{\rhcaMseEnetHarCloseSpan}{}\renewcommand{\rhcaMseEnetHarCloseSpan}{1.97}
\providecommand{\rhcaMaeEnetHarCloseSpan}{}\renewcommand{\rhcaMaeEnetHarCloseSpan}{2.92}
\providecommand{\rhcaRsqEnetHarCloseSpan}{}\renewcommand{\rhcaRsqEnetHarCloseSpan}{77.4}
\providecommand{\rhcaQEnetFullCloseSpan}{}\renewcommand{\rhcaQEnetFullCloseSpan}{0.1276}
\providecommand{\rhcaMseEnetFullCloseSpan}{}\renewcommand{\rhcaMseEnetFullCloseSpan}{3.92}
\providecommand{\rhcaMaeEnetFullCloseSpan}{}\renewcommand{\rhcaMaeEnetFullCloseSpan}{3.28}
\providecommand{\rhcaRsqEnetFullCloseSpan}{}\renewcommand{\rhcaRsqEnetFullCloseSpan}{55.0}
\providecommand{\rhcaPctEnetCloseSpan}{}\renewcommand{\rhcaPctEnetCloseSpan}{0.6}  % QLIKE %, full vs HAR of the same class
\providecommand{\rhcaDmEnetCloseSpan}{}\renewcommand{\rhcaDmEnetCloseSpan}{0.24}
\providecommand{\rhcaPEnetCloseSpan}{}\renewcommand{\rhcaPEnetCloseSpan}{\ensuremath{p = 0.81}}
\providecommand{\rhcaMsePctEnetCloseSpan}{}\renewcommand{\rhcaMsePctEnetCloseSpan}{99.1}  % MSE %, full vs HAR
\providecommand{\rhcaMseDmEnetCloseSpan}{}\renewcommand{\rhcaMseDmEnetCloseSpan}{1.16}
\providecommand{\rhcaMsePEnetCloseSpan}{}\renewcommand{\rhcaMsePEnetCloseSpan}{\ensuremath{p = 0.25}}
\providecommand{\rhcaQLgbmHarCloseSpan}{}\renewcommand{\rhcaQLgbmHarCloseSpan}{0.1235}
\providecommand{\rhcaMseLgbmHarCloseSpan}{}\renewcommand{\rhcaMseLgbmHarCloseSpan}{1.77}
\providecommand{\rhcaMaeLgbmHarCloseSpan}{}\renewcommand{\rhcaMaeLgbmHarCloseSpan}{3.01}
\providecommand{\rhcaRsqLgbmHarCloseSpan}{}\renewcommand{\rhcaRsqLgbmHarCloseSpan}{79.6}
\providecommand{\rhcaQLgbmFullCloseSpan}{}\renewcommand{\rhcaQLgbmFullCloseSpan}{0.1202}
\providecommand{\rhcaMseLgbmFullCloseSpan}{}\renewcommand{\rhcaMseLgbmFullCloseSpan}{2.36}
\providecommand{\rhcaMaeLgbmFullCloseSpan}{}\renewcommand{\rhcaMaeLgbmFullCloseSpan}{3.02}
\providecommand{\rhcaRsqLgbmFullCloseSpan}{}\renewcommand{\rhcaRsqLgbmFullCloseSpan}{72.9}
\providecommand{\rhcaPctLgbmCloseSpan}{}\renewcommand{\rhcaPctLgbmCloseSpan}{\ensuremath{-}2.7}  % QLIKE %, full vs HAR of the same class
\providecommand{\rhcaDmLgbmCloseSpan}{}\renewcommand{\rhcaDmLgbmCloseSpan}{\ensuremath{-}1.76}
\providecommand{\rhcaPLgbmCloseSpan}{}\renewcommand{\rhcaPLgbmCloseSpan}{\ensuremath{p = 0.079}}
\providecommand{\rhcaMsePctLgbmCloseSpan}{}\renewcommand{\rhcaMsePctLgbmCloseSpan}{33.0}  % MSE %, full vs HAR
\providecommand{\rhcaMseDmLgbmCloseSpan}{}\renewcommand{\rhcaMseDmLgbmCloseSpan}{1.10}
\providecommand{\rhcaMsePLgbmCloseSpan}{}\renewcommand{\rhcaMsePLgbmCloseSpan}{\ensuremath{p = 0.27}}
\providecommand{\rhcaQXgbHarCloseSpan}{}\renewcommand{\rhcaQXgbHarCloseSpan}{0.1234}
\providecommand{\rhcaMseXgbHarCloseSpan}{}\renewcommand{\rhcaMseXgbHarCloseSpan}{1.73}
\providecommand{\rhcaMaeXgbHarCloseSpan}{}\renewcommand{\rhcaMaeXgbHarCloseSpan}{2.95}
\providecommand{\rhcaRsqXgbHarCloseSpan}{}\renewcommand{\rhcaRsqXgbHarCloseSpan}{80.1}
\providecommand{\rhcaQXgbFullCloseSpan}{}\renewcommand{\rhcaQXgbFullCloseSpan}{0.1195}
\providecommand{\rhcaMseXgbFullCloseSpan}{}\renewcommand{\rhcaMseXgbFullCloseSpan}{2.19}
\providecommand{\rhcaMaeXgbFullCloseSpan}{}\renewcommand{\rhcaMaeXgbFullCloseSpan}{3.01}
\providecommand{\rhcaRsqXgbFullCloseSpan}{}\renewcommand{\rhcaRsqXgbFullCloseSpan}{74.9}
\providecommand{\rhcaPctXgbCloseSpan}{}\renewcommand{\rhcaPctXgbCloseSpan}{\ensuremath{-}3.1}  % QLIKE %, full vs HAR of the same class
\providecommand{\rhcaDmXgbCloseSpan}{}\renewcommand{\rhcaDmXgbCloseSpan}{\ensuremath{-}2.01}
\providecommand{\rhcaPXgbCloseSpan}{}\renewcommand{\rhcaPXgbCloseSpan}{\ensuremath{p = 0.044}}
\providecommand{\rhcaMsePctXgbCloseSpan}{}\renewcommand{\rhcaMsePctXgbCloseSpan}{26.3}  % MSE %, full vs HAR
\providecommand{\rhcaMseDmXgbCloseSpan}{}\renewcommand{\rhcaMseDmXgbCloseSpan}{1.09}
\providecommand{\rhcaMsePXgbCloseSpan}{}\renewcommand{\rhcaMsePXgbCloseSpan}{\ensuremath{p = 0.28}}
\providecommand{\rhcaQRfHarCloseSpan}{}\renewcommand{\rhcaQRfHarCloseSpan}{0.1266}
\providecommand{\rhcaMseRfHarCloseSpan}{}\renewcommand{\rhcaMseRfHarCloseSpan}{1.63}
\providecommand{\rhcaMaeRfHarCloseSpan}{}\renewcommand{\rhcaMaeRfHarCloseSpan}{2.91}
\providecommand{\rhcaRsqRfHarCloseSpan}{}\renewcommand{\rhcaRsqRfHarCloseSpan}{81.3}
\providecommand{\rhcaQRfFullCloseSpan}{}\renewcommand{\rhcaQRfFullCloseSpan}{0.1246}
\providecommand{\rhcaMseRfFullCloseSpan}{}\renewcommand{\rhcaMseRfFullCloseSpan}{1.95}
\providecommand{\rhcaMaeRfFullCloseSpan}{}\renewcommand{\rhcaMaeRfFullCloseSpan}{2.97}
\providecommand{\rhcaRsqRfFullCloseSpan}{}\renewcommand{\rhcaRsqRfFullCloseSpan}{77.6}
\providecommand{\rhcaPctRfCloseSpan}{}\renewcommand{\rhcaPctRfCloseSpan}{\ensuremath{-}1.6}  % QLIKE %, full vs HAR of the same class
\providecommand{\rhcaDmRfCloseSpan}{}\renewcommand{\rhcaDmRfCloseSpan}{\ensuremath{-}1.03}
\providecommand{\rhcaPRfCloseSpan}{}\renewcommand{\rhcaPRfCloseSpan}{\ensuremath{p = 0.30}}
\providecommand{\rhcaMsePctRfCloseSpan}{}\renewcommand{\rhcaMsePctRfCloseSpan}{19.7}  % MSE %, full vs HAR
\providecommand{\rhcaMseDmRfCloseSpan}{}\renewcommand{\rhcaMseDmRfCloseSpan}{1.11}
\providecommand{\rhcaMsePRfCloseSpan}{}\renewcommand{\rhcaMsePRfCloseSpan}{\ensuremath{p = 0.27}}
\providecommand{\rhcaQOlsHarClosePre}{}\renewcommand{\rhcaQOlsHarClosePre}{0.1365}
\providecommand{\rhcaMseOlsHarClosePre}{}\renewcommand{\rhcaMseOlsHarClosePre}{0.32}
\providecommand{\rhcaMaeOlsHarClosePre}{}\renewcommand{\rhcaMaeOlsHarClosePre}{1.74}
\providecommand{\rhcaRsqOlsHarClosePre}{}\renewcommand{\rhcaRsqOlsHarClosePre}{73.9}
\providecommand{\rhcaQRidgeHarClosePre}{}\renewcommand{\rhcaQRidgeHarClosePre}{0.1365}
\providecommand{\rhcaMseRidgeHarClosePre}{}\renewcommand{\rhcaMseRidgeHarClosePre}{0.32}
\providecommand{\rhcaMaeRidgeHarClosePre}{}\renewcommand{\rhcaMaeRidgeHarClosePre}{1.74}
\providecommand{\rhcaRsqRidgeHarClosePre}{}\renewcommand{\rhcaRsqRidgeHarClosePre}{73.9}
\providecommand{\rhcaQRidgeFullClosePre}{}\renewcommand{\rhcaQRidgeFullClosePre}{0.1309}
\providecommand{\rhcaMseRidgeFullClosePre}{}\renewcommand{\rhcaMseRidgeFullClosePre}{0.28}
\providecommand{\rhcaMaeRidgeFullClosePre}{}\renewcommand{\rhcaMaeRidgeFullClosePre}{1.66}
\providecommand{\rhcaRsqRidgeFullClosePre}{}\renewcommand{\rhcaRsqRidgeFullClosePre}{77.3}
\providecommand{\rhcaPctRidgeClosePre}{}\renewcommand{\rhcaPctRidgeClosePre}{\ensuremath{-}4.1}  % QLIKE %, full vs HAR of the same class
\providecommand{\rhcaDmRidgeClosePre}{}\renewcommand{\rhcaDmRidgeClosePre}{\ensuremath{-}1.42}
\providecommand{\rhcaPRidgeClosePre}{}\renewcommand{\rhcaPRidgeClosePre}{\ensuremath{p = 0.16}}
\providecommand{\rhcaMsePctRidgeClosePre}{}\renewcommand{\rhcaMsePctRidgeClosePre}{\ensuremath{-}13.1}  % MSE %, full vs HAR
\providecommand{\rhcaMseDmRidgeClosePre}{}\renewcommand{\rhcaMseDmRidgeClosePre}{\ensuremath{-}1.79}
\providecommand{\rhcaMsePRidgeClosePre}{}\renewcommand{\rhcaMsePRidgeClosePre}{\ensuremath{p = 0.074}}
\providecommand{\rhcaQLassoHarClosePre}{}\renewcommand{\rhcaQLassoHarClosePre}{0.1365}
\providecommand{\rhcaMseLassoHarClosePre}{}\renewcommand{\rhcaMseLassoHarClosePre}{0.32}
\providecommand{\rhcaMaeLassoHarClosePre}{}\renewcommand{\rhcaMaeLassoHarClosePre}{1.74}
\providecommand{\rhcaRsqLassoHarClosePre}{}\renewcommand{\rhcaRsqLassoHarClosePre}{73.9}
\providecommand{\rhcaQLassoFullClosePre}{}\renewcommand{\rhcaQLassoFullClosePre}{0.1333}
\providecommand{\rhcaMseLassoFullClosePre}{}\renewcommand{\rhcaMseLassoFullClosePre}{0.27}
\providecommand{\rhcaMaeLassoFullClosePre}{}\renewcommand{\rhcaMaeLassoFullClosePre}{1.65}
\providecommand{\rhcaRsqLassoFullClosePre}{}\renewcommand{\rhcaRsqLassoFullClosePre}{78.0}
\providecommand{\rhcaPctLassoClosePre}{}\renewcommand{\rhcaPctLassoClosePre}{\ensuremath{-}2.3}  % QLIKE %, full vs HAR of the same class
\providecommand{\rhcaDmLassoClosePre}{}\renewcommand{\rhcaDmLassoClosePre}{\ensuremath{-}0.67}
\providecommand{\rhcaPLassoClosePre}{}\renewcommand{\rhcaPLassoClosePre}{\ensuremath{p = 0.50}}
\providecommand{\rhcaMsePctLassoClosePre}{}\renewcommand{\rhcaMsePctLassoClosePre}{\ensuremath{-}15.8}  % MSE %, full vs HAR
\providecommand{\rhcaMseDmLassoClosePre}{}\renewcommand{\rhcaMseDmLassoClosePre}{\ensuremath{-}1.78}
\providecommand{\rhcaMsePLassoClosePre}{}\renewcommand{\rhcaMsePLassoClosePre}{\ensuremath{p = 0.075}}
\providecommand{\rhcaQEnetHarClosePre}{}\renewcommand{\rhcaQEnetHarClosePre}{0.1365}
\providecommand{\rhcaMseEnetHarClosePre}{}\renewcommand{\rhcaMseEnetHarClosePre}{0.32}
\providecommand{\rhcaMaeEnetHarClosePre}{}\renewcommand{\rhcaMaeEnetHarClosePre}{1.74}
\providecommand{\rhcaRsqEnetHarClosePre}{}\renewcommand{\rhcaRsqEnetHarClosePre}{73.9}
\providecommand{\rhcaQEnetFullClosePre}{}\renewcommand{\rhcaQEnetFullClosePre}{0.1300}
\providecommand{\rhcaMseEnetFullClosePre}{}\renewcommand{\rhcaMseEnetFullClosePre}{0.28}
\providecommand{\rhcaMaeEnetFullClosePre}{}\renewcommand{\rhcaMaeEnetFullClosePre}{1.64}
\providecommand{\rhcaRsqEnetFullClosePre}{}\renewcommand{\rhcaRsqEnetFullClosePre}{77.6}
\providecommand{\rhcaPctEnetClosePre}{}\renewcommand{\rhcaPctEnetClosePre}{\ensuremath{-}4.8}  % QLIKE %, full vs HAR of the same class
\providecommand{\rhcaDmEnetClosePre}{}\renewcommand{\rhcaDmEnetClosePre}{\ensuremath{-}1.80}
\providecommand{\rhcaPEnetClosePre}{}\renewcommand{\rhcaPEnetClosePre}{\ensuremath{p = 0.071}}
\providecommand{\rhcaMsePctEnetClosePre}{}\renewcommand{\rhcaMsePctEnetClosePre}{\ensuremath{-}14.4}  % MSE %, full vs HAR
\providecommand{\rhcaMseDmEnetClosePre}{}\renewcommand{\rhcaMseDmEnetClosePre}{\ensuremath{-}1.90}
\providecommand{\rhcaMsePEnetClosePre}{}\renewcommand{\rhcaMsePEnetClosePre}{\ensuremath{p = 0.058}}
\providecommand{\rhcaQLgbmHarClosePre}{}\renewcommand{\rhcaQLgbmHarClosePre}{0.1333}
\providecommand{\rhcaMseLgbmHarClosePre}{}\renewcommand{\rhcaMseLgbmHarClosePre}{0.40}
\providecommand{\rhcaMaeLgbmHarClosePre}{}\renewcommand{\rhcaMaeLgbmHarClosePre}{1.84}
\providecommand{\rhcaRsqLgbmHarClosePre}{}\renewcommand{\rhcaRsqLgbmHarClosePre}{68.1}
\providecommand{\rhcaQLgbmFullClosePre}{}\renewcommand{\rhcaQLgbmFullClosePre}{0.1287}
\providecommand{\rhcaMseLgbmFullClosePre}{}\renewcommand{\rhcaMseLgbmFullClosePre}{0.37}
\providecommand{\rhcaMaeLgbmFullClosePre}{}\renewcommand{\rhcaMaeLgbmFullClosePre}{1.74}
\providecommand{\rhcaRsqLgbmFullClosePre}{}\renewcommand{\rhcaRsqLgbmFullClosePre}{70.5}
\providecommand{\rhcaPctLgbmClosePre}{}\renewcommand{\rhcaPctLgbmClosePre}{\ensuremath{-}3.4}  % QLIKE %, full vs HAR of the same class
\providecommand{\rhcaDmLgbmClosePre}{}\renewcommand{\rhcaDmLgbmClosePre}{\ensuremath{-}1.99}
\providecommand{\rhcaPLgbmClosePre}{}\renewcommand{\rhcaPLgbmClosePre}{\ensuremath{p = 0.046}}
\providecommand{\rhcaMsePctLgbmClosePre}{}\renewcommand{\rhcaMsePctLgbmClosePre}{\ensuremath{-}7.5}  % MSE %, full vs HAR
\providecommand{\rhcaMseDmLgbmClosePre}{}\renewcommand{\rhcaMseDmLgbmClosePre}{\ensuremath{-}1.78}
\providecommand{\rhcaMsePLgbmClosePre}{}\renewcommand{\rhcaMsePLgbmClosePre}{\ensuremath{p = 0.076}}
\providecommand{\rhcaQXgbHarClosePre}{}\renewcommand{\rhcaQXgbHarClosePre}{0.1331}
\providecommand{\rhcaMseXgbHarClosePre}{}\renewcommand{\rhcaMseXgbHarClosePre}{0.37}
\providecommand{\rhcaMaeXgbHarClosePre}{}\renewcommand{\rhcaMaeXgbHarClosePre}{1.82}
\providecommand{\rhcaRsqXgbHarClosePre}{}\renewcommand{\rhcaRsqXgbHarClosePre}{70.0}
\providecommand{\rhcaQXgbFullClosePre}{}\renewcommand{\rhcaQXgbFullClosePre}{0.1277}
\providecommand{\rhcaMseXgbFullClosePre}{}\renewcommand{\rhcaMseXgbFullClosePre}{0.33}
\providecommand{\rhcaMaeXgbFullClosePre}{}\renewcommand{\rhcaMaeXgbFullClosePre}{1.74}
\providecommand{\rhcaRsqXgbFullClosePre}{}\renewcommand{\rhcaRsqXgbFullClosePre}{73.0}
\providecommand{\rhcaPctXgbClosePre}{}\renewcommand{\rhcaPctXgbClosePre}{\ensuremath{-}4.0}  % QLIKE %, full vs HAR of the same class
\providecommand{\rhcaDmXgbClosePre}{}\renewcommand{\rhcaDmXgbClosePre}{\ensuremath{-}2.13}
\providecommand{\rhcaPXgbClosePre}{}\renewcommand{\rhcaPXgbClosePre}{\ensuremath{p = 0.033}}
\providecommand{\rhcaMsePctXgbClosePre}{}\renewcommand{\rhcaMsePctXgbClosePre}{\ensuremath{-}10.0}  % MSE %, full vs HAR
\providecommand{\rhcaMseDmXgbClosePre}{}\renewcommand{\rhcaMseDmXgbClosePre}{\ensuremath{-}1.66}
\providecommand{\rhcaMsePXgbClosePre}{}\renewcommand{\rhcaMsePXgbClosePre}{\ensuremath{p = 0.097}}
\providecommand{\rhcaQRfHarClosePre}{}\renewcommand{\rhcaQRfHarClosePre}{0.1364}
\providecommand{\rhcaMseRfHarClosePre}{}\renewcommand{\rhcaMseRfHarClosePre}{0.36}
\providecommand{\rhcaMaeRfHarClosePre}{}\renewcommand{\rhcaMaeRfHarClosePre}{1.81}
\providecommand{\rhcaRsqRfHarClosePre}{}\renewcommand{\rhcaRsqRfHarClosePre}{70.7}
\providecommand{\rhcaQRfFullClosePre}{}\renewcommand{\rhcaQRfFullClosePre}{0.1335}
\providecommand{\rhcaMseRfFullClosePre}{}\renewcommand{\rhcaMseRfFullClosePre}{0.34}
\providecommand{\rhcaMaeRfFullClosePre}{}\renewcommand{\rhcaMaeRfFullClosePre}{1.77}
\providecommand{\rhcaRsqRfFullClosePre}{}\renewcommand{\rhcaRsqRfFullClosePre}{72.3}
\providecommand{\rhcaPctRfClosePre}{}\renewcommand{\rhcaPctRfClosePre}{\ensuremath{-}2.1}  % QLIKE %, full vs HAR of the same class
\providecommand{\rhcaDmRfClosePre}{}\renewcommand{\rhcaDmRfClosePre}{\ensuremath{-}1.14}
\providecommand{\rhcaPRfClosePre}{}\renewcommand{\rhcaPRfClosePre}{\ensuremath{p = 0.25}}
\providecommand{\rhcaMsePctRfClosePre}{}\renewcommand{\rhcaMsePctRfClosePre}{\ensuremath{-}5.3}  % MSE %, full vs HAR
\providecommand{\rhcaMseDmRfClosePre}{}\renewcommand{\rhcaMseDmRfClosePre}{\ensuremath{-}1.39}
\providecommand{\rhcaMsePRfClosePre}{}\renewcommand{\rhcaMsePRfClosePre}{\ensuremath{p = 0.17}}
\providecommand{\rhcaQOlsHarCloseOos}{}\renewcommand{\rhcaQOlsHarCloseOos}{0.1137}
\providecommand{\rhcaMseOlsHarCloseOos}{}\renewcommand{\rhcaMseOlsHarCloseOos}{4.24}
\providecommand{\rhcaMaeOlsHarCloseOos}{}\renewcommand{\rhcaMaeOlsHarCloseOos}{4.55}
\providecommand{\rhcaRsqOlsHarCloseOos}{}\renewcommand{\rhcaRsqOlsHarCloseOos}{77.7}
\providecommand{\rhcaGapOlsHarCloseOos}{}\renewcommand{\rhcaGapOlsHarCloseOos}{10.9}  % QLIKE % vs the closing-interval HAR benchmark
\providecommand{\rhcaGapDmOlsHarCloseOos}{}\renewcommand{\rhcaGapDmOlsHarCloseOos}{2.84}
\providecommand{\rhcaGapPOlsHarCloseOos}{}\renewcommand{\rhcaGapPOlsHarCloseOos}{\ensuremath{p = 0.004}}
\providecommand{\rhcaQRidgeHarCloseOos}{}\renewcommand{\rhcaQRidgeHarCloseOos}{0.1137}
\providecommand{\rhcaMseRidgeHarCloseOos}{}\renewcommand{\rhcaMseRidgeHarCloseOos}{4.24}
\providecommand{\rhcaMaeRidgeHarCloseOos}{}\renewcommand{\rhcaMaeRidgeHarCloseOos}{4.55}
\providecommand{\rhcaRsqRidgeHarCloseOos}{}\renewcommand{\rhcaRsqRidgeHarCloseOos}{77.7}
\providecommand{\rhcaGapRidgeHarCloseOos}{}\renewcommand{\rhcaGapRidgeHarCloseOos}{10.9}  % QLIKE % vs the closing-interval HAR benchmark
\providecommand{\rhcaGapDmRidgeHarCloseOos}{}\renewcommand{\rhcaGapDmRidgeHarCloseOos}{2.84}
\providecommand{\rhcaGapPRidgeHarCloseOos}{}\renewcommand{\rhcaGapPRidgeHarCloseOos}{\ensuremath{p = 0.004}}
\providecommand{\rhcaQRidgeFullCloseOos}{}\renewcommand{\rhcaQRidgeFullCloseOos}{0.1269}
\providecommand{\rhcaMseRidgeFullCloseOos}{}\renewcommand{\rhcaMseRidgeFullCloseOos}{7.94}
\providecommand{\rhcaMaeRidgeFullCloseOos}{}\renewcommand{\rhcaMaeRidgeFullCloseOos}{5.39}
\providecommand{\rhcaRsqRidgeFullCloseOos}{}\renewcommand{\rhcaRsqRidgeFullCloseOos}{58.4}
\providecommand{\rhcaPctRidgeCloseOos}{}\renewcommand{\rhcaPctRidgeCloseOos}{11.6}  % QLIKE %, full vs HAR of the same class
\providecommand{\rhcaDmRidgeCloseOos}{}\renewcommand{\rhcaDmRidgeCloseOos}{2.44}
\providecommand{\rhcaPRidgeCloseOos}{}\renewcommand{\rhcaPRidgeCloseOos}{\ensuremath{p = 0.015}}
\providecommand{\rhcaMsePctRidgeCloseOos}{}\renewcommand{\rhcaMsePctRidgeCloseOos}{87.0}  % MSE %, full vs HAR
\providecommand{\rhcaMseDmRidgeCloseOos}{}\renewcommand{\rhcaMseDmRidgeCloseOos}{1.29}
\providecommand{\rhcaMsePRidgeCloseOos}{}\renewcommand{\rhcaMsePRidgeCloseOos}{\ensuremath{p = 0.20}}
\providecommand{\rhcaGapRidgeFullCloseOos}{}\renewcommand{\rhcaGapRidgeFullCloseOos}{23.8}  % QLIKE % vs the closing-interval HAR benchmark
\providecommand{\rhcaGapDmRidgeFullCloseOos}{}\renewcommand{\rhcaGapDmRidgeFullCloseOos}{3.72}
\providecommand{\rhcaGapPRidgeFullCloseOos}{}\renewcommand{\rhcaGapPRidgeFullCloseOos}{\ensuremath{p < 0.001}}
\providecommand{\rhcaQLassoHarCloseOos}{}\renewcommand{\rhcaQLassoHarCloseOos}{0.1137}
\providecommand{\rhcaMseLassoHarCloseOos}{}\renewcommand{\rhcaMseLassoHarCloseOos}{4.24}
\providecommand{\rhcaMaeLassoHarCloseOos}{}\renewcommand{\rhcaMaeLassoHarCloseOos}{4.55}
\providecommand{\rhcaRsqLassoHarCloseOos}{}\renewcommand{\rhcaRsqLassoHarCloseOos}{77.7}
\providecommand{\rhcaGapLassoHarCloseOos}{}\renewcommand{\rhcaGapLassoHarCloseOos}{10.9}  % QLIKE % vs the closing-interval HAR benchmark
\providecommand{\rhcaGapDmLassoHarCloseOos}{}\renewcommand{\rhcaGapDmLassoHarCloseOos}{2.84}
\providecommand{\rhcaGapPLassoHarCloseOos}{}\renewcommand{\rhcaGapPLassoHarCloseOos}{\ensuremath{p = 0.004}}
\providecommand{\rhcaQLassoFullCloseOos}{}\renewcommand{\rhcaQLassoFullCloseOos}{0.1309}
\providecommand{\rhcaMseLassoFullCloseOos}{}\renewcommand{\rhcaMseLassoFullCloseOos}{8.14}
\providecommand{\rhcaMaeLassoFullCloseOos}{}\renewcommand{\rhcaMaeLassoFullCloseOos}{5.45}
\providecommand{\rhcaRsqLassoFullCloseOos}{}\renewcommand{\rhcaRsqLassoFullCloseOos}{57.3}
\providecommand{\rhcaPctLassoCloseOos}{}\renewcommand{\rhcaPctLassoCloseOos}{15.2}  % QLIKE %, full vs HAR of the same class
\providecommand{\rhcaDmLassoCloseOos}{}\renewcommand{\rhcaDmLassoCloseOos}{2.68}
\providecommand{\rhcaPLassoCloseOos}{}\renewcommand{\rhcaPLassoCloseOos}{\ensuremath{p = 0.007}}
\providecommand{\rhcaMsePctLassoCloseOos}{}\renewcommand{\rhcaMsePctLassoCloseOos}{91.7}  % MSE %, full vs HAR
\providecommand{\rhcaMseDmLassoCloseOos}{}\renewcommand{\rhcaMseDmLassoCloseOos}{1.33}
\providecommand{\rhcaMsePLassoCloseOos}{}\renewcommand{\rhcaMsePLassoCloseOos}{\ensuremath{p = 0.18}}
\providecommand{\rhcaGapLassoFullCloseOos}{}\renewcommand{\rhcaGapLassoFullCloseOos}{27.7}  % QLIKE % vs the closing-interval HAR benchmark
\providecommand{\rhcaGapDmLassoFullCloseOos}{}\renewcommand{\rhcaGapDmLassoFullCloseOos}{3.80}
\providecommand{\rhcaGapPLassoFullCloseOos}{}\renewcommand{\rhcaGapPLassoFullCloseOos}{\ensuremath{p < 0.001}}
\providecommand{\rhcaQEnetHarCloseOos}{}\renewcommand{\rhcaQEnetHarCloseOos}{0.1137}
\providecommand{\rhcaMseEnetHarCloseOos}{}\renewcommand{\rhcaMseEnetHarCloseOos}{4.24}
\providecommand{\rhcaMaeEnetHarCloseOos}{}\renewcommand{\rhcaMaeEnetHarCloseOos}{4.55}
\providecommand{\rhcaRsqEnetHarCloseOos}{}\renewcommand{\rhcaRsqEnetHarCloseOos}{77.7}
\providecommand{\rhcaGapEnetHarCloseOos}{}\renewcommand{\rhcaGapEnetHarCloseOos}{10.9}  % QLIKE % vs the closing-interval HAR benchmark
\providecommand{\rhcaGapDmEnetHarCloseOos}{}\renewcommand{\rhcaGapDmEnetHarCloseOos}{2.84}
\providecommand{\rhcaGapPEnetHarCloseOos}{}\renewcommand{\rhcaGapPEnetHarCloseOos}{\ensuremath{p = 0.004}}
\providecommand{\rhcaQEnetFullCloseOos}{}\renewcommand{\rhcaQEnetFullCloseOos}{0.1245}
\providecommand{\rhcaMseEnetFullCloseOos}{}\renewcommand{\rhcaMseEnetFullCloseOos}{8.96}
\providecommand{\rhcaMaeEnetFullCloseOos}{}\renewcommand{\rhcaMaeEnetFullCloseOos}{5.54}
\providecommand{\rhcaRsqEnetFullCloseOos}{}\renewcommand{\rhcaRsqEnetFullCloseOos}{53.0}
\providecommand{\rhcaPctEnetCloseOos}{}\renewcommand{\rhcaPctEnetCloseOos}{9.5}  % QLIKE %, full vs HAR of the same class
\providecommand{\rhcaDmEnetCloseOos}{}\renewcommand{\rhcaDmEnetCloseOos}{2.09}
\providecommand{\rhcaPEnetCloseOos}{}\renewcommand{\rhcaPEnetCloseOos}{\ensuremath{p = 0.036}}
\providecommand{\rhcaMsePctEnetCloseOos}{}\renewcommand{\rhcaMsePctEnetCloseOos}{111.1}  % MSE %, full vs HAR
\providecommand{\rhcaMseDmEnetCloseOos}{}\renewcommand{\rhcaMseDmEnetCloseOos}{1.25}
\providecommand{\rhcaMsePEnetCloseOos}{}\renewcommand{\rhcaMsePEnetCloseOos}{\ensuremath{p = 0.21}}
\providecommand{\rhcaGapEnetFullCloseOos}{}\renewcommand{\rhcaGapEnetFullCloseOos}{21.4}  % QLIKE % vs the closing-interval HAR benchmark
\providecommand{\rhcaGapDmEnetFullCloseOos}{}\renewcommand{\rhcaGapDmEnetFullCloseOos}{3.48}
\providecommand{\rhcaGapPEnetFullCloseOos}{}\renewcommand{\rhcaGapPEnetFullCloseOos}{\ensuremath{p < 0.001}}
\providecommand{\rhcaQLgbmHarCloseOos}{}\renewcommand{\rhcaQLgbmHarCloseOos}{0.1101}
\providecommand{\rhcaMseLgbmHarCloseOos}{}\renewcommand{\rhcaMseLgbmHarCloseOos}{3.68}
\providecommand{\rhcaMaeLgbmHarCloseOos}{}\renewcommand{\rhcaMaeLgbmHarCloseOos}{4.63}
\providecommand{\rhcaRsqLgbmHarCloseOos}{}\renewcommand{\rhcaRsqLgbmHarCloseOos}{80.7}
\providecommand{\rhcaGapLgbmHarCloseOos}{}\renewcommand{\rhcaGapLgbmHarCloseOos}{7.4}  % QLIKE % vs the closing-interval HAR benchmark
\providecommand{\rhcaGapDmLgbmHarCloseOos}{}\renewcommand{\rhcaGapDmLgbmHarCloseOos}{2.60}
\providecommand{\rhcaGapPLgbmHarCloseOos}{}\renewcommand{\rhcaGapPLgbmHarCloseOos}{\ensuremath{p = 0.009}}
\providecommand{\rhcaQLgbmFullCloseOos}{}\renewcommand{\rhcaQLgbmFullCloseOos}{0.1086}
\providecommand{\rhcaMseLgbmFullCloseOos}{}\renewcommand{\rhcaMseLgbmFullCloseOos}{5.12}
\providecommand{\rhcaMaeLgbmFullCloseOos}{}\renewcommand{\rhcaMaeLgbmFullCloseOos}{4.80}
\providecommand{\rhcaRsqLgbmFullCloseOos}{}\renewcommand{\rhcaRsqLgbmFullCloseOos}{73.1}
\providecommand{\rhcaPctLgbmCloseOos}{}\renewcommand{\rhcaPctLgbmCloseOos}{\ensuremath{-}1.4}  % QLIKE %, full vs HAR of the same class
\providecommand{\rhcaDmLgbmCloseOos}{}\renewcommand{\rhcaDmLgbmCloseOos}{\ensuremath{-}0.49}
\providecommand{\rhcaPLgbmCloseOos}{}\renewcommand{\rhcaPLgbmCloseOos}{\ensuremath{p = 0.62}}
\providecommand{\rhcaMsePctLgbmCloseOos}{}\renewcommand{\rhcaMsePctLgbmCloseOos}{38.9}  % MSE %, full vs HAR
\providecommand{\rhcaMseDmLgbmCloseOos}{}\renewcommand{\rhcaMseDmLgbmCloseOos}{1.22}
\providecommand{\rhcaMsePLgbmCloseOos}{}\renewcommand{\rhcaMsePLgbmCloseOos}{\ensuremath{p = 0.22}}
\providecommand{\rhcaGapLgbmFullCloseOos}{}\renewcommand{\rhcaGapLgbmFullCloseOos}{6.0}  % QLIKE % vs the closing-interval HAR benchmark
\providecommand{\rhcaGapDmLgbmFullCloseOos}{}\renewcommand{\rhcaGapDmLgbmFullCloseOos}{1.41}
\providecommand{\rhcaGapPLgbmFullCloseOos}{}\renewcommand{\rhcaGapPLgbmFullCloseOos}{\ensuremath{p = 0.16}}
\providecommand{\rhcaQXgbHarCloseOos}{}\renewcommand{\rhcaQXgbHarCloseOos}{0.1102}
\providecommand{\rhcaMseXgbHarCloseOos}{}\renewcommand{\rhcaMseXgbHarCloseOos}{3.62}
\providecommand{\rhcaMaeXgbHarCloseOos}{}\renewcommand{\rhcaMaeXgbHarCloseOos}{4.51}
\providecommand{\rhcaRsqXgbHarCloseOos}{}\renewcommand{\rhcaRsqXgbHarCloseOos}{81.0}
\providecommand{\rhcaGapXgbHarCloseOos}{}\renewcommand{\rhcaGapXgbHarCloseOos}{7.5}  % QLIKE % vs the closing-interval HAR benchmark
\providecommand{\rhcaGapDmXgbHarCloseOos}{}\renewcommand{\rhcaGapDmXgbHarCloseOos}{2.70}
\providecommand{\rhcaGapPXgbHarCloseOos}{}\renewcommand{\rhcaGapPXgbHarCloseOos}{\ensuremath{p = 0.007}}
\providecommand{\rhcaQXgbFullCloseOos}{}\renewcommand{\rhcaQXgbFullCloseOos}{0.1083}
\providecommand{\rhcaMseXgbFullCloseOos}{}\renewcommand{\rhcaMseXgbFullCloseOos}{4.76}
\providecommand{\rhcaMaeXgbFullCloseOos}{}\renewcommand{\rhcaMaeXgbFullCloseOos}{4.76}
\providecommand{\rhcaRsqXgbFullCloseOos}{}\renewcommand{\rhcaRsqXgbFullCloseOos}{75.0}
\providecommand{\rhcaPctXgbCloseOos}{}\renewcommand{\rhcaPctXgbCloseOos}{\ensuremath{-}1.7}  % QLIKE %, full vs HAR of the same class
\providecommand{\rhcaDmXgbCloseOos}{}\renewcommand{\rhcaDmXgbCloseOos}{\ensuremath{-}0.63}
\providecommand{\rhcaPXgbCloseOos}{}\renewcommand{\rhcaPXgbCloseOos}{\ensuremath{p = 0.53}}
\providecommand{\rhcaMsePctXgbCloseOos}{}\renewcommand{\rhcaMsePctXgbCloseOos}{31.4}  % MSE %, full vs HAR
\providecommand{\rhcaMseDmXgbCloseOos}{}\renewcommand{\rhcaMseDmXgbCloseOos}{1.23}
\providecommand{\rhcaMsePXgbCloseOos}{}\renewcommand{\rhcaMsePXgbCloseOos}{\ensuremath{p = 0.22}}
\providecommand{\rhcaGapXgbFullCloseOos}{}\renewcommand{\rhcaGapXgbFullCloseOos}{5.7}  % QLIKE % vs the closing-interval HAR benchmark
\providecommand{\rhcaGapDmXgbFullCloseOos}{}\renewcommand{\rhcaGapDmXgbFullCloseOos}{1.40}
\providecommand{\rhcaGapPXgbFullCloseOos}{}\renewcommand{\rhcaGapPXgbFullCloseOos}{\ensuremath{p = 0.16}}
\providecommand{\rhcaQRfHarCloseOos}{}\renewcommand{\rhcaQRfHarCloseOos}{0.1131}
\providecommand{\rhcaMseRfHarCloseOos}{}\renewcommand{\rhcaMseRfHarCloseOos}{3.39}
\providecommand{\rhcaMaeRfHarCloseOos}{}\renewcommand{\rhcaMaeRfHarCloseOos}{4.42}
\providecommand{\rhcaRsqRfHarCloseOos}{}\renewcommand{\rhcaRsqRfHarCloseOos}{82.2}
\providecommand{\rhcaGapRfHarCloseOos}{}\renewcommand{\rhcaGapRfHarCloseOos}{10.4}  % QLIKE % vs the closing-interval HAR benchmark
\providecommand{\rhcaGapDmRfHarCloseOos}{}\renewcommand{\rhcaGapDmRfHarCloseOos}{3.12}
\providecommand{\rhcaGapPRfHarCloseOos}{}\renewcommand{\rhcaGapPRfHarCloseOos}{\ensuremath{p = 0.002}}
\providecommand{\rhcaQRfFullCloseOos}{}\renewcommand{\rhcaQRfFullCloseOos}{0.1123}
\providecommand{\rhcaMseRfFullCloseOos}{}\renewcommand{\rhcaMseRfFullCloseOos}{4.18}
\providecommand{\rhcaMaeRfFullCloseOos}{}\renewcommand{\rhcaMaeRfFullCloseOos}{4.63}
\providecommand{\rhcaRsqRfFullCloseOos}{}\renewcommand{\rhcaRsqRfFullCloseOos}{78.0}
\providecommand{\rhcaPctRfCloseOos}{}\renewcommand{\rhcaPctRfCloseOos}{\ensuremath{-}0.7}  % QLIKE %, full vs HAR of the same class
\providecommand{\rhcaDmRfCloseOos}{}\renewcommand{\rhcaDmRfCloseOos}{\ensuremath{-}0.27}
\providecommand{\rhcaPRfCloseOos}{}\renewcommand{\rhcaPRfCloseOos}{\ensuremath{p = 0.79}}
\providecommand{\rhcaMsePctRfCloseOos}{}\renewcommand{\rhcaMsePctRfCloseOos}{23.3}  % MSE %, full vs HAR
\providecommand{\rhcaMseDmRfCloseOos}{}\renewcommand{\rhcaMseDmRfCloseOos}{1.25}
\providecommand{\rhcaMsePRfCloseOos}{}\renewcommand{\rhcaMsePRfCloseOos}{\ensuremath{p = 0.21}}
\providecommand{\rhcaGapRfFullCloseOos}{}\renewcommand{\rhcaGapRfFullCloseOos}{9.6}  % QLIKE % vs the closing-interval HAR benchmark
\providecommand{\rhcaGapDmRfFullCloseOos}{}\renewcommand{\rhcaGapDmRfFullCloseOos}{2.13}
\providecommand{\rhcaGapPRfFullCloseOos}{}\renewcommand{\rhcaGapPRfFullCloseOos}{\ensuremath{p = 0.033}}
\providecommand{\rhcaQBmCloseOos}{}\renewcommand{\rhcaQBmCloseOos}{0.1025}
\providecommand{\rhcaMseBmCloseOos}{}\renewcommand{\rhcaMseBmCloseOos}{4.31}
\providecommand{\rhcaMaeBmCloseOos}{}\renewcommand{\rhcaMaeBmCloseOos}{4.78}
\providecommand{\rhcaRsqBmCloseOos}{}\renewcommand{\rhcaRsqBmCloseOos}{77.4}
\providecommand{\rhcaQCiRidgeHarCloseOos}{}\renewcommand{\rhcaQCiRidgeHarCloseOos}{0.1033}
\providecommand{\rhcaMseCiRidgeHarCloseOos}{}\renewcommand{\rhcaMseCiRidgeHarCloseOos}{4.15}
\providecommand{\rhcaMaeCiRidgeHarCloseOos}{}\renewcommand{\rhcaMaeCiRidgeHarCloseOos}{4.67}
\providecommand{\rhcaRsqCiRidgeHarCloseOos}{}\renewcommand{\rhcaRsqCiRidgeHarCloseOos}{78.2}
\providecommand{\rhcaGapCiRidgeHarCloseOos}{}\renewcommand{\rhcaGapCiRidgeHarCloseOos}{0.7}  % QLIKE % vs the closing-interval HAR benchmark
\providecommand{\rhcaGapDmCiRidgeHarCloseOos}{}\renewcommand{\rhcaGapDmCiRidgeHarCloseOos}{0.86}
\providecommand{\rhcaGapPCiRidgeHarCloseOos}{}\renewcommand{\rhcaGapPCiRidgeHarCloseOos}{\ensuremath{p = 0.39}}
\providecommand{\rhcaQCiRidgeFullCloseOos}{}\renewcommand{\rhcaQCiRidgeFullCloseOos}{0.1031}
\providecommand{\rhcaMseCiRidgeFullCloseOos}{}\renewcommand{\rhcaMseCiRidgeFullCloseOos}{7.30}
\providecommand{\rhcaMaeCiRidgeFullCloseOos}{}\renewcommand{\rhcaMaeCiRidgeFullCloseOos}{5.27}
\providecommand{\rhcaRsqCiRidgeFullCloseOos}{}\renewcommand{\rhcaRsqCiRidgeFullCloseOos}{61.7}
\providecommand{\rhcaPctCiRidgeCloseOos}{}\renewcommand{\rhcaPctCiRidgeCloseOos}{\ensuremath{-}0.1}  % QLIKE %, full vs HAR of the same class
\providecommand{\rhcaDmCiRidgeCloseOos}{}\renewcommand{\rhcaDmCiRidgeCloseOos}{\ensuremath{-}0.04}
\providecommand{\rhcaPCiRidgeCloseOos}{}\renewcommand{\rhcaPCiRidgeCloseOos}{\ensuremath{p = 0.97}}
\providecommand{\rhcaMsePctCiRidgeCloseOos}{}\renewcommand{\rhcaMsePctCiRidgeCloseOos}{75.8}  % MSE %, full vs HAR
\providecommand{\rhcaMseDmCiRidgeCloseOos}{}\renewcommand{\rhcaMseDmCiRidgeCloseOos}{1.45}
\providecommand{\rhcaMsePCiRidgeCloseOos}{}\renewcommand{\rhcaMsePCiRidgeCloseOos}{\ensuremath{p = 0.15}}
\providecommand{\rhcaGapCiRidgeFullCloseOos}{}\renewcommand{\rhcaGapCiRidgeFullCloseOos}{0.6}  % QLIKE % vs the closing-interval HAR benchmark
\providecommand{\rhcaGapDmCiRidgeFullCloseOos}{}\renewcommand{\rhcaGapDmCiRidgeFullCloseOos}{0.15}
\providecommand{\rhcaGapPCiRidgeFullCloseOos}{}\renewcommand{\rhcaGapPCiRidgeFullCloseOos}{\ensuremath{p = 0.88}}
\providecommand{\rhcaQCiLassoHarCloseOos}{}\renewcommand{\rhcaQCiLassoHarCloseOos}{0.1033}
\providecommand{\rhcaMseCiLassoHarCloseOos}{}\renewcommand{\rhcaMseCiLassoHarCloseOos}{4.15}
\providecommand{\rhcaMaeCiLassoHarCloseOos}{}\renewcommand{\rhcaMaeCiLassoHarCloseOos}{4.67}
\providecommand{\rhcaRsqCiLassoHarCloseOos}{}\renewcommand{\rhcaRsqCiLassoHarCloseOos}{78.2}
\providecommand{\rhcaGapCiLassoHarCloseOos}{}\renewcommand{\rhcaGapCiLassoHarCloseOos}{0.7}  % QLIKE % vs the closing-interval HAR benchmark
\providecommand{\rhcaGapDmCiLassoHarCloseOos}{}\renewcommand{\rhcaGapDmCiLassoHarCloseOos}{0.86}
\providecommand{\rhcaGapPCiLassoHarCloseOos}{}\renewcommand{\rhcaGapPCiLassoHarCloseOos}{\ensuremath{p = 0.39}}
\providecommand{\rhcaQCiLassoFullCloseOos}{}\renewcommand{\rhcaQCiLassoFullCloseOos}{0.0992}
\providecommand{\rhcaMseCiLassoFullCloseOos}{}\renewcommand{\rhcaMseCiLassoFullCloseOos}{5.94}
\providecommand{\rhcaMaeCiLassoFullCloseOos}{}\renewcommand{\rhcaMaeCiLassoFullCloseOos}{5.07}
\providecommand{\rhcaRsqCiLassoFullCloseOos}{}\renewcommand{\rhcaRsqCiLassoFullCloseOos}{68.8}
\providecommand{\rhcaPctCiLassoCloseOos}{}\renewcommand{\rhcaPctCiLassoCloseOos}{\ensuremath{-}4.0}  % QLIKE %, full vs HAR of the same class
\providecommand{\rhcaDmCiLassoCloseOos}{}\renewcommand{\rhcaDmCiLassoCloseOos}{\ensuremath{-}1.91}
\providecommand{\rhcaPCiLassoCloseOos}{}\renewcommand{\rhcaPCiLassoCloseOos}{\ensuremath{p = 0.056}}
\providecommand{\rhcaMsePctCiLassoCloseOos}{}\renewcommand{\rhcaMsePctCiLassoCloseOos}{43.1}  % MSE %, full vs HAR
\providecommand{\rhcaMseDmCiLassoCloseOos}{}\renewcommand{\rhcaMseDmCiLassoCloseOos}{1.34}
\providecommand{\rhcaMsePCiLassoCloseOos}{}\renewcommand{\rhcaMsePCiLassoCloseOos}{\ensuremath{p = 0.18}}
\providecommand{\rhcaGapCiLassoFullCloseOos}{}\renewcommand{\rhcaGapCiLassoFullCloseOos}{\ensuremath{-}3.3}  % QLIKE % vs the closing-interval HAR benchmark
\providecommand{\rhcaGapDmCiLassoFullCloseOos}{}\renewcommand{\rhcaGapDmCiLassoFullCloseOos}{\ensuremath{-}1.44}
\providecommand{\rhcaGapPCiLassoFullCloseOos}{}\renewcommand{\rhcaGapPCiLassoFullCloseOos}{\ensuremath{p = 0.15}}
\providecommand{\rhcaQCiEnetHarCloseOos}{}\renewcommand{\rhcaQCiEnetHarCloseOos}{0.1033}
\providecommand{\rhcaMseCiEnetHarCloseOos}{}\renewcommand{\rhcaMseCiEnetHarCloseOos}{4.15}
\providecommand{\rhcaMaeCiEnetHarCloseOos}{}\renewcommand{\rhcaMaeCiEnetHarCloseOos}{4.67}
\providecommand{\rhcaRsqCiEnetHarCloseOos}{}\renewcommand{\rhcaRsqCiEnetHarCloseOos}{78.2}
\providecommand{\rhcaGapCiEnetHarCloseOos}{}\renewcommand{\rhcaGapCiEnetHarCloseOos}{0.7}  % QLIKE % vs the closing-interval HAR benchmark
\providecommand{\rhcaGapDmCiEnetHarCloseOos}{}\renewcommand{\rhcaGapDmCiEnetHarCloseOos}{0.86}
\providecommand{\rhcaGapPCiEnetHarCloseOos}{}\renewcommand{\rhcaGapPCiEnetHarCloseOos}{\ensuremath{p = 0.39}}
\providecommand{\rhcaQCiEnetFullCloseOos}{}\renewcommand{\rhcaQCiEnetFullCloseOos}{0.1014}
\providecommand{\rhcaMseCiEnetFullCloseOos}{}\renewcommand{\rhcaMseCiEnetFullCloseOos}{6.89}
\providecommand{\rhcaMaeCiEnetFullCloseOos}{}\renewcommand{\rhcaMaeCiEnetFullCloseOos}{5.21}
\providecommand{\rhcaRsqCiEnetFullCloseOos}{}\renewcommand{\rhcaRsqCiEnetFullCloseOos}{63.9}
\providecommand{\rhcaPctCiEnetCloseOos}{}\renewcommand{\rhcaPctCiEnetCloseOos}{\ensuremath{-}1.8}  % QLIKE %, full vs HAR of the same class
\providecommand{\rhcaDmCiEnetCloseOos}{}\renewcommand{\rhcaDmCiEnetCloseOos}{\ensuremath{-}0.55}
\providecommand{\rhcaPCiEnetCloseOos}{}\renewcommand{\rhcaPCiEnetCloseOos}{\ensuremath{p = 0.58}}
\providecommand{\rhcaMsePctCiEnetCloseOos}{}\renewcommand{\rhcaMsePctCiEnetCloseOos}{65.9}  % MSE %, full vs HAR
\providecommand{\rhcaMseDmCiEnetCloseOos}{}\renewcommand{\rhcaMseDmCiEnetCloseOos}{1.45}
\providecommand{\rhcaMsePCiEnetCloseOos}{}\renewcommand{\rhcaMsePCiEnetCloseOos}{\ensuremath{p = 0.15}}
\providecommand{\rhcaGapCiEnetFullCloseOos}{}\renewcommand{\rhcaGapCiEnetFullCloseOos}{\ensuremath{-}1.1}  % QLIKE % vs the closing-interval HAR benchmark
\providecommand{\rhcaGapDmCiEnetFullCloseOos}{}\renewcommand{\rhcaGapDmCiEnetFullCloseOos}{\ensuremath{-}0.32}
\providecommand{\rhcaGapPCiEnetFullCloseOos}{}\renewcommand{\rhcaGapPCiEnetFullCloseOos}{\ensuremath{p = 0.75}}
\providecommand{\rhcaQCiLgbmHarCloseOos}{}\renewcommand{\rhcaQCiLgbmHarCloseOos}{0.1094}
\providecommand{\rhcaMseCiLgbmHarCloseOos}{}\renewcommand{\rhcaMseCiLgbmHarCloseOos}{5.75}
\providecommand{\rhcaMaeCiLgbmHarCloseOos}{}\renewcommand{\rhcaMaeCiLgbmHarCloseOos}{5.02}
\providecommand{\rhcaRsqCiLgbmHarCloseOos}{}\renewcommand{\rhcaRsqCiLgbmHarCloseOos}{69.8}
\providecommand{\rhcaGapCiLgbmHarCloseOos}{}\renewcommand{\rhcaGapCiLgbmHarCloseOos}{6.7}  % QLIKE % vs the closing-interval HAR benchmark
\providecommand{\rhcaGapDmCiLgbmHarCloseOos}{}\renewcommand{\rhcaGapDmCiLgbmHarCloseOos}{3.03}
\providecommand{\rhcaGapPCiLgbmHarCloseOos}{}\renewcommand{\rhcaGapPCiLgbmHarCloseOos}{\ensuremath{p = 0.002}}
\providecommand{\rhcaQCiLgbmFullCloseOos}{}\renewcommand{\rhcaQCiLgbmFullCloseOos}{0.1049}
\providecommand{\rhcaMseCiLgbmFullCloseOos}{}\renewcommand{\rhcaMseCiLgbmFullCloseOos}{17.19}
\providecommand{\rhcaMaeCiLgbmFullCloseOos}{}\renewcommand{\rhcaMaeCiLgbmFullCloseOos}{6.84}
\providecommand{\rhcaRsqCiLgbmFullCloseOos}{}\renewcommand{\rhcaRsqCiLgbmFullCloseOos}{9.8}
\providecommand{\rhcaPctCiLgbmCloseOos}{}\renewcommand{\rhcaPctCiLgbmCloseOos}{\ensuremath{-}4.1}  % QLIKE %, full vs HAR of the same class
\providecommand{\rhcaDmCiLgbmCloseOos}{}\renewcommand{\rhcaDmCiLgbmCloseOos}{\ensuremath{-}1.48}
\providecommand{\rhcaPCiLgbmCloseOos}{}\renewcommand{\rhcaPCiLgbmCloseOos}{\ensuremath{p = 0.14}}
\providecommand{\rhcaMsePctCiLgbmCloseOos}{}\renewcommand{\rhcaMsePctCiLgbmCloseOos}{199.1}  % MSE %, full vs HAR
\providecommand{\rhcaMseDmCiLgbmCloseOos}{}\renewcommand{\rhcaMseDmCiLgbmCloseOos}{1.21}
\providecommand{\rhcaMsePCiLgbmCloseOos}{}\renewcommand{\rhcaMsePCiLgbmCloseOos}{\ensuremath{p = 0.23}}
\providecommand{\rhcaGapCiLgbmFullCloseOos}{}\renewcommand{\rhcaGapCiLgbmFullCloseOos}{2.3}  % QLIKE % vs the closing-interval HAR benchmark
\providecommand{\rhcaGapDmCiLgbmFullCloseOos}{}\renewcommand{\rhcaGapDmCiLgbmFullCloseOos}{0.72}
\providecommand{\rhcaGapPCiLgbmFullCloseOos}{}\renewcommand{\rhcaGapPCiLgbmFullCloseOos}{\ensuremath{p = 0.47}}
\providecommand{\rhcaQCiXgbHarCloseOos}{}\renewcommand{\rhcaQCiXgbHarCloseOos}{0.1076}
\providecommand{\rhcaMseCiXgbHarCloseOos}{}\renewcommand{\rhcaMseCiXgbHarCloseOos}{5.27}
\providecommand{\rhcaMaeCiXgbHarCloseOos}{}\renewcommand{\rhcaMaeCiXgbHarCloseOos}{4.94}
\providecommand{\rhcaRsqCiXgbHarCloseOos}{}\renewcommand{\rhcaRsqCiXgbHarCloseOos}{72.4}
\providecommand{\rhcaGapCiXgbHarCloseOos}{}\renewcommand{\rhcaGapCiXgbHarCloseOos}{5.0}  % QLIKE % vs the closing-interval HAR benchmark
\providecommand{\rhcaGapDmCiXgbHarCloseOos}{}\renewcommand{\rhcaGapDmCiXgbHarCloseOos}{2.45}
\providecommand{\rhcaGapPCiXgbHarCloseOos}{}\renewcommand{\rhcaGapPCiXgbHarCloseOos}{\ensuremath{p = 0.014}}
\providecommand{\rhcaQCiXgbFullCloseOos}{}\renewcommand{\rhcaQCiXgbFullCloseOos}{0.1031}
\providecommand{\rhcaMseCiXgbFullCloseOos}{}\renewcommand{\rhcaMseCiXgbFullCloseOos}{20.94}
\providecommand{\rhcaMaeCiXgbFullCloseOos}{}\renewcommand{\rhcaMaeCiXgbFullCloseOos}{6.98}
\providecommand{\rhcaRsqCiXgbFullCloseOos}{}\renewcommand{\rhcaRsqCiXgbFullCloseOos}{\ensuremath{-}9.9}
\providecommand{\rhcaPctCiXgbCloseOos}{}\renewcommand{\rhcaPctCiXgbCloseOos}{\ensuremath{-}4.1}  % QLIKE %, full vs HAR of the same class
\providecommand{\rhcaDmCiXgbCloseOos}{}\renewcommand{\rhcaDmCiXgbCloseOos}{\ensuremath{-}1.63}
\providecommand{\rhcaPCiXgbCloseOos}{}\renewcommand{\rhcaPCiXgbCloseOos}{\ensuremath{p = 0.1039}}
\providecommand{\rhcaMsePctCiXgbCloseOos}{}\renewcommand{\rhcaMsePctCiXgbCloseOos}{297.6}  % MSE %, full vs HAR
\providecommand{\rhcaMseDmCiXgbCloseOos}{}\renewcommand{\rhcaMseDmCiXgbCloseOos}{1.15}
\providecommand{\rhcaMsePCiXgbCloseOos}{}\renewcommand{\rhcaMsePCiXgbCloseOos}{\ensuremath{p = 0.25}}
\providecommand{\rhcaGapCiXgbFullCloseOos}{}\renewcommand{\rhcaGapCiXgbFullCloseOos}{0.6}  % QLIKE % vs the closing-interval HAR benchmark
\providecommand{\rhcaGapDmCiXgbFullCloseOos}{}\renewcommand{\rhcaGapDmCiXgbFullCloseOos}{0.21}
\providecommand{\rhcaGapPCiXgbFullCloseOos}{}\renewcommand{\rhcaGapPCiXgbFullCloseOos}{\ensuremath{p = 0.83}}
\providecommand{\rhcaQCiRfHarCloseOos}{}\renewcommand{\rhcaQCiRfHarCloseOos}{0.1068}
\providecommand{\rhcaMseCiRfHarCloseOos}{}\renewcommand{\rhcaMseCiRfHarCloseOos}{4.62}
\providecommand{\rhcaMaeCiRfHarCloseOos}{}\renewcommand{\rhcaMaeCiRfHarCloseOos}{4.76}
\providecommand{\rhcaRsqCiRfHarCloseOos}{}\renewcommand{\rhcaRsqCiRfHarCloseOos}{75.8}
\providecommand{\rhcaGapCiRfHarCloseOos}{}\renewcommand{\rhcaGapCiRfHarCloseOos}{4.2}  % QLIKE % vs the closing-interval HAR benchmark
\providecommand{\rhcaGapDmCiRfHarCloseOos}{}\renewcommand{\rhcaGapDmCiRfHarCloseOos}{2.28}
\providecommand{\rhcaGapPCiRfHarCloseOos}{}\renewcommand{\rhcaGapPCiRfHarCloseOos}{\ensuremath{p = 0.022}}
\providecommand{\rhcaQCiRfFullCloseOos}{}\renewcommand{\rhcaQCiRfFullCloseOos}{0.1060}
\providecommand{\rhcaMseCiRfFullCloseOos}{}\renewcommand{\rhcaMseCiRfFullCloseOos}{14.15}
\providecommand{\rhcaMaeCiRfFullCloseOos}{}\renewcommand{\rhcaMaeCiRfFullCloseOos}{6.74}
\providecommand{\rhcaRsqCiRfFullCloseOos}{}\renewcommand{\rhcaRsqCiRfFullCloseOos}{25.8}
\providecommand{\rhcaPctCiRfCloseOos}{}\renewcommand{\rhcaPctCiRfCloseOos}{\ensuremath{-}0.8}  % QLIKE %, full vs HAR of the same class
\providecommand{\rhcaDmCiRfCloseOos}{}\renewcommand{\rhcaDmCiRfCloseOos}{\ensuremath{-}0.28}
\providecommand{\rhcaPCiRfCloseOos}{}\renewcommand{\rhcaPCiRfCloseOos}{\ensuremath{p = 0.78}}
\providecommand{\rhcaMsePctCiRfCloseOos}{}\renewcommand{\rhcaMsePctCiRfCloseOos}{206.6}  % MSE %, full vs HAR
\providecommand{\rhcaMseDmCiRfCloseOos}{}\renewcommand{\rhcaMseDmCiRfCloseOos}{1.22}
\providecommand{\rhcaMsePCiRfCloseOos}{}\renewcommand{\rhcaMsePCiRfCloseOos}{\ensuremath{p = 0.22}}
\providecommand{\rhcaGapCiRfFullCloseOos}{}\renewcommand{\rhcaGapCiRfFullCloseOos}{3.4}  % QLIKE % vs the closing-interval HAR benchmark
\providecommand{\rhcaGapDmCiRfFullCloseOos}{}\renewcommand{\rhcaGapDmCiRfFullCloseOos}{1.11}
\providecommand{\rhcaGapPCiRfFullCloseOos}{}\renewcommand{\rhcaGapPCiRfFullCloseOos}{\ensuremath{p = 0.27}}
% --- every bar: pooled forecasts, Span / Pre / Oos (Table 1 and its subperiods)
\providecommand{\rhcaQOlsHarBarSpan}{}\renewcommand{\rhcaQOlsHarBarSpan}{0.1893}
\providecommand{\rhcaMseOlsHarBarSpan}{}\renewcommand{\rhcaMseOlsHarBarSpan}{1.11}
\providecommand{\rhcaMaeOlsHarBarSpan}{}\renewcommand{\rhcaMaeOlsHarBarSpan}{2.25}
\providecommand{\rhcaRsqOlsHarBarSpan}{}\renewcommand{\rhcaRsqOlsHarBarSpan}{65.4}
\providecommand{\rhcaQRidgeHarBarSpan}{}\renewcommand{\rhcaQRidgeHarBarSpan}{0.1893}
\providecommand{\rhcaMseRidgeHarBarSpan}{}\renewcommand{\rhcaMseRidgeHarBarSpan}{1.11}
\providecommand{\rhcaMaeRidgeHarBarSpan}{}\renewcommand{\rhcaMaeRidgeHarBarSpan}{2.25}
\providecommand{\rhcaRsqRidgeHarBarSpan}{}\renewcommand{\rhcaRsqRidgeHarBarSpan}{65.4}
\providecommand{\rhcaQRidgeFullBarSpan}{}\renewcommand{\rhcaQRidgeFullBarSpan}{0.1778}
\providecommand{\rhcaMseRidgeFullBarSpan}{}\renewcommand{\rhcaMseRidgeFullBarSpan}{2.23}
\providecommand{\rhcaMaeRidgeFullBarSpan}{}\renewcommand{\rhcaMaeRidgeFullBarSpan}{2.38}
\providecommand{\rhcaRsqRidgeFullBarSpan}{}\renewcommand{\rhcaRsqRidgeFullBarSpan}{30.7}
\providecommand{\rhcaPctRidgeBarSpan}{}\renewcommand{\rhcaPctRidgeBarSpan}{\ensuremath{-}6.0}
\providecommand{\rhcaDmRidgeBarSpan}{}\renewcommand{\rhcaDmRidgeBarSpan}{\ensuremath{-}6.28}
\providecommand{\rhcaPRidgeBarSpan}{}\renewcommand{\rhcaPRidgeBarSpan}{\ensuremath{p < 0.001}}
\providecommand{\rhcaMsePctRidgeBarSpan}{}\renewcommand{\rhcaMsePctRidgeBarSpan}{100.6}
\providecommand{\rhcaMseDmRidgeBarSpan}{}\renewcommand{\rhcaMseDmRidgeBarSpan}{1.12}
\providecommand{\rhcaMsePRidgeBarSpan}{}\renewcommand{\rhcaMsePRidgeBarSpan}{\ensuremath{p = 0.26}}
\providecommand{\rhcaQLassoHarBarSpan}{}\renewcommand{\rhcaQLassoHarBarSpan}{0.1893}
\providecommand{\rhcaMseLassoHarBarSpan}{}\renewcommand{\rhcaMseLassoHarBarSpan}{1.11}
\providecommand{\rhcaMaeLassoHarBarSpan}{}\renewcommand{\rhcaMaeLassoHarBarSpan}{2.25}
\providecommand{\rhcaRsqLassoHarBarSpan}{}\renewcommand{\rhcaRsqLassoHarBarSpan}{65.4}
\providecommand{\rhcaQLassoFullBarSpan}{}\renewcommand{\rhcaQLassoFullBarSpan}{0.1723}
\providecommand{\rhcaMseLassoFullBarSpan}{}\renewcommand{\rhcaMseLassoFullBarSpan}{2.39}
\providecommand{\rhcaMaeLassoFullBarSpan}{}\renewcommand{\rhcaMaeLassoFullBarSpan}{2.40}
\providecommand{\rhcaRsqLassoFullBarSpan}{}\renewcommand{\rhcaRsqLassoFullBarSpan}{25.8}
\providecommand{\rhcaPctLassoBarSpan}{}\renewcommand{\rhcaPctLassoBarSpan}{\ensuremath{-}9.0}
\providecommand{\rhcaDmLassoBarSpan}{}\renewcommand{\rhcaDmLassoBarSpan}{\ensuremath{-}6.34}
\providecommand{\rhcaPLassoBarSpan}{}\renewcommand{\rhcaPLassoBarSpan}{\ensuremath{p < 0.001}}
\providecommand{\rhcaMsePctLassoBarSpan}{}\renewcommand{\rhcaMsePctLassoBarSpan}{114.7}
\providecommand{\rhcaMseDmLassoBarSpan}{}\renewcommand{\rhcaMseDmLassoBarSpan}{1.16}
\providecommand{\rhcaMsePLassoBarSpan}{}\renewcommand{\rhcaMsePLassoBarSpan}{\ensuremath{p = 0.25}}
\providecommand{\rhcaQEnetHarBarSpan}{}\renewcommand{\rhcaQEnetHarBarSpan}{0.1893}
\providecommand{\rhcaMseEnetHarBarSpan}{}\renewcommand{\rhcaMseEnetHarBarSpan}{1.11}
\providecommand{\rhcaMaeEnetHarBarSpan}{}\renewcommand{\rhcaMaeEnetHarBarSpan}{2.25}
\providecommand{\rhcaRsqEnetHarBarSpan}{}\renewcommand{\rhcaRsqEnetHarBarSpan}{65.4}
\providecommand{\rhcaQEnetFullBarSpan}{}\renewcommand{\rhcaQEnetFullBarSpan}{0.1766}
\providecommand{\rhcaMseEnetFullBarSpan}{}\renewcommand{\rhcaMseEnetFullBarSpan}{2.42}
\providecommand{\rhcaMaeEnetFullBarSpan}{}\renewcommand{\rhcaMaeEnetFullBarSpan}{2.40}
\providecommand{\rhcaRsqEnetFullBarSpan}{}\renewcommand{\rhcaRsqEnetFullBarSpan}{24.8}
\providecommand{\rhcaPctEnetBarSpan}{}\renewcommand{\rhcaPctEnetBarSpan}{\ensuremath{-}6.7}
\providecommand{\rhcaDmEnetBarSpan}{}\renewcommand{\rhcaDmEnetBarSpan}{\ensuremath{-}7.36}
\providecommand{\rhcaPEnetBarSpan}{}\renewcommand{\rhcaPEnetBarSpan}{\ensuremath{p < 0.001}}
\providecommand{\rhcaMsePctEnetBarSpan}{}\renewcommand{\rhcaMsePctEnetBarSpan}{117.7}
\providecommand{\rhcaMseDmEnetBarSpan}{}\renewcommand{\rhcaMseDmEnetBarSpan}{1.12}
\providecommand{\rhcaMsePEnetBarSpan}{}\renewcommand{\rhcaMsePEnetBarSpan}{\ensuremath{p = 0.26}}
\providecommand{\rhcaQLgbmHarBarSpan}{}\renewcommand{\rhcaQLgbmHarBarSpan}{0.1875}
\providecommand{\rhcaMseLgbmHarBarSpan}{}\renewcommand{\rhcaMseLgbmHarBarSpan}{1.18}
\providecommand{\rhcaMaeLgbmHarBarSpan}{}\renewcommand{\rhcaMaeLgbmHarBarSpan}{2.27}
\providecommand{\rhcaRsqLgbmHarBarSpan}{}\renewcommand{\rhcaRsqLgbmHarBarSpan}{63.3}
\providecommand{\rhcaQLgbmFullBarSpan}{}\renewcommand{\rhcaQLgbmFullBarSpan}{0.1611}
\providecommand{\rhcaMseLgbmFullBarSpan}{}\renewcommand{\rhcaMseLgbmFullBarSpan}{1.42}
\providecommand{\rhcaMaeLgbmFullBarSpan}{}\renewcommand{\rhcaMaeLgbmFullBarSpan}{2.28}
\providecommand{\rhcaRsqLgbmFullBarSpan}{}\renewcommand{\rhcaRsqLgbmFullBarSpan}{55.9}
\providecommand{\rhcaPctLgbmBarSpan}{}\renewcommand{\rhcaPctLgbmBarSpan}{\ensuremath{-}14.1}
\providecommand{\rhcaDmLgbmBarSpan}{}\renewcommand{\rhcaDmLgbmBarSpan}{\ensuremath{-}10.26}
\providecommand{\rhcaPLgbmBarSpan}{}\renewcommand{\rhcaPLgbmBarSpan}{\ensuremath{p < 0.001}}
\providecommand{\rhcaMsePctLgbmBarSpan}{}\renewcommand{\rhcaMsePctLgbmBarSpan}{20.0}
\providecommand{\rhcaMseDmLgbmBarSpan}{}\renewcommand{\rhcaMseDmLgbmBarSpan}{1.03}
\providecommand{\rhcaMsePLgbmBarSpan}{}\renewcommand{\rhcaMsePLgbmBarSpan}{\ensuremath{p = 0.30}}
\providecommand{\rhcaQXgbHarBarSpan}{}\renewcommand{\rhcaQXgbHarBarSpan}{0.1888}
\providecommand{\rhcaMseXgbHarBarSpan}{}\renewcommand{\rhcaMseXgbHarBarSpan}{1.16}
\providecommand{\rhcaMaeXgbHarBarSpan}{}\renewcommand{\rhcaMaeXgbHarBarSpan}{2.26}
\providecommand{\rhcaRsqXgbHarBarSpan}{}\renewcommand{\rhcaRsqXgbHarBarSpan}{63.9}
\providecommand{\rhcaQXgbFullBarSpan}{}\renewcommand{\rhcaQXgbFullBarSpan}{0.1612}
\providecommand{\rhcaMseXgbFullBarSpan}{}\renewcommand{\rhcaMseXgbFullBarSpan}{1.35}
\providecommand{\rhcaMaeXgbFullBarSpan}{}\renewcommand{\rhcaMaeXgbFullBarSpan}{2.27}
\providecommand{\rhcaRsqXgbFullBarSpan}{}\renewcommand{\rhcaRsqXgbFullBarSpan}{58.0}
\providecommand{\rhcaPctXgbBarSpan}{}\renewcommand{\rhcaPctXgbBarSpan}{\ensuremath{-}14.6}
\providecommand{\rhcaDmXgbBarSpan}{}\renewcommand{\rhcaDmXgbBarSpan}{\ensuremath{-}9.94}
\providecommand{\rhcaPXgbBarSpan}{}\renewcommand{\rhcaPXgbBarSpan}{\ensuremath{p < 0.001}}
\providecommand{\rhcaMsePctXgbBarSpan}{}\renewcommand{\rhcaMsePctXgbBarSpan}{16.3}
\providecommand{\rhcaMseDmXgbBarSpan}{}\renewcommand{\rhcaMseDmXgbBarSpan}{1.04}
\providecommand{\rhcaMsePXgbBarSpan}{}\renewcommand{\rhcaMsePXgbBarSpan}{\ensuremath{p = 0.30}}
\providecommand{\rhcaQRfHarBarSpan}{}\renewcommand{\rhcaQRfHarBarSpan}{0.1895}
\providecommand{\rhcaMseRfHarBarSpan}{}\renewcommand{\rhcaMseRfHarBarSpan}{1.19}
\providecommand{\rhcaMaeRfHarBarSpan}{}\renewcommand{\rhcaMaeRfHarBarSpan}{2.27}
\providecommand{\rhcaRsqRfHarBarSpan}{}\renewcommand{\rhcaRsqRfHarBarSpan}{63.0}
\providecommand{\rhcaQRfFullBarSpan}{}\renewcommand{\rhcaQRfFullBarSpan}{0.1681}
\providecommand{\rhcaMseRfFullBarSpan}{}\renewcommand{\rhcaMseRfFullBarSpan}{1.32}
\providecommand{\rhcaMaeRfFullBarSpan}{}\renewcommand{\rhcaMaeRfFullBarSpan}{2.27}
\providecommand{\rhcaRsqRfFullBarSpan}{}\renewcommand{\rhcaRsqRfFullBarSpan}{58.9}
\providecommand{\rhcaPctRfBarSpan}{}\renewcommand{\rhcaPctRfBarSpan}{\ensuremath{-}11.3}
\providecommand{\rhcaDmRfBarSpan}{}\renewcommand{\rhcaDmRfBarSpan}{\ensuremath{-}9.74}
\providecommand{\rhcaPRfBarSpan}{}\renewcommand{\rhcaPRfBarSpan}{\ensuremath{p < 0.001}}
\providecommand{\rhcaMsePctRfBarSpan}{}\renewcommand{\rhcaMsePctRfBarSpan}{10.9}
\providecommand{\rhcaMseDmRfBarSpan}{}\renewcommand{\rhcaMseDmRfBarSpan}{1.10}
\providecommand{\rhcaMsePRfBarSpan}{}\renewcommand{\rhcaMsePRfBarSpan}{\ensuremath{p = 0.27}}
\providecommand{\rhcaQOlsHarBarPre}{}\renewcommand{\rhcaQOlsHarBarPre}{0.1864}
\providecommand{\rhcaMseOlsHarBarPre}{}\renewcommand{\rhcaMseOlsHarBarPre}{0.31}
\providecommand{\rhcaMaeOlsHarBarPre}{}\renewcommand{\rhcaMaeOlsHarBarPre}{1.41}
\providecommand{\rhcaRsqOlsHarBarPre}{}\renewcommand{\rhcaRsqOlsHarBarPre}{61.4}
\providecommand{\rhcaQRidgeHarBarPre}{}\renewcommand{\rhcaQRidgeHarBarPre}{0.1864}
\providecommand{\rhcaMseRidgeHarBarPre}{}\renewcommand{\rhcaMseRidgeHarBarPre}{0.31}
\providecommand{\rhcaMaeRidgeHarBarPre}{}\renewcommand{\rhcaMaeRidgeHarBarPre}{1.41}
\providecommand{\rhcaRsqRidgeHarBarPre}{}\renewcommand{\rhcaRsqRidgeHarBarPre}{61.4}
\providecommand{\rhcaQRidgeFullBarPre}{}\renewcommand{\rhcaQRidgeFullBarPre}{0.1726}
\providecommand{\rhcaMseRidgeFullBarPre}{}\renewcommand{\rhcaMseRidgeFullBarPre}{0.28}
\providecommand{\rhcaMaeRidgeFullBarPre}{}\renewcommand{\rhcaMaeRidgeFullBarPre}{1.35}
\providecommand{\rhcaRsqRidgeFullBarPre}{}\renewcommand{\rhcaRsqRidgeFullBarPre}{64.9}
\providecommand{\rhcaPctRidgeBarPre}{}\renewcommand{\rhcaPctRidgeBarPre}{\ensuremath{-}7.4}
\providecommand{\rhcaDmRidgeBarPre}{}\renewcommand{\rhcaDmRidgeBarPre}{\ensuremath{-}6.82}
\providecommand{\rhcaPRidgeBarPre}{}\renewcommand{\rhcaPRidgeBarPre}{\ensuremath{p < 0.001}}
\providecommand{\rhcaMsePctRidgeBarPre}{}\renewcommand{\rhcaMsePctRidgeBarPre}{\ensuremath{-}8.9}
\providecommand{\rhcaMseDmRidgeBarPre}{}\renewcommand{\rhcaMseDmRidgeBarPre}{\ensuremath{-}1.79}
\providecommand{\rhcaMsePRidgeBarPre}{}\renewcommand{\rhcaMsePRidgeBarPre}{\ensuremath{p = 0.074}}
\providecommand{\rhcaQLassoHarBarPre}{}\renewcommand{\rhcaQLassoHarBarPre}{0.1864}
\providecommand{\rhcaMseLassoHarBarPre}{}\renewcommand{\rhcaMseLassoHarBarPre}{0.31}
\providecommand{\rhcaMaeLassoHarBarPre}{}\renewcommand{\rhcaMaeLassoHarBarPre}{1.41}
\providecommand{\rhcaRsqLassoHarBarPre}{}\renewcommand{\rhcaRsqLassoHarBarPre}{61.4}
\providecommand{\rhcaQLassoFullBarPre}{}\renewcommand{\rhcaQLassoFullBarPre}{0.1667}
\providecommand{\rhcaMseLassoFullBarPre}{}\renewcommand{\rhcaMseLassoFullBarPre}{0.28}
\providecommand{\rhcaMaeLassoFullBarPre}{}\renewcommand{\rhcaMaeLassoFullBarPre}{1.35}
\providecommand{\rhcaRsqLassoFullBarPre}{}\renewcommand{\rhcaRsqLassoFullBarPre}{65.6}
\providecommand{\rhcaPctLassoBarPre}{}\renewcommand{\rhcaPctLassoBarPre}{\ensuremath{-}10.6}
\providecommand{\rhcaDmLassoBarPre}{}\renewcommand{\rhcaDmLassoBarPre}{\ensuremath{-}6.88}
\providecommand{\rhcaPLassoBarPre}{}\renewcommand{\rhcaPLassoBarPre}{\ensuremath{p < 0.001}}
\providecommand{\rhcaMsePctLassoBarPre}{}\renewcommand{\rhcaMsePctLassoBarPre}{\ensuremath{-}10.8}
\providecommand{\rhcaMseDmLassoBarPre}{}\renewcommand{\rhcaMseDmLassoBarPre}{\ensuremath{-}1.81}
\providecommand{\rhcaMsePLassoBarPre}{}\renewcommand{\rhcaMsePLassoBarPre}{\ensuremath{p = 0.070}}
\providecommand{\rhcaQEnetHarBarPre}{}\renewcommand{\rhcaQEnetHarBarPre}{0.1864}
\providecommand{\rhcaMseEnetHarBarPre}{}\renewcommand{\rhcaMseEnetHarBarPre}{0.31}
\providecommand{\rhcaMaeEnetHarBarPre}{}\renewcommand{\rhcaMaeEnetHarBarPre}{1.41}
\providecommand{\rhcaRsqEnetHarBarPre}{}\renewcommand{\rhcaRsqEnetHarBarPre}{61.4}
\providecommand{\rhcaQEnetFullBarPre}{}\renewcommand{\rhcaQEnetFullBarPre}{0.1720}
\providecommand{\rhcaMseEnetFullBarPre}{}\renewcommand{\rhcaMseEnetFullBarPre}{0.28}
\providecommand{\rhcaMaeEnetFullBarPre}{}\renewcommand{\rhcaMaeEnetFullBarPre}{1.35}
\providecommand{\rhcaRsqEnetFullBarPre}{}\renewcommand{\rhcaRsqEnetFullBarPre}{65.2}
\providecommand{\rhcaPctEnetBarPre}{}\renewcommand{\rhcaPctEnetBarPre}{\ensuremath{-}7.7}
\providecommand{\rhcaDmEnetBarPre}{}\renewcommand{\rhcaDmEnetBarPre}{\ensuremath{-}6.97}
\providecommand{\rhcaPEnetBarPre}{}\renewcommand{\rhcaPEnetBarPre}{\ensuremath{p < 0.001}}
\providecommand{\rhcaMsePctEnetBarPre}{}\renewcommand{\rhcaMsePctEnetBarPre}{\ensuremath{-}9.7}
\providecommand{\rhcaMseDmEnetBarPre}{}\renewcommand{\rhcaMseDmEnetBarPre}{\ensuremath{-}1.75}
\providecommand{\rhcaMsePEnetBarPre}{}\renewcommand{\rhcaMsePEnetBarPre}{\ensuremath{p = 0.080}}
\providecommand{\rhcaQLgbmHarBarPre}{}\renewcommand{\rhcaQLgbmHarBarPre}{0.1858}
\providecommand{\rhcaMseLgbmHarBarPre}{}\renewcommand{\rhcaMseLgbmHarBarPre}{0.37}
\providecommand{\rhcaMaeLgbmHarBarPre}{}\renewcommand{\rhcaMaeLgbmHarBarPre}{1.43}
\providecommand{\rhcaRsqLgbmHarBarPre}{}\renewcommand{\rhcaRsqLgbmHarBarPre}{53.3}
\providecommand{\rhcaQLgbmFullBarPre}{}\renewcommand{\rhcaQLgbmFullBarPre}{0.1586}
\providecommand{\rhcaMseLgbmFullBarPre}{}\renewcommand{\rhcaMseLgbmFullBarPre}{0.36}
\providecommand{\rhcaMaeLgbmFullBarPre}{}\renewcommand{\rhcaMaeLgbmFullBarPre}{1.36}
\providecommand{\rhcaRsqLgbmFullBarPre}{}\renewcommand{\rhcaRsqLgbmFullBarPre}{55.5}
\providecommand{\rhcaPctLgbmBarPre}{}\renewcommand{\rhcaPctLgbmBarPre}{\ensuremath{-}14.6}
\providecommand{\rhcaDmLgbmBarPre}{}\renewcommand{\rhcaDmLgbmBarPre}{\ensuremath{-}7.91}
\providecommand{\rhcaPLgbmBarPre}{}\renewcommand{\rhcaPLgbmBarPre}{\ensuremath{p < 0.001}}
\providecommand{\rhcaMsePctLgbmBarPre}{}\renewcommand{\rhcaMsePctLgbmBarPre}{\ensuremath{-}4.7}
\providecommand{\rhcaMseDmLgbmBarPre}{}\renewcommand{\rhcaMseDmLgbmBarPre}{\ensuremath{-}2.59}
\providecommand{\rhcaMsePLgbmBarPre}{}\renewcommand{\rhcaMsePLgbmBarPre}{\ensuremath{p = 0.0097}}
\providecommand{\rhcaQXgbHarBarPre}{}\renewcommand{\rhcaQXgbHarBarPre}{0.1864}
\providecommand{\rhcaMseXgbHarBarPre}{}\renewcommand{\rhcaMseXgbHarBarPre}{0.37}
\providecommand{\rhcaMaeXgbHarBarPre}{}\renewcommand{\rhcaMaeXgbHarBarPre}{1.43}
\providecommand{\rhcaRsqXgbHarBarPre}{}\renewcommand{\rhcaRsqXgbHarBarPre}{53.9}
\providecommand{\rhcaQXgbFullBarPre}{}\renewcommand{\rhcaQXgbFullBarPre}{0.1597}
\providecommand{\rhcaMseXgbFullBarPre}{}\renewcommand{\rhcaMseXgbFullBarPre}{0.36}
\providecommand{\rhcaMaeXgbFullBarPre}{}\renewcommand{\rhcaMaeXgbFullBarPre}{1.37}
\providecommand{\rhcaRsqXgbFullBarPre}{}\renewcommand{\rhcaRsqXgbFullBarPre}{55.5}
\providecommand{\rhcaPctXgbBarPre}{}\renewcommand{\rhcaPctXgbBarPre}{\ensuremath{-}14.3}
\providecommand{\rhcaDmXgbBarPre}{}\renewcommand{\rhcaDmXgbBarPre}{\ensuremath{-}7.48}
\providecommand{\rhcaPXgbBarPre}{}\renewcommand{\rhcaPXgbBarPre}{\ensuremath{p < 0.001}}
\providecommand{\rhcaMsePctXgbBarPre}{}\renewcommand{\rhcaMsePctXgbBarPre}{\ensuremath{-}3.4}
\providecommand{\rhcaMseDmXgbBarPre}{}\renewcommand{\rhcaMseDmXgbBarPre}{\ensuremath{-}2.50}
\providecommand{\rhcaMsePXgbBarPre}{}\renewcommand{\rhcaMsePXgbBarPre}{\ensuremath{p = 0.012}}
\providecommand{\rhcaQRfHarBarPre}{}\renewcommand{\rhcaQRfHarBarPre}{0.1865}
\providecommand{\rhcaMseRfHarBarPre}{}\renewcommand{\rhcaMseRfHarBarPre}{0.37}
\providecommand{\rhcaMaeRfHarBarPre}{}\renewcommand{\rhcaMaeRfHarBarPre}{1.43}
\providecommand{\rhcaRsqRfHarBarPre}{}\renewcommand{\rhcaRsqRfHarBarPre}{53.9}
\providecommand{\rhcaQRfFullBarPre}{}\renewcommand{\rhcaQRfFullBarPre}{0.1646}
\providecommand{\rhcaMseRfFullBarPre}{}\renewcommand{\rhcaMseRfFullBarPre}{0.36}
\providecommand{\rhcaMaeRfFullBarPre}{}\renewcommand{\rhcaMaeRfFullBarPre}{1.37}
\providecommand{\rhcaRsqRfFullBarPre}{}\renewcommand{\rhcaRsqRfFullBarPre}{54.6}
\providecommand{\rhcaPctRfBarPre}{}\renewcommand{\rhcaPctRfBarPre}{\ensuremath{-}11.7}
\providecommand{\rhcaDmRfBarPre}{}\renewcommand{\rhcaDmRfBarPre}{\ensuremath{-}7.21}
\providecommand{\rhcaPRfBarPre}{}\renewcommand{\rhcaPRfBarPre}{\ensuremath{p < 0.001}}
\providecommand{\rhcaMsePctRfBarPre}{}\renewcommand{\rhcaMsePctRfBarPre}{\ensuremath{-}1.6}
\providecommand{\rhcaMseDmRfBarPre}{}\renewcommand{\rhcaMseDmRfBarPre}{\ensuremath{-}1.74}
\providecommand{\rhcaMsePRfBarPre}{}\renewcommand{\rhcaMsePRfBarPre}{\ensuremath{p = 0.083}}
\providecommand{\rhcaQOlsHarBarOos}{}\renewcommand{\rhcaQOlsHarBarOos}{0.1933}
\providecommand{\rhcaMseOlsHarBarOos}{}\renewcommand{\rhcaMseOlsHarBarOos}{2.22}
\providecommand{\rhcaMaeOlsHarBarOos}{}\renewcommand{\rhcaMaeOlsHarBarOos}{3.42}
\providecommand{\rhcaRsqOlsHarBarOos}{}\renewcommand{\rhcaRsqOlsHarBarOos}{66.1}
\providecommand{\rhcaQRidgeHarBarOos}{}\renewcommand{\rhcaQRidgeHarBarOos}{0.1933}
\providecommand{\rhcaMseRidgeHarBarOos}{}\renewcommand{\rhcaMseRidgeHarBarOos}{2.22}
\providecommand{\rhcaMaeRidgeHarBarOos}{}\renewcommand{\rhcaMaeRidgeHarBarOos}{3.42}
\providecommand{\rhcaRsqRidgeHarBarOos}{}\renewcommand{\rhcaRsqRidgeHarBarOos}{66.1}
\providecommand{\rhcaQRidgeFullBarOos}{}\renewcommand{\rhcaQRidgeFullBarOos}{0.1852}
\providecommand{\rhcaMseRidgeFullBarOos}{}\renewcommand{\rhcaMseRidgeFullBarOos}{4.93}
\providecommand{\rhcaMaeRidgeFullBarOos}{}\renewcommand{\rhcaMaeRidgeFullBarOos}{3.80}
\providecommand{\rhcaRsqRidgeFullBarOos}{}\renewcommand{\rhcaRsqRidgeFullBarOos}{24.9}
\providecommand{\rhcaPctRidgeBarOos}{}\renewcommand{\rhcaPctRidgeBarOos}{\ensuremath{-}4.2}
\providecommand{\rhcaDmRidgeBarOos}{}\renewcommand{\rhcaDmRidgeBarOos}{\ensuremath{-}2.52}
\providecommand{\rhcaPRidgeBarOos}{}\renewcommand{\rhcaPRidgeBarOos}{\ensuremath{p = 0.012}}
\providecommand{\rhcaMsePctRidgeBarOos}{}\renewcommand{\rhcaMsePctRidgeBarOos}{121.6}
\providecommand{\rhcaMseDmRidgeBarOos}{}\renewcommand{\rhcaMseDmRidgeBarOos}{1.19}
\providecommand{\rhcaMsePRidgeBarOos}{}\renewcommand{\rhcaMsePRidgeBarOos}{\ensuremath{p = 0.24}}
\providecommand{\rhcaQLassoHarBarOos}{}\renewcommand{\rhcaQLassoHarBarOos}{0.1933}
\providecommand{\rhcaMseLassoHarBarOos}{}\renewcommand{\rhcaMseLassoHarBarOos}{2.22}
\providecommand{\rhcaMaeLassoHarBarOos}{}\renewcommand{\rhcaMaeLassoHarBarOos}{3.42}
\providecommand{\rhcaRsqLassoHarBarOos}{}\renewcommand{\rhcaRsqLassoHarBarOos}{66.1}
\providecommand{\rhcaQLassoFullBarOos}{}\renewcommand{\rhcaQLassoFullBarOos}{0.1801}
\providecommand{\rhcaMseLassoFullBarOos}{}\renewcommand{\rhcaMseLassoFullBarOos}{5.31}
\providecommand{\rhcaMaeLassoFullBarOos}{}\renewcommand{\rhcaMaeLassoFullBarOos}{3.86}
\providecommand{\rhcaRsqLassoFullBarOos}{}\renewcommand{\rhcaRsqLassoFullBarOos}{19.1}
\providecommand{\rhcaPctLassoBarOos}{}\renewcommand{\rhcaPctLassoBarOos}{\ensuremath{-}6.8}
\providecommand{\rhcaDmLassoBarOos}{}\renewcommand{\rhcaDmLassoBarOos}{\ensuremath{-}2.70}
\providecommand{\rhcaPLassoBarOos}{}\renewcommand{\rhcaPLassoBarOos}{\ensuremath{p = 0.007}}
\providecommand{\rhcaMsePctLassoBarOos}{}\renewcommand{\rhcaMsePctLassoBarOos}{138.8}
\providecommand{\rhcaMseDmLassoBarOos}{}\renewcommand{\rhcaMseDmLassoBarOos}{1.23}
\providecommand{\rhcaMsePLassoBarOos}{}\renewcommand{\rhcaMsePLassoBarOos}{\ensuremath{p = 0.22}}
\providecommand{\rhcaQEnetHarBarOos}{}\renewcommand{\rhcaQEnetHarBarOos}{0.1933}
\providecommand{\rhcaMseEnetHarBarOos}{}\renewcommand{\rhcaMseEnetHarBarOos}{2.22}
\providecommand{\rhcaMaeEnetHarBarOos}{}\renewcommand{\rhcaMaeEnetHarBarOos}{3.42}
\providecommand{\rhcaRsqEnetHarBarOos}{}\renewcommand{\rhcaRsqEnetHarBarOos}{66.1}
\providecommand{\rhcaQEnetFullBarOos}{}\renewcommand{\rhcaQEnetFullBarOos}{0.1831}
\providecommand{\rhcaMseEnetFullBarOos}{}\renewcommand{\rhcaMseEnetFullBarOos}{5.38}
\providecommand{\rhcaMaeEnetFullBarOos}{}\renewcommand{\rhcaMaeEnetFullBarOos}{3.85}
\providecommand{\rhcaRsqEnetFullBarOos}{}\renewcommand{\rhcaRsqEnetFullBarOos}{18.0}
\providecommand{\rhcaPctEnetBarOos}{}\renewcommand{\rhcaPctEnetBarOos}{\ensuremath{-}5.3}
\providecommand{\rhcaDmEnetBarOos}{}\renewcommand{\rhcaDmEnetBarOos}{\ensuremath{-}3.58}
\providecommand{\rhcaPEnetBarOos}{}\renewcommand{\rhcaPEnetBarOos}{\ensuremath{p < 0.001}}
\providecommand{\rhcaMsePctEnetBarOos}{}\renewcommand{\rhcaMsePctEnetBarOos}{142.1}
\providecommand{\rhcaMseDmEnetBarOos}{}\renewcommand{\rhcaMseDmEnetBarOos}{1.18}
\providecommand{\rhcaMsePEnetBarOos}{}\renewcommand{\rhcaMsePEnetBarOos}{\ensuremath{p = 0.24}}
\providecommand{\rhcaQLgbmHarBarOos}{}\renewcommand{\rhcaQLgbmHarBarOos}{0.1900}
\providecommand{\rhcaMseLgbmHarBarOos}{}\renewcommand{\rhcaMseLgbmHarBarOos}{2.30}
\providecommand{\rhcaMaeLgbmHarBarOos}{}\renewcommand{\rhcaMaeLgbmHarBarOos}{3.43}
\providecommand{\rhcaRsqLgbmHarBarOos}{}\renewcommand{\rhcaRsqLgbmHarBarOos}{65.0}
\providecommand{\rhcaQLgbmFullBarOos}{}\renewcommand{\rhcaQLgbmFullBarOos}{0.1645}
\providecommand{\rhcaMseLgbmFullBarOos}{}\renewcommand{\rhcaMseLgbmFullBarOos}{2.89}
\providecommand{\rhcaMaeLgbmFullBarOos}{}\renewcommand{\rhcaMaeLgbmFullBarOos}{3.55}
\providecommand{\rhcaRsqLgbmFullBarOos}{}\renewcommand{\rhcaRsqLgbmFullBarOos}{56.0}
\providecommand{\rhcaPctLgbmBarOos}{}\renewcommand{\rhcaPctLgbmBarOos}{\ensuremath{-}13.4}
\providecommand{\rhcaDmLgbmBarOos}{}\renewcommand{\rhcaDmLgbmBarOos}{\ensuremath{-}6.42}
\providecommand{\rhcaPLgbmBarOos}{}\renewcommand{\rhcaPLgbmBarOos}{\ensuremath{p < 0.001}}
\providecommand{\rhcaMsePctLgbmBarOos}{}\renewcommand{\rhcaMsePctLgbmBarOos}{25.5}
\providecommand{\rhcaMseDmLgbmBarOos}{}\renewcommand{\rhcaMseDmLgbmBarOos}{1.13}
\providecommand{\rhcaMsePLgbmBarOos}{}\renewcommand{\rhcaMsePLgbmBarOos}{\ensuremath{p = 0.26}}
\providecommand{\rhcaQXgbHarBarOos}{}\renewcommand{\rhcaQXgbHarBarOos}{0.1921}
\providecommand{\rhcaMseXgbHarBarOos}{}\renewcommand{\rhcaMseXgbHarBarOos}{2.26}
\providecommand{\rhcaMaeXgbHarBarOos}{}\renewcommand{\rhcaMaeXgbHarBarOos}{3.42}
\providecommand{\rhcaRsqXgbHarBarOos}{}\renewcommand{\rhcaRsqXgbHarBarOos}{65.6}
\providecommand{\rhcaQXgbFullBarOos}{}\renewcommand{\rhcaQXgbFullBarOos}{0.1634}
\providecommand{\rhcaMseXgbFullBarOos}{}\renewcommand{\rhcaMseXgbFullBarOos}{2.72}
\providecommand{\rhcaMaeXgbFullBarOos}{}\renewcommand{\rhcaMaeXgbFullBarOos}{3.53}
\providecommand{\rhcaRsqXgbFullBarOos}{}\renewcommand{\rhcaRsqXgbFullBarOos}{58.5}
\providecommand{\rhcaPctXgbBarOos}{}\renewcommand{\rhcaPctXgbBarOos}{\ensuremath{-}14.9}
\providecommand{\rhcaDmXgbBarOos}{}\renewcommand{\rhcaDmXgbBarOos}{\ensuremath{-}6.42}
\providecommand{\rhcaPXgbBarOos}{}\renewcommand{\rhcaPXgbBarOos}{\ensuremath{p < 0.001}}
\providecommand{\rhcaMsePctXgbBarOos}{}\renewcommand{\rhcaMsePctXgbBarOos}{20.8}
\providecommand{\rhcaMseDmXgbBarOos}{}\renewcommand{\rhcaMseDmXgbBarOos}{1.15}
\providecommand{\rhcaMsePXgbBarOos}{}\renewcommand{\rhcaMsePXgbBarOos}{\ensuremath{p = 0.25}}
\providecommand{\rhcaQRfHarBarOos}{}\renewcommand{\rhcaQRfHarBarOos}{0.1938}
\providecommand{\rhcaMseRfHarBarOos}{}\renewcommand{\rhcaMseRfHarBarOos}{2.33}
\providecommand{\rhcaMaeRfHarBarOos}{}\renewcommand{\rhcaMaeRfHarBarOos}{3.44}
\providecommand{\rhcaRsqRfHarBarOos}{}\renewcommand{\rhcaRsqRfHarBarOos}{64.5}
\providecommand{\rhcaQRfFullBarOos}{}\renewcommand{\rhcaQRfFullBarOos}{0.1730}
\providecommand{\rhcaMseRfFullBarOos}{}\renewcommand{\rhcaMseRfFullBarOos}{2.65}
\providecommand{\rhcaMaeRfFullBarOos}{}\renewcommand{\rhcaMaeRfFullBarOos}{3.52}
\providecommand{\rhcaRsqRfFullBarOos}{}\renewcommand{\rhcaRsqRfFullBarOos}{59.6}
\providecommand{\rhcaPctRfBarOos}{}\renewcommand{\rhcaPctRfBarOos}{\ensuremath{-}10.7}
\providecommand{\rhcaDmRfBarOos}{}\renewcommand{\rhcaDmRfBarOos}{\ensuremath{-}6.47}
\providecommand{\rhcaPRfBarOos}{}\renewcommand{\rhcaPRfBarOos}{\ensuremath{p < 0.001}}
\providecommand{\rhcaMsePctRfBarOos}{}\renewcommand{\rhcaMsePctRfBarOos}{13.7}
\providecommand{\rhcaMseDmRfBarOos}{}\renewcommand{\rhcaMseDmRfBarOos}{1.19}
\providecommand{\rhcaMsePRfBarOos}{}\renewcommand{\rhcaMsePRfBarOos}{\ensuremath{p = 0.23}}
% --- calibration, Pre / Oos, every bar and 16:00: Mr mean of RV/sigma2, Rm ratio of means, Under %, MzA (x 10^6), MzB, MzBse, MzRsq (levels MZ), MzlA, MzlB, MzlBse, MzlRsq (logs MZ), ShBias / ShSlope / ShDist (Theil shares of the MSE, %), Ql / Qd (QLIKE level / dispersion terms), OverSe (% of SE from over-predictions)
\providecommand{\rhcaMrRidgeHarBarPre}{}\renewcommand{\rhcaMrRidgeHarBarPre}{1.078}
\providecommand{\rhcaRmRidgeHarBarPre}{}\renewcommand{\rhcaRmRidgeHarBarPre}{1.057}
\providecommand{\rhcaUnderRidgeHarBarPre}{}\renewcommand{\rhcaUnderRidgeHarBarPre}{39.9}
\providecommand{\rhcaMzARidgeHarBarPre}{}\renewcommand{\rhcaMzARidgeHarBarPre}{0.29}
\providecommand{\rhcaMzBRidgeHarBarPre}{}\renewcommand{\rhcaMzBRidgeHarBarPre}{0.97}
\providecommand{\rhcaMzBseRidgeHarBarPre}{}\renewcommand{\rhcaMzBseRidgeHarBarPre}{0.11}
\providecommand{\rhcaMzRsqRidgeHarBarPre}{}\renewcommand{\rhcaMzRsqRidgeHarBarPre}{0.48}
\providecommand{\rhcaMzlARidgeHarBarPre}{}\renewcommand{\rhcaMzlARidgeHarBarPre}{0.04}
\providecommand{\rhcaMzlBRidgeHarBarPre}{}\renewcommand{\rhcaMzlBRidgeHarBarPre}{1.01}
\providecommand{\rhcaMzlBseRidgeHarBarPre}{}\renewcommand{\rhcaMzlBseRidgeHarBarPre}{0.01}
\providecommand{\rhcaMzlRsqRidgeHarBarPre}{}\renewcommand{\rhcaMzlRsqRidgeHarBarPre}{0.76}
\providecommand{\rhcaShBiasRidgeHarBarPre}{}\renewcommand{\rhcaShBiasRidgeHarBarPre}{0.1}
\providecommand{\rhcaShSlopeRidgeHarBarPre}{}\renewcommand{\rhcaShSlopeRidgeHarBarPre}{0.1}
\providecommand{\rhcaShDistRidgeHarBarPre}{}\renewcommand{\rhcaShDistRidgeHarBarPre}{99.8}
\providecommand{\rhcaQlRidgeHarBarPre}{}\renewcommand{\rhcaQlRidgeHarBarPre}{0.0029}
\providecommand{\rhcaQdRidgeHarBarPre}{}\renewcommand{\rhcaQdRidgeHarBarPre}{0.1835}
\providecommand{\rhcaOverSeRidgeHarBarPre}{}\renewcommand{\rhcaOverSeRidgeHarBarPre}{20.0}
\providecommand{\rhcaMrRidgeFullBarPre}{}\renewcommand{\rhcaMrRidgeFullBarPre}{1.096}
\providecommand{\rhcaRmRidgeFullBarPre}{}\renewcommand{\rhcaRmRidgeFullBarPre}{1.049}
\providecommand{\rhcaUnderRidgeFullBarPre}{}\renewcommand{\rhcaUnderRidgeFullBarPre}{41.6}
\providecommand{\rhcaMzARidgeFullBarPre}{}\renewcommand{\rhcaMzARidgeFullBarPre}{0.17}
\providecommand{\rhcaMzBRidgeFullBarPre}{}\renewcommand{\rhcaMzBRidgeFullBarPre}{1.00}
\providecommand{\rhcaMzBseRidgeFullBarPre}{}\renewcommand{\rhcaMzBseRidgeFullBarPre}{0.10}
\providecommand{\rhcaMzRsqRidgeFullBarPre}{}\renewcommand{\rhcaMzRsqRidgeFullBarPre}{0.53}
\providecommand{\rhcaMzlARidgeFullBarPre}{}\renewcommand{\rhcaMzlARidgeFullBarPre}{\ensuremath{-}0.22}
\providecommand{\rhcaMzlBRidgeFullBarPre}{}\renewcommand{\rhcaMzlBRidgeFullBarPre}{0.99}
\providecommand{\rhcaMzlBseRidgeFullBarPre}{}\renewcommand{\rhcaMzlBseRidgeFullBarPre}{0.01}
\providecommand{\rhcaMzlRsqRidgeFullBarPre}{}\renewcommand{\rhcaMzlRsqRidgeFullBarPre}{0.78}
\providecommand{\rhcaShBiasRidgeFullBarPre}{}\renewcommand{\rhcaShBiasRidgeFullBarPre}{0.1}
\providecommand{\rhcaShSlopeRidgeFullBarPre}{}\renewcommand{\rhcaShSlopeRidgeFullBarPre}{0.0}
\providecommand{\rhcaShDistRidgeFullBarPre}{}\renewcommand{\rhcaShDistRidgeFullBarPre}{99.9}
\providecommand{\rhcaQlRidgeFullBarPre}{}\renewcommand{\rhcaQlRidgeFullBarPre}{0.0043}
\providecommand{\rhcaQdRidgeFullBarPre}{}\renewcommand{\rhcaQdRidgeFullBarPre}{0.1683}
\providecommand{\rhcaOverSeRidgeFullBarPre}{}\renewcommand{\rhcaOverSeRidgeFullBarPre}{19.6}
\providecommand{\rhcaMrLassoHarBarPre}{}\renewcommand{\rhcaMrLassoHarBarPre}{1.078}
\providecommand{\rhcaRmLassoHarBarPre}{}\renewcommand{\rhcaRmLassoHarBarPre}{1.057}
\providecommand{\rhcaUnderLassoHarBarPre}{}\renewcommand{\rhcaUnderLassoHarBarPre}{39.9}
\providecommand{\rhcaMzALassoHarBarPre}{}\renewcommand{\rhcaMzALassoHarBarPre}{0.29}
\providecommand{\rhcaMzBLassoHarBarPre}{}\renewcommand{\rhcaMzBLassoHarBarPre}{0.97}
\providecommand{\rhcaMzBseLassoHarBarPre}{}\renewcommand{\rhcaMzBseLassoHarBarPre}{0.11}
\providecommand{\rhcaMzRsqLassoHarBarPre}{}\renewcommand{\rhcaMzRsqLassoHarBarPre}{0.48}
\providecommand{\rhcaMzlALassoHarBarPre}{}\renewcommand{\rhcaMzlALassoHarBarPre}{0.04}
\providecommand{\rhcaMzlBLassoHarBarPre}{}\renewcommand{\rhcaMzlBLassoHarBarPre}{1.01}
\providecommand{\rhcaMzlBseLassoHarBarPre}{}\renewcommand{\rhcaMzlBseLassoHarBarPre}{0.01}
\providecommand{\rhcaMzlRsqLassoHarBarPre}{}\renewcommand{\rhcaMzlRsqLassoHarBarPre}{0.76}
\providecommand{\rhcaShBiasLassoHarBarPre}{}\renewcommand{\rhcaShBiasLassoHarBarPre}{0.1}
\providecommand{\rhcaShSlopeLassoHarBarPre}{}\renewcommand{\rhcaShSlopeLassoHarBarPre}{0.1}
\providecommand{\rhcaShDistLassoHarBarPre}{}\renewcommand{\rhcaShDistLassoHarBarPre}{99.8}
\providecommand{\rhcaQlLassoHarBarPre}{}\renewcommand{\rhcaQlLassoHarBarPre}{0.0029}
\providecommand{\rhcaQdLassoHarBarPre}{}\renewcommand{\rhcaQdLassoHarBarPre}{0.1835}
\providecommand{\rhcaOverSeLassoHarBarPre}{}\renewcommand{\rhcaOverSeLassoHarBarPre}{20.0}
\providecommand{\rhcaMrLassoFullBarPre}{}\renewcommand{\rhcaMrLassoFullBarPre}{1.092}
\providecommand{\rhcaRmLassoFullBarPre}{}\renewcommand{\rhcaRmLassoFullBarPre}{1.045}
\providecommand{\rhcaUnderLassoFullBarPre}{}\renewcommand{\rhcaUnderLassoFullBarPre}{42.1}
\providecommand{\rhcaMzALassoFullBarPre}{}\renewcommand{\rhcaMzALassoFullBarPre}{0.16}
\providecommand{\rhcaMzBLassoFullBarPre}{}\renewcommand{\rhcaMzBLassoFullBarPre}{0.99}
\providecommand{\rhcaMzBseLassoFullBarPre}{}\renewcommand{\rhcaMzBseLassoFullBarPre}{0.10}
\providecommand{\rhcaMzRsqLassoFullBarPre}{}\renewcommand{\rhcaMzRsqLassoFullBarPre}{0.54}
\providecommand{\rhcaMzlALassoFullBarPre}{}\renewcommand{\rhcaMzlALassoFullBarPre}{\ensuremath{-}0.25}
\providecommand{\rhcaMzlBLassoFullBarPre}{}\renewcommand{\rhcaMzlBLassoFullBarPre}{0.99}
\providecommand{\rhcaMzlBseLassoFullBarPre}{}\renewcommand{\rhcaMzlBseLassoFullBarPre}{0.01}
\providecommand{\rhcaMzlRsqLassoFullBarPre}{}\renewcommand{\rhcaMzlRsqLassoFullBarPre}{0.78}
\providecommand{\rhcaShBiasLassoFullBarPre}{}\renewcommand{\rhcaShBiasLassoFullBarPre}{0.1}
\providecommand{\rhcaShSlopeLassoFullBarPre}{}\renewcommand{\rhcaShSlopeLassoFullBarPre}{0.0}
\providecommand{\rhcaShDistLassoFullBarPre}{}\renewcommand{\rhcaShDistLassoFullBarPre}{99.9}
\providecommand{\rhcaQlLassoFullBarPre}{}\renewcommand{\rhcaQlLassoFullBarPre}{0.0040}
\providecommand{\rhcaQdLassoFullBarPre}{}\renewcommand{\rhcaQdLassoFullBarPre}{0.1627}
\providecommand{\rhcaOverSeLassoFullBarPre}{}\renewcommand{\rhcaOverSeLassoFullBarPre}{20.1}
\providecommand{\rhcaMrEnetHarBarPre}{}\renewcommand{\rhcaMrEnetHarBarPre}{1.078}
\providecommand{\rhcaRmEnetHarBarPre}{}\renewcommand{\rhcaRmEnetHarBarPre}{1.057}
\providecommand{\rhcaUnderEnetHarBarPre}{}\renewcommand{\rhcaUnderEnetHarBarPre}{39.9}
\providecommand{\rhcaMzAEnetHarBarPre}{}\renewcommand{\rhcaMzAEnetHarBarPre}{0.29}
\providecommand{\rhcaMzBEnetHarBarPre}{}\renewcommand{\rhcaMzBEnetHarBarPre}{0.97}
\providecommand{\rhcaMzBseEnetHarBarPre}{}\renewcommand{\rhcaMzBseEnetHarBarPre}{0.11}
\providecommand{\rhcaMzRsqEnetHarBarPre}{}\renewcommand{\rhcaMzRsqEnetHarBarPre}{0.48}
\providecommand{\rhcaMzlAEnetHarBarPre}{}\renewcommand{\rhcaMzlAEnetHarBarPre}{0.04}
\providecommand{\rhcaMzlBEnetHarBarPre}{}\renewcommand{\rhcaMzlBEnetHarBarPre}{1.01}
\providecommand{\rhcaMzlBseEnetHarBarPre}{}\renewcommand{\rhcaMzlBseEnetHarBarPre}{0.01}
\providecommand{\rhcaMzlRsqEnetHarBarPre}{}\renewcommand{\rhcaMzlRsqEnetHarBarPre}{0.76}
\providecommand{\rhcaShBiasEnetHarBarPre}{}\renewcommand{\rhcaShBiasEnetHarBarPre}{0.1}
\providecommand{\rhcaShSlopeEnetHarBarPre}{}\renewcommand{\rhcaShSlopeEnetHarBarPre}{0.1}
\providecommand{\rhcaShDistEnetHarBarPre}{}\renewcommand{\rhcaShDistEnetHarBarPre}{99.8}
\providecommand{\rhcaQlEnetHarBarPre}{}\renewcommand{\rhcaQlEnetHarBarPre}{0.0029}
\providecommand{\rhcaQdEnetHarBarPre}{}\renewcommand{\rhcaQdEnetHarBarPre}{0.1835}
\providecommand{\rhcaOverSeEnetHarBarPre}{}\renewcommand{\rhcaOverSeEnetHarBarPre}{20.0}
\providecommand{\rhcaMrEnetFullBarPre}{}\renewcommand{\rhcaMrEnetFullBarPre}{1.091}
\providecommand{\rhcaRmEnetFullBarPre}{}\renewcommand{\rhcaRmEnetFullBarPre}{1.049}
\providecommand{\rhcaUnderEnetFullBarPre}{}\renewcommand{\rhcaUnderEnetFullBarPre}{41.5}
\providecommand{\rhcaMzAEnetFullBarPre}{}\renewcommand{\rhcaMzAEnetFullBarPre}{0.16}
\providecommand{\rhcaMzBEnetFullBarPre}{}\renewcommand{\rhcaMzBEnetFullBarPre}{1.00}
\providecommand{\rhcaMzBseEnetFullBarPre}{}\renewcommand{\rhcaMzBseEnetFullBarPre}{0.10}
\providecommand{\rhcaMzRsqEnetFullBarPre}{}\renewcommand{\rhcaMzRsqEnetFullBarPre}{0.53}
\providecommand{\rhcaMzlAEnetFullBarPre}{}\renewcommand{\rhcaMzlAEnetFullBarPre}{\ensuremath{-}0.16}
\providecommand{\rhcaMzlBEnetFullBarPre}{}\renewcommand{\rhcaMzlBEnetFullBarPre}{0.99}
\providecommand{\rhcaMzlBseEnetFullBarPre}{}\renewcommand{\rhcaMzlBseEnetFullBarPre}{0.01}
\providecommand{\rhcaMzlRsqEnetFullBarPre}{}\renewcommand{\rhcaMzlRsqEnetFullBarPre}{0.78}
\providecommand{\rhcaShBiasEnetFullBarPre}{}\renewcommand{\rhcaShBiasEnetFullBarPre}{0.1}
\providecommand{\rhcaShSlopeEnetFullBarPre}{}\renewcommand{\rhcaShSlopeEnetFullBarPre}{0.0}
\providecommand{\rhcaShDistEnetFullBarPre}{}\renewcommand{\rhcaShDistEnetFullBarPre}{99.9}
\providecommand{\rhcaQlEnetFullBarPre}{}\renewcommand{\rhcaQlEnetFullBarPre}{0.0039}
\providecommand{\rhcaQdEnetFullBarPre}{}\renewcommand{\rhcaQdEnetFullBarPre}{0.1681}
\providecommand{\rhcaOverSeEnetFullBarPre}{}\renewcommand{\rhcaOverSeEnetFullBarPre}{19.7}
\providecommand{\rhcaMrLgbmHarBarPre}{}\renewcommand{\rhcaMrLgbmHarBarPre}{1.083}
\providecommand{\rhcaRmLgbmHarBarPre}{}\renewcommand{\rhcaRmLgbmHarBarPre}{1.075}
\providecommand{\rhcaUnderLgbmHarBarPre}{}\renewcommand{\rhcaUnderLgbmHarBarPre}{40.2}
\providecommand{\rhcaMzALgbmHarBarPre}{}\renewcommand{\rhcaMzALgbmHarBarPre}{0.27}
\providecommand{\rhcaMzBLgbmHarBarPre}{}\renewcommand{\rhcaMzBLgbmHarBarPre}{0.99}
\providecommand{\rhcaMzBseLgbmHarBarPre}{}\renewcommand{\rhcaMzBseLgbmHarBarPre}{0.07}
\providecommand{\rhcaMzRsqLgbmHarBarPre}{}\renewcommand{\rhcaMzRsqLgbmHarBarPre}{0.37}
\providecommand{\rhcaMzlALgbmHarBarPre}{}\renewcommand{\rhcaMzlALgbmHarBarPre}{0.04}
\providecommand{\rhcaMzlBLgbmHarBarPre}{}\renewcommand{\rhcaMzlBLgbmHarBarPre}{1.01}
\providecommand{\rhcaMzlBseLgbmHarBarPre}{}\renewcommand{\rhcaMzlBseLgbmHarBarPre}{0.01}
\providecommand{\rhcaMzlRsqLgbmHarBarPre}{}\renewcommand{\rhcaMzlRsqLgbmHarBarPre}{0.76}
\providecommand{\rhcaShBiasLgbmHarBarPre}{}\renewcommand{\rhcaShBiasLgbmHarBarPre}{0.1}
\providecommand{\rhcaShSlopeLgbmHarBarPre}{}\renewcommand{\rhcaShSlopeLgbmHarBarPre}{0.0}
\providecommand{\rhcaShDistLgbmHarBarPre}{}\renewcommand{\rhcaShDistLgbmHarBarPre}{99.8}
\providecommand{\rhcaQlLgbmHarBarPre}{}\renewcommand{\rhcaQlLgbmHarBarPre}{0.0032}
\providecommand{\rhcaQdLgbmHarBarPre}{}\renewcommand{\rhcaQdLgbmHarBarPre}{0.1826}
\providecommand{\rhcaOverSeLgbmHarBarPre}{}\renewcommand{\rhcaOverSeLgbmHarBarPre}{10.3}
\providecommand{\rhcaMrLgbmFullBarPre}{}\renewcommand{\rhcaMrLgbmFullBarPre}{1.069}
\providecommand{\rhcaRmLgbmFullBarPre}{}\renewcommand{\rhcaRmLgbmFullBarPre}{1.060}
\providecommand{\rhcaUnderLgbmFullBarPre}{}\renewcommand{\rhcaUnderLgbmFullBarPre}{41.0}
\providecommand{\rhcaMzALgbmFullBarPre}{}\renewcommand{\rhcaMzALgbmFullBarPre}{0.29}
\providecommand{\rhcaMzBLgbmFullBarPre}{}\renewcommand{\rhcaMzBLgbmFullBarPre}{0.97}
\providecommand{\rhcaMzBseLgbmFullBarPre}{}\renewcommand{\rhcaMzBseLgbmFullBarPre}{0.06}
\providecommand{\rhcaMzRsqLgbmFullBarPre}{}\renewcommand{\rhcaMzRsqLgbmFullBarPre}{0.40}
\providecommand{\rhcaMzlALgbmFullBarPre}{}\renewcommand{\rhcaMzlALgbmFullBarPre}{\ensuremath{-}0.03}
\providecommand{\rhcaMzlBLgbmFullBarPre}{}\renewcommand{\rhcaMzlBLgbmFullBarPre}{1.00}
\providecommand{\rhcaMzlBseLgbmFullBarPre}{}\renewcommand{\rhcaMzlBseLgbmFullBarPre}{0.01}
\providecommand{\rhcaMzlRsqLgbmFullBarPre}{}\renewcommand{\rhcaMzlRsqLgbmFullBarPre}{0.79}
\providecommand{\rhcaShBiasLgbmFullBarPre}{}\renewcommand{\rhcaShBiasLgbmFullBarPre}{0.1}
\providecommand{\rhcaShSlopeLgbmFullBarPre}{}\renewcommand{\rhcaShSlopeLgbmFullBarPre}{0.1}
\providecommand{\rhcaShDistLgbmFullBarPre}{}\renewcommand{\rhcaShDistLgbmFullBarPre}{99.8}
\providecommand{\rhcaQlLgbmFullBarPre}{}\renewcommand{\rhcaQlLgbmFullBarPre}{0.0023}
\providecommand{\rhcaQdLgbmFullBarPre}{}\renewcommand{\rhcaQdLgbmFullBarPre}{0.1563}
\providecommand{\rhcaOverSeLgbmFullBarPre}{}\renewcommand{\rhcaOverSeLgbmFullBarPre}{12.0}
\providecommand{\rhcaMrXgbHarBarPre}{}\renewcommand{\rhcaMrXgbHarBarPre}{1.083}
\providecommand{\rhcaRmXgbHarBarPre}{}\renewcommand{\rhcaRmXgbHarBarPre}{1.073}
\providecommand{\rhcaUnderXgbHarBarPre}{}\renewcommand{\rhcaUnderXgbHarBarPre}{40.2}
\providecommand{\rhcaMzAXgbHarBarPre}{}\renewcommand{\rhcaMzAXgbHarBarPre}{0.31}
\providecommand{\rhcaMzBXgbHarBarPre}{}\renewcommand{\rhcaMzBXgbHarBarPre}{0.97}
\providecommand{\rhcaMzBseXgbHarBarPre}{}\renewcommand{\rhcaMzBseXgbHarBarPre}{0.09}
\providecommand{\rhcaMzRsqXgbHarBarPre}{}\renewcommand{\rhcaMzRsqXgbHarBarPre}{0.38}
\providecommand{\rhcaMzlAXgbHarBarPre}{}\renewcommand{\rhcaMzlAXgbHarBarPre}{0.06}
\providecommand{\rhcaMzlBXgbHarBarPre}{}\renewcommand{\rhcaMzlBXgbHarBarPre}{1.01}
\providecommand{\rhcaMzlBseXgbHarBarPre}{}\renewcommand{\rhcaMzlBseXgbHarBarPre}{0.01}
\providecommand{\rhcaMzlRsqXgbHarBarPre}{}\renewcommand{\rhcaMzlRsqXgbHarBarPre}{0.76}
\providecommand{\rhcaShBiasXgbHarBarPre}{}\renewcommand{\rhcaShBiasXgbHarBarPre}{0.1}
\providecommand{\rhcaShSlopeXgbHarBarPre}{}\renewcommand{\rhcaShSlopeXgbHarBarPre}{0.0}
\providecommand{\rhcaShDistXgbHarBarPre}{}\renewcommand{\rhcaShDistXgbHarBarPre}{99.8}
\providecommand{\rhcaQlXgbHarBarPre}{}\renewcommand{\rhcaQlXgbHarBarPre}{0.0032}
\providecommand{\rhcaQdXgbHarBarPre}{}\renewcommand{\rhcaQdXgbHarBarPre}{0.1832}
\providecommand{\rhcaOverSeXgbHarBarPre}{}\renewcommand{\rhcaOverSeXgbHarBarPre}{12.2}
\providecommand{\rhcaMrXgbFullBarPre}{}\renewcommand{\rhcaMrXgbFullBarPre}{1.067}
\providecommand{\rhcaRmXgbFullBarPre}{}\renewcommand{\rhcaRmXgbFullBarPre}{1.055}
\providecommand{\rhcaUnderXgbFullBarPre}{}\renewcommand{\rhcaUnderXgbFullBarPre}{40.8}
\providecommand{\rhcaMzAXgbFullBarPre}{}\renewcommand{\rhcaMzAXgbFullBarPre}{0.40}
\providecommand{\rhcaMzBXgbFullBarPre}{}\renewcommand{\rhcaMzBXgbFullBarPre}{0.93}
\providecommand{\rhcaMzBseXgbFullBarPre}{}\renewcommand{\rhcaMzBseXgbFullBarPre}{0.09}
\providecommand{\rhcaMzRsqXgbFullBarPre}{}\renewcommand{\rhcaMzRsqXgbFullBarPre}{0.40}
\providecommand{\rhcaMzlAXgbFullBarPre}{}\renewcommand{\rhcaMzlAXgbFullBarPre}{0.00}
\providecommand{\rhcaMzlBXgbFullBarPre}{}\renewcommand{\rhcaMzlBXgbFullBarPre}{1.01}
\providecommand{\rhcaMzlBseXgbFullBarPre}{}\renewcommand{\rhcaMzlBseXgbFullBarPre}{0.01}
\providecommand{\rhcaMzlRsqXgbFullBarPre}{}\renewcommand{\rhcaMzlRsqXgbFullBarPre}{0.79}
\providecommand{\rhcaShBiasXgbFullBarPre}{}\renewcommand{\rhcaShBiasXgbFullBarPre}{0.1}
\providecommand{\rhcaShSlopeXgbFullBarPre}{}\renewcommand{\rhcaShSlopeXgbFullBarPre}{0.4}
\providecommand{\rhcaShDistXgbFullBarPre}{}\renewcommand{\rhcaShDistXgbFullBarPre}{99.5}
\providecommand{\rhcaQlXgbFullBarPre}{}\renewcommand{\rhcaQlXgbFullBarPre}{0.0021}
\providecommand{\rhcaQdXgbFullBarPre}{}\renewcommand{\rhcaQdXgbFullBarPre}{0.1575}
\providecommand{\rhcaOverSeXgbFullBarPre}{}\renewcommand{\rhcaOverSeXgbFullBarPre}{15.1}
\providecommand{\rhcaMrRfHarBarPre}{}\renewcommand{\rhcaMrRfHarBarPre}{1.085}
\providecommand{\rhcaRmRfHarBarPre}{}\renewcommand{\rhcaRmRfHarBarPre}{1.068}
\providecommand{\rhcaUnderRfHarBarPre}{}\renewcommand{\rhcaUnderRfHarBarPre}{40.1}
\providecommand{\rhcaMzARfHarBarPre}{}\renewcommand{\rhcaMzARfHarBarPre}{0.33}
\providecommand{\rhcaMzBRfHarBarPre}{}\renewcommand{\rhcaMzBRfHarBarPre}{0.96}
\providecommand{\rhcaMzBseRfHarBarPre}{}\renewcommand{\rhcaMzBseRfHarBarPre}{0.10}
\providecommand{\rhcaMzRsqRfHarBarPre}{}\renewcommand{\rhcaMzRsqRfHarBarPre}{0.38}
\providecommand{\rhcaMzlARfHarBarPre}{}\renewcommand{\rhcaMzlARfHarBarPre}{\ensuremath{-}0.04}
\providecommand{\rhcaMzlBRfHarBarPre}{}\renewcommand{\rhcaMzlBRfHarBarPre}{1.00}
\providecommand{\rhcaMzlBseRfHarBarPre}{}\renewcommand{\rhcaMzlBseRfHarBarPre}{0.01}
\providecommand{\rhcaMzlRsqRfHarBarPre}{}\renewcommand{\rhcaMzlRsqRfHarBarPre}{0.76}
\providecommand{\rhcaShBiasRfHarBarPre}{}\renewcommand{\rhcaShBiasRfHarBarPre}{0.1}
\providecommand{\rhcaShSlopeRfHarBarPre}{}\renewcommand{\rhcaShSlopeRfHarBarPre}{0.1}
\providecommand{\rhcaShDistRfHarBarPre}{}\renewcommand{\rhcaShDistRfHarBarPre}{99.8}
\providecommand{\rhcaQlRfHarBarPre}{}\renewcommand{\rhcaQlRfHarBarPre}{0.0034}
\providecommand{\rhcaQdRfHarBarPre}{}\renewcommand{\rhcaQdRfHarBarPre}{0.1831}
\providecommand{\rhcaOverSeRfHarBarPre}{}\renewcommand{\rhcaOverSeRfHarBarPre}{12.8}
\providecommand{\rhcaMrRfFullBarPre}{}\renewcommand{\rhcaMrRfFullBarPre}{1.078}
\providecommand{\rhcaRmRfFullBarPre}{}\renewcommand{\rhcaRmRfFullBarPre}{1.057}
\providecommand{\rhcaUnderRfFullBarPre}{}\renewcommand{\rhcaUnderRfFullBarPre}{41.1}
\providecommand{\rhcaMzARfFullBarPre}{}\renewcommand{\rhcaMzARfFullBarPre}{0.43}
\providecommand{\rhcaMzBRfFullBarPre}{}\renewcommand{\rhcaMzBRfFullBarPre}{0.92}
\providecommand{\rhcaMzBseRfFullBarPre}{}\renewcommand{\rhcaMzBseRfFullBarPre}{0.09}
\providecommand{\rhcaMzRsqRfFullBarPre}{}\renewcommand{\rhcaMzRsqRfFullBarPre}{0.39}
\providecommand{\rhcaMzlARfFullBarPre}{}\renewcommand{\rhcaMzlARfFullBarPre}{\ensuremath{-}0.11}
\providecommand{\rhcaMzlBRfFullBarPre}{}\renewcommand{\rhcaMzlBRfFullBarPre}{1.00}
\providecommand{\rhcaMzlBseRfFullBarPre}{}\renewcommand{\rhcaMzlBseRfFullBarPre}{0.01}
\providecommand{\rhcaMzlRsqRfFullBarPre}{}\renewcommand{\rhcaMzlRsqRfFullBarPre}{0.79}
\providecommand{\rhcaShBiasRfFullBarPre}{}\renewcommand{\rhcaShBiasRfFullBarPre}{0.1}
\providecommand{\rhcaShSlopeRfFullBarPre}{}\renewcommand{\rhcaShSlopeRfFullBarPre}{0.5}
\providecommand{\rhcaShDistRfFullBarPre}{}\renewcommand{\rhcaShDistRfFullBarPre}{99.5}
\providecommand{\rhcaQlRfFullBarPre}{}\renewcommand{\rhcaQlRfFullBarPre}{0.0029}
\providecommand{\rhcaQdRfFullBarPre}{}\renewcommand{\rhcaQdRfFullBarPre}{0.1617}
\providecommand{\rhcaOverSeRfFullBarPre}{}\renewcommand{\rhcaOverSeRfFullBarPre}{14.7}
\providecommand{\rhcaMrRidgeHarClosePre}{}\renewcommand{\rhcaMrRidgeHarClosePre}{1.121}
\providecommand{\rhcaRmRidgeHarClosePre}{}\renewcommand{\rhcaRmRidgeHarClosePre}{1.097}
\providecommand{\rhcaUnderRidgeHarClosePre}{}\renewcommand{\rhcaUnderRidgeHarClosePre}{46.7}
\providecommand{\rhcaMzARidgeHarClosePre}{}\renewcommand{\rhcaMzARidgeHarClosePre}{\ensuremath{-}0.16}
\providecommand{\rhcaMzBRidgeHarClosePre}{}\renewcommand{\rhcaMzBRidgeHarClosePre}{1.14}
\providecommand{\rhcaMzBseRidgeHarClosePre}{}\renewcommand{\rhcaMzBseRidgeHarClosePre}{0.10}
\providecommand{\rhcaMzRsqRidgeHarClosePre}{}\renewcommand{\rhcaMzRsqRidgeHarClosePre}{0.64}
\providecommand{\rhcaMzlARidgeHarClosePre}{}\renewcommand{\rhcaMzlARidgeHarClosePre}{\ensuremath{-}0.68}
\providecommand{\rhcaMzlBRidgeHarClosePre}{}\renewcommand{\rhcaMzlBRidgeHarClosePre}{0.95}
\providecommand{\rhcaMzlBseRidgeHarClosePre}{}\renewcommand{\rhcaMzlBseRidgeHarClosePre}{0.02}
\providecommand{\rhcaMzlRsqRidgeHarClosePre}{}\renewcommand{\rhcaMzlRsqRidgeHarClosePre}{0.78}
\providecommand{\rhcaShBiasRidgeHarClosePre}{}\renewcommand{\rhcaShBiasRidgeHarClosePre}{0.5}
\providecommand{\rhcaShSlopeRidgeHarClosePre}{}\renewcommand{\rhcaShSlopeRidgeHarClosePre}{2.5}
\providecommand{\rhcaShDistRidgeHarClosePre}{}\renewcommand{\rhcaShDistRidgeHarClosePre}{97.0}
\providecommand{\rhcaQlRidgeHarClosePre}{}\renewcommand{\rhcaQlRidgeHarClosePre}{0.0068}
\providecommand{\rhcaQdRidgeHarClosePre}{}\renewcommand{\rhcaQdRidgeHarClosePre}{0.1297}
\providecommand{\rhcaOverSeRidgeHarClosePre}{}\renewcommand{\rhcaOverSeRidgeHarClosePre}{13.1}
\providecommand{\rhcaMrRidgeFullClosePre}{}\renewcommand{\rhcaMrRidgeFullClosePre}{1.137}
\providecommand{\rhcaRmRidgeFullClosePre}{}\renewcommand{\rhcaRmRidgeFullClosePre}{1.061}
\providecommand{\rhcaUnderRidgeFullClosePre}{}\renewcommand{\rhcaUnderRidgeFullClosePre}{47.2}
\providecommand{\rhcaMzARidgeFullClosePre}{}\renewcommand{\rhcaMzARidgeFullClosePre}{\ensuremath{-}0.02}
\providecommand{\rhcaMzBRidgeFullClosePre}{}\renewcommand{\rhcaMzBRidgeFullClosePre}{1.07}
\providecommand{\rhcaMzBseRidgeFullClosePre}{}\renewcommand{\rhcaMzBseRidgeFullClosePre}{0.09}
\providecommand{\rhcaMzRsqRidgeFullClosePre}{}\renewcommand{\rhcaMzRsqRidgeFullClosePre}{0.68}
\providecommand{\rhcaMzlARidgeFullClosePre}{}\renewcommand{\rhcaMzlARidgeFullClosePre}{\ensuremath{-}0.93}
\providecommand{\rhcaMzlBRidgeFullClosePre}{}\renewcommand{\rhcaMzlBRidgeFullClosePre}{0.93}
\providecommand{\rhcaMzlBseRidgeFullClosePre}{}\renewcommand{\rhcaMzlBseRidgeFullClosePre}{0.02}
\providecommand{\rhcaMzlRsqRidgeFullClosePre}{}\renewcommand{\rhcaMzlRsqRidgeFullClosePre}{0.80}
\providecommand{\rhcaShBiasRidgeFullClosePre}{}\renewcommand{\rhcaShBiasRidgeFullClosePre}{0.2}
\providecommand{\rhcaShSlopeRidgeFullClosePre}{}\renewcommand{\rhcaShSlopeRidgeFullClosePre}{0.8}
\providecommand{\rhcaShDistRidgeFullClosePre}{}\renewcommand{\rhcaShDistRidgeFullClosePre}{99.0}
\providecommand{\rhcaQlRidgeFullClosePre}{}\renewcommand{\rhcaQlRidgeFullClosePre}{0.0086}
\providecommand{\rhcaQdRidgeFullClosePre}{}\renewcommand{\rhcaQdRidgeFullClosePre}{0.1222}
\providecommand{\rhcaOverSeRidgeFullClosePre}{}\renewcommand{\rhcaOverSeRidgeFullClosePre}{18.0}
\providecommand{\rhcaMrLassoHarClosePre}{}\renewcommand{\rhcaMrLassoHarClosePre}{1.121}
\providecommand{\rhcaRmLassoHarClosePre}{}\renewcommand{\rhcaRmLassoHarClosePre}{1.097}
\providecommand{\rhcaUnderLassoHarClosePre}{}\renewcommand{\rhcaUnderLassoHarClosePre}{46.7}
\providecommand{\rhcaMzALassoHarClosePre}{}\renewcommand{\rhcaMzALassoHarClosePre}{\ensuremath{-}0.16}
\providecommand{\rhcaMzBLassoHarClosePre}{}\renewcommand{\rhcaMzBLassoHarClosePre}{1.14}
\providecommand{\rhcaMzBseLassoHarClosePre}{}\renewcommand{\rhcaMzBseLassoHarClosePre}{0.10}
\providecommand{\rhcaMzRsqLassoHarClosePre}{}\renewcommand{\rhcaMzRsqLassoHarClosePre}{0.64}
\providecommand{\rhcaMzlALassoHarClosePre}{}\renewcommand{\rhcaMzlALassoHarClosePre}{\ensuremath{-}0.68}
\providecommand{\rhcaMzlBLassoHarClosePre}{}\renewcommand{\rhcaMzlBLassoHarClosePre}{0.95}
\providecommand{\rhcaMzlBseLassoHarClosePre}{}\renewcommand{\rhcaMzlBseLassoHarClosePre}{0.02}
\providecommand{\rhcaMzlRsqLassoHarClosePre}{}\renewcommand{\rhcaMzlRsqLassoHarClosePre}{0.78}
\providecommand{\rhcaShBiasLassoHarClosePre}{}\renewcommand{\rhcaShBiasLassoHarClosePre}{0.5}
\providecommand{\rhcaShSlopeLassoHarClosePre}{}\renewcommand{\rhcaShSlopeLassoHarClosePre}{2.5}
\providecommand{\rhcaShDistLassoHarClosePre}{}\renewcommand{\rhcaShDistLassoHarClosePre}{97.0}
\providecommand{\rhcaQlLassoHarClosePre}{}\renewcommand{\rhcaQlLassoHarClosePre}{0.0068}
\providecommand{\rhcaQdLassoHarClosePre}{}\renewcommand{\rhcaQdLassoHarClosePre}{0.1297}
\providecommand{\rhcaOverSeLassoHarClosePre}{}\renewcommand{\rhcaOverSeLassoHarClosePre}{13.1}
\providecommand{\rhcaMrLassoFullClosePre}{}\renewcommand{\rhcaMrLassoFullClosePre}{1.148}
\providecommand{\rhcaRmLassoFullClosePre}{}\renewcommand{\rhcaRmLassoFullClosePre}{1.059}
\providecommand{\rhcaUnderLassoFullClosePre}{}\renewcommand{\rhcaUnderLassoFullClosePre}{48.6}
\providecommand{\rhcaMzALassoFullClosePre}{}\renewcommand{\rhcaMzALassoFullClosePre}{\ensuremath{-}0.01}
\providecommand{\rhcaMzBLassoFullClosePre}{}\renewcommand{\rhcaMzBLassoFullClosePre}{1.06}
\providecommand{\rhcaMzBseLassoFullClosePre}{}\renewcommand{\rhcaMzBseLassoFullClosePre}{0.09}
\providecommand{\rhcaMzRsqLassoFullClosePre}{}\renewcommand{\rhcaMzRsqLassoFullClosePre}{0.69}
\providecommand{\rhcaMzlALassoFullClosePre}{}\renewcommand{\rhcaMzlALassoFullClosePre}{\ensuremath{-}1.01}
\providecommand{\rhcaMzlBLassoFullClosePre}{}\renewcommand{\rhcaMzlBLassoFullClosePre}{0.92}
\providecommand{\rhcaMzlBseLassoFullClosePre}{}\renewcommand{\rhcaMzlBseLassoFullClosePre}{0.02}
\providecommand{\rhcaMzlRsqLassoFullClosePre}{}\renewcommand{\rhcaMzlRsqLassoFullClosePre}{0.80}
\providecommand{\rhcaShBiasLassoFullClosePre}{}\renewcommand{\rhcaShBiasLassoFullClosePre}{0.2}
\providecommand{\rhcaShSlopeLassoFullClosePre}{}\renewcommand{\rhcaShSlopeLassoFullClosePre}{0.7}
\providecommand{\rhcaShDistLassoFullClosePre}{}\renewcommand{\rhcaShDistLassoFullClosePre}{99.1}
\providecommand{\rhcaQlLassoFullClosePre}{}\renewcommand{\rhcaQlLassoFullClosePre}{0.0100}
\providecommand{\rhcaQdLassoFullClosePre}{}\renewcommand{\rhcaQdLassoFullClosePre}{0.1233}
\providecommand{\rhcaOverSeLassoFullClosePre}{}\renewcommand{\rhcaOverSeLassoFullClosePre}{18.4}
\providecommand{\rhcaMrEnetHarClosePre}{}\renewcommand{\rhcaMrEnetHarClosePre}{1.121}
\providecommand{\rhcaRmEnetHarClosePre}{}\renewcommand{\rhcaRmEnetHarClosePre}{1.097}
\providecommand{\rhcaUnderEnetHarClosePre}{}\renewcommand{\rhcaUnderEnetHarClosePre}{46.7}
\providecommand{\rhcaMzAEnetHarClosePre}{}\renewcommand{\rhcaMzAEnetHarClosePre}{\ensuremath{-}0.16}
\providecommand{\rhcaMzBEnetHarClosePre}{}\renewcommand{\rhcaMzBEnetHarClosePre}{1.14}
\providecommand{\rhcaMzBseEnetHarClosePre}{}\renewcommand{\rhcaMzBseEnetHarClosePre}{0.10}
\providecommand{\rhcaMzRsqEnetHarClosePre}{}\renewcommand{\rhcaMzRsqEnetHarClosePre}{0.64}
\providecommand{\rhcaMzlAEnetHarClosePre}{}\renewcommand{\rhcaMzlAEnetHarClosePre}{\ensuremath{-}0.68}
\providecommand{\rhcaMzlBEnetHarClosePre}{}\renewcommand{\rhcaMzlBEnetHarClosePre}{0.95}
\providecommand{\rhcaMzlBseEnetHarClosePre}{}\renewcommand{\rhcaMzlBseEnetHarClosePre}{0.02}
\providecommand{\rhcaMzlRsqEnetHarClosePre}{}\renewcommand{\rhcaMzlRsqEnetHarClosePre}{0.78}
\providecommand{\rhcaShBiasEnetHarClosePre}{}\renewcommand{\rhcaShBiasEnetHarClosePre}{0.5}
\providecommand{\rhcaShSlopeEnetHarClosePre}{}\renewcommand{\rhcaShSlopeEnetHarClosePre}{2.5}
\providecommand{\rhcaShDistEnetHarClosePre}{}\renewcommand{\rhcaShDistEnetHarClosePre}{97.0}
\providecommand{\rhcaQlEnetHarClosePre}{}\renewcommand{\rhcaQlEnetHarClosePre}{0.0068}
\providecommand{\rhcaQdEnetHarClosePre}{}\renewcommand{\rhcaQdEnetHarClosePre}{0.1297}
\providecommand{\rhcaOverSeEnetHarClosePre}{}\renewcommand{\rhcaOverSeEnetHarClosePre}{13.1}
\providecommand{\rhcaMrEnetFullClosePre}{}\renewcommand{\rhcaMrEnetFullClosePre}{1.138}
\providecommand{\rhcaRmEnetFullClosePre}{}\renewcommand{\rhcaRmEnetFullClosePre}{1.067}
\providecommand{\rhcaUnderEnetFullClosePre}{}\renewcommand{\rhcaUnderEnetFullClosePre}{47.3}
\providecommand{\rhcaMzAEnetFullClosePre}{}\renewcommand{\rhcaMzAEnetFullClosePre}{\ensuremath{-}0.05}
\providecommand{\rhcaMzBEnetFullClosePre}{}\renewcommand{\rhcaMzBEnetFullClosePre}{1.08}
\providecommand{\rhcaMzBseEnetFullClosePre}{}\renewcommand{\rhcaMzBseEnetFullClosePre}{0.09}
\providecommand{\rhcaMzRsqEnetFullClosePre}{}\renewcommand{\rhcaMzRsqEnetFullClosePre}{0.69}
\providecommand{\rhcaMzlAEnetFullClosePre}{}\renewcommand{\rhcaMzlAEnetFullClosePre}{\ensuremath{-}0.90}
\providecommand{\rhcaMzlBEnetFullClosePre}{}\renewcommand{\rhcaMzlBEnetFullClosePre}{0.93}
\providecommand{\rhcaMzlBseEnetFullClosePre}{}\renewcommand{\rhcaMzlBseEnetFullClosePre}{0.02}
\providecommand{\rhcaMzlRsqEnetFullClosePre}{}\renewcommand{\rhcaMzlRsqEnetFullClosePre}{0.80}
\providecommand{\rhcaShBiasEnetFullClosePre}{}\renewcommand{\rhcaShBiasEnetFullClosePre}{0.3}
\providecommand{\rhcaShSlopeEnetFullClosePre}{}\renewcommand{\rhcaShSlopeEnetFullClosePre}{1.2}
\providecommand{\rhcaShDistEnetFullClosePre}{}\renewcommand{\rhcaShDistEnetFullClosePre}{98.5}
\providecommand{\rhcaQlEnetFullClosePre}{}\renewcommand{\rhcaQlEnetFullClosePre}{0.0087}
\providecommand{\rhcaQdEnetFullClosePre}{}\renewcommand{\rhcaQdEnetFullClosePre}{0.1213}
\providecommand{\rhcaOverSeEnetFullClosePre}{}\renewcommand{\rhcaOverSeEnetFullClosePre}{17.3}
\providecommand{\rhcaMrLgbmHarClosePre}{}\renewcommand{\rhcaMrLgbmHarClosePre}{1.068}
\providecommand{\rhcaRmLgbmHarClosePre}{}\renewcommand{\rhcaRmLgbmHarClosePre}{1.062}
\providecommand{\rhcaUnderLgbmHarClosePre}{}\renewcommand{\rhcaUnderLgbmHarClosePre}{42.9}
\providecommand{\rhcaMzALgbmHarClosePre}{}\renewcommand{\rhcaMzALgbmHarClosePre}{0.00}
\providecommand{\rhcaMzBLgbmHarClosePre}{}\renewcommand{\rhcaMzBLgbmHarClosePre}{1.06}
\providecommand{\rhcaMzBseLgbmHarClosePre}{}\renewcommand{\rhcaMzBseLgbmHarClosePre}{0.10}
\providecommand{\rhcaMzRsqLgbmHarClosePre}{}\renewcommand{\rhcaMzRsqLgbmHarClosePre}{0.55}
\providecommand{\rhcaMzlALgbmHarClosePre}{}\renewcommand{\rhcaMzlALgbmHarClosePre}{\ensuremath{-}0.59}
\providecommand{\rhcaMzlBLgbmHarClosePre}{}\renewcommand{\rhcaMzlBLgbmHarClosePre}{0.96}
\providecommand{\rhcaMzlBseLgbmHarClosePre}{}\renewcommand{\rhcaMzlBseLgbmHarClosePre}{0.02}
\providecommand{\rhcaMzlRsqLgbmHarClosePre}{}\renewcommand{\rhcaMzlRsqLgbmHarClosePre}{0.78}
\providecommand{\rhcaShBiasLgbmHarClosePre}{}\renewcommand{\rhcaShBiasLgbmHarClosePre}{0.2}
\providecommand{\rhcaShSlopeLgbmHarClosePre}{}\renewcommand{\rhcaShSlopeLgbmHarClosePre}{0.4}
\providecommand{\rhcaShDistLgbmHarClosePre}{}\renewcommand{\rhcaShDistLgbmHarClosePre}{99.4}
\providecommand{\rhcaQlLgbmHarClosePre}{}\renewcommand{\rhcaQlLgbmHarClosePre}{0.0022}
\providecommand{\rhcaQdLgbmHarClosePre}{}\renewcommand{\rhcaQdLgbmHarClosePre}{0.1310}
\providecommand{\rhcaOverSeLgbmHarClosePre}{}\renewcommand{\rhcaOverSeLgbmHarClosePre}{14.4}
\providecommand{\rhcaMrLgbmFullClosePre}{}\renewcommand{\rhcaMrLgbmFullClosePre}{1.101}
\providecommand{\rhcaRmLgbmFullClosePre}{}\renewcommand{\rhcaRmLgbmFullClosePre}{1.042}
\providecommand{\rhcaUnderLgbmFullClosePre}{}\renewcommand{\rhcaUnderLgbmFullClosePre}{45.9}
\providecommand{\rhcaMzALgbmFullClosePre}{}\renewcommand{\rhcaMzALgbmFullClosePre}{0.26}
\providecommand{\rhcaMzBLgbmFullClosePre}{}\renewcommand{\rhcaMzBLgbmFullClosePre}{0.98}
\providecommand{\rhcaMzBseLgbmFullClosePre}{}\renewcommand{\rhcaMzBseLgbmFullClosePre}{0.07}
\providecommand{\rhcaMzRsqLgbmFullClosePre}{}\renewcommand{\rhcaMzRsqLgbmFullClosePre}{0.58}
\providecommand{\rhcaMzlALgbmFullClosePre}{}\renewcommand{\rhcaMzlALgbmFullClosePre}{\ensuremath{-}0.89}
\providecommand{\rhcaMzlBLgbmFullClosePre}{}\renewcommand{\rhcaMzlBLgbmFullClosePre}{0.93}
\providecommand{\rhcaMzlBseLgbmFullClosePre}{}\renewcommand{\rhcaMzlBseLgbmFullClosePre}{0.02}
\providecommand{\rhcaMzlRsqLgbmFullClosePre}{}\renewcommand{\rhcaMzlRsqLgbmFullClosePre}{0.80}
\providecommand{\rhcaShBiasLgbmFullClosePre}{}\renewcommand{\rhcaShBiasLgbmFullClosePre}{0.1}
\providecommand{\rhcaShSlopeLgbmFullClosePre}{}\renewcommand{\rhcaShSlopeLgbmFullClosePre}{0.1}
\providecommand{\rhcaShDistLgbmFullClosePre}{}\renewcommand{\rhcaShDistLgbmFullClosePre}{99.9}
\providecommand{\rhcaQlLgbmFullClosePre}{}\renewcommand{\rhcaQlLgbmFullClosePre}{0.0048}
\providecommand{\rhcaQdLgbmFullClosePre}{}\renewcommand{\rhcaQdLgbmFullClosePre}{0.1239}
\providecommand{\rhcaOverSeLgbmFullClosePre}{}\renewcommand{\rhcaOverSeLgbmFullClosePre}{17.3}
\providecommand{\rhcaMrXgbHarClosePre}{}\renewcommand{\rhcaMrXgbHarClosePre}{1.064}
\providecommand{\rhcaRmXgbHarClosePre}{}\renewcommand{\rhcaRmXgbHarClosePre}{1.058}
\providecommand{\rhcaUnderXgbHarClosePre}{}\renewcommand{\rhcaUnderXgbHarClosePre}{42.6}
\providecommand{\rhcaMzAXgbHarClosePre}{}\renewcommand{\rhcaMzAXgbHarClosePre}{\ensuremath{-}0.12}
\providecommand{\rhcaMzBXgbHarClosePre}{}\renewcommand{\rhcaMzBXgbHarClosePre}{1.09}
\providecommand{\rhcaMzBseXgbHarClosePre}{}\renewcommand{\rhcaMzBseXgbHarClosePre}{0.09}
\providecommand{\rhcaMzRsqXgbHarClosePre}{}\renewcommand{\rhcaMzRsqXgbHarClosePre}{0.58}
\providecommand{\rhcaMzlAXgbHarClosePre}{}\renewcommand{\rhcaMzlAXgbHarClosePre}{\ensuremath{-}0.58}
\providecommand{\rhcaMzlBXgbHarClosePre}{}\renewcommand{\rhcaMzlBXgbHarClosePre}{0.96}
\providecommand{\rhcaMzlBseXgbHarClosePre}{}\renewcommand{\rhcaMzlBseXgbHarClosePre}{0.02}
\providecommand{\rhcaMzlRsqXgbHarClosePre}{}\renewcommand{\rhcaMzlRsqXgbHarClosePre}{0.78}
\providecommand{\rhcaShBiasXgbHarClosePre}{}\renewcommand{\rhcaShBiasXgbHarClosePre}{0.2}
\providecommand{\rhcaShSlopeXgbHarClosePre}{}\renewcommand{\rhcaShSlopeXgbHarClosePre}{0.9}
\providecommand{\rhcaShDistXgbHarClosePre}{}\renewcommand{\rhcaShDistXgbHarClosePre}{99.0}
\providecommand{\rhcaQlXgbHarClosePre}{}\renewcommand{\rhcaQlXgbHarClosePre}{0.0020}
\providecommand{\rhcaQdXgbHarClosePre}{}\renewcommand{\rhcaQdXgbHarClosePre}{0.1311}
\providecommand{\rhcaOverSeXgbHarClosePre}{}\renewcommand{\rhcaOverSeXgbHarClosePre}{14.2}
\providecommand{\rhcaMrXgbFullClosePre}{}\renewcommand{\rhcaMrXgbFullClosePre}{1.098}
\providecommand{\rhcaRmXgbFullClosePre}{}\renewcommand{\rhcaRmXgbFullClosePre}{1.037}
\providecommand{\rhcaUnderXgbFullClosePre}{}\renewcommand{\rhcaUnderXgbFullClosePre}{45.6}
\providecommand{\rhcaMzAXgbFullClosePre}{}\renewcommand{\rhcaMzAXgbFullClosePre}{0.19}
\providecommand{\rhcaMzBXgbFullClosePre}{}\renewcommand{\rhcaMzBXgbFullClosePre}{0.99}
\providecommand{\rhcaMzBseXgbFullClosePre}{}\renewcommand{\rhcaMzBseXgbFullClosePre}{0.08}
\providecommand{\rhcaMzRsqXgbFullClosePre}{}\renewcommand{\rhcaMzRsqXgbFullClosePre}{0.62}
\providecommand{\rhcaMzlAXgbFullClosePre}{}\renewcommand{\rhcaMzlAXgbFullClosePre}{\ensuremath{-}0.87}
\providecommand{\rhcaMzlBXgbFullClosePre}{}\renewcommand{\rhcaMzlBXgbFullClosePre}{0.94}
\providecommand{\rhcaMzlBseXgbFullClosePre}{}\renewcommand{\rhcaMzlBseXgbFullClosePre}{0.02}
\providecommand{\rhcaMzlRsqXgbFullClosePre}{}\renewcommand{\rhcaMzlRsqXgbFullClosePre}{0.80}
\providecommand{\rhcaShBiasXgbFullClosePre}{}\renewcommand{\rhcaShBiasXgbFullClosePre}{0.1}
\providecommand{\rhcaShSlopeXgbFullClosePre}{}\renewcommand{\rhcaShSlopeXgbFullClosePre}{0.0}
\providecommand{\rhcaShDistXgbFullClosePre}{}\renewcommand{\rhcaShDistXgbFullClosePre}{99.9}
\providecommand{\rhcaQlXgbFullClosePre}{}\renewcommand{\rhcaQlXgbFullClosePre}{0.0045}
\providecommand{\rhcaQdXgbFullClosePre}{}\renewcommand{\rhcaQdXgbFullClosePre}{0.1232}
\providecommand{\rhcaOverSeXgbFullClosePre}{}\renewcommand{\rhcaOverSeXgbFullClosePre}{18.8}
\providecommand{\rhcaMrRfHarClosePre}{}\renewcommand{\rhcaMrRfHarClosePre}{1.091}
\providecommand{\rhcaRmRfHarClosePre}{}\renewcommand{\rhcaRmRfHarClosePre}{1.073}
\providecommand{\rhcaUnderRfHarClosePre}{}\renewcommand{\rhcaUnderRfHarClosePre}{43.6}
\providecommand{\rhcaMzARfHarClosePre}{}\renewcommand{\rhcaMzARfHarClosePre}{\ensuremath{-}0.13}
\providecommand{\rhcaMzBRfHarClosePre}{}\renewcommand{\rhcaMzBRfHarClosePre}{1.11}
\providecommand{\rhcaMzBseRfHarClosePre}{}\renewcommand{\rhcaMzBseRfHarClosePre}{0.10}
\providecommand{\rhcaMzRsqRfHarClosePre}{}\renewcommand{\rhcaMzRsqRfHarClosePre}{0.59}
\providecommand{\rhcaMzlARfHarClosePre}{}\renewcommand{\rhcaMzlARfHarClosePre}{\ensuremath{-}0.68}
\providecommand{\rhcaMzlBRfHarClosePre}{}\renewcommand{\rhcaMzlBRfHarClosePre}{0.95}
\providecommand{\rhcaMzlBseRfHarClosePre}{}\renewcommand{\rhcaMzlBseRfHarClosePre}{0.02}
\providecommand{\rhcaMzlRsqRfHarClosePre}{}\renewcommand{\rhcaMzlRsqRfHarClosePre}{0.78}
\providecommand{\rhcaShBiasRfHarClosePre}{}\renewcommand{\rhcaShBiasRfHarClosePre}{0.2}
\providecommand{\rhcaShSlopeRfHarClosePre}{}\renewcommand{\rhcaShSlopeRfHarClosePre}{1.3}
\providecommand{\rhcaShDistRfHarClosePre}{}\renewcommand{\rhcaShDistRfHarClosePre}{98.5}
\providecommand{\rhcaQlRfHarClosePre}{}\renewcommand{\rhcaQlRfHarClosePre}{0.0039}
\providecommand{\rhcaQdRfHarClosePre}{}\renewcommand{\rhcaQdRfHarClosePre}{0.1325}
\providecommand{\rhcaOverSeRfHarClosePre}{}\renewcommand{\rhcaOverSeRfHarClosePre}{14.1}
\providecommand{\rhcaMrRfFullClosePre}{}\renewcommand{\rhcaMrRfFullClosePre}{1.119}
\providecommand{\rhcaRmRfFullClosePre}{}\renewcommand{\rhcaRmRfFullClosePre}{1.037}
\providecommand{\rhcaUnderRfFullClosePre}{}\renewcommand{\rhcaUnderRfFullClosePre}{46.1}
\providecommand{\rhcaMzARfFullClosePre}{}\renewcommand{\rhcaMzARfFullClosePre}{0.25}
\providecommand{\rhcaMzBRfFullClosePre}{}\renewcommand{\rhcaMzBRfFullClosePre}{0.98}
\providecommand{\rhcaMzBseRfFullClosePre}{}\renewcommand{\rhcaMzBseRfFullClosePre}{0.08}
\providecommand{\rhcaMzRsqRfFullClosePre}{}\renewcommand{\rhcaMzRsqRfFullClosePre}{0.61}
\providecommand{\rhcaMzlARfFullClosePre}{}\renewcommand{\rhcaMzlARfFullClosePre}{\ensuremath{-}1.04}
\providecommand{\rhcaMzlBRfFullClosePre}{}\renewcommand{\rhcaMzlBRfFullClosePre}{0.92}
\providecommand{\rhcaMzlBseRfFullClosePre}{}\renewcommand{\rhcaMzlBseRfFullClosePre}{0.02}
\providecommand{\rhcaMzlRsqRfFullClosePre}{}\renewcommand{\rhcaMzlRsqRfFullClosePre}{0.79}
\providecommand{\rhcaShBiasRfFullClosePre}{}\renewcommand{\rhcaShBiasRfFullClosePre}{0.1}
\providecommand{\rhcaShSlopeRfFullClosePre}{}\renewcommand{\rhcaShSlopeRfFullClosePre}{0.1}
\providecommand{\rhcaShDistRfFullClosePre}{}\renewcommand{\rhcaShDistRfFullClosePre}{99.8}
\providecommand{\rhcaQlRfFullClosePre}{}\renewcommand{\rhcaQlRfFullClosePre}{0.0065}
\providecommand{\rhcaQdRfFullClosePre}{}\renewcommand{\rhcaQdRfFullClosePre}{0.1270}
\providecommand{\rhcaOverSeRfFullClosePre}{}\renewcommand{\rhcaOverSeRfFullClosePre}{18.9}
\providecommand{\rhcaMrRidgeHarBarOos}{}\renewcommand{\rhcaMrRidgeHarBarOos}{1.061}
\providecommand{\rhcaRmRidgeHarBarOos}{}\renewcommand{\rhcaRmRidgeHarBarOos}{1.045}
\providecommand{\rhcaUnderRidgeHarBarOos}{}\renewcommand{\rhcaUnderRidgeHarBarOos}{38.0}
\providecommand{\rhcaMzARidgeHarBarOos}{}\renewcommand{\rhcaMzARidgeHarBarOos}{0.41}
\providecommand{\rhcaMzBRidgeHarBarOos}{}\renewcommand{\rhcaMzBRidgeHarBarOos}{0.99}
\providecommand{\rhcaMzBseRidgeHarBarOos}{}\renewcommand{\rhcaMzBseRidgeHarBarOos}{0.08}
\providecommand{\rhcaMzRsqRidgeHarBarOos}{}\renewcommand{\rhcaMzRsqRidgeHarBarOos}{0.67}
\providecommand{\rhcaMzlARidgeHarBarOos}{}\renewcommand{\rhcaMzlARidgeHarBarOos}{0.15}
\providecommand{\rhcaMzlBRidgeHarBarOos}{}\renewcommand{\rhcaMzlBRidgeHarBarOos}{1.02}
\providecommand{\rhcaMzlBseRidgeHarBarOos}{}\renewcommand{\rhcaMzlBseRidgeHarBarOos}{0.01}
\providecommand{\rhcaMzlRsqRidgeHarBarOos}{}\renewcommand{\rhcaMzlRsqRidgeHarBarOos}{0.79}
\providecommand{\rhcaShBiasRidgeHarBarOos}{}\renewcommand{\rhcaShBiasRidgeHarBarOos}{0.1}
\providecommand{\rhcaShSlopeRidgeHarBarOos}{}\renewcommand{\rhcaShSlopeRidgeHarBarOos}{0.0}
\providecommand{\rhcaShDistRidgeHarBarOos}{}\renewcommand{\rhcaShDistRidgeHarBarOos}{99.9}
\providecommand{\rhcaQlRidgeHarBarOos}{}\renewcommand{\rhcaQlRidgeHarBarOos}{0.0018}
\providecommand{\rhcaQdRidgeHarBarOos}{}\renewcommand{\rhcaQdRidgeHarBarOos}{0.1915}
\providecommand{\rhcaOverSeRidgeHarBarOos}{}\renewcommand{\rhcaOverSeRidgeHarBarOos}{21.5}
\providecommand{\rhcaMrRidgeFullBarOos}{}\renewcommand{\rhcaMrRidgeFullBarOos}{1.117}
\providecommand{\rhcaRmRidgeFullBarOos}{}\renewcommand{\rhcaRmRidgeFullBarOos}{0.977}
\providecommand{\rhcaUnderRidgeFullBarOos}{}\renewcommand{\rhcaUnderRidgeFullBarOos}{42.5}
\providecommand{\rhcaMzARidgeFullBarOos}{}\renewcommand{\rhcaMzARidgeFullBarOos}{3.27}
\providecommand{\rhcaMzBRidgeFullBarOos}{}\renewcommand{\rhcaMzBRidgeFullBarOos}{0.58}
\providecommand{\rhcaMzBseRidgeFullBarOos}{}\renewcommand{\rhcaMzBseRidgeFullBarOos}{0.10}
\providecommand{\rhcaMzRsqRidgeFullBarOos}{}\renewcommand{\rhcaMzRsqRidgeFullBarOos}{0.52}
\providecommand{\rhcaMzlARidgeFullBarOos}{}\renewcommand{\rhcaMzlARidgeFullBarOos}{\ensuremath{-}0.25}
\providecommand{\rhcaMzlBRidgeFullBarOos}{}\renewcommand{\rhcaMzlBRidgeFullBarOos}{0.99}
\providecommand{\rhcaMzlBseRidgeFullBarOos}{}\renewcommand{\rhcaMzlBseRidgeFullBarOos}{0.01}
\providecommand{\rhcaMzlRsqRidgeFullBarOos}{}\renewcommand{\rhcaMzlRsqRidgeFullBarOos}{0.81}
\providecommand{\rhcaShBiasRidgeFullBarOos}{}\renewcommand{\rhcaShBiasRidgeFullBarOos}{0.0}
\providecommand{\rhcaShSlopeRidgeFullBarOos}{}\renewcommand{\rhcaShSlopeRidgeFullBarOos}{35.9}
\providecommand{\rhcaShDistRidgeFullBarOos}{}\renewcommand{\rhcaShDistRidgeFullBarOos}{64.1}
\providecommand{\rhcaQlRidgeFullBarOos}{}\renewcommand{\rhcaQlRidgeFullBarOos}{0.0064}
\providecommand{\rhcaQdRidgeFullBarOos}{}\renewcommand{\rhcaQdRidgeFullBarOos}{0.1788}
\providecommand{\rhcaOverSeRidgeFullBarOos}{}\renewcommand{\rhcaOverSeRidgeFullBarOos}{63.1}
\providecommand{\rhcaMrLassoHarBarOos}{}\renewcommand{\rhcaMrLassoHarBarOos}{1.061}
\providecommand{\rhcaRmLassoHarBarOos}{}\renewcommand{\rhcaRmLassoHarBarOos}{1.045}
\providecommand{\rhcaUnderLassoHarBarOos}{}\renewcommand{\rhcaUnderLassoHarBarOos}{38.0}
\providecommand{\rhcaMzALassoHarBarOos}{}\renewcommand{\rhcaMzALassoHarBarOos}{0.41}
\providecommand{\rhcaMzBLassoHarBarOos}{}\renewcommand{\rhcaMzBLassoHarBarOos}{0.99}
\providecommand{\rhcaMzBseLassoHarBarOos}{}\renewcommand{\rhcaMzBseLassoHarBarOos}{0.08}
\providecommand{\rhcaMzRsqLassoHarBarOos}{}\renewcommand{\rhcaMzRsqLassoHarBarOos}{0.67}
\providecommand{\rhcaMzlALassoHarBarOos}{}\renewcommand{\rhcaMzlALassoHarBarOos}{0.15}
\providecommand{\rhcaMzlBLassoHarBarOos}{}\renewcommand{\rhcaMzlBLassoHarBarOos}{1.02}
\providecommand{\rhcaMzlBseLassoHarBarOos}{}\renewcommand{\rhcaMzlBseLassoHarBarOos}{0.01}
\providecommand{\rhcaMzlRsqLassoHarBarOos}{}\renewcommand{\rhcaMzlRsqLassoHarBarOos}{0.79}
\providecommand{\rhcaShBiasLassoHarBarOos}{}\renewcommand{\rhcaShBiasLassoHarBarOos}{0.1}
\providecommand{\rhcaShSlopeLassoHarBarOos}{}\renewcommand{\rhcaShSlopeLassoHarBarOos}{0.0}
\providecommand{\rhcaShDistLassoHarBarOos}{}\renewcommand{\rhcaShDistLassoHarBarOos}{99.9}
\providecommand{\rhcaQlLassoHarBarOos}{}\renewcommand{\rhcaQlLassoHarBarOos}{0.0018}
\providecommand{\rhcaQdLassoHarBarOos}{}\renewcommand{\rhcaQdLassoHarBarOos}{0.1915}
\providecommand{\rhcaOverSeLassoHarBarOos}{}\renewcommand{\rhcaOverSeLassoHarBarOos}{21.5}
\providecommand{\rhcaMrLassoFullBarOos}{}\renewcommand{\rhcaMrLassoFullBarOos}{1.116}
\providecommand{\rhcaRmLassoFullBarOos}{}\renewcommand{\rhcaRmLassoFullBarOos}{0.978}
\providecommand{\rhcaUnderLassoFullBarOos}{}\renewcommand{\rhcaUnderLassoFullBarOos}{43.2}
\providecommand{\rhcaMzALassoFullBarOos}{}\renewcommand{\rhcaMzALassoFullBarOos}{3.46}
\providecommand{\rhcaMzBLassoFullBarOos}{}\renewcommand{\rhcaMzBLassoFullBarOos}{0.56}
\providecommand{\rhcaMzBseLassoFullBarOos}{}\renewcommand{\rhcaMzBseLassoFullBarOos}{0.10}
\providecommand{\rhcaMzRsqLassoFullBarOos}{}\renewcommand{\rhcaMzRsqLassoFullBarOos}{0.51}
\providecommand{\rhcaMzlALassoFullBarOos}{}\renewcommand{\rhcaMzlALassoFullBarOos}{\ensuremath{-}0.26}
\providecommand{\rhcaMzlBLassoFullBarOos}{}\renewcommand{\rhcaMzlBLassoFullBarOos}{0.98}
\providecommand{\rhcaMzlBseLassoFullBarOos}{}\renewcommand{\rhcaMzlBseLassoFullBarOos}{0.01}
\providecommand{\rhcaMzlRsqLassoFullBarOos}{}\renewcommand{\rhcaMzlRsqLassoFullBarOos}{0.81}
\providecommand{\rhcaShBiasLassoFullBarOos}{}\renewcommand{\rhcaShBiasLassoFullBarOos}{0.0}
\providecommand{\rhcaShSlopeLassoFullBarOos}{}\renewcommand{\rhcaShSlopeLassoFullBarOos}{38.2}
\providecommand{\rhcaShDistLassoFullBarOos}{}\renewcommand{\rhcaShDistLassoFullBarOos}{61.8}
\providecommand{\rhcaQlLassoFullBarOos}{}\renewcommand{\rhcaQlLassoFullBarOos}{0.0062}
\providecommand{\rhcaQdLassoFullBarOos}{}\renewcommand{\rhcaQdLassoFullBarOos}{0.1739}
\providecommand{\rhcaOverSeLassoFullBarOos}{}\renewcommand{\rhcaOverSeLassoFullBarOos}{64.5}
\providecommand{\rhcaMrEnetHarBarOos}{}\renewcommand{\rhcaMrEnetHarBarOos}{1.061}
\providecommand{\rhcaRmEnetHarBarOos}{}\renewcommand{\rhcaRmEnetHarBarOos}{1.045}
\providecommand{\rhcaUnderEnetHarBarOos}{}\renewcommand{\rhcaUnderEnetHarBarOos}{38.0}
\providecommand{\rhcaMzAEnetHarBarOos}{}\renewcommand{\rhcaMzAEnetHarBarOos}{0.41}
\providecommand{\rhcaMzBEnetHarBarOos}{}\renewcommand{\rhcaMzBEnetHarBarOos}{0.99}
\providecommand{\rhcaMzBseEnetHarBarOos}{}\renewcommand{\rhcaMzBseEnetHarBarOos}{0.08}
\providecommand{\rhcaMzRsqEnetHarBarOos}{}\renewcommand{\rhcaMzRsqEnetHarBarOos}{0.67}
\providecommand{\rhcaMzlAEnetHarBarOos}{}\renewcommand{\rhcaMzlAEnetHarBarOos}{0.15}
\providecommand{\rhcaMzlBEnetHarBarOos}{}\renewcommand{\rhcaMzlBEnetHarBarOos}{1.02}
\providecommand{\rhcaMzlBseEnetHarBarOos}{}\renewcommand{\rhcaMzlBseEnetHarBarOos}{0.01}
\providecommand{\rhcaMzlRsqEnetHarBarOos}{}\renewcommand{\rhcaMzlRsqEnetHarBarOos}{0.79}
\providecommand{\rhcaShBiasEnetHarBarOos}{}\renewcommand{\rhcaShBiasEnetHarBarOos}{0.1}
\providecommand{\rhcaShSlopeEnetHarBarOos}{}\renewcommand{\rhcaShSlopeEnetHarBarOos}{0.0}
\providecommand{\rhcaShDistEnetHarBarOos}{}\renewcommand{\rhcaShDistEnetHarBarOos}{99.9}
\providecommand{\rhcaQlEnetHarBarOos}{}\renewcommand{\rhcaQlEnetHarBarOos}{0.0018}
\providecommand{\rhcaQdEnetHarBarOos}{}\renewcommand{\rhcaQdEnetHarBarOos}{0.1915}
\providecommand{\rhcaOverSeEnetHarBarOos}{}\renewcommand{\rhcaOverSeEnetHarBarOos}{21.5}
\providecommand{\rhcaMrEnetFullBarOos}{}\renewcommand{\rhcaMrEnetFullBarOos}{1.109}
\providecommand{\rhcaRmEnetFullBarOos}{}\renewcommand{\rhcaRmEnetFullBarOos}{0.966}
\providecommand{\rhcaUnderEnetFullBarOos}{}\renewcommand{\rhcaUnderEnetFullBarOos}{42.2}
\providecommand{\rhcaMzAEnetFullBarOos}{}\renewcommand{\rhcaMzAEnetFullBarOos}{3.45}
\providecommand{\rhcaMzBEnetFullBarOos}{}\renewcommand{\rhcaMzBEnetFullBarOos}{0.56}
\providecommand{\rhcaMzBseEnetFullBarOos}{}\renewcommand{\rhcaMzBseEnetFullBarOos}{0.10}
\providecommand{\rhcaMzRsqEnetFullBarOos}{}\renewcommand{\rhcaMzRsqEnetFullBarOos}{0.52}
\providecommand{\rhcaMzlAEnetFullBarOos}{}\renewcommand{\rhcaMzlAEnetFullBarOos}{\ensuremath{-}0.22}
\providecommand{\rhcaMzlBEnetFullBarOos}{}\renewcommand{\rhcaMzlBEnetFullBarOos}{0.99}
\providecommand{\rhcaMzlBseEnetFullBarOos}{}\renewcommand{\rhcaMzlBseEnetFullBarOos}{0.01}
\providecommand{\rhcaMzlRsqEnetFullBarOos}{}\renewcommand{\rhcaMzlRsqEnetFullBarOos}{0.81}
\providecommand{\rhcaShBiasEnetFullBarOos}{}\renewcommand{\rhcaShBiasEnetFullBarOos}{0.0}
\providecommand{\rhcaShSlopeEnetFullBarOos}{}\renewcommand{\rhcaShSlopeEnetFullBarOos}{40.8}
\providecommand{\rhcaShDistEnetFullBarOos}{}\renewcommand{\rhcaShDistEnetFullBarOos}{59.1}
\providecommand{\rhcaQlEnetFullBarOos}{}\renewcommand{\rhcaQlEnetFullBarOos}{0.0056}
\providecommand{\rhcaQdEnetFullBarOos}{}\renewcommand{\rhcaQdEnetFullBarOos}{0.1775}
\providecommand{\rhcaOverSeEnetFullBarOos}{}\renewcommand{\rhcaOverSeEnetFullBarOos}{67.0}
\providecommand{\rhcaMrLgbmHarBarOos}{}\renewcommand{\rhcaMrLgbmHarBarOos}{1.067}
\providecommand{\rhcaRmLgbmHarBarOos}{}\renewcommand{\rhcaRmLgbmHarBarOos}{1.028}
\providecommand{\rhcaUnderLgbmHarBarOos}{}\renewcommand{\rhcaUnderLgbmHarBarOos}{38.5}
\providecommand{\rhcaMzALgbmHarBarOos}{}\renewcommand{\rhcaMzALgbmHarBarOos}{0.72}
\providecommand{\rhcaMzBLgbmHarBarOos}{}\renewcommand{\rhcaMzBLgbmHarBarOos}{0.94}
\providecommand{\rhcaMzBseLgbmHarBarOos}{}\renewcommand{\rhcaMzBseLgbmHarBarOos}{0.06}
\providecommand{\rhcaMzRsqLgbmHarBarOos}{}\renewcommand{\rhcaMzRsqLgbmHarBarOos}{0.66}
\providecommand{\rhcaMzlALgbmHarBarOos}{}\renewcommand{\rhcaMzlALgbmHarBarOos}{0.03}
\providecommand{\rhcaMzlBLgbmHarBarOos}{}\renewcommand{\rhcaMzlBLgbmHarBarOos}{1.01}
\providecommand{\rhcaMzlBseLgbmHarBarOos}{}\renewcommand{\rhcaMzlBseLgbmHarBarOos}{0.01}
\providecommand{\rhcaMzlRsqLgbmHarBarOos}{}\renewcommand{\rhcaMzlRsqLgbmHarBarOos}{0.80}
\providecommand{\rhcaShBiasLgbmHarBarOos}{}\renewcommand{\rhcaShBiasLgbmHarBarOos}{0.0}
\providecommand{\rhcaShSlopeLgbmHarBarOos}{}\renewcommand{\rhcaShSlopeLgbmHarBarOos}{0.8}
\providecommand{\rhcaShDistLgbmHarBarOos}{}\renewcommand{\rhcaShDistLgbmHarBarOos}{99.1}
\providecommand{\rhcaQlLgbmHarBarOos}{}\renewcommand{\rhcaQlLgbmHarBarOos}{0.0022}
\providecommand{\rhcaQdLgbmHarBarOos}{}\renewcommand{\rhcaQdLgbmHarBarOos}{0.1878}
\providecommand{\rhcaOverSeLgbmHarBarOos}{}\renewcommand{\rhcaOverSeLgbmHarBarOos}{25.9}
\providecommand{\rhcaMrLgbmFullBarOos}{}\renewcommand{\rhcaMrLgbmFullBarOos}{1.072}
\providecommand{\rhcaRmLgbmFullBarOos}{}\renewcommand{\rhcaRmLgbmFullBarOos}{0.963}
\providecommand{\rhcaUnderLgbmFullBarOos}{}\renewcommand{\rhcaUnderLgbmFullBarOos}{40.6}
\providecommand{\rhcaMzALgbmFullBarOos}{}\renewcommand{\rhcaMzALgbmFullBarOos}{1.86}
\providecommand{\rhcaMzBLgbmFullBarOos}{}\renewcommand{\rhcaMzBLgbmFullBarOos}{0.74}
\providecommand{\rhcaMzBseLgbmFullBarOos}{}\renewcommand{\rhcaMzBseLgbmFullBarOos}{0.07}
\providecommand{\rhcaMzRsqLgbmFullBarOos}{}\renewcommand{\rhcaMzRsqLgbmFullBarOos}{0.64}
\providecommand{\rhcaMzlALgbmFullBarOos}{}\renewcommand{\rhcaMzlALgbmFullBarOos}{\ensuremath{-}0.20}
\providecommand{\rhcaMzlBLgbmFullBarOos}{}\renewcommand{\rhcaMzlBLgbmFullBarOos}{0.99}
\providecommand{\rhcaMzlBseLgbmFullBarOos}{}\renewcommand{\rhcaMzlBseLgbmFullBarOos}{0.01}
\providecommand{\rhcaMzlRsqLgbmFullBarOos}{}\renewcommand{\rhcaMzlRsqLgbmFullBarOos}{0.82}
\providecommand{\rhcaShBiasLgbmFullBarOos}{}\renewcommand{\rhcaShBiasLgbmFullBarOos}{0.0}
\providecommand{\rhcaShSlopeLgbmFullBarOos}{}\renewcommand{\rhcaShSlopeLgbmFullBarOos}{17.7}
\providecommand{\rhcaShDistLgbmFullBarOos}{}\renewcommand{\rhcaShDistLgbmFullBarOos}{82.2}
\providecommand{\rhcaQlLgbmFullBarOos}{}\renewcommand{\rhcaQlLgbmFullBarOos}{0.0025}
\providecommand{\rhcaQdLgbmFullBarOos}{}\renewcommand{\rhcaQdLgbmFullBarOos}{0.1621}
\providecommand{\rhcaOverSeLgbmFullBarOos}{}\renewcommand{\rhcaOverSeLgbmFullBarOos}{49.1}
\providecommand{\rhcaMrXgbHarBarOos}{}\renewcommand{\rhcaMrXgbHarBarOos}{1.069}
\providecommand{\rhcaRmXgbHarBarOos}{}\renewcommand{\rhcaRmXgbHarBarOos}{1.028}
\providecommand{\rhcaUnderXgbHarBarOos}{}\renewcommand{\rhcaUnderXgbHarBarOos}{38.4}
\providecommand{\rhcaMzAXgbHarBarOos}{}\renewcommand{\rhcaMzAXgbHarBarOos}{0.69}
\providecommand{\rhcaMzBXgbHarBarOos}{}\renewcommand{\rhcaMzBXgbHarBarOos}{0.94}
\providecommand{\rhcaMzBseXgbHarBarOos}{}\renewcommand{\rhcaMzBseXgbHarBarOos}{0.06}
\providecommand{\rhcaMzRsqXgbHarBarOos}{}\renewcommand{\rhcaMzRsqXgbHarBarOos}{0.66}
\providecommand{\rhcaMzlAXgbHarBarOos}{}\renewcommand{\rhcaMzlAXgbHarBarOos}{0.04}
\providecommand{\rhcaMzlBXgbHarBarOos}{}\renewcommand{\rhcaMzlBXgbHarBarOos}{1.01}
\providecommand{\rhcaMzlBseXgbHarBarOos}{}\renewcommand{\rhcaMzlBseXgbHarBarOos}{0.01}
\providecommand{\rhcaMzlRsqXgbHarBarOos}{}\renewcommand{\rhcaMzlRsqXgbHarBarOos}{0.80}
\providecommand{\rhcaShBiasXgbHarBarOos}{}\renewcommand{\rhcaShBiasXgbHarBarOos}{0.0}
\providecommand{\rhcaShSlopeXgbHarBarOos}{}\renewcommand{\rhcaShSlopeXgbHarBarOos}{0.8}
\providecommand{\rhcaShDistXgbHarBarOos}{}\renewcommand{\rhcaShDistXgbHarBarOos}{99.2}
\providecommand{\rhcaQlXgbHarBarOos}{}\renewcommand{\rhcaQlXgbHarBarOos}{0.0023}
\providecommand{\rhcaQdXgbHarBarOos}{}\renewcommand{\rhcaQdXgbHarBarOos}{0.1899}
\providecommand{\rhcaOverSeXgbHarBarOos}{}\renewcommand{\rhcaOverSeXgbHarBarOos}{25.9}
\providecommand{\rhcaMrXgbFullBarOos}{}\renewcommand{\rhcaMrXgbFullBarOos}{1.068}
\providecommand{\rhcaRmXgbFullBarOos}{}\renewcommand{\rhcaRmXgbFullBarOos}{0.966}
\providecommand{\rhcaUnderXgbFullBarOos}{}\renewcommand{\rhcaUnderXgbFullBarOos}{40.5}
\providecommand{\rhcaMzAXgbFullBarOos}{}\renewcommand{\rhcaMzAXgbFullBarOos}{1.74}
\providecommand{\rhcaMzBXgbFullBarOos}{}\renewcommand{\rhcaMzBXgbFullBarOos}{0.76}
\providecommand{\rhcaMzBseXgbFullBarOos}{}\renewcommand{\rhcaMzBseXgbFullBarOos}{0.06}
\providecommand{\rhcaMzRsqXgbFullBarOos}{}\renewcommand{\rhcaMzRsqXgbFullBarOos}{0.66}
\providecommand{\rhcaMzlAXgbFullBarOos}{}\renewcommand{\rhcaMzlAXgbFullBarOos}{\ensuremath{-}0.18}
\providecommand{\rhcaMzlBXgbFullBarOos}{}\renewcommand{\rhcaMzlBXgbFullBarOos}{0.99}
\providecommand{\rhcaMzlBseXgbFullBarOos}{}\renewcommand{\rhcaMzlBseXgbFullBarOos}{0.01}
\providecommand{\rhcaMzlRsqXgbFullBarOos}{}\renewcommand{\rhcaMzlRsqXgbFullBarOos}{0.82}
\providecommand{\rhcaShBiasXgbFullBarOos}{}\renewcommand{\rhcaShBiasXgbFullBarOos}{0.0}
\providecommand{\rhcaShSlopeXgbFullBarOos}{}\renewcommand{\rhcaShSlopeXgbFullBarOos}{16.1}
\providecommand{\rhcaShDistXgbFullBarOos}{}\renewcommand{\rhcaShDistXgbFullBarOos}{83.9}
\providecommand{\rhcaQlXgbFullBarOos}{}\renewcommand{\rhcaQlXgbFullBarOos}{0.0022}
\providecommand{\rhcaQdXgbFullBarOos}{}\renewcommand{\rhcaQdXgbFullBarOos}{0.1612}
\providecommand{\rhcaOverSeXgbFullBarOos}{}\renewcommand{\rhcaOverSeXgbFullBarOos}{48.3}
\providecommand{\rhcaMrRfHarBarOos}{}\renewcommand{\rhcaMrRfHarBarOos}{1.077}
\providecommand{\rhcaRmRfHarBarOos}{}\renewcommand{\rhcaRmRfHarBarOos}{1.024}
\providecommand{\rhcaUnderRfHarBarOos}{}\renewcommand{\rhcaUnderRfHarBarOos}{38.8}
\providecommand{\rhcaMzARfHarBarOos}{}\renewcommand{\rhcaMzARfHarBarOos}{0.82}
\providecommand{\rhcaMzBRfHarBarOos}{}\renewcommand{\rhcaMzBRfHarBarOos}{0.92}
\providecommand{\rhcaMzBseRfHarBarOos}{}\renewcommand{\rhcaMzBseRfHarBarOos}{0.07}
\providecommand{\rhcaMzRsqRfHarBarOos}{}\renewcommand{\rhcaMzRsqRfHarBarOos}{0.65}
\providecommand{\rhcaMzlARfHarBarOos}{}\renewcommand{\rhcaMzlARfHarBarOos}{\ensuremath{-}0.07}
\providecommand{\rhcaMzlBRfHarBarOos}{}\renewcommand{\rhcaMzlBRfHarBarOos}{1.00}
\providecommand{\rhcaMzlBseRfHarBarOos}{}\renewcommand{\rhcaMzlBseRfHarBarOos}{0.01}
\providecommand{\rhcaMzlRsqRfHarBarOos}{}\renewcommand{\rhcaMzlRsqRfHarBarOos}{0.80}
\providecommand{\rhcaShBiasRfHarBarOos}{}\renewcommand{\rhcaShBiasRfHarBarOos}{0.0}
\providecommand{\rhcaShSlopeRfHarBarOos}{}\renewcommand{\rhcaShSlopeRfHarBarOos}{1.4}
\providecommand{\rhcaShDistRfHarBarOos}{}\renewcommand{\rhcaShDistRfHarBarOos}{98.6}
\providecommand{\rhcaQlRfHarBarOos}{}\renewcommand{\rhcaQlRfHarBarOos}{0.0028}
\providecommand{\rhcaQdRfHarBarOos}{}\renewcommand{\rhcaQdRfHarBarOos}{0.1909}
\providecommand{\rhcaOverSeRfHarBarOos}{}\renewcommand{\rhcaOverSeRfHarBarOos}{26.7}
\providecommand{\rhcaMrRfFullBarOos}{}\renewcommand{\rhcaMrRfFullBarOos}{1.087}
\providecommand{\rhcaRmRfFullBarOos}{}\renewcommand{\rhcaRmRfFullBarOos}{0.974}
\providecommand{\rhcaUnderRfFullBarOos}{}\renewcommand{\rhcaUnderRfFullBarOos}{41.0}
\providecommand{\rhcaMzARfFullBarOos}{}\renewcommand{\rhcaMzARfFullBarOos}{1.61}
\providecommand{\rhcaMzBRfFullBarOos}{}\renewcommand{\rhcaMzBRfFullBarOos}{0.78}
\providecommand{\rhcaMzBseRfFullBarOos}{}\renewcommand{\rhcaMzBseRfFullBarOos}{0.07}
\providecommand{\rhcaMzRsqRfFullBarOos}{}\renewcommand{\rhcaMzRsqRfFullBarOos}{0.65}
\providecommand{\rhcaMzlARfFullBarOos}{}\renewcommand{\rhcaMzlARfFullBarOos}{\ensuremath{-}0.25}
\providecommand{\rhcaMzlBRfFullBarOos}{}\renewcommand{\rhcaMzlBRfFullBarOos}{0.99}
\providecommand{\rhcaMzlBseRfFullBarOos}{}\renewcommand{\rhcaMzlBseRfFullBarOos}{0.01}
\providecommand{\rhcaMzlRsqRfFullBarOos}{}\renewcommand{\rhcaMzlRsqRfFullBarOos}{0.82}
\providecommand{\rhcaShBiasRfFullBarOos}{}\renewcommand{\rhcaShBiasRfFullBarOos}{0.0}
\providecommand{\rhcaShSlopeRfFullBarOos}{}\renewcommand{\rhcaShSlopeRfFullBarOos}{12.7}
\providecommand{\rhcaShDistRfFullBarOos}{}\renewcommand{\rhcaShDistRfFullBarOos}{87.3}
\providecommand{\rhcaQlRfFullBarOos}{}\renewcommand{\rhcaQlRfFullBarOos}{0.0035}
\providecommand{\rhcaQdRfFullBarOos}{}\renewcommand{\rhcaQdRfFullBarOos}{0.1694}
\providecommand{\rhcaOverSeRfFullBarOos}{}\renewcommand{\rhcaOverSeRfFullBarOos}{44.4}
\providecommand{\rhcaMrRidgeHarCloseOos}{}\renewcommand{\rhcaMrRidgeHarCloseOos}{1.092}
\providecommand{\rhcaRmRidgeHarCloseOos}{}\renewcommand{\rhcaRmRidgeHarCloseOos}{1.103}
\providecommand{\rhcaUnderRidgeHarCloseOos}{}\renewcommand{\rhcaUnderRidgeHarCloseOos}{47.7}
\providecommand{\rhcaMzARidgeHarCloseOos}{}\renewcommand{\rhcaMzARidgeHarCloseOos}{\ensuremath{-}1.23}
\providecommand{\rhcaMzBRidgeHarCloseOos}{}\renewcommand{\rhcaMzBRidgeHarCloseOos}{1.21}
\providecommand{\rhcaMzBseRidgeHarCloseOos}{}\renewcommand{\rhcaMzBseRidgeHarCloseOos}{0.09}
\providecommand{\rhcaMzRsqRidgeHarCloseOos}{}\renewcommand{\rhcaMzRsqRidgeHarCloseOos}{0.80}
\providecommand{\rhcaMzlARidgeHarCloseOos}{}\renewcommand{\rhcaMzlARidgeHarCloseOos}{\ensuremath{-}0.38}
\providecommand{\rhcaMzlBRidgeHarCloseOos}{}\renewcommand{\rhcaMzlBRidgeHarCloseOos}{0.97}
\providecommand{\rhcaMzlBseRidgeHarCloseOos}{}\renewcommand{\rhcaMzlBseRidgeHarCloseOos}{0.02}
\providecommand{\rhcaMzlRsqRidgeHarCloseOos}{}\renewcommand{\rhcaMzlRsqRidgeHarCloseOos}{0.84}
\providecommand{\rhcaShBiasRidgeHarCloseOos}{}\renewcommand{\rhcaShBiasRidgeHarCloseOos}{0.3}
\providecommand{\rhcaShSlopeRidgeHarCloseOos}{}\renewcommand{\rhcaShSlopeRidgeHarCloseOos}{10.7}
\providecommand{\rhcaShDistRidgeHarCloseOos}{}\renewcommand{\rhcaShDistRidgeHarCloseOos}{88.9}
\providecommand{\rhcaQlRidgeHarCloseOos}{}\renewcommand{\rhcaQlRidgeHarCloseOos}{0.0040}
\providecommand{\rhcaQdRidgeHarCloseOos}{}\renewcommand{\rhcaQdRidgeHarCloseOos}{0.1097}
\providecommand{\rhcaOverSeRidgeHarCloseOos}{}\renewcommand{\rhcaOverSeRidgeHarCloseOos}{13.8}
\providecommand{\rhcaMrRidgeFullCloseOos}{}\renewcommand{\rhcaMrRidgeFullCloseOos}{1.169}
\providecommand{\rhcaRmRidgeFullCloseOos}{}\renewcommand{\rhcaRmRidgeFullCloseOos}{0.971}
\providecommand{\rhcaUnderRidgeFullCloseOos}{}\renewcommand{\rhcaUnderRidgeFullCloseOos}{52.1}
\providecommand{\rhcaMzARidgeFullCloseOos}{}\renewcommand{\rhcaMzARidgeFullCloseOos}{3.60}
\providecommand{\rhcaMzBRidgeFullCloseOos}{}\renewcommand{\rhcaMzBRidgeFullCloseOos}{0.70}
\providecommand{\rhcaMzBseRidgeFullCloseOos}{}\renewcommand{\rhcaMzBseRidgeFullCloseOos}{0.06}
\providecommand{\rhcaMzRsqRidgeFullCloseOos}{}\renewcommand{\rhcaMzRsqRidgeFullCloseOos}{0.72}
\providecommand{\rhcaMzlARidgeFullCloseOos}{}\renewcommand{\rhcaMzlARidgeFullCloseOos}{\ensuremath{-}0.95}
\providecommand{\rhcaMzlBRidgeFullCloseOos}{}\renewcommand{\rhcaMzlBRidgeFullCloseOos}{0.92}
\providecommand{\rhcaMzlBseRidgeFullCloseOos}{}\renewcommand{\rhcaMzlBseRidgeFullCloseOos}{0.01}
\providecommand{\rhcaMzlRsqRidgeFullCloseOos}{}\renewcommand{\rhcaMzlRsqRidgeFullCloseOos}{0.84}
\providecommand{\rhcaShBiasRidgeFullCloseOos}{}\renewcommand{\rhcaShBiasRidgeFullCloseOos}{0.0}
\providecommand{\rhcaShSlopeRidgeFullCloseOos}{}\renewcommand{\rhcaShSlopeRidgeFullCloseOos}{32.0}
\providecommand{\rhcaShDistRidgeFullCloseOos}{}\renewcommand{\rhcaShDistRidgeFullCloseOos}{68.0}
\providecommand{\rhcaQlRidgeFullCloseOos}{}\renewcommand{\rhcaQlRidgeFullCloseOos}{0.0128}
\providecommand{\rhcaQdRidgeFullCloseOos}{}\renewcommand{\rhcaQdRidgeFullCloseOos}{0.1140}
\providecommand{\rhcaOverSeRidgeFullCloseOos}{}\renewcommand{\rhcaOverSeRidgeFullCloseOos}{52.6}
\providecommand{\rhcaMrLassoHarCloseOos}{}\renewcommand{\rhcaMrLassoHarCloseOos}{1.092}
\providecommand{\rhcaRmLassoHarCloseOos}{}\renewcommand{\rhcaRmLassoHarCloseOos}{1.103}
\providecommand{\rhcaUnderLassoHarCloseOos}{}\renewcommand{\rhcaUnderLassoHarCloseOos}{47.7}
\providecommand{\rhcaMzALassoHarCloseOos}{}\renewcommand{\rhcaMzALassoHarCloseOos}{\ensuremath{-}1.23}
\providecommand{\rhcaMzBLassoHarCloseOos}{}\renewcommand{\rhcaMzBLassoHarCloseOos}{1.21}
\providecommand{\rhcaMzBseLassoHarCloseOos}{}\renewcommand{\rhcaMzBseLassoHarCloseOos}{0.09}
\providecommand{\rhcaMzRsqLassoHarCloseOos}{}\renewcommand{\rhcaMzRsqLassoHarCloseOos}{0.80}
\providecommand{\rhcaMzlALassoHarCloseOos}{}\renewcommand{\rhcaMzlALassoHarCloseOos}{\ensuremath{-}0.38}
\providecommand{\rhcaMzlBLassoHarCloseOos}{}\renewcommand{\rhcaMzlBLassoHarCloseOos}{0.97}
\providecommand{\rhcaMzlBseLassoHarCloseOos}{}\renewcommand{\rhcaMzlBseLassoHarCloseOos}{0.02}
\providecommand{\rhcaMzlRsqLassoHarCloseOos}{}\renewcommand{\rhcaMzlRsqLassoHarCloseOos}{0.84}
\providecommand{\rhcaShBiasLassoHarCloseOos}{}\renewcommand{\rhcaShBiasLassoHarCloseOos}{0.3}
\providecommand{\rhcaShSlopeLassoHarCloseOos}{}\renewcommand{\rhcaShSlopeLassoHarCloseOos}{10.7}
\providecommand{\rhcaShDistLassoHarCloseOos}{}\renewcommand{\rhcaShDistLassoHarCloseOos}{88.9}
\providecommand{\rhcaQlLassoHarCloseOos}{}\renewcommand{\rhcaQlLassoHarCloseOos}{0.0040}
\providecommand{\rhcaQdLassoHarCloseOos}{}\renewcommand{\rhcaQdLassoHarCloseOos}{0.1097}
\providecommand{\rhcaOverSeLassoHarCloseOos}{}\renewcommand{\rhcaOverSeLassoHarCloseOos}{13.8}
\providecommand{\rhcaMrLassoFullCloseOos}{}\renewcommand{\rhcaMrLassoFullCloseOos}{1.183}
\providecommand{\rhcaRmLassoFullCloseOos}{}\renewcommand{\rhcaRmLassoFullCloseOos}{0.982}
\providecommand{\rhcaUnderLassoFullCloseOos}{}\renewcommand{\rhcaUnderLassoFullCloseOos}{52.6}
\providecommand{\rhcaMzALassoFullCloseOos}{}\renewcommand{\rhcaMzALassoFullCloseOos}{3.74}
\providecommand{\rhcaMzBLassoFullCloseOos}{}\renewcommand{\rhcaMzBLassoFullCloseOos}{0.69}
\providecommand{\rhcaMzBseLassoFullCloseOos}{}\renewcommand{\rhcaMzBseLassoFullCloseOos}{0.05}
\providecommand{\rhcaMzRsqLassoFullCloseOos}{}\renewcommand{\rhcaMzRsqLassoFullCloseOos}{0.71}
\providecommand{\rhcaMzlALassoFullCloseOos}{}\renewcommand{\rhcaMzlALassoFullCloseOos}{\ensuremath{-}0.93}
\providecommand{\rhcaMzlBLassoFullCloseOos}{}\renewcommand{\rhcaMzlBLassoFullCloseOos}{0.92}
\providecommand{\rhcaMzlBseLassoFullCloseOos}{}\renewcommand{\rhcaMzlBseLassoFullCloseOos}{0.01}
\providecommand{\rhcaMzlRsqLassoFullCloseOos}{}\renewcommand{\rhcaMzlRsqLassoFullCloseOos}{0.84}
\providecommand{\rhcaShBiasLassoFullCloseOos}{}\renewcommand{\rhcaShBiasLassoFullCloseOos}{0.0}
\providecommand{\rhcaShSlopeLassoFullCloseOos}{}\renewcommand{\rhcaShSlopeLassoFullCloseOos}{32.1}
\providecommand{\rhcaShDistLassoFullCloseOos}{}\renewcommand{\rhcaShDistLassoFullCloseOos}{67.9}
\providecommand{\rhcaQlLassoFullCloseOos}{}\renewcommand{\rhcaQlLassoFullCloseOos}{0.0150}
\providecommand{\rhcaQdLassoFullCloseOos}{}\renewcommand{\rhcaQdLassoFullCloseOos}{0.1159}
\providecommand{\rhcaOverSeLassoFullCloseOos}{}\renewcommand{\rhcaOverSeLassoFullCloseOos}{51.5}
\providecommand{\rhcaMrEnetHarCloseOos}{}\renewcommand{\rhcaMrEnetHarCloseOos}{1.092}
\providecommand{\rhcaRmEnetHarCloseOos}{}\renewcommand{\rhcaRmEnetHarCloseOos}{1.103}
\providecommand{\rhcaUnderEnetHarCloseOos}{}\renewcommand{\rhcaUnderEnetHarCloseOos}{47.7}
\providecommand{\rhcaMzAEnetHarCloseOos}{}\renewcommand{\rhcaMzAEnetHarCloseOos}{\ensuremath{-}1.23}
\providecommand{\rhcaMzBEnetHarCloseOos}{}\renewcommand{\rhcaMzBEnetHarCloseOos}{1.21}
\providecommand{\rhcaMzBseEnetHarCloseOos}{}\renewcommand{\rhcaMzBseEnetHarCloseOos}{0.09}
\providecommand{\rhcaMzRsqEnetHarCloseOos}{}\renewcommand{\rhcaMzRsqEnetHarCloseOos}{0.80}
\providecommand{\rhcaMzlAEnetHarCloseOos}{}\renewcommand{\rhcaMzlAEnetHarCloseOos}{\ensuremath{-}0.38}
\providecommand{\rhcaMzlBEnetHarCloseOos}{}\renewcommand{\rhcaMzlBEnetHarCloseOos}{0.97}
\providecommand{\rhcaMzlBseEnetHarCloseOos}{}\renewcommand{\rhcaMzlBseEnetHarCloseOos}{0.02}
\providecommand{\rhcaMzlRsqEnetHarCloseOos}{}\renewcommand{\rhcaMzlRsqEnetHarCloseOos}{0.84}
\providecommand{\rhcaShBiasEnetHarCloseOos}{}\renewcommand{\rhcaShBiasEnetHarCloseOos}{0.3}
\providecommand{\rhcaShSlopeEnetHarCloseOos}{}\renewcommand{\rhcaShSlopeEnetHarCloseOos}{10.7}
\providecommand{\rhcaShDistEnetHarCloseOos}{}\renewcommand{\rhcaShDistEnetHarCloseOos}{88.9}
\providecommand{\rhcaQlEnetHarCloseOos}{}\renewcommand{\rhcaQlEnetHarCloseOos}{0.0040}
\providecommand{\rhcaQdEnetHarCloseOos}{}\renewcommand{\rhcaQdEnetHarCloseOos}{0.1097}
\providecommand{\rhcaOverSeEnetHarCloseOos}{}\renewcommand{\rhcaOverSeEnetHarCloseOos}{13.8}
\providecommand{\rhcaMrEnetFullCloseOos}{}\renewcommand{\rhcaMrEnetFullCloseOos}{1.165}
\providecommand{\rhcaRmEnetFullCloseOos}{}\renewcommand{\rhcaRmEnetFullCloseOos}{0.955}
\providecommand{\rhcaUnderEnetFullCloseOos}{}\renewcommand{\rhcaUnderEnetFullCloseOos}{52.9}
\providecommand{\rhcaMzAEnetFullCloseOos}{}\renewcommand{\rhcaMzAEnetFullCloseOos}{3.93}
\providecommand{\rhcaMzBEnetFullCloseOos}{}\renewcommand{\rhcaMzBEnetFullCloseOos}{0.66}
\providecommand{\rhcaMzBseEnetFullCloseOos}{}\renewcommand{\rhcaMzBseEnetFullCloseOos}{0.05}
\providecommand{\rhcaMzRsqEnetFullCloseOos}{}\renewcommand{\rhcaMzRsqEnetFullCloseOos}{0.72}
\providecommand{\rhcaMzlAEnetFullCloseOos}{}\renewcommand{\rhcaMzlAEnetFullCloseOos}{\ensuremath{-}0.94}
\providecommand{\rhcaMzlBEnetFullCloseOos}{}\renewcommand{\rhcaMzlBEnetFullCloseOos}{0.92}
\providecommand{\rhcaMzlBseEnetFullCloseOos}{}\renewcommand{\rhcaMzlBseEnetFullCloseOos}{0.01}
\providecommand{\rhcaMzlRsqEnetFullCloseOos}{}\renewcommand{\rhcaMzlRsqEnetFullCloseOos}{0.84}
\providecommand{\rhcaShBiasEnetFullCloseOos}{}\renewcommand{\rhcaShBiasEnetFullCloseOos}{0.0}
\providecommand{\rhcaShSlopeEnetFullCloseOos}{}\renewcommand{\rhcaShSlopeEnetFullCloseOos}{39.7}
\providecommand{\rhcaShDistEnetFullCloseOos}{}\renewcommand{\rhcaShDistEnetFullCloseOos}{60.3}
\providecommand{\rhcaQlEnetFullCloseOos}{}\renewcommand{\rhcaQlEnetFullCloseOos}{0.0123}
\providecommand{\rhcaQdEnetFullCloseOos}{}\renewcommand{\rhcaQdEnetFullCloseOos}{0.1122}
\providecommand{\rhcaOverSeEnetFullCloseOos}{}\renewcommand{\rhcaOverSeEnetFullCloseOos}{59.7}
\providecommand{\rhcaMrLgbmHarCloseOos}{}\renewcommand{\rhcaMrLgbmHarCloseOos}{1.035}
\providecommand{\rhcaRmLgbmHarCloseOos}{}\renewcommand{\rhcaRmLgbmHarCloseOos}{1.028}
\providecommand{\rhcaUnderLgbmHarCloseOos}{}\renewcommand{\rhcaUnderLgbmHarCloseOos}{42.9}
\providecommand{\rhcaMzALgbmHarCloseOos}{}\renewcommand{\rhcaMzALgbmHarCloseOos}{\ensuremath{-}0.68}
\providecommand{\rhcaMzBLgbmHarCloseOos}{}\renewcommand{\rhcaMzBLgbmHarCloseOos}{1.08}
\providecommand{\rhcaMzBseLgbmHarCloseOos}{}\renewcommand{\rhcaMzBseLgbmHarCloseOos}{0.06}
\providecommand{\rhcaMzRsqLgbmHarCloseOos}{}\renewcommand{\rhcaMzRsqLgbmHarCloseOos}{0.81}
\providecommand{\rhcaMzlALgbmHarCloseOos}{}\renewcommand{\rhcaMzlALgbmHarCloseOos}{\ensuremath{-}0.33}
\providecommand{\rhcaMzlBLgbmHarCloseOos}{}\renewcommand{\rhcaMzlBLgbmHarCloseOos}{0.98}
\providecommand{\rhcaMzlBseLgbmHarCloseOos}{}\renewcommand{\rhcaMzlBseLgbmHarCloseOos}{0.02}
\providecommand{\rhcaMzlRsqLgbmHarCloseOos}{}\renewcommand{\rhcaMzlRsqLgbmHarCloseOos}{0.84}
\providecommand{\rhcaShBiasLgbmHarCloseOos}{}\renewcommand{\rhcaShBiasLgbmHarCloseOos}{0.0}
\providecommand{\rhcaShSlopeLgbmHarCloseOos}{}\renewcommand{\rhcaShSlopeLgbmHarCloseOos}{2.4}
\providecommand{\rhcaShDistLgbmHarCloseOos}{}\renewcommand{\rhcaShDistLgbmHarCloseOos}{97.6}
\providecommand{\rhcaQlLgbmHarCloseOos}{}\renewcommand{\rhcaQlLgbmHarCloseOos}{0.0006}
\providecommand{\rhcaQdLgbmHarCloseOos}{}\renewcommand{\rhcaQdLgbmHarCloseOos}{0.1095}
\providecommand{\rhcaOverSeLgbmHarCloseOos}{}\renewcommand{\rhcaOverSeLgbmHarCloseOos}{22.8}
\providecommand{\rhcaMrLgbmFullCloseOos}{}\renewcommand{\rhcaMrLgbmFullCloseOos}{1.093}
\providecommand{\rhcaRmLgbmFullCloseOos}{}\renewcommand{\rhcaRmLgbmFullCloseOos}{0.969}
\providecommand{\rhcaUnderLgbmFullCloseOos}{}\renewcommand{\rhcaUnderLgbmFullCloseOos}{47.1}
\providecommand{\rhcaMzALgbmFullCloseOos}{}\renewcommand{\rhcaMzALgbmFullCloseOos}{2.03}
\providecommand{\rhcaMzBLgbmFullCloseOos}{}\renewcommand{\rhcaMzBLgbmFullCloseOos}{0.81}
\providecommand{\rhcaMzBseLgbmFullCloseOos}{}\renewcommand{\rhcaMzBseLgbmFullCloseOos}{0.05}
\providecommand{\rhcaMzRsqLgbmFullCloseOos}{}\renewcommand{\rhcaMzRsqLgbmFullCloseOos}{0.77}
\providecommand{\rhcaMzlALgbmFullCloseOos}{}\renewcommand{\rhcaMzlALgbmFullCloseOos}{\ensuremath{-}0.69}
\providecommand{\rhcaMzlBLgbmFullCloseOos}{}\renewcommand{\rhcaMzlBLgbmFullCloseOos}{0.95}
\providecommand{\rhcaMzlBseLgbmFullCloseOos}{}\renewcommand{\rhcaMzlBseLgbmFullCloseOos}{0.01}
\providecommand{\rhcaMzlRsqLgbmFullCloseOos}{}\renewcommand{\rhcaMzlRsqLgbmFullCloseOos}{0.85}
\providecommand{\rhcaShBiasLgbmFullCloseOos}{}\renewcommand{\rhcaShBiasLgbmFullCloseOos}{0.0}
\providecommand{\rhcaShSlopeLgbmFullCloseOos}{}\renewcommand{\rhcaShSlopeLgbmFullCloseOos}{14.8}
\providecommand{\rhcaShDistLgbmFullCloseOos}{}\renewcommand{\rhcaShDistLgbmFullCloseOos}{85.2}
\providecommand{\rhcaQlLgbmFullCloseOos}{}\renewcommand{\rhcaQlLgbmFullCloseOos}{0.0041}
\providecommand{\rhcaQdLgbmFullCloseOos}{}\renewcommand{\rhcaQdLgbmFullCloseOos}{0.1045}
\providecommand{\rhcaOverSeLgbmFullCloseOos}{}\renewcommand{\rhcaOverSeLgbmFullCloseOos}{44.8}
\providecommand{\rhcaMrXgbHarCloseOos}{}\renewcommand{\rhcaMrXgbHarCloseOos}{1.033}
\providecommand{\rhcaRmXgbHarCloseOos}{}\renewcommand{\rhcaRmXgbHarCloseOos}{1.029}
\providecommand{\rhcaUnderXgbHarCloseOos}{}\renewcommand{\rhcaUnderXgbHarCloseOos}{43.2}
\providecommand{\rhcaMzAXgbHarCloseOos}{}\renewcommand{\rhcaMzAXgbHarCloseOos}{\ensuremath{-}0.70}
\providecommand{\rhcaMzBXgbHarCloseOos}{}\renewcommand{\rhcaMzBXgbHarCloseOos}{1.08}
\providecommand{\rhcaMzBseXgbHarCloseOos}{}\renewcommand{\rhcaMzBseXgbHarCloseOos}{0.06}
\providecommand{\rhcaMzRsqXgbHarCloseOos}{}\renewcommand{\rhcaMzRsqXgbHarCloseOos}{0.81}
\providecommand{\rhcaMzlAXgbHarCloseOos}{}\renewcommand{\rhcaMzlAXgbHarCloseOos}{\ensuremath{-}0.30}
\providecommand{\rhcaMzlBXgbHarCloseOos}{}\renewcommand{\rhcaMzlBXgbHarCloseOos}{0.98}
\providecommand{\rhcaMzlBseXgbHarCloseOos}{}\renewcommand{\rhcaMzlBseXgbHarCloseOos}{0.02}
\providecommand{\rhcaMzlRsqXgbHarCloseOos}{}\renewcommand{\rhcaMzlRsqXgbHarCloseOos}{0.84}
\providecommand{\rhcaShBiasXgbHarCloseOos}{}\renewcommand{\rhcaShBiasXgbHarCloseOos}{0.0}
\providecommand{\rhcaShSlopeXgbHarCloseOos}{}\renewcommand{\rhcaShSlopeXgbHarCloseOos}{2.6}
\providecommand{\rhcaShDistXgbHarCloseOos}{}\renewcommand{\rhcaShDistXgbHarCloseOos}{97.4}
\providecommand{\rhcaQlXgbHarCloseOos}{}\renewcommand{\rhcaQlXgbHarCloseOos}{0.0005}
\providecommand{\rhcaQdXgbHarCloseOos}{}\renewcommand{\rhcaQdXgbHarCloseOos}{0.1096}
\providecommand{\rhcaOverSeXgbHarCloseOos}{}\renewcommand{\rhcaOverSeXgbHarCloseOos}{22.7}
\providecommand{\rhcaMrXgbFullCloseOos}{}\renewcommand{\rhcaMrXgbFullCloseOos}{1.091}
\providecommand{\rhcaRmXgbFullCloseOos}{}\renewcommand{\rhcaRmXgbFullCloseOos}{0.965}
\providecommand{\rhcaUnderXgbFullCloseOos}{}\renewcommand{\rhcaUnderXgbFullCloseOos}{47.9}
\providecommand{\rhcaMzAXgbFullCloseOos}{}\renewcommand{\rhcaMzAXgbFullCloseOos}{1.80}
\providecommand{\rhcaMzBXgbFullCloseOos}{}\renewcommand{\rhcaMzBXgbFullCloseOos}{0.83}
\providecommand{\rhcaMzBseXgbFullCloseOos}{}\renewcommand{\rhcaMzBseXgbFullCloseOos}{0.04}
\providecommand{\rhcaMzRsqXgbFullCloseOos}{}\renewcommand{\rhcaMzRsqXgbFullCloseOos}{0.78}
\providecommand{\rhcaMzlAXgbFullCloseOos}{}\renewcommand{\rhcaMzlAXgbFullCloseOos}{\ensuremath{-}0.68}
\providecommand{\rhcaMzlBXgbFullCloseOos}{}\renewcommand{\rhcaMzlBXgbFullCloseOos}{0.95}
\providecommand{\rhcaMzlBseXgbFullCloseOos}{}\renewcommand{\rhcaMzlBseXgbFullCloseOos}{0.01}
\providecommand{\rhcaMzlRsqXgbFullCloseOos}{}\renewcommand{\rhcaMzlRsqXgbFullCloseOos}{0.85}
\providecommand{\rhcaShBiasXgbFullCloseOos}{}\renewcommand{\rhcaShBiasXgbFullCloseOos}{0.0}
\providecommand{\rhcaShSlopeXgbFullCloseOos}{}\renewcommand{\rhcaShSlopeXgbFullCloseOos}{13.1}
\providecommand{\rhcaShDistXgbFullCloseOos}{}\renewcommand{\rhcaShDistXgbFullCloseOos}{86.8}
\providecommand{\rhcaQlXgbFullCloseOos}{}\renewcommand{\rhcaQlXgbFullCloseOos}{0.0039}
\providecommand{\rhcaQdXgbFullCloseOos}{}\renewcommand{\rhcaQdXgbFullCloseOos}{0.1044}
\providecommand{\rhcaOverSeXgbFullCloseOos}{}\renewcommand{\rhcaOverSeXgbFullCloseOos}{46.1}
\providecommand{\rhcaMrRfHarCloseOos}{}\renewcommand{\rhcaMrRfHarCloseOos}{1.063}
\providecommand{\rhcaRmRfHarCloseOos}{}\renewcommand{\rhcaRmRfHarCloseOos}{1.033}
\providecommand{\rhcaUnderRfHarCloseOos}{}\renewcommand{\rhcaUnderRfHarCloseOos}{45.5}
\providecommand{\rhcaMzARfHarCloseOos}{}\renewcommand{\rhcaMzARfHarCloseOos}{\ensuremath{-}0.38}
\providecommand{\rhcaMzBRfHarCloseOos}{}\renewcommand{\rhcaMzBRfHarCloseOos}{1.06}
\providecommand{\rhcaMzBseRfHarCloseOos}{}\renewcommand{\rhcaMzBseRfHarCloseOos}{0.07}
\providecommand{\rhcaMzRsqRfHarCloseOos}{}\renewcommand{\rhcaMzRsqRfHarCloseOos}{0.82}
\providecommand{\rhcaMzlARfHarCloseOos}{}\renewcommand{\rhcaMzlARfHarCloseOos}{\ensuremath{-}0.46}
\providecommand{\rhcaMzlBRfHarCloseOos}{}\renewcommand{\rhcaMzlBRfHarCloseOos}{0.97}
\providecommand{\rhcaMzlBseRfHarCloseOos}{}\renewcommand{\rhcaMzlBseRfHarCloseOos}{0.02}
\providecommand{\rhcaMzlRsqRfHarCloseOos}{}\renewcommand{\rhcaMzlRsqRfHarCloseOos}{0.83}
\providecommand{\rhcaShBiasRfHarCloseOos}{}\renewcommand{\rhcaShBiasRfHarCloseOos}{0.1}
\providecommand{\rhcaShSlopeRfHarCloseOos}{}\renewcommand{\rhcaShSlopeRfHarCloseOos}{1.7}
\providecommand{\rhcaShDistRfHarCloseOos}{}\renewcommand{\rhcaShDistRfHarCloseOos}{98.3}
\providecommand{\rhcaQlRfHarCloseOos}{}\renewcommand{\rhcaQlRfHarCloseOos}{0.0019}
\providecommand{\rhcaQdRfHarCloseOos}{}\renewcommand{\rhcaQdRfHarCloseOos}{0.1112}
\providecommand{\rhcaOverSeRfHarCloseOos}{}\renewcommand{\rhcaOverSeRfHarCloseOos}{25.7}
\providecommand{\rhcaMrRfFullCloseOos}{}\renewcommand{\rhcaMrRfFullCloseOos}{1.113}
\providecommand{\rhcaRmRfFullCloseOos}{}\renewcommand{\rhcaRmRfFullCloseOos}{0.978}
\providecommand{\rhcaUnderRfFullCloseOos}{}\renewcommand{\rhcaUnderRfFullCloseOos}{50.1}
\providecommand{\rhcaMzARfFullCloseOos}{}\renewcommand{\rhcaMzARfFullCloseOos}{1.15}
\providecommand{\rhcaMzBRfFullCloseOos}{}\renewcommand{\rhcaMzBRfFullCloseOos}{0.89}
\providecommand{\rhcaMzBseRfFullCloseOos}{}\renewcommand{\rhcaMzBseRfFullCloseOos}{0.06}
\providecommand{\rhcaMzRsqRfFullCloseOos}{}\renewcommand{\rhcaMzRsqRfFullCloseOos}{0.79}
\providecommand{\rhcaMzlARfFullCloseOos}{}\renewcommand{\rhcaMzlARfFullCloseOos}{\ensuremath{-}0.87}
\providecommand{\rhcaMzlBRfFullCloseOos}{}\renewcommand{\rhcaMzlBRfFullCloseOos}{0.93}
\providecommand{\rhcaMzlBseRfFullCloseOos}{}\renewcommand{\rhcaMzlBseRfFullCloseOos}{0.01}
\providecommand{\rhcaMzlRsqRfFullCloseOos}{}\renewcommand{\rhcaMzlRsqRfFullCloseOos}{0.85}
\providecommand{\rhcaShBiasRfFullCloseOos}{}\renewcommand{\rhcaShBiasRfFullCloseOos}{0.0}
\providecommand{\rhcaShSlopeRfFullCloseOos}{}\renewcommand{\rhcaShSlopeRfFullCloseOos}{5.5}
\providecommand{\rhcaShDistRfFullCloseOos}{}\renewcommand{\rhcaShDistRfFullCloseOos}{94.5}
\providecommand{\rhcaQlRfFullCloseOos}{}\renewcommand{\rhcaQlRfFullCloseOos}{0.0059}
\providecommand{\rhcaQdRfFullCloseOos}{}\renewcommand{\rhcaQdRfFullCloseOos}{0.1064}
\providecommand{\rhcaOverSeRfFullCloseOos}{}\renewcommand{\rhcaOverSeRfFullCloseOos}{43.1}
% --- aggregates for the prose
\providecommand{\rhcaNLowerCloseSpan}{}\renewcommand{\rhcaNLowerCloseSpan}{3}  % classes (of 6) with the lower 16:00 QLIKE on the full set, Span
\providecommand{\rhcaNLowerSigCloseSpan}{}\renewcommand{\rhcaNLowerSigCloseSpan}{1}  % ... and DM p < 0.05
\providecommand{\rhcaNHigherCloseSpan}{}\renewcommand{\rhcaNHigherCloseSpan}{3}  % classes with the higher 16:00 QLIKE on the full set
\providecommand{\rhcaNHigherSigCloseSpan}{}\renewcommand{\rhcaNHigherSigCloseSpan}{0}  % ... and DM p < 0.05
\providecommand{\rhcaNLowerCloseWordSpan}{}\renewcommand{\rhcaNLowerCloseWordSpan}{three}  % as a word
\providecommand{\rhcaNLinHigherSigCloseSpan}{}\renewcommand{\rhcaNLinHigherSigCloseSpan}{0}  % linear classes (of 3): full-set 16:00 QLIKE significantly higher
\providecommand{\rhcaNLinLowerCloseSpan}{}\renewcommand{\rhcaNLinLowerCloseSpan}{0}  % linear classes with a lower 16:00 QLIKE
\providecommand{\rhcaNTreeLowerSigCloseSpan}{}\renewcommand{\rhcaNTreeLowerSigCloseSpan}{1}  % tree classes (of 3): full-set 16:00 QLIKE significantly lower
\providecommand{\rhcaLinPctCloseSpanMin}{}\renewcommand{\rhcaLinPctCloseSpanMin}{0.6}  % linear 16:00 QLIKE %, full vs HAR, Span: min
\providecommand{\rhcaLinPctCloseSpanMax}{}\renewcommand{\rhcaLinPctCloseSpanMax}{4.2}  % ... max
\providecommand{\rhcaLinPctCloseSpanRng}{}\renewcommand{\rhcaLinPctCloseSpanRng}{0.6--4.2}  % ... the range as text
\providecommand{\rhcaTreePctCloseSpanMin}{}\renewcommand{\rhcaTreePctCloseSpanMin}{\ensuremath{-}3.1}  % tree 16:00 QLIKE %, full vs HAR, Span: min
\providecommand{\rhcaTreePctCloseSpanMax}{}\renewcommand{\rhcaTreePctCloseSpanMax}{\ensuremath{-}1.6}  % ... max
\providecommand{\rhcaTreePctCloseSpanRng}{}\renewcommand{\rhcaTreePctCloseSpanRng}{\ensuremath{-}3.1 to \ensuremath{-}1.6}  % ... the range as text
\providecommand{\rhcaLinMsePctCloseSpanMin}{}\renewcommand{\rhcaLinMsePctCloseSpanMin}{77.5}  % linear 16:00 MSE %, full vs HAR, Span: min
\providecommand{\rhcaLinMsePctCloseSpanMax}{}\renewcommand{\rhcaLinMsePctCloseSpanMax}{99.1}  % ... max
\providecommand{\rhcaLinMsePctCloseSpanRng}{}\renewcommand{\rhcaLinMsePctCloseSpanRng}{77.5--99.1}  % ... the range as text
\providecommand{\rhcaTreeMsePctCloseSpanMin}{}\renewcommand{\rhcaTreeMsePctCloseSpanMin}{19.7}  % tree 16:00 MSE %, full vs HAR, Span: min
\providecommand{\rhcaTreeMsePctCloseSpanMax}{}\renewcommand{\rhcaTreeMsePctCloseSpanMax}{33.0}  % ... max
\providecommand{\rhcaTreeMsePctCloseSpanRng}{}\renewcommand{\rhcaTreeMsePctCloseSpanRng}{19.7--33.0}  % ... the range as text
\providecommand{\rhcaNMseHigherCloseSpan}{}\renewcommand{\rhcaNMseHigherCloseSpan}{6}  % classes with a higher 16:00 MSE, full
\providecommand{\rhcaNMseHigherSigCloseSpan}{}\renewcommand{\rhcaNMseHigherSigCloseSpan}{0}  % ... significantly
\providecommand{\rhcaNLowerClosePre}{}\renewcommand{\rhcaNLowerClosePre}{6}  % classes (of 6) with the lower 16:00 QLIKE on the full set, Pre
\providecommand{\rhcaNLowerSigClosePre}{}\renewcommand{\rhcaNLowerSigClosePre}{2}  % ... and DM p < 0.05
\providecommand{\rhcaNHigherClosePre}{}\renewcommand{\rhcaNHigherClosePre}{0}  % classes with the higher 16:00 QLIKE on the full set
\providecommand{\rhcaNHigherSigClosePre}{}\renewcommand{\rhcaNHigherSigClosePre}{0}  % ... and DM p < 0.05
\providecommand{\rhcaNLowerCloseWordPre}{}\renewcommand{\rhcaNLowerCloseWordPre}{six}  % as a word
\providecommand{\rhcaNLinHigherSigClosePre}{}\renewcommand{\rhcaNLinHigherSigClosePre}{0}  % linear classes (of 3): full-set 16:00 QLIKE significantly higher
\providecommand{\rhcaNLinLowerClosePre}{}\renewcommand{\rhcaNLinLowerClosePre}{3}  % linear classes with a lower 16:00 QLIKE
\providecommand{\rhcaNTreeLowerSigClosePre}{}\renewcommand{\rhcaNTreeLowerSigClosePre}{2}  % tree classes (of 3): full-set 16:00 QLIKE significantly lower
\providecommand{\rhcaLinPctClosePreMin}{}\renewcommand{\rhcaLinPctClosePreMin}{\ensuremath{-}4.8}  % linear 16:00 QLIKE %, full vs HAR, Pre: min
\providecommand{\rhcaLinPctClosePreMax}{}\renewcommand{\rhcaLinPctClosePreMax}{\ensuremath{-}2.3}  % ... max
\providecommand{\rhcaLinPctClosePreRng}{}\renewcommand{\rhcaLinPctClosePreRng}{\ensuremath{-}4.8 to \ensuremath{-}2.3}  % ... the range as text
\providecommand{\rhcaTreePctClosePreMin}{}\renewcommand{\rhcaTreePctClosePreMin}{\ensuremath{-}4.0}  % tree 16:00 QLIKE %, full vs HAR, Pre: min
\providecommand{\rhcaTreePctClosePreMax}{}\renewcommand{\rhcaTreePctClosePreMax}{\ensuremath{-}2.1}  % ... max
\providecommand{\rhcaTreePctClosePreRng}{}\renewcommand{\rhcaTreePctClosePreRng}{\ensuremath{-}4.0 to \ensuremath{-}2.1}  % ... the range as text
\providecommand{\rhcaLinMsePctClosePreMin}{}\renewcommand{\rhcaLinMsePctClosePreMin}{\ensuremath{-}15.8}  % linear 16:00 MSE %, full vs HAR, Pre: min
\providecommand{\rhcaLinMsePctClosePreMax}{}\renewcommand{\rhcaLinMsePctClosePreMax}{\ensuremath{-}13.1}  % ... max
\providecommand{\rhcaLinMsePctClosePreRng}{}\renewcommand{\rhcaLinMsePctClosePreRng}{\ensuremath{-}15.8 to \ensuremath{-}13.1}  % ... the range as text
\providecommand{\rhcaTreeMsePctClosePreMin}{}\renewcommand{\rhcaTreeMsePctClosePreMin}{\ensuremath{-}10.0}  % tree 16:00 MSE %, full vs HAR, Pre: min
\providecommand{\rhcaTreeMsePctClosePreMax}{}\renewcommand{\rhcaTreeMsePctClosePreMax}{\ensuremath{-}5.3}  % ... max
\providecommand{\rhcaTreeMsePctClosePreRng}{}\renewcommand{\rhcaTreeMsePctClosePreRng}{\ensuremath{-}10.0 to \ensuremath{-}5.3}  % ... the range as text
\providecommand{\rhcaNMseHigherClosePre}{}\renewcommand{\rhcaNMseHigherClosePre}{0}  % classes with a higher 16:00 MSE, full
\providecommand{\rhcaNMseHigherSigClosePre}{}\renewcommand{\rhcaNMseHigherSigClosePre}{0}  % ... significantly
\providecommand{\rhcaNLowerCloseOos}{}\renewcommand{\rhcaNLowerCloseOos}{3}  % classes (of 6) with the lower 16:00 QLIKE on the full set, Oos
\providecommand{\rhcaNLowerSigCloseOos}{}\renewcommand{\rhcaNLowerSigCloseOos}{0}  % ... and DM p < 0.05
\providecommand{\rhcaNHigherCloseOos}{}\renewcommand{\rhcaNHigherCloseOos}{3}  % classes with the higher 16:00 QLIKE on the full set
\providecommand{\rhcaNHigherSigCloseOos}{}\renewcommand{\rhcaNHigherSigCloseOos}{3}  % ... and DM p < 0.05
\providecommand{\rhcaNLowerCloseWordOos}{}\renewcommand{\rhcaNLowerCloseWordOos}{three}  % as a word
\providecommand{\rhcaNLinHigherSigCloseOos}{}\renewcommand{\rhcaNLinHigherSigCloseOos}{3}  % linear classes (of 3): full-set 16:00 QLIKE significantly higher
\providecommand{\rhcaNLinLowerCloseOos}{}\renewcommand{\rhcaNLinLowerCloseOos}{0}  % linear classes with a lower 16:00 QLIKE
\providecommand{\rhcaNTreeLowerSigCloseOos}{}\renewcommand{\rhcaNTreeLowerSigCloseOos}{0}  % tree classes (of 3): full-set 16:00 QLIKE significantly lower
\providecommand{\rhcaLinPctCloseOosMin}{}\renewcommand{\rhcaLinPctCloseOosMin}{9.5}  % linear 16:00 QLIKE %, full vs HAR, Oos: min
\providecommand{\rhcaLinPctCloseOosMax}{}\renewcommand{\rhcaLinPctCloseOosMax}{15.2}  % ... max
\providecommand{\rhcaLinPctCloseOosRng}{}\renewcommand{\rhcaLinPctCloseOosRng}{9.5--15.2}  % ... the range as text
\providecommand{\rhcaTreePctCloseOosMin}{}\renewcommand{\rhcaTreePctCloseOosMin}{\ensuremath{-}1.7}  % tree 16:00 QLIKE %, full vs HAR, Oos: min
\providecommand{\rhcaTreePctCloseOosMax}{}\renewcommand{\rhcaTreePctCloseOosMax}{\ensuremath{-}0.7}  % ... max
\providecommand{\rhcaTreePctCloseOosRng}{}\renewcommand{\rhcaTreePctCloseOosRng}{\ensuremath{-}1.7 to \ensuremath{-}0.7}  % ... the range as text
\providecommand{\rhcaLinMsePctCloseOosMin}{}\renewcommand{\rhcaLinMsePctCloseOosMin}{87.0}  % linear 16:00 MSE %, full vs HAR, Oos: min
\providecommand{\rhcaLinMsePctCloseOosMax}{}\renewcommand{\rhcaLinMsePctCloseOosMax}{111.1}  % ... max
\providecommand{\rhcaLinMsePctCloseOosRng}{}\renewcommand{\rhcaLinMsePctCloseOosRng}{87.0--111.1}  % ... the range as text
\providecommand{\rhcaTreeMsePctCloseOosMin}{}\renewcommand{\rhcaTreeMsePctCloseOosMin}{23.3}  % tree 16:00 MSE %, full vs HAR, Oos: min
\providecommand{\rhcaTreeMsePctCloseOosMax}{}\renewcommand{\rhcaTreeMsePctCloseOosMax}{38.9}  % ... max
\providecommand{\rhcaTreeMsePctCloseOosRng}{}\renewcommand{\rhcaTreeMsePctCloseOosRng}{23.3--38.9}  % ... the range as text
\providecommand{\rhcaNMseHigherCloseOos}{}\renewcommand{\rhcaNMseHigherCloseOos}{6}  % classes with a higher 16:00 MSE, full
\providecommand{\rhcaNMseHigherSigCloseOos}{}\renewcommand{\rhcaNMseHigherSigCloseOos}{0}  % ... significantly
\providecommand{\rhcaNMseHigherBarSpan}{}\renewcommand{\rhcaNMseHigherBarSpan}{6}  % classes with a higher every-bar MSE, full, Span
\providecommand{\rhcaNMseHigherSigBarSpan}{}\renewcommand{\rhcaNMseHigherSigBarSpan}{0}  % ... significantly
\providecommand{\rhcaNMseLowerBarSpan}{}\renewcommand{\rhcaNMseLowerBarSpan}{0}  % classes with a lower every-bar MSE, full
\providecommand{\rhcaNMseLowerSigBarSpan}{}\renewcommand{\rhcaNMseLowerSigBarSpan}{0}  % ... significantly
\providecommand{\rhcaLinMsePctBarSpanMin}{}\renewcommand{\rhcaLinMsePctBarSpanMin}{100.6}  % linear every-bar MSE %, full vs HAR, Span: min
\providecommand{\rhcaLinMsePctBarSpanMax}{}\renewcommand{\rhcaLinMsePctBarSpanMax}{117.7}  % ... max
\providecommand{\rhcaLinMsePctBarSpanRng}{}\renewcommand{\rhcaLinMsePctBarSpanRng}{100.6--117.7}  % ... the range as text
\providecommand{\rhcaTreeMsePctBarSpanMin}{}\renewcommand{\rhcaTreeMsePctBarSpanMin}{10.9}  % tree every-bar MSE %, full vs HAR, Span: min
\providecommand{\rhcaTreeMsePctBarSpanMax}{}\renewcommand{\rhcaTreeMsePctBarSpanMax}{20.0}  % ... max
\providecommand{\rhcaTreeMsePctBarSpanRng}{}\renewcommand{\rhcaTreeMsePctBarSpanRng}{10.9--20.0}  % ... the range as text
\providecommand{\rhcaNMseHigherBarPre}{}\renewcommand{\rhcaNMseHigherBarPre}{0}  % classes with a higher every-bar MSE, full, Pre
\providecommand{\rhcaNMseHigherSigBarPre}{}\renewcommand{\rhcaNMseHigherSigBarPre}{0}  % ... significantly
\providecommand{\rhcaNMseLowerBarPre}{}\renewcommand{\rhcaNMseLowerBarPre}{6}  % classes with a lower every-bar MSE, full
\providecommand{\rhcaNMseLowerSigBarPre}{}\renewcommand{\rhcaNMseLowerSigBarPre}{2}  % ... significantly
\providecommand{\rhcaLinMsePctBarPreMin}{}\renewcommand{\rhcaLinMsePctBarPreMin}{\ensuremath{-}10.8}  % linear every-bar MSE %, full vs HAR, Pre: min
\providecommand{\rhcaLinMsePctBarPreMax}{}\renewcommand{\rhcaLinMsePctBarPreMax}{\ensuremath{-}8.9}  % ... max
\providecommand{\rhcaLinMsePctBarPreRng}{}\renewcommand{\rhcaLinMsePctBarPreRng}{\ensuremath{-}10.8 to \ensuremath{-}8.9}  % ... the range as text
\providecommand{\rhcaTreeMsePctBarPreMin}{}\renewcommand{\rhcaTreeMsePctBarPreMin}{\ensuremath{-}4.7}  % tree every-bar MSE %, full vs HAR, Pre: min
\providecommand{\rhcaTreeMsePctBarPreMax}{}\renewcommand{\rhcaTreeMsePctBarPreMax}{\ensuremath{-}1.6}  % ... max
\providecommand{\rhcaTreeMsePctBarPreRng}{}\renewcommand{\rhcaTreeMsePctBarPreRng}{\ensuremath{-}4.7 to \ensuremath{-}1.6}  % ... the range as text
\providecommand{\rhcaNMseHigherBarOos}{}\renewcommand{\rhcaNMseHigherBarOos}{6}  % classes with a higher every-bar MSE, full, Oos
\providecommand{\rhcaNMseHigherSigBarOos}{}\renewcommand{\rhcaNMseHigherSigBarOos}{0}  % ... significantly
\providecommand{\rhcaNMseLowerBarOos}{}\renewcommand{\rhcaNMseLowerBarOos}{0}  % classes with a lower every-bar MSE, full
\providecommand{\rhcaNMseLowerSigBarOos}{}\renewcommand{\rhcaNMseLowerSigBarOos}{0}  % ... significantly
\providecommand{\rhcaLinMsePctBarOosMin}{}\renewcommand{\rhcaLinMsePctBarOosMin}{121.6}  % linear every-bar MSE %, full vs HAR, Oos: min
\providecommand{\rhcaLinMsePctBarOosMax}{}\renewcommand{\rhcaLinMsePctBarOosMax}{142.1}  % ... max
\providecommand{\rhcaLinMsePctBarOosRng}{}\renewcommand{\rhcaLinMsePctBarOosRng}{121.6--142.1}  % ... the range as text
\providecommand{\rhcaTreeMsePctBarOosMin}{}\renewcommand{\rhcaTreeMsePctBarOosMin}{13.7}  % tree every-bar MSE %, full vs HAR, Oos: min
\providecommand{\rhcaTreeMsePctBarOosMax}{}\renewcommand{\rhcaTreeMsePctBarOosMax}{25.5}  % ... max
\providecommand{\rhcaTreeMsePctBarOosRng}{}\renewcommand{\rhcaTreeMsePctBarOosRng}{13.7--25.5}  % ... the range as text
\providecommand{\rhcaNCiLowerBm}{}\renewcommand{\rhcaNCiLowerBm}{2}  % closing-interval models (of 12) below the benchmark's QLIKE
\providecommand{\rhcaNCiLowerSigBm}{}\renewcommand{\rhcaNCiLowerSigBm}{0}  % ... significantly
\providecommand{\rhcaNCi}{}\renewcommand{\rhcaNCi}{12}  % closing-interval models other than the benchmark
\providecommand{\rhcaCiHarGapBound}{}\renewcommand{\rhcaCiHarGapBound}{0.8}  % |QLIKE %| of the closing-interval HAR-set linear models vs the benchmark is below
\providecommand{\rhcaLinFullMzBBarPreMin}{}\renewcommand{\rhcaLinFullMzBBarPreMin}{0.99}  % levels MZ slope, linear full, BarPre: min
\providecommand{\rhcaLinFullMzBBarPreMax}{}\renewcommand{\rhcaLinFullMzBBarPreMax}{1.00}  % ... max
\providecommand{\rhcaLinFullMzBBarPreRng}{}\renewcommand{\rhcaLinFullMzBBarPreRng}{0.99--1.00}  % ... the range as text
\providecommand{\rhcaTreeFullMzBBarPreMin}{}\renewcommand{\rhcaTreeFullMzBBarPreMin}{0.92}  % levels MZ slope, tree full, BarPre: min
\providecommand{\rhcaTreeFullMzBBarPreMax}{}\renewcommand{\rhcaTreeFullMzBBarPreMax}{0.97}  % ... max
\providecommand{\rhcaTreeFullMzBBarPreRng}{}\renewcommand{\rhcaTreeFullMzBBarPreRng}{0.92--0.97}  % ... the range as text
\providecommand{\rhcaTreeHarMzBBarPreMin}{}\renewcommand{\rhcaTreeHarMzBBarPreMin}{0.96}  % levels MZ slope, tree HAR, BarPre: min
\providecommand{\rhcaTreeHarMzBBarPreMax}{}\renewcommand{\rhcaTreeHarMzBBarPreMax}{0.99}  % ... max
\providecommand{\rhcaTreeHarMzBBarPreRng}{}\renewcommand{\rhcaTreeHarMzBBarPreRng}{0.96--0.99}  % ... the range as text
\providecommand{\rhcaLinHarMzBBarPre}{}\renewcommand{\rhcaLinHarMzBBarPre}{0.97}  % levels MZ slope, linear HAR (= OLS), BarPre
\providecommand{\rhcaLinFullMzlBBarPreMin}{}\renewcommand{\rhcaLinFullMzlBBarPreMin}{0.99}  % logs MZ slope, linear full, BarPre: min
\providecommand{\rhcaLinFullMzlBBarPreMax}{}\renewcommand{\rhcaLinFullMzlBBarPreMax}{0.99}  % ... max
\providecommand{\rhcaLinFullMzlBBarPreRng}{}\renewcommand{\rhcaLinFullMzlBBarPreRng}{0.99}  % ... the range as text
\providecommand{\rhcaTreeFullMzlBBarPreMin}{}\renewcommand{\rhcaTreeFullMzlBBarPreMin}{1.00}  % logs MZ slope, tree full, BarPre: min
\providecommand{\rhcaTreeFullMzlBBarPreMax}{}\renewcommand{\rhcaTreeFullMzlBBarPreMax}{1.01}  % ... max
\providecommand{\rhcaTreeFullMzlBBarPreRng}{}\renewcommand{\rhcaTreeFullMzlBBarPreRng}{1.00--1.01}  % ... the range as text
\providecommand{\rhcaAllHarMzlBBarPreMin}{}\renewcommand{\rhcaAllHarMzlBBarPreMin}{1.00}  % logs MZ slope, HAR forecasts, BarPre: min
\providecommand{\rhcaAllHarMzlBBarPreMax}{}\renewcommand{\rhcaAllHarMzlBBarPreMax}{1.01}  % ... max
\providecommand{\rhcaAllHarMzlBBarPreRng}{}\renewcommand{\rhcaAllHarMzlBBarPreRng}{1.00--1.01}  % ... the range as text
\providecommand{\rhcaLinFullMzlRsqBarPreMin}{}\renewcommand{\rhcaLinFullMzlRsqBarPreMin}{0.78}  % logs MZ R2, linear full, BarPre: min
\providecommand{\rhcaLinFullMzlRsqBarPreMax}{}\renewcommand{\rhcaLinFullMzlRsqBarPreMax}{0.78}  % ... max
\providecommand{\rhcaLinFullMzlRsqBarPreRng}{}\renewcommand{\rhcaLinFullMzlRsqBarPreRng}{0.78}  % ... the range as text
\providecommand{\rhcaTreeFullMzlRsqBarPreMin}{}\renewcommand{\rhcaTreeFullMzlRsqBarPreMin}{0.79}  % logs MZ R2, tree full, BarPre: min
\providecommand{\rhcaTreeFullMzlRsqBarPreMax}{}\renewcommand{\rhcaTreeFullMzlRsqBarPreMax}{0.79}  % ... max
\providecommand{\rhcaTreeFullMzlRsqBarPreRng}{}\renewcommand{\rhcaTreeFullMzlRsqBarPreRng}{0.79}  % ... the range as text
\providecommand{\rhcaAllHarMzlRsqBarPreMin}{}\renewcommand{\rhcaAllHarMzlRsqBarPreMin}{0.76}  % logs MZ R2, HAR forecasts, BarPre: min
\providecommand{\rhcaAllHarMzlRsqBarPreMax}{}\renewcommand{\rhcaAllHarMzlRsqBarPreMax}{0.76}  % ... max
\providecommand{\rhcaAllHarMzlRsqBarPreRng}{}\renewcommand{\rhcaAllHarMzlRsqBarPreRng}{0.76}  % ... the range as text
\providecommand{\rhcaLinFullMzRsqBarPreMin}{}\renewcommand{\rhcaLinFullMzRsqBarPreMin}{0.53}  % levels MZ R2, linear full, BarPre: min
\providecommand{\rhcaLinFullMzRsqBarPreMax}{}\renewcommand{\rhcaLinFullMzRsqBarPreMax}{0.54}  % ... max
\providecommand{\rhcaLinFullMzRsqBarPreRng}{}\renewcommand{\rhcaLinFullMzRsqBarPreRng}{0.53--0.54}  % ... the range as text
\providecommand{\rhcaAllHarMzRsqBarPreMin}{}\renewcommand{\rhcaAllHarMzRsqBarPreMin}{0.37}  % levels MZ R2, HAR forecasts, BarPre: min
\providecommand{\rhcaAllHarMzRsqBarPreMax}{}\renewcommand{\rhcaAllHarMzRsqBarPreMax}{0.48}  % ... max
\providecommand{\rhcaAllHarMzRsqBarPreRng}{}\renewcommand{\rhcaAllHarMzRsqBarPreRng}{0.37--0.48}  % ... the range as text
\providecommand{\rhcaLinFullShSlopeBarPreMin}{}\renewcommand{\rhcaLinFullShSlopeBarPreMin}{0.0}  % Theil slope share %, linear full, BarPre: min
\providecommand{\rhcaLinFullShSlopeBarPreMax}{}\renewcommand{\rhcaLinFullShSlopeBarPreMax}{0.0}  % ... max
\providecommand{\rhcaLinFullShSlopeBarPreRng}{}\renewcommand{\rhcaLinFullShSlopeBarPreRng}{0.0}  % ... the range as text
\providecommand{\rhcaTreeFullShSlopeBarPreMin}{}\renewcommand{\rhcaTreeFullShSlopeBarPreMin}{0.1}  % Theil slope share %, tree full, BarPre: min
\providecommand{\rhcaTreeFullShSlopeBarPreMax}{}\renewcommand{\rhcaTreeFullShSlopeBarPreMax}{0.5}  % ... max
\providecommand{\rhcaTreeFullShSlopeBarPreRng}{}\renewcommand{\rhcaTreeFullShSlopeBarPreRng}{0.1--0.5}  % ... the range as text
\providecommand{\rhcaAllHarShSlopeBarPreMin}{}\renewcommand{\rhcaAllHarShSlopeBarPreMin}{0.0}  % Theil slope share %, HAR forecasts, BarPre: min
\providecommand{\rhcaAllHarShSlopeBarPreMax}{}\renewcommand{\rhcaAllHarShSlopeBarPreMax}{0.1}  % ... max
\providecommand{\rhcaAllHarShSlopeBarPreRng}{}\renewcommand{\rhcaAllHarShSlopeBarPreRng}{0.0--0.1}  % ... the range as text
\providecommand{\rhcaAllShBiasBarPreMin}{}\renewcommand{\rhcaAllShBiasBarPreMin}{0.1}  % Theil bias share %, all 12, BarPre: min
\providecommand{\rhcaAllShBiasBarPreMax}{}\renewcommand{\rhcaAllShBiasBarPreMax}{0.1}  % ... max
\providecommand{\rhcaAllShBiasBarPreRng}{}\renewcommand{\rhcaAllShBiasBarPreRng}{0.1}  % ... the range as text
\providecommand{\rhcaAllRmBarPreMin}{}\renewcommand{\rhcaAllRmBarPreMin}{1.04}  % ratio of means, all 12, BarPre: min
\providecommand{\rhcaAllRmBarPreMax}{}\renewcommand{\rhcaAllRmBarPreMax}{1.07}  % ... max
\providecommand{\rhcaAllRmBarPreRng}{}\renewcommand{\rhcaAllRmBarPreRng}{1.04--1.07}  % ... the range as text
\providecommand{\rhcaFullRmBarPreMin}{}\renewcommand{\rhcaFullRmBarPreMin}{1.04}  % ratio of means, full set, BarPre: min
\providecommand{\rhcaFullRmBarPreMax}{}\renewcommand{\rhcaFullRmBarPreMax}{1.06}  % ... max
\providecommand{\rhcaFullRmBarPreRng}{}\renewcommand{\rhcaFullRmBarPreRng}{1.04--1.06}  % ... the range as text
\providecommand{\rhcaHarRmBarPreMin}{}\renewcommand{\rhcaHarRmBarPreMin}{1.06}  % ratio of means, HAR, BarPre: min
\providecommand{\rhcaHarRmBarPreMax}{}\renewcommand{\rhcaHarRmBarPreMax}{1.07}  % ... max
\providecommand{\rhcaHarRmBarPreRng}{}\renewcommand{\rhcaHarRmBarPreRng}{1.06--1.07}  % ... the range as text
\providecommand{\rhcaAllMrBarPreMin}{}\renewcommand{\rhcaAllMrBarPreMin}{1.07}  % mean of RV/sigma2, all 12, BarPre: min
\providecommand{\rhcaAllMrBarPreMax}{}\renewcommand{\rhcaAllMrBarPreMax}{1.10}  % ... max
\providecommand{\rhcaAllMrBarPreRng}{}\renewcommand{\rhcaAllMrBarPreRng}{1.07--1.10}  % ... the range as text
\providecommand{\rhcaAllUnderBarPreMin}{}\renewcommand{\rhcaAllUnderBarPreMin}{39.9}  % under-prediction %, all 12, BarPre: min
\providecommand{\rhcaAllUnderBarPreMax}{}\renewcommand{\rhcaAllUnderBarPreMax}{42.1}  % ... max
\providecommand{\rhcaAllUnderBarPreRng}{}\renewcommand{\rhcaAllUnderBarPreRng}{39.9--42.1}  % ... the range as text
\providecommand{\rhcaLinFullUnderBarPreMin}{}\renewcommand{\rhcaLinFullUnderBarPreMin}{41.5}  % under-prediction %, linear full, BarPre: min
\providecommand{\rhcaLinFullUnderBarPreMax}{}\renewcommand{\rhcaLinFullUnderBarPreMax}{42.1}  % ... max
\providecommand{\rhcaLinFullUnderBarPreRng}{}\renewcommand{\rhcaLinFullUnderBarPreRng}{41.5--42.1}  % ... the range as text
\providecommand{\rhcaLinHarUnderBarPre}{}\renewcommand{\rhcaLinHarUnderBarPre}{39.9}  % under-prediction %, linear HAR, BarPre
\providecommand{\rhcaLinFullOverSeBarPreMin}{}\renewcommand{\rhcaLinFullOverSeBarPreMin}{19.6}  % % of SE from over-predictions, linear full, BarPre: min
\providecommand{\rhcaLinFullOverSeBarPreMax}{}\renewcommand{\rhcaLinFullOverSeBarPreMax}{20.1}  % ... max
\providecommand{\rhcaLinFullOverSeBarPreRng}{}\renewcommand{\rhcaLinFullOverSeBarPreRng}{19.6--20.1}  % ... the range as text
\providecommand{\rhcaTreeFullOverSeBarPreMin}{}\renewcommand{\rhcaTreeFullOverSeBarPreMin}{12.0}  % % of SE from over-predictions, tree full, BarPre: min
\providecommand{\rhcaTreeFullOverSeBarPreMax}{}\renewcommand{\rhcaTreeFullOverSeBarPreMax}{15.1}  % ... max
\providecommand{\rhcaTreeFullOverSeBarPreRng}{}\renewcommand{\rhcaTreeFullOverSeBarPreRng}{12.0--15.1}  % ... the range as text
\providecommand{\rhcaAllHarOverSeBarPreMin}{}\renewcommand{\rhcaAllHarOverSeBarPreMin}{10.3}  % % of SE from over-predictions, HAR, BarPre: min
\providecommand{\rhcaAllHarOverSeBarPreMax}{}\renewcommand{\rhcaAllHarOverSeBarPreMax}{20.0}  % ... max
\providecommand{\rhcaAllHarOverSeBarPreRng}{}\renewcommand{\rhcaAllHarOverSeBarPreRng}{10.3--20.0}  % ... the range as text
\providecommand{\rhcaNMaeLowerBarPre}{}\renewcommand{\rhcaNMaeLowerBarPre}{6}  % classes (of 6) with the lower MAE on the full set, BarPre
\providecommand{\rhcaLinFullMrBarPreMin}{}\renewcommand{\rhcaLinFullMrBarPreMin}{1.091}  % mean of RV/sigma2, linear full, BarPre: min
\providecommand{\rhcaLinFullMrBarPreMax}{}\renewcommand{\rhcaLinFullMrBarPreMax}{1.096}  % ... max
\providecommand{\rhcaLinFullMrBarPreRng}{}\renewcommand{\rhcaLinFullMrBarPreRng}{1.091--1.096}  % ... the range as text
\providecommand{\rhcaLinHarMrBarPre}{}\renewcommand{\rhcaLinHarMrBarPre}{1.078}  % mean of RV/sigma2, linear HAR, BarPre
\providecommand{\rhcaLinFullQlBarPreMin}{}\renewcommand{\rhcaLinFullQlBarPreMin}{0.0039}  % QLIKE level term, linear full, BarPre: min
\providecommand{\rhcaLinFullQlBarPreMax}{}\renewcommand{\rhcaLinFullQlBarPreMax}{0.0043}  % ... max
\providecommand{\rhcaLinFullQlBarPreRng}{}\renewcommand{\rhcaLinFullQlBarPreRng}{0.0039--0.0043}  % ... the range as text
\providecommand{\rhcaLinHarQlBarPre}{}\renewcommand{\rhcaLinHarQlBarPre}{0.0029}  % QLIKE level term, linear HAR, BarPre
\providecommand{\rhcaLinFullQdBarPreMin}{}\renewcommand{\rhcaLinFullQdBarPreMin}{0.1627}  % QLIKE dispersion term, linear full, BarPre: min
\providecommand{\rhcaLinFullQdBarPreMax}{}\renewcommand{\rhcaLinFullQdBarPreMax}{0.1683}  % ... max
\providecommand{\rhcaLinFullQdBarPreRng}{}\renewcommand{\rhcaLinFullQdBarPreRng}{0.1627--0.1683}  % ... the range as text
\providecommand{\rhcaLinHarQdBarPre}{}\renewcommand{\rhcaLinHarQdBarPre}{0.1835}  % QLIKE dispersion term, linear HAR, BarPre
\providecommand{\rhcaTreeFullQlBarPreMin}{}\renewcommand{\rhcaTreeFullQlBarPreMin}{0.0021}  % QLIKE level term, tree full, BarPre: min
\providecommand{\rhcaTreeFullQlBarPreMax}{}\renewcommand{\rhcaTreeFullQlBarPreMax}{0.0029}  % ... max
\providecommand{\rhcaTreeFullQlBarPreRng}{}\renewcommand{\rhcaTreeFullQlBarPreRng}{0.0021--0.0029}  % ... the range as text
\providecommand{\rhcaTreeHarQlBarPreMin}{}\renewcommand{\rhcaTreeHarQlBarPreMin}{0.0032}  % QLIKE level term, tree HAR, BarPre: min
\providecommand{\rhcaTreeHarQlBarPreMax}{}\renewcommand{\rhcaTreeHarQlBarPreMax}{0.0034}  % ... max
\providecommand{\rhcaTreeHarQlBarPreRng}{}\renewcommand{\rhcaTreeHarQlBarPreRng}{0.0032--0.0034}  % ... the range as text
\providecommand{\rhcaTreeFullQdBarPreMin}{}\renewcommand{\rhcaTreeFullQdBarPreMin}{0.1563}  % QLIKE dispersion term, tree full, BarPre: min
\providecommand{\rhcaTreeFullQdBarPreMax}{}\renewcommand{\rhcaTreeFullQdBarPreMax}{0.1617}  % ... max
\providecommand{\rhcaTreeFullQdBarPreRng}{}\renewcommand{\rhcaTreeFullQdBarPreRng}{0.1563--0.1617}  % ... the range as text
\providecommand{\rhcaTreeHarQdBarPreMin}{}\renewcommand{\rhcaTreeHarQdBarPreMin}{0.1826}  % QLIKE dispersion term, tree HAR, BarPre: min
\providecommand{\rhcaTreeHarQdBarPreMax}{}\renewcommand{\rhcaTreeHarQdBarPreMax}{0.1832}  % ... max
\providecommand{\rhcaTreeHarQdBarPreRng}{}\renewcommand{\rhcaTreeHarQdBarPreRng}{0.1826--0.1832}  % ... the range as text
\providecommand{\rhcaTreeFullUnderBarPreMin}{}\renewcommand{\rhcaTreeFullUnderBarPreMin}{40.8}  % under-prediction %, tree full, BarPre: min
\providecommand{\rhcaTreeFullUnderBarPreMax}{}\renewcommand{\rhcaTreeFullUnderBarPreMax}{41.1}  % ... max
\providecommand{\rhcaTreeFullUnderBarPreRng}{}\renewcommand{\rhcaTreeFullUnderBarPreRng}{40.8--41.1}  % ... the range as text
\providecommand{\rhcaTreeHarUnderBarPreMin}{}\renewcommand{\rhcaTreeHarUnderBarPreMin}{40.1}  % under-prediction %, tree HAR, BarPre: min
\providecommand{\rhcaTreeHarUnderBarPreMax}{}\renewcommand{\rhcaTreeHarUnderBarPreMax}{40.2}  % ... max
\providecommand{\rhcaTreeHarUnderBarPreRng}{}\renewcommand{\rhcaTreeHarUnderBarPreRng}{40.1--40.2}  % ... the range as text
\providecommand{\rhcaLinFullMzBBarOosMin}{}\renewcommand{\rhcaLinFullMzBBarOosMin}{0.56}  % levels MZ slope, linear full, BarOos: min
\providecommand{\rhcaLinFullMzBBarOosMax}{}\renewcommand{\rhcaLinFullMzBBarOosMax}{0.58}  % ... max
\providecommand{\rhcaLinFullMzBBarOosRng}{}\renewcommand{\rhcaLinFullMzBBarOosRng}{0.56--0.58}  % ... the range as text
\providecommand{\rhcaTreeFullMzBBarOosMin}{}\renewcommand{\rhcaTreeFullMzBBarOosMin}{0.74}  % levels MZ slope, tree full, BarOos: min
\providecommand{\rhcaTreeFullMzBBarOosMax}{}\renewcommand{\rhcaTreeFullMzBBarOosMax}{0.78}  % ... max
\providecommand{\rhcaTreeFullMzBBarOosRng}{}\renewcommand{\rhcaTreeFullMzBBarOosRng}{0.74--0.78}  % ... the range as text
\providecommand{\rhcaTreeHarMzBBarOosMin}{}\renewcommand{\rhcaTreeHarMzBBarOosMin}{0.92}  % levels MZ slope, tree HAR, BarOos: min
\providecommand{\rhcaTreeHarMzBBarOosMax}{}\renewcommand{\rhcaTreeHarMzBBarOosMax}{0.94}  % ... max
\providecommand{\rhcaTreeHarMzBBarOosRng}{}\renewcommand{\rhcaTreeHarMzBBarOosRng}{0.92--0.94}  % ... the range as text
\providecommand{\rhcaLinHarMzBBarOos}{}\renewcommand{\rhcaLinHarMzBBarOos}{0.99}  % levels MZ slope, linear HAR (= OLS), BarOos
\providecommand{\rhcaLinFullMzlBBarOosMin}{}\renewcommand{\rhcaLinFullMzlBBarOosMin}{0.98}  % logs MZ slope, linear full, BarOos: min
\providecommand{\rhcaLinFullMzlBBarOosMax}{}\renewcommand{\rhcaLinFullMzlBBarOosMax}{0.99}  % ... max
\providecommand{\rhcaLinFullMzlBBarOosRng}{}\renewcommand{\rhcaLinFullMzlBBarOosRng}{0.98--0.99}  % ... the range as text
\providecommand{\rhcaTreeFullMzlBBarOosMin}{}\renewcommand{\rhcaTreeFullMzlBBarOosMin}{0.99}  % logs MZ slope, tree full, BarOos: min
\providecommand{\rhcaTreeFullMzlBBarOosMax}{}\renewcommand{\rhcaTreeFullMzlBBarOosMax}{0.99}  % ... max
\providecommand{\rhcaTreeFullMzlBBarOosRng}{}\renewcommand{\rhcaTreeFullMzlBBarOosRng}{0.99}  % ... the range as text
\providecommand{\rhcaAllHarMzlBBarOosMin}{}\renewcommand{\rhcaAllHarMzlBBarOosMin}{1.00}  % logs MZ slope, HAR forecasts, BarOos: min
\providecommand{\rhcaAllHarMzlBBarOosMax}{}\renewcommand{\rhcaAllHarMzlBBarOosMax}{1.02}  % ... max
\providecommand{\rhcaAllHarMzlBBarOosRng}{}\renewcommand{\rhcaAllHarMzlBBarOosRng}{1.00--1.02}  % ... the range as text
\providecommand{\rhcaLinFullMzlRsqBarOosMin}{}\renewcommand{\rhcaLinFullMzlRsqBarOosMin}{0.81}  % logs MZ R2, linear full, BarOos: min
\providecommand{\rhcaLinFullMzlRsqBarOosMax}{}\renewcommand{\rhcaLinFullMzlRsqBarOosMax}{0.81}  % ... max
\providecommand{\rhcaLinFullMzlRsqBarOosRng}{}\renewcommand{\rhcaLinFullMzlRsqBarOosRng}{0.81}  % ... the range as text
\providecommand{\rhcaTreeFullMzlRsqBarOosMin}{}\renewcommand{\rhcaTreeFullMzlRsqBarOosMin}{0.82}  % logs MZ R2, tree full, BarOos: min
\providecommand{\rhcaTreeFullMzlRsqBarOosMax}{}\renewcommand{\rhcaTreeFullMzlRsqBarOosMax}{0.82}  % ... max
\providecommand{\rhcaTreeFullMzlRsqBarOosRng}{}\renewcommand{\rhcaTreeFullMzlRsqBarOosRng}{0.82}  % ... the range as text
\providecommand{\rhcaAllHarMzlRsqBarOosMin}{}\renewcommand{\rhcaAllHarMzlRsqBarOosMin}{0.79}  % logs MZ R2, HAR forecasts, BarOos: min
\providecommand{\rhcaAllHarMzlRsqBarOosMax}{}\renewcommand{\rhcaAllHarMzlRsqBarOosMax}{0.80}  % ... max
\providecommand{\rhcaAllHarMzlRsqBarOosRng}{}\renewcommand{\rhcaAllHarMzlRsqBarOosRng}{0.79--0.80}  % ... the range as text
\providecommand{\rhcaLinFullMzRsqBarOosMin}{}\renewcommand{\rhcaLinFullMzRsqBarOosMin}{0.51}  % levels MZ R2, linear full, BarOos: min
\providecommand{\rhcaLinFullMzRsqBarOosMax}{}\renewcommand{\rhcaLinFullMzRsqBarOosMax}{0.52}  % ... max
\providecommand{\rhcaLinFullMzRsqBarOosRng}{}\renewcommand{\rhcaLinFullMzRsqBarOosRng}{0.51--0.52}  % ... the range as text
\providecommand{\rhcaAllHarMzRsqBarOosMin}{}\renewcommand{\rhcaAllHarMzRsqBarOosMin}{0.65}  % levels MZ R2, HAR forecasts, BarOos: min
\providecommand{\rhcaAllHarMzRsqBarOosMax}{}\renewcommand{\rhcaAllHarMzRsqBarOosMax}{0.67}  % ... max
\providecommand{\rhcaAllHarMzRsqBarOosRng}{}\renewcommand{\rhcaAllHarMzRsqBarOosRng}{0.65--0.67}  % ... the range as text
\providecommand{\rhcaLinFullShSlopeBarOosMin}{}\renewcommand{\rhcaLinFullShSlopeBarOosMin}{35.9}  % Theil slope share %, linear full, BarOos: min
\providecommand{\rhcaLinFullShSlopeBarOosMax}{}\renewcommand{\rhcaLinFullShSlopeBarOosMax}{40.8}  % ... max
\providecommand{\rhcaLinFullShSlopeBarOosRng}{}\renewcommand{\rhcaLinFullShSlopeBarOosRng}{35.9--40.8}  % ... the range as text
\providecommand{\rhcaTreeFullShSlopeBarOosMin}{}\renewcommand{\rhcaTreeFullShSlopeBarOosMin}{12.7}  % Theil slope share %, tree full, BarOos: min
\providecommand{\rhcaTreeFullShSlopeBarOosMax}{}\renewcommand{\rhcaTreeFullShSlopeBarOosMax}{17.7}  % ... max
\providecommand{\rhcaTreeFullShSlopeBarOosRng}{}\renewcommand{\rhcaTreeFullShSlopeBarOosRng}{12.7--17.7}  % ... the range as text
\providecommand{\rhcaAllHarShSlopeBarOosMin}{}\renewcommand{\rhcaAllHarShSlopeBarOosMin}{0.0}  % Theil slope share %, HAR forecasts, BarOos: min
\providecommand{\rhcaAllHarShSlopeBarOosMax}{}\renewcommand{\rhcaAllHarShSlopeBarOosMax}{1.4}  % ... max
\providecommand{\rhcaAllHarShSlopeBarOosRng}{}\renewcommand{\rhcaAllHarShSlopeBarOosRng}{0.0--1.4}  % ... the range as text
\providecommand{\rhcaAllShBiasBarOosMin}{}\renewcommand{\rhcaAllShBiasBarOosMin}{0.0}  % Theil bias share %, all 12, BarOos: min
\providecommand{\rhcaAllShBiasBarOosMax}{}\renewcommand{\rhcaAllShBiasBarOosMax}{0.1}  % ... max
\providecommand{\rhcaAllShBiasBarOosRng}{}\renewcommand{\rhcaAllShBiasBarOosRng}{0.0--0.1}  % ... the range as text
\providecommand{\rhcaAllRmBarOosMin}{}\renewcommand{\rhcaAllRmBarOosMin}{0.96}  % ratio of means, all 12, BarOos: min
\providecommand{\rhcaAllRmBarOosMax}{}\renewcommand{\rhcaAllRmBarOosMax}{1.05}  % ... max
\providecommand{\rhcaAllRmBarOosRng}{}\renewcommand{\rhcaAllRmBarOosRng}{0.96--1.05}  % ... the range as text
\providecommand{\rhcaFullRmBarOosMin}{}\renewcommand{\rhcaFullRmBarOosMin}{0.96}  % ratio of means, full set, BarOos: min
\providecommand{\rhcaFullRmBarOosMax}{}\renewcommand{\rhcaFullRmBarOosMax}{0.98}  % ... max
\providecommand{\rhcaFullRmBarOosRng}{}\renewcommand{\rhcaFullRmBarOosRng}{0.96--0.98}  % ... the range as text
\providecommand{\rhcaHarRmBarOosMin}{}\renewcommand{\rhcaHarRmBarOosMin}{1.02}  % ratio of means, HAR, BarOos: min
\providecommand{\rhcaHarRmBarOosMax}{}\renewcommand{\rhcaHarRmBarOosMax}{1.05}  % ... max
\providecommand{\rhcaHarRmBarOosRng}{}\renewcommand{\rhcaHarRmBarOosRng}{1.02--1.05}  % ... the range as text
\providecommand{\rhcaAllMrBarOosMin}{}\renewcommand{\rhcaAllMrBarOosMin}{1.06}  % mean of RV/sigma2, all 12, BarOos: min
\providecommand{\rhcaAllMrBarOosMax}{}\renewcommand{\rhcaAllMrBarOosMax}{1.12}  % ... max
\providecommand{\rhcaAllMrBarOosRng}{}\renewcommand{\rhcaAllMrBarOosRng}{1.06--1.12}  % ... the range as text
\providecommand{\rhcaAllUnderBarOosMin}{}\renewcommand{\rhcaAllUnderBarOosMin}{38.0}  % under-prediction %, all 12, BarOos: min
\providecommand{\rhcaAllUnderBarOosMax}{}\renewcommand{\rhcaAllUnderBarOosMax}{43.2}  % ... max
\providecommand{\rhcaAllUnderBarOosRng}{}\renewcommand{\rhcaAllUnderBarOosRng}{38.0--43.2}  % ... the range as text
\providecommand{\rhcaLinFullUnderBarOosMin}{}\renewcommand{\rhcaLinFullUnderBarOosMin}{42.2}  % under-prediction %, linear full, BarOos: min
\providecommand{\rhcaLinFullUnderBarOosMax}{}\renewcommand{\rhcaLinFullUnderBarOosMax}{43.2}  % ... max
\providecommand{\rhcaLinFullUnderBarOosRng}{}\renewcommand{\rhcaLinFullUnderBarOosRng}{42.2--43.2}  % ... the range as text
\providecommand{\rhcaLinHarUnderBarOos}{}\renewcommand{\rhcaLinHarUnderBarOos}{38.0}  % under-prediction %, linear HAR, BarOos
\providecommand{\rhcaLinFullOverSeBarOosMin}{}\renewcommand{\rhcaLinFullOverSeBarOosMin}{63.1}  % % of SE from over-predictions, linear full, BarOos: min
\providecommand{\rhcaLinFullOverSeBarOosMax}{}\renewcommand{\rhcaLinFullOverSeBarOosMax}{67.0}  % ... max
\providecommand{\rhcaLinFullOverSeBarOosRng}{}\renewcommand{\rhcaLinFullOverSeBarOosRng}{63.1--67.0}  % ... the range as text
\providecommand{\rhcaTreeFullOverSeBarOosMin}{}\renewcommand{\rhcaTreeFullOverSeBarOosMin}{44.4}  % % of SE from over-predictions, tree full, BarOos: min
\providecommand{\rhcaTreeFullOverSeBarOosMax}{}\renewcommand{\rhcaTreeFullOverSeBarOosMax}{49.1}  % ... max
\providecommand{\rhcaTreeFullOverSeBarOosRng}{}\renewcommand{\rhcaTreeFullOverSeBarOosRng}{44.4--49.1}  % ... the range as text
\providecommand{\rhcaAllHarOverSeBarOosMin}{}\renewcommand{\rhcaAllHarOverSeBarOosMin}{21.5}  % % of SE from over-predictions, HAR, BarOos: min
\providecommand{\rhcaAllHarOverSeBarOosMax}{}\renewcommand{\rhcaAllHarOverSeBarOosMax}{26.7}  % ... max
\providecommand{\rhcaAllHarOverSeBarOosRng}{}\renewcommand{\rhcaAllHarOverSeBarOosRng}{21.5--26.7}  % ... the range as text
\providecommand{\rhcaNMaeLowerBarOos}{}\renewcommand{\rhcaNMaeLowerBarOos}{0}  % classes (of 6) with the lower MAE on the full set, BarOos
\providecommand{\rhcaLinFullMrBarOosMin}{}\renewcommand{\rhcaLinFullMrBarOosMin}{1.109}  % mean of RV/sigma2, linear full, BarOos: min
\providecommand{\rhcaLinFullMrBarOosMax}{}\renewcommand{\rhcaLinFullMrBarOosMax}{1.117}  % ... max
\providecommand{\rhcaLinFullMrBarOosRng}{}\renewcommand{\rhcaLinFullMrBarOosRng}{1.109--1.117}  % ... the range as text
\providecommand{\rhcaLinHarMrBarOos}{}\renewcommand{\rhcaLinHarMrBarOos}{1.061}  % mean of RV/sigma2, linear HAR, BarOos
\providecommand{\rhcaLinFullQlBarOosMin}{}\renewcommand{\rhcaLinFullQlBarOosMin}{0.0056}  % QLIKE level term, linear full, BarOos: min
\providecommand{\rhcaLinFullQlBarOosMax}{}\renewcommand{\rhcaLinFullQlBarOosMax}{0.0064}  % ... max
\providecommand{\rhcaLinFullQlBarOosRng}{}\renewcommand{\rhcaLinFullQlBarOosRng}{0.0056--0.0064}  % ... the range as text
\providecommand{\rhcaLinHarQlBarOos}{}\renewcommand{\rhcaLinHarQlBarOos}{0.0018}  % QLIKE level term, linear HAR, BarOos
\providecommand{\rhcaLinFullQdBarOosMin}{}\renewcommand{\rhcaLinFullQdBarOosMin}{0.1739}  % QLIKE dispersion term, linear full, BarOos: min
\providecommand{\rhcaLinFullQdBarOosMax}{}\renewcommand{\rhcaLinFullQdBarOosMax}{0.1788}  % ... max
\providecommand{\rhcaLinFullQdBarOosRng}{}\renewcommand{\rhcaLinFullQdBarOosRng}{0.1739--0.1788}  % ... the range as text
\providecommand{\rhcaLinHarQdBarOos}{}\renewcommand{\rhcaLinHarQdBarOos}{0.1915}  % QLIKE dispersion term, linear HAR, BarOos
\providecommand{\rhcaTreeFullQlBarOosMin}{}\renewcommand{\rhcaTreeFullQlBarOosMin}{0.0022}  % QLIKE level term, tree full, BarOos: min
\providecommand{\rhcaTreeFullQlBarOosMax}{}\renewcommand{\rhcaTreeFullQlBarOosMax}{0.0035}  % ... max
\providecommand{\rhcaTreeFullQlBarOosRng}{}\renewcommand{\rhcaTreeFullQlBarOosRng}{0.0022--0.0035}  % ... the range as text
\providecommand{\rhcaTreeHarQlBarOosMin}{}\renewcommand{\rhcaTreeHarQlBarOosMin}{0.0022}  % QLIKE level term, tree HAR, BarOos: min
\providecommand{\rhcaTreeHarQlBarOosMax}{}\renewcommand{\rhcaTreeHarQlBarOosMax}{0.0028}  % ... max
\providecommand{\rhcaTreeHarQlBarOosRng}{}\renewcommand{\rhcaTreeHarQlBarOosRng}{0.0022--0.0028}  % ... the range as text
\providecommand{\rhcaTreeFullQdBarOosMin}{}\renewcommand{\rhcaTreeFullQdBarOosMin}{0.1612}  % QLIKE dispersion term, tree full, BarOos: min
\providecommand{\rhcaTreeFullQdBarOosMax}{}\renewcommand{\rhcaTreeFullQdBarOosMax}{0.1694}  % ... max
\providecommand{\rhcaTreeFullQdBarOosRng}{}\renewcommand{\rhcaTreeFullQdBarOosRng}{0.1612--0.1694}  % ... the range as text
\providecommand{\rhcaTreeHarQdBarOosMin}{}\renewcommand{\rhcaTreeHarQdBarOosMin}{0.1878}  % QLIKE dispersion term, tree HAR, BarOos: min
\providecommand{\rhcaTreeHarQdBarOosMax}{}\renewcommand{\rhcaTreeHarQdBarOosMax}{0.1909}  % ... max
\providecommand{\rhcaTreeHarQdBarOosRng}{}\renewcommand{\rhcaTreeHarQdBarOosRng}{0.1878--0.1909}  % ... the range as text
\providecommand{\rhcaTreeFullUnderBarOosMin}{}\renewcommand{\rhcaTreeFullUnderBarOosMin}{40.5}  % under-prediction %, tree full, BarOos: min
\providecommand{\rhcaTreeFullUnderBarOosMax}{}\renewcommand{\rhcaTreeFullUnderBarOosMax}{41.0}  % ... max
\providecommand{\rhcaTreeFullUnderBarOosRng}{}\renewcommand{\rhcaTreeFullUnderBarOosRng}{40.5--41.0}  % ... the range as text
\providecommand{\rhcaTreeHarUnderBarOosMin}{}\renewcommand{\rhcaTreeHarUnderBarOosMin}{38.4}  % under-prediction %, tree HAR, BarOos: min
\providecommand{\rhcaTreeHarUnderBarOosMax}{}\renewcommand{\rhcaTreeHarUnderBarOosMax}{38.8}  % ... max
\providecommand{\rhcaTreeHarUnderBarOosRng}{}\renewcommand{\rhcaTreeHarUnderBarOosRng}{38.4--38.8}  % ... the range as text
\providecommand{\rhcaLinFullMzBClosePreMin}{}\renewcommand{\rhcaLinFullMzBClosePreMin}{1.06}  % levels MZ slope, linear full, ClosePre: min
\providecommand{\rhcaLinFullMzBClosePreMax}{}\renewcommand{\rhcaLinFullMzBClosePreMax}{1.08}  % ... max
\providecommand{\rhcaLinFullMzBClosePreRng}{}\renewcommand{\rhcaLinFullMzBClosePreRng}{1.06--1.08}  % ... the range as text
\providecommand{\rhcaTreeFullMzBClosePreMin}{}\renewcommand{\rhcaTreeFullMzBClosePreMin}{0.98}  % levels MZ slope, tree full, ClosePre: min
\providecommand{\rhcaTreeFullMzBClosePreMax}{}\renewcommand{\rhcaTreeFullMzBClosePreMax}{0.99}  % ... max
\providecommand{\rhcaTreeFullMzBClosePreRng}{}\renewcommand{\rhcaTreeFullMzBClosePreRng}{0.98--0.99}  % ... the range as text
\providecommand{\rhcaTreeHarMzBClosePreMin}{}\renewcommand{\rhcaTreeHarMzBClosePreMin}{1.06}  % levels MZ slope, tree HAR, ClosePre: min
\providecommand{\rhcaTreeHarMzBClosePreMax}{}\renewcommand{\rhcaTreeHarMzBClosePreMax}{1.11}  % ... max
\providecommand{\rhcaTreeHarMzBClosePreRng}{}\renewcommand{\rhcaTreeHarMzBClosePreRng}{1.06--1.11}  % ... the range as text
\providecommand{\rhcaLinHarMzBClosePre}{}\renewcommand{\rhcaLinHarMzBClosePre}{1.14}  % levels MZ slope, linear HAR (= OLS), ClosePre
\providecommand{\rhcaLinFullMzlBClosePreMin}{}\renewcommand{\rhcaLinFullMzlBClosePreMin}{0.92}  % logs MZ slope, linear full, ClosePre: min
\providecommand{\rhcaLinFullMzlBClosePreMax}{}\renewcommand{\rhcaLinFullMzlBClosePreMax}{0.93}  % ... max
\providecommand{\rhcaLinFullMzlBClosePreRng}{}\renewcommand{\rhcaLinFullMzlBClosePreRng}{0.92--0.93}  % ... the range as text
\providecommand{\rhcaTreeFullMzlBClosePreMin}{}\renewcommand{\rhcaTreeFullMzlBClosePreMin}{0.92}  % logs MZ slope, tree full, ClosePre: min
\providecommand{\rhcaTreeFullMzlBClosePreMax}{}\renewcommand{\rhcaTreeFullMzlBClosePreMax}{0.94}  % ... max
\providecommand{\rhcaTreeFullMzlBClosePreRng}{}\renewcommand{\rhcaTreeFullMzlBClosePreRng}{0.92--0.94}  % ... the range as text
\providecommand{\rhcaAllHarMzlBClosePreMin}{}\renewcommand{\rhcaAllHarMzlBClosePreMin}{0.95}  % logs MZ slope, HAR forecasts, ClosePre: min
\providecommand{\rhcaAllHarMzlBClosePreMax}{}\renewcommand{\rhcaAllHarMzlBClosePreMax}{0.96}  % ... max
\providecommand{\rhcaAllHarMzlBClosePreRng}{}\renewcommand{\rhcaAllHarMzlBClosePreRng}{0.95--0.96}  % ... the range as text
\providecommand{\rhcaLinFullMzlRsqClosePreMin}{}\renewcommand{\rhcaLinFullMzlRsqClosePreMin}{0.80}  % logs MZ R2, linear full, ClosePre: min
\providecommand{\rhcaLinFullMzlRsqClosePreMax}{}\renewcommand{\rhcaLinFullMzlRsqClosePreMax}{0.80}  % ... max
\providecommand{\rhcaLinFullMzlRsqClosePreRng}{}\renewcommand{\rhcaLinFullMzlRsqClosePreRng}{0.80}  % ... the range as text
\providecommand{\rhcaTreeFullMzlRsqClosePreMin}{}\renewcommand{\rhcaTreeFullMzlRsqClosePreMin}{0.79}  % logs MZ R2, tree full, ClosePre: min
\providecommand{\rhcaTreeFullMzlRsqClosePreMax}{}\renewcommand{\rhcaTreeFullMzlRsqClosePreMax}{0.80}  % ... max
\providecommand{\rhcaTreeFullMzlRsqClosePreRng}{}\renewcommand{\rhcaTreeFullMzlRsqClosePreRng}{0.79--0.80}  % ... the range as text
\providecommand{\rhcaAllHarMzlRsqClosePreMin}{}\renewcommand{\rhcaAllHarMzlRsqClosePreMin}{0.78}  % logs MZ R2, HAR forecasts, ClosePre: min
\providecommand{\rhcaAllHarMzlRsqClosePreMax}{}\renewcommand{\rhcaAllHarMzlRsqClosePreMax}{0.78}  % ... max
\providecommand{\rhcaAllHarMzlRsqClosePreRng}{}\renewcommand{\rhcaAllHarMzlRsqClosePreRng}{0.78}  % ... the range as text
\providecommand{\rhcaLinFullMzRsqClosePreMin}{}\renewcommand{\rhcaLinFullMzRsqClosePreMin}{0.68}  % levels MZ R2, linear full, ClosePre: min
\providecommand{\rhcaLinFullMzRsqClosePreMax}{}\renewcommand{\rhcaLinFullMzRsqClosePreMax}{0.69}  % ... max
\providecommand{\rhcaLinFullMzRsqClosePreRng}{}\renewcommand{\rhcaLinFullMzRsqClosePreRng}{0.68--0.69}  % ... the range as text
\providecommand{\rhcaAllHarMzRsqClosePreMin}{}\renewcommand{\rhcaAllHarMzRsqClosePreMin}{0.55}  % levels MZ R2, HAR forecasts, ClosePre: min
\providecommand{\rhcaAllHarMzRsqClosePreMax}{}\renewcommand{\rhcaAllHarMzRsqClosePreMax}{0.64}  % ... max
\providecommand{\rhcaAllHarMzRsqClosePreRng}{}\renewcommand{\rhcaAllHarMzRsqClosePreRng}{0.55--0.64}  % ... the range as text
\providecommand{\rhcaLinFullShSlopeClosePreMin}{}\renewcommand{\rhcaLinFullShSlopeClosePreMin}{0.7}  % Theil slope share %, linear full, ClosePre: min
\providecommand{\rhcaLinFullShSlopeClosePreMax}{}\renewcommand{\rhcaLinFullShSlopeClosePreMax}{1.2}  % ... max
\providecommand{\rhcaLinFullShSlopeClosePreRng}{}\renewcommand{\rhcaLinFullShSlopeClosePreRng}{0.7--1.2}  % ... the range as text
\providecommand{\rhcaTreeFullShSlopeClosePreMin}{}\renewcommand{\rhcaTreeFullShSlopeClosePreMin}{0.0}  % Theil slope share %, tree full, ClosePre: min
\providecommand{\rhcaTreeFullShSlopeClosePreMax}{}\renewcommand{\rhcaTreeFullShSlopeClosePreMax}{0.1}  % ... max
\providecommand{\rhcaTreeFullShSlopeClosePreRng}{}\renewcommand{\rhcaTreeFullShSlopeClosePreRng}{0.0--0.1}  % ... the range as text
\providecommand{\rhcaAllHarShSlopeClosePreMin}{}\renewcommand{\rhcaAllHarShSlopeClosePreMin}{0.4}  % Theil slope share %, HAR forecasts, ClosePre: min
\providecommand{\rhcaAllHarShSlopeClosePreMax}{}\renewcommand{\rhcaAllHarShSlopeClosePreMax}{2.5}  % ... max
\providecommand{\rhcaAllHarShSlopeClosePreRng}{}\renewcommand{\rhcaAllHarShSlopeClosePreRng}{0.4--2.5}  % ... the range as text
\providecommand{\rhcaAllShBiasClosePreMin}{}\renewcommand{\rhcaAllShBiasClosePreMin}{0.1}  % Theil bias share %, all 12, ClosePre: min
\providecommand{\rhcaAllShBiasClosePreMax}{}\renewcommand{\rhcaAllShBiasClosePreMax}{0.5}  % ... max
\providecommand{\rhcaAllShBiasClosePreRng}{}\renewcommand{\rhcaAllShBiasClosePreRng}{0.1--0.5}  % ... the range as text
\providecommand{\rhcaAllRmClosePreMin}{}\renewcommand{\rhcaAllRmClosePreMin}{1.04}  % ratio of means, all 12, ClosePre: min
\providecommand{\rhcaAllRmClosePreMax}{}\renewcommand{\rhcaAllRmClosePreMax}{1.10}  % ... max
\providecommand{\rhcaAllRmClosePreRng}{}\renewcommand{\rhcaAllRmClosePreRng}{1.04--1.10}  % ... the range as text
\providecommand{\rhcaFullRmClosePreMin}{}\renewcommand{\rhcaFullRmClosePreMin}{1.04}  % ratio of means, full set, ClosePre: min
\providecommand{\rhcaFullRmClosePreMax}{}\renewcommand{\rhcaFullRmClosePreMax}{1.07}  % ... max
\providecommand{\rhcaFullRmClosePreRng}{}\renewcommand{\rhcaFullRmClosePreRng}{1.04--1.07}  % ... the range as text
\providecommand{\rhcaHarRmClosePreMin}{}\renewcommand{\rhcaHarRmClosePreMin}{1.06}  % ratio of means, HAR, ClosePre: min
\providecommand{\rhcaHarRmClosePreMax}{}\renewcommand{\rhcaHarRmClosePreMax}{1.10}  % ... max
\providecommand{\rhcaHarRmClosePreRng}{}\renewcommand{\rhcaHarRmClosePreRng}{1.06--1.10}  % ... the range as text
\providecommand{\rhcaAllMrClosePreMin}{}\renewcommand{\rhcaAllMrClosePreMin}{1.06}  % mean of RV/sigma2, all 12, ClosePre: min
\providecommand{\rhcaAllMrClosePreMax}{}\renewcommand{\rhcaAllMrClosePreMax}{1.15}  % ... max
\providecommand{\rhcaAllMrClosePreRng}{}\renewcommand{\rhcaAllMrClosePreRng}{1.06--1.15}  % ... the range as text
\providecommand{\rhcaAllUnderClosePreMin}{}\renewcommand{\rhcaAllUnderClosePreMin}{42.6}  % under-prediction %, all 12, ClosePre: min
\providecommand{\rhcaAllUnderClosePreMax}{}\renewcommand{\rhcaAllUnderClosePreMax}{48.6}  % ... max
\providecommand{\rhcaAllUnderClosePreRng}{}\renewcommand{\rhcaAllUnderClosePreRng}{42.6--48.6}  % ... the range as text
\providecommand{\rhcaLinFullUnderClosePreMin}{}\renewcommand{\rhcaLinFullUnderClosePreMin}{47.2}  % under-prediction %, linear full, ClosePre: min
\providecommand{\rhcaLinFullUnderClosePreMax}{}\renewcommand{\rhcaLinFullUnderClosePreMax}{48.6}  % ... max
\providecommand{\rhcaLinFullUnderClosePreRng}{}\renewcommand{\rhcaLinFullUnderClosePreRng}{47.2--48.6}  % ... the range as text
\providecommand{\rhcaLinHarUnderClosePre}{}\renewcommand{\rhcaLinHarUnderClosePre}{46.7}  % under-prediction %, linear HAR, ClosePre
\providecommand{\rhcaLinFullOverSeClosePreMin}{}\renewcommand{\rhcaLinFullOverSeClosePreMin}{17.3}  % % of SE from over-predictions, linear full, ClosePre: min
\providecommand{\rhcaLinFullOverSeClosePreMax}{}\renewcommand{\rhcaLinFullOverSeClosePreMax}{18.4}  % ... max
\providecommand{\rhcaLinFullOverSeClosePreRng}{}\renewcommand{\rhcaLinFullOverSeClosePreRng}{17.3--18.4}  % ... the range as text
\providecommand{\rhcaTreeFullOverSeClosePreMin}{}\renewcommand{\rhcaTreeFullOverSeClosePreMin}{17.3}  % % of SE from over-predictions, tree full, ClosePre: min
\providecommand{\rhcaTreeFullOverSeClosePreMax}{}\renewcommand{\rhcaTreeFullOverSeClosePreMax}{18.9}  % ... max
\providecommand{\rhcaTreeFullOverSeClosePreRng}{}\renewcommand{\rhcaTreeFullOverSeClosePreRng}{17.3--18.9}  % ... the range as text
\providecommand{\rhcaAllHarOverSeClosePreMin}{}\renewcommand{\rhcaAllHarOverSeClosePreMin}{13.1}  % % of SE from over-predictions, HAR, ClosePre: min
\providecommand{\rhcaAllHarOverSeClosePreMax}{}\renewcommand{\rhcaAllHarOverSeClosePreMax}{14.4}  % ... max
\providecommand{\rhcaAllHarOverSeClosePreRng}{}\renewcommand{\rhcaAllHarOverSeClosePreRng}{13.1--14.4}  % ... the range as text
\providecommand{\rhcaNMaeLowerClosePre}{}\renewcommand{\rhcaNMaeLowerClosePre}{6}  % classes (of 6) with the lower MAE on the full set, ClosePre
\providecommand{\rhcaLinFullMrClosePreMin}{}\renewcommand{\rhcaLinFullMrClosePreMin}{1.137}  % mean of RV/sigma2, linear full, ClosePre: min
\providecommand{\rhcaLinFullMrClosePreMax}{}\renewcommand{\rhcaLinFullMrClosePreMax}{1.148}  % ... max
\providecommand{\rhcaLinFullMrClosePreRng}{}\renewcommand{\rhcaLinFullMrClosePreRng}{1.137--1.148}  % ... the range as text
\providecommand{\rhcaLinHarMrClosePre}{}\renewcommand{\rhcaLinHarMrClosePre}{1.121}  % mean of RV/sigma2, linear HAR, ClosePre
\providecommand{\rhcaLinFullQlClosePreMin}{}\renewcommand{\rhcaLinFullQlClosePreMin}{0.0086}  % QLIKE level term, linear full, ClosePre: min
\providecommand{\rhcaLinFullQlClosePreMax}{}\renewcommand{\rhcaLinFullQlClosePreMax}{0.0100}  % ... max
\providecommand{\rhcaLinFullQlClosePreRng}{}\renewcommand{\rhcaLinFullQlClosePreRng}{0.0086--0.0100}  % ... the range as text
\providecommand{\rhcaLinHarQlClosePre}{}\renewcommand{\rhcaLinHarQlClosePre}{0.0068}  % QLIKE level term, linear HAR, ClosePre
\providecommand{\rhcaLinFullQdClosePreMin}{}\renewcommand{\rhcaLinFullQdClosePreMin}{0.1213}  % QLIKE dispersion term, linear full, ClosePre: min
\providecommand{\rhcaLinFullQdClosePreMax}{}\renewcommand{\rhcaLinFullQdClosePreMax}{0.1233}  % ... max
\providecommand{\rhcaLinFullQdClosePreRng}{}\renewcommand{\rhcaLinFullQdClosePreRng}{0.1213--0.1233}  % ... the range as text
\providecommand{\rhcaLinHarQdClosePre}{}\renewcommand{\rhcaLinHarQdClosePre}{0.1297}  % QLIKE dispersion term, linear HAR, ClosePre
\providecommand{\rhcaTreeFullQlClosePreMin}{}\renewcommand{\rhcaTreeFullQlClosePreMin}{0.0045}  % QLIKE level term, tree full, ClosePre: min
\providecommand{\rhcaTreeFullQlClosePreMax}{}\renewcommand{\rhcaTreeFullQlClosePreMax}{0.0065}  % ... max
\providecommand{\rhcaTreeFullQlClosePreRng}{}\renewcommand{\rhcaTreeFullQlClosePreRng}{0.0045--0.0065}  % ... the range as text
\providecommand{\rhcaTreeHarQlClosePreMin}{}\renewcommand{\rhcaTreeHarQlClosePreMin}{0.0020}  % QLIKE level term, tree HAR, ClosePre: min
\providecommand{\rhcaTreeHarQlClosePreMax}{}\renewcommand{\rhcaTreeHarQlClosePreMax}{0.0039}  % ... max
\providecommand{\rhcaTreeHarQlClosePreRng}{}\renewcommand{\rhcaTreeHarQlClosePreRng}{0.0020--0.0039}  % ... the range as text
\providecommand{\rhcaTreeFullQdClosePreMin}{}\renewcommand{\rhcaTreeFullQdClosePreMin}{0.1232}  % QLIKE dispersion term, tree full, ClosePre: min
\providecommand{\rhcaTreeFullQdClosePreMax}{}\renewcommand{\rhcaTreeFullQdClosePreMax}{0.1270}  % ... max
\providecommand{\rhcaTreeFullQdClosePreRng}{}\renewcommand{\rhcaTreeFullQdClosePreRng}{0.1232--0.1270}  % ... the range as text
\providecommand{\rhcaTreeHarQdClosePreMin}{}\renewcommand{\rhcaTreeHarQdClosePreMin}{0.1310}  % QLIKE dispersion term, tree HAR, ClosePre: min
\providecommand{\rhcaTreeHarQdClosePreMax}{}\renewcommand{\rhcaTreeHarQdClosePreMax}{0.1325}  % ... max
\providecommand{\rhcaTreeHarQdClosePreRng}{}\renewcommand{\rhcaTreeHarQdClosePreRng}{0.1310--0.1325}  % ... the range as text
\providecommand{\rhcaTreeFullUnderClosePreMin}{}\renewcommand{\rhcaTreeFullUnderClosePreMin}{45.6}  % under-prediction %, tree full, ClosePre: min
\providecommand{\rhcaTreeFullUnderClosePreMax}{}\renewcommand{\rhcaTreeFullUnderClosePreMax}{46.1}  % ... max
\providecommand{\rhcaTreeFullUnderClosePreRng}{}\renewcommand{\rhcaTreeFullUnderClosePreRng}{45.6--46.1}  % ... the range as text
\providecommand{\rhcaTreeHarUnderClosePreMin}{}\renewcommand{\rhcaTreeHarUnderClosePreMin}{42.6}  % under-prediction %, tree HAR, ClosePre: min
\providecommand{\rhcaTreeHarUnderClosePreMax}{}\renewcommand{\rhcaTreeHarUnderClosePreMax}{43.6}  % ... max
\providecommand{\rhcaTreeHarUnderClosePreRng}{}\renewcommand{\rhcaTreeHarUnderClosePreRng}{42.6--43.6}  % ... the range as text
\providecommand{\rhcaLinFullMzBCloseOosMin}{}\renewcommand{\rhcaLinFullMzBCloseOosMin}{0.66}  % levels MZ slope, linear full, CloseOos: min
\providecommand{\rhcaLinFullMzBCloseOosMax}{}\renewcommand{\rhcaLinFullMzBCloseOosMax}{0.70}  % ... max
\providecommand{\rhcaLinFullMzBCloseOosRng}{}\renewcommand{\rhcaLinFullMzBCloseOosRng}{0.66--0.70}  % ... the range as text
\providecommand{\rhcaTreeFullMzBCloseOosMin}{}\renewcommand{\rhcaTreeFullMzBCloseOosMin}{0.81}  % levels MZ slope, tree full, CloseOos: min
\providecommand{\rhcaTreeFullMzBCloseOosMax}{}\renewcommand{\rhcaTreeFullMzBCloseOosMax}{0.89}  % ... max
\providecommand{\rhcaTreeFullMzBCloseOosRng}{}\renewcommand{\rhcaTreeFullMzBCloseOosRng}{0.81--0.89}  % ... the range as text
\providecommand{\rhcaTreeHarMzBCloseOosMin}{}\renewcommand{\rhcaTreeHarMzBCloseOosMin}{1.06}  % levels MZ slope, tree HAR, CloseOos: min
\providecommand{\rhcaTreeHarMzBCloseOosMax}{}\renewcommand{\rhcaTreeHarMzBCloseOosMax}{1.08}  % ... max
\providecommand{\rhcaTreeHarMzBCloseOosRng}{}\renewcommand{\rhcaTreeHarMzBCloseOosRng}{1.06--1.08}  % ... the range as text
\providecommand{\rhcaLinHarMzBCloseOos}{}\renewcommand{\rhcaLinHarMzBCloseOos}{1.21}  % levels MZ slope, linear HAR (= OLS), CloseOos
\providecommand{\rhcaLinFullMzlBCloseOosMin}{}\renewcommand{\rhcaLinFullMzlBCloseOosMin}{0.92}  % logs MZ slope, linear full, CloseOos: min
\providecommand{\rhcaLinFullMzlBCloseOosMax}{}\renewcommand{\rhcaLinFullMzlBCloseOosMax}{0.92}  % ... max
\providecommand{\rhcaLinFullMzlBCloseOosRng}{}\renewcommand{\rhcaLinFullMzlBCloseOosRng}{0.92}  % ... the range as text
\providecommand{\rhcaTreeFullMzlBCloseOosMin}{}\renewcommand{\rhcaTreeFullMzlBCloseOosMin}{0.93}  % logs MZ slope, tree full, CloseOos: min
\providecommand{\rhcaTreeFullMzlBCloseOosMax}{}\renewcommand{\rhcaTreeFullMzlBCloseOosMax}{0.95}  % ... max
\providecommand{\rhcaTreeFullMzlBCloseOosRng}{}\renewcommand{\rhcaTreeFullMzlBCloseOosRng}{0.93--0.95}  % ... the range as text
\providecommand{\rhcaAllHarMzlBCloseOosMin}{}\renewcommand{\rhcaAllHarMzlBCloseOosMin}{0.97}  % logs MZ slope, HAR forecasts, CloseOos: min
\providecommand{\rhcaAllHarMzlBCloseOosMax}{}\renewcommand{\rhcaAllHarMzlBCloseOosMax}{0.98}  % ... max
\providecommand{\rhcaAllHarMzlBCloseOosRng}{}\renewcommand{\rhcaAllHarMzlBCloseOosRng}{0.97--0.98}  % ... the range as text
\providecommand{\rhcaLinFullMzlRsqCloseOosMin}{}\renewcommand{\rhcaLinFullMzlRsqCloseOosMin}{0.84}  % logs MZ R2, linear full, CloseOos: min
\providecommand{\rhcaLinFullMzlRsqCloseOosMax}{}\renewcommand{\rhcaLinFullMzlRsqCloseOosMax}{0.84}  % ... max
\providecommand{\rhcaLinFullMzlRsqCloseOosRng}{}\renewcommand{\rhcaLinFullMzlRsqCloseOosRng}{0.84}  % ... the range as text
\providecommand{\rhcaTreeFullMzlRsqCloseOosMin}{}\renewcommand{\rhcaTreeFullMzlRsqCloseOosMin}{0.85}  % logs MZ R2, tree full, CloseOos: min
\providecommand{\rhcaTreeFullMzlRsqCloseOosMax}{}\renewcommand{\rhcaTreeFullMzlRsqCloseOosMax}{0.85}  % ... max
\providecommand{\rhcaTreeFullMzlRsqCloseOosRng}{}\renewcommand{\rhcaTreeFullMzlRsqCloseOosRng}{0.85}  % ... the range as text
\providecommand{\rhcaAllHarMzlRsqCloseOosMin}{}\renewcommand{\rhcaAllHarMzlRsqCloseOosMin}{0.83}  % logs MZ R2, HAR forecasts, CloseOos: min
\providecommand{\rhcaAllHarMzlRsqCloseOosMax}{}\renewcommand{\rhcaAllHarMzlRsqCloseOosMax}{0.84}  % ... max
\providecommand{\rhcaAllHarMzlRsqCloseOosRng}{}\renewcommand{\rhcaAllHarMzlRsqCloseOosRng}{0.83--0.84}  % ... the range as text
\providecommand{\rhcaLinFullMzRsqCloseOosMin}{}\renewcommand{\rhcaLinFullMzRsqCloseOosMin}{0.71}  % levels MZ R2, linear full, CloseOos: min
\providecommand{\rhcaLinFullMzRsqCloseOosMax}{}\renewcommand{\rhcaLinFullMzRsqCloseOosMax}{0.72}  % ... max
\providecommand{\rhcaLinFullMzRsqCloseOosRng}{}\renewcommand{\rhcaLinFullMzRsqCloseOosRng}{0.71--0.72}  % ... the range as text
\providecommand{\rhcaAllHarMzRsqCloseOosMin}{}\renewcommand{\rhcaAllHarMzRsqCloseOosMin}{0.80}  % levels MZ R2, HAR forecasts, CloseOos: min
\providecommand{\rhcaAllHarMzRsqCloseOosMax}{}\renewcommand{\rhcaAllHarMzRsqCloseOosMax}{0.82}  % ... max
\providecommand{\rhcaAllHarMzRsqCloseOosRng}{}\renewcommand{\rhcaAllHarMzRsqCloseOosRng}{0.80--0.82}  % ... the range as text
\providecommand{\rhcaLinFullShSlopeCloseOosMin}{}\renewcommand{\rhcaLinFullShSlopeCloseOosMin}{32.0}  % Theil slope share %, linear full, CloseOos: min
\providecommand{\rhcaLinFullShSlopeCloseOosMax}{}\renewcommand{\rhcaLinFullShSlopeCloseOosMax}{39.7}  % ... max
\providecommand{\rhcaLinFullShSlopeCloseOosRng}{}\renewcommand{\rhcaLinFullShSlopeCloseOosRng}{32.0--39.7}  % ... the range as text
\providecommand{\rhcaTreeFullShSlopeCloseOosMin}{}\renewcommand{\rhcaTreeFullShSlopeCloseOosMin}{5.5}  % Theil slope share %, tree full, CloseOos: min
\providecommand{\rhcaTreeFullShSlopeCloseOosMax}{}\renewcommand{\rhcaTreeFullShSlopeCloseOosMax}{14.8}  % ... max
\providecommand{\rhcaTreeFullShSlopeCloseOosRng}{}\renewcommand{\rhcaTreeFullShSlopeCloseOosRng}{5.5--14.8}  % ... the range as text
\providecommand{\rhcaAllHarShSlopeCloseOosMin}{}\renewcommand{\rhcaAllHarShSlopeCloseOosMin}{1.7}  % Theil slope share %, HAR forecasts, CloseOos: min
\providecommand{\rhcaAllHarShSlopeCloseOosMax}{}\renewcommand{\rhcaAllHarShSlopeCloseOosMax}{10.7}  % ... max
\providecommand{\rhcaAllHarShSlopeCloseOosRng}{}\renewcommand{\rhcaAllHarShSlopeCloseOosRng}{1.7--10.7}  % ... the range as text
\providecommand{\rhcaAllShBiasCloseOosMin}{}\renewcommand{\rhcaAllShBiasCloseOosMin}{0.0}  % Theil bias share %, all 12, CloseOos: min
\providecommand{\rhcaAllShBiasCloseOosMax}{}\renewcommand{\rhcaAllShBiasCloseOosMax}{0.3}  % ... max
\providecommand{\rhcaAllShBiasCloseOosRng}{}\renewcommand{\rhcaAllShBiasCloseOosRng}{0.0--0.3}  % ... the range as text
\providecommand{\rhcaAllRmCloseOosMin}{}\renewcommand{\rhcaAllRmCloseOosMin}{0.96}  % ratio of means, all 12, CloseOos: min
\providecommand{\rhcaAllRmCloseOosMax}{}\renewcommand{\rhcaAllRmCloseOosMax}{1.10}  % ... max
\providecommand{\rhcaAllRmCloseOosRng}{}\renewcommand{\rhcaAllRmCloseOosRng}{0.96--1.10}  % ... the range as text
\providecommand{\rhcaFullRmCloseOosMin}{}\renewcommand{\rhcaFullRmCloseOosMin}{0.96}  % ratio of means, full set, CloseOos: min
\providecommand{\rhcaFullRmCloseOosMax}{}\renewcommand{\rhcaFullRmCloseOosMax}{0.98}  % ... max
\providecommand{\rhcaFullRmCloseOosRng}{}\renewcommand{\rhcaFullRmCloseOosRng}{0.96--0.98}  % ... the range as text
\providecommand{\rhcaHarRmCloseOosMin}{}\renewcommand{\rhcaHarRmCloseOosMin}{1.03}  % ratio of means, HAR, CloseOos: min
\providecommand{\rhcaHarRmCloseOosMax}{}\renewcommand{\rhcaHarRmCloseOosMax}{1.10}  % ... max
\providecommand{\rhcaHarRmCloseOosRng}{}\renewcommand{\rhcaHarRmCloseOosRng}{1.03--1.10}  % ... the range as text
\providecommand{\rhcaAllMrCloseOosMin}{}\renewcommand{\rhcaAllMrCloseOosMin}{1.03}  % mean of RV/sigma2, all 12, CloseOos: min
\providecommand{\rhcaAllMrCloseOosMax}{}\renewcommand{\rhcaAllMrCloseOosMax}{1.18}  % ... max
\providecommand{\rhcaAllMrCloseOosRng}{}\renewcommand{\rhcaAllMrCloseOosRng}{1.03--1.18}  % ... the range as text
\providecommand{\rhcaAllUnderCloseOosMin}{}\renewcommand{\rhcaAllUnderCloseOosMin}{42.9}  % under-prediction %, all 12, CloseOos: min
\providecommand{\rhcaAllUnderCloseOosMax}{}\renewcommand{\rhcaAllUnderCloseOosMax}{52.9}  % ... max
\providecommand{\rhcaAllUnderCloseOosRng}{}\renewcommand{\rhcaAllUnderCloseOosRng}{42.9--52.9}  % ... the range as text
\providecommand{\rhcaLinFullUnderCloseOosMin}{}\renewcommand{\rhcaLinFullUnderCloseOosMin}{52.1}  % under-prediction %, linear full, CloseOos: min
\providecommand{\rhcaLinFullUnderCloseOosMax}{}\renewcommand{\rhcaLinFullUnderCloseOosMax}{52.9}  % ... max
\providecommand{\rhcaLinFullUnderCloseOosRng}{}\renewcommand{\rhcaLinFullUnderCloseOosRng}{52.1--52.9}  % ... the range as text
\providecommand{\rhcaLinHarUnderCloseOos}{}\renewcommand{\rhcaLinHarUnderCloseOos}{47.7}  % under-prediction %, linear HAR, CloseOos
\providecommand{\rhcaLinFullOverSeCloseOosMin}{}\renewcommand{\rhcaLinFullOverSeCloseOosMin}{51.5}  % % of SE from over-predictions, linear full, CloseOos: min
\providecommand{\rhcaLinFullOverSeCloseOosMax}{}\renewcommand{\rhcaLinFullOverSeCloseOosMax}{59.7}  % ... max
\providecommand{\rhcaLinFullOverSeCloseOosRng}{}\renewcommand{\rhcaLinFullOverSeCloseOosRng}{51.5--59.7}  % ... the range as text
\providecommand{\rhcaTreeFullOverSeCloseOosMin}{}\renewcommand{\rhcaTreeFullOverSeCloseOosMin}{43.1}  % % of SE from over-predictions, tree full, CloseOos: min
\providecommand{\rhcaTreeFullOverSeCloseOosMax}{}\renewcommand{\rhcaTreeFullOverSeCloseOosMax}{46.1}  % ... max
\providecommand{\rhcaTreeFullOverSeCloseOosRng}{}\renewcommand{\rhcaTreeFullOverSeCloseOosRng}{43.1--46.1}  % ... the range as text
\providecommand{\rhcaAllHarOverSeCloseOosMin}{}\renewcommand{\rhcaAllHarOverSeCloseOosMin}{13.8}  % % of SE from over-predictions, HAR, CloseOos: min
\providecommand{\rhcaAllHarOverSeCloseOosMax}{}\renewcommand{\rhcaAllHarOverSeCloseOosMax}{25.7}  % ... max
\providecommand{\rhcaAllHarOverSeCloseOosRng}{}\renewcommand{\rhcaAllHarOverSeCloseOosRng}{13.8--25.7}  % ... the range as text
\providecommand{\rhcaNMaeLowerCloseOos}{}\renewcommand{\rhcaNMaeLowerCloseOos}{0}  % classes (of 6) with the lower MAE on the full set, CloseOos
\providecommand{\rhcaLinQlShareCloseOosMin}{}\renewcommand{\rhcaLinQlShareCloseOosMin}{64}  % % of the linear full-set QLIKE increase from the level term, CloseOos: min
\providecommand{\rhcaLinQlShareCloseOosMax}{}\renewcommand{\rhcaLinQlShareCloseOosMax}{76}  % ... max
\providecommand{\rhcaLinQlShareCloseOosRng}{}\renewcommand{\rhcaLinQlShareCloseOosRng}{64--76}  % ... the range as text
\providecommand{\rhcaLinFullMrCloseOosMin}{}\renewcommand{\rhcaLinFullMrCloseOosMin}{1.165}  % mean of RV/sigma2, linear full, CloseOos: min
\providecommand{\rhcaLinFullMrCloseOosMax}{}\renewcommand{\rhcaLinFullMrCloseOosMax}{1.183}  % ... max
\providecommand{\rhcaLinFullMrCloseOosRng}{}\renewcommand{\rhcaLinFullMrCloseOosRng}{1.165--1.183}  % ... the range as text
\providecommand{\rhcaLinHarMrCloseOos}{}\renewcommand{\rhcaLinHarMrCloseOos}{1.092}  % mean of RV/sigma2, linear HAR, CloseOos
\providecommand{\rhcaLinFullQlCloseOosMin}{}\renewcommand{\rhcaLinFullQlCloseOosMin}{0.0123}  % QLIKE level term, linear full, CloseOos: min
\providecommand{\rhcaLinFullQlCloseOosMax}{}\renewcommand{\rhcaLinFullQlCloseOosMax}{0.0150}  % ... max
\providecommand{\rhcaLinFullQlCloseOosRng}{}\renewcommand{\rhcaLinFullQlCloseOosRng}{0.0123--0.0150}  % ... the range as text
\providecommand{\rhcaLinHarQlCloseOos}{}\renewcommand{\rhcaLinHarQlCloseOos}{0.0040}  % QLIKE level term, linear HAR, CloseOos
\providecommand{\rhcaLinFullQdCloseOosMin}{}\renewcommand{\rhcaLinFullQdCloseOosMin}{0.1122}  % QLIKE dispersion term, linear full, CloseOos: min
\providecommand{\rhcaLinFullQdCloseOosMax}{}\renewcommand{\rhcaLinFullQdCloseOosMax}{0.1159}  % ... max
\providecommand{\rhcaLinFullQdCloseOosRng}{}\renewcommand{\rhcaLinFullQdCloseOosRng}{0.1122--0.1159}  % ... the range as text
\providecommand{\rhcaLinHarQdCloseOos}{}\renewcommand{\rhcaLinHarQdCloseOos}{0.1097}  % QLIKE dispersion term, linear HAR, CloseOos
\providecommand{\rhcaTreeFullQlCloseOosMin}{}\renewcommand{\rhcaTreeFullQlCloseOosMin}{0.0039}  % QLIKE level term, tree full, CloseOos: min
\providecommand{\rhcaTreeFullQlCloseOosMax}{}\renewcommand{\rhcaTreeFullQlCloseOosMax}{0.0059}  % ... max
\providecommand{\rhcaTreeFullQlCloseOosRng}{}\renewcommand{\rhcaTreeFullQlCloseOosRng}{0.0039--0.0059}  % ... the range as text
\providecommand{\rhcaTreeHarQlCloseOosMin}{}\renewcommand{\rhcaTreeHarQlCloseOosMin}{0.0005}  % QLIKE level term, tree HAR, CloseOos: min
\providecommand{\rhcaTreeHarQlCloseOosMax}{}\renewcommand{\rhcaTreeHarQlCloseOosMax}{0.0019}  % ... max
\providecommand{\rhcaTreeHarQlCloseOosRng}{}\renewcommand{\rhcaTreeHarQlCloseOosRng}{0.0005--0.0019}  % ... the range as text
\providecommand{\rhcaTreeFullQdCloseOosMin}{}\renewcommand{\rhcaTreeFullQdCloseOosMin}{0.1044}  % QLIKE dispersion term, tree full, CloseOos: min
\providecommand{\rhcaTreeFullQdCloseOosMax}{}\renewcommand{\rhcaTreeFullQdCloseOosMax}{0.1064}  % ... max
\providecommand{\rhcaTreeFullQdCloseOosRng}{}\renewcommand{\rhcaTreeFullQdCloseOosRng}{0.1044--0.1064}  % ... the range as text
\providecommand{\rhcaTreeHarQdCloseOosMin}{}\renewcommand{\rhcaTreeHarQdCloseOosMin}{0.1095}  % QLIKE dispersion term, tree HAR, CloseOos: min
\providecommand{\rhcaTreeHarQdCloseOosMax}{}\renewcommand{\rhcaTreeHarQdCloseOosMax}{0.1112}  % ... max
\providecommand{\rhcaTreeHarQdCloseOosRng}{}\renewcommand{\rhcaTreeHarQdCloseOosRng}{0.1095--0.1112}  % ... the range as text
\providecommand{\rhcaTreeFullUnderCloseOosMin}{}\renewcommand{\rhcaTreeFullUnderCloseOosMin}{47.1}  % under-prediction %, tree full, CloseOos: min
\providecommand{\rhcaTreeFullUnderCloseOosMax}{}\renewcommand{\rhcaTreeFullUnderCloseOosMax}{50.1}  % ... max
\providecommand{\rhcaTreeFullUnderCloseOosRng}{}\renewcommand{\rhcaTreeFullUnderCloseOosRng}{47.1--50.1}  % ... the range as text
\providecommand{\rhcaTreeHarUnderCloseOosMin}{}\renewcommand{\rhcaTreeHarUnderCloseOosMin}{42.9}  % under-prediction %, tree HAR, CloseOos: min
\providecommand{\rhcaTreeHarUnderCloseOosMax}{}\renewcommand{\rhcaTreeHarUnderCloseOosMax}{45.5}  % ... max
\providecommand{\rhcaTreeHarUnderCloseOosRng}{}\renewcommand{\rhcaTreeHarUnderCloseOosRng}{42.9--45.5}  % ... the range as text
\providecommand{\rhcaLinTopShareBarOosMin}{}\renewcommand{\rhcaLinTopShareBarOosMin}{92}  % % of the full-minus-HAR squared-error sum (every bar, Oos) on the 5 days of largest difference, linear: min
\providecommand{\rhcaLinTopShareBarOosMax}{}\renewcommand{\rhcaLinTopShareBarOosMax}{94}  % ... max
\providecommand{\rhcaLinTopShareBarOosRng}{}\renewcommand{\rhcaLinTopShareBarOosRng}{92--94}  % ... the range as text
\providecommand{\rhcaTreeTopShareBarOosMin}{}\renewcommand{\rhcaTreeTopShareBarOosMin}{92}  % ... tree ensembles: min
\providecommand{\rhcaTreeTopShareBarOosMax}{}\renewcommand{\rhcaTreeTopShareBarOosMax}{101}  % ... max
\providecommand{\rhcaTreeTopShareBarOosRng}{}\renewcommand{\rhcaTreeTopShareBarOosRng}{92--101}  % ... the range as text
\providecommand{\rhcaTopFirstBarOos}{}\renewcommand{\rhcaTopFirstBarOos}{12 March 2020}  % earliest of those days, any class
\providecommand{\rhcaTopLastBarOos}{}\renewcommand{\rhcaTopLastBarOos}{2 April 2020}  % latest of those days, any class
\providecommand{\rhcaLinMsePctBarOosExCrashMin}{}\renewcommand{\rhcaLinMsePctBarOosExCrashMin}{\ensuremath{-}7.3}  % linear every-bar MSE %, full vs HAR, Oos without the crash window: min
\providecommand{\rhcaLinMsePctBarOosExCrashMax}{}\renewcommand{\rhcaLinMsePctBarOosExCrashMax}{\ensuremath{-}5.2}  % ... max
\providecommand{\rhcaLinMsePctBarOosExCrashRng}{}\renewcommand{\rhcaLinMsePctBarOosExCrashRng}{\ensuremath{-}7.3 to \ensuremath{-}5.2}  % ... the range as text
\providecommand{\rhcaTreeMsePctBarOosExCrashMin}{}\renewcommand{\rhcaTreeMsePctBarOosExCrashMin}{\ensuremath{-}9.0}  % ... tree ensembles: min
\providecommand{\rhcaTreeMsePctBarOosExCrashMax}{}\renewcommand{\rhcaTreeMsePctBarOosExCrashMax}{\ensuremath{-}6.0}  % ... max
\providecommand{\rhcaTreeMsePctBarOosExCrashRng}{}\renewcommand{\rhcaTreeMsePctBarOosExCrashRng}{\ensuremath{-}9.0 to \ensuremath{-}6.0}  % ... the range as text
\providecommand{\rhcaNMseLowerBarOosExCrash}{}\renewcommand{\rhcaNMseLowerBarOosExCrash}{6}  % classes (of 6) with the lower every-bar MSE on the full set, Oos without the crash window
\providecommand{\rhcaNMseLowerSigBarOosExCrash}{}\renewcommand{\rhcaNMseLowerSigBarOosExCrash}{6}  % ... significantly (DM on the daily mean squared error)
\providecommand{\rhcaLinMsePctCloseOosExCrashMin}{}\renewcommand{\rhcaLinMsePctCloseOosExCrashMin}{6.6}  % linear 16:00 MSE %, full vs HAR, Oos without the crash window: min
\providecommand{\rhcaLinMsePctCloseOosExCrashMax}{}\renewcommand{\rhcaLinMsePctCloseOosExCrashMax}{10.0}  % ... max
\providecommand{\rhcaLinMsePctCloseOosExCrashRng}{}\renewcommand{\rhcaLinMsePctCloseOosExCrashRng}{6.6--10.0}  % ... the range as text
\providecommand{\rhcaNMseLowerCloseOosExCrash}{}\renewcommand{\rhcaNMseLowerCloseOosExCrash}{1}  % classes with the lower 16:00 MSE on the full set, Oos without the crash window
\providecommand{\rhcaNMseHigherSigCloseOosExCrash}{}\renewcommand{\rhcaNMseHigherSigCloseOosExCrash}{0}  % ... classes with a significantly higher 16:00 MSE there
\providecommand{\rhcaTreeMsePctCloseOosExCrashMin}{}\renewcommand{\rhcaTreeMsePctCloseOosExCrashMin}{\ensuremath{-}0.7}  % tree 16:00 MSE %, full vs HAR, Oos without the crash window: min
\providecommand{\rhcaTreeMsePctCloseOosExCrashMax}{}\renewcommand{\rhcaTreeMsePctCloseOosExCrashMax}{3.4}  % ... max
\providecommand{\rhcaTreeMsePctCloseOosExCrashRng}{}\renewcommand{\rhcaTreeMsePctCloseOosExCrashRng}{\ensuremath{-}0.7 to 3.4}  % ... the range as text
\providecommand{\rhcaAllMaePctBarOosMin}{}\renewcommand{\rhcaAllMaePctBarOosMin}{2.3}  % every-bar MAE %, full vs HAR, Oos, all six classes: min
\providecommand{\rhcaAllMaePctBarOosMax}{}\renewcommand{\rhcaAllMaePctBarOosMax}{12.8}  % ... max
\providecommand{\rhcaAllMaePctBarOosRng}{}\renewcommand{\rhcaAllMaePctBarOosRng}{2.3--12.8}  % ... the range as text
\providecommand{\rhcaAllMaePctBarOosExCrashMin}{}\renewcommand{\rhcaAllMaePctBarOosExCrashMin}{\ensuremath{-}6.3}  % ... without the crash window: min
\providecommand{\rhcaAllMaePctBarOosExCrashMax}{}\renewcommand{\rhcaAllMaePctBarOosExCrashMax}{\ensuremath{-}4.4}  % ... max
\providecommand{\rhcaAllMaePctBarOosExCrashRng}{}\renewcommand{\rhcaAllMaePctBarOosExCrashRng}{\ensuremath{-}6.3 to \ensuremath{-}4.4}  % ... the range as text
\providecommand{\rhcaNMaeLowerBarOosExCrash}{}\renewcommand{\rhcaNMaeLowerBarOosExCrash}{6}  % classes with the lower every-bar MAE on the full set, Oos without the crash window
\providecommand{\rhcaNMaeLowerSigBarOosExCrash}{}\renewcommand{\rhcaNMaeLowerSigBarOosExCrash}{6}  % ... significantly
\providecommand{\rhcaTopShareRidgeBarOos}{}\renewcommand{\rhcaTopShareRidgeBarOos}{94}  % top-day share, %
\providecommand{\rhcaMsePctRidgeBarOosExCrash}{}\renewcommand{\rhcaMsePctRidgeBarOosExCrash}{\ensuremath{-}5.2}  % every-bar MSE %, ex crash
\providecommand{\rhcaMseDmRidgeBarOosExCrash}{}\renewcommand{\rhcaMseDmRidgeBarOosExCrash}{\ensuremath{-}3.10}
\providecommand{\rhcaMsePRidgeBarOosExCrash}{}\renewcommand{\rhcaMsePRidgeBarOosExCrash}{\ensuremath{p = 0.002}}
\providecommand{\rhcaTopShareLassoBarOos}{}\renewcommand{\rhcaTopShareLassoBarOos}{92}  % top-day share, %
\providecommand{\rhcaMsePctLassoBarOosExCrash}{}\renewcommand{\rhcaMsePctLassoBarOosExCrash}{\ensuremath{-}7.3}  % every-bar MSE %, ex crash
\providecommand{\rhcaMseDmLassoBarOosExCrash}{}\renewcommand{\rhcaMseDmLassoBarOosExCrash}{\ensuremath{-}3.29}
\providecommand{\rhcaMsePLassoBarOosExCrash}{}\renewcommand{\rhcaMsePLassoBarOosExCrash}{\ensuremath{p = 0.001}}
\providecommand{\rhcaTopShareEnetBarOos}{}\renewcommand{\rhcaTopShareEnetBarOos}{94}  % top-day share, %
\providecommand{\rhcaMsePctEnetBarOosExCrash}{}\renewcommand{\rhcaMsePctEnetBarOosExCrash}{\ensuremath{-}6.0}  % every-bar MSE %, ex crash
\providecommand{\rhcaMseDmEnetBarOosExCrash}{}\renewcommand{\rhcaMseDmEnetBarOosExCrash}{\ensuremath{-}3.43}
\providecommand{\rhcaMsePEnetBarOosExCrash}{}\renewcommand{\rhcaMsePEnetBarOosExCrash}{\ensuremath{p < 0.001}}
\providecommand{\rhcaTopShareLgbmBarOos}{}\renewcommand{\rhcaTopShareLgbmBarOos}{101}  % top-day share, %
\providecommand{\rhcaMsePctLgbmBarOosExCrash}{}\renewcommand{\rhcaMsePctLgbmBarOosExCrash}{\ensuremath{-}8.8}  % every-bar MSE %, ex crash
\providecommand{\rhcaMseDmLgbmBarOosExCrash}{}\renewcommand{\rhcaMseDmLgbmBarOosExCrash}{\ensuremath{-}4.30}
\providecommand{\rhcaMsePLgbmBarOosExCrash}{}\renewcommand{\rhcaMsePLgbmBarOosExCrash}{\ensuremath{p < 0.001}}
\providecommand{\rhcaTopShareXgbBarOos}{}\renewcommand{\rhcaTopShareXgbBarOos}{99}  % top-day share, %
\providecommand{\rhcaMsePctXgbBarOosExCrash}{}\renewcommand{\rhcaMsePctXgbBarOosExCrash}{\ensuremath{-}9.0}  % every-bar MSE %, ex crash
\providecommand{\rhcaMseDmXgbBarOosExCrash}{}\renewcommand{\rhcaMseDmXgbBarOosExCrash}{\ensuremath{-}3.67}
\providecommand{\rhcaMsePXgbBarOosExCrash}{}\renewcommand{\rhcaMsePXgbBarOosExCrash}{\ensuremath{p < 0.001}}
\providecommand{\rhcaTopShareRfBarOos}{}\renewcommand{\rhcaTopShareRfBarOos}{92}  % top-day share, %
\providecommand{\rhcaMsePctRfBarOosExCrash}{}\renewcommand{\rhcaMsePctRfBarOosExCrash}{\ensuremath{-}6.0}  % every-bar MSE %, ex crash
\providecommand{\rhcaMseDmRfBarOosExCrash}{}\renewcommand{\rhcaMseDmRfBarOosExCrash}{\ensuremath{-}3.97}
\providecommand{\rhcaMsePRfBarOosExCrash}{}\renewcommand{\rhcaMsePRfBarOosExCrash}{\ensuremath{p < 0.001}}
\providecommand{\rhcaMaxFcRidge}{}\renewcommand{\rhcaMaxFcRidge}{1947}  % largest full-set forecast on a bar, Oos, x 10^6
\providecommand{\rhcaMaxFcRvRidge}{}\renewcommand{\rhcaMaxFcRvRidge}{501}  % the realized variance of that bar, x 10^6
\providecommand{\rhcaMaxFcHarRidge}{}\renewcommand{\rhcaMaxFcHarRidge}{627}  % the HAR forecast of that bar, x 10^6
\providecommand{\rhcaMaxFcRatioRidge}{}\renewcommand{\rhcaMaxFcRatioRidge}{3.9}  % forecast / realized variance on that bar
\providecommand{\rhcaMaxFcLasso}{}\renewcommand{\rhcaMaxFcLasso}{2035}  % largest full-set forecast on a bar, Oos, x 10^6
\providecommand{\rhcaMaxFcRvLasso}{}\renewcommand{\rhcaMaxFcRvLasso}{501}  % the realized variance of that bar, x 10^6
\providecommand{\rhcaMaxFcHarLasso}{}\renewcommand{\rhcaMaxFcHarLasso}{627}  % the HAR forecast of that bar, x 10^6
\providecommand{\rhcaMaxFcRatioLasso}{}\renewcommand{\rhcaMaxFcRatioLasso}{4.1}  % forecast / realized variance on that bar
\providecommand{\rhcaMaxFcEnet}{}\renewcommand{\rhcaMaxFcEnet}{2064}  % largest full-set forecast on a bar, Oos, x 10^6
\providecommand{\rhcaMaxFcRvEnet}{}\renewcommand{\rhcaMaxFcRvEnet}{501}  % the realized variance of that bar, x 10^6
\providecommand{\rhcaMaxFcHarEnet}{}\renewcommand{\rhcaMaxFcHarEnet}{627}  % the HAR forecast of that bar, x 10^6
\providecommand{\rhcaMaxFcRatioEnet}{}\renewcommand{\rhcaMaxFcRatioEnet}{4.1}  % forecast / realized variance on that bar
\providecommand{\rhcaMaxFcDate}{}\renewcommand{\rhcaMaxFcDate}{23 March 2020}  % the bar of the largest linear full-set forecast (all three)
\providecommand{\rhcaMaxFcClock}{}\renewcommand{\rhcaMaxFcClock}{10:00}  % its bar end, ET
\providecommand{\rhcaMaxRvOos}{}\renewcommand{\rhcaMaxRvOos}{1150}  % largest bar RV on Oos, x 10^6
\providecommand{\rhcaYrsUpRidgeCloseOos}{}\renewcommand{\rhcaYrsUpRidgeCloseOos}{4}  % calendar years of Oos (of 5) with a higher 16:00 QLIKE on the full set
\providecommand{\rhcaYrsDownRidgeClosePre}{}\renewcommand{\rhcaYrsDownRidgeClosePre}{4}  % calendar years of Pre (of 6) with a lower 16:00 QLIKE
\providecommand{\rhcaWorstYrRidgeCloseOos}{}\renewcommand{\rhcaWorstYrRidgeCloseOos}{2023}  % the Oos year of the largest full-set increase
\providecommand{\rhcaWorstYrPctRidgeCloseOos}{}\renewcommand{\rhcaWorstYrPctRidgeCloseOos}{32.0}  % ... its QLIKE %
\providecommand{\rhcaYrsUpLassoCloseOos}{}\renewcommand{\rhcaYrsUpLassoCloseOos}{4}  % calendar years of Oos (of 5) with a higher 16:00 QLIKE on the full set
\providecommand{\rhcaYrsDownLassoClosePre}{}\renewcommand{\rhcaYrsDownLassoClosePre}{3}  % calendar years of Pre (of 6) with a lower 16:00 QLIKE
\providecommand{\rhcaWorstYrLassoCloseOos}{}\renewcommand{\rhcaWorstYrLassoCloseOos}{2023}  % the Oos year of the largest full-set increase
\providecommand{\rhcaWorstYrPctLassoCloseOos}{}\renewcommand{\rhcaWorstYrPctLassoCloseOos}{33.0}  % ... its QLIKE %
\providecommand{\rhcaYrsUpEnetCloseOos}{}\renewcommand{\rhcaYrsUpEnetCloseOos}{4}  % calendar years of Oos (of 5) with a higher 16:00 QLIKE on the full set
\providecommand{\rhcaYrsDownEnetClosePre}{}\renewcommand{\rhcaYrsDownEnetClosePre}{5}  % calendar years of Pre (of 6) with a lower 16:00 QLIKE
\providecommand{\rhcaWorstYrEnetCloseOos}{}\renewcommand{\rhcaWorstYrEnetCloseOos}{2023}  % the Oos year of the largest full-set increase
\providecommand{\rhcaWorstYrPctEnetCloseOos}{}\renewcommand{\rhcaWorstYrPctEnetCloseOos}{30.7}  % ... its QLIKE %
\providecommand{\rhcaYrsUpLgbmCloseOos}{}\renewcommand{\rhcaYrsUpLgbmCloseOos}{2}  % calendar years of Oos (of 5) with a higher 16:00 QLIKE on the full set
\providecommand{\rhcaYrsDownLgbmClosePre}{}\renewcommand{\rhcaYrsDownLgbmClosePre}{4}  % calendar years of Pre (of 6) with a lower 16:00 QLIKE
\providecommand{\rhcaWorstYrLgbmCloseOos}{}\renewcommand{\rhcaWorstYrLgbmCloseOos}{2024}  % the Oos year of the largest full-set increase
\providecommand{\rhcaWorstYrPctLgbmCloseOos}{}\renewcommand{\rhcaWorstYrPctLgbmCloseOos}{10.5}  % ... its QLIKE %
\providecommand{\rhcaYrsUpXgbCloseOos}{}\renewcommand{\rhcaYrsUpXgbCloseOos}{2}  % calendar years of Oos (of 5) with a higher 16:00 QLIKE on the full set
\providecommand{\rhcaYrsDownXgbClosePre}{}\renewcommand{\rhcaYrsDownXgbClosePre}{4}  % calendar years of Pre (of 6) with a lower 16:00 QLIKE
\providecommand{\rhcaWorstYrXgbCloseOos}{}\renewcommand{\rhcaWorstYrXgbCloseOos}{2024}  % the Oos year of the largest full-set increase
\providecommand{\rhcaWorstYrPctXgbCloseOos}{}\renewcommand{\rhcaWorstYrPctXgbCloseOos}{11.9}  % ... its QLIKE %
\providecommand{\rhcaYrsUpRfCloseOos}{}\renewcommand{\rhcaYrsUpRfCloseOos}{2}  % calendar years of Oos (of 5) with a higher 16:00 QLIKE on the full set
\providecommand{\rhcaYrsDownRfClosePre}{}\renewcommand{\rhcaYrsDownRfClosePre}{4}  % calendar years of Pre (of 6) with a lower 16:00 QLIKE
\providecommand{\rhcaWorstYrRfCloseOos}{}\renewcommand{\rhcaWorstYrRfCloseOos}{2021}  % the Oos year of the largest full-set increase
\providecommand{\rhcaWorstYrPctRfCloseOos}{}\renewcommand{\rhcaWorstYrPctRfCloseOos}{7.3}  % ... its QLIKE %
\providecommand{\rhcaNYrsOos}{}\renewcommand{\rhcaNYrsOos}{5}  % calendar years of Oos
\providecommand{\rhcaNYrsPre}{}\renewcommand{\rhcaNYrsPre}{6}  % calendar years of Pre
\providecommand{\rhcaLinMedLogRatioClosePreMin}{}\renewcommand{\rhcaLinMedLogRatioClosePreMin}{\ensuremath{-}3.1}  % median 100 ln(full / HAR forecast), linear, 16:00, Pre: min
\providecommand{\rhcaLinMedLogRatioClosePreMax}{}\renewcommand{\rhcaLinMedLogRatioClosePreMax}{\ensuremath{-}2.1}  % ... max
\providecommand{\rhcaLinMedLogRatioClosePreRng}{}\renewcommand{\rhcaLinMedLogRatioClosePreRng}{\ensuremath{-}3.1 to \ensuremath{-}2.1}  % ... the range as text
\providecommand{\rhcaTreeMedLogRatioClosePreMin}{}\renewcommand{\rhcaTreeMedLogRatioClosePreMin}{\ensuremath{-}5.1}  % median 100 ln(full / HAR forecast), trees, 16:00, Pre: min
\providecommand{\rhcaTreeMedLogRatioClosePreMax}{}\renewcommand{\rhcaTreeMedLogRatioClosePreMax}{\ensuremath{-}3.6}  % ... max
\providecommand{\rhcaTreeMedLogRatioClosePreRng}{}\renewcommand{\rhcaTreeMedLogRatioClosePreRng}{\ensuremath{-}5.1 to \ensuremath{-}3.6}  % ... the range as text
\providecommand{\rhcaLinMedLogRatioCloseOosMin}{}\renewcommand{\rhcaLinMedLogRatioCloseOosMin}{\ensuremath{-}6.9}  % median 100 ln(full / HAR forecast), linear, 16:00, Oos: min
\providecommand{\rhcaLinMedLogRatioCloseOosMax}{}\renewcommand{\rhcaLinMedLogRatioCloseOosMax}{\ensuremath{-}5.6}  % ... max
\providecommand{\rhcaLinMedLogRatioCloseOosRng}{}\renewcommand{\rhcaLinMedLogRatioCloseOosRng}{\ensuremath{-}6.9 to \ensuremath{-}5.6}  % ... the range as text
\providecommand{\rhcaTreeMedLogRatioCloseOosMin}{}\renewcommand{\rhcaTreeMedLogRatioCloseOosMin}{\ensuremath{-}8.6}  % median 100 ln(full / HAR forecast), trees, 16:00, Oos: min
\providecommand{\rhcaTreeMedLogRatioCloseOosMax}{}\renewcommand{\rhcaTreeMedLogRatioCloseOosMax}{\ensuremath{-}6.1}  % ... max
\providecommand{\rhcaTreeMedLogRatioCloseOosRng}{}\renewcommand{\rhcaTreeMedLogRatioCloseOosRng}{\ensuremath{-}8.6 to \ensuremath{-}6.1}  % ... the range as text
\providecommand{\rhcaTopDecRatioEnetHarBarPre}{}\renewcommand{\rhcaTopDecRatioEnetHarBarPre}{1.04}  % top forecast decile: mean RV / mean forecast
\providecommand{\rhcaTopDecSeEnetHarBarPre}{}\renewcommand{\rhcaTopDecSeEnetHarBarPre}{85}  % top forecast decile: % of the squared error
\providecommand{\rhcaBotDecRatioEnetHarBarPre}{}\renewcommand{\rhcaBotDecRatioEnetHarBarPre}{1.12}  % bottom forecast decile: mean RV / mean forecast
\providecommand{\rhcaTopDecRatioEnetFullBarPre}{}\renewcommand{\rhcaTopDecRatioEnetFullBarPre}{1.01}  % top forecast decile: mean RV / mean forecast
\providecommand{\rhcaTopDecSeEnetFullBarPre}{}\renewcommand{\rhcaTopDecSeEnetFullBarPre}{84}  % top forecast decile: % of the squared error
\providecommand{\rhcaBotDecRatioEnetFullBarPre}{}\renewcommand{\rhcaBotDecRatioEnetFullBarPre}{1.16}  % bottom forecast decile: mean RV / mean forecast
\providecommand{\rhcaTopDecRatioLgbmHarBarPre}{}\renewcommand{\rhcaTopDecRatioLgbmHarBarPre}{1.08}  % top forecast decile: mean RV / mean forecast
\providecommand{\rhcaTopDecSeLgbmHarBarPre}{}\renewcommand{\rhcaTopDecSeLgbmHarBarPre}{88}  % top forecast decile: % of the squared error
\providecommand{\rhcaBotDecRatioLgbmHarBarPre}{}\renewcommand{\rhcaBotDecRatioLgbmHarBarPre}{1.14}  % bottom forecast decile: mean RV / mean forecast
\providecommand{\rhcaTopDecRatioLgbmFullBarPre}{}\renewcommand{\rhcaTopDecRatioLgbmFullBarPre}{1.06}  % top forecast decile: mean RV / mean forecast
\providecommand{\rhcaTopDecSeLgbmFullBarPre}{}\renewcommand{\rhcaTopDecSeLgbmFullBarPre}{91}  % top forecast decile: % of the squared error
\providecommand{\rhcaBotDecRatioLgbmFullBarPre}{}\renewcommand{\rhcaBotDecRatioLgbmFullBarPre}{1.09}  % bottom forecast decile: mean RV / mean forecast
\providecommand{\rhcaTopDecRatioEnetHarBarOos}{}\renewcommand{\rhcaTopDecRatioEnetHarBarOos}{1.03}  % top forecast decile: mean RV / mean forecast
\providecommand{\rhcaTopDecSeEnetHarBarOos}{}\renewcommand{\rhcaTopDecSeEnetHarBarOos}{91}  % top forecast decile: % of the squared error
\providecommand{\rhcaBotDecRatioEnetHarBarOos}{}\renewcommand{\rhcaBotDecRatioEnetHarBarOos}{1.06}  % bottom forecast decile: mean RV / mean forecast
\providecommand{\rhcaTopDecRatioEnetFullBarOos}{}\renewcommand{\rhcaTopDecRatioEnetFullBarOos}{0.88}  % top forecast decile: mean RV / mean forecast
\providecommand{\rhcaTopDecSeEnetFullBarOos}{}\renewcommand{\rhcaTopDecSeEnetFullBarOos}{97}  % top forecast decile: % of the squared error
\providecommand{\rhcaBotDecRatioEnetFullBarOos}{}\renewcommand{\rhcaBotDecRatioEnetFullBarOos}{1.14}  % bottom forecast decile: mean RV / mean forecast
\providecommand{\rhcaTopDecRatioLgbmHarBarOos}{}\renewcommand{\rhcaTopDecRatioLgbmHarBarOos}{1.00}  % top forecast decile: mean RV / mean forecast
\providecommand{\rhcaTopDecSeLgbmHarBarOos}{}\renewcommand{\rhcaTopDecSeLgbmHarBarOos}{92}  % top forecast decile: % of the squared error
\providecommand{\rhcaBotDecRatioLgbmHarBarOos}{}\renewcommand{\rhcaBotDecRatioLgbmHarBarOos}{1.07}  % bottom forecast decile: mean RV / mean forecast
\providecommand{\rhcaTopDecRatioLgbmFullBarOos}{}\renewcommand{\rhcaTopDecRatioLgbmFullBarOos}{0.89}  % top forecast decile: mean RV / mean forecast
\providecommand{\rhcaTopDecSeLgbmFullBarOos}{}\renewcommand{\rhcaTopDecSeLgbmFullBarOos}{95}  % top forecast decile: % of the squared error
\providecommand{\rhcaBotDecRatioLgbmFullBarOos}{}\renewcommand{\rhcaBotDecRatioLgbmFullBarOos}{1.12}  % bottom forecast decile: mean RV / mean forecast
\providecommand{\rhcaTopDecRatioEnetHarClosePre}{}\renewcommand{\rhcaTopDecRatioEnetHarClosePre}{1.11}  % top forecast decile: mean RV / mean forecast
\providecommand{\rhcaTopDecSeEnetHarClosePre}{}\renewcommand{\rhcaTopDecSeEnetHarClosePre}{89}  % top forecast decile: % of the squared error
\providecommand{\rhcaBotDecRatioEnetHarClosePre}{}\renewcommand{\rhcaBotDecRatioEnetHarClosePre}{1.31}  % bottom forecast decile: mean RV / mean forecast
\providecommand{\rhcaTopDecRatioEnetFullClosePre}{}\renewcommand{\rhcaTopDecRatioEnetFullClosePre}{1.03}  % top forecast decile: mean RV / mean forecast
\providecommand{\rhcaTopDecSeEnetFullClosePre}{}\renewcommand{\rhcaTopDecSeEnetFullClosePre}{88}  % top forecast decile: % of the squared error
\providecommand{\rhcaBotDecRatioEnetFullClosePre}{}\renewcommand{\rhcaBotDecRatioEnetFullClosePre}{1.31}  % bottom forecast decile: mean RV / mean forecast
\providecommand{\rhcaTopDecRatioLgbmHarClosePre}{}\renewcommand{\rhcaTopDecRatioLgbmHarClosePre}{1.08}  % top forecast decile: mean RV / mean forecast
\providecommand{\rhcaTopDecSeLgbmHarClosePre}{}\renewcommand{\rhcaTopDecSeLgbmHarClosePre}{91}  % top forecast decile: % of the squared error
\providecommand{\rhcaBotDecRatioLgbmHarClosePre}{}\renewcommand{\rhcaBotDecRatioLgbmHarClosePre}{1.19}  % bottom forecast decile: mean RV / mean forecast
\providecommand{\rhcaTopDecRatioLgbmFullClosePre}{}\renewcommand{\rhcaTopDecRatioLgbmFullClosePre}{1.01}  % top forecast decile: mean RV / mean forecast
\providecommand{\rhcaTopDecSeLgbmFullClosePre}{}\renewcommand{\rhcaTopDecSeLgbmFullClosePre}{91}  % top forecast decile: % of the squared error
\providecommand{\rhcaBotDecRatioLgbmFullClosePre}{}\renewcommand{\rhcaBotDecRatioLgbmFullClosePre}{1.23}  % bottom forecast decile: mean RV / mean forecast
\providecommand{\rhcaTopDecRatioEnetHarCloseOos}{}\renewcommand{\rhcaTopDecRatioEnetHarCloseOos}{1.13}  % top forecast decile: mean RV / mean forecast
\providecommand{\rhcaTopDecSeEnetHarCloseOos}{}\renewcommand{\rhcaTopDecSeEnetHarCloseOos}{98}  % top forecast decile: % of the squared error
\providecommand{\rhcaBotDecRatioEnetHarCloseOos}{}\renewcommand{\rhcaBotDecRatioEnetHarCloseOos}{1.15}  % bottom forecast decile: mean RV / mean forecast
\providecommand{\rhcaTopDecRatioEnetFullCloseOos}{}\renewcommand{\rhcaTopDecRatioEnetFullCloseOos}{0.87}  % top forecast decile: mean RV / mean forecast
\providecommand{\rhcaTopDecSeEnetFullCloseOos}{}\renewcommand{\rhcaTopDecSeEnetFullCloseOos}{99}  % top forecast decile: % of the squared error
\providecommand{\rhcaBotDecRatioEnetFullCloseOos}{}\renewcommand{\rhcaBotDecRatioEnetFullCloseOos}{1.31}  % bottom forecast decile: mean RV / mean forecast
\providecommand{\rhcaTopDecRatioLgbmHarCloseOos}{}\renewcommand{\rhcaTopDecRatioLgbmHarCloseOos}{1.02}  % top forecast decile: mean RV / mean forecast
\providecommand{\rhcaTopDecSeLgbmHarCloseOos}{}\renewcommand{\rhcaTopDecSeLgbmHarCloseOos}{97}  % top forecast decile: % of the squared error
\providecommand{\rhcaBotDecRatioLgbmHarCloseOos}{}\renewcommand{\rhcaBotDecRatioLgbmHarCloseOos}{1.11}  % bottom forecast decile: mean RV / mean forecast
\providecommand{\rhcaTopDecRatioLgbmFullCloseOos}{}\renewcommand{\rhcaTopDecRatioLgbmFullCloseOos}{0.91}  % top forecast decile: mean RV / mean forecast
\providecommand{\rhcaTopDecSeLgbmFullCloseOos}{}\renewcommand{\rhcaTopDecSeLgbmFullCloseOos}{98}  % top forecast decile: % of the squared error
\providecommand{\rhcaBotDecRatioLgbmFullCloseOos}{}\renewcommand{\rhcaBotDecRatioLgbmFullCloseOos}{1.22}  % bottom forecast decile: mean RV / mean forecast
\providecommand{\rhcaHarLinGap}{}\renewcommand{\rhcaHarLinGap}{5 \times 10^{-6}}  % max relative gap, HAR linear vs OLS forecasts (math mode)
\fi
\begin{table}[tbp]
\centering
\small
\setlength{\tabcolsep}{4pt}
\begin{threeparttable}
\caption{Forecast accuracy for the closing half hour on both evaluation spans}
\label{tab:rh_close_accuracy}
\begin{tabular*}{\linewidth}{@{\extracolsep{\fill}}llrrrrrrr@{}}
\toprule
 & & & & & \multicolumn{2}{c}{Full vs HAR} & \multicolumn{2}{c}{vs closing HAR} \\
\cmidrule(lr){6-7}\cmidrule(l){8-9}
Model & Set & QLIKE & MSE & $R^2_{\mathrm{OOS}}$ (\%) & $\Delta\%$ & $\DM$ & $\Delta\%$ & $\DM$ \\
\midrule
\multicolumn{9}{@{}l}{\textit{Panel A: \rhcaStartSpan{} to \rhcaEndSpan{} ($\rhcaNdaysSpan$ days), pooled intraday models}} \\
\addlinespace[2pt]
Ridge & HAR & $0.1269$ & $1.97$ & $77.4$ &  &  & -- &  \\
 & Full & $0.1279$ & $3.52$ & $59.7$ & $0.8^{\phantom{***}}$ & $(0.31)$ & -- &  \\
Lasso & HAR & $0.1269$ & $1.97$ & $77.4$ &  &  & -- &  \\
 & Full & $0.1310$ & $3.57$ & $59.1$ & $3.3^{\phantom{***}}$ & $(1.04)$ & -- &  \\
Elastic net & HAR & $0.1269$ & $1.97$ & $77.4$ &  &  & -- &  \\
 & Full & $0.1262$ & $3.94$ & $54.8$ & $-0.5^{\phantom{***}}$ & $(-0.21)$ & -- &  \\
LightGBM & HAR & $0.1235$ & $1.77$ & $79.6$ &  &  & -- &  \\
 & Full & $0.1201$ & $2.34$ & $73.1$ & $-2.7^{*\phantom{**}}$ & $(-1.81)$ & -- &  \\
XGBoost & HAR & $0.1234$ & $1.73$ & $80.1$ &  &  & -- &  \\
 & Full & $0.1192$ & $2.20$ & $74.7$ & $-3.4^{**\phantom{*}}$ & $(-2.18)$ & -- &  \\
Random forest & HAR & $0.1266$ & $1.63$ & $81.3$ &  &  & -- &  \\
 & Full & $0.1248$ & $1.99$ & $77.2$ & $-1.4^{\phantom{***}}$ & $(-0.92)$ & -- &  \\
\addlinespace[4pt]
\multicolumn{9}{@{}l}{\textit{Panel B: $\mathcal T_{\mathrm{close}}$, \rhcaStartOos{} to \rhcaEndOos{} ($\rhcaNdaysOos$ days), pooled intraday models}} \\
\addlinespace[2pt]
Ridge & HAR & $0.1137$ & $4.24$ & $77.7$ &  &  & $10.9^{***}$ & $(2.84)$ \\
 & Full & $0.1260$ & $7.99$ & $58.1$ & $10.9^{**\phantom{*}}$ & $(2.31)$ & $22.9^{***}$ & $(3.63)$ \\
Lasso & HAR & $0.1137$ & $4.24$ & $77.7$ &  &  & $10.9^{***}$ & $(2.84)$ \\
 & Full & $0.1301$ & $8.13$ & $57.3$ & $14.4^{**\phantom{*}}$ & $(2.51)$ & $26.9^{***}$ & $(3.68)$ \\
Elastic net & HAR & $0.1137$ & $4.24$ & $77.7$ &  &  & $10.9^{***}$ & $(2.84)$ \\
 & Full & $0.1234$ & $9.01$ & $52.7$ & $8.5^{*\phantom{**}}$ & $(1.91)$ & $20.3^{***}$ & $(3.36)$ \\
LightGBM & HAR & $0.1101$ & $3.68$ & $80.7$ &  &  & $7.4^{***}$ & $(2.60)$ \\
 & Full & $0.1089$ & $5.09$ & $73.3$ & $-1.1^{\phantom{***}}$ & $(-0.39)$ & $6.3^{\phantom{***}}$ & $(1.49)$ \\
XGBoost & HAR & $0.1102$ & $3.62$ & $81.0$ &  &  & $7.5^{***}$ & $(2.70)$ \\
 & Full & $0.1079$ & $4.79$ & $74.9$ & $-2.1^{\phantom{***}}$ & $(-0.76)$ & $5.2^{\phantom{***}}$ & $(1.30)$ \\
Random forest & HAR & $0.1131$ & $3.39$ & $82.2$ &  &  & $10.4^{***}$ & $(3.12)$ \\
 & Full & $0.1123$ & $4.26$ & $77.6$ & $-0.7^{\phantom{***}}$ & $(-0.26)$ & $9.6^{**\phantom{*}}$ & $(2.13)$ \\
\addlinespace[4pt]
\multicolumn{9}{@{}l}{\textit{Panel B, continued: the closing-interval HAR benchmark}} \\
\addlinespace[2pt]
OLS$^{a}$ & HAR & $0.1025$ & $4.31$ & $77.4$ &  &  & -- &  \\
\bottomrule
\end{tabular*}
\begin{tablenotes}[flushleft]
\footnotesize
\item \textit{Notes:} The 15:30 forecast of the closing half hour (the 16:00 bar) of the pooled models of Table~\ref{tab:rh_oos}, scored against the actual realized variance $\RV_t$ of the closing half hour, one forecast a day. Panel~A: the $\rhcaNdaysSpan$ evaluation days of Table~\ref{tab:rh_oos}; Panel~B: $\mathcal T_{\mathrm{close}}$, the \rhcaNdaysOos{} trading days from \rhcaStartOos{} to \rhcaEndOos; models trained on the closing half hour only are in Table~\ref{tab:rh_train_arms}. HAR: HAR benchmark; Full: full predictor set; on the HAR benchmark the pooled ridge, lasso and elastic net reproduce the OLS forecast (Section~\ref{sec:rh_exp_setup}). MSE in units of $10^{-\rhcaMseExp}$. $R^2_{\mathrm{OOS}}$ against the historical mean of the closing-half-hour realized variance from \rhoosBenchStart{} \citep{CampbellThompson2008}. Full vs HAR: $\Delta\% = 100(\bar L^{\mathrm{Full}}/\bar L^{\mathrm{HAR}} - 1)$ for the mean QLIKE of the same model; vs closing HAR: the same change against the closing-interval HAR benchmark ($^{a}$, the OLS on the HAR benchmark estimated on the closing half hour, Appendix~\ref{app:rh_closing_bench}), which has no forecast before \rhcaCloseFirst{} (--: not available on Panel~A's span); the closing-interval ridge, lasso and elastic net on the HAR benchmark are within \rhcaCiHarGapBound\% of it. $\DM$: Diebold--Mariano statistic of the daily QLIKE differential (Harvey--Leybourne--Newbold, $q_{\mathrm{NW}} = \rhcaDMLagSpan$ in Panel~A and \rhcaDMLagOos{} in Panel~B); negative: the first forecast has the lower loss. $^{*}$, $^{**}$, $^{***}$: significance at the 10\%, 5\% and 1\% levels. The penalized linear models' full-set changes differ between the panels because the two halves of Panel~A's span disagree: on \rhcaYearsPre{} the full set lowers their closing-half-hour QLIKE ($\Delta\% = \rhcaPctRidgeClosePre$, \rhcaPctLassoClosePre{} and \rhcaPctEnetClosePre{} for ridge, lasso and elastic net; $\DM = \rhcaDmRidgeClosePre$, \rhcaDmLassoClosePre{} and \rhcaDmEnetClosePre), on $\mathcal T_{\mathrm{close}}$ it raises it, and over the whole span the two offset.
\end{tablenotes}
\end{threeparttable}
\end{table}


% Parts B and C of drafts/2026-10-25_mcs.tex, condensed. The forecast sets and the rejected hypotheses are printed by
% macros (\rhmcs...Set..., ...Rej...), so the text follows a rerun on the rebuilt closing-interval HAR benchmark; the
% remaining words ("every", "all", "none", "alone") are asserted in make_rehaul_mcs.py (draft_claims) and are to be
% rechecked when NUMBERS reruns it on the rebuilt benchmark (A_results.tex flags the same, item PENDING (D9)).
\subsection{Multiple comparisons}
\label{app:rh_closing_mcs}
\begin{itemize}
\item \textbf{What is tested.} Tables~\ref{tab:rh_mcs} and~\ref{tab:rh_mcs_spa} apply the model confidence set, the
  test of superior predictive ability (SPA) and StepM of Appendix~\ref{sec:rh_methods_tests} to the
  \rhmcsNModelsWord{} distinct pooled forecasts of Table~\ref{tab:rh_oos} and, at the closing half hour on
  $\mathcal T_{\mathrm{close}}$, to the closing-interval HAR benchmark as well. Each SPA and StepM family compares
  the \rhmcsNAlt{} full-set forecasts, or strategies, with their comparisons.
\item \textbf{Every bar, QLIKE.} The 90\% model confidence set holds \rhmcsBarQlikeTmaxSetNinety{} under $T_{\max}$
  and \rhmcsBarQlikeTrangeSetNinety{} under $T_R$, at block lengths \rhmcsBlockLo, \rhmcsBlock{} and
  \rhmcsBlockHi; no forecast on the HAR benchmark is in it. The reductions in QLIKE from the full predictor set
  survive the adjustment in all \rhorAfPairs{} model classes: the SPA $p$-value is \rhmcsBarQAfSpaC, and StepM
  rejects all \rhmcsNAlt{} hypotheses, with every adjusted $p$-value \rhmcsBarQAfPMax.
\item \textbf{Every bar, squared error.} The 75\% set holds all \rhmcsBarMseTmaxNSevenFive{} forecasts, and StepM
  rejects none of the \rhmcsNAlt{} hypotheses (SPA $p$-value \rhmcsBarMAfSpaC).
\item \textbf{The closing half hour on $\mathcal T_{\mathrm{oos}}$.} The forecasts are hard to tell apart.
  \begin{itemize}
  \item Under $T_{\max}$ the 90\% set holds \rhmcsAllQlikeTmaxNNinety{} of the \rhmcsNModels{} pooled forecasts,
    under $T_R$ \rhmcsAllQlikeTrangeSetNinety; the 75\% set holds \rhmcsAllQlikeTmaxSetSevenFive. % 2026-10-26 (coordinator, D13): "alone" dropped (two forecasts now)
  \item Against the HAR benchmark of the same class, the SPA $p$-value is \rhmcsAllQAfSpaC; StepM rejects for
    \rhmcsAllQAfRej, and the two-sided StepM for \rhmcsAllQAfRejTwo.
  \end{itemize}
\item \textbf{The closing half hour on $\mathcal T_{\mathrm{close}}$}, with the closing-interval HAR benchmark.
  \begin{itemize}
  \item The 90\% set holds \rhmcsOosQlikeTmaxSetNinety{} under $T_{\max}$ and \rhmcsOosQlikeTrangeSetNinety{}
    under $T_R$; the 75\% set holds \rhmcsOosQlikeTmaxSetSevenFive.
  \item Against the HAR benchmark of the same class, no full-set forecast is more accurate after the adjustment (SPA
    $p$-value \rhmcsOosQAfSpaC; smallest StepM-adjusted $p$-value \rhmcsOosQAfPMin). The two-sided StepM confirms a
    higher loss for \rhmcsOosQAfRejTwo{} (adjusted $p$-value \rhmcsOosQAfPTwoLasso), not for ridge regression
    (\rhmcsOosQAfPTwoRidge).
  \item No full-set forecast is more accurate than the closing-interval HAR benchmark (SPA $p$-value
    \rhmcsOosQBmSpaC). Testing all \rhmcsNModelsWord{} distinct pooled forecasts against it at once, the two-sided
    StepM confirms a higher loss for \rhmcsOosQBmAllRejTwo.
  \item Under squared error every forecast is in the 75\% sets (\rhmcsOosMseTmaxNSevenFive{} of
    \rhmcsNModelsBench), and StepM rejects no hypothesis.
  \end{itemize}
\item \textbf{Strategies} (Table~\ref{tab:rh_mcs_spa}, Panels~C and~D; the \rheSzDays{} post-warmup days).
  \begin{itemize}
  \item Against the unconditional short straddle, the Sharpe ratio gain survives the adjustment for every full-set
    strategy: the SPA $p$-value is \rhmcsMidShortSrSpaC{} at mid-quotes and \rhmcsBaShortSrSpaC{} at quoted bid and
    ask prices, and StepM rejects all \rhmcsNAlt{} hypotheses (two-sided adjusted $p$-values at most
    \rhmcsMidShortSrPTwoMax{} and \rhmcsBaShortSrPTwoMax). For the mean daily return differential the SPA $p$-value
    is \rhmcsMidShortMeanSpaC, and StepM rejects \rhmcsMidShortMeanNRej{} of the \rhmcsNAlt{} hypotheses.
  \item Against the HAR-benchmark strategy of the same class, the Sharpe ratio gain survives for the random forest
    alone: SPA $p$-value \rhmcsMidTwinSrSpaC, StepM-adjusted $p$-value \rhmcsMidTwinSrPRf{} (two-sided
    \rhmcsMidTwinSrPTwoRf).
  \end{itemize}
\end{itemize}
% AUTO-GENERATED by writeup/make_rehaul_mcs.py -- do not edit; re-run the script.
% tab:rh_mcs: MCS p-values (T_max and T_R) of the pooled intraday forecasts and the closing-interval HAR
% benchmark; circular block bootstrap, daily losses (every-bar losses averaged to the day).
\begin{table}[htbp]
\centering
\small
\setlength{\tabcolsep}{4pt}
\begin{threeparttable}
\caption{Model confidence sets of the pooled intraday forecasts}
\label{tab:rh_mcs}
\begin{tabular*}{\linewidth}{@{\extracolsep{\fill}}lrrrrrr@{}}
\toprule
 & \multicolumn{2}{c}{Every bar, 2014--2024} & \multicolumn{2}{c}{16:00 bar, 2014--2024} & \multicolumn{2}{c}{16:00 bar, $\mathcal T_{\mathrm{oos}}$} \\
\cmidrule(lr){2-3}\cmidrule(lr){4-5}\cmidrule(lr){6-7}
Forecast & $T_{\max}$ & $T_R$ & $T_{\max}$ & $T_R$ & $T_{\max}$ & $T_R$ \\
\midrule
\multicolumn{7}{@{}l}{\textit{Panel A: QLIKE}} \\
\addlinespace[2pt]
Linear, HAR$^{a}$ & \ensuremath{<}0.001 & \ensuremath{<}0.001 & \textit{0.126} & \ensuremath{<}0.001 & 0.052 & 0.061 \\
Ridge, Full & \ensuremath{<}0.001 & \ensuremath{<}0.001 & \textit{0.159} & 0.015 & 0.011 & 0.008 \\
Lasso, Full & \ensuremath{<}0.001 & \ensuremath{<}0.001 & \textit{0.126} & 0.011 & 0.011 & 0.011 \\
Elastic net, Full & \ensuremath{<}0.001 & \ensuremath{<}0.001 & \textit{0.246} & 0.030 & 0.036 & 0.016 \\
LightGBM, HAR & \ensuremath{<}0.001 & \ensuremath{<}0.001 & \textit{0.246} & \textit{0.103} & \textit{0.246} & \textit{0.116} \\
LightGBM, Full & \textbf{1.000} & \textbf{1.000} & \textit{0.246} & \textit{0.159} & \textit{0.246} & \textit{0.138} \\
XGBoost, HAR & \ensuremath{<}0.001 & \ensuremath{<}0.001 & \textit{0.246} & \textit{0.103} & \textit{0.246} & \textit{0.116} \\
XGBoost, Full & 0.003 & 0.003 & \textbf{1.000} & \textbf{1.000} & \textbf{0.256} & \textbf{0.256} \\
Random forest, HAR & \ensuremath{<}0.001 & \ensuremath{<}0.001 & \textit{0.159} & \ensuremath{<}0.001 & 0.052 & 0.052 \\
Random forest, Full & \ensuremath{<}0.001 & \ensuremath{<}0.001 & \textit{0.182} & \ensuremath{<}0.001 & \textit{0.129} & 0.016 \\
Closing-interval HAR benchmark$^{b}$ & -- & -- & -- & -- & \textbf{1.000} & \textbf{1.000} \\
\addlinespace[4pt]
\multicolumn{7}{@{}l}{\textit{Panel B: Squared error}} \\
\addlinespace[2pt]
Linear, HAR$^{a}$ & \textbf{1.000} & \textbf{1.000} & \textbf{0.572} & \textbf{0.653} & \textbf{0.537} & \textbf{0.613} \\
Ridge, Full & \textbf{0.489} & \textbf{0.415} & \textbf{0.437} & \textbf{0.653} & \textbf{0.446} & \textbf{0.613} \\
Lasso, Full & \textbf{0.334} & \textbf{0.415} & \textbf{0.371} & \textbf{0.653} & \textbf{0.383} & \textbf{0.613} \\
Elastic net, Full & \textbf{0.422} & \textbf{0.415} & \textbf{0.488} & \textbf{0.653} & \textbf{0.496} & \textbf{0.613} \\
LightGBM, HAR & \textbf{0.489} & \textbf{0.310} & \textbf{0.572} & \textbf{0.369} & \textbf{0.537} & \textbf{0.561} \\
LightGBM, Full & \textbf{0.489} & \textbf{0.415} & \textbf{0.488} & \textbf{0.653} & \textbf{0.537} & \textbf{0.613} \\
XGBoost, HAR & \textbf{0.489} & \textbf{0.415} & \textbf{0.572} & \textbf{0.653} & \textbf{0.537} & \textbf{0.613} \\
XGBoost, Full & \textbf{0.489} & \textbf{0.415} & \textbf{0.488} & \textbf{0.653} & \textbf{0.537} & \textbf{0.613} \\
Random forest, HAR & \textbf{0.489} & \textbf{0.310} & \textbf{1.000} & \textbf{1.000} & \textbf{1.000} & \textbf{1.000} \\
Random forest, Full & \textbf{0.489} & \textbf{0.415} & \textbf{0.558} & \textbf{0.653} & \textbf{0.537} & \textbf{0.613} \\
Closing-interval HAR benchmark$^{b}$ & -- & -- & -- & -- & \textbf{0.537} & \textbf{0.572} \\
\bottomrule
\end{tabular*}
\begin{tablenotes}[flushleft]
\footnotesize
\item \textit{Notes:} MCS $p$-values \citep{HansenLundeNason2011} of the pooled intraday forecasts of Table~\ref{tab:rh_oos} and, on $\mathcal T_{\mathrm{oos}}$, of the closing-interval HAR benchmark. A forecast is in the 90\% (75\%) model confidence set when its $p$-value is at least 0.10 (0.25); \textbf{bold}: in the 75\% set, \textit{italic}: in the 90\% set only. $T_{\max}$ and $T_R$: the two statistics of the elimination rule. Losses are daily (on every bar, the mean loss of the day's \rhorNBars{} bars), so the bootstrap resamples trading days: the \rhmcsDaysAll{} days from \rhorSpanStart{} to \rhorSpanEnd{} in the first four columns and the \rhmcsDaysOos{} days of $\mathcal T_{\mathrm{oos}}$ in the last two. Circular moving-block bootstrap, blocks of \rhmcsBlock{} trading days, \rhmcsB{} replications, seed \rhmcsSeed; the same resamples at every elimination step. $^{a}$The pooled ridge, lasso and elastic net on the HAR benchmark reproduce the pooled OLS (their penalties are numerically zero), so they enter as one model, the OLS forecast. $^{b}$The closing-interval HAR benchmark has forecasts from \rhmcsChStart{} only, so it is not in the 2014--2024 sets.
\end{tablenotes}
\end{threeparttable}
\end{table}

% tab:rh_mcs_spa: SPA p-values and StepM-adjusted p-values (single-test bootstrap p in parentheses).
\begin{table}[htbp]
\centering
\small
\setlength{\tabcolsep}{3pt}
\begin{threeparttable}
\caption{Superior predictive ability and StepM tests of the full predictor set}
\label{tab:rh_mcs_spa}
\begin{tabular*}{\linewidth}{@{\extracolsep{\fill}}lrrrrrrr@{}}
\toprule
 & SPA & \multicolumn{6}{c}{StepM, by model class} \\
\cmidrule(lr){3-8}
Comparison & $p$ & Ridge & Lasso & Elastic net & LightGBM & XGBoost & Random forest \\
\midrule
\multicolumn{8}{@{}l}{\textit{Panel A: QLIKE, full predictor set against the HAR benchmark of the same class}} \\
\addlinespace[2pt]
Every bar, 2014--2024 & \ensuremath{<}0.001 & \textbf{\ensuremath{<}0.001} & \textbf{\ensuremath{<}0.001} & \textbf{\ensuremath{<}0.001} & \textbf{\ensuremath{<}0.001} & \textbf{\ensuremath{<}0.001} & \textbf{\ensuremath{<}0.001} \\
 &  & (\ensuremath{<}0.001) & (\ensuremath{<}0.001) & (\ensuremath{<}0.001) & (\ensuremath{<}0.001) & (\ensuremath{<}0.001) & (\ensuremath{<}0.001) \\
16:00 bar, 2014--2024 & 0.033 & 0.675 & 0.842 & 0.499 & 0.072 & \textbf{0.033} & 0.302 \\
 &  & (0.646) & (0.842) & (0.437) & (0.034) & (0.013) & (0.177) \\
16:00 bar, $\mathcal T_{\mathrm{oos}}$ & 0.380 & 0.987 & 0.987 & 0.973 & 0.581 & 0.415 & 0.599 \\
 &  & (0.984) & (0.986) & (0.961) & (0.345) & (0.219) & (0.402) \\
\addlinespace[4pt]
\multicolumn{8}{@{}l}{\textit{Panel B: QLIKE, full predictor set against the closing-interval HAR benchmark}} \\
\addlinespace[2pt]
16:00 bar, $\mathcal T_{\mathrm{oos}}$ & 0.529 & 1.000 & 1.000 & 0.999 & 0.972 & 0.959 & 0.995 \\
 &  & (0.999) & (0.999) & (0.999) & (0.905) & (0.867) & (0.974) \\
\addlinespace[4pt]
\multicolumn{8}{@{}l}{\textit{Panel C: sign strategies against the unconditional short straddle}} \\
\addlinespace[2pt]
Mean return, mid-quotes & 0.024 & \textbf{0.035} & \textbf{0.035} & \textbf{0.030} & \textbf{0.043} & 0.050 & \textbf{0.024} \\
 &  & (0.023) & (0.025) & (0.020) & (0.038) & (0.050) & (0.014) \\
Mean return, bid/ask & 0.024 & \textbf{0.034} & \textbf{0.034} & \textbf{0.029} & \textbf{0.042} & 0.050 & \textbf{0.024} \\
 &  & (0.023) & (0.024) & (0.019) & (0.038) & (0.050) & (0.014) \\
Sharpe ratio, mid-quotes & 0.002 & \textbf{0.002} & \textbf{0.003} & \textbf{0.002} & \textbf{0.002} & \textbf{0.003} & \textbf{0.002} \\
 &  & (\ensuremath{<}0.001) & (0.003) & (0.002) & (\ensuremath{<}0.001) & (0.002) & (0.001) \\
Sharpe ratio, bid/ask & 0.002 & \textbf{0.003} & \textbf{0.004} & \textbf{0.003} & \textbf{0.002} & \textbf{0.004} & \textbf{0.002} \\
 &  & (\ensuremath{<}0.001) & (0.003) & (0.002) & (0.001) & (0.003) & (0.001) \\
\addlinespace[4pt]
\multicolumn{8}{@{}l}{\textit{Panel D: sign strategies against the HAR-benchmark strategy of the same class}} \\
\addlinespace[2pt]
Mean return, mid-quotes & 0.026 & 0.087 & 0.093 & 0.072 & \textbf{0.048} & 0.138 & \textbf{0.026} \\
 &  & (0.049) & (0.059) & (0.037) & (0.021) & (0.138) & (0.012) \\
Mean return, bid/ask & 0.026 & 0.087 & 0.093 & 0.073 & \textbf{0.048} & 0.138 & \textbf{0.026} \\
 &  & (0.049) & (0.059) & (0.037) & (0.021) & (0.138) & (0.012) \\
Sharpe ratio, mid-quotes & 0.022 & 0.175 & 0.175 & 0.175 & 0.096 & 0.175 & \textbf{0.022} \\
 &  & (0.118) & (0.150) & (0.101) & (0.047) & (0.100) & (0.003) \\
Sharpe ratio, bid/ask & 0.022 & 0.169 & 0.169 & 0.169 & 0.088 & 0.169 & \textbf{0.022} \\
 &  & (0.110) & (0.145) & (0.097) & (0.044) & (0.097) & (0.003) \\
\bottomrule
\end{tabular*}
\begin{tablenotes}[flushleft]
\footnotesize
\item \textit{Notes:} Each row is one family of \rhmcsNAlt{} one-sided hypotheses $H_k$: the full-set forecast (strategy) of class $k$ is no better than its comparison. SPA: the consistent $p$-value of the test of superior predictive ability of \citet{Hansen2005}, $H_0$: no full-set forecast in the row beats its comparison. StepM: $p$-values adjusted for the \rhmcsNAlt{} comparisons of the row by the studentized step-down procedure of \citet{RomanoWolf2005}, which controls the familywise error rate; \textbf{bold}: below 0.05. In parentheses: the one-sided bootstrap $p$-value of the comparison alone, without the adjustment. Panels A and B: daily QLIKE loss differentials (every bar: the mean over the day's \rhorNBars{} bars), circular moving-block bootstrap with blocks of \rhmcsBlock{} trading days, \rhmcsB{} replications, seed \rhmcsSeed. Panels C and D: daily returns of the sign strategy (Section~\ref{sec:rh_econ_strategy}) on the \rhmcsNPw{} post-warmup days, the mean daily return differential and the difference of the Sharpe ratios of Section~\ref{sec:rh_econ_measures}, with the paired circular block bootstrap of Appendix~\ref{sec:rh_methods_tests} (blocks of \rhmcsBlock{} days, \rhmcsB{} replications, seed \rhmcsSeedEcon), each statistic studentized by its bootstrap standard deviation. The single-test $p$-values in parentheses are those of this bootstrap, one-sided; they are not the two-sided Diebold--Mariano and fixed-$b$ $p$-values of Sections~\ref{sec:rh_results} and~\ref{sec:rh_econ}.
\end{tablenotes}
\end{threeparttable}
\end{table}

